\documentclass[12pt]{report}

% ---------------------------------------------------------------
% 字型與中文排版
% ---------------------------------------------------------------
\usepackage{fontspec}
\usepackage{xeCJK}
\setmainfont{Times New Roman}
\setCJKmainfont{Songti TC}
\setCJKsansfont{PingFang TC}
% 等寬字改用 Latin Modern Mono，跟 exp2 參考檔案的排版字體一致
% （它不是系統字型，只存在 TeXLive 樹裡，所以要用 Path 明確指定檔案位置）。
\setmonofont{lmmono10-regular}[
    Path = /usr/local/texlive/2024/texmf-dist/fonts/opentype/public/lm/,
    Extension = .otf,
    AutoFakeBold = 2.5,
]

% ---------------------------------------------------------------
% 本書排版套件（packages/ 目錄）
% ---------------------------------------------------------------
\usepackage{sleek}
\usepackage{sleek-theorems}
\usepackage{sleek-code}
\usepackage{sleek-diagrams}
% 注意：sleek-listings.sty 不載入——它是舊版 lstlistings/inconsolata 排版套件，
% 本書實際用的是 sleek-code.sty 的 minted 版 code/example 環境，兩者同時載入會讓
% minted 抓不到 \setmonofont 設定的等寬字（回退成 Latin Modern Roman），
% 且會讓程式碼裡的直引號被轉成全形彎引號。

% 章節標題中文化：「第 N 章」
\usepackage{titlesec}
\titleformat{\chapter}{\normalfont\huge\bfseries}{第~\thechapter~章}{1em}{}
\titlespacing*{\chapter}{0pt}{-20pt}{20pt}

% 目錄的節編號欄位加寬：report 預設是 2.3em，放不下「10.10」這種兩位數章＋兩位數節的
% 編號，會跟標題黏在一起（例如「10.10finally」）。放寬到 3em。
\makeatletter
\renewcommand{\l@section}{\@dottedtocline{1}{1.5em}{3em}}
\makeatother

% 書末索引：各章在術語第一次正式定義處加 \index{}，書末 \printindex 印出頁碼對照表。
\usepackage{makeidx}
\makeindex
\renewcommand{\indexname}{索引}

\title{Python 入門：從寫程式到看懂軟體世界}
\author{}
\date{2026}

\begin{document}

\maketitle
\romantableofcontents

\chapter*{如何使用這本書}
\addcontentsline{toc}{chapter}{如何使用這本書}

如果你從來沒寫過程式，這本書就是為你寫的。不管你現在是國中生、高中生、大學生，還是已經
在職場工作好幾年的社會人士，這本書都預設你完全沒有寫過程式，也沒有用過任何特定的辦公
軟體或開發工具。書裡需要打比方的地方，會用大家都有的生活經驗（食譜步驟、手機 App 怎麼
運作、日常生活裡的規則跟流程），不會假設你用過某個特定軟體。

這本書有兩個並重的目標，不是只有一個。第一個目標，是讓你能自己動手寫小程式，解決生活、
課業或工作上的瑣事，例如整理一堆檔案、自動化重複性的工作、做簡單的資料分析。終點不是
「精通 Python 語言的每一個角落」，而是看到一件瑣碎的任務時，會反射性想到「這個我可以
寫程式做」，而且真的做得出來。

第二個目標，是讓你理解「軟體」這個東西究竟是怎麼被打造出來的：版本控制在做什麼、一個
專案的檔案通常長什麼樣子、前端、後端、API、部署這些詞在講什麼、工程師團隊的日常運作
方式大概是什麼樣子。這個目標特別是為了讓沒有工程背景、但工作上需要跟工程師對接的人
（最典型的是 PM），有能力聽懂、問對問題，不會被術語唬住。第 \ref{ch:industry}、\ref{ch:api} 章是專門服務這個
目標的章節。

這本書不會把你訓練成職業工程師，也不是那種「系統性教完整個語言」或「給有經驗的工程師
精進」的書。如果你剛好正卡在某堂必修的程式設計課，這本書當然也能幫上忙，但這不是這本書
設定的主要對象，你不會在書裡看到「這在 C 語言裡怎麼寫」這種對照內容。

\section*{每一節大致的節奏}

大部分小節都會照著同一個節奏走：先用一段白話文字說明這個技巧是什麼、為什麼需要它，
接著給你一個通用的片語樣板，然後是一支真正能執行的完整範例程式，範例後面緊接著這支
程式實際的輸入輸出。有些地方會補上一段特別提醒（標成「學習要點」或「風格要點」），
挑的都是容易跟你原本的直覺互相衝突、或是比較道地的 Python 寫法。

很多小節還會示範一次新手常見的錯誤寫法，讓你看看那個錯誤訊息長什麼樣子、該怎麼讀懂它、
怎麼修正。打字打錯一個字元，程式就整個不動，這件事任何人第一次寫程式都會碰到，跟你是
國中生還是已經工作多年沒有關係。看到紅色的錯誤訊息不代表你不是這塊料，只代表 Python
想告訴你哪裡需要修改，讀懂它、修好它，你就往前進了一步。

大部分小節最後會有兩到四題簡短的練習，讓你動手確認自己真的懂了；每一章的最後則集中
放著全章練習的解答，方便你對答案。

\section*{書裡會用到的幾種方塊}

程式碼樣板放在深色標題的方塊裡，標題是這個樣板的用途；真正能執行的完整範例也用類似的
方塊呈現，標題會附上檔名。範例後面常常接著一個標著「input」和「output」的小方塊，那是
這支程式實際的輸入跟輸出內容。「學習要點」跟「風格要點」是灰底的提示段落，通常一句話
就講完一個重點，值得多看一眼。

\section*{兩個會反覆出現的思考方式}

這本書會用兩個貫穿全書的思考方式，幫你建立對 Python 的直覺，而不是要你死背一堆規則。

第一個是「名字牌模型」：把變數想成一張貼在物件上的名字牌，而不是一個裝著值的盒子。這個
想法第 \ref{ch:variables} 章會第一次完整說明，之後在講清單、函式、tuple 的章節都會重新呼應一次，因為
每次呼應的情境不太一樣，能看到的東西也不一樣。

第二個是「先做、錯了再處理」（EAFP），跟「先檢查、確定安全才動手」（LBYL）相對照。
這個想法第 \ref{ch:files-exceptions} 章講檔案與例外處理時會第一次完整說明，之後在講字典、除錯的章節也會
重新呼應。

第一次看到這兩個名字，不用急著完全弄懂，書裡每次呼應時都會簡短提醒一次背景，多看幾次
自然就有感覺了。

\section*{建議的讀法}

這本書大致依照由簡入繁的順序安排，前面幾章（尤其是第 \ref{ch:variables} 到 \ref{ch:files-exceptions} 章）是後面所有內容的
地基，建議照順序讀完，不要跳著看。第 \ref{ch:csv} 章之後的章節彼此之間的關聯性比較低，如果你
已經有明確想解決的問題（例如想學怎麼串接 API），可以考慮直接跳到對應的章節。

不用急著一次讀完一整章，每個小節後面的練習題值得你真的動手做一次，而不是只看過去就好。
程式是動手做出來的技能，不是背出來的知識。

祝你讀得順利，也祝你很快就能親手寫出第一支真正解決問題的程式。


% ---------------------------------------------------------------
% 各章節內容：每個單元的敘述式教材放在對應的
% "chapters/<單元>/chapter.tex"，由 pipeline/chapters/<單元>/draft.tex
% 促升而來（見 pipeline/book-bible.md 與 .claude/skills/python-book-author/）。
% 每完成一章，在這裡加一行 \input，並依 pipeline/chapter-map.yaml
% 的章節順序排列。
% ---------------------------------------------------------------

\chapter{為什麼學 Python、怎麼開始}
\label{ch:introduction}

這一章的目標，是在闔上這一章之前，讓電腦真的照你的話做一件事，哪怕只是印出一行字，至於
把 Python 讀熟，是後面每一章慢慢累積的事。市面上不少入門書一開場會先講一大段「為什麼要
學程式」「為什麼是 Python」，聽起來很有道理，但對一個還沒摸過鍵盤敲出程式碼的人來說，
這些道理其實有點空洞：你連自己在等待什麼都還沒有具體的畫面，講太多動機，反而容易讓人在
真正動手之前就先打退堂鼓。

所以這一章刻意把順序反過來排：先把工具準備好，寫出你人生第一
支程式，建立一個「我真的做到了」的具體經驗當靠山；接著回頭拆解剛剛發生的事，包括你剛才
到底是用哪一種方式在跟 Python 講話、如果打錯字會發生什麼事、那些看起來很嚇人的紅字錯誤
訊息其實在說什麼；最後，帶著這些具體經驗，才回頭談比較大的問題：為什麼值得花時間學程式、
為什麼是 Python、這本書接下來想帶你去的兩個地方。這個順序背後的想法很簡單：例子要先於
定義，成就感要先於道理。

不管你現在是還在準備考試的學生，還是已經在職場上開會時常聽不懂
工程師在講什麼的社會人士，讀完這一章，你都會有一支自己跑得動的程式，也會知道接下來的路
要往哪裡走。

在往下讀之前，先講清楚一件事：程式，就是你寫給電腦的一整串指令，電腦會照著這些指令，
一步一步、原封不動地做。\index{程式}接下來這一整章，你會不斷看到「程式」「程式碼」「指令」這幾個
字，指的都是這件事。

\section{準備工具：安裝 Python 與挑一個編輯器}
\label{sec:setup}
\index{直譯器}
\index{編輯器}

要讓電腦執行 Python 程式，你其實需要兩樣分開的東西：一個是真正看得懂、執行得了 Python
程式碼的軟體，叫做「直譯器」；另一個是你打字、整理程式碼的地方，也就是編輯器。這兩個是
不同的軟體、做不同的事，很多初學者會搞混，以為裝一個就等於裝了另一個。這一節結束時，你
的電腦上會同時有這兩樣東西，接下來整本書、每一次動手寫程式，都會重複用到它們。

如果你覺得這個區分有點抽象，可以想像成準備一份食譜跟找一個真的會做菜的廚師：食譜（你打
的程式碼）只是寫在紙上的文字，紙本身不會自己把菜端出來；真正把食譜上的步驟一一做出來
的，是廚師（直譯器）。編輯器負責讓你把「食譜」寫得整齊、方便修改；直譯器負責照著食譜真
的動手做。

裝了很棒的編輯器，等於有一疊漂亮的空白食譜紙，但沒有廚師，菜還是端不出來；
反過來，只有直譯器沒有編輯器，你也不是完全不能寫程式（甚至可以用最陽春的記事本打字），
只是會很不方便。這本書接下來會請你兩個都裝好。

\textbf{第一步：安裝 Python 直譯器。} 到官方網站 \texttt{python.org/downloads}，網站通常
會自動判斷你的作業系統，直接按下載鈕即可。

\begin{itemize}
\item \textbf{Windows}：下載完成後執行安裝程式，在跳出來的第一個畫面\textbf{務必勾選
「Add python.exe to PATH」}這個選項，再按「Install Now」。這裡的 PATH 是電腦「去哪裡找
程式」的清單，先不用深究，記得勾就好。這一步決定了之後你在終端機（Windows 上這個視窗叫
「命令提示字元」，下面會再說明）打 \texttt{python}，電腦找不找得到它；勾了它，後面的路
會順很多，沒勾的話之後每次打指令都會出現「找不到指令」的錯誤，下一段會示範這個錯誤實際
長什麼樣子。
\item \textbf{macOS}：下載副檔名（檔名最後面 . 後面那幾個字母，用來標記檔案種類）是
\texttt{.pkg} 的安裝檔，直接雙擊執行，照著安裝精靈的預設值一路按下去即可，不需要額外
勾選任何 PATH 選項。
\end{itemize}

\index{終端機}
裝完之後，打開終端機（Windows 叫「命令提示字元」或 PowerShell，macOS 叫「終端機」），
輸入以下指令確認安裝成功：Windows 打 \texttt{python --version}，macOS 打
\texttt{python3 --version}。你應該會看到類似 \texttt{Python 3.13.2} 這樣的一行字（版本
號不用跟這裡完全一樣，只要有印出版本號，就代表裝成功了）。

\textbf{如果打完指令，畫面卻說「找不到這個指令」}：這是 Windows 使用者最容易踩到的一個
小陷阱。如果剛剛安裝時不小心漏勾了「Add python.exe to PATH」，打開終端機輸入
\texttt{python --version} 之後，你會看到類似這樣的訊息：

\begin{verbatim}
'python' is not recognized as an internal or external command,
operable program or batch file.
\end{verbatim}

這句話的意思是：終端機在它認得的那份「清單」（也就是剛才提到的 PATH）裡，找不到一個叫
\texttt{python} 的程式，所以它老實告訴你「不認識這個指令」，不代表你的電腦壞了，也不代表
安裝失敗，只是那份清單裡少了一條路徑。

修法很直接：把剛剛下載的安裝程式重新執行一次，這次
選「Uninstall」移除，再重新安裝一次，記得在第一個畫面把「Add python.exe to PATH」勾起
來。macOS 使用者不會遇到這個問題，因為 macOS 版的安裝過程本來就不需要這一步。

\textbf{第二步：安裝 VS Code，挑一個地方打字。} 到 \texttt{code.visualstudio.com} 下載
安裝，跟安裝任何軟體一樣，一路「下一步」到底即可。裝好之後：

\begin{itemize}
\item 用選單「File $\rightarrow$ Open Folder」開一個資料夾。建議先新建一個空的資料夾，
例如取名 \texttt{python-practice}，之後這本書每一個範例都可以存在這裡。
\item 用選單「View $\rightarrow$ Terminal」（或按下 \texttt{Ctrl+`}）打開一個內建終端機。它其實就是
你電腦本來就有的那個終端機，只是被塞進了 VS Code 的視窗裡，你不用另外切換視窗。
\item 新增一個檔案，存成 \texttt{hello.py}。副檔名 \texttt{.py} 是在告訴電腦、也告訴你
自己，這是一個 Python 程式檔案，下一節就會用到它。
\end{itemize}

先說清楚：這裡刻意只用 VS Code 最陽春的功能，也就是開資料夾、打字、存檔、用終端機執行。
你\textbf{不需要}安裝任何 Python 相關的擴充套件（讓編輯器多長出額外功能的外掛），也不需要
碰任何進階的除錯設定，這些等真的用得上的時候，書裡會再告訴你。

接下來每一章都會重複同一個小動作：在編輯器裡打程式碼、
存檔、切到終端機打指令執行、看結果。現在你已經有了完成這個循環需要的兩樣東西，工具都到
位了，翻到下一節，把你人生的第一支程式真的敲出來。

\begin{drills}[label=drills:setup]
\item 如果你在 Windows 上打開終端機、輸入 \texttt{python --version}，畫面顯示
\texttt{'python' is not recognized as an internal or external command}，你會回頭懷疑
剛才安裝時哪一步沒做？
\item 直譯器跟編輯器，哪一個負責「真正執行」你寫的程式碼？哪一個只負責讓你打字、整理
程式碼？
\end{drills}

\section{寫下你的第一支程式}
\label{sec:hello-world}
\index{print}

電腦本身很安靜，你沒叫它做什麼，它就什麼都不會做。要讓螢幕上出現一行文字，你得用一個
明確的動作告訴 Python：「把這段文字顯示出來」。這個動作在 Python 裡叫做 \texttt{print}，
你可以先把它想成一個「印出來」的指令，把你想顯示的內容放進括號裡，Python 就會把它顯示
在螢幕上。

\begin{code}{python}{印出文字}
print("your text here")
\end{code}

\texttt{print} 後面一定要接一對圓括號，你想印出的內容放在括號裡；如果內容是文字（不是
數字），前後要用引號包起來：雙引號 \texttt{"..."} 或單引號 \texttt{'...'} 都可以，
Python 不在乎你選哪一種，但開頭跟結尾要用同一種，不能一邊雙引號一邊單引號。

\begin{example}{python}{(hello.py) 跟世界打招呼}
print("Hello, world!")
\end{example}

\iosample{}{Hello, world! \\}

把這兩行存成 \texttt{hello.py}（沿用上一節新增的檔案即可），在終端機切到那個資料夾，打
\texttt{python3 hello.py}（macOS）或 \texttt{python hello.py}（Windows），按下 Enter，
螢幕上就會出現上面那行字。

\texttt{print} 能顯示的東西，不只是用引號包起來的文字。數字可以直接放進括號裡，不需要
加任何引號：

\begin{example}{python}{(number-print.py) 印出數字}
print(2024)
\end{example}

\iosample{}{2024 \\}

這裡刻意不加引號：數字本身就是數字，加了引號反而會讓 Python 把它當成一段文字來處理，
效果跟數字不一樣（這件事後面章節會講得更清楚）。目前你只需要記得一個簡單的判斷原則：你
想印出的東西如果是一段文字，前後要包引號；如果是一個數字，直接寫就好，不用包任何東西。

還記得前面提過，程式就是你寫給電腦的一整串指令，電腦會一步一步、原封不動地照做嗎？
\texttt{print} 是最適合拿來體會這件事的動作：寫幾行 \texttt{print}，Python 就會依照你
寫的順序，一行一行執行過去。

\begin{example}{python}{(two-lines.py) 連續印出兩行}
print("Hello")
print("Nice to meet you")
\end{example}

\iosample{}{Hello \\ Nice to meet you \\}

兩個 \texttt{print} 各自獨立，Python 執行完第一行、把「Hello」顯示出來之後，才會接著往
下執行第二行；順序完全按照你寫在檔案裡的先後，不會自己調換，也不會兩行擠在同一行顯示
（除非你自己刻意這樣安排，這件事之後章節會再細講）。

\keypoint{人類讀者常常能自動腦補打錯的字、猜出你想講什麼，但 Python 不會，它只會照著
你打的字元一個不漏地檢查，只要有一個地方不符合它認得的規則，就會停下來告訴你出錯了，不
會幫你腦補。這種「龜毛」不是故意要為難你，說穿了是這門語言（其實幾乎所有程式語言都一樣）
運作的方式而已。}

既然 Python 這麼龜毛，那如果你不小心漏打了一個引號會怎樣？

\begin{example}{python}{(hello-broken.py) 漏了一個引號}
print("Hello, world!)
\end{example}

存檔執行後，螢幕上不會印出任何文字，只會看到這樣的訊息：

\begin{verbatim}
  File "hello-broken.py", line 1
    print("Hello, world!)
          ^
SyntaxError: unterminated string literal (detected at line 1)
\end{verbatim}

這段話在說什麼？拆開來看：
\index{SyntaxError}

\begin{itemize}
\item \texttt{SyntaxError} 代表 Python 在還沒開始執行之前，光是「讀」你寫的
程式碼這一步就卡住了，這種「讀的規則」，我們就叫它語法。
\item \texttt{unterminated string
literal} 是說「你開了一個字串（用引號開頭），但沒看到結尾的引號」。
\item 那個 \texttt{\^{}}
符號指的是 Python 覺得問題大概出在哪個位置附近。
\end{itemize}

修法很直接：把漏掉的引號補回去就好。你
在打字時漏東漏西，是每個人（包括寫了十年程式的人）都會發生的事，說到底只是 Python 很
老實地把它看不懂的地方告訴你而已，不代表你特別笨手笨腳。

漏掉引號是最常見的手誤，但還有一種更容易讓人一頭霧水的變形：頭尾兩個引號的\textbf{種類}
不一樣。

\begin{example}{python}{(quote-mismatch.py) 頭尾引號種類不一致}
print("Hello, world!')
\end{example}

存檔執行後，你會看到跟前面一模一樣的錯誤：

\begin{verbatim}
  File "quote-mismatch.py", line 1
    print("Hello, world!')
          ^
SyntaxError: unterminated string literal (detected at line 1)
\end{verbatim}

乍看之下這行程式明明頭尾都有引號，怎麼還是「沒有結尾」？原因在於 Python 認的是「同一種
引號再出現一次」才算結尾：這一行用雙引號開頭，Python 就會一路往後找，直到看到下一個雙
引號才算數；中間那個單引號，在雙引號眼裡只是一個普通字元，不算數。

所以對 Python 來說，
這一行從頭到尾都沒有真正「關起來」，錯誤訊息才會跟漏打引號時一模一樣。修法一樣直接：
確認頭尾用的是同一種引號。

\begin{drills}[label=drills:hello-world]
\item 不要真的執行，用眼睛看：這段程式執行後，螢幕上會印出什麼？
\begin{verbatim}
print("My favorite programming language is Python")
\end{verbatim}
\item 把 \texttt{hello.py} 改成印出你自己的名字，例如「哈囉，我是〇〇」，存檔後在終端機
重新執行一次，確認畫面真的印出你改過的內容。
\item 如果有人把 \texttt{print("Hello")} 打成 \texttt{print "Hello"}（漏了括號），你覺得
Python 會怎麼抱怨？這是實際的錯誤訊息，讀讀看在說什麼：
\begin{verbatim}
  File "oops.py", line 1
    print "Hello"
    ^^^^^^^^^^^^^
SyntaxError: Missing parentheses in call to 'print'. Did you mean print(...)?
\end{verbatim}
\item 如果你打 \texttt{print('Hello, world!")}（開頭用單引號，結尾用雙引號），你覺得執行
後會出現哪一種錯誤？錯誤訊息會跟本節哪一個例子長得很像？
\end{drills}

\section{編輯器裡的隱藏陷阱：直引號與彎引號}
\label{sec:editor-pitfalls}

你現在已經真的用 VS Code 寫出、也讀懂過一次錯誤訊息了，這裡回頭聊聊「編輯器」這件事，
順便帶你認識一個很多人都踩過、卻常常想不通原因的陷阱。

一個對寫程式來說「好用」的編輯器，通常具備幾個共同的特徵：

\begin{itemize}
\item 它是純文字軟體，存出來的檔案
裡除了你打的字元，不會偷偷夾帶任何額外的排版資訊。
\item 它用等寬字型顯示（每個字元佔的寬度都
一樣），這樣一整排的縮排、對齊才看得整齊。
\item 它不會自作主張「修飾」你打的內容。
\end{itemize}

VS Code 符
合這幾個條件，但它不是唯一的選擇：Sublime Text、PyCharm，甚至終端機裡的 \texttt{nano}、
\texttt{vim}，都是純文字編輯器，效果跟 VS Code 完全一樣，只是操作介面不同。

這本書接下來
的畫面跟指令會統一以 VS Code 示範，方便截圖跟教學一致，但你完全可以換成自己習慣的其他
編輯器，不會有任何差別，因為編輯器只負責讓你打字，真正執行程式碼的一直都是直譯器，它不
在乎那個 \texttt{.py} 檔案是用哪一套編輯器存出來的。

但有一類軟體不適合拿來寫程式，而且很容易在不知不覺間讓你的程式碼壞掉：Word、Pages、
Notion 這類「文書處理」或「筆記」軟體。它們的設計目標是排出漂亮的文件，其中一個常見的
「貼心」功能，是自動把你打的直引號換成造型更好看、頭尾方向不同的彎引號。這在寫一篇報告
時完全沒問題，甚至更好看。

但如果你在這類軟體裡打了一段 Python 程式碼、再貼進
\texttt{.py} 檔案，麻煩就來了：

\begin{example}{python}{(curly-quote.py) 彎引號惹的禍}
print(“Hello, world!”)
\end{example}

存檔執行後，你會看到一個前面沒見過的錯誤：

\begin{verbatim}
  File "curly-quote.py", line 1
    print(“Hello, world!”)
          ^
SyntaxError: invalid character '“' (U+201C)
\end{verbatim}

這行程式碼看起來跟正常版本幾乎一模一樣，肉眼甚至很難一眼看出差別，但 Python 只認得直
引號 \texttt{"} 跟 \texttt{'} 這兩個字元，彎引號對它來說是完全陌生的符號，連字串的開頭
都算不上，所以錯誤訊息連「引號沒關好」都說不上，直接說「這是一個我不認得的字元」。

修法
是把彎引號換回直引號重新打過；更根本的做法，是一開始就不要用 Word、Pages 這類軟體打程式
碼，寫程式全程都在編輯器裡進行，就不會有這個問題。

\keypoint{編輯器決定你打字的地方好不好用，但完全不影響 Python 怎麼看待你的程式碼，那是
直譯器的工作；反過來，如果編輯器（或任何軟體）偷偷「修飾」了你打的字元，像是把直引號換
成彎引號，Python 看到的就不再是你以為自己打的那些字元了。}

還有一個跟編輯器選擇有關、但這一章還用不到的提醒，先幫你打個底：前面說過，縮排在
Python 裡是語法的一部分，不是排版好看而已。\index{縮排}按下 Tab 鍵，有些編輯器會真的插入一個
「Tab 字元」，有些則會自動換算成幾個空白字元，兩種在螢幕上看起來可能一樣整齊，但對
Python 來說是不同的字元。

等到後面章節開始寫需要縮排的程式（例如條件判斷、迴圈）時，
如果同一份檔案裡有的縮排是空白、有的是 Tab，就可能出現縮排相關的錯誤，即使肉眼完全看
不出差異。VS Code 預設會把 Tab 鍵自動換算成空白字元，這正是這本書選它示範的原因之一，
但只要記得「縮排方式全篇要一致」這個原則，用哪一套編輯器都不會有問題。

\begin{drills}[label=drills:editor-pitfalls]
\item 如果你不小心在 Word 裡打完一段 Python 程式碼、直接複製貼上到 \texttt{hello.py}，
存檔執行後最可能先看到哪一種錯誤？
\item VS Code 跟 \texttt{nano}、\texttt{vim} 之類的編輯器，對 Python 直譯器來說有沒有
差別？為什麼？
\end{drills}

\section{兩種跟 Python 對話的方式}
\label{sec:interpreter-vs-file}

目前為止，你都是把程式碼寫進一個 \texttt{.py} 檔案，再叫 Python 執行整個檔案。但其實
還有另一種方式，更即時、更適合你只是想確認「這樣寫對不對」的時候：直接跟 Python 講話，
一行一行來，馬上看到結果，不用先存成檔案。

在終端機打 \texttt{python3}（macOS）或 \texttt{python}（Windows），不加任何檔名，直接
按 Enter，畫面上的提示字元會變成 \texttt{>>>}。這裡的「提示字元」指的是 \texttt{>>>}
這個符號本身，跟前面說的「命令提示字元」（Windows 那個終端機軟體的名字）是兩回事。
\index{提示字元}

看到
\texttt{>>>} 就代表 Python 現在在等你打指令，打完一行按 Enter，它馬上回你結果，不用存
檔、也不用另外下執行指令。
\index{互動模式}

\begin{example}{python}{(互動模式) 兩數相加}
>>> 3 + 5
8
>>> print("Python says hello")
Python says hello
\end{example}

每打一行，Python 就馬上算給你看、或印出結果，你不用做任何額外動作。想離開這個模式，打
\texttt{exit()} 按 Enter 即可。

如果打的那一行本身有問題呢？互動模式不會因此關掉，也不會把你踢出去，它只會針對那一行
抱怨，抱怨完之後照樣顯示 \texttt{>>>}，等你打下一行：

\begin{example}{python}{(互動模式) 出錯後還能繼續}
>>> print("Hello
  File "<stdin>", line 1
    print("Hello
          ^
SyntaxError: unterminated string literal (detected at line 1)
>>> print("This line still works")
This line still works
\end{example}

第一行漏打了結尾的引號，Python 老實地抱怨了一次；但這只影響那一行，緊接著打的下一行完
全不受影響，照樣正常執行。這是互動模式一個很實用的特性：你可以放心地隨手試東西，就算
某一行寫錯了，也不會弄壞什麼、更不用整個重開，改一改繼續打下去就好。

但如果你要做的事情有好幾個步驟，還想留著以後重複執行，互動模式就不太合適了：你打的每
一行，只要一離開這個視窗就全部消失，沒有留下任何檔案。這種時候，把程式碼寫進一個
\texttt{.py} 檔案、一次執行整個檔案，才是你要的做法，這種方式，我們就稱它「檔案模式」，
跟前面的互動模式相對。
\index{檔案模式}

\begin{example}{python}{(add.py) 兩數相加}
first_number = 3
second_number = 5
print(first_number + second_number)
\end{example}

\iosample{}{8 \\}

\texttt{first\_number = 3} 這一行，是把 3 這個值存進一個叫 \texttt{first\_number} 的
名字裡，之後只要打這個名字，就能把值取回來用；\texttt{second\_number = 5} 也是同樣的
道理，這個動作正式的名稱跟細節，後面章節會再講清楚，這裡你只要知道這行程式在做什麼就
好。

存檔後在終端機打 \texttt{python3 add.py}（macOS）或 \texttt{python add.py}
（Windows）。跟互動模式不一樣的地方是：Python 會先把整個檔案讀完、確認沒有問題，才會
一次從頭執行到尾，你只會在最後看到 \texttt{print} 印出來的結果，中間的計算過程不會顯示
（除非你自己也用 \texttt{print} 把它印出來）。

順帶一提，安裝 Python 直譯器時，其實還附贈了一個內建的互動模式視窗，叫做 IDLE，不需要
另外打開終端機也能用。這本書統一用終端機示範互動模式，因為終端機這個工具在後面章節（檔
案、資料夾操作）也會一直用到，先熟悉它比較划算；但如果你之後想試試 IDLE，兩者的行為基本
上是一樣的，只是視窗長相不同。

\stylepoint{互動模式很適合拿來「隨手驗算」：你不確定 \texttt{10 \% 3} 的結果是什麼、想
確認某個寫法到底怎麼運作，打開互動模式試一行就知道，比翻書找答案快多了。但只要是你想留
下來、之後還要重複執行、或超過幾行的程式，就該寫成 \texttt{.py} 檔案，這也是這本書接下
來幾乎所有範例都用檔案模式呈現的原因。}

\begin{drills}[label=drills:interpreter-vs-file]
\item 你想確認「17 除以 5 的餘數」是多少，查完這件事就結束了，不會再用到。你會選互動
模式，還是寫成 \texttt{.py} 檔案？
\item 你要寫一支下星期報告要上台示範、可能還會修改好幾次的小工具，你會選互動模式，還是
寫成 \texttt{.py} 檔案？
\item 在互動模式打錯一行、看到紅字錯誤之後，這個互動模式視窗會不會直接關掉？你需不需要
重新打開一次才能繼續？
\end{drills}

\section{讀懂你的第一個錯誤訊息}
\label{sec:reading-errors}

程式跑不動、螢幕跳出一片紅字，是每個學程式的人都會遇到的日常，不是天賦不夠的證據。這
一節要老實示範一次真正的錯誤長什麼樣子，直接讓你的程式壞掉一次，一行一行帶你讀懂
Python 到底想告訴你什麼，這比空講「以後你會看到錯誤訊息」有用得多。

\begin{example}{python}{(greet.py) 跟指定的人打招呼}
name = "Alice"
pint("Hello, " + name + "!")
\end{example}

這支程式想做的事很單純：把 \texttt{"Hello, "}、\texttt{name} 這個名字目前存的內容、跟
\texttt{"!"} 接在一起印出來。但仔細看第二行，\texttt{print} 被打成了 \texttt{pint}，
少了一個 \texttt{r}。存檔後執行，螢幕上會出現：

\begin{verbatim}
Traceback (most recent call last):
  File "greet.py", line 2, in <module>
    pint("Hello, " + name + "!")
    ^^^^
NameError: name 'pint' is not defined. Did you mean: 'print'?
\end{verbatim}

這段文字看起來很嚇人，但你不需要每個字都懂，只要抓住幾個重點，由下往上讀：
\index{NameError}

\begin{itemize}
\item 最後一行通常才是真正的問題所在：這裡是 \texttt{NameError}，意思是「Python 找不到
一個叫這個名字的東西」；後面的 \texttt{name 'pint' is not defined} 是說「\texttt{pint}
這個名字沒有被定義過」，你沒有教過 Python \texttt{pint} 是什麼，它當然不知道。它甚至
貼心地猜了一句「你是不是想打 \texttt{print}？」，這次剛好被它猜中了，不是每次都這麼
幸運，但這是個很好的起點。
\item 中間那行 \texttt{File "greet.py", line 2} 告訴你問題出在哪個檔案的第幾行，順著
這個線索回去對照原始碼，通常就能找到打錯字的地方。
\item \index{Traceback}第一行 \texttt{Traceback (most recent call last)} 只是在說「以下是程式執行到出錯
為止，經過的路徑」，對只有兩行程式的小範例來說不會太有用，但等你的程式變長、函式互相
呼叫之後，這行會幫上大忙。
\end{itemize}

\keypoint{讀錯誤訊息的訣竅是由下往上讀：最後一行的錯誤類型跟訊息，幾乎都是最直接、最
該先看的線索；上面那些 \texttt{File}、\texttt{line} 之類的資訊，是等你確定「錯在哪一種
問題」之後，拿來定位「錯在哪一行」用的。由上往下讀，反而容易先花時間看那些還不重要的
細節。}

把 \texttt{pint} 改回 \texttt{print}，存檔重跑一次，錯誤就會消失。這種「打錯字、看錯誤
訊息、抓出問題、改掉」的過程，你接下來寫程式的每一天都會一直重複；熟悉了以後，看到紅字
不會再心慌，反射動作會變成直接掃到最後一行找線索。

順帶一提，等你多學一點以後會發現，
Python 面對「事情可能會出錯」這件事，其實有一套特別的處理習慣，這本書後面會慢慢帶你
認識。

錯誤發生的時候，前面已經執行過的部分，會不會也跟著「消失」？答案是不會。Python 是真的
一行一行往下執行，執行到出問題的那一行才會停下來，停下來之前已經做過的事，效果已經生效
了，不會被撤銷：

\begin{example}{python}{(partial-run.py) 錯誤發生前，前面的部分已經執行過了}
print("Starting")
pint("Hello")
print("This line will not run")
\end{example}

存檔執行後，畫面上會是：

\begin{verbatim}
Starting
Traceback (most recent call last):
  File "partial-run.py", line 2, in <module>
    pint("Hello")
    ^^^^
NameError: name 'pint' is not defined. Did you mean: 'print'?
\end{verbatim}

第一行 \texttt{print("Starting")} 順利執行、顯示出來了；第二行才踩到打錯字的
\texttt{pint}，Python 執行到這裡才發現有問題，於是停下來、丟出錯誤；第三行則完全沒有
機會被執行到，畫面上不會出現「This line will not run」這幾個字。

這一點跟前面幾節看到的
\texttt{SyntaxError} 不太一樣：\texttt{SyntaxError} 是 Python 在還沒開始執行任何一行
之前，光是「讀」整個檔案的階段就先卡住了，所以那種情況下畫面上什麼都不會顯示；
\texttt{NameError} 這種錯誤，是真的執行到那一行才會發生，前面已經跑過的程式碼，成果
還在。

\keypoint{有些錯誤（像 \texttt{SyntaxError}）會在 Python 連一行都還沒開始執行之前就被
抓到，整份程式碼完全不會有任何輸出；有些錯誤（像 \texttt{NameError}）要等程式真的執行
到那一行才會發生，前面已經執行過的程式碼，效果不會因為後面出錯而消失。}

還有一種很容易讓人卡關的手誤：大小寫。對 Python 來說，\texttt{name} 跟 \texttt{Name}
是兩個完全不同的名字，就算只差一個字母的大小寫，也不例外：

\begin{example}{python}{(case-sensitive.py) 大小寫也算數}
name = "Alice"
print(Name)
\end{example}

\begin{verbatim}
Traceback (most recent call last):
  File "case-sensitive.py", line 2, in <module>
    print(Name)
          ^^^^
NameError: name 'Name' is not defined. Did you mean: 'name'?
\end{verbatim}

第一行明明已經用小寫的 \texttt{name} 存好了值，第二行卻打成大寫開頭的 \texttt{Name}，
對 Python 來說這是兩個毫不相干的名字，第二行問的那個 \texttt{Name} 它從來沒聽過，於是
一樣丟出 \texttt{NameError}。

這是初學者很容易忽略的細節：中文閱讀習慣不太在意英文字母
大小寫的差異，但 Python 每一個字元都算數，養成全部使用小寫字母、單字之間用底線連接的
命名習慣，可以大幅降低這一類手誤發生的機會（後面章節會再正式介紹這個命名習慣）。

\begin{drills}[label=drills:reading-errors]
\item 承接上面 \texttt{greet.py} 的錯誤訊息：如果你把第二行的 \texttt{pint} 改回
\texttt{print}，但不小心把第一行的變數名稱打成了 \texttt{nmae}（也就是
\texttt{nmae = "Alice"}，但第二行仍然寫 \texttt{name}），你覺得錯誤訊息的類型還會是
\texttt{NameError} 嗎？訊息裡找不到的名字會是哪一個？
\item 如果一段錯誤訊息的最後一行寫著 \texttt{NameError: name 'total' is not defined}，
在你動手改程式之前，你會先去程式碼裡找什麼？
\item 承接上面 \texttt{partial-run.py} 的例子：如果你把第三行也加上一個會出錯的手誤
（例如又打成 \texttt{pint}），你覺得畫面上除了第一行「Starting」以外，還會不會多印出
任何東西？
\item 如果一支程式先前定義了 \texttt{Total = 0}（大寫開頭），後面某一行卻寫
\texttt{print(total)}（小寫開頭），你覺得會發生什麼事？
\end{drills}

\section{為什麼值得學程式}
\label{sec:why-program}

電腦這個東西，拆開來看其實很笨：它不會自己思考「使用者大概想要什麼」，只會照著別人給
的指令，一步一步、原封不動地做。你其實已經做過這件事了：前面幾節你寫下
\texttt{print("Hello, world!")}，電腦真的把這行字顯示出來，那一刻，你已經是在對機器下
指令了，你剛才寫的那幾行，就是一支貨真價實的程式。

換個角度看，「電腦很笨」這件事，反而是它最有用的地方。人與人之間溝通，常常靠上下文、
語氣、常識來腦補對方沒講清楚的部分；電腦完全不會腦補，你講清楚幾分，它就照做幾分，一次
講清楚了，之後同一件事可以一模一樣地重複做一千次、一萬次，不會累、不會分心、不會今天
心情不好就漏做一步。

生活、課業、工作裡有不少瑣事其實正是這種重複、規則明確、只是很花
時間的事，例如：

\begin{itemize}
\item 把一份很長的名單依照某個條件挑出符合的項目。
\item 幫一批檔案依照固定規則改名字。
\item 把
每星期都要重複做一次的整理工作變成一次執行就好。
\end{itemize}

這些事不見得複雜，難的地方通常是「一直
重複做很煩」，剛好就是電腦最擅長的類型。

\keypoint{寫程式的本質，其實是練習把一件事講得夠精確，精確到一個完全不會腦補的對象也能
照做不誤。這跟你平常跟朋友、同事溝通的方式很不一樣：日常對話留一點模糊空間沒關係，對方
會自己補上；跟電腦溝通，模糊的地方就是程式出錯或行為不如預期的地方。這個習慣一旦養成，
其實對「把話講清楚」這件事本身也有幫助，不只是寫程式用得到。}

這本書接下來的旅程，想帶你去兩個地方，兩個一樣重要，沒有誰為誰服務。

\begin{itemize}
\item 第一個地方，是讓
你能自己動手寫小程式，處理這一類瑣碎重複的事，終點在於讓你以後遇到一件瑣事時，腦中會
反射性冒出「這個我可以寫程式做」的念頭，而且真的做得出來，不必先把 Python 這個語言的
每個角落都背熟才敢動手。
\item 第二個地方，是讓你看懂「軟體」這個東西究竟是怎麼被造出來的：
版本控制在幹嘛、一個專案的前端後端分別在講什麼、工程師團隊平常怎麼運作。這個目標特別是
為了你，如果以後工作上要跟工程師對接，開會時聽到「這個 API 還在 review」「這個 sprint
排不進去」會一頭霧水，這本書後段會有專門的章節，帶你把這些詞彙跟運作方式一次搞懂，讓你
能聽懂、也問得出對的問題。
\end{itemize}

這兩個目標都不需要你變成職業工程師：它們的重點，說到底是想讓
你多一種解決問題的能力，也多一種看懂周遭世界的方式。

\begin{drills}[label=drills:why-program]
\item 想一想你自己生活、課業或工作中，有沒有一件重複性很高、規則其實很明確的瑣事？先
寫下來，等讀完這本書後回頭看看，你的想法會不會改變。
\item 如果同事跟你說「這個功能還在 code review，這個 sprint 排不進去」，你現在聽得懂
嗎？如果還聽不懂，記得這本書後段有專門處理這件事的章節。
\end{drills}

\section{為什麼是 Python？這本書要帶你去哪裡}
\label{sec:why-python}

世界上能拿來寫程式的語言有很多種，C、Java、JavaScript 都是，各有各的個性跟用途。如果
把它們想成一群個性各異的人，Python 大概是那個說話最直白、從不拐彎抹角的傢伙。

它的名字來源常常讓人誤會：很多人以為是取自那種會咬人的蟒蛇，其實作者 Guido van Rossum
當年是照著他很喜歡的英國喜劇團體 Monty Python 取的名字，所以你偶爾會在 Python 的官方
文件或範例裡看到跟這個喜劇團體有關的冷笑話，這是個歷史彩蛋，不是巧合。

Python 的個性有幾個特別鮮明的地方：

\begin{itemize}
\item 它好讀，
語法設計得很接近你會寫的英文句子，連縮排本身都是語法的一部分，這代表你的程式碼被強迫
維持一致的排版。
\item 它內建的東西很多，標準函式庫豐富，很多常見任務不用自己從零開始寫，直接
拿現成的來用。
\item 它應用範圍廣，網站後端、資料分析、自動化瑣事、甚至一些遊戲，都看得到
Python 的身影。
\end{itemize}

再加上 Python 完全免費、任何人都能使用，全世界有非常多人在用它、也在
網路上互相幫忙解答問題，之後你卡關時，通常能找到已經有人問過一模一樣的問題，這也是很多
人把它當作「學程式的第一站」的原因。

在正式往下走之前，先跟你打個預防針。這本書會反覆用兩個方式，幫你把 Python 的直覺建立
起來，不只講一次就結束，你在後面很多章節都會重新看到它們。

\begin{itemize}
\item 一個是關於「把一個值指定給
一個名字」這件事，到底在做什麼：你可能已經對「變數」有一套從數學課帶來的直覺，但那套
直覺在 Python 裡有些地方會踩到坑，之後我們會花時間把這件事講清楚。
\item 另一個是關於 Python
遇到「事情可能會出錯」的時候，習慣的處理方式：是先做再說、出了問題再處理，還是先確認
安全才動手，這兩種習慣看似都合理，Python 社群比較偏好其中一種，這本書後面也會帶你認識
為什麼。
\end{itemize}

現在不用急著懂，只要有個印象就好，之後每次重新提起，你就會知道「啊，就是開頭
說的那個」。

接下來的路大致分成三段：

\begin{enumerate}
\item 前面十章左右是地基，變數、條件判斷、迴圈、函式、list 跟 dict
這些最常用的積木，這段會走得比較慢，每個技巧都會讓你真的能自己重現一次，不會匆匆帶過。
\item 中段開始把積木組合起來解決比較真實的任務，處理資料檔案、抓出並讀懂程式裡的錯誤、認識
一些常用的套件。
\item 最後幾章會拉高視角，帶你認識軟體業實際運作的樣子、串接一個真正的網路
服務，並且在闔上這本書之前，指給你看接下來還可以往哪些方向繼續走。
\end{enumerate}

你已經在第一章的
尾聲了，而且已經有一支自己跑得動的程式，這是個不錯的開始。

\section*{練習解答}
\addcontentsline{toc}{section}{練習解答}

\paragraph{第~\ref{sec:setup}~節}
\begin{enumerate}
\item 很可能是安裝 Python 時漏勾了「Add python.exe to PATH」這個選項，重新執行安裝
程式、選擇解除安裝後再重新安裝一次，記得勾選即可。
\item 直譯器負責真正執行程式碼；編輯器只負責讓你打字、整理、存檔，本身不會執行任何
程式碼。
\end{enumerate}

\paragraph{第~\ref{sec:hello-world}~節}
\begin{enumerate}
\item 會印出：My favorite programming language is Python。
\item 開放性練習，只要程式能正確執行、印出你自己的名字即可，記得引號要成對、括號別
漏掉。
\item Python 3 的 \texttt{print} 一定要接括號，漏了括號會出現：
\texttt{SyntaxError: Missing parentheses in call to 'print'. Did you mean print(...)?}，
把文字用括號包起來就能修好，例如改成 \texttt{print("Hello")}。
\item 會出現 \texttt{SyntaxError: unterminated string literal (detected at line 1)}，
跟本節「頭尾引號種類不一致」那個例子一樣，因為 Python 只認雙引號對雙引號、單引號對單
引號，頭尾種類不同時，等於整個字串都沒有真正被關起來。
\end{enumerate}

\paragraph{第~\ref{sec:editor-pitfalls}~節}
\begin{enumerate}
\item 最可能看到 \texttt{SyntaxError: invalid character} 這一類錯誤，因為 Word 會把
直引號自動換成彎引號，Python 完全不認得彎引號這個字元。
\item 沒有差別。編輯器只負責讓你打字、存檔，真正決定程式碼能不能執行的是直譯器；只要
存出來的是純文字、沒有被夾帶額外排版資訊的 \texttt{.py} 檔案，用哪一套編輯器寫都一樣。
\end{enumerate}

\paragraph{第~\ref{sec:interpreter-vs-file}~節}
\begin{enumerate}
\item 互動模式，查完就結束，不需要留下任何檔案。
\item 檔案模式，需要重複執行、之後還會修改，寫成 \texttt{.py} 檔案才留得住。
\item 不會直接關掉，也不需要重新打開。互動模式只會針對出錯的那一行顯示錯誤訊息，之後
照樣顯示 \texttt{>>>}，可以直接接著打下一行。
\end{enumerate}

\paragraph{第~\ref{sec:reading-errors}~節}
\begin{enumerate}
\item 還是 \texttt{NameError}，但這次找不到的名字變成 \texttt{name}：問題出在「兩處
拼法對不起來」，一邊打成 \texttt{nmae}，另一邊還在用 \texttt{name}。
\item 先去程式碼裡找有沒有一個叫 \texttt{total} 的變數被使用了，但從來沒有被指定過值。
可能是名字打錯，也可能是那一行原本該賦值的程式碼被漏掉了。
\item 不會。程式一樣會先印出「Starting」，接著在第二行遇到 \texttt{pint} 就停下來丟出
\texttt{NameError}，第三行不管有沒有手誤都不會被執行到，因為程式根本沒有機會跑到那
一行。
\item 會出現 \texttt{NameError: name 'total' is not defined. Did you mean: 'Total'?}，
因為 \texttt{Total} 跟 \texttt{total} 對 Python 來說是兩個不同的名字，前面定義的是大寫
開頭的 \texttt{Total}，後面卻查詢小寫開頭的 \texttt{total}，等於在問一個從沒定義過的
名字。
\end{enumerate}

\paragraph{第~\ref{sec:why-program}~節}
\begin{enumerate}
\item 開放性練習，沒有固定答案，重點是先具體想過一次「這件事似乎可以用程式解決」的
感覺。
\item 開放性練習：現在聽不懂完全正常，這本書第 \ref{ch:industry} 章會專門帶你認識這些詞彙背後在講
什麼。
\end{enumerate}

\chapter{變數與運算}
\label{ch:variables}

上一章你已經讓電腦照你的話做了一件事：印出一行字。這一章要教的東西聽起來很基礎，但
其實是接下來整本書都會反覆用到的地基，所以我們刻意放慢腳步。先不談任何原理，直接動手：
建立一個變數、把值印出來、學會怎麼幫變數取名字，就跟上一章學 \texttt{print()} 一樣，
先讓你的手熟悉這個動作。

接著回頭拆解剛剛發生的事：當你寫下 \texttt{x = 5} 的時候，
Python 到底做了什麼？很多人對「變數」的第一印象來自數學課或日常生活，習慣把它想成一個
裝東西的盒子，把值放進去。這個直覺可以應付一部分情況，但它會在後面的章節害你踩坑，所以
這一章要花完整篇幅，把更接近 Python 實際運作方式的圖像講清楚。

指定這個動作，做的其實是
把一張名字牌貼到一個值上，跟複製一份值放進盒子完全是兩回事。這是這本書會反覆呼應的兩個
核心比喻之一，這裡是第一次完整登場，之後每一章只要遇到「這裡容易跟直覺打架」的地方，都
會回頭喊一聲「還記得名字牌嗎」。

再來把視角拉開一點：你目前看過的值，其實可以分成幾個
固定的「種類」，包括整數、小數、文字，還有代表對錯的布林值，Python 把每個值歸到哪一類，
可以直接問它。

最後解決一個刻意留下的問題：如果你手上的值種類不對，該怎麼把它轉成你要的
種類，以及轉不過去的時候 Python 會怎麼抗議。整章結束時，你不只會建立變數、分辨型別、
轉換型別，更重要的是，你會多一個比盒子更準確的方式，去理解 Python 裡「賦值」這件事到底
在做什麼，這個理解會在後面好幾章不斷派上用場。

\section{建立你的第一個變數}
\label{sec:first-variable}

寫程式的時候，你常常需要把一個值先記下來，之後可以重複拿出來用，不用每次都重新打一次。
Python 讓你做這件事的方法很直接：幫這個值取一個名字，之後只要打這個名字，就能把值
取回來。這個名字叫做「變數」，取名字、存值的這個動作，叫做「賦值」。
\index{變數}
\index{賦值}

\begin{code}{python}{Create a variable and assign a value}
variable_name = value
\end{code}

等號左邊是你幫值取的名字，右邊是你想存的值。這一行執行完，Python 就會記得這個名字
對應這個值，直到你改變主意為止。

\begin{example}{python}{(today-note.py) Record today's weather}
temperature = 28
is_sunny = "sunny"
print(temperature)
print(is_sunny)

temperature = 31
print(temperature)
\end{example}

\iosample{}{28 \\ sunny \\ 31 \\}

前兩行分別建立了兩個變數：\texttt{temperature} 存一個數字，\texttt{is\_sunny} 存一段
文字（文字前後要用引號包起來，這件事上一章已經看過）。中間兩個 \texttt{print} 把它們
的值印出來確認。

接著第五行很關鍵：\texttt{temperature = 31} 又對同一個變數賦值一次。
效果是把原本存的 \texttt{28} 蓋掉、換成 \texttt{31}，\texttt{temperature} 這個名字
本身沒有變，變的只有它現在對應的值。最後一次 \texttt{print(temperature)} 印出的是
\texttt{31}，因為舊的 \texttt{28} 已經沒有任何東西記得它了。

你可能會好奇，變數重新賦值的時候，Python 到底做了什麼事？是把舊的 \texttt{28} 原地
換成 \texttt{31}，還是有什麼別的事情發生？這個問題先放著，下一段主線就要正式回答它。

\section{幫變數取名字}
\label{sec:variable-naming}

變數名不是隨便打幾個字都可以，Python 對「什麼樣的名字合法」有明確規則；就算合法，
好的變數名跟隨便亂取的變數名，對「未來的你」跟其他要讀你程式碼的人，差別也很大。

\begin{code}{python}{Legal and illegal variable names}
score = 90          # 合法：全部小寫，有意義
student_age = 15    # 合法：多個單字用底線連接
_temp = 100          # 合法：可以用底線開頭，但沒有特別原因不建議這樣做
2nd_place = "silver"   # 不合法：不能以數字開頭
class = "trial"       # 不合法：class 是 Python 保留字
\end{code}

規則整理起來大致有三條：

\begin{itemize}
\item 變數名只能用英文字母、數字、底線這三種字元組成，而且\textbf{不能}以數字開頭。
\item 不能拿 Python 已經保留給自己用的字（像 \texttt{class}、\texttt{for}、
\texttt{if} 這些，統稱「保留字」）當變數名。
\index{保留字}
\item Python 對大小寫敏感，\texttt{score} 跟 \texttt{Score} 是兩個不同的變數。
\end{itemize}

慣例上，這本書接下來的變數名都會用全小寫加底線分隔
多個詞的寫法（例如 \texttt{student\_age}），這只是 Python 社群公認的命名風格，語法上
沒有強制要求，但照著寫會讓你的程式碼一眼就跟其他 Python 開發者的習慣接得上。

寫程式的時候，比起搞不清楚規則，手滑打錯字才是更常見的失誤來源。看看下面這支程式：

\begin{example}{python}{(score-typo.py) Print a student's score}
student_score = 88
print(studnet_score)
\end{example}

第二行的 \texttt{studnet\_score} 把 \texttt{student\_score} 的字母順序打亂了。存檔
執行後，會看到這樣的錯誤：

\begin{verbatim}
Traceback (most recent call last):
  File "score-typo.py", line 2, in <module>
    print(studnet_score)
          ^^^^^^^^^^^^^
NameError: name 'studnet_score' is not defined. Did you mean: 'student_score'?
\end{verbatim}

由下往上讀：最後一行 \texttt{NameError} 是說「Python 找不到一個叫這個名字的變數」。
你在第一行明明建立了 \texttt{student\_score}，但第二行打的是完全不同拼法的
\texttt{studnet\_score}，對 Python 來說這是兩個毫不相干的名字，它不會自動幫你猜你是不是
打錯字（不過這裡它很貼心地猜對了一次）。

中間那行告訴你問題出在第 2 行。修法就是把兩處
拼法對齊，不管你選哪種拼法都行，重點是同一個變數從頭到尾要拼得一致。

\keypoint{Python 不會幫你檢查「這個名字是不是你原本想打的那個」，它只在乎「這個名字有沒
有被定義過」。變數名前後拼法只要有一個字不一樣，Python 就會當成完全不同的兩個東西。這
跟人類讀者會自動腦補、猜你想講什麼，是完全不同的兩套邏輯。}

那如果你取的名字本身就不合法呢？比如說，你想用一個以數字開頭的名字：

\begin{example}{python}{(bad-name.py) An illegal variable name}
1st_place = "Alice"
print(1st_place)
\end{example}

這支程式甚至還沒開始執行就會被擋下來：

\begin{verbatim}
  File "bad-name.py", line 1
    1st_place = "Alice"
    ^
SyntaxError: invalid decimal literal
\end{verbatim}

\texttt{SyntaxError} 代表 Python 光是「讀」你寫的程式碼這一步就卡住了。它看到
\texttt{1st\_place} 開頭的 \texttt{1}，以為你要寫一個數字，結果後面又接了不是數字的
字母，完全看不懂你想表達什麼，直接拒絕往下讀。

同樣的道理，如果你把 \texttt{class} 這種
保留字拿來當變數名（例如 \texttt{class = "初賽"}），也會得到 \texttt{SyntaxError:
invalid syntax}，因為 Python 已經把 \texttt{class} 這個字保留給它自己內部的用途（後面
第 \ref{ch:objects} 章會學到 \texttt{class} 是用來做什麼的），不允許你拿去當普通的變數名使用。這兩種
情況修法都一樣：換一個合法的名字。

\begin{drills}[label=drills:variable-naming]
\item 下面哪些變數名合法？哪些不合法？如果不合法，說出原因。
\texttt{my\_name}、\texttt{2apples}、\texttt{total\_price}、\texttt{if}、
\texttt{User1}。
\item 如果你先寫了 \texttt{favorite\_color = "blue"}，接著寫
\texttt{print(favorite\_colour)}（拼法多了一個 \texttt{u}），執行後你覺得會出現什麼
錯誤訊息？
\item 如果你把 \texttt{total = 100} 寫成 \texttt{3total = 100}，執行後會出現的是
\texttt{NameError} 還是 \texttt{SyntaxError}？為什麼？
\end{drills}

\section{取名字的實務眉角}
\label{sec:naming-practices}

上一節教的是「合法」的底線，這一節要往前一步，談「合法但很糟」的名字會怎麼害到你，以及
幾個實務上值得養成的習慣。

Python 內建了一批像 \texttt{print}、\texttt{type} 這樣隨時可以呼叫的函式，這些名字
本身沒有被保護起來，並不會因為它們是內建的就禁止你拿來當變數名。問題是，一旦你把值賦給
這個名字，Python 就會用你賦的新值蓋掉原本那個函式，之後整份檔案裡只要再用到這個名字，
拿到的都是你賦的新值，跟原本的函式已經沒有關係了。

\begin{example}{python}{(shadow-type.py) Shadowing a built-in function's name}
type = "integer"
print(type)

x = 5
print(type(x))
\end{example}

執行後，前兩行先正常印出，但第五行整個崩潰：

\begin{verbatim}
Traceback (most recent call last):
  File "shadow-type.py", line 5, in <module>
    print(type(x))
          ~~~~^^^
TypeError: 'str' object is not callable
\end{verbatim}

\texttt{integer} 先印出來沒有問題，因為這時候 \texttt{type} 這個名字對應的正是你剛剛賦的
字串。問題出在第五行：你想呼叫「查詢型別」的那個 \texttt{type()}，但 Python 手上的
\texttt{type} 這時候對應的已經是你在第一行賦的字串 \texttt{"integer"}，原本的函式早就被
蓋掉了。字串不能像函式一樣被呼叫（後面加括號），所以噴出 \texttt{TypeError: 'str'
object is not callable}
（「字串物件不能被呼叫」）。

這支程式本身沒有任何一行語法寫錯，Python 也不會在你賦值給
\texttt{type} 的當下警告你「這名字已經被用掉了」，等到你真的想用回原本的 \texttt{type()}
才會出事，而且出事的地方往往離你打錯的地方隔了好幾行，抓起來比 \texttt{NameError}
麻煩得多。

\keypoint{變數名跟函式名共用同一套命名空間，Python 不會替 \texttt{print}、
\texttt{type}、\texttt{int} 這些內建函式加任何保護。幫變數取名字之前，養成習慣避開這些
常見的內建名稱，尤其是本章介紹過的 \texttt{type}、\texttt{int}、\texttt{float}、
\texttt{str}、\texttt{bool}、\texttt{print} 這幾個，之後每一章學到新的內建函式時，也要
記得同一件事。}

除了不要蓋掉內建名稱，變數名本身的可讀性也值得多花幾秒鐘想。回頭看看本章第一個範例：
\texttt{temperature}、\texttt{is\_sunny} 這兩個名字，就算不看任何註解，光看名字也大概
猜得出裡面裝的是什麼。

如果當初把它們取名叫 \texttt{t}、\texttt{x}，程式一樣能執行，但
幾天後回頭看，你會完全忘記 \texttt{t} 到底是溫度、時間，還是別的東西。變數名不需要很長，
但至少要能讓「未來的你」一眼看懂它裝的是什麼。

\begin{code}{python}{Poor names vs. descriptive names}
t = 28          # 看不出 t 裝的是什麼
s = "sunny"       # 同樣看不出來

temperature = 28       # 一看就知道裝的是溫度
weather = "sunny"        # 一看就知道裝的是天氣
\end{code}

\stylepoint{布林值變數習慣用 \texttt{is\_} 或 \texttt{has\_} 開頭命名，像本章前面用過的
\texttt{is\_sunny}、\texttt{is\_weekend}，光看名字就知道這是一個「對／錯」的答案，而不
只是一段文字或一個數字。這屬於命名習慣，Python 語法本身並沒有強制規定，但幾乎所有
Python 開發者都會這樣取名，照著做能讓你的程式碼更容易讓別人看懂。}

最後補充一個大部分人沒想過的合法名字：Python 其實也允許用中文字當變數名。

\begin{code}{python}{Non-English variable names are legal, but this book won't use them}
café = 90  # 帶重音符號或其他非英文字元的名字也合法
print(café)
\end{code}

上面這段程式執行起來完全沒有問題，Python 對變數名裡的重音符號、中文字等非英文字元都
沒有意見。不過這本書接下來還是會統一用英文字命名，原因有兩個：

\begin{itemize}
\item 幾乎所有 Python 的官方文件、套件、其他人寫的
程式碼都是英文命名，習慣英文命名會讓你將來讀別人的程式碼更順手。
\item 中文字之間沒有
空白，名字長一點的時候容易分不清楚是一個詞還是好幾個詞黏在一起，不像英文可以用底線清楚
斷開每個詞。
\end{itemize}

知道「合法」跟「這本書會怎麼寫」是兩碼子事，之後遇到別人用中文命名的程式碼，
至少不會看不懂發生了什麼事。

\begin{drills}[label=drills:naming-practices]
\item 如果你在程式最前面寫了 \texttt{print = 5}，接下來這一行 \texttt{print("hi")}
會發生什麼事？錯誤類型是什麼？
\item 下面兩個變數名，哪一個比較適合拿來存「這位學生是不是及格」這個答案：
\texttt{result} 還是 \texttt{is\_passed}？為什麼？
\item 承接 \texttt{shadow-type.py} 這個範例，如果把第一行 \texttt{type = "integer"} 拿掉，
其他都不改，這支程式還會出錯嗎？
\end{drills}

\section{名字牌：\texttt{=} 到底在做什麼}
\label{sec:name-tags}

上一段你已經自己動手把 \texttt{temperature} 重新賦值過一次，也看過 Python 怎麼抱怨
拼錯的變數名。現在該回頭問一個更根本的問題：\texttt{x = 5} 這一行執行的時候，Python
到底做了什麼事？

多數人對這個問題的第一直覺，來自數學課或日常生活，把變數想成一個「盒子」：\texttt{x = 5}
就是「把 5 放進一個叫 x 的盒子裡」；如果之後又寫 \texttt{y = x}，直覺會覺得「把 x 盒子
裡的東西複製一份，放進 y 盒子」。這個直覺能解釋你目前看過的大部分行為，可惜它跟 Python
實際運作的方式對不上，而且會在後面的章節（尤其是第 \ref{ch:lists} 章的 list）讓你踩到一個解釋不通的
坑。這一節要把更接近事實的圖像講清楚。

Python 實際的運作方式是這樣：\texttt{x = 5} 的意思是「先造一個值是 5 的物件，再把一張
寫著 \texttt{x} 的名字牌，貼到這個物件上」。把變數想成貼在物件上的名字牌，會比想成盒子
更準確：名字牌本身不裝任何東西，它只是指向某個物件，告訴 Python「\texttt{x} 這個名字，
現在對應到的是那個物件」。
\index{名字牌模型}
\index{物件}

\begin{code}{python}{Attach a name tag to an object}
name = value
\end{code}

這跟前面看到的賦值語法長得一模一樣，但現在你知道它背後真正在做的事了。那如果接著寫
\texttt{y = x} 呢？意思是「把另一張名字牌 \texttt{y}，也貼到 \texttt{x} 現在指的那個
物件上」，效果是兩張名字牌指著同一個東西，並沒有另外複製一份新的 5 出來。

\begin{example}{python}{(nametag-demo.py) Two name tags pointing at the same value}
x = 5
y = x
x = 6
print(y)
\end{example}

\iosample{}{5 \\}

如圖~\ref{fig:name-tags} 所示：第二行 \texttt{y = x} 執行完，\texttt{x} 跟 \texttt{y}
這兩張名字牌，指的是同一個 5。接下來第三行 \texttt{x = 6} 執行的時候，Python 造了一個
全新的、值是 6 的物件，然後把 \texttt{x} 這張名字牌\textbf{改貼}到這個新物件上，那個 5
本身完全沒有被「原地改成」6。

\texttt{y} 完全沒有被驚動，它從頭到尾指的都還是原本那個 5。
所以最後 \texttt{print(y)} 印出來的是 \texttt{5}，維持在重新賦值之前的狀態，跟 \texttt{x}
現在指的 \texttt{6} 是兩個不同的物件。

\begin{center}
\begin{tikzpicture}[node distance=8mm and 14mm]
    \node[startend] (xtag1) {\texttt{x}};
    \node[startend, right=20mm of xtag1] (ytag1) {\texttt{y}};
    \node[process, fill=greenFill, draw=greenStroke, below=12mm of xtag1, xshift=10mm] (five1) {\texttt{5}};
    \draw[diagarrow] (xtag1.south) -- (five1.north west);
    \draw[diagarrow] (ytag1.south) -- (five1.north east);
\end{tikzpicture}
\hspace{1.2cm}
\begin{tikzpicture}[node distance=8mm and 14mm]
    \node[startend] (ytag2) {\texttt{y}};
    \node[startend, right=20mm of ytag2] (xtag2) {\texttt{x}};
    \node[process, fill=greenFill, draw=greenStroke, below=12mm of ytag2] (five2) {\texttt{5}};
    \node[process, fill=memFill, draw=memStroke, below=12mm of xtag2] (six2) {\texttt{6}};
    \draw[diagarrow] (ytag2.south) -- (five2.north);
    \draw[diagarrow] (xtag2.south) -- (six2.north);
\end{tikzpicture}
\end{center}
\diagcaption{左圖：\texttt{y = x} 執行完，兩張名字牌指向同一個物件。右圖：
\texttt{x = 6} 執行完，\texttt{x} 改貼到一個全新的物件上，\texttt{y} 仍指著原本那個 5，
兩者從此各自獨立。\label{fig:name-tags}}

這裡要老實講一件事：因為數字這種值不能被「原地修改」，只能整個換掉，沒辦法把 5 直接
改寫成 6，所以在這個例子裡，名字牌模型跟「複製盒子」模型剛好給出一模一樣的結果，你單靠
這個例子分不出兩個模型的差別。這不代表兩個模型是等價的，只是這裡剛好看不出差異。
\index{盒子模型}

等你
學到第 \ref{ch:lists} 章的 list，會看到一種可以「原地修改」的東西，到時候名字牌模型跟盒子模型會給出
完全不同的答案，也是名字牌模型真正派上用場的時候。先把這個印象留著，之後回頭看這一段，
你會更清楚為什麼現在要花這麼多篇幅講一個看起來沒差的東西。

\keypoint{\texttt{x = 6} 執行的效果，是造一個新的物件 6，再把 \texttt{x} 這張名字牌改貼
過去，\texttt{x} 原本指的那個物件並沒有被原地改掉。已經貼在原物件上的其他名字牌（像這裡
的 \texttt{y}）完全不受影響，因為它們指的是另一個物件，跟 \texttt{x} 現在指哪裡無關。}

名字牌模型還能解釋另一種寫法：一行裡同時幫好幾個名字賦同一個值。

\begin{code}{python}{Assign multiple names in one line}
name1 = name2 = name3 = value
\end{code}

\begin{example}{python}{(chained-assign.py) Three name tags attached to the same object}
a = b = c = 0
print(a, b, c)

b = 5
print(a, b, c)
\end{example}

\iosample{}{0 0 0 \\ 0 5 0}

\texttt{a = b = c = 0} 這一行做的事，是先造一個值是 0 的物件，再一口氣把 \texttt{a}、
\texttt{b}、\texttt{c} 三張名字牌全部貼到同一個物件上，效果等同於分開寫三行
\texttt{a = 0}、\texttt{b = 0}、\texttt{c = 0}，只是更精簡。第一次 \texttt{print}
印出 \texttt{0 0 0}，三個名字指的都是同一個 0。

接著 \texttt{b = 5} 只是把 \texttt{b}
這張名字牌改貼到一個新的物件 5 上，\texttt{a} 跟 \texttt{c} 完全不受影響，還是指著原本
那個 0，所以第二次 \texttt{print} 印出 \texttt{0 5 0}。這跟前面 \texttt{y = x} 之後
\texttt{x = 6} 的例子，用的是同一套邏輯：改貼一張名字牌，不會影響貼在同一個物件上的
其他名字牌。

\begin{drills}[label=drills:name-tags-chained]
\item 執行下面這段程式，最後 \texttt{print(x, y, z)} 會印出什麼？
\begin{verbatim}
x = y = z = "start"
x = "changed"
print(x, y, z)
\end{verbatim}
\item 如果把上面這段程式的第二行改成 \texttt{y = "changed"}，其他都不變，
\texttt{print(x, y, z)} 印出來的結果會不會不一樣？為什麼？
\end{drills}

\section{型別是物件的屬性}
\label{sec:type-on-value}

在往下講之前，先把一個詞講清楚：你目前看過的值，其實可以分成幾個固定的「種類」，包括
整數、小數、文字，還有代表對錯的布林值。Python 把每一個值都歸到的這種固定類別，就叫
「型別」。接下來這一節跟下一節都會用到這個詞。
\index{型別}

延續上一節的名字牌模型，還有一件事值得特別拿出來講：同一個變數，可以先後指向完全
不同種類的值。

\begin{example}{python}{(type-on-value.py) The same name tag re-attached to a different kind of value}
x = 5
print(type(x))
x = "hi"
print(type(x))
\end{example}

\iosample{}{<class 'int'> \\ <class 'str'>}

這裡先小小用一下 \texttt{type()} 這個函式（函式：一個有名字、後面帶著括號、把東西
放進括號裡就會做某件事或給你一個結果的積木，下一節會正式介紹它）。
\index{函式}
\index{type() 函式}
這裡的
\texttt{type()} 會告訴你括號裡的值屬於哪一種型別。第一次呼叫 \texttt{type(x)}，
\texttt{x} 指的是數字 5，所以印出 \texttt{<class 'int'>}（\texttt{int} 是「整數」的
意思）。

第三行讓 \texttt{x} 改指向文字 \texttt{"hi"}，這完全合法，Python 不會有任何
意見。第二次呼叫 \texttt{type(x)}，印出的變成 \texttt{<class 'str'>}（\texttt{str}
是「字串」的意思，也就是一段文字）。

這個結果如果套用「盒子」模型會很奇怪，好像同一個盒子，裝的東西的「類型」自己變了。但用
名字牌模型看就完全合理：\texttt{x} 只是一張名字牌，一開始貼在數字 5 這個物件上，後來
改貼到字串 \texttt{"hi"} 這個物件上。

型別其實是每一個物件自己帶的屬性，5 這個物件永遠是
\texttt{int}，\texttt{"hi"} 這個物件永遠是 \texttt{str}，只是 \texttt{x} 這張名字牌
先後指過這兩個不同的物件而已。變數本身沒有型別可言，型別長在物件（也就是值）身上。

\begin{drills}[label=drills:name-tags]
\item 執行下面這段程式之前，先用眼睛預測：\texttt{b} 最後會印出什麼？
\begin{verbatim}
a = 10
b = a
a = 20
print(b)
\end{verbatim}
\item 承上題，如果把最後一行改成 \texttt{print(a)}，答案會不會不一樣？為什麼？
\item 下面這段程式執行完，\texttt{type(x)} 印出來的會是哪個型別？
\begin{verbatim}
x = "3"
x = 3
print(type(x))
\end{verbatim}
\end{drills}

\section{型別家族：\texttt{int}／\texttt{float}／\texttt{str}／\texttt{bool}}
\label{sec:type-family}

你目前已經自然地用過數字（像 \texttt{28}、\texttt{5}）跟文字（像 \texttt{"sunny"}、
\texttt{"hi"}）這兩種值，也隱約感覺到它們是「不一樣的東西」。上一節已經提過，這種
固定的類別就叫「型別」。這一節要正式帶你認識四種最基本的型別。

下面先看看這四種型別的「字面值」長什麼樣子。字面值指的是直接寫在程式碼裡的值本身，
像 \texttt{100} 或 \texttt{"some text"}，跟變數名是不同的東西：
\index{字面值}

\begin{code}{python}{Literals of the four basic types}
100          # int：整數
3.14         # float：帶小數點的數字
"some text"    # str：字串（文字）
True         # bool：布林值（只有 True 或 False）

type(value)     # 查詢一個值的型別
\end{code}

這四種型別分別是：

\begin{itemize}
\item \texttt{int} 是整數，像 \texttt{15}、\texttt{-3}、\texttt{0}，沒有小數點。
\index{int}
\item \texttt{float}
是帶小數點的數字，像 \texttt{3.14}、\texttt{-0.5}，就算小數點後面是 0 也算
\texttt{float}（例如 \texttt{2.0}）。
\index{float}
\item \texttt{str} 是字串，也就是一段文字，一定要用
單引號或雙引號包起來，Python 才會把它讀成文字，跟變數名或程式碼的其他部分區分開來。
\index{str}
\item \texttt{bool} 是布林值，只有 \texttt{True}（真、對、是）跟 \texttt{False}（假、錯、否）
兩個值可選，用來表示「二選一」的答案，注意 \texttt{True}／\texttt{False} 開頭要大寫，
這是 Python 的固定寫法。
\index{bool}
\end{itemize}

如果你不確定手上這個值是哪一種型別，直接問 Python 就好，這就是
\texttt{type()} 函式的用途。

值得一提的是，\texttt{int} 在 Python 裡沒有固定的位數上限，你可以放心存一個非常大的
整數（例如 \texttt{99999999999999999999}），不會像某些程式語言那樣因為數字太大而算錯
或直接出錯，這在後面章節處理大量資料或迴圈累加時，可以少擔心一件事。

\stylepoint{如果數字本身很長，可以在數字中間插入底線讓人眼睛比較好讀，像
\texttt{1\_000\_000}（一百萬），底線不會影響數值本身，Python 讀到的還是同一個數字。}

另外，\texttt{float} 還有一種比較精簡的寫法，叫科學記號：\texttt{3.14e2} 代表
\texttt{3.14} 乘以 \texttt{10} 的 \texttt{2} 次方，也就是 \texttt{314.0}；\texttt{e}
後面的數字也可以是負的，例如 \texttt{5e-3} 代表 \texttt{0.005}。這種寫法平常不一定
用得到，但遇到很大或很小的數字（像天文數字，或是很小的誤差值）時，會比一長串 0 好讀
很多，日後處理科學或工程類的數據時可能會派上用場。

\begin{example}{python}{(type-family.py) Query the four types}
age = 15
price = 39.9
city = "Boston"
is_weekend = True

print(type(age))
print(type(price))
print(type(city))
print(type(is_weekend))
\end{example}

\iosample{}{<class 'int'> \\ <class 'float'> \\ <class 'str'> \\ <class 'bool'>}

四個變數各自存了一種型別的值，四次 \texttt{type()} 呼叫依序印出 \texttt{int}、
\texttt{float}、\texttt{str}、\texttt{bool}，跟你一開始寫下這些字面值時的直覺一致。
\texttt{bool} 這裡你只需要先認得它的存在、知道 \texttt{True}／\texttt{False} 是它的
字面值就夠了，它在條件判斷（第 \ref{ch:conditionals} 章）裡還有更多學問，那時候再深入。

\begin{drills}[label=drills:type-family-literals]
\item \texttt{type(0)} 跟 \texttt{type(0.0)} 印出來的型別一樣嗎？為什麼？
\item \texttt{type(True)} 會印出什麼？
\item \texttt{print(1\_000\_000)} 印出來的內容看起來會是什麼樣子？底線符號會不會被一起
印出來？
\end{drills}

\section{算術：讓數字真的算起來}
\label{sec:arithmetic}

到目前為止，變數只是「把值存起來、再拿出來」。但程式真正有用的地方，是它能替你算東西。
這一節把 Python 的算術運算一次講清楚——其中有兩個運算子（\texttt{//} 跟 \texttt{\%}）
在數學課上沒有專屬符號，卻是寫程式時最常用到的兩個，值得特別花時間。

\begin{code}{python}{Python 的算術運算子}
a + b       # 加
a - b       # 減
a * b       # 乘
a / b       # 除
a // b      # 除完只取整數部分（商）
a % b       # 除完只取餘數
a ** b      # a 的 b 次方
\end{code}

前四個跟數學課上的用法幾乎一樣，只有兩個符號要重新記：乘法是星號 \texttt{*}（不是
\texttt{×}），除法是斜線 \texttt{/}（不是 \texttt{÷}），這兩個符號鍵盤上打得出來。

\begin{example}{python}{(basic-arithmetic.py) 四則運算}
print(7 + 3)
print(7 - 3)
print(7 * 3)
print(7 / 3)
print(6 / 3)
\end{example}

\iosample{}{10 \\ 4 \\ 21 \\ 2.3333333333333335 \\ 2.0 \\}

前三行都符合直覺。第四行 \texttt{7 / 3} 除不盡，Python 給你一長串小數，這也不意外。
真正值得停下來看的是最後一行：\texttt{6 / 3} 明明整整除得盡，答案卻印出
\texttt{2.0}，帶著一個小數點。這不是 Python 算錯，是它的固定規則——用 \texttt{/} 做除法，
結果一律帶小數點。為什麼會這樣，下一節查型別的時候你就會看到答案。

\subsection*{商與餘數：\texttt{//} 跟 \texttt{\%}}

有些時候你要的根本不是小數。假設一段影片長 100 分鐘，你想印出「幾小時又幾分鐘」，
\texttt{100 / 60} 算出來的 \texttt{1.6666666666666667} 完全幫不上忙——你要的是
「1 小時」跟「剩下 40 分鐘」這兩個數字。

這正是 \texttt{//} 跟 \texttt{\%} 的用途：\texttt{//} 給你「除完整整有幾份」，
\texttt{\%} 給你「分完之後剩多少」。
\index{//（整數除法）}\index{\%（取餘數）}

\begin{example}{python}{(minutes-to-hours.py) 把分鐘拆成小時與分鐘}
total_minutes = 100

hours = total_minutes // 60      # 100 分鐘裡面有幾個完整的 60 分鐘
minutes = total_minutes % 60     # 分完之後還剩下幾分鐘

print(hours)
print(minutes)
\end{example}

\iosample{}{1 \\ 40 \\}

\texttt{100 // 60} 得到 \texttt{1}（100 裡面裝得下一個完整的 60），
\texttt{100 \% 60} 得到 \texttt{40}（裝完一個 60 之後剩下的）。這兩個運算子幾乎總是成對
出現，因為它們回答的是同一個問題的兩半：「分成幾份」跟「剩多少」。

\keypoint{\texttt{//} 取商、\texttt{\%} 取餘數，兩個合起來就能把一個數字拆成「幾個整份、
剩下多少」。 只要你遇到「把總數換算成大小單位」的情境——秒換成分、分換成時、總價換成幾
張鈔票——用的都是這一組，不是 \texttt{/}。}

\texttt{\%} 還有一個用途，出現的頻率高到值得單獨記住：判斷一個數字能不能被整除。

\begin{example}{python}{(even-or-odd.py) 判斷奇偶}
number = 14

print(number % 2)
print(number % 2 == 0)
\end{example}

\iosample{}{0 \\ True \\}

一個數字除以 2 如果餘數是 \texttt{0}，代表它整除得盡，也就是偶數；餘數是 \texttt{1}
就是奇數。第二行把這個判斷寫成一個比較式，算出來是 \texttt{True}——這個寫法在第
\ref{ch:conditionals} 章條件判斷會大量用到。 同樣的手法換個數字就是「能不能被 3 整除」
（\texttt{number \% 3 == 0}），是很常見的判斷方式。

\subsection*{次方與兩個好用的內建函式}

\texttt{**} 是次方，兩個星號連在一起：\texttt{2 ** 10} 就是 2 的 10 次方。另外有兩個
內建函式常常跟算術一起出現：\texttt{abs()} 取絕對值（不管正負，只要大小），
\texttt{round()} 四捨五入。
\index{abs()}\index{round()}

\begin{example}{python}{(power-and-helpers.py) 次方、絕對值、四捨五入}
print(2 ** 10)
print(abs(-5))
print(round(3.14159, 2))
\end{example}

\iosample{}{1024 \\ 5 \\ 3.14 \\}

\texttt{round()} 的第二個引數是「要留幾位小數」，寫 \texttt{2} 就是留到小數點後兩位。
如果不寫第二個引數，就是四捨五入到整數。

\subsection*{運算順序，跟數學課一樣}

同一行裡出現好幾個運算子時，Python 的優先順序跟你數學課學的一致：先算次方，再算乘、除、
取商、取餘，最後才算加減；想改變順序就加括號。

\begin{example}{python}{(operator-order.py) 括號改變運算順序}
print(2 + 3 * 4)
print((2 + 3) * 4)
\end{example}

\iosample{}{14 \\ 20 \\}

第一行先算 \texttt{3 * 4} 得到 \texttt{12}，再加 \texttt{2}，所以是 \texttt{14}；
第二行括號強迫先算 \texttt{2 + 3}，再乘 \texttt{4}，變成 \texttt{20}。 不確定的時候就加
括號，多打一對括號不會讓程式變慢，但會讓下一個讀你程式碼的人（很可能是你自己）少猜一次。

\subsection*{常見錯誤：除數是 0}

數學上不能除以 0，Python 也一樣，而且它會直接停下來報錯。

\begin{example}{python}{(divide-by-zero.py) 拿 0 當除數}
seconds = 100

print(seconds // 0)
\end{example}

\begin{verbatim}
Traceback (most recent call last):
  File "divide-by-zero.py", line 3, in <module>
    print(seconds // 0)
          ~~~~~~~~^^~~
ZeroDivisionError: integer division or modulo by zero
\end{verbatim}

錯誤類型是 \texttt{ZeroDivisionError}，
\index{ZeroDivisionError}
訊息直翻是「整數除法或取餘數時除以了 0」。 這個錯誤通常不是你手打了一個
\texttt{0}，而是除數來自某個變數，而那個變數在某次執行時剛好變成了 \texttt{0}——例如
「總分除以人數」的程式，遇到一個人都沒有的情況。 這類「平常都好好的，特定資料才爆炸」
的問題，第 \ref{ch:files-exceptions} 章會教你怎麼接住它。

\begin{drills}[label=drills:arithmetic]
\item \texttt{17 // 5} 跟 \texttt{17 \% 5} 分別算出什麼？
\item 一共有 250 秒，想印成「幾分幾秒」。 分鐘要用哪個運算子算？ 秒數呢？
\item \texttt{print(9 / 3)} 印出來的是 \texttt{3} 還是 \texttt{3.0}？
\item \texttt{7 \% 2} 算出來是多少？ 那 \texttt{8 \% 2} 呢？ 你要怎麼用 \texttt{\%}
判斷一個數字是不是偶數？
\item \texttt{print(1 + 2 * 3 ** 2)} 會印出什麼？（提示：三個運算子，先算哪一個？）
\end{drills}

\section{\texttt{type()} 的更多用法}
\label{sec:type-more-uses}

目前你看過的 \texttt{type()} 都是拿一個變數放進括號裡查詢。其實 \texttt{type()}
接受的括號裡可以放任何一個值，不限於變數本身，包括一個運算式算出來的結果，甚至是另一個
函式呼叫的結果。這一節帶你看幾個 \texttt{type()} 更實用的用法，順便揭穿一個新手很容易
誤會的地方：數字做完運算之後，型別不一定跟你想的一樣。

\begin{code}{python}{Query the type of an expression directly}
type(expression)
\end{code}

\begin{example}{python}{(type-of-expressions.py) Query the type of an expression's result}
print(type(3 + 5))
print(type(6 / 2))
print(type(3 + 2.5))
print(type(int(float("3.14"))))
\end{example}

\iosample{}{<class 'int'> \\ <class 'float'> \\ <class 'float'> \\ <class 'int'>}

這四行呼叫依序看下去：

\begin{itemize}
\item 第一行 \texttt{3 + 5} 兩邊都是 \texttt{int}，結果 \texttt{8} 也是 \texttt{int}，這符合
直覺。
\item 第二行 \texttt{6 / 2} 用 \texttt{/} 做除法，雖然算出來的值剛好是整數 \texttt{3}，
但 \texttt{type()} 告訴你它其實是 \texttt{float}，也就是 \texttt{3.0}。這是 Python 的
固定規則：只要用 \texttt{/} 做除法，不管兩邊除不除得盡，結果永遠是 \texttt{float}，跟
數學上「6 除以 2 就是整數 3」的直覺不一樣，寫程式的時候要特別留意。
\item 第三行 \texttt{3 +
2.5} 是 \texttt{int} 加 \texttt{float}，Python 會自動把整數那一邊當成浮點數來算，所以
結果也是 \texttt{float}：只要加減乘除任何一邊出現 \texttt{float}，整個結果就會是
\texttt{float}。
\item 第四行示範函式可以一層包一層：最裡面 \texttt{float("3.14")} 先把字串
轉成 \texttt{3.14}，外面的 \texttt{int(...)} 再把 \texttt{3.14} 轉成整數 \texttt{3}，
最外層的 \texttt{type()} 查的是最終這個 \texttt{3} 的型別，結果是 \texttt{int}。
\end{itemize}

\keypoint{只要運算式裡任何一個地方出現 \texttt{float}（包括除法 \texttt{/} 本身自動
產生的 \texttt{float}），最後算出來的結果就會是 \texttt{float}，就算數值上剛好是整數
也一樣。這個規則在後面章節做數學運算、判斷結果型別時會一直用到，先在這裡建立印象。}

\stylepoint{在互動模式（第 \ref{ch:introduction} 章介紹過的、逐行輸入馬上看結果的那個環境）裡除錯時，不用
先把運算式存進一個變數才能查型別，直接把整段運算式包進 \texttt{type()} 就好，例如
\texttt{type(6 / 2)}，省下一個變數名，用完即丟。}

\begin{drills}[label=drills:type-more-uses]
\item \texttt{type(10 - 4)} 會印出什麼？\texttt{type(10 / 4)} 呢？
\item \texttt{type(2 * 3.0)} 會印出什麼？為什麼？
\item \texttt{type(str(3) + str(5))} 會印出什麼？（提示：先想清楚 \texttt{str(3) +
str(5)} 這一步算出來的是什麼。）
\end{drills}

\section{同一個符號，不同型別不同意思}
\label{sec:type-operators}

型別不只影響一個值「看起來」是什麼，也影響某個符號用在它身上的時候，實際會發生什麼事。
最明顯的例子就是加號 \texttt{+}。

\begin{example}{python}{(plus-behavior.py) \texttt{+} used on numbers and strings}
print(3 + 5)
print("Py" + "thon")
\end{example}

\iosample{}{8 \\ Python}

同一個 \texttt{+}，用在兩個數字上是「相加」，算出 \texttt{8}；用在兩個字串上卻是
「串接」，把兩段文字直接接在一起，變成 \texttt{"Python"}，完全不是在做數學加法。
\index{串接}
這是
Python 針對不同型別、讓同一個符號做不同事情的刻意設計（正式名稱之後會學到，叫「運算子
多載」，這裡你只需要記得「同一個符號，看型別決定意思」這個現象就好）。

那如果你把一個數字跟一個字串直接用 \texttt{+} 接在一起呢？

\begin{example}{python}{(plus-error.py) Adding a number to a string causes trouble}
count = 3
print("There are currently " + count + " people")
\end{example}

執行後，程式會停下來，印出：

\begin{verbatim}
Traceback (most recent call last):
  File "plus-error.py", line 2, in <module>
    print("There are currently " + count + " people")
          ~~~~~~~~~~~~~~~~~~~~~~~^~~~~~~
TypeError: can only concatenate str (not "int") to str
\end{verbatim}

\texttt{TypeError} 代表 Python 看懂你的程式碼在講什麼，只是型別不對。
\index{TypeError}
它知道你想把
字串跟 \texttt{count} 接在一起，但 \texttt{count} 的型別是 \texttt{int}，跟 \texttt{str}
不同，\texttt{+} 用在字串上只認得另一個字串，不會自動幫你把數字轉成文字。這個問題怎麼
解決？下一節就要處理它。

\texttt{+} 不是唯一一個會這樣變臉的符號，\texttt{*} 也有一樣的故事。

\begin{example}{python}{(star-behavior.py) \texttt{*} used on numbers and strings}
print(4 * 3)
print("ab" * 3)
\end{example}

\iosample{}{12 \\ ababab}

\texttt{4 * 3} 是「相乘」，算出 \texttt{12}；\texttt{"ab" * 3} 用在字串上卻變成
「重複」，把 \texttt{"ab"} 這段文字連續接三次，變成 \texttt{"ababab"}。這再次印證：
同一個符號實際的意思，會隨著旁邊值的型別而改變，判斷規則要看型別，不能只看符號本身
長什麼樣子。

\begin{drills}[label=drills:type-family]
\item \texttt{type(2.5)} 會印出什麼？\texttt{type("2.5")} 呢？兩者一樣嗎？
\item 下面兩行分別會印出什麼？如果其中一行會出錯，錯誤類型是什麼？
\begin{verbatim}
print("age: " + 20)
print("age: " + "20")
\end{verbatim}
\item \texttt{"ha" * 0} 會印出什麼？\texttt{"ha" * -1} 呢？
\end{drills}

\section{型別轉換：把值換成你要的種類}
\label{sec:type-conversion}

上一節留下一個問題：想把數字接進字串裡，直接用 \texttt{+} 行不通。這一節要介紹的四個
函式，可以把一個值從一種型別，轉換成另一種型別。
\index{型別轉換}

\begin{code}{python}{Type conversion functions}
str(value)      # 轉成字串
int(value)      # 轉成整數
float(value)    # 轉成浮點數
bool(value)     # 轉成布林值
\end{code}

先回頭解決上一節的 \texttt{TypeError}：只要把 \texttt{count} 用 \texttt{str()} 包起
來，先轉成字串，再用 \texttt{+} 接，就不會有型別不合的問題了。

\begin{example}{python}{(fix-concat.py) Fixing the type error with \texttt{str()}}
count = 3
print("There are currently " + str(count) + " people")
\end{example}

\iosample{}{There are currently 3 people}

\texttt{str(count)} 造出一個全新的字串物件 \texttt{"3"}，內容看起來跟數字 3 一樣，但
型別已經是 \texttt{str}，可以安心跟其他字串用 \texttt{+} 接在一起了。注意這裡的動作是
造一個新物件、把結果貼上去，\texttt{count} 本身這個物件仍然是原本那個整數 \texttt{3}，
型別轉換並沒有「原地」把它改掉，這跟前面「指定是造新物件、貼名字牌」的說法是同一件事。

反過來的情境也很常見：你手上已經有一段文字形式的資料，其實是數字的意思，但因為存的時候
是字串，沒辦法直接拿去算。

\begin{example}{python}{(text-to-number.py) Convert a numeric string back into a number}
age_text = "25"
age = int(age_text)
print(age + 1)

price_text = "39.9"
price = float(price_text)
print(price * 2)

print(bool(0))
print(bool(1))
\end{example}

\iosample{}{26 \\ 79.8 \\ False \\ True}

\texttt{age\_text} 存的是字串 \texttt{"25"}，看起來像數字，但型別是 \texttt{str}，不能
直接拿去做加法。\texttt{int(age\_text)} 把它轉成真正的整數 \texttt{25}，之後
\texttt{age + 1} 才能正常運算，印出 \texttt{26}。\texttt{float()} 的道理一樣，只是轉成
帶小數點的浮點數。

最後兩行順便看一下 \texttt{bool()}：\texttt{bool(0)} 是
\texttt{False}，\texttt{bool(1)} 是 \texttt{True}，你可以先記住這兩個最單純的例子，
更完整的規則（哪些值算「真」、哪些算「假」）留到第 \ref{ch:conditionals} 章條件判斷再深入。

型別轉換不是每次都會成功。如果你要求 Python 把一段「根本不像數字」的文字轉成整數，它
沒辦法憑空猜出你想要的數字：

\begin{example}{python}{(convert-fail.py) A failed conversion raises an error}
score_text = "hello"
score = int(score_text)
print(score)
\end{example}

執行後會看到：

\begin{verbatim}
Traceback (most recent call last):
  File "convert-fail.py", line 2, in <module>
    score = int(score_text)
ValueError: invalid literal for int() with base 10: 'hello'
\end{verbatim}

這次遇到的是 \texttt{ValueError}，跟前面的 \texttt{TypeError} 屬於不同種類的問題：
\texttt{"hello"} 本身型別就是 \texttt{str}，型別上完全沒問題，\texttt{int()} 也確實
可以接受字串當引數（引數就是放進括號裡的那個值）
\index{引數}
；問題出在這個字串的\textbf{內容}並非
一串看起來像數字的文字，Python 沒辦法把 \texttt{"hello"} 解讀成任何一個整數，抱怨的
重點是「值不對」，跟「型別不對」是兩回事。
\index{ValueError}

修法要回到問題的源頭：確認你要轉換的字串，
內容真的是看起來像數字的文字；如果它本來就是使用者手誤或資料本身有問題，轉換這一步
沒辦法自己生出一個合理答案。

\begin{drills}[label=drills:type-conversion]
\item \texttt{int("42")} 執行後會得到什麼？型別是什麼？
\item \texttt{str(3.14)} 執行後會得到什麼？型別是什麼？
\item 下面這行程式碼執行會出錯嗎？如果會，是哪一種錯誤，為什麼？
\begin{verbatim}
float("3.14 dollars")
\end{verbatim}
\end{drills}

\section{型別轉換的邊界情況}
\label{sec:conversion-edge-cases}

上一節看過型別轉換成功跟失敗的基本樣子。這一節整理幾個新手特別容易誤判的邊界情況：
\texttt{int()} 把小數轉成整數的時候，實際做的事跟「四捨五入」不一樣；還有 \texttt{bool()}
對字串的判斷方式，常常跟直覺完全相反。

先看 \texttt{int()} 把 \texttt{float} 轉成 \texttt{int} 的規則。很多人直覺會以為
\texttt{int()} 會幫你四捨五入，但實際上它做的是「無條件捨去小數部分」，而且不管正負號，
一律往 0 的方向捨去。

\begin{example}{python}{(int-truncation.py) How \texttt{int()} truncates decimals}
print(int(3.9))
print(int(3.1))
print(int(-3.9))
\end{example}

\iosample{}{3 \\ 3 \\ -3}

\texttt{int(3.9)} 結果是 \texttt{3}，四捨五入原本會得到的 \texttt{4} 並沒有出現；
\texttt{int(3.1)} 也是 \texttt{3}，這個看起來符合「捨去」的直覺。真正容易誤判的是第
三行：\texttt{int(-3.9)} 結果是 \texttt{-3}，很多人會直覺猜成 \texttt{-4}，猜錯的原因
就出在這裡。

\texttt{int()} 一律朝 0 的方向捨去小數部分，正數往下捨、負數往上捨，跟先
取絕對值再無條件捨去的效果不一樣。如果你真的需要四捨五入，那是另一個內建函式的工作，
這本書會在用得到的時候再介紹。

\keypoint{\texttt{int()} 對 \texttt{float} 做的事情永遠是捨去小數、朝 0 靠攏，跟四捨
五入是兩套完全不同的規則。這個規則對負數特別容易讓人猜錯方向，寫程式時如果剛好在處理
負數的小數轉換，值得多留意一秒。}

還有一個常見的踩雷點：如果字串裡本身帶著小數點，直接丟給 \texttt{int()} 會失敗，就算
內容看起來完全像一個合理的數字。

\begin{example}{python}{(decimal-string-fail.py) A string with a decimal point can't go straight into \texttt{int()}}
age_text = "25.0"
age = int(age_text)
print(age)
\end{example}

執行後會看到：

\begin{verbatim}
Traceback (most recent call last):
  File "decimal-string-fail.py", line 2, in <module>
    age = int(age_text)
ValueError: invalid literal for int() with base 10: '25.0'
\end{verbatim}

\texttt{"25.0"} 這段文字裡有一個小數點，\texttt{int()} 只認得「純數字」形式的字串
（像 \texttt{"25"}），只要多了小數點就會直接丟出 \texttt{ValueError}，不會自動幫你把
小數部分捨去。想把這種字串轉成整數，要先繞道 \texttt{float()}，讓 Python 先讀懂它是
一個帶小數的數字，再用 \texttt{int()} 把結果轉成整數，兩步分開做，Python 才處理得來：
\texttt{int(float(age\_text))} 就會得到 \texttt{25}。

最後看 \texttt{bool()}。上一節你看過 \texttt{bool(0)} 是 \texttt{False}、
\texttt{bool(1)} 是 \texttt{True}，這符合直覺。但把字串丟進 \texttt{bool()}，判斷規則
只看這段字串裡有沒有字元，完全不會去讀字元的內容講的是什麼。

\begin{example}{python}{(bool-conversions.py) Converting between \texttt{bool()} and numbers}
print(bool(""))
print(bool(" "))
print(bool("False"))
print(bool("0"))
print(int(True))
print(float(True))
\end{example}

\iosample{}{False \\ True \\ True \\ True \\ 1 \\ 1.0}

只有第一行 \texttt{bool("")} 是 \texttt{False}，因為這是一個完全沒有任何字元的空
字串。第二行 \texttt{bool(" ")} 是 \texttt{True}，因為裡面有一個空格字元，就算這個
字元用肉眼看幾乎是空的，對 Python 來說它仍然是一個字元，字串本身並不是空的。

第三行、
第四行最容易讓人誤會：\texttt{bool("False")} 跟 \texttt{bool("0")} 印出來都是
\texttt{True}。原因是 \texttt{bool()} 完全不會去讀字串裡寫的內容究竟是「False」還是
「0」，它只在乎這個字串是不是空字串，只要不是空字串，一律當成 \texttt{True}。想要判斷
一段文字內容代表的是真是假，得先把文字轉成正確的型別，或是自己寫規則去比對內容，不能
只靠 \texttt{bool()} 硬轉。

最後兩行反過來看：\texttt{bool} 轉成數字的時候，
\texttt{True} 當成 \texttt{1}、\texttt{False} 當成 \texttt{0}，所以 \texttt{int(True)}
是 \texttt{1}，\texttt{float(True)} 是 \texttt{1.0}。這跟字串那三行的規則完全不一樣：
字串轉 \texttt{bool()} 看的是「是不是空的」，\texttt{bool} 轉數字看的是「是
\texttt{True} 還是 \texttt{False}」，兩個方向不能用同一套邏輯套用。

\keypoint{\texttt{bool()} 用在字串上，判斷依據只看這個字串是不是空的，字串裡寫的文字
內容完全不影響結果。\texttt{bool("False")}、\texttt{bool("0")} 這種看起來像是「假」的
字串，因為都不是空字串，轉換結果都會是 \texttt{True}，這是新手最容易誤踩的一個地雷。}

\begin{drills}[label=drills:conversion-edge-cases]
\item \texttt{int(-2.1)} 跟 \texttt{int(2.9)} 分別會得到什麼？
\item \texttt{int("3.0")} 執行會不會出錯？如果會，要怎麼修？
\item \texttt{bool("0.0")} 會印出 \texttt{True} 還是 \texttt{False}？為什麼？
\end{drills}

\section*{練習解答}
\addcontentsline{toc}{section}{練習解答}

\paragraph{第~\ref{sec:variable-naming}~節}
\begin{enumerate}
\item \texttt{my\_name} 合法；\texttt{2apples} 不合法（以數字開頭）；
\texttt{total\_price} 合法；\texttt{if} 不合法（是保留字）；\texttt{User1} 合法（數字
不在開頭就可以）。
\item \texttt{NameError: name 'favorite\_colour' is not defined}（可能還會貼心猜一句
「是不是想打 \texttt{favorite\_color}」），因為 Python 找不到叫這個名字的變數，兩處
拼法對不起來。
\item \texttt{SyntaxError}，因為變數名不能以數字開頭，Python 光是讀這一行就看不懂，
連執行都還沒開始，跟「找不到變數」的 \texttt{NameError} 是不同階段的問題。
\end{enumerate}

\paragraph{第~\ref{sec:naming-practices}~節}
\begin{enumerate}
\item \texttt{print = 5} 之後再呼叫 \texttt{print("hi")} 會出錯，因為這時候
\texttt{print} 這個名字對應的已經是數字 5，原本的函式被蓋掉了，數字不能像函式一樣被
呼叫，錯誤類型是 \texttt{TypeError}（\texttt{'int' object is not callable}）。
\item \texttt{is\_passed} 比較適合，因為布林值變數習慣用 \texttt{is\_} 開頭命名，光看
名字就知道這是一個對／錯的答案；\texttt{result} 太籠統，看不出裝的究竟是分數、名次還是
對錯。
\item 拿掉第一行之後，\texttt{type} 這個名字沒有被蓋掉，第五行的 \texttt{type(x)}
呼叫的就是原本查詢型別的內建函式，程式會正常執行，印出 \texttt{<class 'int'>}，不會
出錯。
\end{enumerate}

\paragraph{第~\ref{sec:name-tags}~節}
\begin{enumerate}
\item \texttt{x}、\texttt{y}、\texttt{z} 一開始都指向同一個字串 \texttt{"start"}。
\texttt{x = "changed"} 只是把 \texttt{x} 改貼到一個新的物件上，\texttt{y}、\texttt{z}
不受影響，所以 \texttt{print(x, y, z)} 印出 \texttt{changed start start}。
\item 會不一樣，印出的是 \texttt{start changed start}：改的名字不同，被改貼的名字牌就
不同，跟哪一個名字被重新賦值有關，跟賦值的順序無關。
\end{enumerate}

\paragraph{第~\ref{sec:type-on-value}~節}
\begin{enumerate}
\item \texttt{b} 會印出 \texttt{10}，因為 \texttt{b = a} 時 \texttt{b} 指向跟 \texttt{a}
一樣的物件（值 10），之後 \texttt{a = 20} 只是讓 \texttt{a} 改指到新的物件 20，
\texttt{b} 完全沒被影響。
\item 會不一樣，\texttt{print(a)} 會印出 \texttt{20}，因為 \texttt{a} 這張名字牌已經
改貼到新的物件 20 上面了。
\item \texttt{<class 'int'>}，雖然 \texttt{x} 一開始指向字串 \texttt{"3"}，但第二行
\texttt{x = 3} 讓 \texttt{x} 改指向整數 3，型別長在物件身上，跟著目前指的那個物件走。
\end{enumerate}

\paragraph{第~\ref{sec:arithmetic}~節}
\begin{enumerate}
\item \texttt{17 // 5} 是 \texttt{3}（17 裡面裝得下三個完整的 5），\texttt{17 \% 5} 是
\texttt{2}（裝完三個 5 之後剩下的）。
\item 分鐘用 \texttt{250 // 60}，算出 \texttt{4}；秒數用 \texttt{250 \% 60}，算出
\texttt{10}。也就是 4 分 10 秒。
\item 印出 \texttt{3.0}。只要用 \texttt{/} 做除法，就算除得盡，結果一樣帶小數點。
\item \texttt{7 \% 2} 是 \texttt{1}，\texttt{8 \% 2} 是 \texttt{0}。判斷偶數的寫法是
\texttt{number \% 2 == 0}：餘數為 \texttt{0} 就代表除得盡，是偶數。
\item 印出 \texttt{19}。先算次方 \texttt{3 ** 2} 得到 \texttt{9}，再算乘法
\texttt{2 * 9} 得到 \texttt{18}，最後才加 \texttt{1}。
\end{enumerate}

\paragraph{第~\ref{sec:type-more-uses}~節}
\begin{enumerate}
\item \texttt{type(10 - 4)} 印出 \texttt{<class 'int'>}（兩邊都是整數，結果也是
整數）；\texttt{type(10 / 4)} 印出 \texttt{<class 'float'>}（用 \texttt{/} 做除法，
結果一律是 \texttt{float}）。
\item 印出 \texttt{<class 'float'>}，因為 \texttt{2 * 3.0} 裡有一邊是 \texttt{float}
（\texttt{3.0}），只要運算式裡任何一邊出現 \texttt{float}，結果就會是 \texttt{float}。
\item 印出 \texttt{<class 'str'>}：\texttt{str(3)} 是字串 \texttt{"3"}，\texttt{str(5)}
是字串 \texttt{"5"}，兩個字串用 \texttt{+} 接起來是串接，結果是字串 \texttt{"35"}，
型別是 \texttt{str}。
\end{enumerate}

\paragraph{第~\ref{sec:type-family}~節}
\begin{enumerate}
\item \texttt{type(2.5)} 印出 \texttt{<class 'float'>}；\texttt{type("2.5")} 印出
\texttt{<class 'str'>}，兩者不一樣：被引號包住的 \texttt{"2.5"} 屬於文字，就算內容
看起來像數字，型別仍然是 \texttt{str}。
\item 第一行 \texttt{print("age: " + 20)} 會出錯，錯誤類型是 \texttt{TypeError}（字串
不能直接跟 \texttt{int} 用 \texttt{+} 接在一起）；第二行 \texttt{print("age: " +
"20")} 會正常印出 \texttt{age: 20}，因為兩邊都是字串，\texttt{+} 做的是串接。
\item \texttt{"ha" * 0} 跟 \texttt{"ha" * -1} 都印出一個空字串 \texttt{""}：重複零次
或負數次，效果都是「什麼都不留」，Python 不會因為次數是負的而報錯。
\item 不一樣，\texttt{type(0)} 是 \texttt{<class 'int'>}，\texttt{type(0.0)} 是
\texttt{<class 'float'>}：就算數值上兩者相等，有沒有寫小數點會決定字面值的型別。
\item \texttt{type(True)} 印出 \texttt{<class 'bool'>}。
\item 印出 \texttt{1000000}，底線只是幫人眼閱讀用的排版符號，不會出現在印出來的內容
裡，Python 讀到的跟直接寫 \texttt{1000000} 是同一個數字。
\end{enumerate}

\paragraph{第~\ref{sec:type-conversion}~節}
\begin{enumerate}
\item 得到整數 \texttt{42}，型別是 \texttt{<class 'int'>}。
\item 得到字串 \texttt{"3.14"}，型別是 \texttt{<class 'str'>}。
\item 會出錯，錯誤類型是 \texttt{ValueError}，因為 \texttt{"3.14 dollars"} 這段文字裡除了
數字跟小數點，還混了「dollars」這幾個字跟一個空格，不是一段乾淨的、看起來像浮點數的文字，
\texttt{float()} 沒辦法把它解讀成任何一個數字。
\end{enumerate}

\paragraph{第~\ref{sec:conversion-edge-cases}~節}
\begin{enumerate}
\item \texttt{int(-2.1)} 得到 \texttt{-2}，\texttt{int(2.9)} 得到 \texttt{2}：
\texttt{int()} 一律朝 0 的方向捨去小數，跟四捨五入是兩套不同的規則。
\item 會出錯，因為 \texttt{"3.0"} 裡帶著小數點，\texttt{int()} 不接受字串裡出現小數點，
會丟出 \texttt{ValueError}；要修的話先用 \texttt{float("3.0")} 轉成浮點數，再用
\texttt{int()} 轉成整數，也就是 \texttt{int(float("3.0"))}。
\item 印出 \texttt{True}，因為 \texttt{"0.0"} 雖然內容看起來像數字 0，但它是一個非空
字串，\texttt{bool()} 用在字串上只看是不是空字串，不會去讀內容代表的數值，所以結果是
\texttt{True}。
\end{enumerate}

\chapter{字串}
\label{ch:strings}

上一章你已經知道怎麼幫一個值貼上名字牌，這一章要正式認識目前為止你貼過最有趣的一種
值，字串。你可以把一串字想像成電影院裡一整排緊靠在一起、依序編號的座位，每個座位坐著
一個字元，你可以直接說「我要第幾號座位」跳過去看，也可以說「我要看第幾號到第幾號之間
的座位」一次看一整段，這就是這一章前半段要學的\textbf{索引}跟\textbf{切片}。跟你平常
「從 1 開始數」的習慣不一樣，Python 的座位編號是從 0 開始算的，這是幾乎每個人第一次接觸
字串都會愣一下的地方，我們會花點篇幅把這個「為什麼」講清楚，讓你理解背後的邏輯，這樣才
記得住，死背數字很快就會忘記。

接著，光是能「看」一段字還不夠，你常常需要把一長串文字拆開、把好幾段文字接起來、
把多餘的空白清掉、把某個詞換成另一個詞，這是這一章中段的重點：\textbf{字串方法}，也是你
第一次看到「在一個值後面加一個點、呼叫它自己內建的動作」這種新語法。認識完最基本的幾個
方法之後，我們還會多花一點篇幅，補齊幾個你之後在真實情境裡幾乎一定會用到的方法：判斷
大小寫、判斷一段文字的開頭或結尾、在一長串文字裡找出某段內容出現在哪裡、出現了幾次。

最後，光會拆解、清理文字還不夠，你常常需要反過來，把幾個變數的值組回一句話說給別人聽，
像「你好，小明，你今年 15 歲」這種句子，這就是這一章最後要學的 \textbf{f-string}，把
變數的值直接嵌進一段文字裡，還會多學一點數字的格式化寫法，讓輸出的金額、人數看起來更像
正式報表。

這一章的順序背後有一個簡單的邏輯：先學會「看」一段既有的字串，再學會「拆解、
清理」它，最後學會「組回」一段新的字串，從最具體的操作一路推進到最貼近你實際會用到的
目的。中間會有一個小地方值得停下來想一想：字串跟上一章教的名字牌模型有一點微妙的互動，
我們會在講到「字串不能被原地修改」的時候簡短回頭呼應一次。

\section{索引：抓出字串裡的某一個字元}
\label{sec:string-indexing}

還記得第 \ref{ch:variables} 章的 \texttt{str} 型別嗎？字串就是這種型別的值。從現在開始，我們要正式認識
這種資料能做的事。字串本質上是一串排好順序的字元，就像前面說的電影院座位：每個字元都坐在一個固定的位置
上，你可以用中括號指名要看第幾號座位上的字元，這個動作叫做「索引」（indexing）。
\index{索引}

跟你
平常數東西「從 1 開始」的習慣不一樣，Python 的座位編號是從 0 開始算的：第一個字元是位置
0，第二個字元是位置 1，以此類推。這聽起來有點反直覺，但你很快就會習慣，之後你會發現，
Python 裡幾乎所有「數東西」的地方都是從 0 開始算，這裡先建立起這個習慣，後面會輕鬆很多。

除了從前面數，你也可以從後面數：用負數索引，\texttt{-1} 永遠代表「最後一個字元」，
\texttt{-2} 是倒數第二個，不用先算出字串總共有幾個字元才能拿到最後一個。

\begin{code}{python}{用中括號取出單一字元}
s[i]
\end{code}

\begin{example}{python}{(index-letters.py) 取出水果名字的某個字元}
fruit = "banana"
print(fruit[0])   # 第一個字元
print(fruit[1])   # 第二個字元
print(fruit[-1])  # 最後一個字元，從尾端數
print(fruit[-2])  # 倒數第二個字元
\end{example}

\iosample{}{b \\ a \\ a \\ n \\}

\keypoint{Python 的索引從 0 開始，跟你數學課、日常生活「第一個就是 1 號」的習慣不
一樣，是新手最容易在這裡卡住的地方。記住一個簡單的換算方式：字串的第 $n$ 個字元，索引
是 $n - 1$。}

那如果指定了一個字串裡根本不存在的位置呢？\texttt{fruit = "banana"} 總共有 6 個字元，
有效的索引是 0 到 5，很多人會直覺覺得「總共 6 個字元，那 \texttt{fruit[6]} 應該是最後
一個」，但其實 \texttt{fruit[5]} 才是最後一個字元，\texttt{fruit[6]} 已經超出範圍了。

\begin{example}{python}{(index-error.py) 索引超出範圍}
fruit = "banana"
print(len(fruit))   # 6，這個字串總共有 6 個字元
print(fruit[6])      # 錯誤示範：把字元總數當成索引來用，超出範圍
\end{example}

存檔執行後，畫面上會出現：

\begin{verbatim}
6
Traceback (most recent call last):
  File "index-error.py", line 3, in <module>
    print(fruit[6])
          ~~~~~^^^
IndexError: string index out of range
\end{verbatim}

\texttt{IndexError: string index out of range} 是說「你指定的索引超出了這個字串的
範圍」。\index{IndexError}這是一個經典的「差一」錯誤（off-by-one）：\index{差一錯誤}
\texttt{len(fruit)} 算出來的是字元
\textit{總數}，但最後一個字元的\textit{索引}永遠比總數少 1，因為索引是從 0 開始算的。
之後你會在很多地方一再碰到這個陷阱，記住這個換算關係就能避開它。

負索引一樣有邊界，只是方向反過來數：\texttt{fruit = "banana"} 只有 6 個字元，最前面
對應到的負索引是 \texttt{-6}，如果寫 \texttt{fruit[-7]}，一樣會得到
\texttt{IndexError}，道理跟正索引超出範圍完全一樣。

\begin{drills}[label=drills:string-indexing]
\item 給定 \texttt{word = "orange"}，\texttt{word[2]} 的值是什麼？
\item 給定 \texttt{word = "orange"}，\texttt{word[-3]} 的值是什麼？
\item \texttt{word = "orange"} 總共有 6 個字元。如果寫 \texttt{word[6]}，執行時會發生
什麼事？為什麼？
\end{drills}

\section{切片：一次抓出一段字元}
\label{sec:string-slicing}

只能一次抓一個字元有點不夠用，很多時候你想要的是一整段，例如「前三個字」或「從某個
位置到結尾」。這時候一樣用中括號，只是裡面多放一個冒號，變成「從哪裡開始、到哪裡結束」，
這個動作叫做「切片」（slicing）。
\index{切片}

切片有一個和索引一樣容易讓人卡住的地方：\texttt{start}
指定的位置\textbf{會}被包含進來，但 \texttt{end} 指定的位置\textbf{不會}，也就是常說的
「含頭不含尾」。你可以把它想成：冒號前後的數字指的是座位跟座位之間的「縫隙」，不是座位
本身，\texttt{s[0:3]} 意思是「從縫隙 0 切到縫隙 3」，剛好切出 3 個字元。

\begin{code}{python}{用中括號取出一段字元}
s[start:end]
\end{code}

\begin{example}{python}{(slice-fruit.py) 切出水果名字的一段}
fruit = "banana"
print(fruit[0:3])   # 從位置 0 到位置 3（不含），也就是前 3 個字元
print(fruit[3:6])   # 從位置 3 到位置 6（不含）
print(fruit[:3])     # 省略 start，代表「從頭開始」
print(fruit[3:])     # 省略 end，代表「到最後為止」
print(fruit[:])       # start 跟 end 都省略，等於整個字串
\end{example}

\iosample{}{ban \\ ana \\ ban \\ ana \\ banana \\}

\stylepoint{省略 \texttt{start} 或 \texttt{end}，是很常見、也很值得養成的寫法。
\texttt{s[:n]} 讀起來就是「前 $n$ 個」、\texttt{s[n:]} 讀起來就是「從第 $n$ 個開始的
剩下部分」，比每次都手動算出完整的起訖位置更不容易出錯，也更好讀。}

切片還有一個和索引很不一樣的地方：如果你指定的位置超出字串範圍，索引會直接報錯，但切片
不會，它會很寬容地幫你「夾」到最接近的有效範圍，例如 \texttt{fruit[3:100]} 不會出錯，
只會切到字串結尾為止。除了「從哪裡切到哪裡」，切片其實還可以指定「每隔幾個字元取一次」，
甚至反過來從尾巴讀到頭，這是下一節要正式介紹的進階用法。

\begin{drills}[label=drills:string-slicing]
\item 給定 \texttt{text = "computer"}，\texttt{text[2:5]} 的值是什麼？
\item 給定 \texttt{text = "computer"}，\texttt{text[:4]} 的值是什麼？
\item 給定 \texttt{text = "computer"}，\texttt{text[-3:]} 的值是什麼？
\end{drills}

\section{切片進階：用 step 指定間隔，反轉整個字串}
\label{sec:string-slicing-step}

上一節的切片只寫了「從哪裡到哪裡」兩個數字，其實中括號裡還可以放第三個數字，用另一個
冒號隔開，變成 \texttt{s[start:end:step]}。第三個數字叫做 step，意思是「每隔幾個位置
取一次」，預設值是 1，也就是「每個都取」，這也是為什麼前面兩節你完全沒感覺到它的存在。
把 step 改成 2，就會變成「每兩個取一個」，也就是只留下位置 0、2、4 等等上的字元。

\begin{code}{python}{用第三個數字指定間隔}
s[start:end:step]
\end{code}

\begin{example}{python}{(every-other-letter.py) 每隔一個字元取一次}
fruit = "banana"
print(fruit[::2])     # 從頭到尾，每隔一個字元取一次
print(fruit[1:5:2])   # 從位置 1 到位置 5（不含），每隔一個字元取一次
\end{example}

\iosample{}{bnn \\ aa \\}

start 跟 end 一樣可以照上一節學過的規則省略，\texttt{fruit[::2]} 省略了 start 跟
end，只指定 step，意思是「整個字串，每隔一個取一次」。

step 還可以是負數，負數的 step 代表「反過來走」，從後面往前一個一個（或每隔幾個）取。
把 start、end 都省略，只留下 step 是 \texttt{-1}，就會把整個字串反過來走一遍，剛好
組成一個反轉後的新字串，這是一個很常見、也很值得記住的寫法。

\begin{example}{python}{(reverse-string.py) 用 step 為 -1 反轉字串}
fruit = "banana"
print(fruit[::-1])       # 整個字串反轉過來
print(fruit[5:1:-1])     # 從位置 5 往回走到位置 2（不含位置 1）
\end{example}

\iosample{}{ananab \\ anan \\}

用負數 step 時，start 跟 end 的「含頭不含尾」規則還在，只是方向反過來：
\texttt{fruit[5:1:-1]} 是「從位置 5 開始，往回走到位置 1 之前為止（不含位置 1）」，
所以拿到的是位置 5、4、3、2 上的字元，也就是 \texttt{anan}。

\keypoint{\texttt{s[::-1]} 是 Python 裡反轉字串最常見的寫法。很多人會先入為主地覺得
應該有一個專門叫 \texttt{reverse()} 的字串方法，但字串其實沒有這個方法，因為字串是
不可變的，沒有「原地反轉」這回事。想要反轉字串，切片配合負數 step 就是最直接的做法。}

\begin{drills}[label=drills:string-slicing-step]
\item 給定 \texttt{text = "computer"}，\texttt{text[::2]} 的值是什麼？
\item 給定 \texttt{text = "computer"}，\texttt{text[::-1]} 的值是什麼？
\item 給定 \texttt{text = "computer"}，\texttt{text[1:7:2]} 的值是什麼？
\end{drills}

\section{字串不可變：為什麼不能直接改掉一個字元}
\label{sec:string-immutable}

目前為止你都是在「讀」字串裡的內容，沒有真的去「改」它。如果你想把一個字串裡的某個字元
換掉，你可能會想，既然可以用中括號取出某個位置的字元，那應該也能用中括號把某個位置換成
別的字元吧？可惜這個直覺在字串上行不通。

\begin{example}{python}{(broken-name.py) 想直接修改字串的其中一個字元}
name = "Jenny"
name[0] = "P"
\end{example}

存檔執行後，畫面上會出現：

\begin{verbatim}
Traceback (most recent call last):
  File "broken-name.py", line 2, in <module>
    name[0] = "P"
    ~~~~^^^
TypeError: 'str' object does not support item assignment
\end{verbatim}

這個錯誤訊息在告訴你：字串這種物件，本來就不支援「指定某一項的值」這個動作，跟你程式
碼有沒有打錯字完全無關。這是 Python 字串一個很基本的特性：\textbf{字串一旦造出來，內容
就不能被原地修改}，這個特性叫做「不可變」（immutable）。
\index{不可變}想「修改」一個字串，正確的做法
是用切片組出一個全新的字串，再把它整個指定給原本的變數：

\begin{example}{python}{(fixed-name.py) 用切片重組出一個新字串}
name = "Jenny"
name = "P" + name[1:]   # 把第一個字元換成 P，保留剩下的部分，組成一個新字串
print(name)
\end{example}

\iosample{}{Penny \\}

為什麼 Python 要把字串設計成不可變的？因為程式裡常常有好幾個名字牌一起指到同一個
字串物件。如果字串可以被原地修改，某個名字牌偷偷把內容改掉，其他指著同一個物件的名字牌
會在完全不知情的狀況下被連帶影響，程式的行為會變得很難預測。不可變性保證了「只要沒有
重新賦值，這個名字牌指到的內容就一定不會變」，讀程式碼時可以放心很多。

\keypoint{還記得上一章的名字牌模型嗎？\texttt{x = "hi"; y = x} 之後改 \texttt{x} 不會
影響 \texttt{y}，看起來很像「複製了一份」，但事實上兩張名字牌一樣貼在同一個字串物件上，
只是因為字串不可變，沒有任何辦法能「原地」改掉那個物件的內容，所以你永遠看不出兩張名字
牌共用同一個物件的差異。字串這裡表面上像是複製，其實是不可變性造成的巧合。}

\section{方法呼叫語法與 \texttt{strip()}：清除頭尾多餘的空白}
\label{sec:strip-method}

前面你用中括號「讀」字串的內容，這一節開始，你要學會主動「處理」字串，第一次看到一種新
語法：在一個值後面加一個點，呼叫它自己內建的動作，叫做「方法」（method）。
\index{方法}

這跟你之前看
過的函式呼叫不一樣：函式呼叫是 \texttt{f(x)}，把值放進括號裡交給函式；方法呼叫是
\texttt{x.f()}，值本身站到前面，後面接一個點、再接方法名稱，意思是「叫這個值自己去做
這件事」。

字串身上內建了很多這種方法，這一節先從最直覺的一個開始：\texttt{strip()}，
清除字串頭尾多餘的空白（包含空格、換行）。
\index{strip()}它特別實用的場景是處理使用者的輸入，你請人
輸入名字，對方常常會不小心在前後多打一個空格，或者按 Enter 時多留了一個換行，
\texttt{strip()} 就是拿來收拾這些小雜訊的。

\begin{code}{python}{呼叫字串自己的方法}
s.strip()
\end{code}

\texttt{input()} 是 Python 內建的一個動作，會在螢幕上顯示你給的提示文字、暫停程式、等你
打字並按下 Enter，然後把你打的內容變成一個字串交給程式。
\index{input()}

\begin{example}{python}{(clean-input.py) 清除使用者輸入頭尾的空白}
raw = input("Enter your name: ")
name = raw.strip()               # 清除頭尾多餘的空白
print("Hello, " + name + "!")
\end{example}

下面用 \texttt{␣} 代表一個空白字元，方便你看清楚使用者實際打了什麼：

\iosample{␣␣Alice␣␣ \\}{Enter your name: Hello, Alice! \\}

\texttt{strip()} 呼叫完之後，回傳的是一個\textbf{新的}字串，原本的 \texttt{raw} 完全
沒有被改動，這跟上一節「字串不可變」是同一件事的另一種說法：字串方法從來不會原地修改
呼叫它的那個字串，只會算出一個新字串回傳給你。這代表如果你忘了把回傳值重新指定回變數，
原本的字串會維持原樣，不會被「就地清乾淨」：

\begin{example}{python}{(strip-mistake.py) 忘記把結果重新指定回變數}
raw = "   Alice   "
raw.strip()                  # 算出一個新字串，但沒有存到任何地方
print("[" + raw + "]")       # raw 本身完全沒有被改動，空白還在
\end{example}

\iosample{}{[␣␣␣Alice␣␣␣] \\}

\keypoint{呼叫一個字串方法，永遠不會改變原本的字串本身，只會回傳一個新字串。如果你想
保留處理過的結果，一定要把回傳值重新指定回一個變數，例如
\texttt{raw = raw.strip()}，只呼叫 \texttt{raw.strip()} 卻不接回傳值，這行程式碼幾乎
等於什麼都沒做。}

\texttt{strip()} 預設清除的是空白字元，但其實它也能接受一個引數，指定要清除哪些字元。
傳進去的字串裡列出的每一個字元，只要出現在頭尾，都會一路被剝掉，直到剝到一個不在清單裡
的字元為止。

\begin{example}{python}{(strip-symbols.py) 清除頭尾指定的符號}
notice = "###IMPORTANT NOTICE###"
print(notice.strip("#"))
\end{example}

\iosample{}{IMPORTANT NOTICE \\}

\texttt{strip("\#")} 只會處理\textit{頭尾}的 \texttt{"\#"}，字串中間的內容完全不會被
動到，這一點跟後面要學的 \texttt{replace()}（會處理字串裡任何位置符合的內容）不一樣，
先在這裡留一個印象，之後比較會更清楚。

如果你只想清掉一邊，還有兩個近親：\texttt{lstrip()} 只清左邊（開頭），
\texttt{rstrip()} 只清右邊（結尾），用法跟 \texttt{strip()} 一模一樣，也同樣可以帶一個
參數指定要清哪些字元。
\index{lstrip()／rstrip()}

\begin{example}{python}{(strip-one-side.py) 只清掉一邊的空白}
raw = "   ID-2024   "

print("[" + raw.strip() + "]")
print("[" + raw.lstrip() + "]")
print("[" + raw.rstrip() + "]")
\end{example}

\iosample{}{[ID-2024] \\{} [ID-2024   ] \\{} [   ID-2024] \\}

三行都用中括號把結果框起來，才看得出空白到底被清掉了沒有。 實務上
\texttt{strip()} 用得最多，但遇到「行首的縮排要保留、行尾的換行要清掉」這類情況，
\texttt{rstrip()} 就派得上用場了。

\begin{drills}[label=drills:strip-method]
\item 預測輸出：
\begin{verbatim}
msg = "  hi  ".strip()
print("[" + msg + "]")
\end{verbatim}
\item 承接上一題的錯誤示範：如果把第一行改成 \texttt{msg = "  hi  "}、第二行改成
\texttt{msg.strip()}（沒有重新指定回 \texttt{msg}），第三行還是
\texttt{print("[" + msg + "]")}，這次輸出的中括號裡面，頭尾還會有空白嗎？為什麼？
\item 預測輸出：\texttt{"xxpythonxx".strip("x")} 的結果是什麼？
\item 預測輸出：\texttt{"xxpythonxx".lstrip("x")} 的結果是什麼？跟上一題差在哪裡？
\end{drills}

\section{\texttt{lower()} 與 \texttt{upper()}：轉換大小寫}
\label{sec:case-methods}

英文字母有大小寫之分，但很多時候你並不在乎使用者輸入的到底是大寫還是小寫，例如比對
使用者輸入的指令、比對信箱帳號，通常都希望「Alice」跟「alice」被當成一樣的東西。
\texttt{lower()} 把字串裡所有的英文字母轉成小寫，\index{lower()}
\texttt{upper()} 反過來全部轉成大寫，
\index{upper()}
非英文字母（數字、標點、中文字）不受影響，會原封不動保留。

\begin{code}{python}{轉換大小寫}
s.lower()
s.upper()
\end{code}

\begin{example}{python}{(case-convert.py) 轉換一段英文的大小寫}
word = "Python"
print(word.lower())
print(word.upper())
\end{example}

\iosample{}{python \\ PYTHON \\}

跟前面學過的所有字串方法一樣，\texttt{lower()}、\texttt{upper()} 都會回傳一個新字串，
原本的字串不會被改動，這是字串不可變的老規矩，之後每個字串方法都適用，不會再重複提醒。

大小寫轉換最常見的用途，是在比較兩段文字「內容是不是一樣」的時候，先統一成同一種大小寫，
避免只因為大小寫不同就被判斷成不相等。

\begin{example}{python}{(case-insensitive-check.py) 忽略大小寫比對兩個字串}
name1 = "Alice"
name2 = "alice"
print(name1 == name2)                  # 直接比較，大小寫不同就算不相等
print(name1.lower() == name2.lower())  # 兩邊都轉成小寫後再比較
\end{example}

\iosample{}{False \\ True \\}

\keypoint{\texttt{==} 是逐字元比對的，\texttt{"Alice"} 跟 \texttt{"alice"} 因為第一個
字母大小寫不同，直接用 \texttt{==} 比較會是 \texttt{False}。如果你想要的是「內容一樣就
算相同，不管大小寫」，養成習慣先各自呼叫 \texttt{.lower()}（或都呼叫 \texttt{.upper()}）
再比較，這是判斷使用者輸入時很常用的技巧。}

\begin{drills}[label=drills:case-methods]
\item 預測輸出：\texttt{"Hello World".lower()} 的結果是什麼？
\item 預測輸出：\texttt{"Hello World".upper()} 的結果是什麼？
\item 預測輸出：\texttt{"PYTHON".lower() == "python"} 的結果是什麼？
\end{drills}

\section{\texttt{startswith()} 與 \texttt{endswith()}：判斷開頭與結尾}
\label{sec:startswith-endswith}

除了判斷兩段文字是不是完全一樣，你常常只想知道「這段文字是不是以某個特定的內容開頭
（或結尾）」，例如判斷一個檔名是不是圖片檔、判斷一個網址是不是用安全的 \texttt{https}
開頭。\texttt{startswith()} 檢查字串是不是以指定的內容開頭，\index{startswith()}
\texttt{endswith()} 檢查
是不是以指定的內容結尾，兩者都回傳 \texttt{bool}。
\index{endswith()}

\begin{code}{python}{判斷開頭與結尾}
s.startswith(text)
s.endswith(text)
\end{code}

\begin{example}{python}{(check-filename.py) 判斷檔名的副檔名}
filename = "report.pdf"
print(filename.endswith(".pdf"))
print(filename.endswith(".docx"))
\end{example}

\iosample{}{True \\ False \\}

\begin{example}{python}{(check-url.py) 判斷網址開頭}
url = "https://example.com"
print(url.startswith("https"))
print(url.startswith("http://"))
\end{example}

\iosample{}{True \\ False \\}

第二個例子容易讓人誤會：\texttt{url} 明明是用 \texttt{https} 開頭的，為什麼
\texttt{url.startswith("http://")} 是 \texttt{False}？因為 \texttt{startswith()} 是
逐字元比對，\texttt{"http://"} 這串文字裡的第五個字元是 \texttt{:}，但 \texttt{url}
對應位置上的字元是 \texttt{s}（來自 \texttt{https}），兩者對不上，所以整體判斷不成立。
\texttt{startswith()} 不會幫你猜測你的意圖，你給它什麼內容，它就照字面一個字元一個字元
去對。

\stylepoint{判斷「這個字串是不是某個開頭或結尾」，優先用 \texttt{startswith()} 或
\texttt{endswith()}，不要土法煉鋼用切片自己算，例如寫成
\texttt{filename[-4:] == ".pdf"}。這種寫法要自己數副檔名有幾個字元，容易數錯，也比較
難讀；\texttt{endswith(".pdf")} 一眼就看得懂在做什麼，而且不用擔心字數算錯。}

\begin{drills}[label=drills:startswith-endswith]
\item 給定 \texttt{fname = "photo.png"}，\texttt{fname.endswith(".jpg")} 的值是什麼？
\item 給定 \texttt{fname = "photo.png"}，\texttt{fname.endswith(".png")} 的值是什麼？
\item 給定 \texttt{fname = "photo.png"}，\texttt{fname.startswith("photo")} 的值是
什麼？
\end{drills}

\section{\texttt{split()}：把一段文字拆成好幾段}
\label{sec:split-method}

有了「呼叫方法」這個新語法，接下來幾節要認識幾個你之後會反覆用到的字串方法。第一個是
\texttt{split()}：把一個字串依照指定的分隔符號拆開，變成好幾段。想像你手上有一整行用
逗號隔開的清單文字，想把每一項分開處理，\texttt{split()} 就是做這件事的工具。
\index{split()}

\begin{code}{python}{依分隔符拆開字串}
s.split(separator)
\end{code}

\begin{example}{python}{(split-list.py) 把逗號分隔的清單文字拆開}
line = "apple,banana,cherry"
fruits = line.split(",")
print(fruits)
\end{example}

\iosample{}{['apple', 'banana', 'cherry'] \\}

你可能注意到輸出前後多了一對中括號，每一項也各自被單引號包住，這是下一個大主題 list
的外觀，第 \ref{ch:lists} 章會正式介紹它。這裡你只要知道：\texttt{split()} 會把字串拆成好幾段，
裝進這種容器裡，之後可以用第 \ref{ch:lists} 章教的方式一個一個拿出來用即可。

\texttt{split()} 也可以完全不帶任何引數呼叫，這時候它會依「空白」拆開，而且不管中間有
幾個連續的空白，都只會當成一個分隔點，不會因為中間多空了幾格就多切出空字串。

\begin{example}{python}{(split-sentence.py) 不帶引數，依空白拆開句子}
sentence = "the quick   brown fox"
words = sentence.split()
print(words)
\end{example}

\iosample{}{['the', 'quick', 'brown', 'fox'] \\}

\keypoint{\texttt{split()} 不帶引數時，會把連續好幾個空白都當成同一個分隔點，不會產生
空字串；但如果你自己明確指定分隔符（像 \texttt{split(" ")}），Python 就會照字面上的
規則機械式處理，字串裡連續兩個空格中間會多切出一個空字串。不確定資料裡的空白會不會連續
出現時，不帶引數呼叫 \texttt{split()} 通常比較安全。}

\section{\texttt{join()}：把好幾段文字接回一句話}
\label{sec:join-method}

\texttt{split()} 把一個字串拆成好幾段，\texttt{join()} 做的是完全相反的事：把好幾段
文字，用一個指定的分隔符號接回成一個字串。
\index{join()}這兩個方法方向相反，但語法的「接法」也剛好
相反，非常容易搞混：\texttt{split} 是接在\textbf{要拆的字串}上呼叫；\texttt{join} 反過來，
是接在\textbf{分隔符字串}上呼叫，把裝了好幾段文字的容器當成參數傳進去。

\begin{code}{python}{用分隔符把好幾段文字接回去}
separator.join(strings)
\end{code}

\begin{example}{python}{(rejoin-fruits.py) 把拆開的水果清單重新接回去}
fruits = "apple,banana,cherry".split(",")
print(fruits)                    # 先拆開，跟上一節一樣

sentence = ", ".join(fruits)      # 用逗號接回成一句話
print(sentence)

lines = ["Line 1", "Line 2", "Line 3"]
print("\n".join(lines))           # \n 代表換行，接成三行
\end{example}

\iosample{}{['apple', 'banana', 'cherry'] \\ apple, banana, cherry \\ Line 1 \\ Line 2 \\ Line 3 \\}

\keypoint{\texttt{split} 接在\textbf{要拆的字串}上（\texttt{s.split(分隔符)}）；
\texttt{join} 接在\textbf{分隔符}上（\texttt{分隔符.join(一堆字串)}），方向剛好相反，
新手非常容易把兩者搞反，寫成 \texttt{fruits.join(", ")} 這種錯誤寫法。記憶方式：誰站在
點的前面，誰就是「主詞」；\texttt{split} 的主詞是被拆的字串，\texttt{join} 的主詞是
分隔符。}

\texttt{join()} 還有一個常見誤解：它要求容器裡裝的每一項都得是字串，如果不小心裝的是
數字，Python 不會貼心地幫你自動轉型，只會直接出錯。

\begin{example}{python}{(join-numbers-oops.py) 試著用 join 接數字組成的清單}
numbers = [1, 2, 3]
print(",".join(numbers))
\end{example}

存檔執行後，畫面上會出現：

\begin{verbatim}
Traceback (most recent call last):
  File "join-numbers-oops.py", line 2, in <module>
    print(",".join(numbers))
          ~~~~~~~~^^^^^^^^^
TypeError: sequence item 0: expected str instance, int found
\end{verbatim}

\texttt{join()} 只接受「裝著字串的容器」，清單裡只要有一項不是字串（這裡是 \texttt{int}
型別的數字），就會出現 \texttt{TypeError}，訊息說「本來期待字串，卻拿到了 int」。想接
一份裝著數字的清單，得先把每個數字自己轉成字串，這裡可以用第 \ref{ch:variables} 章學過的 \texttt{str()}
型別轉換：

\begin{example}{python}{(join-numbers-fixed.py) 先用 str() 把每個數字轉成字串}
a = 1
b = 2
c = 3
strings = [str(a), str(b), str(c)]
print(",".join(strings))
\end{example}

\iosample{}{1,2,3 \\}

\begin{drills}[label=drills:split-join]
\item 預測輸出：\texttt{"a-b-c".split("-")} 的結果是什麼？
\item 預測輸出：\texttt{"-".join(["a", "b", "c"])} 的結果是什麼？
\item 給定
\begin{verbatim}
parts = "2024/09/01".split("/")
result = "-".join(parts)
\end{verbatim}
\texttt{result} 的值會是什麼？（注意 split 跟 join 用的分隔符號不一樣）
\item 給定 \texttt{"a  b   c".split()}（沒有指定分隔符），結果會是什麼？
\item 如果直接寫 \texttt{",".join([1, 2, 3])}（清單裡是數字，不是字串），會發生什麼事？
為什麼？
\end{drills}

\section{\texttt{replace()}：搜尋並替換文字}
\label{sec:replace-method}

最後一個要認識的方法是 \texttt{replace()}：在字串裡找出所有符合的子字串，全部換成另一段
文字。
\index{replace()}它的用法很直覺，告訴它要找什麼、換成什麼，兩個都是字串。

\begin{code}{python}{搜尋並替換文字}
s.replace(old, new)
\end{code}

\begin{example}{python}{(fix-typo.py) 訂正重複打錯的字}
sentence = "I am studing Python, and I really enjoy studing it."
fixed = sentence.replace("studing", "studying")
print(fixed)
\end{example}

\iosample{}{I am studying Python, and I really enjoy studying it. \\}

要特別小心的是：\texttt{replace()} 不會判斷「這個位置是不是一個完整的詞」，它只是機械式
地找出每一個符合的子字串、逐一換掉。如果你要換掉的字串太短，剛好也是其他詞的一部分，就
可能誤傷到不相關的地方：

\begin{example}{python}{(replace-oops.py) 替換太短的子字串，誤傷其他地方}
text = "cat catalog"
result = text.replace("cat", "dog")
print(result)
\end{example}

\iosample{}{dog dogalog \\}

原本只想把單獨的 \texttt{"cat"} 換成 \texttt{"dog"}，結果連 \texttt{"catalog"} 裡開頭
的那個 \texttt{"cat"} 也一起被換掉，變成了怪怪的 \texttt{"dogalog"}。使用
\texttt{replace()} 時，養成習慣想一下：要換掉的這段文字，會不會也出現在其他不想被動到
的地方。

\texttt{replace()} 其實還可以傳第三個引數，限制最多替換幾次，超過這個次數的部分會維持
原樣，不會被繼續替換下去。想只改掉「前幾個」符合的地方、又不想全部一次換光，這個引數
就派得上用場。

\begin{example}{python}{(replace-limit.py) 限制替換次數}
text = "cat cat cat"
print(text.replace("cat", "dog", 1))   # 只替換第一個
print(text.replace("cat", "dog", 2))   # 替換前兩個
\end{example}

\iosample{}{dog cat cat \\ dog dog cat \\}

\begin{drills}[label=drills:replace-method]
\item 預測輸出：\texttt{"banana".replace("an", "AN")} 的結果是什麼？
\item 預測輸出，並判斷有沒有誤傷其他地方：\texttt{"cats and cathedrals".replace("cat",
"dog")} 的結果是什麼？
\item 預測輸出：\texttt{"cat cat cat".replace("cat", "dog", 2)} 的結果是什麼？
\end{drills}

\section{\texttt{find()} 與 \texttt{count()}：搜尋子字串的位置與次數}
\label{sec:find-count}

\texttt{replace()} 會直接動手把符合的內容換掉，但有時候你只是想「知道」某段文字有沒有
出現、出現在哪裡、出現了幾次，還不想真的改動任何東西。這時候可以用 \texttt{find()} 跟
\texttt{count()}：\texttt{find()} 回傳子字串第一次出現的位置（索引），如果根本沒出現，
回傳 \texttt{-1}；\index{find()}
\texttt{count()} 回傳子字串總共出現了幾次。
\index{count()}

\begin{code}{python}{搜尋子字串的位置與次數}
s.find(substring)
s.count(substring)
\end{code}

\begin{example}{python}{(find-substring.py) 找出子字串第一次出現的位置}
text = "banana"
print(text.find("na"))    # 找到了，回傳第一次出現的位置
print(text.find("z"))     # 沒找到，回傳 -1
\end{example}

\iosample{}{2 \\ -1 \\}

\begin{example}{python}{(count-substring.py) 數子字串出現的次數}
text = "banana"
print(text.count("a"))    # 單一字元出現了幾次
print(text.count("na"))   # 兩個字元的子字串出現了幾次
\end{example}

\iosample{}{3 \\ 2 \\}

\keypoint{\texttt{find()} 找不到東西的時候，回傳的是 \texttt{-1}，不是報錯，也不是
\texttt{None}。這是新手很容易忽略的地方：如果你直接把 \texttt{find()} 的回傳值當成一個
有效的位置去用（例如拿去切片），\texttt{-1} 剛好是合法的負數索引，代表「倒數第一個」，
程式不會報錯，但結果會完全不對，是一種不容易被發現的邏輯錯誤。想判斷「到底有沒有找到」，
要先檢查回傳值是不是 \texttt{-1}，不能只看程式有沒有報錯。}

\texttt{count()} 有一個和 \texttt{replace()} 類似的地方：它一樣是機械式地逐一比對，
不會判斷「這是不是一個完整的詞」。\texttt{"mississippi".count("is")} 之所以是 2，是因為
\texttt{"is"} 這段文字在整串裡剛好出現了兩次，跟它是不是獨立的一個詞完全無關。

\begin{drills}[label=drills:find-count]
\item 給定 \texttt{text = "mississippi"}，\texttt{text.count("s")} 的值是什麼？
\item 給定 \texttt{text = "mississippi"}，\texttt{text.find("ss")} 的值是什麼？
\item 給定 \texttt{text = "mississippi"}，\texttt{text.find("z")} 的值是什麼？
\end{drills}

\section{f-string：把變數的值嵌進一句話裡}
\label{sec:f-string-basics}

前面幾節都是在拆解、清理既有的文字，這一節要反過來：把幾個變數的值，組回一句人看得懂的
話。你之前想印出一句包含變數的話，大概是用 \texttt{+} 把好幾段字串跟變數接起來，這樣寫
其實不算太難讀，但變數一多、標點符號一多，很快就會變得很難維護。Python 提供了更直接的
寫法，叫做 f-string：在字串的引號前面加一個 \texttt{f}，字串裡面用大括號包住的部分，會
被直接換成變數的值。
\index{f-string}

\begin{code}{python}{把變數的值嵌進字串裡}
f"text{variable}text"
\end{code}

\begin{example}{python}{(introduce-myself.py) 組一句自我介紹}
name = "Taylor"
age = 15
print(f"Hello, I'm {name} and I'm {age} years old.")
\end{example}

\iosample{}{Hello, I'm Taylor and I'm 15 years old. \\}

f-string 有一個特別容易漏掉、卻不會讓程式報錯的地方：字串引號前面的那個 \texttt{f}。少
打這一個字母，Python 不會把它當成錯誤，只會把整個字串當成普通文字，大括號跟裡面的
變數名稱會被原封不動地印出來：

\begin{example}{python}{(missing-f.py) 忘記加 f 前綴}
name = "Taylor"
age = 15
print("Hello, I'm {name} and I'm {age} years old.")
\end{example}

\iosample{}{Hello, I'm \{name\} and I'm \{age\} years old. \\}

\keypoint{忘記加 \texttt{f} 前綴，是一個「安靜的錯」，程式不會報任何錯誤訊息，只是
輸出的內容不對，大括號跟變數名稱原封不動地被印了出來。這種錯誤特別容易讓人卡關，因為
螢幕上沒有紅字可以看，你得靠自己發現「怎麼變數的值沒有被換進去」。看到輸出裡出現不該
出現的大括號，第一個該檢查的就是：引號前面是不是漏了 \texttt{f}。}

\section{f-string 進階：在大括號裡放運算式}
\label{sec:f-string-advanced}

大括號裡不是只能放一個變數名稱，其實可以放任何一段能算出結果的運算式，Python 會先把
它算出來，再把結果嵌進字串裡。一個字串裡也可以放好幾組大括號，一次組合好幾個值。

\begin{code}{python}{在大括號裡放運算式}
f"{expression}"
\end{code}

\begin{example}{python}{(shopping-total.py) 計算購物金額並印出}
price = 35
quantity = 3
print(f"Unit price ${price}, quantity {quantity}, total ${price * quantity}.")
\end{example}

\iosample{}{Unit price \$35, quantity 3, total \$105. \\}

\stylepoint{與其先算出一個新變數（例如 \texttt{total = price * quantity}）再放進
f-string，直接把運算式寫在大括號裡（\texttt{\{price * quantity\}}）通常更常見、也更
省事，前提是這段運算式夠簡短、一眼就看得懂在算什麼；如果運算式變得又長又複雜，先拆成
一個有名字的變數，可讀性反而更好。}

如果你想更講究一點，例如金額要固定顯示小數點後兩位、文字要靠右對齊，f-string 大括號裡
其實還可以加上額外的格式化語法，用來控制小數位數、千分位、對齊方式，這是下一節要正式
介紹的內容。

\begin{drills}[label=drills:f-string]
\item 有人把 \texttt{f} 忘記加了，寫成下面這樣：
\begin{verbatim}
x = 5
print("Value is {x}")
\end{verbatim}
這段程式的輸出會是什麼？跟正確加上 \texttt{f} 的版本比起來，差在哪裡？
\item 預測輸出：
\begin{verbatim}
a = 4
b = 6
print(f"{a}+{b}={a + b}")
\end{verbatim}
\item 預測輸出：
\begin{verbatim}
name = "Bob"
score = 88
print(f"{name} scored {score} points")
\end{verbatim}
\end{drills}

\section{f-string 進階：控制數字要怎麼顯示}
\label{sec:f-string-format}

大括號裡除了放變數或運算式，還可以在後面加一個冒號，接著寫「格式規格」，控制這個值最後
要怎麼被排版顯示，例如小數只留幾位、千位要不要加逗號分隔、文字要靠左還是靠右對齊。
\index{f-string 格式化選項}語法是
在大括號裡，運算式後面接一個冒號，冒號後面接格式規格。

\begin{code}{python}{在大括號裡加上格式規格}
f"{expression:format_spec}"
\end{code}

最常用的格式規格是控制小數位數：\texttt{.2f} 代表「當成有小數點的數字處理，固定顯示到
小數點後兩位」。這裡的 \texttt{f} 跟字串前面的 f-string 的 \texttt{f} 是兩件不相干的事，
只是剛好長得一樣，這裡的 \texttt{f} 指的是 float（浮點數）。

\begin{example}{python}{(format-decimal.py) 固定顯示小數點後兩位}
price = 35.5
quantity = 3
total = price * quantity
print(f"Unit price ${price}, quantity {quantity}, total ${total:.2f}.")
\end{example}

\iosample{}{Unit price \$35.5, quantity 3, total \$106.50. \\}

沒有格式規格的時候，\texttt{price * quantity} 算出來是 \texttt{106.5}，只有一位小數；
加上 \texttt{:.2f}，Python 會幫你補滿到兩位小數，變成 \texttt{106.50}，這在處理金額的
時候特別有用，畢竟現實中的金額習慣上都會顯示到小數點後兩位，就算剛好是整數也一樣。

\keypoint{\texttt{:.2f} 只是改變「顯示」出來的樣子，不會真的把變數本身的值四捨五入存
回去。\texttt{total} 這個變數本身還是原本的 \texttt{106.5}，之後如果你直接用
\texttt{total} 去做別的運算，用的還是完整精度的原始值，不是螢幕上看到的那個四捨五入
後的版本。}

除了小數位數，數字太大不好一眼看出位數的時候，可以加上千分位逗號分隔，格式規格寫一個
逗號 \texttt{,} 就可以。

\begin{example}{python}{(format-thousands.py) 幫大數字加上千分位逗號}
population = 23894394
print(f"Population is about {population:,} people.")
\end{example}

\iosample{}{Population is about 23,894,394 people. \\}

格式規格也可以控制文字（不只是數字）要用多寬的欄位顯示、要靠左、靠右還是置中對齊，這在
排整齊的表格輸出時很好用：欄位寬度寫在冒號後面的數字，前面加 \texttt{>} 代表靠右對齊、
\texttt{<} 代表靠左對齊、\texttt{\textasciicircum} 代表置中，如果不指定對齊符號，數字
預設靠右、文字預設靠左。

\begin{example}{python}{(format-align.py) 控制欄位寬度與對齊方式}
name = "Bob"
print(f"[{name:>10}]")
print(f"[{name:<10}]")
print(f"[{name:^10}]")
\end{example}

下面用 \texttt{␣} 代表一個空白字元，方便你數清楚每個欄位實際留了多少空白：

\iosample{}{[␣␣␣␣␣␣␣Bob] \\{} [Bob␣␣␣␣␣␣␣] \\{} [␣␣␣Bob␣␣␣␣] \\}

三個輸出都被塞進總共 10 個字元寬的欄位裡，缺的部分用空白補上，只是補在不同邊：
\texttt{>} 把空白補在左邊（讓內容貼右），\texttt{<} 把空白補在右邊（讓內容貼左），
\texttt{\textasciicircum} 儘量把空白平均補在兩邊，補不平均的時候，多出來的一格會補
在右邊。

\begin{drills}[label=drills:f-string-format]
\item 給定 \texttt{pi = 3.14159265}，\texttt{f"{pi:.3f}"} 的值是什麼？
\item 給定 \texttt{big = 1234567}，\texttt{f"{big:,}"} 的值是什麼？
\item 給定 \texttt{name = "Bo"}，\texttt{f"[{name:>6}]"} 的值是什麼？（用 \texttt{␣}
代表空白）
\end{drills}

\section*{練習解答}
\addcontentsline{toc}{section}{練習解答}

\paragraph{第~\ref{sec:string-indexing}~節}
\begin{enumerate}
\item \texttt{'a'}，因為 \texttt{word[2]} 是第三個字元。
\item \texttt{'n'}，從尾端數回來第三個字元。
\item 會出現 \texttt{IndexError: string index out of range}。因為
\texttt{word} 總共有 6 個字元，有效索引是 0 到 5，\texttt{word[6]} 已經超出範圍，
字元總數跟最後一個字元的索引永遠差 1。
\end{enumerate}

\paragraph{第~\ref{sec:string-slicing}~節}
\begin{enumerate}
\item \texttt{'mpu'}，從位置 2 切到位置 5（不含）。
\item \texttt{'comp'}，省略 start，從頭切到位置 4（不含）。
\item \texttt{'ter'}，從倒數第 3 個字元切到結尾。
\end{enumerate}

\paragraph{第~\ref{sec:string-slicing-step}~節}
\begin{enumerate}
\item \texttt{'cmue'}，從頭到尾每隔一個字元取一次。
\item \texttt{'retupmoc'}，step 為 \texttt{-1}，整個字串反過來。
\item \texttt{'opt'}，從位置 1 開始、每隔一個取一次，取到位置 7（不含）之前。
\end{enumerate}

\paragraph{第~\ref{sec:strip-method}~節}
\begin{enumerate}
\item \texttt{[hi]}，\texttt{strip()} 清掉了頭尾的空白，中括號會緊貼著 \texttt{hi}。
\item 還會有，因為呼叫 \texttt{msg.strip()} 只是算出一個新字串，沒有把它重新指定回
\texttt{msg}，原本的 \texttt{msg} 完全沒有被改動。
\item \texttt{'python'}，頭尾的 \texttt{x} 一路被剝掉，直到剝到不是 \texttt{x} 的字元
為止。
\item \texttt{'pythonxx'}。\texttt{lstrip()} 只清左邊，所以開頭的兩個 \texttt{x} 被
剝掉了，結尾的兩個原封不動留著。
\end{enumerate}

\paragraph{第~\ref{sec:case-methods}~節}
\begin{enumerate}
\item \texttt{'hello world'}。
\item \texttt{'HELLO WORLD'}。
\item \texttt{True}，兩邊都先轉成小寫後內容一樣。
\end{enumerate}

\paragraph{第~\ref{sec:startswith-endswith}~節}
\begin{enumerate}
\item \texttt{False}，\texttt{fname} 是以 \texttt{.png} 結尾，不是 \texttt{.jpg}。
\item \texttt{True}。
\item \texttt{True}。
\end{enumerate}

\paragraph{第~\ref{sec:join-method}~節}
\begin{enumerate}
\item \texttt{['a', 'b', 'c']}。
\item \texttt{'a-b-c'}。
\item \texttt{'2024-09-01'}，先用 \texttt{/} 拆開，再用 \texttt{-} 接回去，兩次用的
分隔符可以不一樣。
\item \texttt{['a', 'b', 'c']}，不帶引數的 \texttt{split()} 會把連續的空白都當成一個
分隔點。
\item 會出現 \texttt{TypeError}，因為 \texttt{join()} 要求容器裡每一項都得是字串，
清單裡是 \texttt{int} 型別的數字，得先用 \texttt{str()} 把每個數字轉成字串才能接起來。
\end{enumerate}

\paragraph{第~\ref{sec:replace-method}~節}
\begin{enumerate}
\item \texttt{'bANANa'}，字串裡兩處符合 \texttt{"an"} 的地方都被換成了 \texttt{"AN"}。
\item \texttt{'dogs and doghedrals'}，不只獨立的 \texttt{"cat"} 被換掉，
\texttt{"cathedrals"} 開頭的 \texttt{"cat"} 也一起被誤傷，變成了不存在的字
\texttt{"doghedrals"}。
\item \texttt{'dog dog cat'}，第三個引數限制最多替換兩次，第三個 \texttt{"cat"} 維持
原樣沒被換掉。
\end{enumerate}

\paragraph{第~\ref{sec:find-count}~節}
\begin{enumerate}
\item \texttt{4}，\texttt{"mississippi"} 裡總共有 4 個 \texttt{s}。
\item \texttt{2}，\texttt{"ss"} 第一次出現的位置是索引 2。
\item \texttt{-1}，\texttt{"mississippi"} 裡完全沒有 \texttt{z}。
\end{enumerate}

\paragraph{第~\ref{sec:f-string-advanced}~節}
\begin{enumerate}
\item 輸出會是 \texttt{Value is \{x\}}，因為少了 \texttt{f} 前綴，Python 把整個字串當成
普通文字，大括號跟裡面的 \texttt{x} 原封不動被印出來，不會被換成變數的值 5。
\item \texttt{4+6=10}。
\item \texttt{Bob scored 88 points}。
\end{enumerate}

\paragraph{第~\ref{sec:f-string-format}~節}
\begin{enumerate}
\item \texttt{'3.142'}，固定顯示到小數點後三位。
\item \texttt{'1,234,567'}，每三位數字加一個逗號分隔。
\item \texttt{'[␣␣␣␣Bo]'}，\texttt{"Bo"} 只有 2 個字元，靠右對齊塞進寬度 6 的欄位，
前面補 4 個空白。
\end{enumerate}

\chapter{條件判斷}
\label{ch:conditionals}

每個人每天都在做「如果……就……否則……」這種判斷，只是平常不會意識到自己在做。外面
在下雨，出門就帶傘，不帶傘就淋雨；紅燈亮著，就停下來等，燈一變綠才走。這一章要教的，
就是怎麼把這種你已經很熟悉的判斷方式，翻譯成電腦看得懂的語言。

我們分幾段推進：
\begin{itemize}
\item 先從最小的判斷開始，一次只問一個「是不是」的問題，讓電腦根據這一個答案決定
做不做某件事。
\item 接著把範圍擴大：一個問題往往不夠用，你可能要在好幾條路裡選一條，或是同時考慮
好幾個條件才能下決定，這時候需要更多分支，需要把好幾個「是不是」串在一起，也會碰到
「條件裡面還藏著條件」這種更貼近真實生活的情境，以及好幾個條件擠在同一行時，怎麼寫
才不會讓自己（或別人）看不懂。
\item 最後，我們要介紹 Python 一個有點特別的脾氣：不只 \texttt{True}、\texttt{False}
算是非黑即白的答案，Python 在 \texttt{if} 裡看很多東西時，會把「空空如也」的東西直接
當成 \texttt{False}，把「有內容」的東西當成 \texttt{True}，不用你特別去檢查「是不是
空的」。
\end{itemize}

整章下來，你會從「只會判斷一件事」慢慢累積到「能把好幾個日常生活裡稀鬆平常的判斷，原封
不動翻譯成 Python 看得懂的條件」。

\section{比較是不是：認識比較運算子}
\label{sec:comparison-operators}

要讓電腦幫你做判斷，第一步是先教它怎麼問一個「是不是」的問題，例如「這個數字是不是比
那個數字大」「這兩個值是不是相等」。這種問句在 Python 裡叫做比較，用的是比較運算子，問完
之後得到的答案，就是你在第 \ref{ch:variables} 章認識過的 \texttt{bool} 型別：不是 \texttt{True} 就是
\texttt{False}，沒有第三種答案。
\index{比較}
\index{比較運算子}

\begin{code}{python}{比較兩個值}
value1 operator value2
\end{code}

Python 提供的比較運算子有這幾個：\texttt{==}（相等）、\texttt{!=}（不相等）、
\texttt{<}（小於）、\texttt{>}（大於）、\texttt{<=}（小於等於）、\texttt{>=}（大於等於）。

\begin{example}{python}{(compare-scores.py) 比較兩科分數哪個高}
math_score = 85
english_score = 92

print(math_score > english_score)
print(math_score == english_score)
print(math_score != english_score)
\end{example}

\iosample{}{False \\ False \\ True \\}

\keypoint{比較運算子的結果永遠是一個 \texttt{bool} 值，只有 \texttt{True} 或
\texttt{False} 兩種可能。這正是第 \ref{ch:variables} 章介紹過的四種基本型別之一，這裡是型別介紹完之後，
你第一次真正拿它來做事。}

要特別注意 \texttt{==} 跟 \texttt{=} 是兩個完全不同的東西：\texttt{=} 是「把值指定給
一個名字」（例如 \texttt{score = 60}），\texttt{==} 才是「問這兩個值是不是相等」。很多人
在數學課習慣把等號當成「相等」的意思，寫 Python 時反而容易把這兩個搞混，下一節會實際
示範打錯會發生什麼事。

比較運算子不只能用在數字身上。字串也可以用 \texttt{==} 比較內容是不是完全一樣，例如
\texttt{"apple" == "apple"} 會是 \texttt{True}，\texttt{"apple" == "Apple"} 因為大小寫
不同會是 \texttt{False}。至於 \texttt{<}、\texttt{>} 用在字串上，牽涉到字母排列順序的
規則，這本書之後真正需要的時候才會提到，這一章你只要先把數字的比較用熟就好，那是最直覺、
也最常用的情況。

還有一個小細節值得留意：\texttt{int} 跟 \texttt{float} 這兩種數字型別互相比較時，Python
只在乎數值本身，不在乎型別是不是完全一樣。

\begin{example}{python}{(mixed-type-compare.py) int 跟 float 互相比較}
price = 100
budget = 100.0

print(price == budget)
print(price < 100.5)
\end{example}

\iosample{}{True \\ True \\}

\texttt{price} 是 \texttt{int} 的 100，\texttt{budget} 是 \texttt{float} 的
100.0，兩者型別不一樣，但值相等，\texttt{price == budget} 一樣是 \texttt{True}。這跟
第 \ref{ch:variables} 章「型別長在物件上，不是長在變數上」的概念不衝突：比較運算子看的是兩邊的值，型別
不同只要數值相等，比較的結果照樣成立。

\begin{drills}[label=drills:comparison-operators]
\item \texttt{price = 120}，這個比較式的結果是 \texttt{True} 還是 \texttt{False}？
\texttt{price >= 120}
\item \texttt{name = "Alice"}，這個比較式的結果是 \texttt{True} 還是 \texttt{False}？
\texttt{name == "alice"}
\end{drills}

\section{如果……就……：if 陳述句}
\label{sec:if-basics}

光是問出「是不是」還不夠，你通常還想根據這個答案，決定電腦要不要做某件事：分數及格才印出
「通過」，天氣冷才提醒帶外套。這種「只有條件成立才做」的結構，在 Python 裡叫做
\texttt{if} 陳述句。
\index{if 陳述句}

\begin{code}{python}{if 陳述句}
if condition:
    action
\end{code}

\texttt{if} 後面接一個條件（通常是一個比較運算式），再接一個冒號；下一行開始，只要是
縮排過的程式碼，都算是「條件成立才做」的區塊，Python 靠縮排（通常是 4 個空白）來判斷
一段程式碼屬不屬於這個區塊，不是靠大括號或其他符號。

\begin{example}{python}{(check-passing.py) 分數及格才印出通過}
score = 75
if score >= 60:
    print("Pass")
\end{example}

\iosample{}{Pass \\}

如果你把第一行的 \texttt{score} 改成 \texttt{50}，重新執行這支程式，螢幕上會完全沒有
任何輸出。因為 \texttt{50 >= 60} 是 \texttt{False}，縮排的那一行 \texttt{print} 根本
不會被執行，Python 會直接跳過整個 \texttt{if} 區塊，往下面（如果還有的話）繼續執行。

這一節開始，你要正式打字寫出 \texttt{if} 這個結構，剛好是新手最容易在三個地方打錯的時候，
接下來一個一個示範打錯會發生什麼事。

\textbf{漏掉冒號：}

\begin{example}{python}{(check-passing-broken-colon.py) 漏掉冒號}
score = 75
if score >= 60
    print("Pass")
\end{example}

執行後會看到：

\begin{verbatim}
  File "check-passing-broken-colon.py", line 2
    if score >= 60
                  ^
SyntaxError: expected ':'
\end{verbatim}

\texttt{expected ':'} 就是字面上的意思。Python 讀到這一行，發現該有冒號的地方沒有冒號，
直接告訴你它期待看到什麼。補上冒號就能修好。

\textbf{縮排不對：}

\begin{example}{python}{(check-passing-broken-indent.py) 少了縮排}
score = 75
if score >= 60:
print("Pass")
\end{example}

執行後會看到：

\begin{verbatim}
  File "check-passing-broken-indent.py", line 3
    print("Pass")
    ^^^^^
IndentationError: expected an indented block after 'if' statement on line 2
\end{verbatim}

\texttt{IndentationError} 專門處理縮排出錯的情況；這裡的意思是「第 2 行的 if 陳述句
後面，應該要接一個縮排過的區塊，但沒有」。把 \texttt{print} 那一行往右縮排 4 個空白
即可。
\index{IndentationError}

\textbf{把 \texttt{==} 打成 \texttt{=}：}

\begin{example}{python}{(check-passing-broken-equals.py) 等號打錯}
score = 75
if score = 60:
    print("Exactly 60 points")
\end{example}

執行後會看到：

\begin{verbatim}
  File "check-passing-broken-equals.py", line 2
    if score = 60:
       ^^^^^^^^^^
SyntaxError: invalid syntax. Maybe you meant '==' or ':=' instead of '='?
\end{verbatim}

這個錯誤訊息特別貼心，直接猜到你大概是想打 \texttt{==}。\texttt{if} 後面的條件要問的是
「是不是相等」，不是「把值指定給誰」，這種時候一定要用 \texttt{==}。

這三個錯誤有一個共同的地方：Python 都是在還沒有真正開始執行程式之前，光是「讀」你的
程式碼這一步就卡住了，所以看到的都是 \texttt{SyntaxError} 或 \texttt{IndentationError}，
而不是程式執行到一半才出的錯。之後你會漸漸讀得懂這類錯誤訊息的固定套路：先看錯誤類型，
再看它指到程式碼的哪個位置、用什麼符號標出問題所在，通常就能猜到八九不離十的修法。

\begin{drills}[label=drills:if-basics]
\item 這段程式的輸出是什麼？
\begin{verbatim}
score = 88
if score > 90:
    print("Excellent")
print("Done")
\end{verbatim}
\item 這段程式漏了一個符號，Python 執行時會出現什麼類型的錯誤？請指出漏了什麼、怎麼修：
\begin{verbatim}
temperature = 30
if temperature > 28
    print("Hot")
\end{verbatim}
\item 這段程式的作者本來想「氣溫超過 28 度才印出 Hot」，但寫錯了一個符號，導致每次執行
都會印出「Hot」，不管氣溫多少度。請找出錯在哪裡：
\begin{verbatim}
temperature = 20
if temperature = 28:
    print("Hot")
\end{verbatim}
\end{drills}

\section{不然呢：if/else 二選一}
\label{sec:if-else}

上一節的 \texttt{if} 只處理了「條件成立才做」，但如果條件不成立，你通常也想做點別的事，
不是什麼都不做。分數不及格，也想印出「不及格」提醒一下。這種「兩條路一定會走其中一條」
的結構，要用 \texttt{if}/\texttt{else} 搭配。
\index{if/else}

\begin{code}{python}{if/else}
if condition:
    actionA
else:
    actionB
\end{code}

\texttt{else} 不需要再寫一次條件，它的意思就是「上面那個條件不成立的話」，跟
\texttt{if} 一樣，冒號後面接縮排過的區塊。

\begin{example}{python}{(pass-or-fail.py) 及格或不及格}
score = 55
if score >= 60:
    print("Pass")
else:
    print("Fail")
\end{example}

\iosample{}{Fail \\}

\keypoint{\texttt{if}/\texttt{else} 保證兩條路裡一定會走其中一條，不會有「兩條都不
執行」或「兩條都執行」的情況。這跟下一節要介紹的、一連串各自獨立的 \texttt{if}
不一樣，那種寫法反而可能同時執行好幾條路，或一條都不執行。}

也要注意 \texttt{else} 本身沒辦法再加條件。如果你想要的其實是「分數 90 以上、80 到 89、
其他」這種好幾條路的分法，光靠一個 \texttt{if}/\texttt{else} 是做不到的，這正是接下來
兩節要解決的問題：先看「條件裡面還可以再包一層條件」，再看「怎麼把好幾個條件串成一條鏈」。

\begin{drills}[label=drills:if-else]
\item \texttt{score = 60}，這段程式會印出什麼？
\begin{verbatim}
if score >= 60:
    print("Pass")
else:
    print("Fail")
\end{verbatim}
\item 下面這段程式想做的是「及格印通過，不及格印不及格」，但完全沒有反應，執行後畫面
什麼都沒有。請指出問題在哪裡：
\begin{verbatim}
score = 45
if score >= 60
    print("Pass")
else:
    print("Fail")
\end{verbatim}
\end{drills}

\section{條件裡面還有條件：巢狀 if}
\label{sec:nested-if}

現實生活裡的判斷，常常一層接著一層，很少問一次就結束：先確認有沒有票，有票才
接著看年齡夠不夠買成人票；先確認今天有沒有下雨，有下雨才接著看有沒有帶傘。前面一個
條件先成立，後面那個條件才有意義去問。

這種「條件裡面還有條件」的結構，在 Python 裡叫做
巢狀 if（\texttt{if} 裡面再包一個 \texttt{if}），寫法是把整個內層的 \texttt{if}（或
\texttt{if}/\texttt{else}）當成外層 \texttt{if} 底下「條件成立才做」的那個動作。
\index{巢狀 if}

\begin{code}{python}{巢狀 if}
if condition1:
    if condition2:
        actionA
    else:
        actionB
else:
    actionC
\end{code}

內層的 \texttt{if}/\texttt{else} 要比外層的 \texttt{if} 多縮排一層（通常是再多 4 個
空白，變成 8 個空白），這樣 Python 才知道它屬於外層條件成立時的那個區塊。縮排的層數
沒有上限，你可以視需要再往裡面包第三層、第四層，但巢狀包太多層通常代表可以想辦法簡化，
這一點在後面幾節會再回頭討論。

\begin{example}{python}{(party-entry.py) 有票才看年齡夠不夠}
age = 16
has_ticket = True

if has_ticket:
    if age >= 18:
        print("Adult ticket, entry allowed")
    else:
        print("Minor, parental accompaniment required")
else:
    print("No ticket, entry denied")
\end{example}

\iosample{}{Minor, parental accompaniment required \\}

一步一步看這段程式在做什麼：外層先問 \texttt{has\_ticket} 是不是 \texttt{True}。這裡
\texttt{has\_ticket} 是 \texttt{True}，於是往內層走，接著問 \texttt{age >= 18}；
\texttt{age} 是 16，不滿 18，所以走進內層的 \texttt{else}，印出「Minor, parental
accompaniment required」。如果一開始 \texttt{has\_ticket} 就是 \texttt{False}，內層那整段 \texttt{if}/
\texttt{else}（不管年齡是多少）根本不會被檢查，會直接跳到最外層的 \texttt{else}，
印出「No ticket, entry denied」。

\keypoint{巢狀 \texttt{if} 表達的是「後面的條件，要在前面的條件先成立的情況下才有意義」，
這跟下一節要介紹的 \texttt{elif} 鏈不一樣：\texttt{elif} 鏈是好幾條互斥的路，彼此地位
平等，巢狀 \texttt{if} 則是一層一層有先後關係，內層條件完全依賴外層條件先成立。}

\begin{drills}[label=drills:nested-if]
\item 這段程式的輸出是什麼？
\begin{verbatim}
weather = "Rainy"
has_umbrella = True

if weather == "Rainy":
    if has_umbrella:
        print("Can go out")
    else:
        print("Wait for rain to stop")
else:
    print("Go out directly")
\end{verbatim}
\item 承上題，如果把 \texttt{has\_ticket} 換成 \texttt{False}，其他都不變，
\texttt{party-entry.py} 這段程式會印出什麼？內層的 \texttt{if}/\texttt{else} 有沒有
被執行到？
\item 為什麼外層條件不成立時，內層的 \texttt{if}/\texttt{else} 整段都不會被檢查？
\end{drills}

\section{好幾條路選一條：if/elif/else 鏈}
\label{sec:elif-chain}

現實中的判斷常常不只兩條路。用分數給等第就是個好例子：90 分以上是 A、80 到 89 是 B、
70 到 79 是 C，這種情況是好幾個區間選一個，不是「及格/不及格」那種二選一。這時候需要
\texttt{elif}（是 else if 的縮寫）把好幾個條件串成一條鏈。
\index{elif}

\begin{code}{python}{if/elif/else}
if condition1:
    action1
elif condition2:
    action2
else:
    action3
\end{code}

一條 \texttt{if}/\texttt{elif}/\texttt{else} 鏈由上往下檢查，只要有一個條件成立，
就執行那一段、然後直接跳過整條鏈剩下的部分，不會再往下檢查，即使後面的條件其實也成立。
\index{if/elif/else 鏈}

\begin{example}{python}{(grade-letter.py) 依分數給等第}
score = 82

if score >= 90:
    grade = "A"
elif score >= 80:
    grade = "B"
elif score >= 70:
    grade = "C"
elif score >= 60:
    grade = "D"
else:
    grade = "F"

print(grade)
\end{example}

\iosample{}{B \\}

\texttt{score} 是 82，第一個條件 \texttt{score >= 90} 不成立，往下檢查
\texttt{score >= 80}，這個成立，於是 \texttt{grade} 被設成 \texttt{"B"}，接下來
\texttt{elif score >= 70}、\texttt{elif score >= 60} 就算也成立（82 確實也大於等於
70 跟 60），一律被跳過、完全不會執行。

\keypoint{一條 \texttt{elif} 鏈裡，最多只有一段會被執行，第一個成立的那一段。這跟一連串
各自獨立的 \texttt{if} 完全不一樣，接下來直接示範兩者的差別。}

如果把 \texttt{elif} 錯打成一連串各自獨立的 \texttt{if}，會發生什麼事？

\begin{example}{python}{(grade-letter-wrong.py) 誤把 elif 鏈拆成獨立 if}
score = 82

grade = ""
if score >= 90:
    grade = "A"
if score >= 80:
    grade = "B"
if score >= 70:
    grade = "C"

print(grade)
\end{example}

\iosample{}{C \\}

同樣是 82 分，這裡卻印出了 \texttt{C}，不是預期的 \texttt{B}，問題出在每一個
\texttt{if} 都是獨立檢查、互不相關，Python 會把三個條件從頭到尾都檢查一遍：
\begin{itemize}
\item \texttt{score >= 90} 不成立，跳過。
\item \texttt{score >= 80} 成立，\texttt{grade} 被設成 \texttt{"B"}。
\item 接著 \texttt{score >= 70} 也成立，\texttt{grade} 又被改成 \texttt{"C"}，把剛剛
設好的 \texttt{"B"} 蓋掉了。
\end{itemize}

三段獨立的 \texttt{if} 各自都會被檢查、只要成立就
執行，不會因為前一段已經執行過就跳過後面，這跟 \texttt{elif} 鏈「只執行第一個成立的
那一段」剛好相反。

判斷的原則很簡單：
\index{互斥的區間}
\begin{itemize}
\item 如果好幾個條件本來就是「互斥的區間」（一個值不可能同時落在兩個區間裡），像分數
等第這種情況，要用 \texttt{elif} 鏈。
\item 如果好幾個條件彼此獨立、互不相關，每一個都想各自檢查一次，才用一連串各自獨立的
\texttt{if}（例如「分數及不及格」跟「是不是週末」是兩件不相干的事，各自用一個
\texttt{if} 才對）。
\end{itemize}

這一點也可以拿來跟上一節
的巢狀 \texttt{if} 對照：巢狀 \texttt{if} 是「內層要靠外層先成立」，\texttt{elif} 鏈是
「好幾條路彼此互斥、只能走一條」，各自條件之間沒有誰依賴誰的關係。

還有一件事值得特別提醒：\texttt{elif} 鏈最後不一定要接 \texttt{else}。如果你只想在
少數幾種情況下做點什麼，其他情況一律不用理會，可以只寫 \texttt{if} 跟幾個
\texttt{elif}，不接 \texttt{else}。

\begin{example}{python}{(elif-no-else.py) elif 鏈不接 else}
score = 55

if score >= 90:
    print("A")
elif score >= 80:
    print("B")
elif score >= 70:
    print("C")

print("Score processed")
\end{example}

\iosample{}{Score processed \\}

\texttt{score} 是 55，三個 \texttt{elif} 條件全部不成立，因為沒有 \texttt{else}，整條
鏈裡沒有任何一段被執行，直接跳到鏈外面繼續往下執行 \texttt{print("Score processed")}。這跟
本章稍早提到的 \texttt{if}/\texttt{else} 保證「兩條路一定會走其中一條」不一樣：一條
沒接 \texttt{else} 的 \texttt{elif} 鏈，是有可能「一條路都不走」的，要不要接
\texttt{else}，取決於「沒有任何條件成立時，是不是還想做點什麼」這件事。

\begin{drills}[label=drills:elif-chain]
\item 這段程式的輸出是什麼？
\begin{verbatim}
score = 65
if score >= 90:
    grade = "A"
elif score >= 80:
    grade = "B"
elif score >= 60:
    grade = "D"
else:
    grade = "F"
print(grade)
\end{verbatim}
\item 這段程式想做的是「依分數印出等第」，但作者不小心把 \texttt{elif} 全部改成獨立的
\texttt{if}。\texttt{score} 是 95 的話，這段程式實際上會印出什麼？
\begin{verbatim}
score = 95
grade = ""
if score >= 60:
    grade = "Pass"
if score >= 90:
    grade = "Excellent"
print(grade)
\end{verbatim}
\item 你要判斷「今天是不是週末」跟「氣溫是不是超過 30 度」這兩件事，兩者互不相關、
你都想各自知道結果，這種情況該用 \texttt{elif} 鏈，還是兩個各自獨立的 \texttt{if}？
\item \texttt{score = 40}，這段程式沒有接 \texttt{else}，會印出什麼？
\begin{verbatim}
if score >= 90:
    print("A")
elif score >= 80:
    print("B")
print("Done")
\end{verbatim}
\end{drills}

\section{同時看好幾個條件：and、or、not}
\label{sec:logical-operators}

有時候一個判斷需要同時看好幾個條件才能下定論：「兩科都要及格才算真的及格」「只要有一科
拿高分就值得慶祝」。Python 提供三個邏輯運算子，把好幾個「是不是」串成一個更大的條件：
\texttt{and}（兩邊都成立才算成立）、\texttt{or}（有一邊成立就算成立）、\texttt{not}
（把結果反過來）。
\index{邏輯運算子}
\index{and}
\index{or}
\index{not}

\begin{code}{python}{and／or／not}
condition1 and condition2
condition1 or condition2
not condition
\end{code}

先看 \texttt{and}：兩科都及格，才算真的及格。

\begin{example}{python}{(both-subjects-pass.py) 兩科都及格才算}
math_score = 75
english_score = 68

if math_score >= 60 and english_score >= 60:
    print("Both subjects passed")
else:
    print("At least one subject failed")
\end{example}

\iosample{}{Both subjects passed \\}

再看 \texttt{or}：只要有一科拿高分，就值得注意。

\begin{example}{python}{(either-subject-high.py) 任一科拿高分就算}
math_score = 95
english_score = 68

if math_score >= 90 or english_score >= 90:
    print("At least one subject scored high")
else:
    print("Neither subject scored high")
\end{example}

\iosample{}{At least one subject scored high \\}

最後看 \texttt{not}：把一個條件的結果整個反過來。

\begin{example}{python}{(not-passing.py) 用 not 反過來問「沒有及格」}
math_score = 45
print(not (math_score >= 60))
print(math_score < 60)
\end{example}

\iosample{}{True \\ True \\}

\texttt{not (math\_score >= 60)} 的意思就是「math\_score 不是大於等於 60」，效果跟直接
寫 \texttt{math\_score < 60} 一樣，兩行印出的結果完全相同。有些情況（例如條件本身已經
很長）用 \texttt{not} 包起來會比重寫一次相反的比較式更清楚，但也不是每次都比較好讀，
拿捏的原則留到本章最後一節再細談。

除了 \texttt{and}，Python 還提供一種更精簡的寫法，專門用在「一個值是不是落在某個範圍
內」這種情境。假設你想知道兩科分數的差距是不是在 10 分以內（代表兩科表現差不多），有
三種寫法可以做到同一件事，一種比一種精簡：

\begin{code}{python}{同一個判斷的三種寫法}
if 0 < gap:
    if gap < 10:
        ...

if 0 < gap and gap < 10:
    ...

if 0 < gap < 10:
    ...
\end{code}

第一種是你剛剛學過的巢狀 \texttt{if}（\texttt{if} 裡面再包一個 \texttt{if}）；第二種用
\texttt{and} 把兩層壓成一層；第三種是 Python 特有的連鎖比較寫法，直接把兩個比較式接在
一起寫成 \texttt{0 < gap < 10}，效果跟前兩種完全一樣，但更接近你在數學課上寫
「$0 < x < 10$」的方式。實際跑一次，確認三種寫法真的都成立：
\index{連鎖比較}

\begin{example}{python}{(score-gap-check.py) 三段式演進：分數差距是否在 10 分以內}
math_score = 85
english_score = 92
gap = english_score - math_score

if gap > 0:
    if gap < 10:
        print("Nested version: gap within 10 points")

if gap > 0 and gap < 10:
    print("Flattened with and: gap within 10 points")

if 0 < gap < 10:
    print("Chained comparison: gap within 10 points")
\end{example}

\iosample{}{Nested version: gap within 10 points \\ Flattened with and: gap within 10 points \\ Chained comparison: gap within
10 points \\}

\stylepoint{碰到「一個值是不是落在某個範圍內」這種判斷，直接寫連鎖比較
\texttt{0 < gap < 10}，比巢狀 \texttt{if} 或 \texttt{and} 更精簡、也更接近你原本的
數學直覺，這是 Python 少數幾個其他語言通常沒有的方便寫法。}

\begin{drills}[label=drills:logical-operators]
\item \texttt{age = 16}，這個條件式的結果是 \texttt{True} 還是 \texttt{False}？
\texttt{age >= 12 and age < 18}
\item \texttt{score = 45}，這個條件式的結果是 \texttt{True} 還是 \texttt{False}？
\texttt{score < 60 or score == 100}
\item 把 \texttt{score > 0 and score < 60} 改寫成連鎖比較的寫法會是什麼？
\end{drills}

\section{and／or 不見得會兩邊都看：短路求值}
\label{sec:short-circuit}

上一節看到的 \texttt{and}／\texttt{or}，有一個很多人第一次接觸時會覺得意外的行為：
Python 不一定會把兩邊的條件都算過一遍才給你答案。這叫短路求值：如果從左邊的條件就已經
能確定整條式子的結果，Python 會直接把答案交出來，右邊那個條件根本不會被執行到。
\index{短路求值}

\begin{code}{python}{短路求值的判斷順序}
condition1 and condition2   # condition1 是 False，condition2 會被跳過
condition1 or condition2    # condition1 是 True，condition2 會被跳過
\end{code}

\texttt{and} 的邏輯是「兩邊都成立才算成立」，所以只要左邊已經是 \texttt{False}，不管
右邊是什麼，整條式子注定是 \texttt{False}，Python 沒有理由浪費力氣去算右邊。同樣地，
\texttt{or} 的邏輯是「有一邊成立就算成立」，只要左邊已經是 \texttt{True}，整條式子注定
是 \texttt{True}，右邊也會被跳過。

這個行為不只是省一點執行時間而已，它還有一個很實用的地方：可以用來避免一個本來會出錯的
操作真的被執行到。想像你想知道一個字串是不是以「A」開頭，如果這個字串本身是空的，你根本
不可能去看它的第一個字元，硬要看反而會出錯。

\begin{example}{python}{(starts-with-a-broken.py) 沒有先確認字串是不是空的}
name = ""
if name[0] == "A":
    print("Starts with A")
else:
    print("Doesn't start with A")
\end{example}

執行後會看到：

\begin{verbatim}
Traceback (most recent call last):
  File "starts-with-a-broken.py", line 2, in <module>
    if name[0] == "A":
       ~~~~^^^
IndexError: string index out of range
\end{verbatim}

\texttt{name} 是空字串，根本沒有第 0 個字元，\texttt{name[0]} 這一步自己就先出錯了，
連比較都還沒開始比。第 \ref{ch:strings} 章學過的 \texttt{IndexError} 在這裡又出現了，只是這次不是索引
超出字串長度，而是字串本身長度就是 0，任何索引都超出範圍。用短路求值可以先把「字串是不是
空的」擋在前面：

\begin{example}{python}{(starts-with-a-fixed.py) 用 and 先確認字串不是空的}
name = ""
if name and name[0] == "A":
    print("Starts with A")
else:
    print("Empty string, or doesn't start with A")
\end{example}

\iosample{}{Empty string, or doesn't start with A \\}

這裡的 \texttt{name} 是空字串，在 \texttt{if} 裡會被當成 \texttt{False}（下一節會完整
解釋這件事），所以 \texttt{name and name[0] == "A"} 左邊已經是 \texttt{False}，短路
求值讓右邊的 \texttt{name[0] == "A"} 根本不會被執行，也就不會有機會出錯。

順序在這裡很
重要：一定要把「用來擋錯誤的條件」寫在左邊，才能靠短路求值保護右邊；如果順序寫反，變成
\texttt{name[0] == "A" and name}，Python 一樣由左至右求值，容易出錯的那一步反而會先
被執行、先出錯，短路求值完全幫不上忙。

\texttt{or} 也有一樣的效果，只是換成「左邊已經是 \texttt{True} 就不看右邊」：

\begin{example}{python}{(starts-with-h-or.py) 用 or 提前判斷}
text = ""
if not text or text[0] == "H":
    print("Empty string, or starts with H")
else:
    print("Has content, but doesn't start with H")
\end{example}

\iosample{}{Empty string, or starts with H \\}

\texttt{text} 是空字串，\texttt{not text} 的結果是 \texttt{True}（空字串本身是
\texttt{False}，\texttt{not} 把它反過來變成 \texttt{True}），左邊已經是 \texttt{True}，
\texttt{or} 直接判定整條式子成立，右邊的 \texttt{text[0] == "H"} 完全不會被執行，同樣
避開了原本會出現的 \texttt{IndexError}。

\keypoint{短路求值是 Python 允許你把「先確認安全，再做可能出錯的事」寫在同一行的方式，
不只是「順便發生」的小細節，重點是順序：確認安全的那個條件一定要放在左邊。之後你會學到
Python 針對「這件事可能會失敗」還有另一種更專門的處理方式，等學到那裡，你可以回頭比較
兩者的差異。}

\begin{drills}[label=drills:short-circuit]
\item \texttt{score = 0}，這個條件式的結果是 \texttt{True} 還是 \texttt{False}？
執行過程中，右邊的 \texttt{100 / score > 1} 有沒有真的被算出來？
\texttt{score != 0 and 100 / score > 1}
\item \texttt{name = ""}，如果把 \texttt{name and name[0] == "A"} 的順序寫反，變成
\texttt{name[0] == "A" and name}，還能不能避免 \texttt{IndexError}？為什麼？
\item \texttt{has\_data = True}，\texttt{count = 5}，這個條件式的結果是
\texttt{True} 還是 \texttt{False}？\texttt{has\_data or count > 100}
\end{drills}

\section{把好幾個條件寫得清楚：可讀性小技巧}
\label{sec:readable-conditions}

\texttt{and}、\texttt{or}、\texttt{not} 一旦混在同一行裡用，很容易碰到一個新問題：
Python 判斷這些運算子的順序，不一定跟你讀這句話時腦中浮現的順序一樣。Python 會先看
\texttt{not}，再看 \texttt{and}，最後才看 \texttt{or}，這個優先順序跟英文文法裡
「先乘除、後加減」的數學規則類似，但很少有人真的把它背起來，寫程式的時候也很容易忘記，
結果寫出跟自己原本想的不一樣的條件。

\begin{example}{python}{(party-entry-precedence.py) 優先順序造成的判斷結果不同}
membership = True
has_invitation = False
age = 16

if membership or has_invitation and age >= 18:
    print("Entry allowed (no parens)")
else:
    print("Entry denied (no parens)")

if (membership or has_invitation) and age >= 18:
    print("Entry allowed (with parens)")
else:
    print("Entry denied (with parens)")
\end{example}

\iosample{}{Entry allowed (no parens) \\ Entry denied (with parens) \\}

同樣的三個變數，兩段判斷卻給出不一樣的答案。原意是「要嘛是會員、要嘛有邀請函，而且必須
年滿 18 歲才能入場」。加了括號的版本，也就是 \texttt{(membership or has\_invitation)}
搭配 \texttt{and age >= 18}，才真的照這個順序做事：先看是不是符合資格，再看年齡。

但如果沒加括號，直接寫成 \texttt{membership or has\_invitation and age >= 18}，
Python 會先算 \texttt{and}。這一整條式子等於自動被理解成
\texttt{has\_invitation and age >= 18} 先算出一個結果，再拿去跟 \texttt{membership}
做 \texttt{or}。這裡 \texttt{membership} 是 \texttt{True}，\texttt{or} 左邊已經成立，
短路求值直接判定整條式子成立，完全沒有機會檢查年齡，未成年的會員因此被誤放行了。

\stylepoint{只要一行條件式裡同時用到 \texttt{and} 跟 \texttt{or}，一律加括號寫清楚你要
的優先順序，不要依賴預設的優先權規則。加括號幾乎不花額外力氣，卻能讓自己跟之後看程式碼
的人都不用去記、也不用去猜 Python 到底會先算哪一半。}

除了加括號，另一個讓長條件變好讀的辦法，是把每一段子條件拆成一個有名字的布林變數，讓
最後那行 \texttt{if} 讀起來幾乎像一句白話：

\begin{example}{python}{(party-entry-named.py) 拆成有名字的布林變數}
membership = True
has_invitation = False
age = 16

is_adult = age >= 18
qualifies_by_status = membership or has_invitation

if qualifies_by_status and is_adult:
    print("Entry allowed")
else:
    print("Entry denied")
\end{example}

\iosample{}{Entry denied \\}

這段程式先把「符不符合資格」跟「夠不夠年齡」各自算成一個獨立的 \texttt{bool} 值，存進
名字清楚的變數，最後 \texttt{if qualifies\_by\_status and is\_adult:} 這一行完全不用
擔心運算子的優先順序，因為每個變數本身已經是算好的 \texttt{True}／\texttt{False}，
彼此之間只剩一個 \texttt{and}，不會再有混著 \texttt{or} 的判讀問題。

條件愈複雜，這種
拆解方式的好處愈明顯：與其在腦中一次解一整條長算式，不如先解決一小段、給它一個名字，
最後再把幾個名字組起來。

\begin{drills}[label=drills:readable-conditions]
\item 有一段折扣判斷式想表達的是「會員或持有優惠券，而且消費滿 100 元才能打折」，作者
寫成 \texttt{is\_member or has\_coupon and total >= 100}，沒有加括號。
\texttt{is\_member = True}，\texttt{has\_coupon = False}，\texttt{total = 50} 的情況
下，這樣寫算出來的結果，跟原本想表達的意思一不一樣？
\item 承上題，把條件式改成 \texttt{(is\_member or has\_coupon) and total >= 100}，
同樣的三個變數值，結果會不會不一樣？為什麼？
\item 把 \texttt{not (age >= 18 and has\_ticket)} 拆成兩個有名字的布林變數
（例如 \texttt{is\_adult}、\texttt{has\_ticket}）再組合，你會怎麼寫？
\end{drills}

\section{空的東西當作沒有}
\label{sec:truthy-falsy}

目前為止，你放進 \texttt{if} 的條件都是比較運算式的結果，答案不是 \texttt{True} 就是
\texttt{False}。但 Python 其實還有一個特別的脾氣：數字、字串這些值本身也可以直接放進
\texttt{if} 裡，不用先寫成比較式。\texttt{0}、空字串 \texttt{""} 會被當成
\texttt{False}，非零的數字、有內容的字串會被當成 \texttt{True}。
\index{truthy/falsy}

\begin{code}{python}{把值直接放進 if}
if x:
    action
\end{code}

先看數字的情況：

\begin{example}{python}{(bonus-check.py) 有沒有加分}
bonus_points = 0

if bonus_points:
    print("Bonus points:", bonus_points)
else:
    print("No bonus points")
\end{example}

\iosample{}{No bonus points \\}

\texttt{bonus\_points} 是 \texttt{0}，放進 \texttt{if} 裡直接被當成 \texttt{False}，
跟寫 \texttt{if bonus\_points != 0:} 的效果一樣，但不用你自己動手比較。再看字串的情況：

\begin{example}{python}{(name-check.py) 使用者有沒有輸入名字}
name = ""

if name:
    print("Hello, " + name)
else:
    print("You haven't entered a name")
\end{example}

\iosample{}{You haven't entered a name \\}

空字串 \texttt{""} 一樣被當成 \texttt{False}；如果 \texttt{name} 是 \texttt{"Alice"}
這種有內容的字串，\texttt{if name:} 就會是 \texttt{True}。

\stylepoint{比起特地寫一段程式碼去檢查「這個字串是不是空的」，Python 慣用的寫法直接寫
\texttt{if name:} 就好，空的東西自然就是 \texttt{False}，不用你多此一舉再比一次。}

很多人第一次看到「非零的數字都是 \texttt{True}」，會直覺以為負數比較特別，可能被當成
\texttt{False}，其實 Python 的規則很單純：只有 \texttt{0}（不管是整數的 \texttt{0}
還是後面會學到的小數 \texttt{0.0}）才是 \texttt{False}，非零的數字不管是正是負，一律
都是 \texttt{True}。

\begin{example}{python}{(temperature-change-check.py) 負數也算是有內容}
temperature_change = -3

if temperature_change:
    print("Temperature changed")
else:
    print("Temperature unchanged")
\end{example}

\iosample{}{Temperature changed \\}

\texttt{temperature\_change} 是 \texttt{-3}，雖然是負數，但它不是 \texttt{0}，放進
\texttt{if} 裡一樣被當成 \texttt{True}，所以印出「Temperature changed」。判斷的關鍵永遠只有一個：
「是不是等於 0（或空的）」，跟正負號、跟數字大小完全無關。

這個「空的當 \texttt{False}、有東西當 \texttt{True}」的脾氣，不只數字、字串適用，之後
你會學到的容器型別，空的時候也是一樣的邏輯，到時候你會發現這個習慣完全用得上。

這也是
本章多次提到、卻還沒完全講開的伏筆：Python 除了用 \texttt{if} 檢查「東西存不存在、
是不是空的」，之後你還會學到另一套更 Python 的做法，直接先試著去用一個東西，等真的
出問題了再處理，而不是每次動手之前都先檢查一輪。

這一章看到的 truthy／falsy 判斷跟短路
求值，都是幫你先打好地基，等真正學到那套做法時，會更容易感覺到兩者的差異跟各自適合
的場合。

\begin{drills}[label=drills:truthy-falsy]
\item \texttt{count = 0}，\texttt{if count:} 裡面那段程式碼會不會被執行？
\item \texttt{message = "Done"}，\texttt{if message:} 裡面那段程式碼會不會被執行？
\item \texttt{score = -10}，\texttt{if score:} 裡面那段程式碼會不會被執行？
\end{drills}

\section*{練習解答}
\addcontentsline{toc}{section}{練習解答}

\paragraph{第~\ref{sec:comparison-operators}~節}
\begin{enumerate}
\item \texttt{True}，120 大於等於 120。
\item \texttt{False}，\texttt{"Alice"} 跟 \texttt{"alice"} 大小寫不同，\texttt{==} 比較
字串時會連大小寫一起比，兩者不相等。
\end{enumerate}

\paragraph{第~\ref{sec:if-basics}~節}
\begin{enumerate}
\item 會印出：Done。\texttt{score > 90} 不成立，所以「Excellent」不會被印出，但
\texttt{print("Done")} 沒有縮排在 \texttt{if} 裡面，不管條件成不成立都會執行。
\item \texttt{SyntaxError: expected ':'}，\texttt{if} 那一行結尾漏了冒號，補上
\texttt{if temperature > 28:} 即可。
\item 把 \texttt{=} 誤打成該用 \texttt{==} 的地方，把 \texttt{if temperature = 28:}
改成 \texttt{if temperature == 28:} 就能修好（不過原意應該是「超過 28 度」，可能要改成
\texttt{if temperature > 28:} 才對）。
\end{enumerate}

\paragraph{第~\ref{sec:if-else}~節}
\begin{enumerate}
\item 會印出：Pass。\texttt{score >= 60} 成立（60 大於等於 60）。
\item 少了冒號。\texttt{if score >= 60} 那一行結尾要補上冒號，變成
\texttt{if score >= 60:}，程式才會正常執行。
\end{enumerate}

\paragraph{第~\ref{sec:nested-if}~節}
\begin{enumerate}
\item 會印出：Can go out。外層 \texttt{weather == "Rainy"} 成立，往內層走，
\texttt{has\_umbrella} 是 \texttt{True}，印出「Can go out」。
\item 會印出：No ticket, entry denied。內層的 \texttt{if}/\texttt{else}（判斷年齡的那一段）
完全沒有被執行到，因為外層 \texttt{has\_ticket} 已經是 \texttt{False}，Python 直接
跳到最外層的 \texttt{else}。
\item 因為內層的 \texttt{if}/\texttt{else} 整段都縮排在外層 \texttt{if} 的區塊裡，是
外層條件成立時才會做的「動作」；外層條件一旦不成立，這整個縮排區塊（不管裡面有幾層）都
會被跳過，不會有任何一部分被執行。
\end{enumerate}

\paragraph{第~\ref{sec:elif-chain}~節}
\begin{enumerate}
\item 會印出：D。\texttt{score >= 90} 不成立，\texttt{score >= 80} 不成立，
\texttt{score >= 60} 成立（65 大於等於 60），\texttt{grade} 被設成 \texttt{"D"}，
之後不會再檢查 \texttt{else}。
\item 會印出：Excellent。兩個 \texttt{if} 都是獨立檢查：\texttt{score >= 60} 成立，
\texttt{grade} 先變成 \texttt{"Pass"}；接著 \texttt{score >= 90} 也成立，
\texttt{grade} 又被改成 \texttt{"Excellent"}，蓋掉前面的結果。
\item 兩個各自獨立的 \texttt{if}，兩件事彼此無關，不是同一組互斥的區間，各自檢查、
各自印出結果才對，用 \texttt{elif} 鏈反而會讓其中一個判斷被跳過。
\item 只會印出：Done。\texttt{score} 是 40，兩個條件都不成立，這條 \texttt{elif} 鏈
沒有接 \texttt{else}，所以沒有任何一段被執行，直接往下跳到 \texttt{print("Done")}。
\end{enumerate}

\paragraph{第~\ref{sec:logical-operators}~節}
\begin{enumerate}
\item \texttt{True}，16 大於等於 12，也小於 18，兩個條件都成立，\texttt{and} 的結果是
\texttt{True}。
\item \texttt{True}，45 小於 60，\texttt{score < 60} 已經成立，\texttt{or} 只要有一邊
成立就是 \texttt{True}，不用管另一邊。
\item \texttt{0 < score < 60}。
\end{enumerate}

\paragraph{第~\ref{sec:short-circuit}~節}
\begin{enumerate}
\item \texttt{True}。\texttt{score != 0} 已經成立，\texttt{and} 左邊是
\texttt{True}，還不能確定整條式子的結果，所以右邊的 \texttt{100 / score > 1} 確實會
被算出來（這裡跟左邊已經是 \texttt{False} 就會跳過右邊的情況不一樣，因為左邊是
\texttt{True} 時 \texttt{and} 沒辦法提早下結論，一定要看右邊）。
\item 不能。\texttt{and} 由左至右求值，把容易出錯的比較放在左邊，等於在還沒機會被短路
之前就已經先執行、也已經先出錯，順序寫反，短路求值完全派不上用場。
\item \texttt{True}。\texttt{has\_data} 已經是 \texttt{True}，\texttt{or} 左邊已經
成立，直接判定整條式子成立，右邊的 \texttt{count > 100} 不會被執行到。
\end{enumerate}

\paragraph{第~\ref{sec:readable-conditions}~節}
\begin{enumerate}
\item 不一樣。因為 \texttt{and} 的優先順序比 \texttt{or} 高，這行等於
\texttt{is\_member or (has\_coupon and total >= 100)}，\texttt{is\_member} 已經是
\texttt{True}，短路求值讓 \texttt{or} 直接判定整條式子成立，完全沒有檢查到
\texttt{total} 夠不夠 100，跟原本「消費滿 100 元才能打折」的意思不一樣。
\item 會不一樣。加上括號後，Python 會先算 \texttt{(is\_member or has\_coupon)}
得到 \texttt{True}，再跟 \texttt{total >= 100}（\texttt{False}，因為 50 不到 100）
做 \texttt{and}，整條式子變成 \texttt{False}，這次才真的照原意判斷出「沒有打折」。
\item 例如寫成 \texttt{is\_adult = age >= 18}、接著
\texttt{if not (is\_adult and has\_ticket):}，或是進一步拆成 \texttt{qualifies =
is\_adult} \texttt{and has\_ticket} 再寫 \texttt{if not qualifies:}，效果都跟原本
那一整行相同，但讀起來更清楚每一段在檢查什麼。
\end{enumerate}

\paragraph{第~\ref{sec:truthy-falsy}~節}
\begin{enumerate}
\item 不會，\texttt{0} 在 \texttt{if} 裡被當成 \texttt{False}。
\item 會，\texttt{"Done"} 是非空字串，在 \texttt{if} 裡被當成 \texttt{True}。
\item 會，\texttt{-10} 雖然是負數，但不等於 \texttt{0}，在 \texttt{if} 裡一樣被當成
\texttt{True}。
\end{enumerate}

\chapter{迴圈}
\label{ch:loops}
\index{迴圈}

生活裡有兩種「重複做同一件事」：

\begin{itemize}
\item 第一種你一開始就知道要做幾次，摺完一籃衣服，你知道今天總共要摺幾件；打掃前，你
知道家裡有幾扇窗戶要擦。
\item 第二種你不知道要做幾次，只知道「做到某個狀態出現為止」，洗碗要洗到碗槽空了才能
停，存錢要存到存夠目標金額才能停，你沒辦法在動手之前就先算出精確的次數。
\end{itemize}

這兩種重複，光靠你之前學過的工具都解決不了，你不可能把\texttt{print} 打五十次來擦五十
扇窗戶，更不可能在動手存錢之前，先算好到底要存幾天。這一章要教的，就是 Python 怎麼分別
處理這兩種重複，讓你只寫一次程式碼，就能重複執行任意多次。

這一章分成幾段推進：先教「已知次數」的 \texttt{for}，再教「不知道要做幾次、要看狀態
才知道該不該繼續」的 \texttt{while}。中間會穿插兩個實用的補充技巧：怎麼在一個迴圈裡面
再放一個迴圈，處理「每一組裡的每一個」這種雙層重複；以及怎麼在走訪一串東西的同時，順便
拿到「這是第幾個」的位置資訊。

接著回頭處理幾件更細節的事：怎麼在迴圈進行到一半時喊停或
跳過某一輪，以及怎麼知道一個迴圈到底是正常跑完、還是被提早喊停結束的。最後，把前面所有
例子拉高一層來看，你會發現不管是擦窗戶、存錢還是翻抽屜找鑰匙，背後其實只有三種常見的
迴圈任務「長相」：累加、計數、搜尋。

這三種長相會在你之後寫的幾乎每一支程式裡反覆出現，
提早認出它們的骨架，之後遇到新的問題時，會更容易想到「這其實就是搜尋的模式」，不用每次
都從零開始想解法。這一章跟前面幾章一樣是地基章節，技巧點不少，但每一個都值得你自己動手
重現一次，光是看得懂並不算數。

\section{已知次數的重複：\texttt{for} 搭配 \texttt{range()}}
\label{sec:for-range}

先從最單純的情境開始：你動手之前就知道要做幾次，不涉及任何「這次要不要繼續」的判斷，
只是把同一個動作重複套用固定的次數。Python 用 \texttt{for} 搭配一個叫 \texttt{range()}
的工具處理這種情境。\texttt{range(n)} 會產生一串數字，讓 \texttt{for} 迴圈依序拿去用，
剛好重複 \texttt{n} 次。
\index{for（搭配 range()）}

\begin{code}{python}{for 搭配 range() 重複已知次數}
for i in range(n):
    repeated_action
\end{code}

\texttt{range(n)}裡的\texttt{n}就是你想重複的次數；\texttt{i}是每一輪迴圈都會拿到的
一個數字，你可以在迴圈裡用它，也可以完全不用它。重點是縮排在\texttt{for}底下的那幾行，
會被原封不動地重複執行\texttt{n}次。

\begin{example}{python}{(windows.py) 擦五扇窗戶}
for i in range(5):
    print(f"Cleaning window {i + 1}")

print("All windows cleaned!")
\end{example}

\iosample{}{Cleaning window 1 \\ Cleaning window 2 \\ Cleaning window 3 \\ Cleaning window 4
\\ Cleaning window 5 \\ All windows cleaned! \\}

留意這裡多寫了一個\texttt{+ 1}。這不是隨便加的，是為了修正一個新手幾乎都會踩到的坑。
把上面例子裡的\texttt{i + 1}換回單純的\texttt{i}，實際印出來的會是這樣：

\begin{example}{python}{(range-check.py) 直接印出 i}
for i in range(5):
    print(i)
\end{example}

\iosample{}{0 \\ 1 \\ 2 \\ 3 \\ 4 \\}

\keypoint{\texttt{range(5)}確實會讓迴圈跑滿 5 次，但拿到的數字是從 0 開始、到 4 為止，
不是從 1 到 5。很多人第一次看到會誤以為最後應該印出 5，其實 5 從來沒有出現過。這是
\texttt{range()}最容易讓人誤會的地方：次數是對的，起點卻不是你直覺以為的 1。想要顯示
「第幾次」給人看的時候（像上面擦窗戶的例子），記得自己補一個\texttt{+ 1}。}

除了只給一個數字，\texttt{range()}還可以指定起點、終點、跟每次要跳幾格（間隔），間隔
也可以是負的，讓數字倒著數。

\begin{code}{python}{range() 指定起點、終點與間隔}
for i in range(start, stop, step):
    repeated_action
\end{code}

如果只給兩個數字\texttt{range(start, stop)}，Python 會假設間隔是 1；如果間隔是負的，
數字就會從\texttt{start}開始往下數，一路數到比\texttt{stop}大一點點的地方為止。

\begin{example}{python}{(squats.py) 深蹲倒數}
for remaining in range(10, 0, -1):
    print(f"Reps remaining: {remaining}")

print("Squats complete!")
\end{example}

\iosample{}{Reps remaining: 10 \\ Reps remaining: 9 \\ Reps remaining: 8 \\ Reps remaining: 7
\\ Reps remaining: 6 \\ Reps remaining: 5 \\ Reps remaining: 4 \\ Reps remaining: 3 \\
Reps remaining: 2 \\ Reps remaining: 1 \\ Squats complete! \\}

這裡\texttt{range(10, 0, -1)}是「從 10 開始，每次減 1，數到比 0 大的地方為止」，所以
最後印出來的是 1，不是 0。跟前面\texttt{range(n)}的道理一樣，終點那個數字本身永遠不會
被包含進去。

\begin{drills}[label=drills:for-range]
\item 不要真的執行，用眼睛看：\texttt{for i in range(3): print(i)}會印出哪三行？
\item \texttt{for i in range(2, 9, 2): print(i)}最後一行印出的數字是多少？
\item 填空：想要一個迴圈從 9 開始、每次減 2、數到 1（含）為止，\texttt{range()}該怎麼
填？（提示：終點那個數字不會被包含進去，所以終點要填比 1 還小的數字。）
\item \texttt{for i in range(5, 5): print(i)}這段程式碼，迴圈本體總共會執行幾次？
\end{drills}

\section{對一串東西的每一個都做一次：\texttt{for} 走訪字串}
\label{sec:for-string}

前一節的\texttt{for}是「把同一個動作重複做\texttt{n}次」，這一節要教的\texttt{for}
語感不太一樣，重點換成「對一串東西裡的每一個都做一次」。這串東西可以是你在第 \ref{ch:strings} 章學過的
字串。字串本身就是一串字元排在一起，\texttt{for}可以直接一個一個走訪它們，不用先算出
字串有幾個字元、也不用自己管索引數字。
\index{for 走訪序列（例如字串）}

\begin{code}{python}{for 走訪字串的每個字元}
for char in text:
    repeated_action
\end{code}

每一輪迴圈，\texttt{char}會依序變成\texttt{text}裡的下一個字元，直到走完整個字串為止。

\begin{example}{python}{(count-spaces.py) 數一句話裡有幾個空白字元}
sentence = "nice weather today"
space_count = 0

for char in sentence:
    if char == " ":
        space_count += 1

print(f"This sentence has {space_count} spaces")
\end{example}

\iosample{}{This sentence has 2 spaces \\}

留意這裡的\texttt{space\_count += 1}，這是\texttt{space\_count = space\_count + 1}
的簡寫，意思是「把\texttt{space\_count}加 1 的結果，重新貼回\texttt{space\_count}這個
名字牌」，呼應第 \ref{ch:variables} 章教過的名字牌模型。之後你會常看到的\texttt{-=}則是減法版本，效果
一樣是簡寫，不是新的運算方式。
\index{複合賦值運算子（+=／-=）}

\begin{example}{python}{(check-letter.py) 檢查一個字裡有沒有特定字母}
word = "python"
target = "y"
found = False

for letter in word:
    if letter == target:
        found = True

if found:
    print(f"The word '{word}' contains the letter '{target}'")
else:
    print(f"The word '{word}' does not contain the letter '{target}'")
\end{example}

\iosample{}{The word 'python' contains the letter 'y' \\}

\stylepoint{如果你只是想拿到字元本身，直接\texttt{for char in word}走訪就好，不需要先
繞一圈去算索引數字（例如\texttt{for i in range(len(word))}再用\texttt{word[i]}取出
字元）。除非你真的需要知道「這是第幾個字元」這個位置資訊，不然直接走訪內容本身才是更
Python 的寫法；如果真的需要位置資訊，本章稍後會介紹更好的做法（\texttt{enumerate()}），
不用回頭走\texttt{range(len(word))}這條路。}

之後你學到 list，會發現這個\texttt{for}寫法幾乎不用改，list 走訪的方式跟字串一模一樣，
只是走訪的對象換成另一種容器。

\begin{drills}[label=drills:for-string]
\item 不要真的執行，用眼睛看：\texttt{for char in "cat": print(char.upper())}會印出
哪三行？
\item 抓蟲：下面這段程式碼執行後會出現錯誤，錯誤訊息會是什麼？該怎麼修？
\begin{verbatim}
for char in "abc"
    print(char)
\end{verbatim}
\end{drills}

\section{迴圈裡面還有迴圈：巢狀迴圈}
\label{sec:nested-loops}

前面兩節的\texttt{for}走訪的都是「一層」的重複，不管是固定次數還是走訪字串，每一輪
迴圈做的都是一件單純的事。但生活裡的重複性任務常常不只一層。

打掃房子的時候，你面對的
其實是「對每一個房間，重複擦每一扇窗戶」這種雙層任務：外層是一個一個房間輪流處理，內層
是處理每個房間時，一扇一扇窗戶擦過去。這種「迴圈裡面還有一個迴圈」的寫法，叫做巢狀迴圈，
寫法就是把一個\texttt{for}迴圈整段，直接放進另一個\texttt{for}迴圈的縮排底下。
\index{巢狀迴圈}

\begin{code}{python}{巢狀迴圈：for 裡面再放一個 for}
for outer_item in outer_sequence:
    repeated_action
    for inner_item in inner_sequence:
        repeated_action
\end{code}

外層迴圈每跑一輪，內層迴圈就會完整跑過一次自己全部的輪次；外層跑幾輪，內層那一整套動作
就會被重複幾次。縮排在這裡特別重要：內層\texttt{for}縮排比外層多一層，Python 才認得出
「這是外層迴圈裡面的東西」，不是外層迴圈外面的獨立程式碼。

\begin{example}{python}{(clean-rooms.py) 打掃每個房間的每一扇窗戶}
for room in range(1, 4):
    print(f"Entering room {room}")
    for window in range(1, 3):
        print(f"  Cleaning room {room}, window {window}")
\end{example}

\iosample{}{Entering room 1 \\ ~~Cleaning room 1, window 1 \\ ~~Cleaning room 1, window 2
\\ Entering room 2 \\ ~~Cleaning room 2, window 1 \\ ~~Cleaning room 2, window 2
\\ Entering room 3 \\ ~~Cleaning room 3, window 1 \\ ~~Cleaning room 3, window 2 \\}

外層\texttt{room}先變成 1，進入第 1 個房間之後，內層\texttt{window}迴圈完整跑了 1、2
兩輪，把第 1 個房間的兩扇窗戶都擦完，才會回到外層，讓\texttt{room}變成 2，重新進入內層
迴圈再跑一次。整個迴圈總共會執行 3 乘 2 等於 6 次擦窗戶的動作，不是 3 加 2 次。

巢狀迴圈有一個新手很容易誤會的地方：\texttt{break}或\texttt{continue}只會作用在離它
最近的那一層迴圈上，不會連帶影響外層迴圈。

\begin{example}{python}{(clean-rooms-break.py) break 只跳出內層迴圈}
for room in range(1, 3):
    print(f"Cleaning room {room}")
    for window in range(1, 4):
        if window == 2:
            break
        print(f"  Finished window {window}")
    print(f"Room {room} cleaning done")
\end{example}

\iosample{}{Cleaning room 1 \\ ~~Finished window 1 \\ Room 1 cleaning done \\ Cleaning room 2
\\ ~~Finished window 1 \\ Room 2 cleaning done \\}

留意每個房間都只印出「Finished window 1」，第 2、3 扇窗戶完全沒被擦到，因為\texttt{window}
一等於 2，\texttt{break}就讓內層\texttt{for}迴圈提早結束了。但外層\texttt{for}迴圈完全
沒受影響，兩個房間都照常輪到，各自印出「cleaning done」。如果你原本期待\texttt{break}能讓整個
巢狀迴圈（包含外層）都停下來，這裡就是一個容易踩到的坑。

\keypoint{\texttt{break}／\texttt{continue}的作用範圍，只到離它最近的那一層迴圈為止。
想讓外層迴圈也跟著提早結束，得在外層迴圈自己的層級再加一次判斷（例如另外設一個變數記錄
「要不要繼續跑外層」），單靠內層的\texttt{break}辦不到。}

巢狀迴圈也可以搭配第~\ref{sec:accumulate-count}~節會正式介紹的累加骨架一起用：外層跟
內層各跑各的輪次，但用同一個變數，把每一輪的結果一路加總起來。

\begin{example}{python}{(total-windows.py) 算出整棟房子總共有幾扇窗戶}
total_windows = 0

for room in range(1, 4):
    for window in range(1, 3):
        total_windows += 1

print(f"The whole house has {total_windows} windows total")
\end{example}

\iosample{}{The whole house has 6 windows total \\}

這裡\texttt{total\_windows}在兩層迴圈外面先初始化成 0，不管走到外層第幾個房間、內層第
幾扇窗戶，每一輪都對同一個\texttt{total\_windows}做\texttt{+= 1}，兩層迴圈跑完之後，
\texttt{total\_windows}就會是 3 個房間乘 2 扇窗戶，總共 6 扇窗戶的總和，不用另外拆開來
一個房間一個房間手動加。

\begin{drills}[label=drills:nested-loops]
\item 不要真的執行，用眼睛看：下面這段程式碼總共會印出幾行「check」？
\begin{verbatim}
for i in range(2):
    for j in range(3):
        print("check")
\end{verbatim}
\item 承接上面\texttt{clean-rooms.py}：如果外層\texttt{range(1, 4)}改成
\texttt{range(1, 6)}，內層\texttt{range(1, 3)}不變，整個迴圈總共會印出幾次「Cleaning」
開頭的訊息？
\item 抓蟲：如果你想讓\texttt{clean-rooms-break.py}裡的\texttt{break}讓兩層迴圈都一起
結束，只在內層寫\texttt{break}是不夠的，你還需要什麼？（不用寫出完整程式碼，說明思路
即可）
\end{drills}

\section{同時要位置又要內容：\texttt{enumerate()} 搭配 \texttt{for}}
\label{sec:enumerate}

第~\ref{sec:for-string}~節教過，\texttt{for char in word}這種寫法能直接拿到字串裡的
每個字元本身，不用自己管索引數字。但有些情況你兩個都要：不只要知道「這個字元是什麼」，
還要知道「它排在第幾個位置」。例如你想知道某個字母第一次出現在字串的哪個位置，只走訪
字元本身是不夠的，你還得同時追蹤位置。Python 提供一個現成的工具處理這種「內容跟位置都
要」的情境，叫\texttt{enumerate()}。
\index{enumerate()}

\begin{code}{python}{for 搭配 enumerate() 同時取得位置與內容}
for index, item in enumerate(sequence):
    repeated_action
\end{code}

\texttt{enumerate(sequence)}把\texttt{sequence}包成一種每一輪都會吐出「位置、內容」
這一組配對的東西，\texttt{for}迴圈這裡一次接住兩個名字：\texttt{index}是這一輪走到第幾
個位置（從 0 開始算，跟\texttt{range()}的起點一樣），\texttt{item}則是這個位置上的內容
本身。之後在講 tuple 的那一章，你會學到這種一次接住好幾個值的寫法背後的完整機制，這裡
先掌握「\texttt{enumerate()}同時給你位置跟內容」這個用法就夠了。

\begin{example}{python}{(find-letter.py) 找出某個字母第一次出現的位置}
word = "banana"
target = "n"

for index, char in enumerate(word):
    if char == target:
        print(f"'{target}' first appears at index {index}")
        break
\end{example}

\iosample{}{'n' first appears at index 2 \\}

迴圈依序走到索引 0（\texttt{b}）、1（\texttt{a}）都不符合，走到索引 2（\texttt{n}）
才符合\texttt{target}，印出訊息後就用\texttt{break}提早結束，後面的索引 3、4、5 都不會
再檢查。留意這裡把\texttt{enumerate()}跟第~\ref{sec:break-search}~節的線性搜尋合在
一起用：一邊走訪、一邊記錄位置，一找到符合的字元就用\texttt{break}提早結束，是很常見的
組合。

\begin{example}{python}{(enumerate-oops.py) 忘記同時接住兩個值}
word = "python"

for index in enumerate(word):
    print(index)
\end{example}

\iosample{}{(0, 'p') \\ (1, 'y') \\ (2, 't') \\ (3, 'h') \\ (4, 'o') \\ (5, 'n') \\}

如果只寫一個名字\texttt{index}去接\texttt{enumerate(word)}吐出來的東西，Python 不會
報錯，但\texttt{index}接到的其實是「位置跟內容綁在一起」的一整組配對，印出來會是像
\texttt{(0, 'p')}這樣一組一組的東西，跟原本想要的乾淨數字或字元不一樣。想要兩個分開拿，
就得像前面\texttt{find-letter.py}那樣，一次寫兩個名字\texttt{index}、\texttt{char}
去接。

\stylepoint{只要你需要「這是第幾個」這個位置資訊，優先用\texttt{enumerate()}，不要自己
另外造一個變數手動累加計數。手動維護一個計數器正是下一節要比較的\texttt{while}常見的
負擔，\texttt{enumerate()}讓\texttt{for}也能拿到位置資訊，卻不用揹上這個負擔。}

\begin{drills}[label=drills:enumerate]
\item 不要真的執行，用眼睛看：下面這段程式碼會印出哪幾行？
\begin{verbatim}
for index, char in enumerate("cat"):
    print(f"{index}: {char}")
\end{verbatim}
\item 填空：想要印出\texttt{word = "hello"}裡每個字母的「第幾個」（從 1 開始算，不是
從 0），迴圈該怎麼寫？（提示：跟第~\ref{sec:for-range}~節的\texttt{+ 1}是同樣的手法。）
\item 抓蟲：下面這段程式碼想找出\texttt{"apple"}裡字母\texttt{"p"}第一次出現的位置，
但實際執行會印出兩個索引，問題出在哪裡？
\begin{verbatim}
word = "apple"

for index, char in enumerate(word):
    if char == "p":
        print(index)
\end{verbatim}
（提示：對照\texttt{find-letter.py}，找到之後少了什麼？）
\end{drills}

\section{不確定次數的重複：\texttt{while}}
\label{sec:while-basics}

前面幾節的\texttt{for}，不管是重複固定次數、走訪一串東西，還是巢狀著用、搭配
\texttt{enumerate()}一起拿位置資訊，都有一個共同點：Python 在迴圈開始之前，其實就已經
知道要跑幾輪。

但生活裡更常見的重複其實不是這樣，洗碗要洗到碗槽空了才能停，存錢要存到
存夠目標金額才能停，你沒辦法在動手之前就先算出精確的次數，只知道「什麼狀態出現時該停
下來」。這種「不確定要做幾次，但知道停止條件」的重複，要用\texttt{while}處理。
\index{while}

\begin{code}{python}{while 迴圈：先檢查、再決定要不要做}
while condition:
    repeated_action
    update_variable_used_in_condition
\end{code}

\texttt{while}每一輪開始前都會先檢查\texttt{condition}，條件成立才執行縮排底下的內容，
執行完再回頭檢查一次，條件不成立了才停下來離開迴圈。

\begin{example}{python}{(saving.py) 存錢存到存夠目標金額}
savings = 0
daily_amount = 45
target = 200
days = 0

while savings < target:
    savings += daily_amount
    days += 1

print(f"Saved for {days} days, total saved: {savings} dollars")
\end{example}

\iosample{}{Saved for 5 days, total saved: 225 dollars \\}

留意最後存到的金額是 225 元，比目標的 200 元還多。這很正常，跟現實生活一樣：
\texttt{while}只在乎「這一輪開始前，條件成不成立」，不會幫你精準停在剛好等於目標的那
一刻。存到第 5 天時，存款從 155 元加到 200 元，已經到達目標了，但這一輪動作已經做完，
程式才會在下一輪開始前重新檢查，發現\texttt{200 < 200}不成立，於是停下來，不是在存錢的
當下就精準喊停。

\texttt{while}迴圈最容易出的問題，是忘記更新\texttt{condition}用到的變數，導致條件
永遠成立、迴圈永遠不會停。
\index{無窮迴圈}

\begin{example}{python}{(saving-bug.py) 忘記更新條件變數}
savings = 0
target = 100

while savings < target:
    print("Saved a little money")
\end{example}

這段程式碼裡，\texttt{savings}從頭到尾都是 0，\texttt{savings < target}永遠成立，
迴圈永遠不會停下來。實際執行的話，螢幕會被「Saved a little money」這行字灌爆，程式看起來像當機了
一樣，其實只是陷入了無窮迴圈。想要停下來，只能按下\texttt{Ctrl+C}（或是在編輯器裡按下
停止鍵）強制中斷程式。修法是在迴圈裡把\texttt{savings}真的加上去，例如補上
\texttt{savings += 50}這一行，讓條件變數每一輪都真的往目標靠近。

\keypoint{\texttt{while}迴圈能不能停下來，完全取決於迴圈裡有沒有確實更新
\texttt{condition}用到的變數。寫\texttt{while}迴圈時，養成習慣問自己一句：「這個條件
會不會在某一輪之後變成不成立？」如果答案是不會，這個迴圈就永遠不會停。}

忘記更新變數只是無窮迴圈最常見的一種寫法，還有一種更隱蔽的版本：變數確實有更新，但更新
的方式讓條件永遠不會變成不成立。

\begin{example}{python}{(never-zero.py) 條件用 != 卻永遠等不到那個值}
n = 5

while n != 0:
    print(n)
    n -= 2
\end{example}

\texttt{n}從 5 開始，每一輪減 2：5、3、1，再減 2 變成 -1、-3、-5……\texttt{n}永遠不會
剛好等於 0，因為它是用「跳 2 格」的方式往下數，直接跳過了 0 這個值。條件\texttt{n != 0}
（\texttt{n}不等於 0）就永遠成立，即使變數確實每一輪都有更新，這仍然是一個貨真價實的
無窮迴圈。

\keypoint{拿\texttt{!=}當\texttt{while}的停止條件時要特別小心：只要變數更新的方式
（例如每次加 2、減 3）有可能整個跳過你比較的那個值，條件就永遠不會變成不成立。想確保
迴圈一定停得下來，優先用\texttt{<}或\texttt{>}這種「範圍」型的比較，不要用\texttt{!=}
這種「剛好等於」型的比較。像\texttt{savings < target}或\texttt{n > 0}，就算變數不小心
跳過某個精確數字，範圍式的條件還是會在下一輪就抓到「已經不成立」。}

\begin{drills}[label=drills:while-basics]
\item 不要真的執行，用眼睛看：下面這段程式碼會印出幾次「chances left」？
\begin{verbatim}
n = 3
while n > 0:
    print(f"{n} chances left")
    n -= 1
\end{verbatim}
\item 抓蟲：下面這段程式碼會陷入無窮迴圈，問題出在哪？該怎麼改？
\begin{verbatim}
attempts = 3

while attempts > 0:
    print(f"{attempts} chances left")
\end{verbatim}
\item 抓蟲：下面這段程式碼也會陷入無窮迴圈，問題出在哪？（提示：注意迴圈裡到底更新的是
哪一個變數）
\begin{verbatim}
remaining = 3
count = 0

while remaining > 0:
    print(f"{remaining} chances left")
    count += 1
\end{verbatim}
\item 承接\texttt{never-zero.py}：如果把\texttt{n = 5}改成\texttt{n = 6}，其他都不變，
這個迴圈還會是無窮迴圈嗎？為什麼？
\end{drills}

\section{不知道要問幾次：\texttt{while} 搭配 \texttt{input()}}
\label{sec:while-input}

前一節的存錢範例，\texttt{while}要不要繼續是靠程式自己算出來的；但很多時候，決定迴圈
要跑幾輪的人其實是使用者，你不知道他今天要記錄幾件家事，只知道他打完之後，會告訴你
「好了」。還記得第 \ref{ch:strings} 章用過的\texttt{input()}嗎？這裡我們要反覆用它讀取使用者輸入，
每一輪迴圈都問一次，直到使用者輸入約定好的結束關鍵字為止。

\begin{code}{python}{while 搭配 input() 反覆詢問}
value = input(prompt_text)

while value != stop_keyword:
    repeated_action
    value = input(prompt_text)
\end{code}

留意這裡\texttt{input()}被呼叫了兩次：一次在迴圈開始之前，先問一輪、拿到第一個答案；
一次在迴圈裡面，每做完一輪就再問一次、拿新的答案回來，讓下一輪開始前的條件檢查看到的是
最新的輸入。

\begin{example}{python}{(chores.py) 反覆記錄今天做了哪些家事}
chore = input("What chore did you do today? (Type 'done' to finish)")

while chore != "done":
    print(f"Logged: {chore}")
    chore = input("Anything else? (Type 'done' to finish)")

print("Chore log complete!")
\end{example}

\iosample{wash dishes \\ take out trash \\ done \\}{Logged: wash dishes \\ Logged: take out trash \\ Chore log complete! \\}

這裡只用到\texttt{input()}最基本的能力，拿到使用者打的一行字，拿去跟一個固定的字串
比較。至於提示文字要怎麼設計得更友善、或是把使用者輸入的內容轉成數字來用，這本書之後
需要的時候會再回頭補充，這裡先掌握「\texttt{while}的停止條件可以來自使用者輸入」這件
事就夠了。

\begin{drills}[label=drills:while-input]
\item 承接上面的\texttt{chores.py}：如果使用者依序輸入「mop floor」、「done」，程式總共會
印出幾行「Logged:」？
\item 如果使用者一開始就直接輸入「done」，程式還會印出任何一行「Logged:」嗎？
\end{drills}

\section{\texttt{for} 還是 \texttt{while}？怎麼選}
\label{sec:for-vs-while}

現在你已經看過\texttt{for}的幾種用法（固定次數、走訪字串、巢狀、搭配
\texttt{enumerate()}），也看過\texttt{while}怎麼處理不確定次數的重複。兩種語法都能寫出
「重複做某件事」的程式，新手常見的疑惑是：這兩個工具，到底該挑哪一個？

判斷的關鍵只有一個問題：動手之前，你知不知道要重複幾次，或者知不知道要處理的是哪一批
東西？知道的話，用\texttt{for}；只知道「什麼狀態出現時該停」，但不知道要跑幾輪，才用
\texttt{while}。擦五扇窗戶、走訪一句話裡的每個字元，這些你在動手之前就已經確定範圍，
用\texttt{for}最直接。存錢存到存夠、反覆詢問使用者直到他說「好了」，這些你不知道要跑
幾輪，只能用\texttt{while}邊做邊檢查。

實務上可以用兩個問題快速判斷：

\begin{itemize}
\item 動手之前，「要跑幾輪」或「要處理哪些東西」是不是已經確定？確定的話，用
\texttt{for}。
\item 停止的判斷，是不是要看執行過程中才會出現、事先算不出來的狀態（使用者說了什麼、
存款有沒有到達目標）？是的話，用\texttt{while}。
\end{itemize}

這兩個問題的答案通常會一致：已知範圍幾乎都對應「用\texttt{for}」，看狀態才能決定幾乎都
對應「用\texttt{while}」。真正會讓人猶豫的情境並不多，多半是因為把一個其實已知的範圍，
誤看成了不確定的狀態，或是反過來。

\begin{example}{python}{(windows-while.py) 把擦窗戶的例子改寫成 while}
window = 0

while window < 5:
    print(f"Cleaning window {window + 1}")
    window += 1

print("All windows cleaned!")
\end{example}

\iosample{}{Cleaning window 1 \\ Cleaning window 2 \\ Cleaning window 3 \\ Cleaning window 4
\\ Cleaning window 5 \\ All windows cleaned! \\}

這段程式碼跟第~\ref{sec:for-range}~節的\texttt{windows.py}做的事完全一樣，但多寫了兩
行：一行負責在迴圈開始前，先把\texttt{window}設成 0；一行負責在迴圈裡手動把
\texttt{window}加 1。多這兩行不只是麻煩，還多了一個風險：如果哪天不小心漏寫了
\texttt{window += 1}這一行，這段程式碼就會變成無窮迴圈，跟第~\ref{sec:while-basics}~
節看過的坑一模一樣。同樣的重複，用\texttt{for}搭配\texttt{range()}既不用自己管理計數
變數，也不用擔心忘記更新它。

\keypoint{已知次數、或是走訪一批固定內容的重複，優先用\texttt{for}：\texttt{range()}
和直接走訪序列都會替你處理好「這一輪走到哪、還有沒有下一輪」，不會有忘記更新變數導致
無窮迴圈的風險。只有在停止條件真的要看執行過程中才會變化的狀態（使用者輸入、還沒存夠
的金額）時，才需要\texttt{while}。}

這也是為什麼你會發現，越到後面的章節，Python 程式碼裡\texttt{for}出現的頻率遠高於
\texttt{while}。第 \ref{ch:lists} 章要學的 list、第 \ref{ch:dict-set} 章的 dict，只要是「有一批東西，要對每一個都
做一件事」的情境，幾乎都是先想\texttt{for}，很少會拿\texttt{while}硬幹。

原因跟這一節
比較\texttt{windows.py}／\texttt{windows-while.py}時看到的一樣：\texttt{for}走訪容器
不用你自己管索引、不用自己設計停止條件，程式碼更短，也更不容易寫出無窮迴圈。
\texttt{while}真正無可取代的地方，是那些連你自己都不知道要跑幾輪的情境，像使用者輸入、
或是之後會遇到的網路請求重試這一類。

\begin{drills}[label=drills:for-vs-while]
\item 下面這個任務比較適合用\texttt{for}還是\texttt{while}？「檢查購物清單上的每一項，
在清單上打勾」（提示：你在開始檢查之前，知不知道清單上有幾項？）
\item 下面這個任務比較適合用\texttt{for}還是\texttt{while}？「打電話給客服，一直等到
有人接聽為止」
\item 抓蟲：把\texttt{windows-while.py}裡的\texttt{window += 1}拿掉，會發生什麼事？
\item 第~\ref{sec:while-basics}~節的\texttt{saving.py}用\texttt{while}存錢存到存夠為
止。如果需求換成「固定存 30 天，不管有沒有存到目標」，這時候比較適合用\texttt{for}還是
\texttt{while}？為什麼？
\end{drills}

\section{提早結束：\texttt{break} 與線性搜尋}
\label{sec:break-search}

前面幾節的迴圈，都是乖乖跑完該跑的次數，或是跑到條件不成立才停。但有些情境不一樣。你在
翻找家裡的抽屜找鑰匙時，只要找到了，就不用繼續翻剩下的抽屜；反過來，如果翻遍了所有抽屜
都沒找到，你才會下結論說「鑰匙不在這裡」。這種「找到就提早結束，找完都沒找到才下結論」
的動作，Python 用\texttt{break}處理。
\index{break}

\begin{code}{python}{for 搭配 break 提早結束}
for item in sequence:
    if condition:
        break
    repeated_action
\end{code}

一旦執行到\texttt{break}，Python 會立刻跳出目前這個迴圈，後面還沒跑的輪次全部跳過，
直接接著執行迴圈後面的程式碼。

\begin{example}{python}{(find-key.py) 翻找抽屜找鑰匙}
key_location = 3
found = False

for drawer in range(1, 6):
    print(f"Searching drawer {drawer}")
    if drawer == key_location:
        found = True
        break

if found:
    print(f"Key found in drawer {key_location}!")
else:
    print("Searched every drawer, no key found")
\end{example}

\iosample{}{Searching drawer 1 \\ Searching drawer 2 \\ Searching drawer 3 \\ Key found in
drawer 3! \\}

這種「從頭開始一個一個找，找到就停，找完都沒找到才說沒有」的做法，有個正式的名字叫
\textbf{線性搜尋}。
\index{線性搜尋}
上面的例子用了一個\texttt{found}變數記錄「到底有沒有找到」，等迴圈
結束（不管是提早\texttt{break}結束，還是正常跑完）才根據\texttt{found}的值決定要印哪一
句話。這個順序很重要，接下來要示範一個新手很常見的寫法錯誤。

\begin{example}{python}{(find-key-bug.py) 把兩種結果都寫在迴圈裡面判斷}
key_location = 5

for drawer in range(1, 6):
    if drawer == key_location:
        print("Key found!")
    else:
        print("No key in this drawer, search ended")
    break
\end{example}

\iosample{}{No key in this drawer, search ended \\}

鑰匙明明放在第 5 個抽屜，這段程式碼卻在翻第 1 個抽屜就說「search ended」，而且完全沒提到
鑰匙到底找到了沒有。問題出在\texttt{break}的縮排位置：它跟\texttt{if}／\texttt{else}
同一層，不管這一輪翻到的是不是目標抽屜，迴圈都只跑一輪就會被\texttt{break}停掉。

正確的
寫法是像上面\texttt{find-key.py}那樣，只有真的符合條件（\texttt{drawer == key\_location}）
的那一輪，才把\texttt{break}放進\texttt{if}裡面。沒找到的那幾輪只該正常結束這一輪、
繼續翻下一個抽屜，不該提前喊停，也不該在迴圈裡面就直接下「沒找到」的結論。

\keypoint{\texttt{break}只會讓你跳出目前正在跑的這一層迴圈，不會讓整支程式停下來。
程式會接著執行迴圈後面的程式碼。}

順帶一提，Python 內建的\texttt{in}運算子，其實就是幫你把整個「一個一個找、找到就算」
的搜尋迴圈包好了：\texttt{5 in range(1, 6)}會直接告訴你 True，不用自己手寫迴圈。之後
遇到單純只想知道「有沒有」而不需要知道「翻了幾次」的情境，可以直接用\texttt{in}；但
搜尋迴圈本身的骨架，還是值得你自己動手寫過一次，之後才看得懂\texttt{in}幫你省了什麼。
\index{in 運算子}

\begin{drills}[label=drills:break-search]
\item 不要真的執行，用眼睛看：下面這段程式碼會印出哪幾行？
\begin{verbatim}
target = 4

for number in range(1, 7):
    if number == target:
        print(f"Found {number}")
        break
    print(f"{number} is not the target")
\end{verbatim}
\item 抓蟲：下面這段程式碼哪裡有問題？為什麼鑰匙明明放在第 2 個抽屜，程式卻說沒找到？
\begin{verbatim}
key_location = 2

for drawer in range(1, 4):
    if drawer == key_location:
        print("Key found!")
    else:
        print("No key in this drawer, search ended")
    break
\end{verbatim}
\item 如果找鑰匙的情境改成巢狀的（例如每個房間裡都有好幾個抽屜，要先走訪房間、再走訪
抽屜），只在最內層（抽屜那層）寫\texttt{break}，找到鑰匙後真的能讓兩層迴圈都停下來嗎？
\end{drills}

\section{跳過本輪：\texttt{continue}}
\label{sec:continue}

\texttt{break}是「整個迴圈都不要繼續了」，但有時候你只是想跳過某一輪，其他輪照常進行。
檢查一週七天的運動紀錄時，週日是休息日，遇到週日就跳過不用記錄，但接下來的每一週還是要
繼續檢查。這種「跳過本輪、直接進到下一輪」的動作，Python 用\texttt{continue}處理：一旦
執行到\texttt{continue}，這一輪縮排在它後面的程式碼全部不執行，直接跳到下一輪重新開始。
\index{continue}

\begin{example}{python}{(weekly-check.py) 檢查一週七天的運動紀錄}
for day in range(1, 8):
    if day == 7:
        continue
    print(f"Day {day}: workout complete")

print("Weekly check complete")
\end{example}

\iosample{}{Day 1: workout complete \\ Day 2: workout complete \\ Day 3: workout complete
\\ Day 4: workout complete \\ Day 5: workout complete \\ Day 6: workout complete
\\ Weekly check complete \\}

留意第 7 天完全沒有出現在輸出裡。不是程式漏掉了，是\texttt{continue}讓那一輪直接跳過
了底下的\texttt{print}，但迴圈本身並沒有結束，第 7 天結束後迴圈也正常跑完，接著印出
「Weekly check complete」。

\begin{drills}[label=drills:continue]
\item 不要真的執行，用眼睛看：\texttt{for n in range(1, 6):}底下接
\texttt{if n == 3: continue}，再接\texttt{print(n)}，會印出哪幾個數字？
\item 如果把上面\texttt{weekly-check.py}裡判斷的條件從\texttt{day == 7}改成
\texttt{day == 3}，\texttt{continue}換成\texttt{break}之後，輸出會有什麼不一樣？
\item 抓蟲（延伸）：如果一個巢狀迴圈的內層寫了\texttt{continue}，外層那一輪會不會也
被跳過？
\end{drills}

\section{迴圈正常跑完才做的事：for-else 與 while-else}
\label{sec:loop-else}

第~\ref{sec:break-search}~節翻抽屜找鑰匙的例子，用了一個\texttt{found}變數記錄「到底
有沒有找到」，等迴圈結束才根據\texttt{found}的值決定要印哪一句話。Python 其實提供一個
內建的語法，可以省掉這個額外的變數：\texttt{for}和\texttt{while}都可以接一個\texttt{else}
區塊，這個\texttt{else}只有在迴圈正常跑完、沒有被\texttt{break}提早結束的情況下才會
執行。
\index{for-else／while-else}

\begin{code}{python}{for-else：迴圈正常跑完才執行 else}
for item in sequence:
    if condition:
        break
else:
    action_if_not_broken
\end{code}

留意\texttt{else}的縮排跟\texttt{for}同一層，不是跟裡面的\texttt{if}同一層。這裡的
\texttt{else}問的其實是「整個迴圈跑完的過程中，有沒有被\texttt{break}打斷過」，跟裡面
那個\texttt{if}有沒有成立是兩回事。只要迴圈從頭到尾都沒碰到\texttt{break}，不管是正常
跑滿所有輪次，還是\texttt{sequence}根本是空的、一次都沒跑，\texttt{else}都會執行；只要
中途執行到\texttt{break}，\texttt{else}就會被跳過。

\begin{example}{python}{(find-key-else.py) 用 for-else 改寫翻抽屜找鑰匙}
key_location = 3

for drawer in range(1, 6):
    print(f"Searching drawer {drawer}")
    if drawer == key_location:
        print(f"Key found in drawer {drawer}!")
        break
else:
    print("Searched every drawer, no key found")
\end{example}

\iosample{}{Searching drawer 1 \\ Searching drawer 2 \\ Searching drawer 3 \\ Key found in
drawer 3! \\}

跟第~\ref{sec:break-search}~節的\texttt{find-key.py}比較一下：這裡完全不需要
\texttt{found}這個變數，迴圈找到鑰匙就\texttt{break}、印出訊息；找不到的話，迴圈正常
跑完六輪，\texttt{else}就會接手印出「Searched every drawer, no key found」。少了一個額外變數，
邏輯卻一樣清楚。

\begin{example}{python}{(find-target-else.py) 迴圈正常跑完，else 被執行}
for number in range(1, 4):
    if number == 5:
        print("Found 5")
        break
else:
    print("5 not found")
\end{example}

\iosample{}{5 not found \\}

\texttt{range(1, 4)}裡從頭到尾都不會出現 5，迴圈跑完三輪，一次\texttt{break}都沒碰到，
於是\texttt{else}被執行，印出「5 not found」。

\keypoint{很多人第一次看到\texttt{for-else}會誤以為\texttt{else}是「迴圈一次都沒跑」
才執行，其實不是。\texttt{else}只跟「有沒有被\texttt{break}打斷」有關，跟迴圈本身跑了
幾輪、甚至有沒有跑完全沒關係，只要正常跑完（哪怕跑了零輪），\texttt{else}就會執行。}

\texttt{while}也可以接同樣的\texttt{else}，規則完全一樣：\texttt{while}的條件變成不
成立、正常結束迴圈，\texttt{else}就會執行；中途被\texttt{break}提早結束，\texttt{else}
就會被跳過。

\begin{code}{python}{while-else：迴圈正常結束才執行 else}
while condition:
    if early_stop_condition:
        break
else:
    action_if_not_broken
\end{code}

\begin{example}{python}{(login-attempts.py) 密碼嘗試三次後鎖定帳號}
attempts = 0
password = "1234"
guess = "0000"

while attempts < 3:
    if guess == password:
        print("Password correct")
        break
    print("Wrong password, try again")
    attempts += 1
else:
    print("Max attempts reached, account locked")
\end{example}

\iosample{}{Wrong password, try again \\ Wrong password, try again \\ Wrong password, try again
\\ Max attempts reached, account locked \\}

這裡\texttt{guess}是先寫死的示範值，只是要展示整個流程；如果要讓使用者真的每一次都
重新輸入，可以把\texttt{guess = input(...)}搬進迴圈裡面，就像第~\ref{sec:while-input}~
節做過的那樣。迴圈跑了三輪，\texttt{guess}始終不等於\texttt{password}，
\texttt{attempts}從 0 累加到 3，\texttt{while attempts < 3}的條件變成不成立，迴圈正常
結束，不是被\texttt{break}打斷，於是\texttt{else}執行，印出鎖定訊息。

\stylepoint{\texttt{for-else}／\texttt{while-else}是 Python 一個相對少見的語法，很多
其他程式語言完全沒有這個概念，所以就算你之後查資料查到別的語言的範例，也不用太意外看
不到它。適合用在「找到就\texttt{break}，沒找到才要做某件事」這種場景，能省掉一個額外
的旗標變數；但如果你或合作的人對這個語法不熟，用一個明確的旗標變數（像
\texttt{find-key.py}那樣）有時候反而更好讀，兩種寫法都要看得懂。}

\begin{drills}[label=drills:loop-else]
\item 不要真的執行，用眼睛看：下面這段程式碼會印出哪些行？
\begin{verbatim}
for number in range(1, 5):
    if number == 3:
        print("Found 3")
        break
    print(f"{number} is not the target")
else:
    print("Not found")
\end{verbatim}
\item 如果把上面的\texttt{number == 3}改成\texttt{number == 9}，輸出會有什麼不同？
\item 抓蟲：下面這段\texttt{while-else}，想在猜對密碼時印出「Password correct」，猜錯三次印出
「Account locked」，但實際執行會發現兩句話都印出來了，問題出在哪？
\begin{verbatim}
attempts = 0
while attempts < 3:
    if guess == password:
        print("Password correct")
    attempts += 1
else:
    print("Account locked")
\end{verbatim}
（提示：對照\texttt{login-attempts.py}，這裡少了什麼關鍵字？）
\end{drills}

\section{三種常見的迴圈任務長相：累加、計數、搜尋}
\label{sec:accumulate-count}

回頭看看這一章寫過的幾個例子：

\begin{itemize}
\item 第~\ref{sec:while-basics}~節的存錢範例，是把每天存的錢一次一次加總起來，直到存
夠為止。
\item 第~\ref{sec:for-string}~節數空白字元的範例，是每遇到一次符合條件的字元，就把計
數往上加一。
\item 第~\ref{sec:break-search}~節翻抽屜找鑰匙，是一個一個找，找到就停。
\end{itemize}

這三個例子表面上做的事情完全不一樣，但拆開來看，骨架其實長得很像：

\begin{code}{python}{迴圈裡的常見任務骨架}
result = initial_value    # 迴圈開始前，先準備好需要的變數

for item in sequence:
    if condition:
        result = new_result_value   # 迴圈中，視情況每一輪更新它

print(result)              # 迴圈結束後，檢查最終結果
\end{code}

三種長相的差別，只在「迴圈中每一輪具體做了什麼」：

\begin{itemize}
\item \textbf{累加}：\index{累加}每一輪把值加到\texttt{result}上，像第~\ref{sec:while-basics}~節
的存錢範例，\texttt{savings += daily\_amount}。
\item \textbf{計數}：\index{計數}每一輪遇到符合條件的情況，就把\texttt{result}加 1，像
第~\ref{sec:for-string}~節數空白字元，\texttt{space\_count += 1}。
\item \textbf{搜尋}：\index{搜尋（迴圈任務長相）}每一輪檢查是不是找到目標，找到就用\texttt{break}提早結束，像
第~\ref{sec:break-search}~節翻抽屜找鑰匙。
\end{itemize}

再用一個新例子練一次計數模式：

\begin{example}{python}{(count-vowels.py) 數一句話裡有幾個母音字母}
sentence = "today is a good day"
vowels = "aeiou"
vowel_count = 0

for char in sentence:
    if char in vowels:
        vowel_count += 1

print(f"This sentence has {vowel_count} vowels")
\end{example}

\iosample{}{This sentence has 7 vowels \\}

這三種骨架都是「先準備一個變數，每一輪視情況更新它」，但\texttt{result}不一定要是
數字。同一套邏輯換成字串一樣能用：先準備一個空字串，每一輪把新的內容用串接的方式接
上去，就能一路把過程中產生的東西收集起來。

\begin{example}{python}{(mask-spaces.py) 把一句話裡的空白換成底線，重新組成新字串}
sentence = "today is a good day"
report = ""

for char in sentence:
    if char == " ":
        report += "_"
    else:
        report += char

print(report)
\end{example}

\iosample{}{today\_is\_a\_good\_day \\}

\texttt{report}一開始是空字串，每一輪迴圈都用\texttt{+=}把一個字元（或代替空白的底線）
接到\texttt{report}後面，跑完整個\texttt{sentence}，\texttt{report}就變成一句話從頭到
尾都處理過一遍的新字串。這正是累加骨架的另一種長相，只是\texttt{result}累加的不是數字，
是字串。

之後學到 list，你會看到同一個「先準備一個空的容器，每一輪塞一點東西進去」的
邏輯，只是換成 list 搭配\texttt{append()}方法。不管容器是數字、字串，還是之後的 list，
先準備、逐輪更新、最後看結果，這個節奏都是一樣的。

\keypoint{累加、計數、搜尋這三種骨架，會在你之後幾乎每一章、每一支程式裡反覆出現。
「迴圈前先準備好一個變數，迴圈中每一輪視情況更新它，迴圈後看最終結果」這個順序，值得你
牢牢記住，不管\texttt{result}放的是數字還是字串，節奏都一樣。之後遇到一個新問題時，先
問自己一句：「這比較像累加、計數、還是搜尋？」通常就能很快想出迴圈該怎麼寫。}

\begin{drills}[label=drills:accumulate-count]
\item 不要真的執行，用眼睛看：下面這段程式碼最後印出的\texttt{total}是多少？
\begin{verbatim}
total = 0
for i in range(1, 5):
    total += i
print(total)
\end{verbatim}
\item 填空：讓下面這段程式碼算出\texttt{word}這個字串裡有幾個字母\texttt{"e"}。
\begin{verbatim}
word = "elephant"
count = ____

for letter in word:
    if letter == "e":
        count ____

print(count)
\end{verbatim}
\item 填空：想要用累加字串的方式，把\texttt{word = "python"}裡的每個字母都重複兩次接
起來（結果應該是\texttt{"ppyytthhoonn"}），迴圈該怎麼寫？（提示：可以用兩次\texttt{+}
把同一個字母接兩遍。）
\end{drills}

\section*{練習解答}
\addcontentsline{toc}{section}{練習解答}

\paragraph{第~\ref{sec:for-range}~節}
\begin{enumerate}
\item 印出：0、1、2（三行）。
\item 8，\texttt{range(2, 9, 2)}依序是 2、4、6、8，9 本身不會被包含進去。
\item \texttt{range(9, 0, -2)}，從 9 開始，每次減 2，數到比 1 大的地方為止，最後一個
數字剛好是 1。
\item 0 次。\texttt{range(5, 5)}起點跟終點相同，代表沒有任何數字落在這個範圍內，迴圈
本體一次都不會執行。
\end{enumerate}

\paragraph{第~\ref{sec:for-string}~節}
\begin{enumerate}
\item 印出：C、A、T（三行），\texttt{.upper()}把每個字元轉成大寫。
\item 缺少冒號，錯誤訊息會是\texttt{SyntaxError: expected ':'}。修法是在
\texttt{"abc"}後面補上冒號：\texttt{for char in "abc":}。
\end{enumerate}

\paragraph{第~\ref{sec:nested-loops}~節}
\begin{enumerate}
\item 6 行。外層迴圈跑 2 輪，內層迴圈每一輪都完整跑 3 輪，2 乘 3 等於 6。
\item 10 次。外層\texttt{range(1, 6)}是 5 個房間，內層\texttt{range(1, 3)}是每個房間
2 扇窗戶，5 乘 2 等於 10。
\item 只在內層寫\texttt{break}不夠，因為\texttt{break}只會跳出離它最近的那一層迴圈
（抽屜／窗戶那一層）。想讓外層也跟著停，需要額外設一個變數（例如\texttt{found}）記錄
「已經找到了」，在外層迴圈裡也檢查這個變數，符合條件時外層自己也執行一次
\texttt{break}。
\end{enumerate}

\paragraph{第~\ref{sec:enumerate}~節}
\begin{enumerate}
\item 印出：\texttt{0: c}、\texttt{1: a}、\texttt{2: t}（三行）。
\item \texttt{for index, char in enumerate(word): print(f"Letter \{index + 1\} is
\{char\}")}，跟\texttt{range()}一樣，索引從 0 開始算，要顯示「第幾個」給人看時，記得
補上\texttt{+ 1}。
\item 少了\texttt{break}。找到\texttt{"p"}就印出索引之後沒有提早結束迴圈，所以
\texttt{"apple"}裡兩個\texttt{"p"}（索引 1 跟 2）都會被印出來，不只印出第一次出現的
位置。
\end{enumerate}

\paragraph{第~\ref{sec:while-basics}~節}
\begin{enumerate}
\item 印出 3 次「chances left」（3 chances left、2 chances left、1 chances left），因為
\texttt{n}每輪都有確實減 1，跑到\texttt{n == 0}時條件就不成立了。
\item 這是無窮迴圈，迴圈裡從來沒有更新\texttt{attempts}，條件永遠成立。修法是在迴圈裡
補上一行讓\texttt{attempts}遞減，例如\texttt{attempts -= 1}。
\item 這也是無窮迴圈：條件用的是\texttt{remaining}，但迴圈裡實際更新的卻是
\texttt{count}，\texttt{remaining}從頭到尾都沒變過，\texttt{remaining > 0}永遠成立。
修法是把\texttt{count += 1}改成\texttt{remaining -= 1}，或是額外補上更新
\texttt{remaining}的那一行。
\item 不會是無窮迴圈。6 每次減 2 會剛好經過 0（6、4、2、0），條件\texttt{n != 0}在
\texttt{n}變成 0 的那一輪開始前就不成立了，迴圈正常停止，印出 6、4、2 三行。
\end{enumerate}

\paragraph{第~\ref{sec:while-input}~節}
\begin{enumerate}
\item 1 行，「Logged: mop floor」。「mop floor」不是結束關鍵字，所以迴圈印出這一行、再問一次；
下一次輸入「done」，條件不成立，迴圈結束，不會再印出「Logged:」開頭的行。
\item 不會，因為第一次讀到的輸入就已經是「done」，\texttt{while}的條件一開始就不成立，
迴圈本體一次都不會執行，只會印出「Chore log complete!」。
\end{enumerate}

\paragraph{第~\ref{sec:for-vs-while}~節}
\begin{enumerate}
\item 適合用\texttt{for}，因為開始檢查之前，購物清單上有幾項是已知的、固定的。
\item 適合用\texttt{while}，因為你不知道要等幾秒、等幾次才會有人接聽，只知道「有人接
才停」這個條件。
\item 條件永遠成立，\texttt{window < 5}變成無窮迴圈，程式會不斷印出「Cleaning window
1」，因為\texttt{window}從頭到尾都是 0，永遠小於 5。
\item 適合用\texttt{for}。「固定存 30 天」在動手之前就已經確定次數，屬於已知範圍，改用
\texttt{for i in range(30):}最直接；換成\texttt{while}反而得自己額外維護一個計數變數，
多繞一圈。
\end{enumerate}

\paragraph{第~\ref{sec:break-search}~節}
\begin{enumerate}
\item 印出：1 is not the target、2 is not the target、3 is not the target、Found 4，迴圈
跑到\texttt{number == 4}時符合條件，印出訊息後用\texttt{break}提早結束，後面的 5、6
都不會再檢查。
\item \texttt{break}的縮排跟\texttt{if}／\texttt{else}同一層，不管有沒有符合條件，
迴圈都只跑一輪就會被強制停掉，翻第 1 個抽屜沒找到，就直接被\texttt{break}停住，根本
沒機會翻到第 2 個抽屜。修法是把\texttt{break}移進\texttt{if}裡面，只有真的找到才提早
結束。
\item 不能。只在最內層（抽屜層）寫\texttt{break}，只會讓那一個房間的抽屜迴圈提早結束，
外層的房間迴圈完全不受影響，還是會繼續進到下一個房間重新翻抽屜。想讓兩層都停，得額外
處理外層（參考第~\ref{sec:nested-loops}~節的做法）。
\end{enumerate}

\paragraph{第~\ref{sec:continue}~節}
\begin{enumerate}
\item 印出：1、2、4、5，\texttt{n == 3}那一輪被\texttt{continue}跳過，沒有執行到
\texttt{print(n)}。
\item 用\texttt{continue}的話，只有第 3 天被跳過，第 4 到第 7 天照常印出，輸出還是有
六行「workout complete」；換成\texttt{break}的話，迴圈跑到第 3 天就整個結束，第 4 到第 7 天
根本不會再檢查，只會印出「Day 1」「Day 2」兩行。\texttt{continue}是「跳過這一輪，
其他輪照常進行」，\texttt{break}是「後面全部都不要了」，兩者不一樣。「Weekly check
complete」那一行則不管哪種寫法都會印出來，因為它在迴圈外面，迴圈結束（不管是正常跑完還是被
\texttt{break}提早結束）之後都會執行到它。
\item 不會。\texttt{continue}只會跳過內層迴圈自己這一輪剩下的動作，外層迴圈完全不知道
內層發生了什麼事，會繼續照常進行自己的輪次。
\end{enumerate}

\paragraph{第~\ref{sec:loop-else}~節}
\begin{enumerate}
\item 印出：1 is not the target、2 is not the target、Found 3。因為找到 3 就
\texttt{break}，\texttt{else}被跳過，不會印出「Not found」。
\item 迴圈會跑完 1 到 4 全部四輪，一次都沒有符合\texttt{number == 9}，也就是一次
\texttt{break}都沒碰到，於是\texttt{else}執行，最後多印出「Not found」這一行，其餘四行
「is not the target」照常印出。
\item 少了\texttt{break}。猜對密碼只有印出「Password correct」，但迴圈沒有提早結束，還是會把
剩下的輪次跑完，跑完之後（沒有被\texttt{break}打斷）\texttt{else}一定會執行，所以
「Account locked」永遠都會印出來，不管密碼有沒有猜對。
\end{enumerate}

\paragraph{第~\ref{sec:accumulate-count}~節}
\begin{enumerate}
\item 10，\texttt{range(1, 5)}依序是 1、2、3、4，加總起來是 1+2+3+4=10。
\item \texttt{count = 0}，以及\texttt{count += 1}，先把\texttt{count}初始化成 0，
每遇到一個字母\texttt{"e"}就加 1。
\item \texttt{result = ""}，接著\texttt{for letter in word: result += letter +
letter}，每一輪把同一個字母接兩次，串到\texttt{result}後面。
\end{enumerate}

\chapter{List}
\label{ch:lists}

前面幾章，你手上能用的變數一次只能裝一個東西：一個數字、一個字串、一個真假值。可是生活
裡要記的事情很少是單獨一件，你要記的是「這禮拜要買的東西」，是「今天要做的三件事」，是
「一班同學的名字」。

這一章要教的 list（中文有時翻成「清單」或「串列」，這本書直接用
list 這個字），就是 Python 用來裝「一整批、有順序的東西」的方式。整條路線大致是這樣：

\begin{itemize}
\item 先學會建立一份清單、用位置把裡面的東西一個一個拿出來、用迴圈把整份清單走一遍。
\item 接著讓你動手改這份清單，加東西進去、把東西拿掉，這時你會第一次真正碰到「list 是
可以被改動的」這件事，跟前面學過的數字、字串完全不一樣。
\item 再來是這一章真正的重頭戲：當你把同一份清單「借」給另一個名字的時候，會發生一件
一開始很反直覺的事——改其中一個，另一個也跟著變了。這件事不是 bug，是 Python 處理
「指定」這件事本來的運作方式，第 \ref{ch:variables} 章已經先預告過一次「名字牌」的比喻，這一章你會第一次
在真實情境裡撞見它，也學會怎麼在真的需要「另外複製一份」時做到。
\item 中段會補齊排序、常用工具方法，以及清單裡再放一份清單的寫法。
\item 最後初步認識一種更精簡的寫法，叫 list comprehension，讓你看得懂、用得出來就好，
還不用追求寫得又快又漂亮。
\end{itemize}

\section{建立一份清單、用位置拿東西}
\label{sec:list-basics}

想像你手上有這一整週要做的事情，通常一次會有好幾件，不會只有一件孤零零的事。Python 讓你
用一對方括號把這些東西通通裝進同一個變數裡，這就是 list。你在第 \ref{ch:strings} 章字串已經學過索引
（index）這個概念，這裡 list 一樣可以用同樣的中括號語法取值。
\index{list}

\begin{code}{python}{建立 list、用索引取值}
my_list = [item1, item2, item3]
my_list[i]
\end{code}

\texttt{my\_list} 只是名字，你可以取任何合法的變數名；\texttt{item1}、\texttt{item2}
這些位置放什麼型別都可以，數字、字串、甚至混著放也合法。\texttt{[i]} 裡的 \texttt{i}
就是索引，第一個東西的索引是 0，不是 1，這一點常常讓剛學的人卡關，多練幾次就會習慣。

\begin{example}{python}{(todo-list.py) 建立代辦清單、取出項目}
todo_list = ["do homework", "take out trash", "wash dishes", "mop floor"]
print(todo_list[0])
print(todo_list[2])
print(todo_list[-1])
\end{example}

\iosample{}{do homework \\ wash dishes \\ mop floor \\}

\texttt{todo\_list[0]} 拿到第一項「do homework」，\texttt{todo\_list[2]} 拿到第三項「wash dishes」。
最後一行的 \texttt{todo\_list[-1]} 用的是負索引：\texttt{-1} 代表「倒數第一個」，
\texttt{-2} 代表「倒數第二個」，以此類推。

這個寫法特別方便的地方是：你常常不確定一份
清單到底有幾項，但只要你想拿「最後一個」，直接寫 \texttt{[-1]} 就好，不用先去算清單的
長度是多少。空清單也是合法的寫法，例如 \texttt{empty\_list = []}，代表「目前還沒有任何
東西，之後會陸續加進去」。

清單裡放的東西也不用是同一種型別。前面幾章學過的 int、str、bool，全部都可以混在同一份
清單裡，Python 不會因此抱怨。

\begin{example}{python}{(mixed-info.py) 一份清單裝不同型別的東西}
info = ["Alice", 15, True]
print(info)
print(info[0], info[1], info[2])
print(type(info[0]))
print(type(info[1]))
\end{example}

\iosample{}{['Alice', 15, True] \\{} Alice 15 True \\{} <class 'str'> \\{} <class 'int'> \\}

\texttt{info} 這份清單同時裝了字串、整數、布林值，Python 完全不介意。每個位置存的東西，
型別是記在物件自己身上的，不是記在清單這個容器上，還記得第 \ref{ch:variables} 章「型別長在物件上，不是長在
變數上」這件事嗎？list 只是把好幾個各自帶著自己型別的物件，依照順序收在同一個容器裡而已。

如果你打的索引超出了清單的範圍呢？

\begin{example}{python}{(todo-list-error.py) 索引超出範圍}
todo_list = ["do homework", "take out trash", "wash dishes", "mop floor"]
print(todo_list[4])
\end{example}

這份清單一共只有 4 項，合法的索引是 0 到 3，\texttt{[4]} 已經超出去了。執行後你會看到：

\begin{verbatim}
Traceback (most recent call last):
  File "todo-list-error.py", line 2, in <module>
    print(todo_list[4])
          ~~~~~~~~~^^^
IndexError: list index out of range
\end{verbatim}

由下往上讀：最後一行 \texttt{IndexError: list index out of range}，還記得第 \ref{ch:strings} 章字串
也會有的 \texttt{IndexError} 嗎？list 一樣會發生，道理相同，只是訊息換成了
\texttt{list index out of range}。之後只要清單長度沒你想的長，就很容易撞見它，記得
回頭數一下清單裡到底有幾項。

\begin{drills}[label=drills:list-indexing]
\item 給定 \texttt{fruits = ["apple", "banana", "orange", "grape"]}，\texttt{print(fruits[1])}
會印出什麼？
\item 承上題，\texttt{print(fruits[-2])} 會印出什麼？
\item 如果你想拿到 \texttt{fruits} 的最後一項，但不確定這份清單一共有幾項，你會怎麼寫
索引，才能不用先數清單長度就拿到它？
\item 給定 \texttt{info = ["Emma", 20, False]}，\texttt{type(info[2])} 的結果會是哪一種
型別？
\end{drills}

\section{切片：取一段清單}
\label{sec:list-slicing}

這個語法你在第 \ref{ch:strings} 章字串那裡已經學過一次了，list 的切片用法跟字串完全一樣，這裡只給
一兩個例子喚醒你的記憶，不重講切片的基本概念本身。

\begin{code}{python}{list 切片}
my_list[start:end]
my_list[:]
\end{code}

\texttt{my\_list[start:end]} 取出從 \texttt{start} 開始、到 \texttt{end}（不含）為止
的一段；把 \texttt{start} 或 \texttt{end} 留空，就代表「從最開頭」或「到最後」；兩個都
留空，也就是 \texttt{my\_list[:]}，代表「取整份清單」。

\begin{example}{python}{(weekly-menu.py) 複習切片}
week_menu = ["egg pancake", "rice ball", "sandwich", "egg pancake", "bread", "sushi", "dumplings"]
print(week_menu[1:4])
print(week_menu[:3])

copy_menu = week_menu[:]
copy_menu.append("pizza")
print(week_menu)
print(copy_menu)
\end{example}

\iosample{}{['rice ball', 'sandwich', 'egg pancake'] \\{} ['egg pancake', 'rice ball', 'sandwich'] \\{} ['egg pancake', 'rice ball', 'sandwich', 'egg pancake', 'bread', 'sushi', 'dumplings'] \\{} ['egg pancake', 'rice ball', 'sandwich', 'egg pancake', 'bread', 'sushi', 'dumplings', 'pizza'] \\}

這裡先埋一個伏筆：\texttt{copy\_menu = week\_menu[:]} 這一行做完之後，接下來不管你怎麼改
\texttt{copy\_menu}（例如上面用 \texttt{append} 多加一項披薩），原本的 \texttt{week\_menu}
都不會被影響，這兩份清單做完切片之後就是「各自獨立」的。

至於為什麼需要「特地」用這種
寫法才能做到「各自獨立」，這件事聽起來理所當然，但其實不是每種寫法都會有這個效果，完整
的解釋留到第~\ref{sec:aliasing}~節，你會在那裡看到不特地複製會發生什麼事。

切片還有一個跟索引不一樣的地方：如果 \texttt{start} 或 \texttt{end} 給的數字超出了清單
的範圍，並不會像單一索引那樣丟出 \texttt{IndexError}，而是自動夾到有效範圍內。

\begin{example}{python}{(slice-out-of-range.py) 切片的範圍超過清單長度}
scores = [70, 85, 90, 60, 100]
print(scores[2:100])
print(scores[-100:3])
print(scores[100:200])
\end{example}

\iosample{}{[90, 60, 100] \\{} [70, 85, 90] \\{} [] \\}

\texttt{scores[2:100]} 裡的 \texttt{100} 遠遠超過清單長度，但 Python 不會抱怨，直接切到
清單結尾為止；\texttt{scores[-100:3]} 裡的 \texttt{-100} 也一樣，直接從最開頭切起；第三
行 \texttt{scores[100:200]} 起點跟終點都超出範圍，切出來的結果是一份空清單
\texttt{[]}，也不是錯誤。這跟單一索引很不一樣，切片對「範圍抓太大」比較寬容。

\begin{drills}[label=drills:list-slicing]
\item 給定 \texttt{names = ["Alice", "Bob", "Charlie", "Taylor", "David"]}，
\texttt{names[1:3]} 會印出什麼？
\item 承上題，\texttt{names[:2]} 會印出什麼？
\item 給定 \texttt{names} 跟上題一樣，\texttt{names[2:100]} 會印出什麼？會不會報錯？
\end{drills}

\section{用迴圈走訪整份清單}
\label{sec:list-loop}

第 \ref{ch:loops} 章教過的 \texttt{for} 迴圈，那時候你多半是拿它配 \texttt{range()} 產生一串數字。
其實 \texttt{for} 迴圈真正拿手的，是直接對一份清單裡「真正的東西」跑一遍，一個一個處理，
不用透過索引繞一圈。

\begin{example}{python}{(todo-list-loop.py) 走訪代辦清單}
todo_list = ["do homework", "take out trash", "wash dishes"]

for task in todo_list:
    print("Todo:", task)

print(len(todo_list))
print("wash dishes" in todo_list)
print("mop floor" in todo_list)
\end{example}

\iosample{}{Todo: do homework \\ Todo: take out trash \\ Todo: wash dishes \\ 3 \\ True \\ False \\}

\texttt{for task in todo\_list:} 這一行，每繞一圈，\texttt{task} 就會依序變成清單裡的
下一項，直到走完整份清單為止，你不用自己管索引跑到第幾個了，Python 幫你處理。
\texttt{len(todo\_list)} 告訴你這份清單一共有幾項；
\index{len()}
還記得第 \ref{ch:loops} 章的 \texttt{in} 運算子
嗎？一樣可以用在 list 上，問「這個東西在不在清單裡」，回傳的 \texttt{True}／\texttt{False}
就是第 \ref{ch:conditionals} 章學過的布林值。

迴圈除了走訪、印東西出來，另一個很常見的用法是邊走訪、邊挑、邊把符合條件的東西收進一份
新的清單。這個組合之後會反覆出現，這裡先看一次、感覺一下它的長相。

\begin{example}{python}{(steps-over-8000.py) 用迴圈篩出符合條件的項目}
weekday_steps = [8000, 10500, 6200, 7300, 9100]

over_8000 = []
for steps in weekday_steps:
    if steps > 8000:
        over_8000.append(steps)

print(over_8000)
\end{example}

\iosample{}{[10500, 9100] \\}

這裡先造一份空清單 \texttt{over\_8000}，用第 \ref{ch:conditionals} 章教過的 \texttt{if} 判斷「這一步數是不是
超過 8000」，符合的才用 \texttt{append} 收進去，不符合就什麼都不做、直接進下一輪。這種
「造空清單、跑迴圈、視情況 append」的組合非常常見，本章最後一節會教你一種更精簡的寫法，把
這整個組合濃縮成一行。

清單裡除了放數字、字串，也可以再放一份清單進去，變成清單裡面裝著清單。這種寫法之後常常
拿來表示表格狀的資料，本章稍後會有一整節專門介紹它，這裡先知道有這回事就好。

\begin{drills}[label=drills:list-loop]
\item 這段程式執行後，螢幕上依序會印出什麼？
\begin{verbatim}
colors = ["red", "green", "blue"]
for c in colors:
    print(c + " color")
\end{verbatim}
\item 這段程式執行後，最後 \texttt{print(total)} 會印出什麼？
\begin{verbatim}
nums = [3, 1, 4, 1, 5]
total = 0
for n in nums:
    total = total + n
print(total)
\end{verbatim}
\item 給定 \texttt{nums = [3, 8, 15, 4, 20]}，如果你要用迴圈篩出所有大於 10 的數字、
收進一份新清單 \texttt{big\_nums}，程式碼要怎麼寫？\texttt{big\_nums} 最後會是什麼？
\end{drills}

\section{修改與新增元素}
\label{sec:list-add}

前面幾章學過的數字、字串、布林值，一旦造出來，內容就不能被「原地」改掉，你只能造一個
新的值，重新指定給同一個名字。

list 完全不同：它是\textbf{可變的（mutable）}，你可以直接
在原本那份清單上動手，加東西、改東西，不用造一份新的出來取代它。這是這一章要花時間立
起來的一個新性質，之後很多行為都建立在這個性質上。
\index{可變（mutable）}

\begin{code}{python}{修改與新增元素}
my_list[i] = new_value
my_list.append(new_item)
my_list.insert(i, new_item)
\end{code}

\texttt{my\_list[i] = new\_value} 把第 \texttt{i} 個位置的東西整個換掉；
\index{索引賦值}
\texttt{my\_list.append(new\_item)} 把新東西加到清單最尾端；
\index{append()}
\texttt{my\_list.insert(i,
new\_item)} 把新東西插進第 \texttt{i} 個位置，原本在那個位置（以及後面）的東西都往後
挪一格。
\index{insert()}

\begin{example}{python}{(todo-list-update.py) 修改與新增代辦事項}
todo_list = ["do homework", "take out trash", "wash dishes"]

todo_list[0] = "do math homework"
print(todo_list)

todo_list.append("take out recycling")
print(todo_list)

todo_list.insert(0, "reply to client email")
print(todo_list)
\end{example}

\iosample{}{['do math homework', 'take out trash', 'wash dishes'] \\{} ['do math homework', 'take out trash', 'wash dishes', 'take out recycling'] \\{} ['reply to client email', 'do math homework', 'take out trash', 'wash dishes', 'take out recycling'] \\}

第一步先把「do homework」改成更精確的「do math homework」，這是索引賦值；接著臨時多想到一件事要辦
（take out recycling），用 \texttt{append} 加到清單最後面；最後又想到一件更急、得插隊處理的事
（reply to client email），用 \texttt{insert(0, ...)} 直接塞到最前面。三種寫法各自對應「換掉」
「加到尾巴」「插進指定位置」三種不同的需求。

跟切片一樣，\texttt{insert()} 對「位置給太大」也很寬容：如果你給的索引超過清單目前的
長度，Python 不會報錯，只會直接把新東西放到最後面。

\begin{example}{python}{(insert-clamp.py) insert 位置超出範圍}
todo_list = ["do homework", "take out trash", "wash dishes"]
todo_list.insert(100, "mop floor")
print(todo_list)
\end{example}

\iosample{}{['do homework', 'take out trash', 'wash dishes', 'mop floor'] \\}

這份清單明明只有 3 項，\texttt{insert(100, "拖地")} 給的索引 100 遠遠超過，但 Python 把
它當成「反正就是放最後面」處理，不會像單一索引取值那樣丟出 \texttt{IndexError}。

\begin{drills}[label=drills:list-add]
\item 給定 \texttt{cart = ["apple", "banana"]}，執行 \texttt{cart.append("orange")} 之後，
\texttt{cart} 會變成什麼？
\item 接續上題，再執行 \texttt{cart.insert(1, "watermelon")}，\texttt{cart} 會變成什麼？
\item 接續上題，再執行 \texttt{cart[0] = "guava"}，\texttt{cart} 最後會變成什麼？
\end{drills}

\section{一次加入很多項：extend()}
\label{sec:list-extend}

\texttt{append()} 每次只能加一個東西進去。如果你手上已經有另一份清單，想把它裡面的每一項
都個別加進來（而不是把整份清單當成「一個東西」塞進去），\texttt{append()} 就不是你要的
工具了。

\begin{code}{python}{extend()：展開另一份清單，逐項加入}
my_list.extend(other_list)
\end{code}

\begin{example}{python}{(extend-vs-append.py) append 跟 extend 的差別}
cart_a = ["apple", "banana"]
cart_b = ["apple", "banana"]

new_items = ["orange", "grape"]

cart_a.append(new_items)
print(cart_a)

cart_b.extend(new_items)
print(cart_b)
\end{example}

\iosample{}{['apple', 'banana', ['orange', 'grape']] \\{} ['apple', 'banana', 'orange', 'grape'] \\}

\texttt{cart\_a.append(new\_items)} 把整份 \texttt{new\_items} 當成「一個東西」加進去，
所以 \texttt{cart\_a} 最後變成三項，其中第三項本身又是一份清單，也就是清單裡面裝著清單；
\texttt{cart\_b.extend(new\_items)} 則是把 \texttt{new\_items} 裡的每一項一個一個拆開，
分別加進 \texttt{cart\_b}，結果是四個平行的項目，沒有多出一層清單。

\keypoint{\texttt{append(x)} 永遠只增加「一項」，就算 \texttt{x} 本身是一份清單也一樣，
會整份塞進去變成清單裡的清單；\texttt{extend(x)} 則是把 \texttt{x} 裡的每一項逐一攤開、
分別加進來。想合併兩份清單的內容、又不想多一層，要用 \texttt{extend()}，不能用
\texttt{append()}。}
\index{extend()}

\begin{drills}[label=drills:list-extend]
\item 給定 \texttt{a = [1, 2]}、\texttt{b = [3, 4]}，執行 \texttt{a.append(b)} 之後，
\texttt{a} 會變成什麼？
\item 承上題，如果一開始執行的是 \texttt{a.extend(b)}，\texttt{a} 會變成什麼？
\end{drills}

\section{刪除元素}
\label{sec:list-remove}

清單裡的東西除了加，當然也要能刪。Python 給了三種刪除的寫法，選哪一種要看你手上有的
資訊是什麼。

\begin{code}{python}{刪除元素}
del my_list[i]
my_list.pop(i)
my_list.remove(value)
\end{code}

判斷的口訣很簡單：

\begin{itemize}
\item 知道位置、也不需要留著那個值，用 \texttt{del}。
\index{del}
\item 知道位置、但還想留著被移除的那個值（例如接下來要用它），用 \texttt{pop()}。
\index{pop()}
\item 只知道值本身、不知道它在清單裡的位置，用 \texttt{remove()}。
\index{remove()}
\end{itemize}

（其實 \texttt{del} 不只能刪清單裡的元素，也能整個刪掉一個
變數本身，這裡我們只示範刪清單元素的用法，不展開這個旁支。）

\begin{example}{python}{(todo-list-delete.py) 三種刪除方式}
todo_list = ["do homework", "take out trash", "wash dishes", "mop floor"]

del todo_list[1]
print(todo_list)

removed_task = todo_list.pop(0)
print(todo_list)
print(removed_task)

todo_list.remove("mop floor")
print(todo_list)
\end{example}

\iosample{}{['do homework', 'wash dishes', 'mop floor'] \\{} ['wash dishes', 'mop floor'] \\ do homework \\{} ['wash dishes'] \\}

\texttt{del todo\_list[1]} 直接把「take out trash」從清單上抹掉，不留下任何值；
\texttt{todo\_list.pop(0)} 把「do homework」拿掉，同時把它回傳給 \texttt{removed\_task}，
所以你還留著這個值可以用；\texttt{todo\_list.remove("mop floor")} 則是直接告訴 Python
「把值是『mop floor』的那一項刪掉」，不用管它現在排第幾個。

\texttt{pop()} 其實還有一個很常用的省略寫法：如果括號裡什麼都不寫，預設拿掉的是「最後一
項」，效果等於 \texttt{pop(-1)}。

\begin{example}{python}{(pop-default.py) pop() 不給位置，預設拿最後一項}
stack = ["Plate 1", "Plate 2", "Plate 3"]
top = stack.pop()
print(top)
print(stack)
\end{example}

\iosample{}{Plate 3 \\{} ['Plate 1', 'Plate 2'] \\}

\texttt{stack.pop()} 沒有給任何位置，Python 預設幫你拿掉並回傳「最後一項」，這裡是最後
疊上去的「Plate 3」。這種「每次都處理最後加進來的那一個」的用法很常見，例如疊盤子、瀏覽器
的上一頁紀錄，之後你只要看到 \texttt{.pop()} 沒帶引數，就知道它動的是清單最尾端。

這裡有一個新手很容易踩到的坑：\texttt{pop()}、\texttt{append()}、\texttt{insert()}、
\texttt{remove()} 這些會修改清單的方法，大多數都直接回傳 \texttt{None}（也就是「沒有
值」），不是回傳一份改好的新清單。如果你把它們的結果重新指定回原本的變數，會發生什麼事？
\index{None}

\begin{example}{python}{(todo-list-remove-bug.py) 常見錯誤：把方法結果指定回原變數}
todo_list = ["do homework", "take out trash", "wash dishes"]
todo_list = todo_list.remove("take out trash")
print(todo_list)
todo_list.append("mop floor")
\end{example}

執行後你會看到：

\begin{verbatim}
None
Traceback (most recent call last):
  File "todo-list-remove-bug.py", line 4, in <module>
    todo_list.append("mop floor")
    ^^^^^^^^^^^^^^^^
AttributeError: 'NoneType' object has no attribute 'append'
\end{verbatim}

第一行先印出 \texttt{None}，因為 \texttt{remove()} 把清單裡的「take out trash」刪掉之後，
回傳的是 \texttt{None}，而 \texttt{todo\_list = todo\_list.remove("take out trash")} 這一行，
把這個 \texttt{None} 重新指定回了 \texttt{todo\_list}，原本那份「刪完之後」的清單，
就這樣被 \texttt{None} 蓋掉、再也拿不回來了。

下一行想對 \texttt{todo\_list} 呼叫
\texttt{append}，但這時候 \texttt{todo\_list} 已經是 \texttt{None}，不是清單了，
所以 Python 抱怨 \texttt{'NoneType' object has no attribute 'append'}，意思是
「\texttt{None} 這種東西身上根本沒有 \texttt{append} 這個方法可以叫」。
\index{AttributeError}

\keypoint{\texttt{remove()}、\texttt{append()}、\texttt{insert()}、\texttt{pop()}
這些會直接修改清單的方法，大多數回傳的都是 \texttt{None}，不是一份改好的新清單。正確的
用法是直接呼叫它們（\texttt{todo\_list.remove("take out trash")}），不用、也不能把結果重新指定
回原本的變數。}

你會發現，很多 list 方法動的是清單本身、不是回傳一份新的，這跟接下來第~\ref{sec:aliasing}~
節要講的「名字牌」有關，先記住這個感覺就好。

\begin{drills}[label=drills:list-remove]
\item 給定 \texttt{items = ["A", "B", "C", "D"]}，執行 \texttt{del items[0]} 之後，
\texttt{items} 會變成什麼？
\item 接續上題，執行 \texttt{x = items.pop(1)} 之後，\texttt{items} 跟 \texttt{x}
分別會是什麼？
\item 給定 \texttt{plates = ["Plate A", "Plate B", "Plate C"]}，執行
\texttt{last = plates.pop()} 之後，\texttt{last} 跟 \texttt{plates} 分別會是什麼？
\item 抓蟲：這段程式碼哪裡有問題？\texttt{data} 最後會是什麼？
\begin{verbatim}
data = [5, 2, 8]
data = data.remove(2)
print(data)
\end{verbatim}
\end{drills}

\section{兩個名字，同一份清單：別名行為}
\label{sec:aliasing}

前一節你已經確實看到 list 是「可變的」，現在要把這件事再往前推一步，這也是這一章概念
強度最高的地方，值得放慢腳步看清楚。

想像你把手上的代辦清單「借」給室友看，不是影印一份給他，是直接把同一張紙遞過去。如果
室友在上面多寫了一項，你把紙拿回來的時候，那一項當然也在上面，因為根本是同一張紙，只是
一時之間有兩個人在看它。Python 裡的變數賦值，做的正是這件事。

\begin{code}{python}{別名賦值}
other_name = my_list
\end{code}

這一行看起來像是「把 \texttt{my\_list} 複製一份給 \texttt{other\_name}」，但 Python
實際做的完全不是這樣。

\begin{example}{python}{(borrow-list.py) 借給室友看的代辦清單}
my_list = ["wash dishes", "take out trash"]
roommate_list = my_list
roommate_list.append("buy soy sauce")
print(my_list)
print(roommate_list)
\end{example}

\iosample{}{['wash dishes', 'take out trash', 'buy soy sauce'] \\{} ['wash dishes', 'take out trash', 'buy soy sauce'] \\}

你只對 \texttt{roommate\_list} 呼叫了 \texttt{append}，明明沒有動 \texttt{my\_list}
半根手指頭，但 \texttt{print(my\_list)} 卻也印出多了一項「buy soy sauce」的結果，兩份清單
變成一模一樣，因為它們根本就是同一份。

\keypoint{\texttt{=} 不是複製一份新的放進另一個變數，是幫同一個物件多貼一張名字牌。
\texttt{roommate\_list = my\_list} 的意思是「把另一張名字牌 \texttt{roommate\_list}
也貼到 \texttt{my\_list} 指的那份清單上」，兩張名字牌指著同一個東西，透過任何一張
名字牌改動內容，另一張名字牌看到的也是改過的結果。}
\index{別名（aliasing）}

\begin{center}
\begin{tikzpicture}[node distance=1.1cm and 3.4cm]
\node[rectangle, rounded corners=4pt, draw=coreStroke, thick, fill=coreFill,
      minimum width=2.9cm, minimum height=0.9cm, font=\ttfamily\small] (nameA) {my\_list};
\node[rectangle, rounded corners=4pt, draw=coreStroke, thick, fill=coreFill,
      minimum width=2.9cm, minimum height=0.9cm, font=\ttfamily\small,
      below=1.1cm of nameA] (nameB) {roommate\_list};
\node[rectangle, rounded corners=4pt, draw=memStroke, thick, fill=memFill,
      minimum width=6.4cm, minimum height=1.3cm, font=\ttfamily\small,
      right=3.4cm of nameA] (obj) {["wash dishes", "take out trash", "buy soy sauce"]};
\draw[diagarrow] (nameA.east) -- ++(0.7,0) |- (obj.west);
\draw[diagarrow] (nameB.east) -- ++(0.7,0) |- (obj.west);
\end{tikzpicture}
\end{center}
\diagcaption{\texttt{my\_list} 跟 \texttt{roommate\_list} 是貼在同一份清單物件上的
兩張名字牌，不是兩份各自獨立的清單，改其中一張名字牌看到的內容，等於改了另一張。}

那如果你真的想要「另外造一份」，兩份清單各自獨立、互不影響呢？第~\ref{sec:list-slicing}~
節已經先示範過 \texttt{[:]} 切片可以做到這件事；除此之外，Python 還提供了 \texttt{.copy()}
跟 \texttt{list()} 兩種寫法，效果是一樣的。
\index{copy()}

\begin{example}{python}{(copy-list.py) 真正複製一份獨立的清單}
my_list = ["wash dishes", "take out trash"]
copy_list = my_list.copy()
copy_list.append("buy soy sauce")
print(my_list)
print(copy_list)
\end{example}

\iosample{}{['wash dishes', 'take out trash'] \\{} ['wash dishes', 'take out trash', 'buy soy sauce'] \\}

這次 \texttt{copy\_list} 是用 \texttt{.copy()} 造出來的一份全新清單，跟 \texttt{my\_list}
不再是同一個物件，所以對 \texttt{copy\_list} 呼叫 \texttt{append} 之後，\texttt{my\_list}
完全沒被影響。\texttt{copy\_list = my\_list.copy()}、\texttt{copy\_list = list(my\_list)}、
\texttt{copy\_list = my\_list[:]} 這三種寫法，效果完全一樣，挑你順手的用即可。

既然「兩份清單內容一樣」跟「兩份清單根本是同一份」是兩件不同的事，Python 也提供了兩種不
一樣的問法：\texttt{==} 問的是「內容一不一樣」，\texttt{is} 問的是「是不是貼在同一個物件
上」。

\begin{example}{python}{(equal-vs-is.py) == 跟 is 的差別}
my_list = ["wash dishes", "take out trash"]
roommate_list = my_list
copy_list = my_list.copy()

print(my_list == roommate_list)
print(my_list is roommate_list)
print(my_list == copy_list)
print(my_list is copy_list)
\end{example}

\iosample{}{True \\ True \\ True \\ False \\}

\texttt{my\_list} 跟 \texttt{roommate\_list} 不只內容一樣（\texttt{==} 是 True），根本
是同一個物件（\texttt{is} 也是 True）；\texttt{my\_list} 跟 \texttt{copy\_list} 內容也
一樣（\texttt{==} 是 True），但因為 \texttt{copy\_list} 是另外造出來的，兩者不是同一個
物件（\texttt{is} 是 False）。大多數時候你想問的是「內容一不一樣」，該用 \texttt{==}；
\texttt{is} 通常只用來確認「是不是同一份」這種比較少見、但這一章正好用得上的情境，之後在
第 \ref{ch:functions} 章、第 \ref{ch:tuples} 章還會再遇到它。

這也是為什麼把 list 傳給函式之後，函式裡面對它做的修改，你在外面也會看到，第 \ref{ch:functions} 章教函式
的時候會完整說明，這裡先知道有這回事就好。

\begin{drills}[label=drills:aliasing]
\item 這段程式碼執行後，\texttt{print(a)} 會印出什麼？
\begin{verbatim}
a = [1, 2, 3]
b = a
b.append(4)
print(a)
\end{verbatim}
\item 這段程式碼執行後，\texttt{print(x)} 會印出什麼？
\begin{verbatim}
x = ["red", "green"]
y = x
y[0] = "yellow"
print(x)
\end{verbatim}
\item 改寫題：\texttt{original = [1, 2, 3]} 之後執行 \texttt{new\_list = original}，
如果你希望之後改動 \texttt{new\_list} 不會影響到 \texttt{original}，這一行程式碼要
怎麼改？
\item 給定 \texttt{p = [1, 2]}、\texttt{q = p}、\texttt{r = p.copy()}，
\texttt{p is q}、\texttt{p is r} 這兩個問句的結果分別會是什麼？
\end{drills}

\section{\texttt{sort()} 跟 \texttt{sorted()}：原地排序 vs 造一份新的}
\label{sec:sort-vs-sorted}

排序是這一章對「別名／可變」這個概念的第二次具體練習。Python 給了兩種排序的寫法，差別
剛好就是上一節談的「改動原本那份」跟「另外造一份新的」。

\begin{code}{python}{排序 list}
my_list.sort()
sorted(my_list)
\end{code}

\texttt{my\_list.sort()} 改的是 \texttt{my\_list} 這份清單本身，跟第~\ref{sec:aliasing}~
節「同一個物件被改動」是同一件事，排序完之後原本的順序就回不去了；
\index{sort()}
\texttt{sorted(my\_list)}
則是造一份新的、排好序的清單回傳給你，原本那份 \texttt{my\_list} 完全不受影響。
\index{sorted()}

\begin{example}{python}{(sort-vs-sorted.py) 依優先順序排序代辦清單}
tasks = ["wash dishes", "take out trash", "do homework"]
tasks.sort()
print(tasks)

scores = [82, 95, 67]
sorted_scores = sorted(scores)
print(scores)
print(sorted_scores)
\end{example}

\iosample{}{['do homework', 'take out trash', 'wash dishes'] \\{} [82, 95, 67] \\{} [67, 82, 95] \\}

\texttt{tasks.sort()} 之後，\texttt{tasks} 直接變成排好序的樣子；但 \texttt{sorted(scores)}
不一樣，呼叫完之後 \texttt{scores} 還是原本的順序沒有變，排好序的結果被存進了另一個名字
\texttt{sorted\_scores} 裡。想留住原本的順序、又要用到排序後的結果，就該用
\texttt{sorted()}；確定不需要原本的順序了，直接 \texttt{sort()} 比較省事。

\begin{drills}[label=drills:sort-vs-sorted]
\item 這段程式碼執行後，\texttt{print(nums)} 會印出什麼？
\begin{verbatim}
nums = [5, 3, 8, 1]
nums.sort()
print(nums)
\end{verbatim}
\item 這段程式碼執行後，\texttt{nums} 跟 \texttt{result} 分別會是什麼？
\begin{verbatim}
nums = [5, 3, 8, 1]
result = sorted(nums)
print(nums)
print(result)
\end{verbatim}
\end{drills}

\section{進階排序：reverse 與 key 參數}
\label{sec:sort-key}

\texttt{sort()} 跟 \texttt{sorted()} 預設是由小到大排序，但實際上兩者都還接受兩個很常
用的選擇性引數：\texttt{reverse} 決定要不要反過來排，\texttt{key} 決定排序時真正拿來比較
的是什麼。

\begin{code}{python}{reverse 與 key 參數}
my_list.sort(reverse=True)
my_list.sort(key=function_name)
sorted(my_list, key=function_name, reverse=True)
\end{code}

\texttt{reverse=True} 讓排序結果由大到小，不寫的話預設是 \texttt{reverse=False}（由小到
大）；
\index{sort()／sorted() 的 reverse 參數}
\texttt{key=function\_name} 則是告訴 Python：「排序的時候，先對清單裡每一項呼叫
這個函式一次，拿它的回傳值來決定順序，不要直接比較元素本身」。這裡 \texttt{key} 接的是
函式這個東西本身，不是呼叫函式之後的結果，所以名字後面不加括號。
\index{sort()／sorted() 的 key 參數}

\begin{example}{python}{(sort-key.py) 用 reverse 跟 key 客製化排序}
words = ["banana", "kiwi", "fig", "watermelon"]
words.sort(key=len)
print(words)

scores = [82, 95, 67, 71]
scores.sort(reverse=True)
print(scores)

places = ["home", "school", "Central Park", "Times Square"]
by_length = sorted(places, key=len, reverse=True)
print(by_length)
\end{example}

\iosample{}{['fig', 'kiwi', 'banana', 'watermelon'] \\{} [95, 82, 71, 67] \\{} ['Central Park', 'Times Square', 'school', 'home'] \\}

上面這三段各自示範一種用法：

\begin{itemize}
\item \texttt{words.sort(key=len)} 排序時真正拿來比較的不是字串本身的字母順序，是每個
字串先被算出 \texttt{len()}（字元數）之後的結果，再依照字元數由少到多排。
\item \texttt{scores.sort(reverse=True)} 單純把預設的由小到大反過來，變成由大到小。
\item 第三段把 \texttt{key} 跟 \texttt{reverse} 合起來用，依照地點名稱的字數由多到少
排序。
\end{itemize}

如果兩個地點名稱字數剛好一樣（「Central Park」跟「Times Square」都是
十二個字元），Python 會維持它們原本在清單裡的相對順序，不會隨便調換。

\keypoint{\texttt{key=len} 給的是 \texttt{len} 這個函式本身，不是 \texttt{len()} 呼叫完
的結果，中間不能加括號。如果不小心寫成 \texttt{key=len()}，Python 會在呼叫 \texttt{sort}
之前就先執行一次 \texttt{len()}，而 \texttt{len()} 缺引數會直接丟出 \texttt{TypeError}。}

\begin{example}{python}{(sort-key-bug.py) 常見錯誤：key 後面多加了括號}
words = ["banana", "kiwi", "fig"]
words.sort(key=len())
print(words)
\end{example}

執行後你會看到：

\begin{verbatim}
Traceback (most recent call last):
  File "sort-key-bug.py", line 2, in <module>
    words.sort(key=len())
                   ~~~^^
TypeError: len() takes exactly one argument (0 given)
\end{verbatim}

\texttt{len()} 這個函式一定要給它一個東西才能算長度，\texttt{key=len()} 卻在還沒排序前
就單獨呼叫了一次 \texttt{len()}，沒給任何引數，Python 當場抱怨「\texttt{len()} 剛好需要一
個引數（但你給了 0 個）」。正確寫法是拿掉括號，只把函式名稱本身交給 \texttt{key}。

\begin{drills}[label=drills:sort-key]
\item 給定 \texttt{nums = [4, 1, 9, 3]}，\texttt{nums.sort(reverse=True)} 之後，
\texttt{nums} 會變成什麼？
\item 給定 \texttt{words = ["cat", "elephant", "dog"]}，
\texttt{sorted(words, key=len)} 的結果會是什麼？
\end{drills}

\section{常用工具方法一瞥}
\label{sec:list-tools}

list 還有不少好用的小工具，這裡快速瀏覽過去就好，你不用把它們全部背下來，知道「有這些
東西存在」，要用的時候查得到、看得懂輸出，就已經足夠。

\begin{example}{python}{(steps-toolkit.py) list 常用工具方法}
weekday_steps = [8000, 10500, 6200, 7300, 9100]
weekend_steps = [12000, 4500]

all_steps = weekday_steps + weekend_steps
print(all_steps)

repeated = [0] * 5
print(repeated)

print(min(weekday_steps))
print(max(weekday_steps))
print(sum(weekday_steps))

print(weekday_steps.count(8000))
print(weekday_steps.index(7300))
\end{example}

\iosample{}{[8000, 10500, 6200, 7300, 9100, 12000, 4500] \\{} [0, 0, 0, 0, 0] \\ 6200 \\ 10500 \\ 41100 \\ 1 \\ 3 \\}

這幾個工具各自做的事情是：

\begin{itemize}
\item \texttt{+} 把兩份清單接成一份新的（跟字串接字串很像），\texttt{weekday\_steps} 跟
\texttt{weekend\_steps} 自己都不會被改動。
\index{+（清單相加）}
\item \texttt{*} 把清單內容重複幾次，\texttt{[0] * 5}
就是「造一份裝了 5 個 0 的清單」，常拿來一次生出一份佔位用的清單。
\index{*（清單重複）}
\item \texttt{min()}、\texttt{max()}、\texttt{sum()} 分別拿最小值、最大值、總和，清單裡
要是數字才能這樣用。
\index{min()}
\index{max()}
\index{sum()}
\item \texttt{count(8000)} 數清單裡有幾個 8000。
\index{count()}
\item \texttt{index(7300)} 找出 7300 第一次出現在第幾個位置。
\index{index()（清單方法）}
\end{itemize}

\texttt{+} 跟 \texttt{*} 都是「造一份新清單」，不會動到原本的清單，這裡順便確認一次。

\begin{example}{python}{(list-add-mul.py) + 跟 * 不會改動原本的清單}
a = [1, 2, 3]
b = [4, 5]
print(a + b)
print(a)
print(a * 3)
print(a)
\end{example}

\iosample{}{[1, 2, 3, 4, 5] \\{} [1, 2, 3] \\{} [1, 2, 3, 1, 2, 3, 1, 2, 3] \\{} [1, 2, 3] \\}

\texttt{a + b} 跟 \texttt{a * 3} 執行完之後，\texttt{a} 都還是原本的 \texttt{[1, 2, 3]}，
沒有變。這跟第~\ref{sec:list-add}~節、第~\ref{sec:list-remove}~節那些直接修改清單本身的
方法（\texttt{append}、\texttt{insert}、\texttt{remove} 等）不一樣，這裡的 \texttt{+}
跟 \texttt{*} 是運算子，做的事情是「算出一個新的值」，不是「修改原本這個物件」，跟第
~\ref{sec:sort-vs-sorted}~節 \texttt{sorted()} 那種「造新的、不動舊的」是同一類行為。

如果 \texttt{index()} 要找的值根本不在清單裡呢？

\begin{example}{python}{(index-not-found.py) index() 找不到值}
fruits = ["apple", "banana", "orange"]
print(fruits.index("guava"))
\end{example}

\begin{verbatim}
Traceback (most recent call last):
  File "index-not-found.py", line 2, in <module>
    print(fruits.index("guava"))
          ~~~~~~~~~~~~^^^^^^^^^
ValueError: 'guava' is not in list
\end{verbatim}

清單裡沒有「guava」，\texttt{index()} 找不到就丟出 \texttt{ValueError}，還記得第 \ref{ch:variables} 章的
\texttt{ValueError} 嗎？意思都是「值本身沒問題，但內容不合理」，這裡就是「你要找的值根本
不存在」。相較之下，\texttt{count()} 對「找不到」寬容得多：\texttt{fruits.count("guava")}
不會報錯，只會老實印出 \texttt{0}，因為「數出現幾次」本來就允許答案是零次。

\begin{drills}[label=drills:list-tools]
\item 給定 \texttt{a = [1, 2]}、\texttt{b = [3]}，\texttt{a + b} 的結果是什麼？執行完
之後 \texttt{a} 本身會不會被改變？
\item \texttt{[1, 2] * 3} 的結果會是什麼？
\item 給定 \texttt{nums = [4, 8, 15, 16, 23, 42]}，\texttt{sum(nums)} 跟
\texttt{max(nums)} 分別是多少？
\item 給定 \texttt{fruits = ["apple", "banana", "orange"]}，\texttt{fruits.index("guava")}
會發生什麼事？\texttt{fruits.count("guava")} 呢？兩者的行為為什麼不一樣？
\end{drills}

\section{清單中的清單}
\label{sec:nested-list}

清單裡的元素除了數字、字串，也可以是另一份清單，變成清單裡面裝著清單。這種寫法很適合拿
來表示「表格狀」的資料，例如一個班級的成績單：每個學生是一列，列裡面裝著姓名跟好幾科的
分數。
\index{清單中的清單}

\begin{code}{python}{清單中的清單}
outer_list = [inner_list1, inner_list2, inner_list3]
outer_list[i][j]
\end{code}

要拿到最外層清單裡的某一項，用一個中括號就跟平常一樣；如果那一項本身又是一份清單，你可以
再接一個中括號，取出那份清單裡的某個位置，也就是 \texttt{outer\_list[i][j]} 這種寫法。

\begin{example}{python}{(nested-list.py) 用清單中的清單表示成績表}
gradebook = [
    ["Alice", 88, 92],
    ["Bob", 76, 81],
    ["Charlie", 95, 90],
]

print(gradebook[0])
print(gradebook[0][0])
print(gradebook[1][2])

for student in gradebook:
    name = student[0]
    avg = (student[1] + student[2]) / 2
    print(name, avg)
\end{example}

\iosample{}{['Alice', 88, 92] \\ Alice \\ 81 \\ Alice 90.0 \\ Bob 78.5 \\ Charlie 92.5 \\}

\texttt{gradebook[0]} 取出最外層清單的第一項，也就是「Alice」那一整列，本身還是一份清單；
\texttt{gradebook[0][0]} 再往裡面取一層，拿到那一列的第一個東西，也就是姓名；
\texttt{gradebook[1][2]} 拿到第二個學生（Bob）那一列的第三個東西（索引 2，也就是第二科
分數）。

最後的 \texttt{for} 迴圈示範一個很常見的組合：外層迴圈的每一輪 \texttt{student}
本身就是一整列（一份清單），你可以直接對它再用中括號取值，不用另外寫一層巢狀迴圈才能拿到
姓名跟分數。

走訪清單中的清單，除了像上面那樣先取出整列再各自取值，也可以用兩層 \texttt{for} 迴圈，
一層負責走外層的每一列，另一層負責走該列裡的每一個數字，這種「迴圈裡面再包一層迴圈」的
寫法叫巢狀迴圈，跟清單中的清單正好是天生一對。

\begin{example}{python}{(nested-sum.py) 用巢狀迴圈加總所有數字}
weekly_steps = [
    [8000, 10500, 6200],
    [7300, 9100, 12000],
]

total = 0
for person_steps in weekly_steps:
    for steps in person_steps:
        total = total + steps

print(total)
\end{example}

\iosample{}{53100 \\}

外層迴圈每一輪 \texttt{person\_steps} 拿到的是一整列（一份清單），內層迴圈再對這一列裡的
每一個數字各自跑一次、把它加進 \texttt{total}。兩層迴圈合起來，效果就是把整份清單中的
清單、每一個最裡層的數字都走過一遍，跟第 \ref{ch:loops} 章單層迴圈的道理完全一樣，只是多包了一層。

\begin{drills}[label=drills:nested-list]
\item 給定 \texttt{matrix = [[1, 2], [3, 4], [5, 6]]}，\texttt{matrix[2]} 會是什麼？
\item 承上題，\texttt{matrix[2][0]} 會是什麼？
\item 給定 \texttt{board = [["X", "O"], ["O", "X"]]}，要拿到第二列第二欄的
\texttt{"X"}，索引要怎麼寫？
\item 給定 \texttt{matrix = [[1, 2], [3, 4]]}，用巢狀 \texttt{for} 迴圈把裡面所有數字
加總，\texttt{total} 最後會是多少？
\end{drills}

\section{清單相乘的陷阱：巢狀清單的別名}
\label{sec:nested-alias}

第~\ref{sec:list-tools}~節看過 \texttt{[0] * 5} 可以快速造出一份裝滿 0 的清單，很多人
會很自然地想拿同樣的招數，造出一個表格狀的「格子」，例如一個 3 列 3 欄、全部填 0 的
方格。可惜這裡藏著一個經典的坑，正好是第~\ref{sec:aliasing}~節「名字牌」概念的延伸應用。

\begin{example}{python}{(nested-alias-trap.py) 用 * 重複巢狀清單的陷阱}
grid = [[0] * 3] * 3
print(grid)

grid[0][0] = 1
print(grid)
\end{example}

\iosample{}{[[0, 0, 0], [0, 0, 0], [0, 0, 0]] \\{} [[1, 0, 0], [1, 0, 0], [1, 0, 0]] \\}

一開始看起來沒問題，\texttt{grid} 確實印出三列、每列三個 0。但只改了 \texttt{grid[0][0]}
一個位置，結果三列全部都變成 \texttt{[1, 0, 0]}，不是只有第一列。

原因跟第~
\ref{sec:aliasing}~節的別名行為一模一樣：\texttt{[0] * 3} 先造出「一份」內層清單，接著
外層的 \texttt{* 3} 做的事情，其實是把同一份內層清單的名字牌重複貼了三次進外層清單，
沒有「造三份各自獨立的內層清單」。三個位置指向的其實是同一個物件，改其中一個，另外兩個
當然也跟著變。

\keypoint{\texttt{[x] * n} 這種寫法只會造出「一份」\texttt{x}，然後把指向它的名字牌在外層
清單裡放 \texttt{n} 次。如果 \texttt{x} 本身是不可變的（像數字、字串），你完全感覺不出
差異；但如果 \texttt{x} 是像 list 這種可變物件，就會變成 \texttt{n} 個名字牌全部指著同一
份東西，改一個、全部一起變。}

正確的做法是每一列都要各自獨立造一次，這正好可以用本章最後一節要教的 list comprehension
做到。

\begin{example}{python}{(fixed-grid.py) 用 comprehension 造出各自獨立的每一列}
fixed_grid = [[0] * 3 for _ in range(3)]
fixed_grid[0][0] = 1
print(fixed_grid)
\end{example}

\iosample{}{[[1, 0, 0], [0, 0, 0], [0, 0, 0]] \\}

這次結果符合預期：只有第一列被改動。差別在於 \texttt{[0] * 3 for \_ in range(3)} 每繞一
圈迴圈，都會重新執行一次 \texttt{[0] * 3}，造出一份全新、獨立的內層清單，而不是像
\texttt{[[0] * 3] * 3} 那樣只造一次、重複貼三張名字牌。\texttt{\_} 這個變數名字是 Python
的一個小慣例，代表「這一輪迴圈的值我不需要用到，只是要讓迴圈重複跑幾次而已」。

\begin{drills}[label=drills:nested-alias]
\item 給定 \texttt{rows = [[]] * 3}，執行 \texttt{rows[0].append("a")} 之後，
\texttt{rows} 會變成什麼？為什麼？
\item 如果想要三列各自獨立的空清單，\texttt{rows = [[]] * 3} 這一行該怎麼改寫？
\end{drills}

\section{list comprehension 初探}
\label{sec:list-comprehension}

前面你已經很熟悉「造一個空清單、跑迴圈、每次把結果 \texttt{append} 進去」這個組合了。
Python 有一種更精簡的寫法，可以把這整個組合濃縮成一行，叫做 list comprehension。
\index{list comprehension}

\begin{code}{python}{list comprehension（最簡形式）}
[expression for item in iterable]
\end{code}

\texttt{iterable} 就是你可以放進 \texttt{for} 迴圈裡一個一個處理的東西，像 list、字串、
\texttt{range()} 都是；\texttt{expression} 是你想對每個 \texttt{item} 做的運算，運算
完的結果會被收進一份新的清單裡。

\begin{example}{python}{(squares-comprehension.py) 傳統寫法 vs comprehension 寫法}
numbers = [1, 2, 3, 4, 5]

# 傳統寫法：造一個空清單、跑迴圈、逐一 append
squares = []
for n in numbers:
    squares.append(n ** 2)
print(squares)

# comprehension 寫法：效果一樣，濃縮成一行
squares_v2 = [n ** 2 for n in numbers]
print(squares_v2)
\end{example}

\iosample{}{[1, 4, 9, 16, 25] \\{} [1, 4, 9, 16, 25] \\}

兩段程式做的事完全一樣，結果也完全一樣，差別只在後面那一行把「造空清單、跑迴圈、
append」三個步驟濃縮成了一個表達式。這是很多人寫 Python 時的慣用寫法，你不用馬上自己
流利地寫出來，但要看得懂它在做什麼，之後每一章只要用得上，我們會繼續自然地用到它，
多看幾次就會熟。

這裡只教這種最簡單的形式，不帶條件式、也不是雙層迴圈，那些之後有需要
再說。

comprehension 也可以直接拿來造上一節那種「每一列各自獨立」的巢狀清單，效果就跟
\texttt{fixed-grid.py} 那個例子一樣：外層再包一層 comprehension，每繞一圈都重新造一份
全新的內層清單。

\begin{example}{python}{(comprehension-grid.py) 用巢狀 comprehension 造格子}
grid = [[0 for _ in range(3)] for _ in range(3)]
print(grid)
grid[1][1] = 9
print(grid)
\end{example}

\iosample{}{[[0, 0, 0], [0, 0, 0], [0, 0, 0]] \\{} [[0, 0, 0], [0, 9, 0], [0, 0, 0]] \\}

外層的 \texttt{for \_ in range(3)} 跑三輪，每一輪都各自重新執行一次
\texttt{[0 for \_ in range(3)]}，也就是各自重新造一份新的內層清單，所以只改動
\texttt{grid[1][1]} 這一個位置時，其他列完全不受影響，不會再發生上一節那種「改一個、
全部一起變」的陷阱。

\stylepoint{看到「造空清單 + for 迴圈 + append」這個組合時，如果邏輯夠單純（就是對每一項
做一次運算、收進新清單），改寫成一行 list comprehension 會更接近 Python 慣用的寫法；
但如果邏輯複雜到要拆好幾行才看得懂，寫回傳統的 \texttt{for} 迴圈反而更清楚，不用為了
「看起來很 Python」硬塞進一行。}

List 可以隨時被改動，但 Python 裡也有一種內容一旦定了就不能再改的「親戚」，叫做 tuple，
後面章節會看到它，以及為什麼有時候你會想要一份「定死了」的清單。

\begin{drills}[label=drills:list-comprehension]
\item 改寫題：把下面這段迴圈改寫成一行 list comprehension。
\begin{verbatim}
words = ["cat", "dog", "bird"]
lengths = []
for w in words:
    lengths.append(len(w))
\end{verbatim}
\item 填空：下面這行 comprehension 要填什麼，才能讓 \texttt{doubled} 變成
\texttt{[2, 4, 6]}？
\begin{verbatim}
doubled = [______ for n in [1, 2, 3]]
\end{verbatim}
\end{drills}

\section*{練習解答}
\addcontentsline{toc}{section}{練習解答}

\paragraph{第~\ref{sec:list-basics}~節}
\begin{enumerate}
\item \texttt{fruits[1]} 是索引 1（第二項），會印出 \texttt{banana}。
\item \texttt{fruits[-2]} 是倒數第二項，會印出 \texttt{orange}。
\item 寫 \texttt{fruits[-1]}，負索引 \texttt{-1} 永遠代表「最後一項」，不用先知道清單
長度。
\item \texttt{info[2]} 是 \texttt{False}，型別是 \texttt{<class 'bool'>}。
\end{enumerate}

\paragraph{第~\ref{sec:list-slicing}~節}
\begin{enumerate}
\item \texttt{names[1:3]} 取索引 1 到 2（不含 3），會印出
\texttt{['Bob', 'Charlie']}。
\item \texttt{names[:2]} 取開頭到索引 1（不含 2），會印出
\texttt{['Alice', 'Bob']}。
\item \texttt{names[2:100]} 會從索引 2 切到清單結尾，印出
\texttt{['Charlie', 'Taylor', 'David']}，不會報錯，因為切片對超出範圍的邊界很寬容。
\end{enumerate}

\paragraph{第~\ref{sec:list-loop}~節}
\begin{enumerate}
\item 依序印出：red color、green color、blue color（各自一行）。
\item \texttt{total} 會是 \texttt{3 + 1 + 4 + 1 + 5 = 14}。
\item 寫法是：
\begin{verbatim}
big_nums = []
for n in nums:
    if n > 10:
        big_nums.append(n)
\end{verbatim}
\texttt{big\_nums} 最後會是 \texttt{[15, 20]}。
\end{enumerate}

\paragraph{第~\ref{sec:list-add}~節}
\begin{enumerate}
\item \texttt{cart} 會變成 \texttt{['apple', 'banana', 'orange']}。
\item 承上，\texttt{insert(1, "watermelon")} 把「watermelon」插進索引 1 的位置，變成
\texttt{['apple', 'watermelon', 'banana', 'orange']}。
\item 承上，\texttt{cart[0] = "guava"} 把第一項換掉，最後變成
\texttt{['guava', 'watermelon', 'banana', 'orange']}。
\end{enumerate}

\paragraph{第~\ref{sec:list-extend}~節}
\begin{enumerate}
\item \texttt{a.append(b)} 把整份 \texttt{b} 當成一個東西塞進去，\texttt{a} 會變成
\texttt{[1, 2, [3, 4]]}。
\item \texttt{a.extend(b)} 把 \texttt{b} 的每一項逐一攤開加入，\texttt{a} 會變成
\texttt{[1, 2, 3, 4]}。
\end{enumerate}

\paragraph{第~\ref{sec:list-remove}~節}
\begin{enumerate}
\item \texttt{del items[0]} 之後，\texttt{items} 變成 \texttt{['B', 'C', 'D']}。
\item \texttt{x = items.pop(1)} 把索引 1 的 \texttt{'C'} 拿掉並存進 \texttt{x}，所以
\texttt{x} 是 \texttt{'C'}，\texttt{items} 變成 \texttt{['B', 'D']}。
\item \texttt{plates.pop()} 沒給位置，預設拿掉最後一項，\texttt{last} 是
\texttt{'Plate C'}，\texttt{plates} 變成 \texttt{['Plate A', 'Plate B']}。
\item 問題出在 \texttt{data = data.remove(2)}，\texttt{remove()} 回傳的是
\texttt{None}，把它重新指定回 \texttt{data} 之後，原本刪完的清單就不見了，\texttt{data}
最後會是 \texttt{None}。正確寫法是不要重新指定，直接寫 \texttt{data.remove(2)} 就好。
\end{enumerate}

\paragraph{第~\ref{sec:aliasing}~節}
\begin{enumerate}
\item \texttt{b = a} 之後兩個名字指著同一份清單，\texttt{b.append(4)} 之後
\texttt{print(a)} 會印出 \texttt{[1, 2, 3, 4]}。
\item 同樣的道理，\texttt{y[0] = "yellow"} 改的是 \texttt{x} 跟 \texttt{y} 共用的那份清單，
\texttt{print(x)} 會印出 \texttt{['yellow', 'green']}。
\item 把 \texttt{new\_list = original} 改成 \texttt{new\_list = original.copy()}
（或 \texttt{list(original)}、\texttt{original[:]}，三種寫法效果一樣），才會另外造出
一份獨立的清單。
\item \texttt{q = p} 是別名，兩者是同一份清單，\texttt{p is q} 是 \texttt{True}；
\texttt{r = p.copy()} 是另外造出來的獨立清單，\texttt{p is r} 是 \texttt{False}。
\end{enumerate}

\paragraph{第~\ref{sec:sort-vs-sorted}~節}
\begin{enumerate}
\item \texttt{nums.sort()} 直接改動 \texttt{nums} 本身，\texttt{print(nums)} 會印出
\texttt{[1, 3, 5, 8]}。
\item \texttt{sorted(nums)} 不會動到 \texttt{nums}，所以 \texttt{nums} 還是
\texttt{[5, 3, 8, 1]}，\texttt{result} 才是排好序的 \texttt{[1, 3, 5, 8]}。
\end{enumerate}

\paragraph{第~\ref{sec:sort-key}~節}
\begin{enumerate}
\item \texttt{nums.sort(reverse=True)} 由大到小排序，\texttt{nums} 會變成
\texttt{[9, 4, 3, 1]}。
\item \texttt{sorted(words, key=len)} 依字元數由少到多排，結果是
\texttt{['cat', 'dog', 'elephant']}（\texttt{cat} 跟 \texttt{dog} 字數一樣，維持原本
的相對順序）。
\end{enumerate}

\paragraph{第~\ref{sec:list-tools}~節}
\begin{enumerate}
\item \texttt{a + b} 的結果是 \texttt{[1, 2, 3]}；\texttt{+} 造出的是一份新清單，
\texttt{a} 本身不會被改變，還是 \texttt{[1, 2]}。
\item \texttt{[1, 2] * 3} 的結果是 \texttt{[1, 2, 1, 2, 1, 2]}。
\item \texttt{sum(nums)} 是 \texttt{108}，\texttt{max(nums)} 是 \texttt{42}。
\item \texttt{fruits.index("guava")} 因為清單裡沒有「guava」會丟出 \texttt{ValueError}；
\texttt{fruits.count("guava")} 不會報錯，會印出 \texttt{0}。兩者行為不同是因為「找位置」
在找不到時沒有合理的答案可以回傳，但「數出現幾次」允許答案是零次。
\end{enumerate}

\paragraph{第~\ref{sec:nested-list}~節}
\begin{enumerate}
\item \texttt{matrix[2]} 是最外層第三項，會是 \texttt{[5, 6]}。
\item \texttt{matrix[2][0]} 再往裡取第一個，會是 \texttt{5}。
\item 寫 \texttt{board[1][1]}：先用 \texttt{[1]} 拿到第二列 \texttt{["O", "X"]}，再用
\texttt{[1]} 拿到那一列的第二欄 \texttt{"X"}。
\item \texttt{total} 會是 \texttt{1 + 2 + 3 + 4 = 10}。
\end{enumerate}

\paragraph{第~\ref{sec:nested-alias}~節}
\begin{enumerate}
\item \texttt{rows} 會變成 \texttt{[['a'], ['a'], ['a']]}。因為 \texttt{[[]] * 3}
只造了一份空清單，三個位置的名字牌都指著同一份，對其中一個 \texttt{append}，另外兩個
看到的也是同一份、當然一起變。
\item 改寫成 \texttt{rows = [[] for \_ in range(3)]}，每一輪迴圈都重新造一份全新、獨立
的空清單，三列才會真正各自獨立。
\end{enumerate}

\paragraph{第~\ref{sec:list-comprehension}~節}
\begin{enumerate}
\item \texttt{lengths = [len(w) for w in words]}。
\item 填 \texttt{n * 2}：\texttt{doubled = [n * 2 for n in [1, 2, 3]]}。
\end{enumerate}

\chapter{Dict 與 Set：用「線索」而不是「位置」找答案}
\label{ch:dict-set}

想找一個朋友的電話號碼時，你腦中不會想「這是我認識的第幾個人」，而是直接用他的名字去
想：名字才是你真正拿來查東西的線索，「第幾個」根本不重要。前面幾章學過的 list，只能用
「第幾個」這種位置去查資料，遇到「我想用名字查電話」這種需求就沒轍了。

這一章要教的 dict（字典），就是專門解決這件事的容器：讓你用任何一個有意義的「線索」去
對應一個「答案」，而不是死板地用位置查東西。你會先學會 dict 的骨架，也就是怎麼造一個
dict、怎麼用線索查到答案、查不到的時候會發生什麼事、以及怎麼把資料放進去、改掉、拿掉。
\index{dict}

接著會往上疊一層，教 dict 真正好用的地方：

\begin{itemize}
\item 怎麼把整個 dict 掃過一遍（不只查單一個線索）
\item 怎麼一次合併大量資料進去
\item 怎麼一邊拿資料一邊把它刪掉
\item 怎麼用 dict 做一件很常見的事，也就是數東西出現幾次
\end{itemize}

還有一個之後會反覆看到的寫法習慣：與其每次查之前都先問「這個線索存在嗎？」，Python
更常見的做法是直接假設它存在，用 \texttt{get()} 順便準備好一個備用答案。這個習慣現在
只是先讓你眼熟，完整的道理要等到第 \ref{ch:files-exceptions} 章講到檔案與例外處理時才會攤開來講。

最後會轉到一個新角色：set（集合）。如果把 dict 想成「線索對答案」的兩欄表格，set 就是
把「答案」那一欄整個拿掉，只留下線索本身，你不在乎每個東西對應什麼，只在乎「這個東西
在不在裡面」，還有「兩群東西裡面，哪些是共同的、哪些只有一邊有、哪些是任何一邊都算數
的」。

這種需求生活裡也很常見，比如兩個社團的成員名單，想知道誰兩邊都參加、誰只參加了
其中一邊、或是把兩份名單合併成一份不重複的名單。set 的操作方式跟 dict 有很多地方相通，
畢竟它本質上就是「只留 key 的 dict」，所以放在 dict 之後教，你可以直接把剛學過的直覺
搬過來用，不用重新建立一套心智模型。

\section{Dict 是什麼：建立字典、用線索查資料}
\label{sec:dict-basics}

要建立一個 dict，你在大括號裡放進一組一組的「線索：答案」，中間用冒號分開，好幾組之間
用逗號隔開。查資料的方式跟 list 很像，一樣用中括號，只是裡面放的線索是你自己定義的，
不再是位置：正式的說法，這個線索叫做 key（鍵），對應的答案叫做 value（值）。
\index{key}
\index{value}

而且用 key 查 value 的速度通常很快，不太受 dict 裡東西數量的影響，這跟 list 得一個一個
位置比對過去很不一樣，但背後的原理這裡先不深究。

\begin{code}{python}{建立與存取 dict}
my_dict = {key1: value1, key2: value2}
my_dict[key1]
\end{code}

如果想從一個空字典開始，之後再一筆一筆加資料進去，可以寫 \texttt{my\_dict = \{\}}，或是
用 \texttt{dict()} 這個函式，效果一樣。

\begin{example}{python}{(phonebook.py) 查電話本}
phonebook = {"Alice": "0912-345-678", "Bob": "0933-222-111"}
print(phonebook["Alice"])
\end{example}

\iosample{}{0912-345-678 \\}

如果你查的線索根本不在 dict 裡面，會發生什麼事？

\begin{example}{python}{(phonebook-error.py) 查不存在的名字}
phonebook = {"Alice": "0912-345-678"}
print(phonebook["Bob"])
\end{example}

存檔執行後，螢幕上不會印出電話號碼，而是會看到這樣的訊息：

\begin{verbatim}
Traceback (most recent call last):
  File "phonebook-error.py", line 2, in <module>
    print(phonebook["Bob"])
          ~~~~~~~~~^^^^^^^^
KeyError: 'Bob'
\end{verbatim}

\texttt{KeyError} 的意思很直接：這個 key 在 dict 裡找不到。\texttt{phonebook} 這個字典
裡根本沒有 \texttt{"Bob"} 這個線索，Python 也不會幫你猜一個相近的答案，就直接讓程式
停下來、告訴你出錯了。

\index{KeyError}
查不到答案時，Python 選擇直接讓程式停下來報錯，而不是默默給你一個奇怪的預設值，這種
「先做、錯了才處理」的態度，之後你會在 Python 裡看到更多次，第 \ref{ch:files-exceptions} 章會把整套完整的
道理講清楚，這裡你只要先認得 \texttt{KeyError} 長什麼樣子就好。

\begin{drills}[label=drills:dict-basics]
\item 這段程式執行後，螢幕上會印出什麼？
\begin{verbatim}
grades = {"Alice": 90, "Bob": 85}
print(grades["Bob"])
\end{verbatim}
\item 這段程式執行後會發生什麼事？如果會出錯，是哪一種錯誤？
\begin{verbatim}
grades = {"Alice": 90, "Bob": 85}
print(grades["Taylor"])
\end{verbatim}
\item 如果你想從一個空字典開始，之後再一筆一筆加資料進去，除了 \texttt{\{\}}，還可以用
哪一個函式建立空 dict？
\end{drills}

\section{新增、修改、刪除 Dict 裡的資料}
\label{sec:dict-modify}

光會查還不夠，dict 建好之後你多半會想繼續維護它：加一筆新的資料進去、把舊資料改掉、
不要的資料整筆刪掉。這三件事都很直接，用的語法你其實都已經見過。

\begin{code}{python}{新增、修改、刪除 dict 裡的資料}
my_dict[key] = value   # 如果 key 不存在就新增；如果 key 已存在就覆蓋原本的值
del my_dict[key]       # 整筆刪除這個 key
\end{code}

\begin{example}{python}{(picnic-list.py) 野餐要帶的東西}
picnic_items = {
    "Alice": ["drinks", "fruit"],
    "Bob": ["napkins", "hula hoop"],
}

picnic_items["Alice"].append("sandwich")   # 幫 Alice 的清單再加一項
picnic_items["Taylor"] = ["grill", "charcoal"]  # 新增一個人，以及他要帶的東西
del picnic_items["Bob"]                  # Bob 說他不來了，把他的資料刪掉

print(picnic_items)
\end{example}

\iosample{}{\{'Alice': ['drinks', 'fruit', 'sandwich'], 'Taylor': ['grill', 'charcoal']\} \\}

留意 \texttt{picnic\_items["小美"]} 本身是一份 list，所以可以直接對它呼叫
\texttt{.append()}。這也順便提醒了你一件事：dict 的 value 可以是任何型別，不只是數字
或字串，list、甚至另一個 dict 都可以，之後處理比較複雜的資料時會常常用到這一點。

\begin{drills}[label=drills:dict-modify]
\item 這段程式執行後，\texttt{scores} 最後會是什麼？
\begin{verbatim}
scores = {"Alice": 70}
scores["Alice"] = 85
scores["Bob"] = 92
del scores["Alice"]
print(scores)
\end{verbatim}
\item 這段程式執行後會印出什麼？
\begin{verbatim}
cart = {"apple": ["red", "green"]}
cart["apple"].append("yellow")
print(cart)
\end{verbatim}
\item 如果你想在 dict 裡新增一筆從來沒出現過的 key，要用哪一種語法（不是 \texttt{del}，
也不是 \texttt{print}）？
\end{drills}

\section{一次合併多筆資料：update()}
\label{sec:dict-update}

上一節的 \texttt{my\_dict[key] = value} 一次只能新增或改掉一筆。如果你手上已經有另一份
dict，想把它整個併進來（有重複的 key 就用新的蓋掉舊的），一筆一筆用迴圈慢慢寫也可以，
但 Python 幫你把這件事包成一個方法：\texttt{update()}。
\index{update()}

\begin{code}{python}{update()：把另一份 dict 併進來}
my_dict.update(other_dict)
\end{code}

\begin{example}{python}{(merge-settings.py) 合併預設設定與使用者設定}
default_settings = {"volume": 50, "brightness": 80, "language": "English"}
user_settings = {"volume": 70, "font_size": 16}

default_settings.update(user_settings)
print(default_settings)
\end{example}

\iosample{}{\{'volume': 70, 'brightness': 80, 'language': 'English', 'font\_size': 16\} \\}

\texttt{update()} 的邏輯很直覺：\texttt{user\_settings} 裡有的 key，如果
\texttt{default\_settings} 也有（例如「volume」），就用 \texttt{user\_settings} 的值蓋過去；
如果是 \texttt{default\_settings} 原本沒有的 key（例如「font\_size」），就直接新增進來。
沒被 \texttt{user\_settings} 提到的 key（例如「brightness」「language」）維持原樣不變。留意這個
方法是直接原地修改 \texttt{default\_settings} 本身，不會回傳一份新的 dict，跟第 \ref{ch:lists} 章
list 的 \texttt{.append()}、\texttt{.sort()} 是同一種脾氣。

\begin{drills}[label=drills:dict-update]
\item 這段程式執行後，\texttt{a} 最後會是什麼？
\begin{verbatim}
a = {"x": 1, "y": 2}
b = {"y": 20, "z": 30}
a.update(b)
print(a)
\end{verbatim}
\item \texttt{update()} 呼叫完之後，\texttt{b} 這份 dict 本身會不會被改動？
\end{drills}

\section{一邊拿資料一邊刪除：pop()}
\label{sec:dict-pop}

\texttt{del} 可以整筆刪除一個 key，但刪掉之後那個值就直接消失了，你沒機會再用到它。很多
情境下你其實想要「把這筆資料拿出來用，同時把它從 dict 裡移除」，例如處理完一筆訂單就把
它從待處理清單拿掉。這正是 \texttt{pop()} 在做的事：刪除的同時，把被刪除的 value
回傳給你。
\index{pop()}

\begin{code}{python}{pop()：刪除並拿回這個 key 對應的值}
my_dict.pop(key)            # 如果 key 不存在會丟出 KeyError，跟中括號查詢一樣
my_dict.pop(key, default)   # 如果 key 不存在不會報錯，會回傳你給的預設值
\end{code}

\begin{example}{python}{(serve-orders.py) 依序處理點餐佇列}
orders = {"table3": "beef noodles", "table5": "braised pork rice"}

dish = orders.pop("table3")
print(dish)
print(orders)

print(orders.pop("table9", "no such table"))
\end{example}

\iosample{}{beef noodles \\ \{'table5': 'braised pork rice'\} \\ no such table \\}

\texttt{orders.pop("table3")} 做了兩件事：把 \texttt{"table3"} 這筆從 \texttt{orders} 裡
刪掉，同時把它的 value（\texttt{"beef noodles"}）回傳出來存進 \texttt{dish}。之後
\texttt{orders} 裡就只剩下 \texttt{"table5"} 這一筆。

第二次呼叫 \texttt{orders.pop("table9",
"no such table")} 因為 \texttt{"table9"} 根本不存在，但因為有給第二個引數當預設值，程式不會
出錯，直接拿到 \texttt{"no such table"}。如果沒給預設值又查一個不存在的 key，
\texttt{pop()} 一樣會丟出 \texttt{KeyError}，行為跟中括號查詢一致，跟你在
\ref{sec:dict-get} 節會看到的 \texttt{get()} 剛好相反：\texttt{get()} 不給預設值時預設
值是 \texttt{None} 不會報錯，\texttt{pop()} 不給預設值則會直接報錯。

\begin{drills}[label=drills:dict-pop]
\item 這段程式執行後，依序會印出什麼？
\begin{verbatim}
stock = {"apple": 10, "banana": 5}
amount = stock.pop("banana")
print(amount)
print(stock)
\end{verbatim}
\item 這段程式執行後會印出什麼？
\begin{verbatim}
stock = {"apple": 10}
print(stock.pop("orange", 0))
\end{verbatim}
\item \texttt{stock.pop("orange")}（沒有給第二個引數，且 \texttt{"orange"} 不在
\texttt{stock} 裡）執行後會發生什麼事？
\end{drills}

\section{走訪整個 Dict：keys()、values()、items()}
\label{sec:dict-iterate}

現在你已經會對單一線索查、改 dict 了。但很多時候你想要的是把整個 dict 掃過一遍，
不只是查一筆，例如列出所有人跟他們各自要帶的東西。這裡會用到三個方法：\texttt{keys()}
拿到全部的線索、\texttt{values()} 拿到全部的答案、\texttt{items()} 拿到「線索、答案」
成對的組合。
\index{keys()}
\index{values()}
\index{items()}

\begin{code}{python}{走訪整個 dict}
for key in my_dict:
    ...

for key, value in my_dict.items():
    ...
\end{code}

直接對 dict 做 \texttt{for} 迴圈，預設走的是 key，效果跟寫
\texttt{for key in my\_dict.keys()} 一樣，只是 \texttt{.keys()} 這幾個字通常可以省略。
想要同時拿到 key 跟 value，就用 \texttt{.items()}，配合兩個變數名字一次接住一組。

\begin{example}{python}{(picnic-checklist.py) 列出每個人要帶的東西}
picnic_items = {
    "Alice": ["drinks", "fruit", "sandwich"],
    "Taylor": ["grill", "charcoal"],
}

for name, items in picnic_items.items():
    print(f"{name} is bringing: {', '.join(items)}")
\end{example}

\iosample{}{Alice is bringing: drinks, fruit, sandwich \\ Taylor is bringing: grill, charcoal \\}

順帶一提，dict 的賦值也是貼名字牌的行為，跟第 \ref{ch:lists} 章 list 的別名行為一模一樣：
\texttt{other = picnic\_items} 之後改 \texttt{other} 的內容，\texttt{picnic\_items}
也會跟著變，因為兩個名字指的是同一份 dict，不是各自獨立的一份。真的想要另外一份互不
影響的複製，要用 \texttt{.copy()}。

\begin{drills}[label=drills:dict-iterate]
\item 這段程式執行後，螢幕上會依序印出什麼？
\begin{verbatim}
prices = {"apple": 30, "banana": 15}
for name in prices:
    print(name)
\end{verbatim}
\item 這段程式執行後，螢幕上會依序印出什麼？
\begin{verbatim}
prices = {"apple": 30, "banana": 15}
for name, price in prices.items():
    print(f"{name}: {price}")
\end{verbatim}
\end{drills}

\section{用 Dict 數東西出現幾次}
\label{sec:dict-counter}

dict 有一個特別常見的用途：數東西出現幾次。例如統計一段文字裡每個字元各出現幾次。你可能
會很直覺地想：先檢查這個字元有沒有出現過，沒出現過就先設成 0，出現過就直接加一。

\begin{example}{python}{(word-count.py) 統計文字裡每個字元出現幾次}
text = "banana"
counts = {}

for ch in text:
    if ch not in counts:
        counts[ch] = 0
    counts[ch] += 1

print(counts)
\end{example}

\iosample{}{\{'b': 1, 'a': 3, 'n': 2\} \\}

這種「先檢查、再決定」的寫法沒有錯，能算出正確答案。但下一節你會看到一個更常見的
Python 寫法，可以把中間那段判斷式縮成一行。

\section{get()：與其先檢查，不如直接查}
\label{sec:dict-get}

剛才那段程式，每次都要先問一次「這個字元出現過嗎？」，才決定要不要先把值設成 0。Python
比較常見的寫法，是省略掉「先問」這一步，直接呼叫 \texttt{get(key, default)}：如果 key
存在，就拿到它目前的值；如果不存在，就直接給你準備好的預設值，不會噴 \texttt{KeyError}。
\index{get()}

\begin{code}{python}{get()：查不到就用預設值}
my_dict.get(key, default)
\end{code}

\begin{example}{python}{(word-count-get.py) 用 get() 精簡統計程式}
text = "banana"
counts = {}

for ch in text:
    counts[ch] = counts.get(ch, 0) + 1

print(counts)
\end{example}

\iosample{}{\{'b': 1, 'a': 3, 'n': 2\} \\}

跟上一節的版本比對一下，答案完全一樣，但省掉了一整段 \texttt{if} 判斷。

\stylepoint{與其先寫 \texttt{if key not in my\_dict} 再決定要不要處理，Python 更常見的
做法是直接用 \texttt{get(key, default)} 把「查不到怎麼辦」一次講清楚，這個習慣你會在
後面章節一直看到，先讓自己眼熟這個寫法就好，完整的道理留到第 \ref{ch:files-exceptions} 章講檔案與例外處理時
再攤開。}

如果呼叫 \texttt{get()} 時完全不給第二個引數，查不到的話會回傳 \texttt{None}，不會報錯，
也不會像不給預設值的 \texttt{pop()} 那樣丟出 \texttt{KeyError}。這個差異值得記住：
\texttt{get()} 天生就是「查不到也沒關係」的個性，\texttt{pop()} 則是預設「查不到就該讓
你知道」，除非你自己明確給了第二個引數。

\begin{drills}[label=drills:dict-get]
\item 這段程式執行後會印出什麼？
\begin{verbatim}
text = "apple"
counts = {}
for ch in text:
    if ch not in counts:
        counts[ch] = 0
    counts[ch] += 1
print(counts)
\end{verbatim}
\item 把下面這段「先檢查再存取」的程式碼，改寫成用 \texttt{get()} 的版本：
\begin{verbatim}
if "Alice" in grades:
    score = grades["Alice"]
else:
    score = 0
\end{verbatim}
\item 這段程式執行後會印出什麼？
\begin{verbatim}
grades = {"Alice": 90}
print(grades.get("Bob", "not registered"))
\end{verbatim}
\end{drills}

\section{setdefault()：查不到就先設好、順便拿回來}
\label{sec:dict-setdefault}

\texttt{get()} 適合用在「拿到值之後想直接拿去用」（例如加 1 累計次數），但如果每個
key 對應的 value 本身是一份 list，你想要的往往不只是「查到就用、查不到就給預設值」，
而是「查不到的話，順便把這個預設值真的存進 dict 裡」，這樣才能接著對它呼叫
\texttt{.append()}。這正是 \texttt{setdefault()} 要處理的情境。
\index{setdefault()}

\begin{code}{python}{setdefault()：查不到就先存進去，再回傳}
my_dict.setdefault(key, default)
\end{code}

\texttt{setdefault(key, default)} 的行為分兩種情況：如果 \texttt{key} 已經在
\texttt{my\_dict} 裡，直接回傳它目前的 value，\texttt{default} 完全不會被用到；如果
\texttt{key} 不在裡面，就先把 \texttt{my\_dict[key] = default} 這件事做掉，再把
\texttt{default} 回傳給你。跟 \texttt{get()} 最大的差別在這裡：\texttt{get()} 查不到時
只是「借你看一下」預設值，並不會真的存進 dict；\texttt{setdefault()} 查不到時會真的把
預設值寫進去。

\begin{example}{python}{(group-by-letter.py) 依照開頭第一個字分組}
words = ["apple", "banana", "blueberry", "watermelon", "cherry"]
groups = {}

for word in words:
    first = word[0]
    groups.setdefault(first, []).append(word)

print(groups)
\end{example}

\iosample{}{\{'a': ['apple'], 'b': ['banana', 'blueberry'], 'w': ['watermelon'], 'c': ['cherry']\} \\}

每一輪迴圈裡，\texttt{groups.setdefault(first, [])} 先確保 \texttt{groups[first]} 這個
key 一定存在（第一次遇到某個開頭字時，會自動存進一份空 list），再把回傳值（不管是原本
就有的 list，還是剛存進去的空 list）直接接上 \texttt{.append(word)}。整個過程只需要一行，
不用像前面 \nameref{sec:dict-counter} 節那樣先寫一段 \texttt{if} 判斷 key 存不存在。

跟 \texttt{get()} 用在「數字累加」的情境相比，\texttt{setdefault()} 更常用在這種「value 是
可變容器、要先確保它存在才能對它動手」的情境，兩者是同一種精神在不同場景下的變化。

\begin{drills}[label=drills:dict-setdefault]
\item 這段程式執行後，\texttt{groups} 最後會是什麼？
\begin{verbatim}
groups = {"sweet": ["apple"]}
groups.setdefault("sweet", []).append("banana")
groups.setdefault("sour", []).append("lemon")
print(groups)
\end{verbatim}
\item \texttt{setdefault()} 跟 \texttt{get()} 最主要的差別是什麼？（提示：查不到的時候，
兩者對 dict 本身分別做了什麼事）
\end{drills}

\section{Set 是什麼：只留線索，不留答案}
\label{sec:set-basics}

如果把 dict 想成「線索對答案」的兩欄表格，set 就是把答案那一欄整個拿掉，只留下線索
本身，你不在乎每個東西對應什麼，只在乎「這個東西在不在裡面」。
\index{set}

\begin{code}{python}{建立與操作 set}
my_set = {item1, item2, item3}
my_set.add(item)      # 把 item 加進 set 裡
my_set.remove(item)   # 從 set 裡把 item 刪掉
item in my_set
\end{code}

這裡要特別小心一件事：空的大括號 \texttt{\{\}} 造出來的是空 dict，不是空 set，這是很多人
第一次用 set 時會踩到的陷阱。

\begin{example}{python}{(empty-set-trap.py) 空大括號其實是空 dict}
empty = {}
print(type(empty))
\end{example}

\iosample{}{<class 'dict'> \\}

想要一個真正的空 set，要寫 \texttt{empty\_set = set()}。

\begin{example}{python}{(club-members.py) 社團成員名單}
book_club = {"Alice", "Bob", "Taylor"}
book_club.add("Charlie")
book_club.remove("Bob")

print("Alice" in book_club)
print(book_club)
\end{example}

\iosample{}{True \\ \{'Alice', 'Charlie', 'Taylor'\} \\}

\keypoint{如果你自己執行這段程式，印出來的 \texttt{book\_club} 順序可能跟這裡不一
樣，這是正常現象，跟到目前為止看過的 dict 不同，set 不會記得元素被加進去的順序。你只
要在乎某個東西在不在裡面（用 \texttt{in} 檢查）、以及裡面實際有哪些元素就好，順序本身
不重要。之後這一章只要再出現印出 set 本身的例子，都會有同樣的情況，不會每次都重複
提醒。}

set 造出來之後如果放進重複的元素，Python 也只會留一份，不會噴錯、也不會出現兩個一模
一樣的東西：\texttt{\{"Alice", "Alice", "Bob"\}} 造出來的 set，內容跟
\texttt{\{"Alice", "Bob"\}} 完全一樣。這個特性也直接呼應了下一節要講的第一個用途。

\begin{drills}[label=drills:set-basics]
\item 這段程式執行後會印出什麼？
\begin{verbatim}
colors = {"red", "red", "blue"}
print(len(colors))
\end{verbatim}
\item 空大括號 \texttt{\{\}} 造出來的是什麼型別？想造一個真正的空 set 要怎麼寫？
\end{drills}

\section{Set 的兩個常見用途：去重複、比較兩群東西}
\label{sec:set-ops}

認識了 set 的基本操作後，來看 set 最常見的兩個用途：把一份可能有重複的資料去掉重複、
以及拿兩個 set 互相比較。

\begin{example}{python}{(remove-duplicates.py) 去除清單裡重複的名字}
signups = ["Alice", "Bob", "Alice", "Taylor", "Bob"]
unique_names = set(signups)
print(unique_names)
\end{example}

\iosample{}{\{'Bob', 'Taylor', 'Alice'\} \\}

把一份 list 直接丟進 \texttt{set()}，重複的東西就自動只留一份，這是 Python 裡最常見
的去重複寫法（同樣的，順序可能因人而異）。

\begin{example}{python}{(club-overlap.py) 兩個社團的共同成員}
book_club = {"Alice", "Bob", "Taylor"}
hiking_club = {"Bob", "Charlie", "Alice"}

both = book_club & hiking_club
either = book_club | hiking_club

print(both)
print(either)
\end{example}

\iosample{}{\{'Alice', 'Bob'\} \\ \{'Charlie', 'Taylor', 'Bob', 'Alice'\} \\}

\texttt{\&} 算出兩邊都有的東西（交集）、\texttt{|} 算出兩邊合起來、不重複的東西（聯
集）。想知道「兩個社團都參加的人」用 \texttt{\&}，想把「兩份名單合併成一份不重複的名
單」用 \texttt{|}，兩種需求生活裡都很常見。
\index{交集}
\index{聯集}

\begin{drills}[label=drills:set-ops]
\item 這段程式執行後會依序印出什麼？
\begin{verbatim}
club = {"Alice", "Bob", "Taylor"}
print("Bob" in club)
print("Charlie" in club)
\end{verbatim}
\item 這段程式執行後會印出什麼？
\begin{verbatim}
votes = ["apple", "banana", "apple", "orange", "banana"]
choices = set(votes)
print(len(choices))
\end{verbatim}
\item 這段程式執行後會依序印出什麼？
\begin{verbatim}
a = {"Alice", "Bob", "Taylor"}
b = {"Bob", "Alice", "Charlie"}
print(len(a & b))
print(len(a | b))
\end{verbatim}
\end{drills}

\section{Set 的另外兩種比較：差集與對稱差集}
\label{sec:set-diff}

上一節的交集（\texttt{\&}）跟聯集（\texttt{|}）已經能回答「兩邊都有誰」跟「兩邊合起來
有誰」，但還有兩種同樣常見的問法：「只在這一邊、不在另一邊的是誰」，還有「只在其中一邊、
不是兩邊都有的是誰」。這兩個問法分別對應差集跟對稱差集。
\index{差集}
\index{對稱差集}

\begin{code}{python}{差集與對稱差集}
my_set - other_set   # 差集：在 my_set 裡、但不在 other_set 裡的東西
my_set ^ other_set    # 對稱差集：剛好只出現在其中一個 set 裡的東西
\end{code}

\begin{example}{python}{(club-diff.py) 只參加書友會、以及只參加一邊的人}
book_club = {"Alice", "Bob", "Taylor"}
hiking_club = {"Bob", "Charlie", "Alice"}

only_book = book_club - hiking_club
only_one_club = book_club ^ hiking_club

print(only_book)
print(only_one_club)
\end{example}

\iosample{}{\{'Taylor'\} \\ \{'Taylor', 'Charlie'\} \\}

\texttt{book\_club - hiking\_club} 只留下在 \texttt{book\_club} 裡、但沒出現在
\texttt{hiking\_club} 裡的人，也就是「只參加書友會、沒參加健行社」的 Taylor。留意減法的
方向會影響結果：\texttt{hiking\_club - book\_club} 算出來的會是只參加健行社的 Charlie，
跟前面算出來的不是同一組人，這點跟數字的減法一樣，交換順序答案就不一樣。

\texttt{book\_club \^{} hiking\_club} 則是把兩邊都算進來，但排除掉兩邊都有的人（Alice、
Bob），只留下「剛好只在其中一邊出現」的 Taylor 跟 Charlie，效果等於聯集扣掉交集。

\begin{drills}[label=drills:set-diff]
\item 這段程式執行後會印出什麼？
\begin{verbatim}
a = {"apple", "banana", "orange"}
b = {"banana", "watermelon"}
print(a - b)
\end{verbatim}
\item 這段程式執行後會印出什麼？
\begin{verbatim}
a = {"apple", "banana", "orange"}
b = {"banana", "watermelon"}
print(len(a ^ b))
\end{verbatim}
\item \texttt{a - b} 跟 \texttt{b - a} 算出來的結果會不會一樣？
\end{drills}

\section{安全地刪除、快速地判斷有沒有}
\label{sec:set-membership}

前面用過的 \texttt{remove()}，如果要刪除的東西根本不在 set 裡，會直接報錯：

\begin{example}{python}{(remove-error.py) 刪除一個不存在的成員}
book_club = {"Alice", "Bob", "Taylor"}
book_club.remove("David")
\end{example}

\begin{verbatim}
Traceback (most recent call last):
  File "remove-error.py", line 2, in <module>
    book_club.remove("David")
    ~~~~~~~~~~~~~~~~^^^^^^^^
KeyError: 'David'
\end{verbatim}

有意思的是，即使你刪的是 set 裡的元素，錯誤訊息顯示的仍然是 \texttt{KeyError}，因為
set 內部運作的機制跟 dict 的 key 很接近（只是沒有對應的 value）。如果你不確定要刪除的
東西存不存在、又不想每次都先寫一次 \texttt{if ... in ...} 判斷，可以改用
\texttt{discard()}：找不到就什麼事都不做，不會報錯。
\index{discard()}

\begin{code}{python}{discard()：找不到也不會報錯的刪除}
my_set.discard(item)
\end{code}

\begin{example}{python}{(safe-remove.py) 用 discard() 安全地刪除}
book_club = {"Alice", "Bob", "Taylor"}
book_club.discard("Bob")
book_club.discard("David")

print(book_club)
\end{example}

\iosample{}{\{'Alice', 'Taylor'\} \\}

第一次 \texttt{discard("Bob")} 因為「Bob」真的在 \texttt{book\_club} 裡，效果跟
\texttt{remove()} 一樣，直接刪掉；第二次 \texttt{discard("David")} 因為「David」根本不在
裡面，程式照樣往下跑，不會像 \texttt{remove()} 那樣中斷。什麼時候該用哪一個，看你要不要
被提醒「這個東西其實不存在」：確定一定存在、想在少了它時被提醒就用 \texttt{remove()}；
只是想確保它「最後一定不在裡面」、不在乎它原本有沒有，就用 \texttt{discard()}。

除了刪除，set 另一個常被用到的地方是快速判斷「這個東西在不在一大堆東西裡面」。用
\texttt{in} 檢查一個東西在不在 set 裡，速度跟 dict 用 key 查 value 一樣快，不太受 set
裡東西數量的影響，這跟 list 得從頭一個一個比對過去很不一樣。

如果你手上有一份很長的
名單，需要反覆檢查「這個名字有沒有出現過」，比起每次都對 list 做 \texttt{in} 檢查，先
把名單轉成 set 再檢查，會是更常見、也更划算的寫法：

\begin{example}{python}{(fast-membership.py) 用 set 快速檢查名字有不有效}
valid_names = ["Alice", "Bob", "Taylor", "Charlie"]
valid_set = set(valid_names)

print("Taylor" in valid_set)
print("David" in valid_set)
\end{example}

\iosample{}{True \\ False \\}

\begin{drills}[label=drills:set-membership]
\item 這段程式執行後會發生什麼事？如果會出錯，是哪一種錯誤？
\begin{verbatim}
fruits = {"apple", "banana"}
fruits.remove("orange")
\end{verbatim}
\item 把上一題的 \texttt{remove("orange")} 換成 \texttt{discard("orange")}，執行後會發生
什麼事？
\item 手上有一份很長的名單，需要反覆檢查某個名字在不在裡面，為什麼通常會先把它轉成
set，而不是直接對原本的 list 做 \texttt{in} 檢查？
\end{drills}

\section{一行造出 dict 或 set：dict comprehension、set comprehension}
\label{sec:comprehension}
\index{dict comprehension}
\index{set comprehension}

第 \ref{ch:lists} 章教過的 list comprehension，把「造一份空清單、跑迴圈、逐一加東西進去」濃縮成一行。
同樣的邏輯，其實也可以直接套用在 dict 跟 set 上，只是把方括號換成大括號：造 set 的寫法
省略冒號，跟 list comprehension 幾乎一樣；造 dict 的寫法則需要冒號分開 key 跟 value。

\begin{code}{python}{set comprehension 與 dict comprehension}
{expression for item in iterable}              # 這是 set comprehension（集合生成式）
{key_expr: value_expr for item in iterable}    # 這是 dict comprehension（字典生成式）
\end{code}

\begin{example}{python}{(comprehension.py) 用一行分別造出 set 與 dict}
signups = ["Alice", "Bob", "Alice", "Taylor", "Bob"]

unique_names = {name for name in signups}
name_lengths = {name: len(name) for name in signups}

print(unique_names)
print(name_lengths)
\end{example}

\iosample{}{\{'Taylor', 'Alice', 'Bob'\} \\ \{'Alice': 5, 'Bob': 3, 'Taylor': 6\} \\}

\texttt{\{name for name in signups\}} 跟 \nameref{sec:set-ops} 節用
\texttt{set(signups)} 去重複，效果完全一樣，只是寫法換成了 comprehension；
\texttt{\{name: len(name) for name in signups\}} 則是幫每個名字造出一筆「名字對應到它
的字數」的資料，重複出現的名字（例如「Alice」出現兩次）在 dict 裡自然只會留一份，因為
同一個 key 後面出現的賦值會直接蓋掉前一次。

跟 list comprehension 一樣，這兩種寫法你
現在看得懂、知道長什麼樣子就好，之後章節用得上時會愈來愈熟。

\begin{drills}[label=drills:comprehension]
\item 這段程式執行後會印出什麼？
\begin{verbatim}
nums = [1, 2, 2, 3, 3, 3]
unique = {n for n in nums}
print(len(unique))
\end{verbatim}
\item 把下面這段用迴圈寫的程式，改寫成一行 dict comprehension：
\begin{verbatim}
squares = {}
for n in [1, 2, 3, 4]:
    squares[n] = n * n
\end{verbatim}
\end{drills}

\section*{練習解答}
\addcontentsline{toc}{section}{練習解答}

\paragraph{第~\ref{sec:dict-basics}~節}
\begin{enumerate}
\item 會印出：85。
\item 會出現 \texttt{KeyError: 'Taylor'}，因為 \texttt{grades} 這個字典裡沒有這個 key。
\item \texttt{dict()}。
\end{enumerate}

\paragraph{第~\ref{sec:dict-modify}~節}
\begin{enumerate}
\item \texttt{\{'Bob': 92\}}，因為「Alice」先被改成 85，後來又被 \texttt{del} 整筆刪掉，
只剩「Bob」。
\item \texttt{\{'apple': ['red', 'green', 'yellow']\}}，因為 value 本身是一份 list，
\texttt{.append()} 直接改的是那份 list。
\item 用 \texttt{my\_dict[key] = value} 這種索引賦值語法，key 不存在時就會新增一筆。
\end{enumerate}

\paragraph{第~\ref{sec:dict-update}~節}
\begin{enumerate}
\item \texttt{\{'x': 1, 'y': 20, 'z': 30\}}，「y」這個 key 兩邊都有，被 \texttt{b} 裡的
值蓋過去；「z」是 \texttt{b} 才有的 key，直接新增進來；「x」沒被提到，維持原樣。
\item 不會，\texttt{update()} 只會修改呼叫它的那份 dict（例子裡是 \texttt{a}），
\texttt{b} 本身完全沒有變動。
\end{enumerate}

\paragraph{第~\ref{sec:dict-pop}~節}
\begin{enumerate}
\item 依序印出：5、\texttt{\{'apple': 10\}}，「banana」對應的值被拿出來存進
\texttt{amount}，同時從 \texttt{stock} 裡刪掉。
\item 印出：0，因為 \texttt{"orange"} 不在 \texttt{stock} 裡，\texttt{pop()} 回傳給的
第二個引數當作預設值。
\item 會出現 \texttt{KeyError: 'orange'}，因為沒有給預設值、\texttt{pop()} 又找不到這個
key。
\end{enumerate}

\paragraph{第~\ref{sec:dict-iterate}~節}
\begin{enumerate}
\item 依序印出：apple、banana（dict 會記得新增的順序）。
\item 依序印出：apple: 30、banana: 15。
\end{enumerate}

\paragraph{第~\ref{sec:dict-get}~節}
\begin{enumerate}
\item \texttt{\{'a': 1, 'p': 2, 'l': 1, 'e': 1\}}。
\item \texttt{score = grades.get("Alice", 0)}。
\item not registered，因為 \texttt{"Bob"} 不在 \texttt{grades} 裡，\texttt{get()} 回傳給的
預設值。
\end{enumerate}

\paragraph{第~\ref{sec:dict-setdefault}~節}
\begin{enumerate}
\item \texttt{\{'sweet': ['apple', 'banana'], 'sour': ['lemon']\}}，「sweet」這個 key 本來就存在，
\texttt{setdefault()} 直接回傳原本的 list 讓 \texttt{.append()} 接上去；「sour」不存在，
先被存成一份空 list 再接上 \texttt{.append()}。
\item \texttt{get()} 查不到時只是「借你看一下」預設值，不會真的把它存進 dict 裡；
\texttt{setdefault()} 查不到時會把預設值真的寫進 dict、之後同一個 key 再查就查得到。
\end{enumerate}

\paragraph{第~\ref{sec:set-basics}~節}
\begin{enumerate}
\item 印出：2，因為 set 裡重複的「red」只會留一份。
\item 空大括號造出來的是空 dict；真正的空 set 要寫 \texttt{set()}。
\end{enumerate}

\paragraph{第~\ref{sec:set-ops}~節}
\begin{enumerate}
\item 依序印出：True、False，「Bob」在 \texttt{club} 裡，「Charlie」不在。
\item 印出：3，\texttt{votes} 裡不重複的東西只有 apple、banana、orange 三種，跟印出來的
順序無關。
\item 依序印出：2、4，\texttt{a} 跟 \texttt{b} 共同的成員是「Alice」跟「Bob」兩個
人，兩邊合起來不重複總共有「Alice」「Bob」「Taylor」「Charlie」四個人。
\end{enumerate}

\paragraph{第~\ref{sec:set-diff}~節}
\begin{enumerate}
\item \texttt{\{'orange', 'apple'\}}，這兩樣只在 \texttt{a} 裡，沒出現在 \texttt{b} 裡。
\item 印出：3，兩邊合起來不重複的東西是 apple、banana、orange、watermelon 共 4 個，扣掉兩邊都有的
banana 1 個，剩下 3 個。
\item 不會一樣，\texttt{a - b} 是只在 \texttt{a} 的東西，\texttt{b - a} 是只在
\texttt{b} 的東西，除非 \texttt{a} 跟 \texttt{b} 完全沒有交集，否則兩者不同。
\end{enumerate}

\paragraph{第~\ref{sec:set-membership}~節}
\begin{enumerate}
\item 會出現 \texttt{KeyError: 'orange'}，因為 \texttt{"orange"} 不在 \texttt{fruits}
裡，\texttt{remove()} 找不到要刪除的東西就會報錯。
\item 什麼事都不會發生，程式照常往下執行，\texttt{fruits} 內容不變，因為
\texttt{discard()} 找不到要刪除的東西時不會報錯。
\item 因為對 set 用 \texttt{in} 檢查東西在不在裡面，速度不太受 set 裡東西數量影響；對
list 用 \texttt{in}，內容越多、檢查得越久，需要一個一個比對過去。
\end{enumerate}

\paragraph{第~\ref{sec:comprehension}~節}
\begin{enumerate}
\item 印出：3，\texttt{nums} 裡不重複的數字只有 1、2、3 三種。
\item \texttt{squares = \{n: n * n for n in [1, 2, 3, 4]\}}。
\end{enumerate}

\chapter{函式}
\label{ch:functions}

前面六章，你寫的每一支程式都是「從頭到尾一路往下唸」。就算中間有 \texttt{if} 分岔、有
\texttt{for} 迴圈繞圈子，整支程式終究是同一塊、攤開來看得到全貌的敘述。這一章要教你的函式
（function），第一次讓你的程式碼可以被「摺起來」：把一段有名字的動作收進一個盒子，之後
只要喊它的名字，不用把內容再抄一次。

這一章分六段推進：

\begin{itemize}
\item 第一段先讓你體驗函式最基本的樣子：怎麼定義一個函式、怎麼呼叫它、怎麼讓它把算好的
東西「送回來」給你而不是只印在螢幕上，以及一次送回不只一個值長什麼樣子
（第 \ref{sec:function-def}～\ref{sec:function-multi-return} 節）。
\item 第二段回頭處理一個很實際的問題：函式如果要接收不只一項資訊，你有哪些方式可以把這些
資訊餵進去、每種方式各自的講究是什麼；這一段的最後兩節還會處理一個更麻煩的情況——連「會
收到幾項資訊」都事先不知道的時候該怎麼寫
（第 \ref{sec:function-positional}～\ref{sec:function-kwargs} 節）。
\item 第三段拉開視角，先談為什麼值得把程式碼拆成函式，再談函式本身其實也可以像數字、字串
一樣被傳來傳去（第 \ref{sec:function-why}～\ref{sec:function-first-class} 節）。
\item 第四段是這一章真正的核心，也是整本書第二次正式呼應「名字牌」這個比喻的地方：你會
看到，呼叫函式、傳參數，其實跟之前學過的「指定變數」是同一套邏輯，而且這套邏輯會直接導出
一個所有 Python 使用者遲早都會踩到的經典陷阱：預設參數如果是一個 list，會發生什麼詭異的事
（第 \ref{sec:function-nametag}～\ref{sec:mutable-default} 節）。
\item 第五段把視角轉向「函式裡面的變數」跟「函式外面的變數」是兩個互不干擾的世界，這件事
既是這一章的最後一塊核心拼圖，也是為什麼函式能被放心「摺起來」重複使用的真正原因
（第 \ref{sec:function-scope} 節）。
\item 第六段是一個簡短的收尾：怎麼替函式留一句自我介紹，讓下一個讀到這段程式碼的人（很
可能是三個月後的你自己）知道它在做什麼（第 \ref{sec:function-docstring} 節）。
\end{itemize}

在往下讀之前，先講清楚一件事：函式，就是一段被取了名字、可以重複呼叫的程式碼。你其實已經
用過別人寫好的函式了。第 \ref{ch:introduction} 章的 \texttt{print}、之後幾章陸續出現的 \texttt{len}、
\texttt{int} 都是函式。這一章要教的，是怎麼自己動手定義一個。

\section{函式的最小樣子：定義與呼叫}
\label{sec:function-def}

要定義一個函式，你需要 \texttt{def}（英文 define 的縮寫）這個關鍵字開頭，接函式名稱、
一對圓括號、一個冒號，底下縮排的部分就是函式被呼叫時會執行的內容，這一段叫做函式的
「主體」（body）。函式名稱的命名規則跟變數一樣，只能用英文字母、數字、底線，不能用數字
開頭；習慣上也會用 snake\_case，也就是全部小寫、單字之間用底線連接的命名風格，這是
Python 函式／變數命名最常見的慣例。
\index{def}
\index{主體}
\index{snake\_case}

\begin{code}{python}{定義並呼叫一個沒有參數的函式}
def function_name():
    code
    code

function_name()
\end{code}

光是寫了 \texttt{def} 這幾行，並不會讓裡面的程式碼真的執行，這幾行只是在告訴 Python
「以後如果有人喊這個名字，要做的事情是這些」。真正的執行，要等到你在下面另外寫一行
「函式名稱\texttt{()}」，也就是呼叫（call）這個函式的那一刻，才會發生。
\index{呼叫}

\begin{example}{python}{(say-hello.py) 定義並呼叫一個打招呼函式}
def say_hello():
    print("Hello! Welcome to this program.")


print("Program starting...")
say_hello()
say_hello()
print("Program finished.")
\end{example}

\iosample{}{Program starting... \\ Hello! Welcome to this program. \\ Hello! Welcome to this program. \\ Program finished. \\}

看這段程式的輸出順序，會發現一件有意思的事：\texttt{say\_hello} 的定義明明寫在最上面，
但「Program starting...」這行字卻先被印出來，直到程式跑到 \texttt{say\_hello()} 這一行，
函式主體裡的那句打招呼才真正出現。而且因為呼叫了兩次，那句話也真的印了兩次。

\keypoint{函式定義的時候，裡面的程式碼並不會執行；只有在函式被呼叫的那一刻，主體裡的內容
才會真正跑起來，而且可以被呼叫任意多次。前面六章你寫的程式都是「寫一行、跑一行」，函式是
你第一次遇到「寫了程式碼，但它先按兵不動，等被叫到才動作」的情形，一開始不習慣很正常。}

\section{帶一個參數的函式}
\label{sec:function-param-basic}

上一節的 \texttt{say\_hello} 每次都只會印出同一句話，沒辦法跟不同的人打招呼。如果想讓
函式在被呼叫的時候，可以帶一項資訊進去、依這項資訊做不同的事，就要在 \texttt{def} 那一行
的括號裡放一個名字，這個名字叫做參數（parameter）；呼叫的時候括號裡實際放進去的值，叫做
引數（argument）。這兩個詞很多人會混著講，不用太計較，但看到 Python 的錯誤訊息裡出現
「argument」，知道它在講的是「呼叫時給的值」就好。
\index{參數}

\begin{code}{python}{定義並呼叫一個帶單一參數的函式}
def function_name(parameter_name):
    code (can use the parameter_name variable)

function_name(argument)
\end{code}

\begin{example}{python}{(greet-name.py) 跟指定的人打招呼}
def say_hello(name):
    print(f"Hello, {name}! Welcome to this program.")


say_hello("Alice")
say_hello("Bob")
\end{example}

\iosample{}{Hello, Alice! Welcome to this program. \\ Hello, Bob! Welcome to this program. \\}

呼叫函式的時候，引數 \texttt{"Alice"} 被對應給了參數 \texttt{name}，函式主體裡用到
\texttt{name} 的地方，用的就是這次呼叫傳進來的值。下一次呼叫換成 \texttt{"Bob"}，
\texttt{name} 裡裝的東西也跟著換了。

如果呼叫函式的時候忘了帶引數，會怎麼樣？

\begin{example}{python}{(greet-name-broken.py) 忘記帶引數}
def say_hello(name):
    print(f"Hello, {name}! Welcome to this program.")


say_hello()
\end{example}

\begin{verbatim}
Traceback (most recent call last):
  File "greet-name-broken.py", line 5, in <module>
    say_hello()
    ~~~~~~~~~^^
TypeError: say_hello() missing 1 required positional argument: 'name'
\end{verbatim}

最後一行說得很白話：\texttt{missing 1 required positional argument: 'name'}，意思是
「少了一個必要的引數，名字叫 \texttt{name}」。\texttt{say\_hello} 定義的時候要求呼叫者
一定要給 \texttt{name}，這裡卻一個都沒給，Python 沒辦法幫你猜這個空缺該填什麼，只好直接
停下來告訴你缺了什麼。這是這一章第一次示範函式相關的錯誤訊息，之後只要引數的數量跟函式
要求的對不上，都會看到類似的 \texttt{TypeError}。

\section{把值送回來：\texttt{return}}
\label{sec:function-return}

前面兩節的函式都只是把結果直接印在螢幕上。如果你想把函式算出來的東西存進一個變數、拿去
做別的事，光靠 \texttt{print} 是不夠的：\texttt{print} 只是把文字顯示出來，不代表這個值
可以被程式的其他部分接住。要讓函式把算好的東西「送回去」給呼叫它的地方，要用
\texttt{return}。
\index{return}

\begin{code}{python}{用 return 把值送回呼叫端}
def function_name(parameter):
    return value
\end{code}

\begin{example}{python}{(add-numbers.py) 把兩個數字加起來並回傳}
def add_numbers(a, b):
    return a + b


result = add_numbers(3, 5)
print(result)
\end{example}

\iosample{}{8 \\}

這裡跟前兩節不一樣的地方是：\texttt{add\_numbers} 本身完全沒有呼叫 \texttt{print}，它
只是把 \texttt{a + b} 算好的結果透過 \texttt{return} 交還給呼叫它的那一行程式碼；真正把
結果印出來的，是外面的 \texttt{print(result)}。\texttt{return} 一執行，函式就會立刻
結束，把值帶回呼叫端，你可以把它存進變數（像這裡的 \texttt{result}），也可以直接拿去做
其他計算，函式本身不會自己決定要拿這個值做什麼。

順帶一提：如果一個函式從頭到尾都沒有寫 \texttt{return}（就像前兩節的 \texttt{say\_hello}
一樣），呼叫它其實還是會得到一個值，只是那個值是 Python 一個特別的東西，叫做
\texttt{None}，代表「什麼都沒有」。如果你把第 \ref{sec:function-param-basic} 節那個沒有 \texttt{return} 的
\texttt{say\_hello} 拿來寫成 \texttt{result = say\_hello("Alice")}，再
\texttt{print(result)}，螢幕上除了打招呼那行字，還會多印一行 \texttt{None}。這件事現在
知道就好，\texttt{None} 完整的用法留到後面章節再細講。

\section{一次送回不只一個值}
\label{sec:function-multi-return}

有時候一個函式一次要報告的資訊不只一項。舉例來說，如果你想寫一個函式，一次幫你找出一份
成績清單裡最低分跟最高分，這兩個數字都是這次呼叫的成果，你不會想為了拿到兩個值就得呼叫
兩次函式（一次找最低、一次找最高，等於把清單重複掃過兩遍）。Python 允許
\texttt{return} 後面用逗號隔開好幾個值，一次全部送回去。

\begin{code}{python}{用 return 一次送回多個值}
def function_name(parameter):
    return value1, value2
\end{code}

\begin{example}{python}{(score-range.py) 一次找出成績清單的最低分跟最高分}
def score_range(scores):
    return min(scores), max(scores)


lowest, highest = score_range([88, 72, 95, 60, 100])
print(lowest)
print(highest)
\end{example}

\iosample{}{60 \\ 100 \\}

呼叫端 \texttt{lowest, highest = score\_range(...)} 這一行，用逗號隔開兩個變數名稱，
一次接住函式送回來的兩個值：第一個對應 \texttt{lowest}，第二個對應 \texttt{highest}，
順序跟 \texttt{return} 裡寫的順序一致。

如果不特地用兩個變數去接，直接把回傳值印出來，會看到什麼？

\begin{example}{python}{(score-range-tuple.py) 不拆開接、直接印出回傳值}
def score_range(scores):
    return min(scores), max(scores)


print(score_range([88, 72, 95, 60, 100]))
\end{example}

\iosample{}{(60, 100) \\}

答案是一組用圓括號包起來的東西：\texttt{(60, 100)}。這代表
\texttt{return min(scores), max(scores)} 這一行做的事，其實是先把 \texttt{min(scores)}
跟 \texttt{max(scores)} 兩個值打包成一份小包裹，再把整份包裹送回去；呼叫端寫
\texttt{lowest, highest = ...} 的時候，Python 才把這份包裹拆開，按照順序分別塞進兩個
變數。這份「打包起來的東西」有一個正式名字叫 tuple，完整的規則留到第 \ref{ch:tuples} 章再詳細介紹，
這裡你只要記得「函式可以一次送回不只一個值，呼叫端也可以用同樣數量的變數一次接住」就好。

如果呼叫端接的變數數量跟函式實際送回來的值不一樣多，會怎麼樣？

\begin{example}{python}{(score-range-broken.py) 接的變數數量對不上}
def score_range(scores):
    return min(scores), max(scores)


lowest, highest, average = score_range([88, 72, 95, 60, 100])
\end{example}

\begin{verbatim}
Traceback (most recent call last):
  File "score-range-broken.py", line 5, in <module>
    lowest, highest, average = score_range([88, 72, 95, 60, 100])
    ^^^^^^^^^^^^^^^^^^^^^^^^^
ValueError: not enough values to unpack (expected 3, got 2)
\end{verbatim}

這裡想用三個變數接住只送回了兩個值的包裹，Python 沒辦法幫你補出第三個值，只好丟出
\texttt{ValueError}，訊息直翻是「可以拆開的值不夠：預期 3 個，實際拿到 2 個」。這跟
\texttt{TypeError} 不一樣：不是型別不對，是數量對不上。

\keypoint{\texttt{return} 後面用逗號隔開好幾個值，可以一次把多個結果送回呼叫端；呼叫端
只要用同樣數量、逗號隔開的變數名稱去接，就會依照順序各自對應。這其實是同一套「打包再拆開」
的機制，正式規則會在第 \ref{ch:tuples} 章介紹 tuple 時完整說明。}

\begin{drills}[label=drills:function-multi-return]
\item 有一個函式 \texttt{def minmax(nums): return min(nums), max(nums)}，如果呼叫
\texttt{small, big = minmax([4, 9, 1, 7])}，\texttt{small} 跟 \texttt{big} 各會是
多少？
\item 承上題，如果改成只用一個變數接：\texttt{result = minmax([4, 9, 1, 7])}，
\texttt{print(result)} 會印出什麼？
\item 如果 \texttt{minmax} 函式呼叫時只給一個數字的清單，例如 \texttt{minmax([5])}，
拆開接的兩個變數 \texttt{small}、\texttt{big} 會不會相等？
\end{drills}

\section{一次餵進不只一項資訊：positional 引數}
\label{sec:function-positional}

前面的函式都只需要一項資訊，但很多時候一個函式需要好幾項資訊才能做完一件事，比如描述
一隻寵物，至少要知道牠是什麼動物、叫什麼名字。這種時候，只要在 \texttt{def} 的括號裡
用逗號隔開多個參數即可；呼叫的時候依照\textbf{同樣的順序}把引數一個個放進去，Python 會
按照位置把引數對應到參數，這種方式叫做 positional（位置）引數。
\index{positional 引數}

\begin{code}{python}{定義並呼叫一個帶多個 positional 參數的函式}
def function_name(param1, param2):
    code

function_name(arg1, arg2)
\end{code}

\begin{example}{python}{(describe-pet.py) 描述一隻寵物}
def describe_pet(animal_type, pet_name):
    print(f"I have a pet {animal_type} named {pet_name}.")


describe_pet("dog", "Rex")
\end{example}

\iosample{}{I have a pet dog named Rex. \\}

因為是依照順序對應，第一個引數 \texttt{"dog"} 對到第一個參數 \texttt{animal\_type}，
第二個引數 \texttt{"Rex"} 對到第二個參數 \texttt{pet\_name}。這也代表順序如果寫反了，
Python 不會幫你檢查「這個字看起來不像動物名稱」，它只管照順序對應，其他的事一概不管。

\begin{example}{python}{(describe-pet-swapped.py) 引數順序寫反了}
def describe_pet(animal_type, pet_name):
    print(f"I have a pet {animal_type} named {pet_name}.")


# 引數順序寫反了：第一個位置現在對應到 animal_type
describe_pet("Rex", "dog")
\end{example}

\iosample{}{I have a pet Rex named dog. \\}

程式沒有跳出任何紅字錯誤，正常執行到底，但印出來的內容完全不合理：「I have a pet Rex」。
這種「沒有報錯，但結果荒謬」的情形，比丟出例外訊息的錯誤更容易被忽略，如果你不特別留意
輸出內容，很可能以為程式一切正常，實際上已經錯了。

\keypoint{positional 引數的對應完全靠順序，Python 不會替你檢查「這個引數放在這個位置合
不合理」。引數一多，順序寫反是很容易發生、卻又不會報錯的錯誤，讀完下一節你會看到另一種
可以避開這個問題的呼叫方式。}

\section{用名字指定：keyword 引數}
\label{sec:function-keyword}

如果不想冒著順序寫反的風險，或是引數一多、光靠位置很難一眼看出誰對應誰，Python 讓你可以
改用參數的名字來指定引數，順序就不再重要了，這種方式叫做 keyword（關鍵字）引數，寫法是
在呼叫時用 \texttt{參數名稱=值} 的形式。
\index{keyword 引數}

\begin{code}{python}{用 keyword 引數呼叫函式}
function_name(param1=value1, param2=value2)
\end{code}

\begin{example}{python}{(describe-pet-keyword.py) 用參數名稱呼叫}
def describe_pet(animal_type, pet_name):
    print(f"I have a pet {animal_type} named {pet_name}.")


describe_pet(pet_name="Rex", animal_type="dog")
\end{example}

\iosample{}{I have a pet dog named Rex. \\}

雖然呼叫的時候 \texttt{pet\_name} 寫在 \texttt{animal\_type} 前面，跟函式定義的參數
順序不一樣，結果依然正確，因為 Python 是照名字對應，不是照位置。

\stylepoint{引數一多、或是好幾個引數型別相同（例如都是字串）容易搞混順序的時候，習慣用
keyword 引數會讓呼叫端更好讀、也更不容易在順序上出錯，這是比較 Python 慣用的寫法，尤其
是在呼叫別人寫的函式、你記不清楚參數順序的時候特別好用。}

\section{省略引數：default 參數值}
\label{sec:function-params}

如果一個函式大部分時候都會用同一個值呼叫，只有少數情況才需要換成別的，每次呼叫都把這個
值重複打一遍很麻煩。這種時候可以在定義函式的時候，替參數設一個預設值，呼叫時如果省略這個
引數，就自動套用預設值；想要不同的值，照樣可以在呼叫時自己指定。
\index{default 參數值}

\begin{code}{python}{定義帶預設值的參數}
def function_name(param1, param2=default_value):
    code
\end{code}

\begin{example}{python}{(describe-pet-default.py) 大部分的寵物都是狗}
def describe_pet(pet_name, animal_type="dog"):
    print(f"I have a pet {animal_type} named {pet_name}.")


describe_pet("Rex")
describe_pet("Whiskers", "cat")
\end{example}

\iosample{}{I have a pet dog named Rex. \\ I have a pet cat named Whiskers. \\}

第一次呼叫只給了 \texttt{pet\_name}，\texttt{animal\_type} 就自動套用了預設值
\texttt{"dog"}；第二次呼叫額外給了第二個引數 \texttt{"cat"}，蓋掉了預設值。

留意這個範例裡參數的順序：有預設值的 \texttt{animal\_type} 被放在沒有預設值的
\texttt{pet\_name} 後面，這不是隨便排的，是 Python 語法的硬性規定。

\begin{example}{python}{(describe-pet-default-broken.py) 預設值參數放錯位置}
def describe_pet(animal_type="dog", pet_name):
    print(f"I have a pet {animal_type} named {pet_name}.")
\end{example}

\begin{verbatim}
  File "describe-pet-default-broken.py", line 1
    def describe_pet(animal_type="dog", pet_name):
                                        ^^^^^^^^
SyntaxError: parameter without a default follows parameter with a default
\end{verbatim}

這段錯誤甚至在函式被呼叫之前、程式一讀到這行定義就發作了：\texttt{SyntaxError} 代表
Python 光是「讀」這行程式碼的規則就過不了關，訊息直翻是「一個沒有預設值的參數，出現在
一個有預設值的參數後面」。

\keypoint{有預設值的參數，一定要放在沒有預設值的參數後面。原因不難想：如果
\texttt{describe\_pet(animal\_type="dog", pet\_name)} 這樣的定義是合法的，那呼叫
\texttt{describe\_pet("Rex")} 時，Python 沒辦法判斷 \texttt{"Rex"} 這個唯一的引數，
到底該對應到有預設值的 \texttt{animal\_type}，還是該去填沒有預設值、非給不可的
\texttt{pet\_name}。為了避免這種歧義，Python 乾脆直接禁止這樣寫。}

\begin{drills}[label=drills:function-params]
\item 承 \texttt{describe-pet-default.py}：如果呼叫 \texttt{describe\_pet(animal\_type="bird",
pet\_name="Tweety")}，用的是 keyword 引數，兩個都自己指定了值，輸出會是什麼？
\item 如果呼叫 \texttt{describe\_pet("Max", "cat")}，用的是 positional 引數，
\texttt{animal\_type} 會是 \texttt{"cat"} 還是保留預設值 \texttt{"dog"}？
\item 有一個函式定義成
\texttt{def make\_coffee(size, sugar=True, milk=False):}，如果只想指定
\texttt{milk=True}、\texttt{size} 跟 \texttt{sugar} 都用預設值/正常給，你會怎麼呼叫這個
函式？（提示：\texttt{size} 沒有預設值，一定要給。）
\end{drills}

\section{收下不確定個數的引數：\texttt{*args}}
\label{sec:function-args}

到目前為止，你寫的每一個函式，都有一件事是事先講好的：它會收到幾項資訊。
\texttt{describe\_pet} 就是要兩項，多給少給都會出事。但你其實一直在用一個不是這樣的
函式——\texttt{print}。你可以寫 \texttt{print("hello")}，也可以寫
\texttt{print("hello", "world", 123)}，它照單全收，從來不抱怨。這種函式是怎麼寫出來的？

答案是在參數名稱前面加一個星號。加了星號的參數，會把呼叫時給的所有引數通通收成一包，
不管給了幾個。

\begin{code}{python}{收下不確定個數的引數}
def function_name(*parameter_name):
    code that uses parameter_name
\end{code}

\begin{example}{python}{(total-price.py) 把任意個數的價格加總起來}
def total_price(*prices):
    total = 0
    for price in prices:
        total = total + price
    return total


print(total_price(100))
print(total_price(100, 250, 80))
print(total_price())
\end{example}

\iosample{}{100 \\ 430 \\ 0 \\}

同一個 \texttt{total\_price}，第一次呼叫給一個價格、第二次給三個、第三次一個都不給，
三次都正常執行。函式主體裡的 \texttt{prices}，就是那一包被收起來的引數，可以直接放進
\texttt{for} 迴圈一個一個處理。第三次呼叫沒給任何引數，那一包是空的，迴圈一次都不會跑，
\texttt{total} 維持一開始的 \texttt{0}——「零個」也算合法，不會報錯。

星號收進來的那一包，型別叫做 tuple。tuple 是下一章的主題，這裡你只要先知道兩件事：它可以
放進 \texttt{for} 迴圈一個一個走過去，也可以用 \texttt{len()} 數出裡面有幾個東西。剩下的
細節留到第 \ref{ch:tuples} 章。

\keypoint{星號寫在\textbf{參數名稱前面}，意思是「把呼叫時給的引數通通收成一包」。
\texttt{args} 只是這個參數最常見的慣例名字，不是關鍵字——叫 \texttt{prices}、
\texttt{items}、\texttt{numbers} 都完全可以，取一個講得出它裝什麼的名字反而更好。}

這個寫法有一個新手幾乎都會踩一次的誤解：既然它可以收下一堆東西，那我把一個現成的 list
整包丟進去，應該也一樣吧？

\begin{example}{python}{(show-all-wrong.py) 把一個 list 整包傳進去}
def show_all(*items):
    print(items)
    print(len(items))


show_all("apple", "banana")
show_all(["apple", "banana"])
\end{example}

\iosample{}{('apple', 'banana') \\ 2 \\ (['apple', 'banana'],) \\ 1 \\}

前兩行是預期中的結果：兩個引數被收成一包，裡面有兩個東西。但後兩行不太對勁——
\texttt{(['apple', 'banana'],)} 是「一包裡面裝著一個 list」，\texttt{len()} 數出來是
\texttt{1} 不是 \texttt{2}。原因是第二次呼叫只給了\textbf{一個}引數，那個引數剛好是一個
list，星號老老實實地把「這一個東西」收成一包，它沒辦法知道你其實是想把 list 裡面的兩個
元素分開來給它。

那如果真的手上就是一個 list，想把裡面的東西拆開來餵進去呢？這件事有專門的寫法，就在
下一節。

\begin{drills}[label=drills:function-args]
\item 承 \texttt{total-price.py}：呼叫 \texttt{total\_price(10, 20)} 會回傳什麼？
\item 如果把 \texttt{total\_price} 的定義改成 \texttt{def total\_price(prices):}
（拿掉星號），再呼叫 \texttt{total\_price(100, 250, 80)}，會發生什麼事？
（提示：回頭看第 \ref{sec:function-param-basic} 節那個引數數量對不上的錯誤訊息。）
\end{drills}

\section{收下不確定名字的引數：\texttt{**kwargs}}
\label{sec:function-kwargs}

一個星號收下的是「一堆沒有名字的值」。但你在第 \ref{sec:function-keyword} 節學過，引數也可以帶名字
（\texttt{animal\_type="dog"} 這種寫法）。如果連「會收到哪些名字」都事先不知道呢？這時候
用兩個星號。

\begin{code}{python}{收下不確定名字的引數}
def function_name(**parameter_name):
    code that uses parameter_name
\end{code}

\begin{example}{python}{(build-profile.py) 收下任意個欄位，組成一份資料}
def build_profile(**fields):
    profile = {}
    for key, value in fields.items():
        profile[key] = value
    return profile


print(build_profile(name="Amy", city="Taipei"))
print(build_profile(name="Ben", city="Tainan", age=30))
\end{example}

\iosample{}{\{'name': 'Amy', 'city': 'Taipei'\} \\ \{'name': 'Ben', 'city': 'Tainan', 'age': 30\} \\}

兩個星號收進來的那一包，是一個 dict——這個你在第 \ref{ch:dict-set} 章已經完整學過了，
引數的名字變成 key，引數的值變成 value，所以主體裡可以直接用 \texttt{.items()} 走訪它。
第一次呼叫給了兩個欄位、第二次給三個，函式完全不用改。

\keypoint{一個函式如果同時要用一般參數、\texttt{*} 參數、\texttt{**} 參數，
兩個星號的那一個一定放在最後面，因為它是「剩下所有帶名字的引數」的收容所，後面不可能還有
別的東西。}

\subsection*{反過來用：把一包拆開餵進去}

真正容易混淆的來了。同樣是這兩個符號，寫在\textbf{呼叫端}的時候，做的事情剛好相反——不是
把散的收成一包，而是把一包拆成散的。上一節結尾那個「手上就是一個 list」的問題，答案就在
這裡。

\begin{example}{python}{(unpack-call.py) 把 list 與 dict 拆開餵進參數}
def describe_pet(animal_type, pet_name):
    print(f"I have a {animal_type} named {pet_name}.")


pet_list = ["cat", "Mochi"]
pet_dict = {"animal_type": "dog", "pet_name": "Lucky"}

describe_pet(pet_list[0], pet_list[1])
describe_pet(*pet_list)
describe_pet(**pet_dict)
\end{example}

\iosample{}{I have a cat named Mochi. \\ I have a cat named Mochi. \\ I have a dog named Lucky. \\}

前兩行輸出一模一樣，這正是重點：\texttt{describe\_pet(*pet\_list)} 跟自己動手寫
\texttt{describe\_pet(pet\_list[0], pet\_list[1])} 是同一件事，只是不用一個一個把索引
打出來。一個星號放在呼叫端的 list 前面，意思是「把這個 list 攤開，一個元素對應一個
位置引數」。

第三行則是 dict 版本：\texttt{describe\_pet(**pet\_dict)} 等於
\texttt{describe\_pet(animal\_type="dog", pet\_name="Lucky")}，dict 的每一組 key/value
被攤成一個 keyword 引數。這也是為什麼 \texttt{pet\_dict} 的 key 必須剛好叫
\texttt{animal\_type} 跟 \texttt{pet\_name}——它們要對得上函式的參數名字。

\keypoint{同一個星號，位置不同意思剛好相反：寫在 \texttt{def} 的\textbf{參數}前面是
「把散的收成一包」，寫在呼叫時的\textbf{引數}前面是「把一包拆成散的」。判斷的方法很單純，
看它出現在哪一邊——定義那一行，還是呼叫那一行。}

\begin{drills}[label=drills:function-varargs]
\item 有一個函式 \texttt{def show(*args): print(len(args))}。分別呼叫
\texttt{show(1, 2, 3)} 與 \texttt{show([1, 2, 3])}，各會印出什麼？
\item 承上題，如果想讓 \texttt{show([1, 2, 3])} 也印出 \texttt{3}，呼叫時要怎麼改？
\item 有一個函式 \texttt{def greet(greeting, name): print(greeting, name)}，
以及一個 dict \texttt{info = \{"greeting": "Hi", "name": "Amy"\}}。
要怎麼用一行呼叫，把 \texttt{info} 裡的兩個值餵給 \texttt{greet}？
\end{drills}

\section{為什麼要拆函式？}
\label{sec:function-why}

到這裡你已經知道函式的所有基本零件：怎麼定義、怎麼傳引數、怎麼送回值。但你可能會問：
明明不用函式，一路把程式碼從頭寫到尾也能跑，為什麼要多繞這一圈？看一段沒有函式、但重複
做了三次幾乎一樣的事的程式碼，會比較容易感受到差別。

\begin{example}{python}{(discount-before.py) 沒有函式，三段幾乎一樣的程式碼}
price1 = 100
discount1 = 0.9
print(f"Discounted price: {price1 * discount1}")

price2 = 250
discount2 = 0.8
print(f"Discounted price: {price2 * discount2}")

price3 = 60
discount3 = 0.95
print(f"Discounted price: {price3 * discount3}")
\end{example}

\iosample{}{Discounted price: 90.0 \\ Discounted price: 200.0 \\ Discounted price: 57.0 \\}

這段程式能動，但同一句「算出折扣後價格、印出來」的邏輯被抄了三遍。如果之後想改顯示格式
（比如要四捨五入到整數），得記得同時改三個地方，只要漏改一處，三段的行為就會不一致，而且
這種不一致往往要等到有人回報「奇怪，這裡的價格看起來跟別的地方格式不一樣」才會被發現。
改寫成函式之後，同一段邏輯只需要寫一次：

\begin{example}{python}{(discount-after.py) 用函式收起重複的邏輯}
def show_discounted_price(price, discount):
    print(f"Discounted price: {price * discount}")


show_discounted_price(100, 0.9)
show_discounted_price(250, 0.8)
show_discounted_price(60, 0.95)
\end{example}

\iosample{}{Discounted price: 90.0 \\ Discounted price: 200.0 \\ Discounted price: 57.0 \\}

輸出跟前一個版本一模一樣，但這次「怎麼算、怎麼顯示」的規則只出現在
\texttt{show\_discounted\_price} 定義的那三行裡，之後三個呼叫端都是共用同一份規則。
之後如果要改顯示格式，只要改函式定義那裡一次，三個呼叫端全部自動套用新規則，不會有改一個
忘一個的風險。

除了避免重複，拆函式還有一個更直接的好處：讓程式讀起來更像一份有意義的步驟清單。呼叫端
寫著 \texttt{show\_discounted\_price(100, 0.9)}，光看名字就知道這行在做什麼，不用一路
讀進函式主體的細節才搞懂意圖；細節被摺進函式裡面，只有真的需要修改邏輯的時候才需要打開來
看。

\stylepoint{看到同一段程式碼（或幾乎一樣、只有幾個值不同的程式碼）出現兩次以上，通常就是
該考慮拆成函式的訊號：把不一樣的部分變成參數，一樣的部分變成函式主體。}

\section{函式也是一種值：把函式傳來傳去}
\label{sec:function-first-class}

還記得第 \ref{ch:variables} 章的名字牌比喻嗎？\texttt{x = 5} 這一行做的事，是造一個值是 5 的物件，再把
名字牌 \texttt{x} 貼上去，並不是把 5 複製一份塞進一個叫 \texttt{x} 的盒子。這件事其實對
函式也成立：用 \texttt{def} 定義出來的函式本身也是一個物件，\texttt{def} 那一行裡寫的
函式名稱，就是貼在這個函式物件上的第一張名字牌。既然函式也是物件，跟數字、字串、list 一樣，
它就可以被指定給另一個變數、當作引數傳給另一個函式。

\begin{example}{python}{(function-as-value.py) 把函式指定給另一個變數}
def square(n):
    return n * n


operation = square
print(operation(5))
\end{example}

\iosample{}{25 \\}

留意這裡的 \texttt{operation = square} 沒有加括號。加了括號的 \texttt{square()} 是
「呼叫這個函式，把結果拿回來」；沒有括號的 \texttt{square} 單獨寫，指的是函式物件本身，
也就是那個貼著 \texttt{square} 這張名字牌的東西。\texttt{operation = square} 做的事，
是把另一張名字牌 \texttt{operation} 也貼到同一個函式物件上，之後 \texttt{operation(5)}
跟 \texttt{square(5)} 呼叫的其實是同一個函式，只是用了不同的名字牌喊它。

如果不小心把括號加上去，會發生什麼事？

\begin{example}{python}{(function-as-value-broken.py) 忘記拿掉括號}
def square(n):
    return n * n


operation = square()
\end{example}

\begin{verbatim}
Traceback (most recent call last):
  File "function-as-value-broken.py", line 5, in <module>
    operation = square()
                ^^^^^^^^
TypeError: square() missing 1 required positional argument: 'n'
\end{verbatim}

\texttt{square()} 這樣寫是「立刻呼叫 square，把結果指定給 operation」，但
\texttt{square} 需要一個引數 \texttt{n}，這裡卻一個都沒給，於是跟 \ref{sec:function-param-basic} 節看過的錯誤一樣，
丟出 \texttt{TypeError}。這是分辨「函式物件本身」跟「呼叫函式、拿到的結果」最直接的方式：
有沒有加括號。

函式物件不只可以存進變數，還可以直接當作引數，傳給另一個函式使用：

\begin{example}{python}{(apply-twice.py) 把函式當作引數傳給另一個函式}
def apply_twice(func, value):
    return func(func(value))


def add_ten(n):
    return n + 10


print(apply_twice(add_ten, 5))
\end{example}

\iosample{}{25 \\}

\texttt{apply\_twice} 的第一個參數 \texttt{func}，接住的是 \texttt{add\_ten} 這個函式
物件本身，不是呼叫它拿到的結果。\texttt{apply\_twice} 主體裡的
\texttt{func(func(value))} 先用 \texttt{value}（也就是 5）呼叫一次 \texttt{func}，
算出 15，再拿 15 當引數呼叫一次 \texttt{func}，算出 25。整個過程 \texttt{apply\_twice}
完全不需要知道傳進來的函式實際上在做加法還是別的運算，它只負責「呼叫兩次、把上一次的結果
餵給下一次」，真正的邏輯是由呼叫端決定要傳哪一個函式進去。

\keypoint{函式本身也是一個物件，函式名稱只是貼在這個物件上的一張名字牌。把函式指定給另一個
變數、或當作引數傳給另一個函式，用的是同一套名字牌邏輯，只是這次貼名字牌的對象是一個函式
物件，不是數字或字串。加不加括號差很多：不加括號指的是函式物件本身，加了括號才是「呼叫」
它。}

\begin{drills}[label=drills:function-first-class]
\item 有 \texttt{def double(n): return n * 2}，如果寫
\texttt{f = double; print(f(6))}，會印出什麼？
\item 承 \texttt{apply-twice.py}：如果改成呼叫
\texttt{apply\_twice(square, 3)}（\texttt{square} 是本節前面定義的平方函式），
會印出什麼？（提示：先算 \texttt{square(3)}，再把結果拿去算一次 \texttt{square}。）
\item 如果把 \texttt{apply\_twice(add\_ten, 5)} 不小心寫成
\texttt{apply\_twice(add\_ten(), 5)}，會發生什麼事？
\end{drills}

\section{呼叫函式時，你的東西是被「借」走的}
\label{sec:function-nametag}

上一節看到函式物件本身可以被傳來傳去，這一節要看更常見的情況：呼叫函式時，傳進去的引數
發生了什麼事。還記得第 \ref{ch:variables} 章講過的名字牌比喻嗎？\texttt{x = 5} 是造一個值是 5 的物件，把
名字牌 \texttt{x} 貼上去；\texttt{y = x} 則是把另一張名字牌也貼到同一個物件上，兩張名字
牌指的是同一個東西。

第 \ref{ch:lists} 章你也看過，這件事在 list 上特別明顯：改動一個名字看到的內容，另一個指著同一個 list
的名字也會跟著變。

呼叫函式、傳引數的時候，發生的其實是同一件事：Python 是把引數對應的那個物件「借」給函式
裡的參數這張名字牌，不是複製一份新的給它。

\begin{example}{python}{(cart-append.py) 函式內對 list 動手腳，函式外看得到}
def add_apple(cart):
    cart.append("apple")


my_cart = ["banana"]
add_apple(my_cart)
print(my_cart)
\end{example}

\iosample{}{['banana', 'apple'] \\}

呼叫 \texttt{add\_apple(my\_cart)} 的時候，\texttt{my\_cart} 指著的那個 list 物件，被
「借」給了函式裡的參數 \texttt{cart}：\texttt{cart} 跟外面的 \texttt{my\_cart} 現在是
兩張名字牌，指著同一個物件。函式主體裡 \texttt{cart.append("apple")} 是對這個物件本身
「動手腳」（原地修改），所以不管你是透過 \texttt{cart} 這張名字牌，還是外面的
\texttt{my\_cart} 這張名字牌去看，看到的都是同一份改過的內容。

但如果函式做的事情，是把參數整個重新賦值呢？

\begin{example}{python}{(cart-reassign.py) 對參數重新賦值，函式外不受影響}
def replace_cart(cart):
    # 這行把 cart 這張名字牌重新貼到一個全新的 list 上，
    # 跟外面的 my_cart 已經沒有關係了
    cart = ["new list"]


my_cart = ["banana"]
replace_cart(my_cart)
print(my_cart)
\end{example}

\iosample{}{['banana'] \\}

這次 \texttt{my\_cart} 完全沒變。\texttt{cart = ["new list"]} 做的事，是把
\texttt{cart} 這張名字牌從原本借來的物件上撕下來，重新貼到一個全新造出來的 list 上：這個
動作只影響 \texttt{cart} 這張名字牌貼在哪裡，跟外面的 \texttt{my\_cart} 完全無關，
\texttt{my\_cart} 從頭到尾都貼在它原本的那個物件上，沒有被撕過。

下面這張圖把 \texttt{add\_apple(my\_cart)} 呼叫當下的狀態畫出來：\texttt{my\_cart} 跟
函式裡的參數 \texttt{cart} 是兩張分開的名字牌，但貼的是同一個 list 物件，所以
\texttt{cart.append("apple")} 改的東西，兩張名字牌都看得到。

\begin{center}
\begin{tikzpicture}
\node[process] (myvar) at (0,2) {my\_cart（呼叫端）};
\node[process] (param) at (4.6,2) {cart（函式參數）};
\node[startend, minimum width=5.2cm] (obj) at (2.3,0)
    {list 物件：\texttt{['banana', 'apple']}};
\draw[diagarrow] (myvar.south) -- (obj.north -| myvar);
\draw[diagarrow] (param.south) -- (obj.north -| param);
\end{tikzpicture}
\end{center}

\keypoint{函式呼叫時，引數是把物件「借」給參數這張名字牌，不是複製一份給它。函式內對參數
「動手腳」（例如 \texttt{.append()}、\texttt{.sort()} 這類原地修改的操作）會影響到函式
外面看到的同一個物件；對參數「重新貼名字牌」（重新賦值，例如 \texttt{cart = ...}）則只是
讓參數這張名字牌換個地方貼，不會影響函式外面原本的名字牌指向哪裡。}

\section{可變預設參數陷阱}
\label{sec:mutable-default}

上一節的名字牌模型，可以直接解釋一個所有 Python 使用者遲早都會踩到的經典陷阱：如果一個
函式的預設參數值是 list 這種可變（能被原地修改）的物件，會發生非常違反直覺的事。

先看一個「有問題」的版本：一個函式接收一個商品名稱，把它加進購物車 list，如果呼叫時沒有
自己帶購物車，就用預設值。

\begin{code}{python}{危險寫法：可變物件當預設參數值}
def function_name(item, cart=[]):
    cart.append(item)
    return cart
\end{code}

\begin{example}{python}{(shopping-cart-bug.py) 連續呼叫兩次，購物車卻黏在一起}
def add_item(item, cart=[]):      # 危險：cart 只在 def 執行時被建立一次
    cart.append(item)
    return cart


print(add_item("apple"))
print(add_item("banana"))
\end{example}

\iosample{}{['apple'] \\{} ['apple', 'banana'] \\}

這裡兩次呼叫都沒有自己帶 \texttt{cart}，理論上你可能會預期第二次呼叫應該也是拿到一個
「全新的空購物車」，只裝著 \texttt{"banana"}，但實際看到的是 \texttt{['apple', 'banana']}，
第一次呼叫加進去的 \texttt{"apple"} 竟然還留著。

原因跟 \ref{sec:function-nametag} 節的名字牌模型完全一樣：\texttt{cart=[]} 這個預設值，只在 \texttt{def} 那
一行被\textbf{執行的當下}造了一次 list 物件、把預設用的名字牌貼上去，這件事只發生一次，
不是每次呼叫都重新造一個新的空 list。之後每一次呼叫，只要你沒有自己傳 \texttt{cart} 進
去，用的都是同一個物件；而函式主體裡的 \texttt{cart.append(item)} 又是原地修改，於是每次
呼叫都在同一個 list 上疊加內容，跟「每次呼叫都應該拿到一個乾淨的新盒子」的直覺完全相反。

修法是把預設值換成 \texttt{None}，\texttt{None} 是不可變的，不會被意外累積內容，再讓
函式自己在主體裡判斷：如果沒收到 \texttt{cart}，才自己造一個新的空 list。

\begin{code}{python}{安全寫法：預設值用 None，函式內自己造新物件}
def function_name(item, cart=None):
    if cart is None:
        cart = []
    cart.append(item)
    return cart
\end{code}

這裡的 \texttt{cart is None} 用 \texttt{is} 而不是 \texttt{==}，是因為 \texttt{None}
在 Python 裡全世界只有一個，\texttt{is} 問的是「是不是同一個東西」，這是檢查一個值是不是
\texttt{None} 時，Python 慣用的寫法。
\index{is}

\begin{example}{python}{(shopping-cart-fixed.py) 每次呼叫都拿到乾淨的購物車}
def add_item(item, cart=None):
    if cart is None:        # 每次呼叫都會檢查，並自己建立一個新的空 list
        cart = []
    cart.append(item)
    return cart


print(add_item("apple"))
print(add_item("banana"))
\end{example}

\iosample{}{['apple'] \\{} ['banana'] \\}

這次兩次呼叫互不影響，因為 \texttt{if cart is None} 這一行，每次沒收到 \texttt{cart}
時都會重新執行 \texttt{cart = []}，也就是每次都貼一張新的名字牌到一個全新造出來的空 list
上，不會再共用同一個物件。

下面這張圖畫出「有問題」的版本（\texttt{cart=[]}）裡發生的事：\texttt{def} 那一行只在
程式讀到它的當下執行一次，造出一個預設用的 list 物件；之後不管呼叫幾次、只要沒自己帶
\texttt{cart}，用的都是同一個物件。

\begin{center}
\begin{tikzpicture}
\node[process] (def) at (2.3,3.4) {\texttt{def} 執行時（只發生一次）};
\node[startend, minimum width=5.6cm] (obj) at (2.3,1.9)
    {預設 list 物件（一開始是 \texttt{[]}）};
\node[process] (call1) at (0,0) {第一次呼叫\\\texttt{add\_item("apple")}};
\node[process] (call2) at (4.6,0) {第二次呼叫\\\texttt{add\_item("banana")}};
\draw[diagarrow] (def.south) -- (obj.north);
\draw[diagarrow] (call1.north) -- (obj.south -| call1);
\draw[diagarrow] (call2.north) -- (obj.south -| call2);
\end{tikzpicture}
\end{center}

\keypoint{函式的預設參數值只在 \texttt{def} 執行的那一刻被建立一次，不是每次呼叫都重新
建立。如果預設值是 list 這類可變物件，又在函式內對它做原地修改，不同次呼叫之間就會共用、
累積同一個物件的內容，這是 Python 一個非常有名的坑。安全的寫法是把預設值換成
\texttt{None}，再讓函式自己在主體裡判斷、自己建立新的物件。}

\begin{drills}[label=drills:mutable-default]
\item 承 \texttt{shopping-cart-bug.py}：如果再加一行
\texttt{print(add\_item("orange"))}（第三次呼叫，一樣沒帶 \texttt{cart}），這一行會印出
什麼？
\item 承 \texttt{shopping-cart-bug.py}：如果第二次呼叫改成自己帶一個購物車，寫成
\texttt{print(add\_item("banana", ["watermelon"]))}，這一行會印出什麼？（提示：這次有自己傳
\texttt{cart} 進去，用的還是同一個共用物件嗎？）
\item 承 \texttt{shopping-cart-fixed.py}：如果連續呼叫三次
\texttt{add\_item("apple")}、\texttt{add\_item("banana")}、\texttt{add\_item("orange")}，
分別 \texttt{print} 每一次的回傳值，三行輸出各會是什麼？
\end{drills}

\section{函式裡的世界跟函式外的世界：變數作用域}
\label{sec:function-scope}

到目前為止，你可能已經注意到一件事：函式裡面定義的變數，好像沒辦法在函式外面直接用。這不
是巧合，是 Python 刻意設計的規則：一個變數「在哪裡看得到、在哪裡看不到」，這個範圍叫做
作用域（scope）。
\index{作用域}

在函式主體裡定義的變數，叫做區域變數（local variable），「區域」的意思是，它只在這個
函式執行的期間存在，函式一旦執行完畢、把值 \texttt{return} 出去，這個變數就消失了，外面
沒有辦法再存取它。
\index{區域變數}

\begin{example}{python}{(order-total-scope.py) 函式外看不到函式內的變數}
def calculate_total(price, quantity):
    total = price * quantity
    return total


result = calculate_total(50, 3)
print(result)
print(total)
\end{example}

\begin{verbatim}
150
Traceback (most recent call last):
  File "order-total-scope.py", line 8, in <module>
    print(total)
          ^^^^^
NameError: name 'total' is not defined
\end{verbatim}

\texttt{print(result)} 正常印出了 \texttt{150}，\texttt{result} 是在函式外面定義的
變數，接住了 \texttt{calculate\_total} 回傳的值。但緊接著的 \texttt{print(total)} 卻
報錯了：\texttt{total} 是在 \texttt{calculate\_total} 主體裡定義的區域變數，函式一
\texttt{return} 完就消失了，函式外面根本沒有一個叫 \texttt{total} 的名字存在過，Python
當然找不到。這是這一章第二次看到 \texttt{NameError}，跟第 \ref{ch:introduction} 章看到的一樣，都是「這個
名字，Python 找不到」的意思，差別只在於，這次不是打錯字，是你想用一個原本就不存在於這個
範圍的變數。

那函式外面的變數呢？函式外定義的變數叫做全域變數（global variable），跟區域變數的規則
不一樣：函式\textbf{可以}讀取函式外面已經存在的全域變數，只是不能直接（不用特殊語法）在
函式內修改它。
\index{全域變數}

\begin{example}{python}{(order-total-tax.py) 函式讀取函式外的變數}
tax_rate = 0.05


def calculate_total(price):
    return price * (1 + tax_rate)


print(calculate_total(100))
\end{example}

\iosample{}{105.0 \\}

\texttt{tax\_rate} 是在函式外面定義的全域變數，\texttt{calculate\_total} 主體裡直接用
了它，沒有另外當作參數傳進去，程式依然正常執行，這說明「函式外面看不到函式裡面」是單向的
規則，不是完全隔絕：函式裡面看得到函式外面已經存在的東西，只是反過來不行。

到這裡為止都還算符合直覺：「函式可以偷看外面，但不能直接動外面的東西」。但這條規則有一個
非常容易誤解的地方，很值得單獨拿出來看：如果函式主體裡\textbf{有任何一行}打算對一個跟
全域變數同名的名字賦值，Python 會把這個名字在整段函式主體裡都當成區域變數，就算那一行
賦值寫在後面，前面想「先讀一下」都不行。

\begin{example}{python}{(counter-broken.py) 想在函式裡直接更新全域變數}
counter = 0


def increment():
    counter = counter + 1
    print(counter)


increment()
\end{example}

\begin{verbatim}
Traceback (most recent call last):
  File "counter-broken.py", line 9, in <module>
    increment()
    ~~~~~~~~~^^
  File "counter-broken.py", line 5, in increment
    counter = counter + 1
              ^^^^^^^
UnboundLocalError: cannot access local variable 'counter' where it is not associated with a value
\end{verbatim}

這個錯誤第一次看到通常會很困惑：明明外面就有一個 \texttt{counter}，為什麼函式裡讀不到？
原因是 \texttt{counter = counter + 1} 這一行，等號左邊出現了 \texttt{counter}，Python
一看到函式主體裡有任何一行在對某個名字賦值，就會把這個名字在整段主體裡都當成區域變數，
跟外面那個同名的全域變數是兩個完全不同的東西。

問題是，右邊的 \texttt{counter + 1} 要先讀取 \texttt{counter} 目前的值才能算，可是這個
「區域版本」的 \texttt{counter} 在這一行之前根本還沒被賦值過，於是 Python 丟出
\texttt{UnboundLocalError}：「這個區域變數還沒有對應到任何值，不能讀取」。
\index{UnboundLocalError}

如果真的想在函式裡修改一個函式外面的全域變數，要用 \texttt{global} 這個關鍵字，明確告訴
Python「這個名字不要當成新的區域變數，去對應外面那個全域變數」。
\index{global}

\begin{code}{python}{用 global 明確表示要修改函式外的變數}
variable = initial_value


def function_name():
    global variable
    variable = variable + 1
\end{code}

\begin{example}{python}{(counter-fixed.py) 用 global 修改函式外的變數}
counter = 0


def increment():
    global counter
    counter = counter + 1
    print(counter)


increment()
increment()
\end{example}

\iosample{}{1 \\ 2 \\}

有了 \texttt{global counter} 這一行，函式裡的 \texttt{counter} 不再被當成一個新的
區域變數，而是明確對應到外面那個全域變數。兩次呼叫都直接修改同一個物件，所以第二次呼叫
印出的是累加過的 \texttt{2}，不是每次都從頭開始的 \texttt{1}。

\texttt{global} 雖然能解決問題，但多數 Python 教學（包括這本書）都不建議把它當成常用
工具：一個函式如果偷偷改了外面某個變數，呼叫端光看函式怎麼被呼叫，很難猜到它還動了「看
不見」的東西。更常見、也更清楚的做法，還是把需要的值當作參數傳進去、用 \texttt{return}
把結果送出來，函式的行為只要看它接收什麼、送回什麼就懂，不用去猜它背後偷偷改了哪個全域
變數。

\stylepoint{看到函式需要用到 \texttt{global} 才能完成工作，通常是個訊號：這段邏輯改寫成
「接收參數、回傳結果」的形式，會比讓函式直接動全域變數更容易讀懂、也更不容易出錯。}

順帶一提：不同函式各自的區域變數是彼此獨立的，就算兩個函式都用了同一個變數名字（例如都
叫 \texttt{total}），也不會互相干擾，每次呼叫函式，Python 都會替它準備一份全新、只屬於
這次呼叫的區域變數空間。這也回應了 \ref{sec:function-nametag} 節的名字牌模型：區域變數消失的時候，它身上貼的
名字牌被撕掉了，如果那個名字牌是唯一指向某個物件的名字牌，那個物件也就跟著找不到主人了。

函式主體裡除了可以呼叫別的函式，其實也可以呼叫函式自己，這種寫法叫做遞迴。它能成立，
靠的正是這一節講清楚的規則——每一次呼叫都有自己獨立的一份區域變數，所以同一個函式一層一層
疊著呼叫下去，彼此的變數也不會打架。遞迴完整的討論（怎麼讓它停下來、什麼時候值得用它、
什麼時候用迴圈更好）留到第 \ref{ch:recursion} 章。

\begin{drills}[label=drills:scope]
\item 承 \texttt{order-total-scope.py}：如果把最後一行 \texttt{print(total)} 刪掉，
只留 \texttt{print(result)}，這支程式還會不會報錯？
\item 有兩個函式 \texttt{def area(length): result = length * length; return result} 跟
\texttt{def double(n): result = n * 2; return result}，如果依序呼叫
\texttt{print(area(3))} 跟 \texttt{print(double(5))}，兩個函式各自的 \texttt{result}
會不會互相干擾、印出錯誤的值？
\item 有一段程式 \texttt{total = 100}，接著定義
\texttt{def show(): print(total)}，最後呼叫 \texttt{show()}。這支程式會不會報錯？會印
出什麼？（提示：函式主體裡有沒有任何一行對 \texttt{total} 賦值？）
\item 如果把上一題的函式改成
\texttt{def show(): total = total + 1; print(total)}，再呼叫 \texttt{show()}，會發生
什麼事？
\end{drills}

\section{幫函式留下說明：docstring 簡短一瞥}
\label{sec:function-docstring}

函式名稱取得好，通常就能看出它在做什麼，但有些函式的用途沒辦法只靠名字說清楚，例如參數
的單位、回傳值代表什麼。Python 讓你可以在函式主體的第一行，寫一段用三個引號包起來的
說明文字，緊接在 \texttt{def} 那一行後面。

\begin{example}{python}{(temp-convert.py) 附上說明文字的函式}
def celsius_to_fahrenheit(celsius):
    """Convert Celsius temperature to Fahrenheit."""
    return celsius * 9 / 5 + 32


print(celsius_to_fahrenheit(20))
print(celsius_to_fahrenheit.__doc__)
\end{example}

\iosample{}{68.0 \\ Convert Celsius temperature to Fahrenheit. \\}

這段三引號字串不會被當成程式的一部分執行，也不影響函式的回傳值，它單純是說明文字。但
Python 會把它記住，存在函式的 \texttt{\_\_doc\_\_} 這個特殊屬性裡，之後不管是你自己用
\texttt{print(函式名稱.\_\_doc\_\_)} 查詢，還是編輯器在你打字時跳出來的提示視窗，用的
都是這段文字。

這種寫在函式主體第一行的說明字串，有一個正式名字叫 docstring，完整的寫法規則（例如函式
有好幾個參數時該怎麼分段說明）留到第 \ref{ch:modules} 章再細談，這裡你只要知道「函式可以帶一段自我介紹，
而且 Python 有辦法把它查回來」就好。


最後補一句這一章沒有走進去、但值得先知道的事：函式除了「被定義、被呼叫」之外，還有另一面
——它本身也可以像資料一樣被傳來傳去、被塞進別的函式裡、甚至被包上一層改變行為。第
\ref{sec:function-first-class} 節你已經摸到這件事的邊了，完整的討論留到第
\ref{ch:advanced-functions} 章。

\section*{練習解答}
\addcontentsline{toc}{section}{練習解答}

\paragraph{第~\ref{sec:function-multi-return}~節}
\begin{enumerate}
\item \texttt{small} 是 \texttt{1}，\texttt{big} 是 \texttt{9}。
\item 印出 \texttt{(1, 9)}，也就是兩個值打包起來、還沒被拆開的樣子。
\item \texttt{small} 跟 \texttt{big} 都會是 \texttt{5}：清單裡只有一個數字，最小值跟
最大值剛好是同一個數。
\end{enumerate}

\paragraph{第~\ref{sec:function-params}~節}
\begin{enumerate}
\item 印出「I have a pet bird named Tweety.」。
\item \texttt{animal\_type} 會是 \texttt{"cat"}，用 positional 引數呼叫時，第二個引數
一樣會依順序蓋掉預設值。
\item \texttt{make\_coffee(size, milk=True)}，只指定想改的 \texttt{milk}，
\texttt{sugar} 用預設值，\texttt{size} 因為沒有預設值、一定要自己給。
\end{enumerate}

\paragraph{第~\ref{sec:function-args}~節}
\begin{enumerate}
\item 回傳 \texttt{30}：兩個引數被收成一包，迴圈把 \texttt{10} 跟 \texttt{20} 加起來。
\item 會丟出 \texttt{TypeError: total\_price() takes 1 positional argument but 3 were given}。
拿掉星號之後，\texttt{prices} 就變回一個普通參數，只肯收一個引數，這裡卻給了三個。
\end{enumerate}

\paragraph{第~\ref{sec:function-kwargs}~節}
\begin{enumerate}
\item \texttt{show(1, 2, 3)} 印出 \texttt{3}；\texttt{show([1, 2, 3])} 印出
\texttt{1}——後者只給了一個引數（那個 list 本身），星號把「這一個東西」收成一包。
\item 呼叫時在 list 前面加一個星號：\texttt{show(*[1, 2, 3])}，先把 list 拆成三個
獨立的引數再餵進去，就會印出 \texttt{3}。
\item \texttt{greet(**info)}：兩個星號把 dict 攤成
\texttt{greeting="Hi"}、\texttt{name="Amy"} 兩個 keyword 引數，剛好對上函式的兩個參數
名字。
\end{enumerate}

\paragraph{第~\ref{sec:function-first-class}~節}
\begin{enumerate}
\item 印出 \texttt{12}：\texttt{f} 是貼在 \texttt{double} 這個函式物件上的另一張名字
牌，\texttt{f(6)} 跟 \texttt{double(6)} 呼叫的是同一個函式。
\item 印出 \texttt{81}：先算 \texttt{square(3)} 得到 \texttt{9}，再算
\texttt{square(9)} 得到 \texttt{81}。
\item 會在呼叫 \texttt{apply\_twice} 之前就先出錯：\texttt{add\_ten()} 在傳進
\texttt{apply\_twice} 之前就已經被呼叫了，但 \texttt{add\_ten} 需要一個引數
\texttt{n}，這裡沒給，丟出 \texttt{TypeError}。
\end{enumerate}

\paragraph{第~\ref{sec:mutable-default}~節}
\begin{enumerate}
\item 印出 \texttt{['apple', 'banana', 'orange']}，每次沒帶 \texttt{cart} 都在用同一個
共用物件，內容持續疊加。
\item 印出 \texttt{['watermelon', 'banana']}，這次自己傳了 \texttt{cart} 進去，用的是全新的
\texttt{["watermelon"]} 這個物件，不是預設值共用的那個，所以不會看到第一次呼叫留下的
\texttt{"apple"}。
\item 依序印出 \texttt{['apple']}、\texttt{['banana']}、\texttt{['orange']}，因為每次呼叫
都自己判斷、自己造一個新的空 list，彼此互不影響。
\end{enumerate}

\paragraph{第~\ref{sec:function-scope}~節}
\begin{enumerate}
\item 不會報錯：\texttt{result} 是在函式外定義的變數，會正常印出 \texttt{150}；真正會
觸發 \texttt{NameError} 的是嘗試在函式外存取函式內的區域變數 \texttt{total}，這裡已經
被刪掉了。
\item 不會互相干擾：每次呼叫函式，Python 都會替那次呼叫準備一份全新、獨立的區域變數
空間，就算兩個函式都用了 \texttt{result} 這個名字，也是各自獨立的物件，
\texttt{area(3)} 印出 \texttt{9}，\texttt{double(5)} 印出 \texttt{10}。
\item 不會報錯，會印出 \texttt{100}：\texttt{show} 函式主體裡沒有任何一行對
\texttt{total} 賦值，只有讀取，所以 Python 會直接去讀函式外面那個全域變數。
\item 會丟出 \texttt{UnboundLocalError}：主體裡出現了 \texttt{total = total + 1}，
Python 因此把整段主體裡的 \texttt{total} 都當成區域變數，但等號右邊要讀取
\texttt{total} 的時候，這個區域版本還沒有被賦值過。
\end{enumerate}

\chapter{Tuple 與 unpacking}
\label{ch:tuples}

前面幾章你已經學過怎麼用 list 把一堆東西裝在一起、怎麼用 dict 把「名字」對應到「值」。
這一章要認識另一種裝東西的容器，叫 tuple（唸作「thu-pl」或「too-pl」都有人講，沒有標準
答案）。

乍看之下 tuple 很像 list：一樣可以用中括號的方式取出裡面第幾個東西，一樣可以用
\texttt{for} 走訪，一樣可以串接，但它有一個關鍵的差別：\textbf{造出來之後，裡面的內容
就不能再改了}。這聽起來像是缺點（少了一個功能），但正是這個「不能改」的特性，讓 tuple
在很多情境下比 list 更好用，你會在這一章慢慢看到為什麼。
\index{tuple}

這一章分三段推進：

\begin{itemize}
\item 第一段先讓你認識 tuple 這個新容器本身：怎麼造一個 tuple、它跟 list
有哪些操作是共通的（包括串接與重複這兩個你在字串、list 章節都見過的運算子）、然後親手
撞一次「不能修改」這道牆，看看 Python 會丟出什麼樣的錯誤訊息，順便簡單認識一個以 tuple
為基礎、可以額外用名字取值的進階變體。
\item 第二段是這一章真正的重頭戲：把 tuple 跟賦值語法
結合在一起的技巧，統稱「unpacking」（解包）。一行程式碼同時把好幾個值分別存進好幾個變數
裡，這個技巧會讓你少寫很多先存成一個容器、再一個一個挑出來的繁瑣步驟，日常寫程式會一直
用到它，包括等一下你會學到的「函式同時回傳好幾個結果」也是靠它撐起來的，你還會學到用星
號一次接住「剩下的」那些值，以及用底線忽略用不到的值這兩個實用的小技巧。
\item 第三段把視野拉
高一層，正式回答一個從第 \ref{ch:variables} 章就埋下的問題：如果「等號是幫物件貼名字牌，不是複製」這件事
是真的，那為什麼 tuple 看起來「整個換掉」的時候，感覺就像被複製了一份新的一樣？這一段會
用具體的例子（外加一個能親眼看見「這是不是同一個物件」的小工具）證明，其實 tuple 背後遵
守的規則跟 list 一模一樣，只是因為它不能被原地修改，所以任何「看起來像修改」的操作，骨
子裡都只能是「造一個新物件，把名字牌重新貼上去」。
\end{itemize}

\section{認識 Tuple：先像 List，再撞見「不能改」}
\label{sec:tuple-basics}

要造一個 tuple，最常見的寫法是用一對圓括號把幾個值包起來，中間用逗號分隔；其實圓括號
不是必要的，逗號本身才是關鍵，但為了可讀性，大多數時候還是會加上括號。如果手上已經有一
個 list 或其他容器，也可以直接用 \texttt{tuple()} 把它轉換成 tuple。
\index{tuple()}

\begin{code}{python}{建立 tuple 的幾種寫法}
# 用括號建立（最常見的寫法）
t = (x, y, z)

# 逗號才是重點，括號其實可以省略
t = x, y, z

# 從其他序列轉換（例如 list）
t = tuple(some_container)
\end{code}

\begin{example}{python}{(make-coordinate.py) 建立一組座標}
latitude = 25.03
longitude = 121.56
coordinate = (latitude, longitude)
print(coordinate)
print(type(coordinate))

# 也可以把現有的 list 轉換成 tuple
scores_list = [88, 92, 76]
scores_tuple = tuple(scores_list)
print(scores_tuple)
\end{example}

\iosample{}{(25.03, 121.56) \\ <class 'tuple'> \\ (88, 92, 76) \\}

留意 \texttt{print(coordinate)} 印出來的樣子：Python 用圓括號顯示 tuple，這也是為什麼
大部分人習慣用括號來建立它，看起來跟印出來的樣子一致，比較好對照。

\subsection*{tuple 也支援你已經學過的操作}

你會發現，tuple 支援的操作你幾乎都學過，因為第 \ref{ch:lists} 章介紹 list 的時候，用索引取值、切
片、\texttt{in}、\texttt{len()}、\texttt{for} 走訪，這些操作在 tuple 上的寫法跟意義完
全一樣，不需要重新學一次。

\begin{example}{python}{(tuple-like-list.py) tuple 也能索引、切片、走訪}
colors = ("red", "green", "blue")

print(colors[0])          # 索引，跟 list 一樣
print(colors[1:])         # 切片，跟 list 一樣
print("green" in colors)  # 用 in 檢查成員關係
print(len(colors))        # 元素數量

for color in colors:      # for 迴圈走訪，跟 list 一樣
    print(color)
\end{example}

\iosample{}{red \\ ('green', 'blue') \\ True \\ 3 \\ red \\ green \\ blue \\}

\subsection*{tuple 也可以用 \texttt{+} 串接、用 \texttt{*} 重複}

除了上面這幾個操作，第 \ref{ch:strings} 章字串、第 \ref{ch:lists} 章 list 都教過另外兩個運算子，tuple 一樣支援：
用 \texttt{+} 把兩個 tuple 接成一個新的 tuple，用 \texttt{*} 把一個 tuple 的內容重複
好幾次。

\begin{code}{python}{tuple 的串接與重複}
combined = t1 + t2   # 把兩個 tuple 合併成一個新的 tuple
repeated = t * n      # 把 t 的內容重複 n 次
\end{code}

\begin{example}{python}{(tuple-concat-repeat.py) 串接時段、預留固定長度的佔位內容}
morning = ("09:00", "10:00", "11:00")
afternoon = ("14:00", "15:00")
schedule = morning + afternoon
print(schedule)

placeholder = ("-",) * 5
print(placeholder)
\end{example}

\iosample{}{('09:00', '10:00', '11:00', '14:00', '15:00') \\ ('-', '-', '-', '-', '-') \\}

這兩個運算子的行為跟第 \ref{ch:lists} 章 list 版本的 \texttt{+}、\texttt{*} 完全對應：\texttt{+}
不會修改任何一邊原本的 tuple，會造出一個全新的 tuple 把兩邊接起來；\texttt{*} 也不是
把每個元素乘上一個數字，是把整份內容重複貼上好幾次。\texttt{("-",) * 5} 這種寫法在需
要先準備好一份固定長度、內容之後才會確定的容器時很好用。

\subsection*{tuple 之間也可以直接比較大小}

tuple 除了可以用 \texttt{==} 比較兩個 tuple 的內容是否完全一樣，也可以直接用
\texttt{<}、\texttt{>} 這些比較運算子比較「大小」。比較的規則跟比對兩個字串很像：先比
第一個元素，如果分得出高下就直接得到答案，分不出來才往下比第二個、第三個，以此類推，這
種比法有個名字叫「字典順序比較」，因為概念上跟查字典時比較兩個詞先後順序的方式一樣。

\begin{example}{python}{(compare-tuples.py) 比較版本號、依完賽時間排序}
version_a = (1, 4, 2)
version_b = (1, 5, 0)
print(version_a < version_b)

runner_x = ("11:32", "Alice")
runner_y = ("11:28", "Bob")
print(sorted([runner_x, runner_y]))
\end{example}

\iosample{}{True \\{} [('11:28', 'Bob'), ('11:32', 'Alice')] \\}

第一個例子裡，\texttt{version\_a} 和 \texttt{version\_b} 的第一個元素都是
\texttt{1}，分不出高下，於是往下比第二個元素，\texttt{4 < 5}，比較到這裡就有答案了，
後面的元素完全不會再看。第二個例子把「完賽時間、姓名」包成 tuple 放進 list 裡，直接呼
叫第 \ref{ch:lists} 章學過的 \texttt{sorted()}，因為第一個元素（時間字串）就已經能分出順序，
\texttt{sorted()} 自然依完賽時間由先到後排好，這正是把好幾筆相關資料包成 tuple、再交
給排序函式處理的常見寫法。

\stylepoint{把好幾筆相關資料包成 tuple 放進 list 裡，再交給 \texttt{sorted()} 排序，是
很常見的資料處理寫法。因為 tuple 之間的比較會自動照順序一路比下去，你不用自己寫比較邏
輯，只要把最重要的排序依據放在 tuple 的第一個位置就好，這個小技巧之後處理表格型資料時
會很有用。}

\subsection*{撞牆：tuple 不能被修改}

既然 tuple 這麼像 list，你可能會想，那對它做 \texttt{append()} 或直接改某個位置的值，
應該也可以吧？試試看就知道了。

\begin{example}{python}{(change-tuple.py) 嘗試修改 tuple 的內容}
colors = ("red", "green", "blue")
colors[0] = "purple"
\end{example}

存檔執行後，Python 不會讓你得逞：

\begin{verbatim}
Traceback (most recent call last):
  File "change-tuple.py", line 2, in <module>
    colors[0] = "purple"
    ~~~~~~^^^
TypeError: 'tuple' object does not support item assignment
\end{verbatim}

\texttt{TypeError} 是說「這個型別不支援你要做的這個操作」；\texttt{does not support
item assignment} 說得更白：「不支援用索引賦值來改變裡面某一個位置的內容」。這句話你要
牢牢記住，這件事跟「變數到底是什麼」有關，我們留到這章最後一節再回頭講清楚。

那如果改用 \texttt{.append()} 呢？結果甚至不是同一種錯誤：tuple 根本沒有
\texttt{.append()} 這個方法，Python 連「有沒有這個操作」都判斷不出來，直接告訴你這個
方法不存在：

\begin{verbatim}
Traceback (most recent call last):
  File "change-tuple.py", line 2, in <module>
    colors.append("purple")
    ^^^^^^^^^^^^^
AttributeError: 'tuple' object has no attribute 'append'
\end{verbatim}

\keypoint{list 有 \texttt{.append()}、\texttt{.remove()}、\texttt{.sort()} 這些「會改
變自己」的方法，tuple 完全沒有這些方法。這不是「有但被鎖住」，是打從一開始就沒有設計這
些方法，因為 tuple 造出來之後本來就不打算被改變。}

\subsection*{單元素 tuple 的陷阱：別忘了逗號}

還有一個很容易踩到的坑：如果你只想放一個元素進 tuple，光加括號是不夠的。

\begin{example}{python}{(single-element-tuple.py) 忘記逗號的陷阱}
not_a_tuple = ("p")
print(type(not_a_tuple))

really_a_tuple = ("p",)
print(type(really_a_tuple))
\end{example}

\iosample{}{<class 'str'> \\ <class 'tuple'> \\}

原因是：圓括號在 Python 裡本來就有另一個用途，用來改變運算的優先順序（例如
\texttt{(1 + 2) * 3}），所以單獨一個 \texttt{("p")} 會被 Python 當成括號裡包著一個字
串，跟普通的 \texttt{"p"} 沒有兩樣。真正讓 Python 認出「這是 tuple」的關鍵，其實是那
個逗號，\texttt{("p",)} 裡的逗號才是重點，括號只是順便加上去、讓人看得更清楚而已。

\subsection*{簡單認識一個進階變體：namedtuple}

Python 標準函式庫裡有一個叫 \texttt{namedtuple} 的工具，可以幫 tuple 裡的每一個位置
額外取一個名字，讓你除了用 \texttt{coordinate[0]} 這種數字索引，也能用
\texttt{coordinate.latitude} 這種點加名字的方式取值，讀起來更清楚，也比較不容易搞錯
順序。

這個工具需要透過之後幾章才會教到的「把現成工具叫進來使用」的語法才能用，這裡先
不深入，只需要知道有這麼一個東西存在：它本質上還是一個 tuple（一樣不能修改內容），只
是多了一層「幫每個位置取名字」的方便包裝，等你之後真的碰到需要「一組固定欄位、但又想
用名字讀取」的情境，可以回頭來查它怎麼用。

\begin{drills}[label=drills:tuple-basics]
\item 這段程式碼會印出什麼？
\begin{verbatim}
value = ("hello")
print(type(value))
\end{verbatim}
\item 這段程式碼執行後會發生什麼事？寫出你預期看到的錯誤類型（不用寫出完整訊息）。
\begin{verbatim}
sizes = ("S", "M", "L")
sizes[1] = "XL"
\end{verbatim}
\item 這段程式碼執行後會發生什麼事？寫出你預期看到的錯誤類型，並說說看它跟上一題的錯誤
類型為什麼不一樣。
\begin{verbatim}
sizes = ("S", "M", "L")
sizes.append("XL")
\end{verbatim}
\item 這段程式碼會印出什麼？
\begin{verbatim}
weekday = ("Mon", "Tue")
weekend = ("Sat", "Sun")
print(weekday + weekend)
print(weekday * 2)
\end{verbatim}
\item 這段程式碼的執行結果會是 \texttt{True} 還是 \texttt{False}？
\begin{verbatim}
print((2, 9) < (2, 10))
\end{verbatim}
\end{drills}

\section{把 Tuple 跟賦值結合起來：Unpacking}
\label{sec:tuple-unpacking}

知道 tuple 是什麼之後，這一節要介紹這一章真正的重頭戲：把 tuple 跟賦值語法結合起來的
技巧，叫做 unpacking（解包）。概念很直接：如果右邊是一個裝著好幾個值的容器，左邊可以
一次列出好幾個變數名字，Python 會照順序把每個值分別存進對應的變數，一行程式碼就搞定，
不用再寫「先存成一個容器、再用索引一個一個挑出來」的繁瑣步驟。
\index{unpacking}

\begin{code}{python}{多重賦值的基本語法}
name1, name2, name3 = a_container_with_matching_length
\end{code}

\begin{example}{python}{(parse-date.py) 拆解日期字串}
date_string = "2026-09-01"
year, month, day = date_string.split("-")
print(f"Year: {year}, Month: {month}, Day: {day}")
\end{example}

\iosample{}{Year: 2026, Month: 09, Day: 01 \\}

這裡 \texttt{date\_string.split("-")} 會回傳一個裝著三個字串的 list（\texttt{["2026",
"09", "01"]}），unpacking 對 list 也一樣有效，不是只有 tuple 才能用，重點在「右邊的東
西裝了幾個值」，跟它是 tuple 還是 list 無關。跟原本要寫
\texttt{year = date\_string.split("-")[0]} 這種一個個挑出來的寫法比起來，unpacking 一
次把三個有意義的名字都準備好了，可讀性明顯好上不少。

\subsection*{常見錯誤：數量對不上}

unpacking 有一個硬性規定：左邊變數的個數，一定要跟右邊容器裡的元素個數完全一樣，多一
個少一個都不行。

\begin{example}{python}{(unpack-mismatch.py) 賦值數量對不上}
point = (10, 20, 30)
x, y = point
\end{example}

\begin{verbatim}
Traceback (most recent call last):
  File "unpack-mismatch.py", line 2, in <module>
    x, y = point
    ^^^^
ValueError: too many values to unpack (expected 2)
\end{verbatim}

\texttt{too many values to unpack} 是說「右邊的值太多了，左邊的變數裝不下」。反過來，
如果左邊變數比右邊值還多，會得到相反方向的抱怨：

\begin{verbatim}
Traceback (most recent call last):
  File "unpack-mismatch2.py", line 2, in <module>
    x, y, z = point
    ^^^^^^^
ValueError: not enough values to unpack (expected 3, got 2)
\end{verbatim}

兩則訊息都直接把「期望的個數」跟「實際拿到的個數」告訴你，通常一看就知道該加還是該減。

\subsection*{用底線忽略用不到的值}

有時候右邊容器裡的某個值你根本用不到，只是因為數量要對齊，不得不放一個變數名字接住它。
這種情況 Python 社群習慣用底線 \texttt{\_} 當變數名，代表「這個值我知道存在，但故意不
打算用它」。

\begin{example}{python}{(ignore-with-underscore.py) 用底線接住不需要的值}
student = ("Taylor", 16, "Denver")
name, _, city = student
print(f"{name} lives in {city}")
\end{example}

\iosample{}{Taylor lives in Denver \\}

\texttt{\_} 本身其實沒有任何特殊語法地位，它就是一個普通的變數名字，你甚至可以正常
印出 \texttt{print(\_)} 拿到那個被跳過的值。它會被拿來當「這裡故意不用」的訊號，純粹
是 Python 寫作者之間約定俗成的習慣，讓讀你程式碼的人一眼就看出「這個位置的值不重要，
不用花力氣去追它後面有沒有被用到」。

\subsection*{用星號一次接住「剩下的」值}

如果你只在乎前面幾個值，剩下不管有幾個都想整包收進一個 list 裡呢？在某個變數名前面加
一個星號 \texttt{*}，就能把「剩下的所有值」打包成一個 list，交給那個變數。

\begin{code}{python}{星號 unpacking 的三種常見寫法}
first, *rest = some_container      # 單獨接住第一個元素，其餘打包成 list
*init, last = some_container       # 單獨接住最後一個元素，前面打包成 list
first, *middle, last = some_container  # 接住第一個和最後一個，中間打包成 list
\end{code}

\begin{example}{python}{(star-unpacking.py) 用星號拆解一份成績清單}
scores = [95, 88, 76, 82, 91]

top, *others = scores
print(top)
print(others)

*rest, last = scores
print(rest)
print(last)

first, *middle, last = scores
print(first)
print(middle)
print(last)
\end{example}

\iosample{}{95 \\{} [88, 76, 82, 91] \\{} [95, 88, 76, 82] \\ 91 \\ 95 \\{} [88, 76, 82] \\ 91 \\}

留意每一次印出來的 \texttt{others}、\texttt{rest}、\texttt{middle} 都是一個 list，不
是 tuple。星號變數負責收留下來的所有值，數量不管是零個、一個還是很多個都沒問題，Python
會自動打包成一個 list 給你，不需要事先知道剩下確切有幾個值，這正是星號 unpacking 跟前
面「數量一定要對齊」規則最大的不同：一個星號變數可以吃下任意數量的值，讓左右兩邊不再需
要完全一樣長。

一個賦值裡最多只能出現一個星號變數，Python 沒辦法判斷兩個「剩下的」到底
該怎麼切分：

\begin{verbatim}
File "two-stars.py", line 1
    a, *b, *c = [1, 2, 3, 4]
    ^^^^^^^^^
SyntaxError: multiple starred expressions in assignment
\end{verbatim}

星號變數雖然可以吃下任意數量的值，但左邊固定位置的變數還是至少要有值可以對應，如果右邊
的元素數量連固定位置都湊不滿，一樣會噴出 \texttt{ValueError}：

\begin{verbatim}
Traceback (most recent call last):
  File "not-enough-for-star.py", line 1, in <module>
    a, b, *rest = [1]
    ^^^^^^^^^^^
ValueError: not enough values to unpack (expected at least 2, got 1)
\end{verbatim}

留意這裡的措辭是 \texttt{expected at least 2}，跟前面沒有星號時的
\texttt{expected 3, got 2} 不太一樣：因為星號變數願意接受「零個或更多」，Python 只要
求固定位置的變數（這裡是 \texttt{a}、\texttt{b}）一定要湊得滿，星號那個位置本身沒有下
限。

星號 unpacking 在處理「第一列是標題、後面每一列才是真正資料」這種常見的資料格式時特別
好用，之後第 \ref{ch:csv} 章處理 CSV 檔案時會常常遇到這種形狀的資料。

\begin{example}{python}{(skip-header-row.py) 用星號拆出標題與其餘資料列}
raw_data = [
    "Name,Score",
    "Taylor,95",
    "Jordan,88",
]

header, *records = raw_data
print(header)
print(records)
\end{example}

\iosample{}{Name,Score \\{} ['Taylor,95', 'Jordan,88'] \\}

\texttt{header} 接住第一列標題，\texttt{*records} 把剩下不管幾列，全部打包成一個
list，不需要事先知道這份資料到底有幾列，這正是星號 unpacking 比「數清楚長度、再用切片
手動拆」方便的地方。

\subsection*{用 unpacking 交換兩個變數}

unpacking 有一個很受歡迎的小應用：交換兩個變數的值。

\begin{example}{python}{(swap.py) 交換兩個變數的值}
a = "mango"
b = "lemon"
print(f"Before swap: a = {a}, b = {b}")

a, b = b, a

print(f"After swap: a = {a}, b = {b}")
\end{example}

\iosample{}{Before swap: a = mango, b = lemon \\ After swap: a = lemon, b = mango \\}

\stylepoint{很多語言要交換兩個變數的值，得先借一個第三個暫存變數：把 \texttt{a} 存到暫
存變數，再把 \texttt{b} 存進 \texttt{a}，最後把暫存變數存進 \texttt{b}。Python 直接一
行 \texttt{a, b = b, a} 就做完了。右邊的 \texttt{b, a} 先被組成一個 tuple，記住兩個舊
值，再一次 unpack 回左邊的 \texttt{a} 和 \texttt{b}，完全不需要暫存變數，這是很多人覺
得「這樣寫比較 Python」的經典例子。}

這個技巧不只限於兩個變數，三個（或更多）變數要「輪流交換」時，同樣的寫法一樣管用。

\begin{example}{python}{(rotate-three.py) 三個變數輪流交換}
a = 1
b = 2
c = 3
print(f"Before swap: a = {a}, b = {b}, c = {c}")

a, b, c = c, a, b

print(f"After swap: a = {a}, b = {b}, c = {c}")
\end{example}

\iosample{}{Before swap: a = 1, b = 2, c = 3 \\ After swap: a = 3, b = 1, c = 2 \\}

原理跟兩個變數的 swap 完全一樣：右邊的 \texttt{c, a, b} 先被組成一個 tuple，記住三個舊
值，再一次 unpack 回左邊的 \texttt{a}、\texttt{b}、\texttt{c}，每個變數拿到的都是「上
一個位置原本的值」，一次到位，不需要暫存變數，也不需要拆成三行分開賦值。

\subsection*{\texttt{for} 迴圈中的 unpacking}

第 \ref{ch:dict-set} 章你學過用 \texttt{.items()} 走訪字典的 key 跟 value，當時的寫法
\texttt{for key, value in some\_dict.items():} 其實就是 unpacking，\texttt{.items()}
每一輪給你的都是一個 \texttt{(key, value)} 的 tuple，\texttt{for} 迴圈把它直接 unpack
成兩個變數，原來那個語法背後就是這一章教的東西。

\begin{example}{python}{(dict-items-unpacking.py) 走訪字典的 key 與 value}
prices = {"apple": 35, "banana": 20, "orange": 45}

for name, price in prices.items():
    print(f"{name} costs {price} dollars each")
\end{example}

\iosample{}{apple costs 35 dollars each \\ banana costs 20 dollars each \\ orange costs 45 dollars each \\}

\subsection*{巢狀 unpacking：tuple 裡面還裝著 tuple}

如果容器裡的元素本身又是一個 tuple，unpacking 的寫法也可以跟著「巢狀」，只要讓左邊的
形狀跟右邊對得上，Python 就能一路拆到底。

\begin{example}{python}{(nested-unpacking.py) 走訪一份「城市配座標」的清單}
points = [
    ("Boston", (25.03, 121.56)),
    ("Chicago", (24.15, 120.67)),
]

for city, (lat, lon) in points:
    print(f"{city}: latitude {lat}, longitude {lon}")
\end{example}

\iosample{}{Boston: latitude 25.03, longitude 121.56 \\ Chicago: latitude 24.15, longitude 120.67 \\}

\texttt{points} 裡的每一個元素，都是一個 \texttt{(城市名稱, (緯度, 經度))} 這樣兩層的
tuple。\texttt{for} 迴圈的 \texttt{city, (lat, lon)} 這個寫法，形狀跟右邊完全對應：外
層的兩個位置分別對到 \texttt{city} 跟一整個座標 tuple，而座標 tuple 本身又立刻被拆成
\texttt{lat}、\texttt{lon} 兩個變數。左右兩邊的「形狀」要吻合，Python 才知道該怎麼拆，
這跟第 1 節提到的「數量要對齊」是同一個規則，只是這裡多了一層。

\subsection*{函式同時回傳多個值}

第 \ref{ch:functions} 章你學過函式用 \texttt{return} 回傳一個結果。如果想讓函式同時給出好幾個結果呢？
做法是讓 \texttt{return} 後面接好幾個用逗號分隔的值，這其實是在回傳一個 tuple，函式本
質上還是只能回傳一個東西，只是這個東西剛好是一個 tuple。呼叫端再用 unpacking 把它拆開，
分別存進有意義名字的變數。

\begin{example}{python}{(min-max.py) 一次算出最大值和最小值}
def min_max(numbers):
    return min(numbers), max(numbers)


data = [8, 3, 15, 1, 9]
lo, hi = min_max(data)
print(f"Min: {lo}, Max: {hi}")
\end{example}

\iosample{}{Min: 1, Max: 15 \\}

如果呼叫端忘記 unpack，只用一個變數接住呼叫結果，會發生什麼事？並不會出錯，只是你拿到
的會是整個 tuple 本身，沒有被拆開。

\begin{example}{python}{(forget-to-unpack.py) 忘記 unpack 會拿到整個 tuple}
def min_max(numbers):
    return min(numbers), max(numbers)


data = [8, 3, 15, 1, 9]
result = min_max(data)
print(result)
print(type(result))
\end{example}

\iosample{}{(1, 15) \\{} <class 'tuple'> \\}

這正好提醒你，函式回傳多個值這件事，從函式的角度看，它從頭到尾都只回傳了一個東西
（一個 tuple），「同時拿到兩個變數」是呼叫端 unpack 之後才有的結果，不是函式本身的
魔法。忘記 unpack 不會報錯，但你接下來想用 \texttt{lo}、\texttt{hi} 這兩個名字的地方，
就會因為它們根本不存在而噴出 \texttt{NameError}。

其實不用等到你自己寫函式，Python 內建函式裡就有現成的例子。\texttt{divmod(a, b)} 會
一次算出 \texttt{a} 除以 \texttt{b} 的商跟餘數，回傳的正是一個裝著兩個值的 tuple。

\begin{example}{python}{(divmod-example.py) 用內建函式 divmod() 一次拿到商與餘數}
quotient, remainder = divmod(17, 5)
print(f"Quotient is {quotient}, remainder is {remainder}")
\end{example}

\iosample{}{Quotient is 3, remainder is 2 \\}

這正好說明「函式回傳多個值」不是什麼特殊魔法，只要函式回傳的是一個 tuple，呼叫端就能
用這一節學到的 unpacking 語法把它拆開，不管這個函式是你自己寫的，還是 Python 內建的。

\begin{drills}[label=drills:tuple-unpacking]
\item 這段程式碼會印出什麼？
\begin{verbatim}
first = "cat"
second = "dog"
first, second = second, first
print(first, second)
\end{verbatim}
\item 這段程式碼會印出什麼？
\begin{verbatim}
inventory = {"notebook": 3, "pen": 12}
for item, count in inventory.items():
    print(f"{item}: {count}")
\end{verbatim}
\item 這段程式碼會噴出 \texttt{ValueError}。要怎麼修改（不改變 \texttt{record} 那一
行），才能讓它正常執行？
\begin{verbatim}
record = ("Alice", 15, "Boston")
name, age = record
\end{verbatim}
\item 一個函式定義如下，呼叫端要怎麼寫，才能同時拿到 \texttt{total} 跟
\texttt{average} 兩個變數？
\begin{verbatim}
def total_and_average(numbers):
    return sum(numbers), sum(numbers) / len(numbers)
\end{verbatim}
\item 這段程式碼會印出什麼？
\begin{verbatim}
values = [10, 20, 30, 40]
first, *rest = values
print(rest)
\end{verbatim}
\item 這段程式碼會印出什麼？
\begin{verbatim}
pair = ("Boston Station", (25.05, 121.52))
name, (lat, lon) = pair
print(name, lat)
\end{verbatim}
\item 這段程式碼會印出什麼？
\begin{verbatim}
a, b, c = 1, 2, 3
a, b, c = b, c, a
print(a, b, c)
\end{verbatim}
\end{drills}

\section{不可變 vs 可變：名字牌模型的收尾對比}
\label{sec:tuple-vs-list}

還記得第 1 節撞牆的時候留下的那個問題嗎？如果「等號是幫物件貼名字牌，不是複製」這件事
是真的，那為什麼 tuple 看起來「整個換掉」的時候，感覺就像被複製了一份新的一樣？這一節
就要把這件事講清楚。

先看一段程式碼，雖然 tuple 不能被原地修改，但你還是可以寫出「看起來像修改」的程式碼，
例如把某個位置的值換掉、拼接一個新的 tuple 上去：

\begin{example}{python}{(replace-in-tuple.py) 看起來像修改，其實是造新物件}
scores = ("A", "B", "C")
print(scores, id(scores))

scores = ("Z",) + scores[1:]
print(scores, id(scores))
\end{example}

\iosample{}{('A', 'B', 'C') 4349906432 \\ ('Z', 'B', 'C') 4349912128 \\}

（你自己實際跑出來的那兩個數字，幾乎一定跟這裡不一樣，\texttt{id()} 回傳的是物件在記
憶體裡的位址，每次執行、每台機器都可能不同。重點不是數字本身，是\textbf{前後兩個數字不
一樣}。）

\texttt{id()} 是 Python 內建的函式，能告訴你一個物件在記憶體裡的「門牌號碼」，兩個物
件的 \texttt{id()} 一樣，就代表它們是同一個物件。\index{id()}上面這段程式碼裡，\texttt{scores}
「看起來」被改成了 \texttt{('Z', 'B', 'C')}，但 \texttt{id()} 告訴你，第二行印出來的
\texttt{scores} 已經是一個全新的物件了，跟第一行那個完全不是同一個東西。

這正好回答了開頭的問題：\textbf{不是名字牌模型在 tuple 這裡失效了，是 tuple 不可變，所
以每一次「看起來像修改」的操作，都逼著 Python 造一個全新的物件，再把 \texttt{scores}
這張名字牌重新貼上去}。第二行的 \texttt{scores = ("Z",) + scores[1:]} 其實分兩步：先用
\texttt{("Z",) + scores[1:]} 造出一個全新的 tuple，再把 \texttt{scores} 這張名字牌從舊
物件撕下來、貼到這個新物件上，原本那個 \texttt{('A', 'B', 'C')} 並沒有被改變，它只是
不再被任何名字牌指著而已。

拿第 \ref{ch:lists} 章 list 的別名行為回頭對照，會看得更清楚。如果兩張名字牌指著同一個 list，透過其
中一個改動內容，另一個看到的也是改過的結果，因為 list 可以被原地修改，兩張名字牌指的
從頭到尾都是同一個物件。

但如果換成 tuple，同樣「兩張名字牌指著同一個物件」的情況，你不
可能用某張名字牌去「原地修改」它，因為 tuple 根本沒有能原地修改自己的操作，任何看起來
像修改的動作，都只能是造一個新物件、把其中一張名字牌重新貼上去，另一張名字牌仍然貼在原
本那個沒被動過的舊物件上。

\keypoint{對 int、字串、tuple 這類不可變型別，「名字牌模型」跟你原本從數學課或日常生活
帶來的「盒子」直覺，在這裡剛好會給出一樣的結果，因為東西不能被原地修改，每次「看起來
像修改」都只能整個換掉，感覺就很像「盒子裡的內容被換掉了」。但這只是巧合，不是名字牌模
型的例外：底層規則從頭到尾都沒有變，變的只是「這個型別到底允不允許被原地修改」。}

\subsection*{一個更細微的對照：\texttt{+=} 到底有沒有「原地修改」}

第 \ref{ch:loops} 章你學過 \texttt{+=} 是「把右邊的值加到變數原本的值上，再把結果重新貼回同一個名
字牌」的簡寫。這句話聽起來好像 \texttt{+=} 一定會造出新物件，但事情沒有這麼絕對：對
list 使用 \texttt{+=}，Python 其實會直接在原本那個物件上把新內容接上去，物件本身沒有
換過；對 tuple 使用 \texttt{+=}，因為 tuple 沒有任何原地修改自己的操作，只能乖乖造一
個新物件。用 \texttt{id()} 可以親眼看到這個差別。

\begin{example}{python}{(plus-equals-tuple-vs-list.py) 用 id() 比較 += 對 list 和 tuple 的差別}
numbers_list = [1, 2, 3]
before_list_id = id(numbers_list)
numbers_list += [4]
print(numbers_list, id(numbers_list) == before_list_id)

numbers_tuple = (1, 2, 3)
before_tuple_id = id(numbers_tuple)
numbers_tuple += (4,)
print(numbers_tuple, id(numbers_tuple) == before_tuple_id)
\end{example}

\iosample{}{[1, 2, 3, 4] True \\ (1, 2, 3, 4) False \\}

第一行的 \texttt{True} 告訴你，\texttt{numbers\_list} 用 \texttt{+=} 之後，還是原本
那個物件，只是內容被原地擴充了；第二行的 \texttt{False} 告訴你，\texttt{numbers\_tuple}
用 \texttt{+=} 之後，已經是一個全新的物件了。同一個運算子、幾乎一模一樣的寫法，因為
兩邊型別的可變性不同，骨子裡做的事情完全不一樣：這正是這一章從頭到尾一直在強調的重點，
「型別允不允許原地修改」才是決定 Python 背後到底做了什麼事的關鍵，不是看程式碼表面長
得像不像。

\subsection*{tuple 不可變，指的是「這一層」不能換，不是裡面的東西永遠凍結}

這裡有一個很容易被誤會的細節：tuple 保證的是「這個 tuple 本身裝著哪些物件」不能被改
變，但如果裡面裝的某個元素本身是可變的（例如一個 list），那個元素自己還是可以被原地修
改，跟外層的 tuple 完全是兩回事。

\begin{example}{python}{(mutable-inside-tuple.py) tuple 裡面的 list 還是可以被修改}
data = (1, 2, [3, 4])
print(data)

data[2].append(5)
print(data)
\end{example}

\iosample{}{(1, 2, [3, 4]) \\ (1, 2, [3, 4, 5]) \\}

\texttt{data[2] = "something"} 這種「把 tuple 裡第三個位置整個換掉」的寫法一樣會踩到
第 1 節那道牆，觸發 \texttt{TypeError}。但 \texttt{data[2].append(5)} 完全是另一回事：
它沒有動 \texttt{data} 這個 tuple 裡「第三個位置放的是哪個物件」這件事，第三個位置從
頭到尾都指著同一個 list，只是那個 list 自己被原地修改了內容。tuple 只管住自己這一層，
管不到（也沒打算管）裡面的物件要不要允許被修改。

這個細節也解釋了為什麼 hashable（可雜湊）的判斷要看得更仔細一點：如果一個 tuple 裡裝
了像 list 這種不可雜湊的東西，這個 tuple 整體也會變成不可雜湊，一樣不能拿來當 dict 的
key。

\begin{verbatim}
Traceback (most recent call last):
  File "mutable-inside-tuple.py", line 3, in <module>
    d = {data: "some value"}
        ^^^^^^^^^^^^^^^^^^^^
TypeError: unhashable type: 'list'
\end{verbatim}

錯誤訊息抱怨的是 \texttt{'list'} 而不是 \texttt{'tuple'}，因為 Python 檢查一個 tuple
能不能被雜湊時，會連著檢查它裡面每一個元素能不能被雜湊，只要有一個不行，整個 tuple 就
不行。

\keypoint{tuple 保證的不可變是「淺層」的：只保證它自己這一層的內容（裝著哪些物件）不會
被換掉，不保證每個裝進去的物件本身也不能被修改。想要一份資料整個都不能被動到，裡面每
一層也都要用不可變型別組成，例如全部用 tuple、字串、數字組成，不能有任何一層是 list
或 dict。}

\subsection*{不可變的實際好處：tuple 可以當 dict 的 key}

「內容不會變」這件事，除了避免你意外改壞資料，還有一個很實際的好處：tuple 可以拿來當
dict 的 key，而 list 不行。

\begin{example}{python}{(tuple-as-dict-key.py) 用座標當字典的 key}
locations = {
    (25.03, 121.56): "Boston Tower",
    (22.63, 120.30): "Seattle Port",
}
print(locations[(25.03, 121.56)])
\end{example}

\iosample{}{Boston Tower \\}

如果改用 list 當 key，Python 會直接拒絕：

\begin{verbatim}
Traceback (most recent call last):
  File "tuple-as-dict-key.py", line 2, in <module>
    locations = {bad_key: "Boston Tower"}
                ^^^^^^^^^^^^^^^^^^^^^
TypeError: unhashable type: 'list'
\end{verbatim}

\texttt{unhashable} 的意思是「沒辦法被雜湊」，dict 的 key 內部依賴一個叫「雜湊值」的
機制才能快速查找，而這個機制要求 key 的內容永遠不會變，這樣雜湊值才不會失效。list 的內
容隨時可能被改，Python 乾脆不允許它當 key；tuple 造出來就不能再改，正好滿足這個要求，
所以可以放心拿來當 key。\index{hashable}這也是前面提到的 \texttt{namedtuple} 一個常被拿來用的場景：既
要一組固定的欄位當 key，又想讓讀程式碼的人看得出每個位置代表什麼意思。

\begin{drills}[label=drills:tuple-vs-list]
\item 這段程式碼執行完，\texttt{a is b} 會是 \texttt{True} 還是 \texttt{False}？
\begin{verbatim}
a = [1, 2, 3]
b = a
b.append(4)
\end{verbatim}
\item 這段程式碼執行完，\texttt{t is t2} 會是 \texttt{True} 還是 \texttt{False}？
\begin{verbatim}
t = (1, 2, 3)
t2 = t
t2 = t2 + (4,)
\end{verbatim}
\item 承上題，如果把最後一行改成 \texttt{t2 = t}（不做任何加法），\texttt{t is t2}
還會是同樣的答案嗎？為什麼？
\item 這段程式碼執行完，\texttt{original\_id == id(nums)} 會是 \texttt{True} 還是
\texttt{False}？
\begin{verbatim}
nums = [1, 2, 3]
original_id = id(nums)
nums += [4, 5]
\end{verbatim}
\item 這段程式碼執行完，\texttt{t} 的內容會是什麼？tuple 本身有沒有被「修改」？
\begin{verbatim}
t = (1, [2, 3])
t[1].append(4)
\end{verbatim}
\end{drills}

\section*{練習解答}
\addcontentsline{toc}{section}{練習解答}

\paragraph{第~\ref{sec:tuple-basics}~節}
\begin{enumerate}
\item 印出 \texttt{<class 'str'>}，單獨一個括號包住的值不會變成 tuple，少了逗號，
Python 只當它是一個字串。
\item \texttt{TypeError}，tuple 不支援用索引賦值來改變內容
（\texttt{'tuple' object does not support item assignment}）。
\item \texttt{AttributeError}，因為 tuple 從一開始就沒有 \texttt{.append()} 這個方
法，跟上一題「有這個操作但不准做」的 \texttt{TypeError} 不一樣，這裡是「這個操作根本不
存在」。
\item 印出：
\begin{verbatim}
('Mon', 'Tue', 'Sat', 'Sun')
('Mon', 'Tue', 'Mon', 'Tue')
\end{verbatim}
\texttt{+} 把兩個 tuple 接成一個新 tuple，\texttt{*} 把整份內容重複兩次。
\item \texttt{True}，兩個 tuple 第一個元素都是 \texttt{2}，分不出高下，往下比第二個
元素，\texttt{9 < 10} 成立，比較到這裡就有答案了。
\end{enumerate}

\paragraph{第~\ref{sec:tuple-unpacking}~節}
\begin{enumerate}
\item 印出 \texttt{dog cat}，\texttt{first, second = second, first} 先把 \texttt{dog}、
\texttt{cat} 組成一個暫時的 tuple，再 unpack 回去，達成交換。
\item 印出：
\begin{verbatim}
notebook: 3
pen: 12
\end{verbatim}
\item \texttt{record} 有三個元素，但左邊只列了兩個變數，數量對不上。修法是讓左邊也列
出三個變數，例如 \texttt{name, age, city = record}（如果不需要用到 \texttt{city}，也
可以用一個底線當佔位名字：\texttt{name, age, \_ = record}）。
\item \texttt{total, average = total\_and\_average(某個數字 list)}，函式回傳的是一個
裝著兩個值的 tuple，呼叫端用 unpacking 把它拆成兩個變數接住。
\item 印出 \texttt{[20, 30, 40]}，星號變數 \texttt{*rest} 會把除了第一個以外剩下的所
有值打包成一個 list。
\item 印出 \texttt{Boston Station 25.05}，\texttt{name} 對到外層第一個位置，\texttt{(lat, lon)}
這個形狀對到外層第二個位置（本身又是一個座標 tuple），再拆成 \texttt{lat}、
\texttt{lon} 兩個變數，這裡只印出 \texttt{name} 和 \texttt{lat}。
\item 印出 \texttt{2 3 1}，右邊的 \texttt{b, c, a} 先依照舊值組成一個 tuple
\texttt{(2, 3, 1)}，再一次 unpack 回左邊的 \texttt{a}、\texttt{b}、\texttt{c}，等於
把三個值往前輪轉一格。
\end{enumerate}

\paragraph{第~\ref{sec:tuple-vs-list}~節}
\begin{enumerate}
\item \texttt{True}，\texttt{b = a} 只是幫同一個 list 多貼一張名字牌，\texttt{.append()}
是原地修改，兩張名字牌指的還是同一個物件。
\item \texttt{False}，\texttt{t2 = t2 + (4,)} 造出一個全新的 tuple，再把 \texttt{t2}
這張名字牌重新貼到新物件上，\texttt{t} 仍然貼在原本沒被動過的舊物件上。
\item 不會，答案會變成 \texttt{True}，這次沒有做任何「看起來像修改」的操作，
\texttt{t2 = t} 單純是幫同一個物件多貼一張名字牌 \texttt{t2}，跟 list 的別名行為（第一
題）道理完全一樣，只是這裡剛好用在 tuple 上而已。
\item \texttt{True}，對 list 使用 \texttt{+=} 是原地修改，\texttt{nums} 這張名字牌從
頭到尾都貼在同一個物件上，跟前一節 \texttt{plus-equals-tuple-vs-list.py} 範例裡 list
那一半的結果一致。
\item \texttt{t} 會變成 \texttt{(1, [2, 3, 4])}；tuple 本身沒有被修改，它從頭到尾都是
同一個物件、一樣裝著兩個位置，只是第二個位置指著的那個 list 被原地修改了內容。
\end{enumerate}

\chapter{檔案與例外處理}
\label{ch:files-exceptions}

這一章要教的，表面上是兩件不相干的事：讀寫檔案，還有處理錯誤。但骨子裡它們是同一套邏輯的
兩面，所以我們把它們放進同一章，而且刻意用同一個順序講給你聽。

前面九章你寫的每一支程式，只要視窗一關、程式一停，變數、清單、字典裡存的東西就通通消失，
下次執行又是從頭開始。這一章第一段要解決的就是這件事：怎麼把資料寫進一個檔案裡，讓它撐過
程式結束、撐過電腦重開機，下次程式再打開來還在原地。我們會先教怎麼讀一個已經存在的檔案，
再教怎麼自己寫一個，也會教你怎麼一次同時打開兩個檔案、一邊讀一邊寫。

只要程式開始跟外面的世界打交道，不管是打開一個檔案，還是之後章節會遇到的網路連線、使用者
亂打的輸入，就一定會遇到「這件事可能失敗」的狀況：檔案可能根本不存在、路徑打錯字、硬碟臨時
被拔掉。這一章的第二段就是在處理「失敗」這件事，而且會正式告訴你 Python 社群面對「可能失敗
的操作」有一套特別的處理哲學，跟大多數人「先確認安全才動手」的直覺剛好相反。

這就是這本書從第一章就預告過的兩個貫穿全書心智模型之一，這裡是它第一次完整登場，之後每一章
遇到「可能出錯的操作」都會回頭呼應這裡教的東西。你也會學到同一段程式碼怎麼同時防範好幾種不
同的失敗原因，以及不管成功或失敗都一定要做的收尾動作要怎麼寫。

最後，這一章會把視角拉得更高：這套「先做、錯了再處理」的邏輯，其實不只是檔案操作專屬的做
法。回頭看看第 \ref{ch:conditionals} 章講條件判斷、第 \ref{ch:dict-set} 章講字典查詢時，我們已經偷偷埋了兩個相關的伏筆，這一章
最後一節會把它們一次收尾，讓你看到這一套邏輯貫穿的範圍，比你以為的更大。

\section{讀一份檔案：\texttt{open()} 與 \texttt{read()}}
\label{sec:reading-files}

到目前為止，你程式裡用到的資料，不是你自己在程式碼裡打好的，就是使用者當下打進去的。程式
一結束，這些東西就沒了。但現實世界裡，資料常常早就已經存在某個地方，例如電腦裡的一份文字
檔，這一節要教你怎麼把這種已經存在的檔案讀進程式裡。

想像你去一個有寄物櫃的地方寄放行李：櫃檯會給你一張號碼牌，之後你要拿行李，靠的是這張號碼
牌，號碼牌本身不是行李。Python 的 \texttt{open()} 函式做的事很像：呼叫它會回傳一個「檔案
物件」，你可以把它想成那張號碼牌，之後要讀檔案內容，都要透過這個物件進行操作，物件本身不是
內容。
\index{檔案物件}

\begin{code}{python}{開檔並讀出全部內容}
with open(filename) as file_object:
    content = file_object.read()
\end{code}

\texttt{open(filename)} 打開檔案、回傳檔案物件；\texttt{read()} 才是真正把整份內容讀出來，變成
一個字串。
\index{open()}
\index{read()}
前面的 \texttt{with ... as ...:} 這個寫法，是在告訴 Python「這段縮排區塊做完之
後，自動幫我把檔案關掉」。延續寄物櫃的比喻，就像「用完自動幫你把號碼牌繳回去、置物櫃自動關
好」，你不用自己記得要不要去櫃檯歸還。你之後在很多情境都會看到 \texttt{with} 這個寫法，不只
用在檔案上。
\index{with ... as ...}

\keypoint{\texttt{open()} 回傳的東西是一個代表「已經打開的這個檔案」的物件，不是檔案裡的
文字內容本身，就像寄物櫃的號碼牌不是你的行李本身。真正要拿到檔案裡的文字，要另外呼叫
\texttt{.read()} 之類的方法才行。}

這裡有一個小提醒：範例裡 \texttt{open("shopping\_list.txt")} 這樣只寫檔名、沒有夾帶資料夾
路徑，Python 預設會去「目前執行程式的那個資料夾」裡找，也就是你的 \texttt{.py} 檔案所在的
那個資料夾。如果你把 \texttt{shopping\_list.txt} 放在別的資料夾，或者用編輯器從不同位置執行
程式，\texttt{open()} 就可能找不到它。這一章後面談到「檔案可能不存在」時，這也是最常見的原因
之一：檔案不一定真的不見了，很可能只是程式在錯的地方找。

假設你的資料夾裡有一份 \texttt{shopping\_list.txt}，內容是一行一項要買的東西：

\begin{verbatim}
Milk
Eggs
Toast
Apple
\end{verbatim}

\begin{example}{python}{(read-shopping-list.py) 讀一份購物清單}
with open("shopping_list.txt") as shopping_list_file:
    content = shopping_list_file.read()

print(content)
print("---")
print(content.rstrip())
\end{example}

\iosample{}{Milk \\ Eggs \\ Toast \\ Apple \\~\\ --- \\ Milk \\ Eggs \\ Toast \\ Apple \\}

第一次 \texttt{print(content)} 之後，「Apple」下面多了一行空白，這是因為
\texttt{shopping\_list.txt} 這份檔案本身結尾就有一個看不見的換行符號，\texttt{read()} 把
它原封不動讀進來，\texttt{print()} 自己又會再加一次換行，兩個疊在一起就變成一行空白。
\texttt{content.rstrip()} 會把字串「結尾」的空白字元（包含換行）去掉，所以第二次印出來就
乾淨俐落，沒有多餘的空行。
\index{rstrip()}

你可能已經在想：如果 \texttt{shopping\_list.txt} 這個檔名根本打錯字、或者這個檔案壓根不存
在，\texttt{open()} 會怎樣？先賣個關子，這個問題我們晚一點會正式回來處理。

\section{逐行讀取檔案}
\label{sec:reading-lines}

上一節整份讀出來很方便，但很多檔案的資料是「一行一筆」的形狀：一行一個名字、一行一筆紀錄，
這時候你通常想要的是能一行一行分開處理，不是一大坨字串。

\begin{code}{python}{逐行讀取檔案}
with open(filename) as file_object:
    for line in file_object:
        process(line.rstrip())
\end{code}

檔案物件本身就是可以放進 \texttt{for} 迴圈裡的東西：每一輪迴圈，\texttt{line} 會是檔案裡的
「下一行」，而且包含結尾那個看不見的換行符號，所以幾乎每次都要搭配 \texttt{rstrip()} 把它去
掉。

假設你有一份 \texttt{names.txt}，一行一個名字：

\begin{verbatim}
Alice
Bob
Taylor
\end{verbatim}

\begin{example}{python}{(greet-each-name.py) 跟名單裡每個人打招呼}
with open("names.txt") as names_file:
    for line in names_file:
        name = line.rstrip()
        print(f"Hello, {name}!")
\end{example}

\iosample{}{Hello, Alice! \\ Hello, Bob! \\ Hello, Taylor! \\}

如果忘記呼叫 \texttt{rstrip()}，把上面的 \texttt{name = line.rstrip()} 改成
\texttt{name = line}，輸出會在每句問候後面多出一堆空行，因為 \texttt{line} 裡其實藏著一個
看不見的換行符號，用括號把每一行框起來就能看清楚它到底藏在哪裡：

\begin{example}{python}{(show-hidden-newline.py) 看清楚每一行藏著的換行符號}
with open("names.txt") as names_file:
    for line in names_file:
        print(f"[{line}]")
\end{example}

\begin{verbatim}
[Alice
]
[Bob
]
[Taylor
]
\end{verbatim}

方括號的結尾跑到了下一行，這就是那個藏在 \texttt{line} 裡、\texttt{rstrip()} 沒被呼叫就會
一直跟著你的換行符號。

除了逐行處理，Python 還有一個方法叫 \texttt{readlines()}，
\index{readlines()}
會一次把每一行變成一個字串，存進
一個 list 裡回傳給你（例如 \texttt{names\_file.readlines()} 會得到
\texttt{['Alice\textbackslash n', 'Bob\textbackslash n', 'Taylor\textbackslash n']}），跟
\texttt{for line in file\_object} 做的事情類似，只是拿到的資料形狀不一樣，一個是「邊讀邊處
理」，一個是「先整包存成 list，之後再慢慢用」。這裡先讓你認得這個名字，之後真正需要一次拿到
所有行的 list 時，你會知道有這個選項。

這裡值得多想一步：\texttt{for line in file\_object} 是一行一行來，處理完這一行才去拿下一行，
同一時間電腦記憶體裡只需要放著「目前這一行」；\texttt{readlines()} 則是一口氣把整份檔案的每
一行都讀出來、變成 list 存好，才把整個 list 交給你。

檔案不大的時候兩者感覺不出差別，但如果檔案有幾百萬行，\texttt{readlines()} 要先把這麼多行全
部塞進記憶體才能開始用，\texttt{for line in file\_object} 卻完全不需要，因為它從頭到尾都只
處理「目前這一行」。這也是為什麼很多 Python 使用者的預設習慣是逐行處理，只有真的需要「一次
拿到所有行的清單」（例如要整份倒著讀、或要跳來跳去存取指定的某幾行）才會特地呼叫
\texttt{readlines()}。

\begin{drills}[label=drills:reading-files]
\item 假設 \texttt{fruits.txt} 的內容是：
\begin{verbatim}
Orange
Banana
\end{verbatim}
這段程式執行後，螢幕上會印出什麼？
\begin{verbatim}
with open("fruits.txt") as fruits_file:
    for line in fruits_file:
        print(line)
\end{verbatim}
\item 如果把上面的 \texttt{print(line)} 改成 \texttt{print(line.rstrip())}，輸出會有什麼
不同？
\item \texttt{readlines()} 會把整份檔案的每一行變成 list 裡的一個元素。如果你要做的事情是
「統計檔案總共有幾行」，你覺得用 \texttt{readlines()} 拿到的 list 直接算 \texttt{len()}，
跟用 \texttt{for line in file\_object} 自己數一個計數器，哪一種寫起來比較直接？
\end{drills}

\section{把資料寫進檔案}
\label{sec:writing-files}

讀資料是把別人（或你自己之前）留下的東西拿進來，這一節反過來，讓你的程式自己在檔案裡留下
東西，之後不管是你自己還是別的程式，都能再讀回來。

在動手學新模式之前，先看一次如果用第~\ref{sec:reading-files}~節那種預設的讀取方式，硬要呼
叫 \texttt{write()} 會發生什麼事：

\begin{verbatim}
with open("shopping_list.txt") as f:
    f.write("new item\n")
\end{verbatim}

執行後：

\begin{verbatim}
Traceback (most recent call last):
  File "add-item-wrong-mode.py", line 2, in <module>
    f.write("new item\n")
io.UnsupportedOperation: not writable
\end{verbatim}

\texttt{open()} 沒有另外指定模式時，預設是唯讀（只能讀，不能寫）。這個檔案物件上雖然也有
\texttt{write()} 這個方法可以呼叫，但作業系統一開始就是用「唯讀」的權限打開這份檔案，硬要
寫入，就會被擋下來、丟出 \texttt{io.UnsupportedOperation}。
\index{io.UnsupportedOperation}
要真正能寫入檔案，得在
\texttt{open()} 明確指定寫入用的模式，這就是這一節要教的事。

這裡換一個更貼切的類比：白板。\texttt{'w'} 模式就像先把白板整個擦乾淨，才開始寫新的內容，
不管白板上原本寫了什麼，通通不見。
\index{'w' 模式}

\begin{code}{python}{寫入模式開檔並寫入內容}
with open(filename, "w") as file_object:
    file_object.write(content)
\end{code}

跟 \texttt{print()} 不一樣，\texttt{write()} 不會自動幫你加上換行符號。想要換行，要自己在
字串裡寫 \texttt{\textbackslash n}。
\index{write()}

\begin{example}{python}{(write-diary-entry.py) 把兩句話存進日記檔案}
with open("diary.txt", "w") as diary_file:
    diary_file.write("Today I learned how to open files")
    diary_file.write("I also learned how to write files")

with open("diary.txt") as diary_file:
    print(diary_file.read())
\end{example}

\iosample{}{Today I learned how to open filesI also learned how to write files \\}

兩次 \texttt{write()} 呼叫寫出來的字串直接黏在一起，中間完全沒有分隔，這就是「\texttt{write()}
不會自動加換行」的具體長相。如果想讓它們分成兩行，要自己在第一句字串尾巴加上
\texttt{\textbackslash n}，例如把第一個 \texttt{write()} 改成
\texttt{diary\_file.write("Today I learned how to open files\textbackslash n")}。

\texttt{'w'} 模式還有一個新手很容易踩到的坑：如果你之後又用 \texttt{'w'} 模式打開同一個檔
案，不管上次寫了什麼，都會被整個清空重寫。很多人會誤以為 \texttt{'w'} 是「接著寫」，其實它
每次打開都是先擦乾淨再寫。

\begin{example}{python}{(overwrite-diary-mistake.py) 常見錯誤：以為 'w' 是「加寫」}
with open("diary.txt", "w") as diary_file:
    diary_file.write("Today I was busy and did not get anything done")

with open("diary.txt") as diary_file:
    print(diary_file.read())
\end{example}

\iosample{}{Today I was busy and did not get anything done \\}

把這段程式接在剛才 \texttt{write-diary-entry.py} 之後執行，打開 \texttt{diary.txt} 就只剩
下「Today I was busy and did not get anything done」這一句。「Today I learned how to
open filesI also learned how to write files」已經被整個蓋
掉，沒有任何警告，也救不回來。

\keypoint{\texttt{'w'} 模式每次打開檔案，都會先把整個檔案清空，不管你原本寫了什麼。它的行為
是「重寫」，不是「加寫」。如果你想在原有內容後面繼續添加，需要的是下一節要教的 \texttt{'a'}
模式，不是 \texttt{'w'}。}

\section{一次寫入 vs 逐行寫入：\texttt{writelines()} 與寫入迴圈}
\label{sec:writing-styles}

上一節看到的 \texttt{write-diary-entry.py}，兩次 \texttt{write()} 呼叫各自寫一句話。如果
要寫進檔案的內容本來就是一份 list，例如處理過的購物清單，你有好幾種寫法可以選，這一節把常見
的幾種放在一起比較。

第一種寫法，是在迴圈裡逐項呼叫 \texttt{write()}：

\begin{code}{python}{在迴圈裡逐項寫入}
with open(filename, "w") as file_object:
    for item in items:
        file_object.write(item + "\n")
\end{code}

每一輪迴圈呼叫一次 \texttt{write()}，等於總共呼叫了跟清單長度一樣多次；因為 \texttt{write()}
不會自動換行，每一輪都要自己接上 \texttt{\textbackslash n}。

第二種寫法，是用檔案物件的另一個方法 \texttt{writelines()}，一次接收一份「字串構成的
list」，把每一個字串依序寫進檔案：
\index{writelines()}

\begin{example}{python}{(write-cart-with-writelines.py) 用 writelines() 一次寫進整份清單}
items = ["Milk\n", "Eggs\n", "Toast\n"]

with open("cart.txt", "w") as cart_file:
    cart_file.writelines(items)

with open("cart.txt") as cart_file:
    print(cart_file.read())
\end{example}

\iosample{}{Milk \\ Eggs \\ Toast \\}

\texttt{writelines()} 這個名字容易讓人誤會它會自動幫每一行加上換行，其實不會。它做的事情就
只是把 list 裡每一個字串照順序接在一起寫進檔案，跟 \texttt{write()} 一樣完全不會自動插入任何
分隔符號。上面的例子如果把 \texttt{items} 換成 \texttt{["Milk", "Eggs", "Toast"]}（每個字串
結尾都沒有 \texttt{\textbackslash n}），寫出來的 \texttt{cart.txt} 就會是「MilkEggsToast」
黏成一坨，跟第~\ref{sec:writing-files}~節 \texttt{write-diary-entry.py} 踩到的坑一模一樣。

第三種寫法，是先把整份清單組成一個大字串，只呼叫一次 \texttt{write()}：

\begin{example}{python}{(write-cart-joined.py) 先組好整份內容，一次寫入}
items = ["Milk", "Eggs", "Toast"]
content = "\n".join(items) + "\n"

with open("cart.txt", "w") as cart_file:
    cart_file.write(content)

with open("cart.txt") as cart_file:
    print(cart_file.read())
\end{example}

\iosample{}{Milk \\ Eggs \\ Toast \\}

三種寫法產生的結果完全一樣，差別在「什麼時候真的去動檔案」。迴圈裡逐項 \texttt{write()}，
每一輪都跟檔案系統打一次交道；\texttt{writelines()} 跟「先 \texttt{join} 好、一次
\texttt{write()}」則是先把要寫的內容準備好，只跟檔案系統打一次交道。

對很小的檔案，這個差異感覺不出來；當清單長到成千上萬筆，每一輪都個別呼叫 \texttt{write()}
累積起來的額外開銷就會比較明顯，這時候先準備好整份內容、一次寫入，通常比較有效率。但如果清
單本身太大，大到記憶體放不下一整份組好的字串，逐項寫反而是唯一的選擇。這呼應
第~\ref{sec:reading-lines}~節讀檔案時「逐行讀」跟「\texttt{readlines()} 一次全讀」的取捨，
讀跟寫是同一個道理的兩面。

\keypoint{\texttt{writelines()} 跟 \texttt{write()} 一樣，不會自動幫你在每一項之間加上換
行或任何分隔符號。list 裡的每個字串要不要換行，得自己準備好。}

\begin{drills}[label=drills:writing-styles]
\item 如果 \texttt{write-cart-with-writelines.py} 裡的 \texttt{items} 沒有在每個字串結尾
加 \texttt{\textbackslash n}（也就是寫成 \texttt{items = ["Milk", "Eggs", "Toast"]}），
\texttt{writelines()} 寫出來的 \texttt{cart.txt} 內容會是什麼？
\item 如果有一百萬行資料要寫進同一個檔案，你會選擇「迴圈裡逐項 \texttt{write()}」，還是「先
組好一個超大字串再一次 \texttt{write()}」？兩種各自要注意什麼？
\end{drills}

\section{附加寫入：\texttt{'a'} 模式}
\label{sec:appending-files}

第~\ref{sec:writing-files}~節看到 \texttt{'w'} 模式的坑：每次打開都會清空重寫。如果你要的
其實是「接著上次寫的地方繼續往下寫」，像日記一樣，今天寫的不該把昨天寫的洗掉，你要的是
\texttt{'a'}（append，附加）模式。延續白板的類比，這就是「接著上次寫的地方，往下繼續寫，不
擦掉原本的內容」。
\index{'a' 模式}

不需要重新學一套片語，把 \texttt{'w'} 換成 \texttt{'a'} 就好：

\begin{example}{python}{(append-diary-entry.py) 用附加模式接續寫日記}
diary_entry = "Today I learned how to write files"

with open("diary.txt", "a") as diary_file:
    diary_file.write(diary_entry + "\n")

with open("diary.txt") as diary_file:
    print(diary_file.read())
\end{example}

第一次執行（假設 \texttt{diary.txt} 原本不存在，\texttt{'a'} 模式會自動幫你新建一個）：

\iosample{}{Today I learned how to write files \\}

隔天你把 \texttt{diary\_entry} 改成另一句話，重新執行同一支程式（\texttt{diary.txt} 沒有
被刪掉）：

\iosample{}{Today I learned how to write files \\ Today I reviewed for loops \\}

兩次執行的內容都留下來了。這就是 \texttt{'a'} 跟 \texttt{'w'} 唯一的差別：如果檔案已經存
在，\texttt{'w'} 清空重寫，\texttt{'a'} 接著寫；如果檔案原本就不存在，兩者的行為其實一樣，
都會自動幫你新建一個。

\begin{drills}[label=drills:appending-files]
\item 你想在使用者每次執行程式時都保留之前的紀錄，不要把上次寫的東西洗掉，該用哪個模式，
\texttt{'w'} 還是 \texttt{'a'}？
\item 你想確保每次執行都是全新的一份紀錄，不要保留上次的內容，該用哪個模式，\texttt{'w'}
還是 \texttt{'a'}？
\item 如果你用 \texttt{'a'} 模式打開一個原本根本不存在的檔案，會發生什麼事？跟 \texttt{'w'}
模式遇到「檔案不存在」時的行為一樣嗎？
\end{drills}

\section{同時開兩個檔案：一次讀、一次寫}
\label{sec:multiple-files}

讀檔案、寫檔案分開學完之後，很多實際任務其實是「讀一份、依照內容做點加工、寫進另一份」，例
如把一份購物清單轉成帶編號的檢查清單。你當然可以先用一個 \texttt{with} 讀完存進一個變數，讀
完之後再另外用一個 \texttt{with} 把結果寫出去，但 Python 也允許你用同一個 \texttt{with}
同時打開好幾個檔案，中間用逗號隔開。
\index{with 逗號同時開多個檔案}

\begin{code}{python}{同一個 with 打開兩個檔案}
with open(source_filename) as source_file, open(target_filename, "w") as target_file:
    for line in source_file:
        target_file.write(processed_content)
\end{code}

逗號前後各是一次完整的 \texttt{open(...) as ...}，兩個檔案物件在同一個縮排區塊裡都用得到；
區塊結束時，Python 會把兩個檔案都關掉，不用分別寫兩層 \texttt{with}。

\begin{example}{python}{(number-shopping-list.py) 讀一份清單，寫出帶編號的檢查清單}
number = 0

with open("shopping_list.txt") as source_file, open("shopping_checklist.txt", "w") as target_file:
    for line in source_file:
        item = line.rstrip()
        number += 1
        target_file.write(f"{number}. {item}\n")

with open("shopping_checklist.txt") as checklist_file:
    print(checklist_file.read())
\end{example}

\iosample{}{1. Milk \\ 2. Eggs \\ 3. Toast \\ 4. Apple \\}

\texttt{number} 變數用的是第 \ref{ch:loops} 章教過的「計數」寫法，每讀到一行就加 1；f-string 把數字跟項
目接在一起組成一行文字。整段迴圈同時在做兩件事：一邊用 \texttt{source\_file} 逐行讀，一邊用
\texttt{target\_file} 逐行寫，兩個檔案物件互不干擾。

\keypoint{千萬不要把「\texttt{source\_filename}」跟「\texttt{target\_filename}」設成同一個檔案。\texttt{'w'} 模式的清空動作，
發生在 \texttt{open()} 這一行執行的當下，不是等到第一次呼叫 \texttt{write()} 才清空。如果來
源、目標是同一個檔案，目標那邊的 \texttt{open(filename, "w")} 一執行就把整個檔案清空了，等
\texttt{for line in source\_file} 想讀取內容時，能讀到的只剩下空的檔案，什麼都讀不到，也就什麼
都寫不出來。想「原地修改」一份檔案，正確做法是先讀出全部內容存進一個變數，確定來源檔案的
\texttt{with} 區塊已經讀取完成，再另外開一次同一個檔名、用 \texttt{'w'} 模式寫回去。}

\begin{drills}[label=drills:multiple-files]
\item 如果把 \texttt{number-shopping-list.py} 裡的 \texttt{"shopping\_checklist.txt"} 也
換成 \texttt{"shopping\_list.txt"}（跟來源檔案同一個檔名），執行後 \texttt{shopping\_list.txt}
會變成什麼？
\item 承上題，如果真的想把 \texttt{shopping\_list.txt} 原地改成帶編號的版本（讀跟寫都是同一
個檔案），應該怎麼修改這段程式，才能避免上一題的問題？
\end{drills}

\section{Python 怎麼面對「可能失敗」：LBYL 與 EAFP}
\label{sec:lbyl-vs-eafp}

從第一節到現在，我們示範的檔案都乖乖存在，一切順利。但真實世界不會這麼聽話。檔案可能根本不
存在、路徑打錯字，甚至檔案在你打開的那一刻剛好被別人搬走或刪掉。這一節要教你 Python 怎麼面
對「這件事可能失敗」的狀況，而且會告訴你，Python 社群面對這種狀況的直覺，可能跟你原本以為
的不太一樣。

想像你去自動提款機領錢。大多數人的直覺是「先確認安全才行動」：出門前先打電話問銀行，帳戶裡
還有沒有夠的錢，確定夠了才出門去領。這種「先檢查、確定安全才做」的習慣，有一個專門的名字，
叫 LBYL（look before you leap，先看清楚再跳）。
\index{LBYL}

但你實際去操作提款機的時候，通常不會這樣：你直接把卡插進去、輸入金額，如果餘額真的不夠，
提款機會直接告訴你「餘額不足」，你再想辦法，不會因為多打一通電話確認，事情就變得更好。
Python 面對「可能失敗的操作」時，比較習慣的做法就是後面這一種：直接做，如果真的出了問題，
程式會告訴你哪裡出錯，你再處理。這種「先做，錯了再處理」的習慣，也有專門的名字，叫 EAFP
（easier to ask forgiveness than permission，先做了再說比先問許可容易）。
\index{EAFP}

開檔案的時候道理也一樣。如果完全不管檔案存不存在，直接打開一個不存在的檔案，會發生什麼事？

\begin{example}{python}{(open-todo-list.py) 直接打開一個不存在的檔案}
todo_file = open("todo.txt")
print(todo_file.read())
\end{example}

假設資料夾裡根本沒有 \texttt{todo.txt} 這個檔案，執行後會看到：

\begin{verbatim}
Traceback (most recent call last):
  File "open-todo-list.py", line 1, in <module>
    todo_file = open("todo.txt")
FileNotFoundError: [Errno 2] No such file or directory: 'todo.txt'
\end{verbatim}

由下往上讀（呼應第 \ref{ch:introduction} 章教過的讀法）：最後一行 \texttt{FileNotFoundError} 就是 Python 專門
用來表示「找不到檔案」的例外類型；後面 \texttt{[Errno 2] No such file or directory} 是作業
系統層級更詳細的說明。程式沒有處理過這種情況，所以整支程式直接當掉。
\index{FileNotFoundError}

你原本的直覺可能是：那先檢查檔案在不在，確定存在才打開它，不就不會當掉了嗎？這是很合理的想
法，寫成程式碼像這樣：

\begin{code}{python}{LBYL：先檢查檔案在不在，才敢打開}
import os

if os.path.exists(filename):
    with open(filename) as file_object:
        content = file_object.read()
else:
    handle_missing_file()
\end{code}

\texttt{import os} 這一行，是在告訴 Python「我要用到 \texttt{os} 這批工具」。\texttt{os}
是 Python 內建、專門處理作業系統與檔案路徑的一批現成程式碼，不用你自己寫，\texttt{import}
就是把這批工具叫進來給你用；這裡你只需要先照著打就好，\texttt{import} 真正的運作方式，第
\ref{ch:modules} 章會完整介紹。
\index{import}
\texttt{os.path.exists(filename)} 傳入一個檔名（路徑字串），檔案存在就回傳
\texttt{True}，不存在就回傳 \texttt{False}，這就是為什麼上面可以直接把它放進 \texttt{if}
裡判斷。
\index{os.path.exists()}

\begin{example}{python}{(check-file-before-reading.py) 先檢查再開檔}
import os

filename = "todo.txt"

if os.path.exists(filename):
    with open(filename) as todo_file:
        print(todo_file.read())
else:
    print("This file was not found, skipping for now")
\end{example}

\iosample{}{This file was not found, skipping for now \\}

這段程式確實不會當掉了。但它藏著一個容易被忽略的坑：\texttt{os.path.exists()} 檢查完、到真
正呼叫 \texttt{open()} 之間，其實有一段很短暫的空檔，如果剛好在這個空檔裡，這個檔案被搬走、
被刪掉，或者換了另一個程式在寫入，檢查時看到「存在」，真正打開時它卻已經不在了。

單獨一個人在自己的電腦上跑這種小程式，這個空檔幾乎不會出問題；但只要牽涉到多個程式、多個人
同時讀寫同一批檔案（例如伺服器上好幾個服務同時處理同一批上傳的檔案），這種「檢查完到動手之
間狀態已經變了」的情況，就沒有想像中罕見。這就跟提款機的比喻一樣，先打電話問餘額，到真正插
卡領錢之間，帳戶餘額也可能已經變了（例如剛好有一筆扣款進來），「檢查過」不代表「等一下動手
時還是一樣的狀態」。

現在改用 EAFP 的寫法，做的事情完全一樣，但邏輯反過來，不先問，直接做，失敗了再處理：

\begin{code}{python}{EAFP：直接開檔，失敗了再處理}
try:
    with open(filename) as file_object:
        content = file_object.read()
except FileNotFoundError:
    handle_missing_file()
\end{code}

\begin{example}{python}{(open-file-safely.py) 直接開檔，失敗了再處理}
filename = "todo.txt"

try:
    with open(filename) as todo_file:
        print(todo_file.read())
except FileNotFoundError:
    print("This file was not found, skipping for now")
\end{example}

\iosample{}{This file was not found, skipping for now \\}

兩段程式做的事情一模一樣、印出來的結果也一模一樣，差別只在「怎麼面對可能失敗這件事」：LBYL
先問、EAFP 先做。Python 比較習慣後者，原因不是 LBYL 有什麼錯，關鍵在於「檢查完到真正動手之
間，狀態可能已經變了」這種坑，EAFP 完全不會踩到。不用先確認任何事，直接做，Python 會在真正
失敗的那一刻老實告訴你發生了什麼事。

EAFP 還有一個實際的好處：LBYL 常常等於把同一件事做兩次，先用 \texttt{os.path.exists()} 問
一次「這個檔案存在嗎」，通過了才真的呼叫 \texttt{open()} 再問一次「這個檔案真的能打開嗎」
（作業系統實際檢查的動作，兩次都要重新來一遍）。EAFP 只做一次：直接 \texttt{open()}，成功就
是成功，失敗就用 \texttt{except} 接住，沒有「先問一次、真正做的時候再確認一次」這種重複的
成本。

\texttt{try} 後面接的縮排區塊，是「可能會失敗」的那段程式碼；\texttt{except} 後面接的
\texttt{FileNotFoundError}，是「你想接住的那一種錯誤類型」。只有真的發生這一種錯誤，下面的
區塊才會執行，其他種類的錯誤不會被這裡攔下來
\index{try/except}
（第~\ref{sec:try-except-else}~節的練習會讓你
親眼看到這件事）。

\keypoint{LBYL 檢查完到真正動手之間，永遠存在一個「狀態可能已經變了」的空檔，任何「先確認再
行動」的邏輯都躲不掉這個坑，不只是 Python 才有這個問題。EAFP 直接做、失敗了再處理，剛好把這
個空檔完全消掉，這是 Python 社群偏好 EAFP 的核心原因，不只是「比較潮」而已。}

\section{try/except/else：把「只有成功才該做」的事分開放}
\label{sec:try-except-else}

上一節的 \texttt{except} 區塊，放的是「失敗了要做的事」。但如果想接著做一些「只有成功才該
做」的事，例如剛剛讀到的內容還要拿去做進一步的處理，直接把這些程式碼也塞進 \texttt{try}
區塊裡雖然能動，但不是 Python 建議的寫法。這一節教你第三塊：\texttt{else}。

\begin{code}{python}{try/except/else}
try:
    value = operation_that_might_fail()
except ExceptionType:
    handle_failure()
else:
    do_something_on_success(value)
\end{code}

\texttt{else} 區塊只有在 \texttt{try} 裡的程式碼完全沒有出錯才會執行，如果 \texttt{try}
裡出錯了，Python 會跳去執行 \texttt{except}，\texttt{else} 根本不會被碰到。
\index{else（搭配 try/except）}把「依賴
\texttt{try} 成功才該做的事」放進 \texttt{else}，而不是直接接著寫在 \texttt{try} 本體裡，
是有意義的：\texttt{try} 的範圍應該盡量只包住「真正可能失敗、你想保護的那一小段」，範圍抓
太大，之後如果在 \texttt{try} 裡不小心寫錯了別的東西，程式一樣會被 \texttt{except} 攔下
來，卻讓人誤以為問題出在檔案讀取上，反而更難抓出真正的問題。

\begin{example}{python}{(count-characters-in-file.py) 讀取檔案並統計字數}
filename = "shopping_list.txt"

try:
    with open(filename) as file:
        content = file.read()
except FileNotFoundError:
    print("Could not find this file, so character count is unavailable")
else:
    character_count = len(content.replace("\n", ""))
    print(f"This file has {character_count} characters")
\end{example}

\iosample{}{This file has 18 characters \\}

\keypoint{\texttt{try} 區塊要盡量「瘦」，只包住真正可能失敗的那幾行；任何「依賴這件事成功才
該做」的後續處理，應該放進 \texttt{else}，不要直接寫在 \texttt{try} 本體裡。這樣未來如果那
段後續程式碼本身出了 bug，才能一眼看出問題出在統計字數這一段，跟檔案讀取本身無關，不會被
\texttt{except} 的訊息誤導。}

\begin{drills}[label=drills:try-except]
\item 把上面 \texttt{count-characters-in-file.py} 的 \texttt{filename} 換成
\texttt{"names.txt"}（內容是第~\ref{sec:reading-lines}~節用過的三個名字），這段程式的輸
出會是什麼？
\item 如果把 \texttt{filename} 換成一個不存在的檔案，例如 \texttt{"todo.txt"}，輸出會是
什麼？
\item 假設有人把統計字數的邏輯改成這樣，放進了 \texttt{try} 本體，而不是 \texttt{else}：
\begin{verbatim}
try:
    with open(filename) as file:
        content = file.read()
    number = int(content)
    print(f"The number in the file is {number}")
except FileNotFoundError:
    print("Could not find this file")
\end{verbatim}
如果 \texttt{filename} 是 \texttt{"shopping\_list.txt"}（內容是英文單字組成的購物清單，
不是數字），這段程式執行後會發生什麼事？\texttt{except FileNotFoundError} 接得住這個問題嗎？
\end{drills}

\section{同時處理多個 except：分別處理不同的失敗原因}
\label{sec:multiple-except}

上一節練習的第 3 題已經先讓你踩到這個坑：\texttt{try} 區塊裡如果做了不只一件可能失敗的事
（先開檔案、再把內容轉成數字），只寫一個 \texttt{except FileNotFoundError} 只能接住「檔案
不存在」這一種問題，內容轉數字失敗丟出的 \texttt{ValueError} 完全不會被接住，程式照樣當
掉。這一節教你怎麼一次防範好幾種不同的失敗原因。

寫法很像第 \ref{ch:conditionals} 章的 \texttt{if/elif} 鏈：一個 \texttt{try} 後面可以接好幾個 \texttt{except}，
Python 由上往下依序比對，符合哪一個就執行哪一段，其餘全部跳過。

\begin{code}{python}{try 後面接好幾個 except}
try:
    operation_that_might_fail()
except ExceptionTypeA:
    handle_a_failure()
except ExceptionTypeB:
    handle_b_failure()
else:
    do_something_on_success()
\end{code}

Python 從第一個 \texttt{except} 開始比對，實際發生的例外類型只要符合其中一個，就執行那一
段，其餘的 \texttt{except} 全部跳過不看；如果一個都不符合，例外會繼續往外傳，程式跟完全沒寫
\texttt{except} 時一樣，帶著 traceback 中止。

假設有一份 \texttt{score.txt}，內容只有一個數字：

\begin{verbatim}
95
\end{verbatim}

另外還有一份格式打錯的 \texttt{bad\_score.txt}：

\begin{verbatim}
ninety-five points
\end{verbatim}

\begin{example}{python}{(read-score-multi-except.py) 讀取分數檔案，分別處理兩種可能的失敗}
def read_score(filename):
    try:
        with open(filename) as score_file:
            content = score_file.read()
        score = int(content.strip())
    except FileNotFoundError:
        print(f"{filename} does not exist")
    except ValueError:
        print(f"The content in {filename} is not a valid number")
    else:
        print(f"The score read is {score}")

read_score("score.txt")
read_score("bad_score.txt")
read_score("missing.txt")
\end{example}

\iosample{}{The score read is 95 \\ The content in bad\_score.txt is not a valid number \\ missing.txt does not exist \\}

三次呼叫分別踩中三種情況：\texttt{score.txt} 內容是合法數字，\texttt{try} 裡兩行都順利，
掉進 \texttt{else}；\texttt{bad\_score.txt} 內容是文字不是數字，\texttt{int()} 轉換失敗，丟出
\texttt{ValueError}，被第二個 \texttt{except} 接住；\texttt{missing.txt} 根本不存在，
\texttt{open()} 這一步就丟出 \texttt{FileNotFoundError}，被第一個 \texttt{except} 接
住，後面 \texttt{int()} 轉換那一行根本不會被執行到。

\keypoint{\texttt{except} 的比對順序由上往下，跟 \texttt{if/elif} 鏈一樣，只要前面一個
\texttt{except} 已經接住了，後面的 \texttt{except} 就不會再被檢查。同一個 \texttt{try}
裡如果有好幾個可能出錯的地方，把每一種例外類型各自寫成一個獨立的 \texttt{except}，可以讓
你針對不同的失敗原因給出不同的處理方式，不用把它們全部混在一起、只給一種籠統的錯誤訊息。}

這也代表 \texttt{except} 的擺放順序真的有影響，跟 \texttt{if/elif} 鏈的道理完全一樣。Python
提供了一個範圍非常廣的例外類型，叫 \texttt{Exception}，幾乎所有你在這本書裡遇到過的例外類
型（\texttt{FileNotFoundError}、\texttt{ValueError}、\texttt{KeyError} 之類）都算是它的一
種。
\index{Exception}如果不小心把 \texttt{except Exception:} 這種範圍很廣的寫法擺在前面，後面比較具體的
\texttt{except} 就永遠沒有機會被執行到：

\begin{example}{python}{(except-order-mistake.py) 常見錯誤：把範圍廣的例外類型寫在前面}
def read_score(filename):
    try:
        with open(filename) as score_file:
            content = score_file.read()
        score = int(content.strip())
    except Exception:
        print("An error occurred")
    except FileNotFoundError:
        print(f"{filename} does not exist")

read_score("missing.txt")
\end{example}

\iosample{}{An error occurred \\}

\texttt{missing.txt} 真正丟出的是 \texttt{FileNotFoundError}，但因為 \texttt{except
Exception:} 排在前面，Python 由上往下比對，第一個就符合了，直接執行「An error occurred」，後
面那個寫得更精確的 \texttt{except FileNotFoundError} 永遠輪不到它上場，寫了也是白寫。
\texttt{except} 的擺放順序，應該是範圍愈小、愈明確的類型放愈前面，範圍最廣的放最後（如果真
的需要它的話）。

\begin{drills}[label=drills:multiple-except]
\item 如果把 \texttt{read-score-multi-except.py} 裡兩個 \texttt{except} 的順序對調（先寫
\texttt{except ValueError}，再寫 \texttt{except FileNotFoundError}），呼叫
\texttt{read\_score("missing.txt")} 的結果會不會改變？為什麼？
\item 如果只留下 \texttt{except FileNotFoundError}，刪掉 \texttt{except ValueError}
那一段，呼叫 \texttt{read\_score("bad\_score.txt")} 會發生什麼事？
\item 如果把 \texttt{except-order-mistake.py} 裡兩個 \texttt{except} 的順序對調
（\texttt{FileNotFoundError} 寫在前面，\texttt{Exception} 寫在後面），呼叫
\texttt{read\_score("missing.txt")} 的輸出會變成什麼？
\end{drills}

\section{except 抓多種例外類型：用 tuple 一次接住好幾種}
\label{sec:except-tuple}

上一節兩種例外分開處理，各自印出不同訊息。但有時候你根本不在乎是哪一種原因失敗，只要「反正
失敗了」就做同一件事，這時候每種例外分別寫一個 \texttt{except}、卻塞進一模一樣的處理內容，
就顯得囉唆。Python 讓你把好幾種例外類型包成一個 tuple（第 \ref{ch:tuples} 章教過，好幾個值用逗號隔開、通
常用括號包起來），放進同一個 \texttt{except}。
\index{except 的 tuple 語法}

\begin{code}{python}{except 用 tuple 一次接住好幾種例外}
try:
    operation_that_might_fail()
except (ExceptionTypeA, ExceptionTypeB):
    handle_either_failure()
\end{code}

括號裡用逗號列出好幾種例外類型，只要實際發生的例外符合其中任何一種，這個 \texttt{except}
就會被觸發，處理內容完全共用同一段。

\begin{example}{python}{(read-score-tuple-except.py) 用同一段處理接住兩種可能的失敗}
def read_score(filename):
    try:
        with open(filename) as score_file:
            content = score_file.read()
        score = int(content.strip())
    except (FileNotFoundError, ValueError):
        print(f"Could not read a valid score from {filename}")
    else:
        print(f"The score read is {score}")

read_score("score.txt")
read_score("bad_score.txt")
read_score("missing.txt")
\end{example}

\iosample{}{The score read is 95 \\ Could not read a valid score from bad\_score.txt \\ Could not read a valid score from missing.txt \\}

這段程式跟上一節做的事情很像，差別在於：上一節兩種失敗給出不同訊息（分別知道是「檔案不存
在」還是「內容不是數字」），這一節不在乎是哪一種原因，一律印出同一句「沒辦法讀到有效的分
數」。要選哪一種寫法，取決於呼叫端之後拿到這個結果，需不需要分辨失敗的原因：

\begin{itemize}
\item 如果不同的失敗原因之後會有不同的後續動作（例如「檔案不存在」要建議使用者檢查路徑，
「內容不是數字」要建議使用者檢查檔案格式），分開寫成兩個 \texttt{except} 比較好。
\item 如果不管哪種原因、後續要做的事情都一樣，tuple 版本比較精簡，不用寫兩段一模一樣的處理
內容。
\end{itemize}

寫 tuple 版本的 \texttt{except} 時，括號不能省。忘記加括號是這裡最容易踩到的坑：

\begin{verbatim}
except FileNotFoundError, ValueError:
\end{verbatim}

這樣寫，Python 直接在讀進這行程式碼的當下就報錯：

\begin{verbatim}
  File "read-score-tuple-except.py", line 5
    except FileNotFoundError, ValueError:
SyntaxError: multiple exception types must be parenthesized
\end{verbatim}

這是 \texttt{SyntaxError}（第 \ref{ch:introduction} 章認識過的「程式還沒開始執行、光讀就卡住」的錯誤類型），因
為 Python 看到逗號、卻沒看到括號，會誤以為你想用一種很舊、現在已經不能用的寫法（把捕捉到的
例外物件存進一個變數）。括號不能省。

\keypoint{\texttt{except (TypeA, TypeB):} 括號裡的東西是一個 tuple，只要實際發生的例外符合
裡面任何一種類型，就會被這個 \texttt{except} 接住。這跟寫兩個各自獨立的 \texttt{except}
效果不完全一樣：tuple 版本只能對這幾種例外做「完全相同」的處理；如果不同例外類型需要不同的
處理內容，還是要分開寫成好幾個 \texttt{except}。}

\begin{drills}[label=drills:except-tuple]
\item 把 \texttt{except (FileNotFoundError, ValueError):} 改寫成兩個各自獨立的
\texttt{except}，但兩段都印出一模一樣的「Could not read a valid score」，這樣寫法跟 tuple 版本比起
來，你覺得哪一種比較適合這個情境？為什麼？
\item 如果把 tuple 的括號拿掉，寫成 \texttt{except FileNotFoundError, ValueError:}，你覺
得會發生什麼事？動手試試看，跟你的猜測一不一樣？
\end{drills}

\section{finally：不論成功或失敗都要做的事}
\label{sec:finally}

\texttt{try/except}（搭配 \texttt{else}）已經處理了「成功」跟「失敗」兩條路。但有些事情，
不管成功還是失敗，都想要它一定被執行，例如留一筆「這次操作結束了」的紀錄。這就是
\texttt{finally} 要解決的事。

\begin{code}{python}{try/except/finally}
try:
    operation_that_might_fail()
except ExceptionType:
    handle_failure()
finally:
    always_do_this()
\end{code}

\texttt{finally} 區塊不管 \texttt{try} 裡有沒有出錯、有沒有被 \texttt{except} 接住，永遠都
會執行。
\index{finally}就算 \texttt{except} 裡面本身又出了另一個錯，\texttt{finally} 也會在程式真正結束之
前先跑一次。

\begin{example}{python}{(read-file-with-finally.py) 不論成功或失敗都留一筆紀錄}
filename = "shopping_list.txt"

try:
    with open(filename) as file:
        content = file.read()
    print(content)
except FileNotFoundError:
    print("This file was not found")
finally:
    print("This read operation has finished")
\end{example}

\iosample{}{Milk \\ Eggs \\ Toast \\ Apple \\~\\ This read operation has finished \\}

如果把 \texttt{filename} 換成一個不存在的檔案，例如 \texttt{"todo.txt"}，結果會變成：

\iosample{}{This file was not found \\ This read operation has finished \\}

不管是哪一種結果，「This read operation has finished」這一行永遠都會印出來，這就是 \texttt{finally} 的意
義：它不在乎 \texttt{try} 裡發生了什麼事，只在乎「這個收尾動作一定要做」。

檔案的關閉這件事，第~\ref{sec:reading-files}~節教的 \texttt{with} 其實已經自動幫你處理掉
了。\texttt{with} 背後的機制，精神上跟 \texttt{finally} 很像：不管區塊裡的程式順不順利，該
做的收尾動作都會做。所以用 \texttt{with} 開檔案時，你不會常常需要自己額外再寫一個
\texttt{finally} 去關檔案。但接下來你會看到，\texttt{finally} 真正常用的場合，是那些沒有用
\texttt{with}、需要自己手動管理收尾動作的情境。

什麼時候你會真的自己動手寫 \texttt{finally}？最常見的情境是：你手動呼叫 \texttt{open()}
（不透過 \texttt{with}），因為某個理由需要在同一個變數上，貫穿好幾段程式碼一直用同一個檔案
物件，這時候如果你不確定中間會不會出錯，一定要把「關閉檔案」這個動作放進 \texttt{finally}，
才能保證不管中間出不出錯，檔案最後都會被正確關閉。

\begin{example}{python}{(write-log-manually.py) 手動管理檔案開關，用 finally 保證一定會關閉}
log_file = open("log.txt", "a")

try:
    log_file.write("Starting to process the shopping list\n")
    with open("shopping_list.txt") as source_file:
        content = source_file.read()
    log_file.write("Read successful\n")
finally:
    log_file.close()

print("Log written to log.txt")
\end{example}

\iosample{}{Log written to log.txt \\}

\texttt{log\_file = open("log.txt", "a")} 這一行沒有搭配 \texttt{with}，所以 Python 不會
自動幫你關閉它；中間的 \texttt{try} 區塊如果任何一行出錯（例如 \texttt{shopping\_list.txt}
真的不存在），Python 會立刻跳出 \texttt{try}，但在程式真正往外傳遞這個錯誤之前，
\texttt{finally} 裡的 \texttt{log\_file.close()} 一定會先執行一次，\texttt{log\_file} 不
會被遺忘在「打開卻沒關閉」的狀態。

這跟第~\ref{sec:reading-files}~節提過的「\texttt{with}
自動關檔案」做的是同一件事，只是這裡改成自己手動做，因為 \texttt{log\_file} 這個檔案物件要
橫跨好幾行程式碼繼續使用，沒辦法只用一段小小的 \texttt{with} 縮排區塊包住它。沒有明確關閉的
檔案，占用的資源不會馬上釋放，長時間執行、開很多檔案的程式可能因此出問題，這也是為什麼前面
幾節都建議優先用 \texttt{with}，只有像這裡「檔案要橫跨很長一段程式碼」這種不方便用
\texttt{with} 的情境，才手動 \texttt{open()} 加 \texttt{finally} 管理。

\texttt{finally} 保證一定會執行，這件事比很多人以為的還要徹底：就算 \texttt{try} 裡發生的
錯誤，根本沒有被任何一個 \texttt{except} 接住，\texttt{finally} 還是會先執行一次，才讓這個
錯誤繼續往外傳、讓程式中止。

\begin{example}{python}{(convert-score-finally.py) finally 在例外沒被接住時依然先執行}
def convert_to_number(text):
    try:
        number = int(text)
        print(f"The converted number is {number}")
    except FileNotFoundError:
        print("File not found")
    finally:
        print("Conversion process finished")

convert_to_number("ninety-five")
\end{example}

執行後會先印出「Conversion process finished」，接著才是 Python 印出的 traceback，程式最終還是當掉：

\begin{verbatim}
Conversion process finished
Traceback (most recent call last):
  File "convert-score-finally.py", line 9, in <module>
    convert_to_number("ninety-five")
  File "convert-score-finally.py", line 3, in convert_to_number
    number = int(text)
ValueError: invalid literal for int() with base 10: 'ninety-five'
\end{verbatim}

\texttt{text} 傳入的是「ninety-five」，\texttt{int()} 沒辦法把它轉成數字，丟出的是
\texttt{ValueError}；但這裡的 \texttt{except} 只接 \texttt{FileNotFoundError}，接不住
\texttt{ValueError}，所以錯誤繼續往外傳。即使如此，\texttt{finally} 裡的
\texttt{print("Conversion process finished")} 依然先執行了一次，這一行印出來之後，Python 才真正把這個沒被
接住的錯誤往外送，最終讓整支程式中止並印出 traceback。這正是 \texttt{finally}「不論成功或
失敗」裡最徹底的那一種情況：連例外都沒被接住，\texttt{finally} 照樣先跑。

\keypoint{\texttt{finally} 保證一定執行的範圍，比 \texttt{except} 涵蓋的範圍更大：不管
\texttt{try} 有沒有出錯、出的錯有沒有被某個 \texttt{except} 接住，\texttt{finally} 都會先
執行過一次，才讓程式繼續往下（沒出錯，或錯誤被接住的情況），或是讓錯誤繼續往外傳、程式中止
（錯誤沒被接住的情況）。}

\begin{drills}[label=drills:finally]
\item 把 \texttt{convert-score-finally.py} 的 \texttt{except FileNotFoundError:} 改成
\texttt{except ValueError:}，呼叫 \texttt{convert\_to\_number("ninety-five")} 這次還會不會印
出 traceback？「Conversion process finished」還會不會印出來？
\item 如果 \texttt{write-log-manually.py} 裡的 \texttt{shopping\_list.txt} 其實不存在，這
段程式最後 \texttt{log\_file} 會不會被正確關閉？為什麼？
\end{drills}

\section{自己丟出例外：\texttt{raise}}
\label{sec:raise}

到目前為止，例外都是 Python 幫你丟的：你打開一個不存在的檔案，它丟
\texttt{FileNotFoundError}；你把一段不是數字的文字餵給 \texttt{int()}，它丟
\texttt{ValueError}。但有些「不合理」是 Python 看不出來的——它不知道你的程式裡年齡不該是
負數、成績不該超過一百分。 這種只有你才知道的規則，你可以自己丟一個例外出來。

\begin{code}{python}{自己丟出一個例外}
raise ExceptionType("說明這次為什麼出錯的訊息")
\end{code}

\begin{example}{python}{(set-age.py) 遇到不合理的值就丟出例外}
def set_age(age):
    if age < 0:
        raise ValueError("age cannot be negative")
    print(f"age is set to {age}")


set_age(20)
set_age(-3)
\end{example}

\begin{verbatim}
age is set to 20
Traceback (most recent call last):
  File "set-age.py", line 8, in <module>
    set_age(-3)
    ~~~~~~~^^^^
  File "set-age.py", line 3, in set_age
    raise ValueError("age cannot be negative")
ValueError: age cannot be negative
\end{verbatim}

第一次呼叫正常印出訊息；第二次傳進 \texttt{-3}，\texttt{raise} 這一行執行之後，函式立刻
停下來，程式跟著中斷，畫面上出現的 traceback 跟 Python 自己丟的例外長得一模一樣——因為
它本來就是同一種東西。 你選用的 \texttt{ValueError} 也不是新發明的，是第
\ref{sec:multiple-except} 節就見過的那一個，意思是「型別對，但值不合理」。
\index{raise}

\keypoint{\texttt{raise} 是「主動宣告這裡出錯了」，跟 \texttt{return} 一樣會讓函式立刻
結束，差別在 \texttt{return} 是正常交差、\texttt{raise} 是把問題往外丟。 括號裡的訊息會
原封不動出現在 traceback 最後一行，所以要寫得夠具體——寫
\texttt{"age cannot be negative"} 比寫 \texttt{"error"} 有用得多。}

自己丟出來的例外，一樣可以被 \texttt{try}／\texttt{except} 接住，規則跟前面幾節完全相同。
這才是 \texttt{raise} 真正的價值：函式負責「發現問題並老實說出來」，呼叫端負責「決定遇到
問題要怎麼辦」，兩件事分開。

\begin{example}{python}{(validate-ages.py) 自己丟的例外，自己接住}
def set_age(age):
    if age < 0:
        raise ValueError("age cannot be negative")
    return age


for value in [20, -3, 45]:
    try:
        result = set_age(value)
    except ValueError as error:
        print("skipped:", error)
    else:
        print("accepted:", result)
\end{example}

\iosample{}{accepted: 20 \\ skipped: age cannot be negative \\ accepted: 45 \\}

\texttt{except ValueError as error} 這個寫法你在第 \ref{sec:multiple-except} 節看過，
\texttt{error} 接住的就是被丟出來的那個例外物件，直接 \texttt{print} 它會印出當初寫在括號
裡的那句訊息。 整個迴圈沒有因為中間那筆壞資料而中斷，這正是第 \ref{sec:try-except-else}
節那套「一筆失敗不該讓整批當機」的想法，只是這次連「失敗」的定義都是你自己訂的。

\stylepoint{什麼時候該 \texttt{raise}、什麼時候該回傳一個特別的值（例如 \texttt{None}）？
判準是「呼叫端如果忘記處理，出事比較好還是默默錯下去比較好」。 年齡是負數這種明顯不該
發生的事，丟例外比較安全——忘記處理會當場炸給你看；如果默默回傳 \texttt{None}，錯誤會
一路飄到很後面才爆，那時候就很難查了。 這正是這一章 EAFP 精神的另一面。}

\begin{drills}[label=drills:raise]
\item \texttt{set\_age(0)} 會不會丟出例外？為什麼？
\item 如果把 \texttt{validate-ages.py} 的 \texttt{except ValueError} 改成
\texttt{except TypeError}，跑起來會發生什麼事？
\item 你要寫一個 \texttt{set\_score(score)}，規定分數必須介於 0 到 100 之間。 判斷條件
該怎麼寫？
\end{drills}

\section{回頭看：EAFP 這套邏輯有多大}
\label{sec:eafp-recap}

這一節帶你回頭看幾個之前已經用過、但沒有完全解釋原因的地方，沒有新語法要學。你會看到，EAFP
這套邏輯，遠比「檔案讀寫」這一件事還要大。

還記得第 \ref{ch:dict-set} 章你學過用 dict 存資料、用 \texttt{[]} 查值嗎？如果想查一個可能不存在的 key，其
實可以用今天學到的方式重新想這個問題：

\begin{example}{python}{(dict-lookup-lbyl.py) 查電話簿：先檢查 key 在不在}
phonebook = {"Alice": "0912-345-678", "Bob": "0933-222-111"}

if "Taylor" in phonebook:
    number = phonebook["Taylor"]
else:
    number = "Not found"

print(number)
\end{example}

\iosample{}{Not found \\}

\begin{example}{python}{(dict-lookup-eafp.py) 查電話簿：直接查，查不到再處理}
phonebook = {"Alice": "0912-345-678", "Bob": "0933-222-111"}

try:
    number = phonebook["Taylor"]
except KeyError:
    number = "Not found"

print(number)
\end{example}

\iosample{}{Not found \\}

第 \ref{ch:dict-set} 章其實已經說過：與其每次查之前都先問「這個線索存在嗎？」，不如直接假設它存在，用
\texttt{get()} 順便準備好一個備用答案，現在你知道，這句話講的正是 EAFP。

\begin{example}{python}{(dict-lookup-get.py) 查電話簿：用 get() 一行搞定}
phonebook = {"Alice": "0912-345-678", "Bob": "0933-222-111"}

number = phonebook.get("Taylor", "Not found")
print(number)
\end{example}

\iosample{}{Not found \\}

三段程式做的事完全一樣、輸出也完全一樣，但份量差很多，這正是 EAFP 這套邏輯的價值所在：
Python 常常會把「先做、錯了再處理」這件事包裝成更精簡的寫法，\texttt{get()} 就是一個很好的
例子，它就是 EAFP 精神的另一種展現，只是 Python 把「查不到就用預設值」這個常見情境包成一個
現成的方法，連 \texttt{try/except} 都不用自己寫。

再往前翻，第 \ref{ch:variables} 章教型別轉換時，你已經看過 \texttt{int("hello")} 會丟出 \texttt{ValueError}，
那時候還沒教 \texttt{try/except}，只能眼睜睜看著程式當掉，沒有別的辦法。現在你手上有完整的
工具了，可以用同一套 EAFP 邏輯，把「使用者亂打的輸入」這種最貼近日常的情境也接住。

\begin{example}{python}{(ask-age-eafp.py) 接住使用者打錯格式的輸入}
text = input("Please enter your age: ")

try:
    age = int(text)
except ValueError:
    print("That is not a valid number, treating it as blank")
    age = None

if age is not None:
    print(f"Your age is {age}")
\end{example}

輸入 \texttt{17}：

\iosample{17 \\}{Your age is 17 \\}

輸入 \texttt{seventeen}：

\iosample{seventeen \\}{That is not a valid number, treating it as blank \\}

\texttt{text} 是 \texttt{input()} 讀進來的字串，只要不是純數字，\texttt{int()} 就會丟出
\texttt{ValueError}；\texttt{except} 接住它、印出提醒，順便把 \texttt{age} 設成
\texttt{None}（第 \ref{ch:lists} 章認識過的「沒有值」的特殊值）；最後用 \texttt{if age is not None}
判斷輸入到底有沒有成功轉換。這跟 LBYL 版本（自己寫一堆字串方法檢查「是不是每個字元都是數
字」）比起來，程式碼更短，也不用自己重新發明 \texttt{int()} 內部早就做好的判斷邏輯。

再回頭看第 \ref{ch:conditionals} 章。你當時學到一件事：Python 的 \texttt{if} 不只認 \texttt{True}／
\texttt{False}，連數字、字串本身都可以直接放進 \texttt{if} 裡問「是不是有內容」，\texttt{0}
跟空字串會被當成 \texttt{False}，非零、非空會被當成 \texttt{True}，不用特別寫
\texttt{len(x) > 0}。這個「空的當 \texttt{False}」的邏輯，對 list、dict 這些容器一樣成立：

\begin{example}{python}{(check-shopping-cart.py) 購物車空了沒}
shopping_cart = []

if shopping_cart:
    print("The cart has items, starting checkout")
else:
    print("The cart is empty")
\end{example}

\iosample{}{The cart is empty \\}

這跟 EAFP 不是同一個語法機制：\texttt{if shopping\_cart:} 用的不是 \texttt{try/except}，
但背後在意的事情是一樣的：不要先小題大作地驗證狀態（例如特地寫
\texttt{len(shopping\_cart) > 0} 確認「長度大於 0」），讓語言本身的行為直接替你簡化判斷。
EAFP 是「先做，出錯了 Python 會告訴你」；空的當 \texttt{False} 則是「Python 自己就把『有沒
有內容』這件事，直接兌換成一個你能拿來判斷的答案」。兩者機制不同，方向卻一致：不要繞遠路先
確認，讓語言本身的行為替你簡化判斷。

\stylepoint{從這一章往後，只要腦中冒出「要不要先檢查一下」的念頭，查字典前、開檔案前，甚至
任何「這個東西有沒有內容」的判斷，都可以先反問自己一句：Python 有沒有更直接的寫法，讓程式
直接做、讓它自己告訴你結果？\texttt{try/except}、\texttt{dict.get()}、\texttt{if items:}
都是同一種精神的不同展現。}

這一章教會你 Python 面對「事情可能會出錯」的兩種態度，也讓你看清楚 Python 比較偏好哪一種、
為什麼，還讓你學會怎麼同時防範好幾種不同的失敗原因、怎麼保證收尾動作一定會被執行。之後每一
章只要遇到可能失敗的操作，第 \ref{ch:debugging} 章讀懂錯誤訊息、以後接觸資料庫或網路連線，都會回頭用到這裡
打下的地基。

\begin{drills}[label=drills:eafp-recap]
\item 下面這段程式是 LBYL 還是 EAFP？為什麼？
\begin{verbatim}
import os

if os.path.exists("data.txt"):
    with open("data.txt") as f:
        print(f.read())
\end{verbatim}
\item 把下面這段 LBYL 風格的字典查詢，改寫成用 \texttt{try/except} 的 EAFP 風格：
\begin{verbatim}
scores = {"Alice": 90, "Bob": 85}
if "Taylor" in scores:
    score = scores["Taylor"]
else:
    score = 0
\end{verbatim}
\item 上一題的程式，如果改用 \texttt{get()} 寫，一行可以怎麼寫？
\item 想像你要把 \texttt{ask-age-eafp.py} 改寫成 LBYL 風格：不用 \texttt{try/except}，改
成先想辦法檢查這個字串「看起來像不像數字」，才呼叫 \texttt{int()}。你覺得自己寫的檢查邏
輯，真的能完全覆蓋 \texttt{int()} 內建判斷的所有情況嗎（提示：想想看負數、前後夾帶空白這些
情況）？
\end{drills}

\section*{練習解答}
\addcontentsline{toc}{section}{練習解答}

\paragraph{第~\ref{sec:reading-lines}~節}
\begin{enumerate}
\item 會印出：
\begin{verbatim}
Orange

Banana

\end{verbatim}
每一行結尾藏著的換行符號，加上 \texttt{print()} 自己又加一次換行，兩個疊在一起，讓每一行
後面都多出一行空白。
\item 不會再有多餘的空行，「Orange」直接接著「Banana」印出來，因為 \texttt{rstrip()} 把藏在
\texttt{line} 結尾的換行符號去掉了。
\item 開放性問題，兩種都可以：用 \texttt{readlines()} 拿到 list 之後直接 \texttt{len()}，
程式碼短，但要先把整份檔案讀進記憶體；用 \texttt{for line in file\_object} 自己數計數器要多寫
兩三行，但不用一次把整份檔案存進 list，檔案很大時比較省記憶體。這裡沒有標準答案，重點是能
講出兩種寫法各自的取捨。
\end{enumerate}

\paragraph{第~\ref{sec:writing-styles}~節}
\begin{enumerate}
\item \texttt{cart.txt} 會變成「MilkEggsToast」黏成一坨、沒有分行，因為 \texttt{writelines()}
不會自動幫每個項目加換行，list 裡的字串結尾沒有 \texttt{\textbackslash n}，就完全不會有分
行效果。
\item 開放性問題：資料量很大時，通常會偏向「先組好一個大字串，一次 \texttt{write()}」，因
為只跟檔案系統打一次交道，額外開銷比較小；但要注意這代表要先把一百萬行全部組成一個字串放進
記憶體，如果這份字串大到記憶體放不下，就得改回「迴圈裡逐項 \texttt{write()}」，用比較多次
的檔案操作換取比較低的記憶體用量。
\end{enumerate}

\paragraph{第~\ref{sec:appending-files}~節}
\begin{enumerate}
\item \texttt{'a'}，附加模式才不會把之前的紀錄洗掉。
\item \texttt{'w'}，寫入模式每次都會先清空，確保是全新的一份。
\item 會自動新建一個空檔案，再開始寫入。這跟 \texttt{'w'} 模式的行為相同：兩者在「檔案原本
不存在」這件事上行為一樣，差別只在「檔案已經存在」時，\texttt{'w'} 清空重寫，\texttt{'a'}
接著寫。
\end{enumerate}

\paragraph{第~\ref{sec:multiple-files}~節}
\begin{enumerate}
\item \texttt{shopping\_list.txt} 會變成空檔案，什麼內容都沒有。因為
\texttt{open("shopping\_list.txt", "w")} 一執行就先把這個檔案清空，等 \texttt{for line
in source\_file} 想讀取內容時，能讀到的只剩下空的檔案。
\item 應該分成兩段：先用一個 \texttt{with open("shopping\_list.txt")} 把全部內容讀出來、
處理好，存進一個變數（例如一份組好的 list 或字串），確定第一個 \texttt{with} 區塊已經結束、
確定讀取完成，再另外用 \texttt{with open("shopping\_list.txt", "w")} 把處理好的結果寫回
去，不能讓兩個 \texttt{open()} 同時對著同一個檔名。
\end{enumerate}

\paragraph{第~\ref{sec:try-except-else}~節}
\begin{enumerate}
\item 會印出：This file has 14 characters（Alice、Bob、Taylor 三個名字接在一起，拿掉換行
符號後總共 14 個字元）。
\item 會印出：Could not find this file, so character count is unavailable。
\item 程式會當掉，出現未被接住的 \texttt{ValueError}。\texttt{int(content)} 想把購
物清單裡的文字內容轉成整數，這件事本身跟「檔案找不到」無關，\texttt{except FileNotFoundError} 只
接 \texttt{FileNotFoundError} 這一種例外類型，不會接住 \texttt{ValueError}，程式會帶著一
段 \texttt{ValueError} 的 traceback 直接中止。這正是把不相關的邏輯塞進 \texttt{try} 本
體、不放進 \texttt{else} 的風險：一個跟檔案讀取無關的 bug，一樣會讓程式看起來像是「檔案操
作」出了問題。
\end{enumerate}

\paragraph{第~\ref{sec:multiple-except}~節}
\begin{enumerate}
\item 結果不會改變。\texttt{missing.txt} 不存在，\texttt{open()} 丟出的是
\texttt{FileNotFoundError}，不管這個 \texttt{except} 寫在第一個還是第二個位置，只要
Python 由上往下比對到符合的類型，就會執行對應的區塊；順序只在「同一個例外可能同時符合好幾
個 \texttt{except}」時才會影響結果，這裡兩種例外類型完全不重疊，所以順序對調不影響輸出。
\item 程式會當掉。\texttt{bad\_score.txt} 讓 \texttt{int()} 丟出 \texttt{ValueError}，但
剩下的 \texttt{except} 只接 \texttt{FileNotFoundError}，接不住 \texttt{ValueError}，程式
會帶著 \texttt{ValueError} 的 traceback 直接中止。
\item 輸出會變成「missing.txt does not exist」。因為 \texttt{except FileNotFoundError}
現在排在前面，Python 由上往下比對，先符合這個比較具體的類型，就不會再往下看到
\texttt{except Exception}，這正是「範圍小的類型放前面」的正確順序。
\end{enumerate}

\paragraph{第~\ref{sec:except-tuple}~節}
\begin{enumerate}
\item 開放性問題：這個情境裡兩種失敗的後續處理完全一樣（都只是印出同一句提醒），用 tuple
版本比較精簡，不用寫兩段重複的內容；如果之後兩種失敗需要不同的後續動作，才會傾向拆成兩個獨
立的 \texttt{except}。
\item Python 在讀進這行程式碼的當下就會報 \texttt{SyntaxError}：
\texttt{multiple exception types must be parenthesized}，因為沒有括號的
\texttt{except Type, name:} 是 Python 舊版一種已經不能用的寫法，Python 看到逗號會誤以為
你想用那種寫法，而不是列出好幾種例外類型。
\end{enumerate}

\paragraph{第~\ref{sec:finally}~節}
\begin{enumerate}
\item 這次不會印出 traceback，因為 \texttt{except ValueError:} 接得住 \texttt{int("ninety-five")}
丟出的 \texttt{ValueError}。輸出會先印「File not found」（訊息文字雖然跟實際情況不符，但
這一段程式本來就只是示範捕捉行為，不代表訊息內容一定正確），接著印「Conversion process
finished」，程式正常結束。
\item 會被正確關閉。不管 \texttt{shopping\_list.txt} 存不存在，\texttt{try} 裡任何一行出
錯，Python 都會先執行 \texttt{finally} 裡的 \texttt{log\_file.close()}，才讓錯誤繼續往外
傳，這正是手動管理檔案開關時一定要把關閉動作放進 \texttt{finally} 的原因。
\end{enumerate}

\paragraph{第~\ref{sec:raise}~節}
\begin{enumerate}
\item 不會。判斷條件是 \texttt{age < 0}，\texttt{0} 不小於 \texttt{0}，所以不成立，
函式正常往下執行、印出訊息。
\item \texttt{ValueError} 不會被 \texttt{except TypeError} 接住，所以跑到 \texttt{-3}
那一輪時程式直接中斷、印出 traceback，後面的 \texttt{45} 完全不會被處理。
\item \texttt{if score < 0 or score > 100:}，兩個條件任一成立就丟出 \texttt{ValueError}。
\end{enumerate}

\paragraph{第~\ref{sec:eafp-recap}~節}
\begin{enumerate}
\item LBYL，先用 \texttt{os.path.exists()} 檢查檔案在不在，確認過後才真正 \texttt{open()}，
這正是「先檢查、確定安全才做」的邏輯。
\item
\begin{verbatim}
scores = {"Alice": 90, "Bob": 85}
try:
    score = scores["Taylor"]
except KeyError:
    score = 0
\end{verbatim}
\item \texttt{score = scores.get("Taylor", 0)}
\item 開放性問題：很難完全覆蓋。\texttt{int()} 本身能接受負數（\texttt{"-5"}）、正負號、
前後夾帶空白（\texttt{" 5 "}）等好幾種合法格式，自己重新寫一套檢查邏輯很容易漏掉某些情況，
或者太嚴格擋掉了原本合法的輸入。這正是 EAFP 的優勢：直接讓 \texttt{int()} 自己去判斷，不用
自己重新發明一套判斷規則。
\end{enumerate}

\chapter{CSV 與簡單資料檔}
\label{ch:csv}

資料要在不同系統、不同程式之間搬來搬去，最可靠的做法通常是把資料寫成「誰都看得懂、誰都
打得開」的純文字格式。這就像你把一串電話號碼寫在一張紙條上交給朋友：不管對方用什麼手機、
什麼記事習慣，攤開紙條都看得懂內容。但如果你把號碼鎖進一個只有你自己那支手機才能打開的
備忘錄裡，紙條換成別人手上就成了一堆看不懂的東西。

程式的世界也是同一個道理，一份資料如果
存成「純文字、有固定規則」的格式，不管換到哪一台電腦、哪一種程式語言、哪一個系統，都能
被讀懂、被處理。

這一章要學的兩種格式是 CSV（逗號分隔的表格資料）跟 JSON（有層次、有欄位
名稱的資料），它們就是程式世界裡最常見的兩種「誰都讀得懂」的純文字資料格式。
\index{純文字資料格式}

這一章分三段推進：

\begin{itemize}
\item 第一段先不急著學新工具，用你已經會的字串方法跟迴圈，自己動手把一份
逗號分隔的純文字檔案解析成看得懂的資料，順便親身感受一下「自己手動處理資料」在哪裡會卡
關、會出包。
\item 第二段換上 Python 專門為表格式資料準備的 \texttt{csv} 模組，解決第一段手動
處理時遇到的坑，學會用欄位名稱而不是死背第幾欄的方式存取資料，也會碰到幾個實際處理 CSV
檔案時常見的眉角：檔案有沒有標題列、檔案的編碼是不是跟你以為的一樣。
\item 第三段介紹另一種更
有彈性、可以有層次（欄位裡面還能包欄位）的純文字格式 JSON，包括怎麼把巢狀資料直接寫成
一份 JSON 檔案。
\end{itemize}

三段合起來，你會建立起「純文字資料格式」這個概念，知道 CSV 和 JSON
各自適合什麼場合，也會親手體驗一次「手動處理資料」跟「用對工具處理資料」之間的落差，這種
落差感，正是後面認識更專門的資料工具時最實在的參考點。

\section{讀一份逗號分隔的純文字檔案}
\label{sec:csv-manual-read}

想像你手上有一份班級或社團匯出的成績清單，存成一個純文字檔案，每一行是一筆資料，欄位
之間用逗號隔開，這種格式就叫 CSV（comma-separated values，逗號分隔值）。
\index{CSV（comma-separated values）}

要把這種檔案
讀進程式裡，你其實已經具備所有需要的工具了：第 \ref{ch:files-exceptions} 章學過的 \texttt{with open()} 負責把
檔案打開、逐行讀出來；字串的 \texttt{split()} 方法可以把一行文字按逗號拆成一份份欄位。
這一節要做的事，就是把這兩樣東西接在一起，看看一份 CSV 檔案讀進 Python 之後長什麼樣子。

\begin{code}{python}{讀取一份逗號分隔的純文字檔案}
with open("filename.csv") as f:
    for line in f:
        fields = line.strip().split(",")
        # fields 是這一列的欄位清單，用 fields[0]、fields[1]... 依序取出每個欄位
\end{code}

\texttt{line.strip()} 先把這一行頭尾的空白字元去掉，包含結尾看不見的換行符號。這裡刻意
換用 \texttt{strip()}，而不是第 \ref{ch:files-exceptions} 章讀檔案時用的 \texttt{rstrip()}：\texttt{rstrip()}
只清掉字串「尾端」的空白，\texttt{strip()} 則是頭尾兩端都清，因為 CSV 欄位除了結尾的換行
之外，前後也可能夾帶多餘空白。再用 \texttt{split(",")} 把整行文字按逗號切開，變成一份字
串構成的清單。假設有
一份 \texttt{quiz\_scores.csv}，內容如下：

\begin{verbatim}
name,score
Alice,88
Bob,76
Taylor,92
Charlie,65
David,79
\end{verbatim}

第一行是欄位標題，之後每一行是一筆資料。把它讀進來、每一行都拆成欄位清單再印出來看看：

\begin{example}{python}{(read-quiz-scores.py) 讀出一份成績清單}
with open("quiz_scores.csv") as f:
    for line in f:
        fields = line.strip().split(",")
        print(fields)
\end{example}

\iosample{}{['name', 'score'] \\{} ['Alice', '88'] \\{} ['Bob', '76'] \\{} ['Taylor', '92'] \\{} ['Charlie', '65'] \\{} ['David', '79'] \\}

留意一件事：連第一行的欄位標題「name」「score」也被當成一般資料讀進來、印出來了。對
Python 來說，這份檔案裡每一行看起來都一樣，它並不知道第一行其實是「標題」，這件事要靠
你自己的程式邏輯去處理（例如跳過它，或是拿它當欄位名稱用，後面 CSV 模組那一節會示範更
乾淨的做法）。

另外也留意每個欄位都是字串，就算內容看起來像數字（例如 \texttt{'88'}），
拆出來之後仍然是字串，不是整數，這一點在下一節要動手做計算的時候會很重要。

\section{動手寫迴圈：算出總分與平均}
\label{sec:csv-manual-average}

光是把資料讀出來、印出來還不夠，你通常想拿這些資料做點計算，例如算出全班的平均分數。
這裡要用到的技巧，其實你在第 \ref{ch:loops} 章學迴圈的時候就已經練過了：用一個「累加器」變數，在迴圈
裡一次次把值加上去，跑完迴圈之後這個變數就存著總和。差別只在於，這次的數字不是憑空打出
來的，是從檔案裡讀出來、再親手轉型過的。

這一段程式碼務必自己動手重現一次，不只是因為
「累加器迴圈算總和、算平均」是一個到處都用得到的基本功，更因為下一章要認識的資料工具，
一開始就是拿這一段手寫的迴圈當對照組，讓你看看有專門工具之後可以省掉多少麻煩。

\begin{example}{python}{(average-quiz-score.py) 手動算出總分與平均}
with open("quiz_scores.csv") as f:
    lines = f.readlines()

scores = []
for line in lines[1:]:                      # 第一行是標題，不是資料，跳過
    fields = line.strip().split(",")
    score = int(fields[1])                   # 索引 1 是分數欄位，把字串轉成整數
    scores.append(score)

total = 0
for score in scores:
    total += score

average = total / len(scores)

print("Total:", total)
print("Average score:", average)
\end{example}

\iosample{}{Total: 400 \\{} Average score: 80.0 \\}

把這段程式拆開來看三個步驟：

\begin{enumerate}
\item 把每一行（跳過第一行標題）拆出分數欄位、用
\texttt{int()} 轉成整數，收進一份叫 \texttt{scores} 的清單。這一步結束後，你手上有的
是一份「乾淨的數字清單」，不再是字串、也不再夾雜標題行。
\item 也是全章最值得記住的一
段：先在迴圈外把累加器 \texttt{total} 設成 \texttt{0}，再用 \texttt{for} 迴圈把
\texttt{scores} 裡每一個數字依序加進 \texttt{total}，\texttt{total += score} 這一行，
意思是「把 \texttt{total} 現在的值，加上這次迴圈拿到的 \texttt{score}，再存回
\texttt{total}」。
\item 迴圈跑完之後，\texttt{total} 就是所有分數的總和，再除以
\texttt{len(scores)}（清單裡有幾個數字），就是平均值。注意這裡除的是
\texttt{len(scores)}，不是檔案的行數，因為 \texttt{scores} 已經是跳過標題行之後、只
剩下真正資料的乾淨清單。
\end{enumerate}

\keypoint{「\texttt{total = 0} → 迴圈裡 \texttt{total += 這次的值} → 迴圈外
\texttt{total / len(數量)}」這個三段式寫法，是用迴圈算總和、算平均最標準的樣板，幾乎
任何「把一堆數字合併成一個結果」的需求都是同一套邏輯的變形。把這個樣板記熟，比背任何
特定的函式名稱都值得，因為不是每種語言、每個情境都有現成的函式可以用，但這個手寫的邏輯
到哪裡都通。}

\begin{drills}[label=drills:csv-manual-average]
\item 預測輸出：如果把 \texttt{average-quiz-score.py} 裡的 \texttt{scores} 換成
\texttt{[10, 20, 30, 40, 50]}，程式印出的「總分」跟「平均分數」會是多少？
\item 填空：下面這段程式想計算一份數字清單 \texttt{numbers} 的總和與平均，請把三個空格
填上正確的內容：
\begin{verbatim}
total = ______
for n in numbers:
    total ______ n
average = total / ______
\end{verbatim}
\item 抓蟲：如果 \texttt{average-quiz-score.py} 裡忘記寫 \texttt{lines[1:]}，直接對
整個 \texttt{lines} 清單跑迴圈（也就是把標題那一行也當成資料處理），程式執行時會發生
什麼事？
\end{drills}

\section{手動 split 的極限：值裡面有逗號怎麼辦}
\label{sec:csv-manual-limits}

前面兩節手動處理 CSV 都很順利，但那是因為資料剛好乾淨，每個欄位裡都沒有藏著逗號。如果
資料本身就長了一個逗號在裡面呢？比如說一份報名資料裡，「地址」這個欄位寫成「Springfield,
Downtown」，中間就有一個逗號。假設有一份 \texttt{signup.csv}，內容如下：

\begin{verbatim}
name,address,phone
Charlie,"Springfield, Downtown",0912345678
Emma,Riverside East,0922333444
\end{verbatim}

留意第二行。因為地址裡面本來就有逗號，正確的 CSV 格式會在這種欄位外面加一對雙引號
\texttt{"..."}，標記「這整段都算同一個欄位，裡面的逗號不要拆」。如果你還是用上一節那種
最簡單的 \texttt{split(",")}，並不會知道這對雙引號代表什麼特殊意義，只會傻傻地看到逗號
就拆：

\begin{example}{python}{(broken-address-split.py) 地址欄裡的逗號讓資料錯位}
with open("signup.csv") as f:
    for line in f:
        fields = line.strip().split(",")
        print(fields)
\end{example}

\iosample{}{['name', 'address', 'phone'] \\{} ['Charlie', '"Springfield', ' Downtown"', '0912345678'] \\{} ['Emma', 'Riverside East', '0922333444'] \\}

第二筆資料本來應該是 3 個欄位（name、address、phone），結果被拆成了 4 個，而且雙引號還原封
不動地留在字串裡面。如果你以為 \texttt{fields[2]} 是電話號碼，實際拿到的卻是
\texttt{' Downtown"'} 這種殘缺的地址片段，電話號碼反而被擠到 \texttt{fields[3]} 去了。

這種錯位很難用肉眼一眼看穿，尤其資料筆數一多，你根本不會知道哪幾行悄悄錯位了。逗號分隔
的欄位裡本身又能出現逗號、甚至換行，這類「格式裡的特殊規則」要靠手動 \texttt{split()}
自己想清楚每一種狀況，既麻煩又容易漏掉，這正是需要一個專門工具的時候。

\section{用 csv 模組正確讀取資料}
\label{sec:csv-module-read}

Python 標準函式庫內建了一個 \texttt{csv} 模組，專門處理「表格式、逗號分隔」這種資料格
式的所有細節，包括上一節卡住的「欄位裡面有逗號」這種情況。

這種一批批寫好、用
\texttt{import} 叫進來就能直接用的現成程式碼，就叫「模組」，
\index{模組}
第 \ref{ch:files-exceptions} 章 \texttt{import os}
那裡已經先遇過一次，模組系統完整的運作方式要等第 \ref{ch:modules} 章才會介紹；而 Python 本身內建的一
大批這樣的模組，合起來就叫「標準函式庫」，
\index{標準函式庫}
不用另外安裝，直接 \texttt{import} 就能用，
\texttt{csv}、\texttt{json}、\texttt{os} 都屬於這一類。用法跟你已經熟悉的
\texttt{with open()} 差不多，只是多包一層 \texttt{csv.reader()}。
\index{csv 模組／csv.reader()}

\begin{code}{python}{用 csv.reader() 逐列讀取}
import csv

with open("filename.csv") as f:
    reader = csv.reader(f)
    for row in reader:
        # row 是這一列的欄位清單，每個元素都是字串
        ...
\end{code}

\texttt{import csv} 這一行，就是把 \texttt{csv} 這個模組叫進來給你用，跟第 \ref{ch:files-exceptions} 章
\texttt{import os} 的道理一樣。用同一份 \texttt{signup.csv} 試試看：

\begin{example}{python}{(fix-address-split.py) 用 csv 模組正確讀出地址欄}
import csv

with open("signup.csv") as f:
    reader = csv.reader(f)
    for row in reader:
        print(row)
\end{example}

\iosample{}{['name', 'address', 'phone'] \\{} ['Charlie', 'Springfield, Downtown', '0912345678'] \\{} ['Emma', 'Riverside East', '0922333444'] \\}

這次每一列都乖乖地是 3 個欄位，地址欄裡的逗號被完整保留在同一個欄位裡，外層的雙引號也
自動被拿掉了，剛剛手動處理會出包的地方，\texttt{csv} 模組已經幫你想好了。也留意一下
電話號碼那一欄：\texttt{'0912345678'} 開頭的 \texttt{0} 完整保留著，因為它本來就是
字串，不是數字。

\keypoint{CSV 檔案裡讀出來的每一個值，不管看起來多像數字，一律都是字串。\texttt{csv}
模組不會替你自動判斷、自動轉型。這剛好救了電話號碼一命：如果它自動被轉成整數，開頭的
\texttt{0} 就會憑空消失，變成一個少一位數的號碼。要不要轉型、轉成什麼型別，永遠是你自己
決定的事。}

那如果你真的需要把某個欄位轉成數字呢？資料常常不會乾乾淨淨，可能有人那天缺考，分數欄
位是空的。假設有一份 \texttt{scores\_with\_gaps.csv}：

\begin{verbatim}
name,score
Alice,88
Bob,
Taylor,92
\end{verbatim}

Bob 那一行分數欄位是空字串，直接用 \texttt{int()} 轉型會失敗。呼應第 \ref{ch:files-exceptions} 章學過的
EAFP（先做、錯了再處理），與其自己先寫一堆檢查字串是不是「看起來像數字」，不如直接
\texttt{int()} 硬轉，轉失敗了再用 \texttt{try/except} 接住：

\begin{example}{python}{(sum-scores-safely.py) 安全地把分數欄位轉成整數}
import csv

with open("scores_with_gaps.csv") as f:
    reader = csv.reader(f)
    next(reader)                              # 跳過標題列
    for row in reader:
        name, score_text = row
        try:
            score = int(score_text)
            print(name, "score:", score)
        except ValueError:
            print(name, "has no score recorded, skipping")
\end{example}

\iosample{}{Alice score: 88 \\{} Bob has no score recorded, skipping \\{} Taylor score: 92 \\}

這裡也順便介紹一個小技巧：\texttt{next(reader)} 會從 \texttt{reader} 裡拿走「下一列」
資料。放在 \texttt{for} 迴圈之前呼叫一次，效果就是把第一列（也就是標題列）先讀掉、丟在
一邊，之後的 \texttt{for row in reader} 就只會拿到真正的資料列。
\index{next()}

\stylepoint{比起自己先檢查字串「是不是每個字元都是數字」，直接嘗試 \texttt{int()} 轉型、
失敗了再用 \texttt{except ValueError} 接住，是更 Python 的寫法。這正是第 \ref{ch:files-exceptions} 章鋪陳過
的 EAFP：先做、錯了再處理，而不是先確認安全才動手。}

\section{用 csv 模組把資料寫出去}
\label{sec:csv-module-write}

讀資料學會了，接下來換成反方向：把程式裡算出來的結果寫成一份 CSV 檔案，讓它可以被
Excel、Google Sheets，或是任何其他程式打開。用的是 \texttt{csv.writer()}，搭配
\texttt{writerow()} 一次寫入一列。
\index{csv.writer()／writerow()}

\begin{code}{python}{用 csv.writer() 逐列寫入}
import csv

with open("filename.csv", "w", newline="") as f:
    writer = csv.writer(f)
    writer.writerow([field1, field2, ...])
\end{code}

用第~\ref{sec:csv-manual-average}~節的成績資料，幫每個人多算一個「及格與否」欄位，
寫成一份新的 CSV：

\begin{example}{python}{(write-quiz-results.py) 寫出成績與及格結果}
import csv

students = [
    ("Alice", 88),
    ("Bob", 76),
    ("Taylor", 92),
    ("Charlie", 65),
    ("David", 79),
]

with open("quiz_results.csv", "w", newline="") as f:
    writer = csv.writer(f)
    writer.writerow(["name", "score", "passed"])
    for name, score in students:
        passed = "Yes" if score >= 70 else "No"
        writer.writerow([name, score, passed])

print("Wrote quiz_results.csv")
\end{example}

\iosample{}{Wrote quiz\_results.csv \\}

程式執行完之後，\texttt{quiz\_results.csv} 這個檔案的內容會是：

\begin{verbatim}
name,score,passed
Alice,88,Yes
Bob,76,Yes
Taylor,92,Yes
Charlie,65,No
David,79,Yes
\end{verbatim}

留意 \texttt{open()} 那一行多寫了 \texttt{newline=""} 這個參數。這是使用
\texttt{csv.writer()} 時官方文件明講的固定寫法，接下來就要老實示範一次，如果漏寫它會
發生什麼事。

\texttt{csv} 模組自己輸出的每一列，結尾用的是 \texttt{\textbackslash r\textbackslash n}
這兩個字元。如果你打開檔案時沒有加 \texttt{newline=""}，Python 的文字模式可能會把
\texttt{csv} 模組已經寫好的換行符號「再翻譯一次」，導致每一列後面多出一個看不見的
\texttt{\textbackslash r}，在某些系統上，這個多出來的字元會讓 Excel 或某些編輯器把每
一列資料之間顯示成多了一行空白。這裡直接動手檢查一次，看看在這台寫這本書用的電腦上，
這件事實際發生了沒有：

\begin{example}{python}{(check-newline-bug.py) 檢查漏寫 newline='' 的後果}
import csv

rows = [("name", "score"), ("Alice", 88), ("Bob", 76)]

with open("with_newline.csv", "w", newline="") as f:
    writer = csv.writer(f)
    for row in rows:
        writer.writerow(row)

with open("without_newline.csv", "w") as f:          # 忘記寫 newline=""
    writer = csv.writer(f)
    for row in rows:
        writer.writerow(row)

with_bytes = open("with_newline.csv", "rb").read()
without_bytes = open("without_newline.csv", "rb").read()

print("With newline='':", with_bytes.count(b"\r\r\n"), "extra newlines")
print("Without newline='':", without_bytes.count(b"\r\r\n"), "extra newlines")
\end{example}

\iosample{}{With newline='': 0 extra newlines \\{} Without newline='': 0 extra newlines \\}

老實說：在這台測試用的 macOS 電腦上，兩個檔案幾乎一模一樣，完全看不出差異。原因出在
「多一個 \texttt{\textbackslash r}」這個坑，根源是作業系統本身的換行慣例跟 \texttt{csv}
模組自己輸出的換行符號重複翻譯。macOS 跟 Linux 系統本身的換行字元本來就只有
\texttt{\textbackslash n}，不會被文字模式的換行轉換再多加一個 \texttt{\textbackslash r}；
但 Windows 系統的換行慣例是 \texttt{\textbackslash r\textbackslash n}，同一段程式如果
在 Windows 上執行，就真的會讓每一列資料之間多出一行看得見的空白。

\keypoint{\texttt{newline=""} 這個坑不是每台電腦都看得到，但這不代表可以放心省略它。
你寫的程式不會只在你自己的電腦上跑，你的讀者、你的同事很可能就在 Windows 上打開你寫的
程式。每次用 \texttt{csv.writer()} 都養成加上 \texttt{newline=""} 的習慣，比每次都去
記「我現在是不是在 Windows 上」更可靠，這也是 Python 官方文件明講建議的固定寫法。}

\begin{drills}[label=drills:csv-writer-bugs]
\item 抓蟲：一段程式先用 \texttt{writerow(["name", "score", "passed"])} 寫入標題列，接著
對某一位學生寫 \texttt{writerow(["Alice", 88])}（只給了兩個值），又對另一位學生寫
\texttt{writerow(["Bob", 76, "Yes", "extra"])}（給了四個值）。這段程式執行時會報錯
嗎？寫出來的 csv 檔案內容大概會是什麼樣子？
\item 概念題：為什麼在這一節實際測試用的電腦上，忘記寫 \texttt{newline=""} 沒有看到
「每列資料多一行空白」的症狀？這樣是不是代表以後可以不用寫它？
\item 如果你想在程式每次執行時，把新的一筆資料「接著」寫進已經存在的 csv 檔案後面，而
不是每次都把整個檔案清空重寫，\texttt{open()} 那一行該怎麼改？
\end{drills}

\section{用欄位名稱存取：DictReader 與 DictWriter}
\label{sec:csv-dict}

前面用 \texttt{csv.reader()} 讀出來的每一列，都是靠數字索引取值，像 \texttt{row[0]}、
\texttt{row[1]} 這樣。這樣寫有個隱患：萬一以後欄位的順序調整了（例如「地址」欄搬到最
前面），所有靠索引取值的程式碼就通通要跟著改，而且改錯了也不會有任何錯誤訊息提醒你，只
會默默拿錯欄位。呼應第 \ref{ch:dict-set} 章學過的字典操作，\texttt{csv} 模組提供了
\texttt{csv.DictReader}，可以直接用欄位名稱取值。
\index{csv.DictReader／csv.DictWriter}

\begin{code}{python}{用 csv.DictReader() 依欄位名稱讀取}
import csv

with open("filename.csv") as f:
    reader = csv.DictReader(f)
    for row in reader:
        # row 是一個 dict，key 是欄位標題，例如 row["name"]
        ...
\end{code}

用第~\ref{sec:csv-module-write}~節寫出來的 \texttt{quiz\_results.csv}，篩選出及格的
學生，再用 \texttt{csv.DictWriter} 把篩選結果寫成一份新檔案：

\begin{example}{python}{(filter-and-write-passed.py) 篩選及格名單並寫出}
import csv

with open("quiz_results.csv") as f:
    reader = csv.DictReader(f)
    passed_students = []
    for row in reader:
        if row["passed"] == "Yes":
            passed_students.append(row)

with open("passed.csv", "w", newline="") as f:
    writer = csv.DictWriter(f, fieldnames=["name", "score", "passed"])
    writer.writeheader()
    for student in passed_students:
        writer.writerow(student)

with open("passed.csv") as f:
    print(f.read())
\end{example}

\iosample{}{name,score,passed \\{} Alice,88,Yes \\{} Bob,76,Yes \\{} Taylor,92,Yes \\{} David,79,Yes \\}

\texttt{DictReader} 讀出來的每一列都是一個 dict，用 \texttt{row["passed"]} 取值，比
\texttt{row[2]} 一眼就看得懂在取什麼；\texttt{DictWriter} 反過來，寫入前要先用
\texttt{fieldnames} 告訴它欄位順序，再用 \texttt{writeheader()} 寫出標題列，之後每次
\texttt{writerow()} 都直接餵一個 dict 進去即可（不需要自己排欄位順序，\texttt{DictWriter}
會照 \texttt{fieldnames} 幫你排好）。

\stylepoint{用欄位名稱存取比用數字索引更好讀、也更不容易因為欄位順序調整而悄悄出錯。這
跟第 \ref{ch:dict-set} 章學到的道理是一樣的：字典用 key 存取，永遠比「猜第幾個位置」更清楚。}

到目前為止的範例，餵給 \texttt{writerow()} 的 dict，key 都跟 \texttt{fieldnames} 完全
對得上。如果哪天資料剛好缺了一欄，或是多了一欄呢？這兩種情況 \texttt{DictWriter} 的處理
方式並不對稱，值得認識清楚。

先看「缺欄位」的情況：

\begin{example}{python}{(write-missing-field.py) dict 少了 fieldnames 裡的一個欄位}
import csv

with open("missing_field.csv", "w", newline="", encoding="utf-8") as f:
    writer = csv.DictWriter(f, fieldnames=["name", "score", "passed"])
    writer.writeheader()
    writer.writerow({"name": "Alice", "score": 88})   # 少了 "passed"

with open("missing_field.csv", encoding="utf-8") as f:
    print(f.read())
\end{example}

\iosample{}{name,score,passed \\{} Alice,88, \\}

「passed」這一欄沒有出現在 dict 裡，\texttt{DictWriter} 並不會報錯，而是安靜地把這一欄
填成空字串，寫出來的那一列結尾多了一個逗號、後面什麼都沒有。如果想讓它填一個自己指定的
預設值，而不是空字串，可以在建立 \texttt{DictWriter} 時多給一個 \texttt{restval} 參數，
例如 \texttt{csv.DictWriter(f, fieldnames=[...], restval="Unknown")}。
\index{restval／extrasaction}

再看「多欄位」的情況：

\begin{example}{python}{(write-extra-field.py) dict 多了一個 fieldnames 沒有的欄位}
import csv

with open("extra_field.csv", "w", newline="", encoding="utf-8") as f:
    writer = csv.DictWriter(f, fieldnames=["name", "score"])
    writer.writeheader()
    try:
        writer.writerow({"name": "Alice", "score": 88, "note": "transfer student"})
    except ValueError:
        print("Write failed: dict has a key not in fieldnames")
\end{example}

\iosample{}{Write failed: dict has a key not in fieldnames \\}

這次反而直接丟出 \texttt{ValueError}，因為 \texttt{DictWriter} 不知道多出來的「note」
這個 key 該塞進哪一欄，寧可直接告訴你有問題，也不會自己偷偷決定要不要保留它。如果確定這
些多出來的欄位可以直接丟掉，可以在建立 \texttt{DictWriter} 時加上
\texttt{extrasaction="ignore"}，讓它安靜地忽略掉 \texttt{fieldnames} 以外的 key。

\keypoint{\texttt{DictWriter} 面對「dict 跟 fieldnames 對不上」這件事，兩種方向的處理
並不對稱：少寫的欄位會被默默填成空字串（除非用 \texttt{restval} 指定別的預設值），多出
來的欄位則會直接丟出 \texttt{ValueError}（除非用 \texttt{extrasaction="ignore"} 明確
允許忽略）。這種不對稱剛好反映一種設計上的態度：資料「不夠」還能用預設值撐過去，資料
「莫名其妙多了東西」卻可能代表哪裡搞錯了，值得讓程式直接告訴你。}

\section{如果分隔符號不是逗號：delimiter 參數}
\label{sec:csv-delimiter}

CSV 的全名雖然是「逗號分隔值」，但實務上收到的「表格式純文字檔案」不見得都乖乖用逗號
分隔欄位。最常見的另一種是用 Tab（鍵盤上的跳格鍵）當分隔符號，這種格式通常存成
\texttt{.tsv} 檔（tab-separated values，跟 CSV 是同一套邏輯，只是分隔符號不同）。假設
有一份用 Tab 分隔的 \texttt{scores.tsv}（下面用一個中括號 \texttt{[Tab]} 代表實際檔案
裡看不見的跳格字元）：

\begin{verbatim}
name[Tab]score
Alice[Tab]88
Bob[Tab]76
\end{verbatim}

如果直接用 \texttt{csv.reader(f)}，它預設只認得逗號，遇到這種檔案只會把整行當成一個
欄位，完全拆不開。\texttt{csv.reader()} 和 \texttt{csv.writer()} 都接受一個
\texttt{delimiter} 參數，讓你自己指定分隔符號：
\index{delimiter 參數}

\begin{code}{python}{用 delimiter 參數指定分隔符號}
import csv

with open("filename.tsv") as f:
    reader = csv.reader(f, delimiter="\t")
    for row in reader:
        ...
\end{code}

\begin{example}{python}{(read-tsv.py) 讀取一份用 Tab 分隔的檔案}
import csv

with open("scores.tsv", encoding="utf-8") as f:
    reader = csv.reader(f, delimiter="\t")
    for row in reader:
        print(row)
\end{example}

\iosample{}{['name', 'score'] \\{} ['Alice', '88'] \\{} ['Bob', '76'] \\}

\texttt{delimiter="\textbackslash t"} 裡的 \texttt{\textbackslash t}，用法跟第
~\ref{sec:csv-module-write}~節看過的 \texttt{\textbackslash n}、\texttt{\textbackslash r}
是同一套規則：用一個反斜線加一個字母，代表一個打字打不出來的特殊符號，
\texttt{\textbackslash t} 代表的就是一個 Tab 字元。\texttt{csv.writer()} 也接受同樣的
\texttt{delimiter}
參數，用法完全對稱：建立 \texttt{writer} 的時候多給這個參數，之後每次 \texttt{writerow()}
寫出來的每一列，欄位之間用的就會是你指定的分隔符號，而不是預設的逗號。

\stylepoint{不管分隔符號是逗號、Tab，還是其他符號，只要交給 \texttt{csv} 模組並用
\texttt{delimiter} 參數說清楚，前面學過的 \texttt{DictReader}、\texttt{DictWriter}、
\texttt{fieldnames}、\texttt{encoding} 這些技巧全部都還是同一套，不需要為了「分隔符號
不是逗號」重新學一遍。「表格式純文字檔案」的核心概念不會因為分隔符號改變，只有那一個
參數需要換。}

\section{沒有標題列的 CSV：用 fieldnames 自己指定欄位名稱}
\label{sec:csv-no-header}

前面幾節看到的 CSV 檔案，第一行永遠是欄位標題，\texttt{csv.DictReader} 也理所當然地把
第一行當成標題來用。但不是所有 CSV 檔案都長這樣：很多系統匯出的資料（例如感測器記錄、
交易紀錄）為了節省空間，直接省略標題行，每一行從頭到尾都是純資料。假設有一份溫度紀錄
\texttt{temp\_log.csv}，內容如下：

\begin{verbatim}
2024-01-01,Springfield,18.5
2024-01-02,Springfield,17.2
2024-01-03,Lakeside,22.0
\end{verbatim}

這份檔案完全沒有標題行，三個欄位分別是日期、城市、氣溫，但這件事只有你自己知道，
Python 不會幫你猜。如果你直接用上一節學到的 \texttt{csv.DictReader(f)}，它還是會照
原本的規矩，把第一行（也就是這裡的第一筆真正資料 \texttt{2024-01-01,Springfield,18.5}）當成
標題列，這樣一來，第一筆資料就白白被當成欄位名稱用掉了，實際的資料反而少了一筆。

要避免這個問題，\texttt{csv.DictReader} 提供了一個 \texttt{fieldnames} 參數，讓你自己
把欄位名稱列出來，告訴它「這份檔案沒有標題行，全部都是資料，欄位名稱由我來決定」：
\index{fieldnames 參數}

\begin{code}{python}{用 fieldnames 參數自己指定欄位名稱}
import csv

with open("filename.csv") as f:
    reader = csv.DictReader(f, fieldnames=[field1, field2, ...])
    for row in reader:
        # 每一列（包含原本的第一列）都當成資料處理，不會有任何一列被當成標題吃掉
        ...
\end{code}

\begin{example}{python}{(read-temp-log.py) 讀取一份沒有標題列的 CSV}
import csv

with open("temp_log.csv") as f:
    reader = csv.DictReader(f, fieldnames=["date", "city", "temperature"])
    for row in reader:
        print(row)
\end{example}

\iosample{}{\{'date': '2024-01-01', 'city': 'Springfield', 'temperature': '18.5'\} \\{} \{'date': '2024-01-02', 'city': 'Springfield', 'temperature': '17.2'\} \\{} \{'date': '2024-01-03', 'city': 'Lakeside', 'temperature': '22.0'\} \\}

三筆資料一筆都沒少，每一列都乖乖對應到 \texttt{fieldnames} 裡指定的欄位名稱。反過來，
如果這份檔案明明沒有標題列，你卻忘記傳 \texttt{fieldnames}，會發生什麼事？

\begin{example}{python}{(read-temp-log-buggy.py) 忘記傳 fieldnames 的後果}
import csv

with open("temp_log.csv") as f:
    reader = csv.DictReader(f)                # 忘記傳入 fieldnames
    rows = list(reader)
    print("Read", len(rows), "rows")
    for row in rows:
        print(row)
\end{example}

\iosample{}{Read 2 rows \\{} \{'2024-01-01': '2024-01-02', 'Springfield': 'Springfield', '18.5': '17.2'\} \\{} \{'2024-01-01': '2024-01-03', 'Springfield': 'Lakeside', '18.5': '22.0'\} \\}

這個結果比手動 \texttt{split()} 錯位還更難察覺：程式沒有報任何錯誤，執行起來一切正常，
但實際上只讀到 2 筆資料，不是原本的 3 筆，而且每一筆的 key 都變成第一筆資料裡的值
（\texttt{'2024-01-01'}、\texttt{'Springfield'}、\texttt{'18.5'}），完全不是想要的欄位名稱。
原因跟猜的一樣：\texttt{csv.DictReader} 沒有拿到 \texttt{fieldnames}，只好按照它原本
的規矩，把它看到的第一列資料當成標題列，這一列資料也就從「資料」變成「欄位名稱」，從此
消失在輸出結果裡。

\keypoint{\texttt{csv.DictReader} 預設一定會把讀到的第一列當成標題列，這件事不會因為
「這份檔案其實沒有標題列」而自動改變，你必須自己知道手上這份檔案有沒有標題行。有標題行
就直接用 \texttt{csv.DictReader(f)}；沒有標題行，一定要透過 \texttt{fieldnames} 參數
自己把欄位名稱列出來，否則第一筆真正的資料會被誤當成標題列，悄悄消失不見。}

\stylepoint{同樣是讀 CSV，普通的 \texttt{csv.reader()} 完全不會有這個問題，因為它從
頭到尾都只回傳一份份 list，不會特別把任何一列當成標題來對待，「這一列是不是標題」永遠
是程式設計者自己要處理的事。只有 \texttt{DictReader} 因為要幫你把每一列包成 dict、需要
key 的名稱，才會預設借用第一列當標題，這也是它比 \texttt{csv.reader()} 好用、卻也多一
個地方要留意的原因。}

反過來，如果你要用 \texttt{csv.DictWriter} 寫出一份「不要標題列」的檔案（例如要接在
既有的無標題檔案後面），只要不呼叫 \texttt{writeheader()} 就好，\texttt{DictWriter}
本身不會強迫你一定要寫標題列。

\begin{drills}[label=drills:csv-no-header]
\item 預測輸出：如果 \texttt{temp\_log.csv} 多加一行 \texttt{2024-01-04,Springfield,19.0}，
用範例 \texttt{read-temp-log.py} 正確傳入 \texttt{fieldnames} 的寫法重新執行，會印出
幾筆資料？最後一筆的內容是什麼？
\item 填空：一份沒有標題列的 CSV 檔案有「name」「phone」兩個欄位，想用 \texttt{DictReader}
讀取，請把下面的空格填上正確內容：
\begin{verbatim}
reader = csv.DictReader(f, fieldnames=______)
\end{verbatim}
\item 抓蟲：如果對第~\ref{sec:csv-module-read}~節的 \texttt{quiz\_scores.csv}（有標題
列「name,score」）誤用 \texttt{csv.DictReader(f, fieldnames=["a", "b"])}，讀出來的第一
筆資料會是什麼？
\end{drills}

\section{開檔案時的編碼問題：encoding="utf-8"}
\label{sec:csv-encoding}

除了「有沒有標題列」，另一個處理 CSV（或任何純文字檔案）時常踩到的坑是編碼。你的電腦
儲存文字檔案時，實際上存的是一串串數字（位元組），「編碼」就是這串數字跟螢幕上看到的
文字之間的對應規則。

UTF-8 是目前最通用的編碼方式，幾乎所有現代系統都看得懂；但有些軟體
（尤其是舊版 Windows 上的部分程式）預設會用其他編碼存檔，例如 Big5（台灣早期常見的中文
編碼）。如果打開檔案時使用的編碼，跟這份檔案實際存檔時用的編碼不一樣，輕則看到一堆亂
碼，重則程式直接報錯中斷。
\index{encoding 參數}

\begin{code}{python}{開檔案時明確指定編碼}
with open("filename.csv", encoding="utf-8") as f:
    ...
\end{code}

\begin{example}{python}{(write-utf8-scores.py) 先寫一份 UTF-8 編碼的 CSV 檔案}
with open("utf8_scores.csv", "w", encoding="utf-8") as f:
    f.write("name,score\n")
    f.write("Alice Ångström,88\n")

print("Wrote utf8_scores.csv")
\end{example}

\iosample{}{Wrote utf8\_scores.csv \\}

這份檔案本身完全沒有問題，寫入時明確用了 \texttt{encoding="utf-8"}。麻煩出在讀取的
那一端：如果之後打開它的人（可能是你自己換了一台電腦，也可能是別人）不小心指定了另一
種編碼，會發生什麼事？為了讓這個示範不管在哪台電腦上執行結果都一樣，這裡故意用
\texttt{encoding="big5"}（一種較舊的中文編碼）打開這份用 UTF-8 存的檔案，模擬編碼不
一致的情況：

\begin{example}{python}{(read-wrong-encoding.py) 用錯誤的編碼打開檔案}
with open("utf8_scores.csv", encoding="big5") as f:
    content = f.read()
print(content)
\end{example}

執行後程式並不會印出檔案內容，而是直接當掉：

\begin{verbatim}
Traceback (most recent call last):
  File "read-wrong-encoding.py", line 2, in <module>
    content = f.read()
UnicodeDecodeError: 'big5' codec can't decode byte 0xc3 in position 17:
illegal multibyte sequence
\end{verbatim}

\texttt{UnicodeDecodeError} 是「用來解碼這些位元組的編碼方式不對」時噴出的例外類型。
\index{UnicodeDecodeError}
訊息裡的 \texttt{position 17} 是說，從第 17 個位元組（從 0 算起）開始，Big5 這套規則沒
辦法把接下來的位元組拼回一個合法的字元（就是姓名裡 \texttt{Å} 這個字母編碼出來的位元組），
程式因此直接停下來。

\keypoint{每次用 \texttt{open()} 打開文字檔案，都應該明確寫上 \texttt{encoding="utf-8"}，
不要依賴電腦的「預設編碼」。道理跟第~\ref{sec:csv-module-write}~節 \texttt{newline=""}
的坑一模一樣：預設編碼在不同作業系統、不同地區設定的電腦上可能不一樣，自己的電腦能正常
打開檔案，不代表換一台電腦、換一個人的電腦也一樣行得通。}

還有一種更難抓的編碼陷阱，即使兩邊都乖乖寫著 \texttt{encoding="utf-8"} 也可能中招。
假設同事用 Excel「另存新檔」成「CSV UTF-8（逗號分隔）」格式匯出一份成績單
\texttt{excel\_export.csv}，內容看起來完全正常：

\begin{verbatim}
name,score
Alice,88
Bob,76
\end{verbatim}

用 \texttt{encoding="utf-8"} 打開、讀出來看看：

\begin{example}{python}{(read-excel-export.py) 讀取 Excel 匯出的 CSV UTF-8 檔案}
import csv

with open("excel_export.csv", encoding="utf-8") as f:
    reader = csv.DictReader(f)
    for row in reader:
        print(row)
        print(row["name"])
\end{example}

執行後會發現，印出來的第一個欄位名稱竟然多了一段看不懂的字元，緊接著就報錯：

\begin{verbatim}
{'\ufeffname': 'Alice', 'score': '88'}
Traceback (most recent call last):
  File "read-excel-export.py", line 7, in <module>
    print(row["name"])
          ~~~^^^^^^^^
KeyError: 'name'
\end{verbatim}

明明檔案內容看起來就是「name」四個字母，字典的 key 卻變成
\texttt{'\textbackslash ufeffname'}，多了一段看不見的 \texttt{\textbackslash ufeff}。
這並非檔案本身打錯字，原因是 Excel 存成「CSV UTF-8」格式時，習慣在檔案最開頭偷偷加上一個
特殊、完全看不見的隱藏字元，用來告訴其他軟體「這份檔案是用 UTF-8 編碼寫的」。這個隱藏
字元有個名稱，叫 BOM（byte order mark，位元組順序記號）。
\index{BOM／utf-8-sig}

單純的 \texttt{encoding="utf-8"}
不會主動幫你把這個記號拿掉，它會被當成資料的一部分，黏在第一個欄位名稱前面，讓
\texttt{"name"} 跟 \texttt{"\textbackslash ufeffname"} 變成兩個不一樣的字串，難怪
\texttt{row["name"]} 找不到。

要正確處理這種檔案，把 \texttt{encoding} 從 \texttt{"utf-8"} 換成 \texttt{"utf-8-sig"}
即可，這個變體版本會自動辨認並拿掉開頭的 BOM：

\begin{example}{python}{(read-excel-export-fixed.py) 用 utf-8-sig 正確處理 BOM}
import csv

with open("excel_export.csv", encoding="utf-8-sig") as f:
    reader = csv.DictReader(f)
    for row in reader:
        print(row["name"], "score:", row["score"])
\end{example}

\iosample{}{Alice score: 88 \\{} Bob score: 76 \\}

\stylepoint{如果打開的 CSV 檔案欄位名稱看起來一模一樣，用 \texttt{row["name"]} 卻查
不到、甚至丟出 \texttt{KeyError}，第一個該懷疑的對象就是檔案開頭那個看不見的 BOM，
尤其是這份檔案來自 Excel「另存新檔」匯出的時候。把 \texttt{encoding} 從
\texttt{"utf-8"} 換成 \texttt{"utf-8-sig"} 通常就能解決。}

\begin{drills}[label=drills:csv-encoding]
\item 預測輸出：如果把 \texttt{read-wrong-encoding.py} 裡的 \texttt{encoding="big5"}
改成跟寫入時一致的 \texttt{encoding="utf-8"}，這段程式還會出現 \texttt{UnicodeDecodeError}
嗎？
\item 填空：想明確指定用 UTF-8 編碼開啟一份純文字檔案，請把下面的空格填上正確內容：
\begin{verbatim}
f = open(filename, encoding=______)
\end{verbatim}
\item 抓蟲：一份從 Excel 存成「CSV UTF-8」格式匯出的檔案，用 \texttt{encoding="utf-8"}
打開，\texttt{DictReader} 讀出來的第一個欄位名稱卻多了一段看不懂的字元，導致
\texttt{row["name"]} 查不到。這是什麼問題？該怎麼解決？
\end{drills}

\section{JSON：有層次的純文字資料}
\label{sec:json-basics}

CSV 很適合「每一列長得都一樣、欄位固定」的表格資料，但如果資料本身有層次呢？比如一筆
個人資料裡，「hobbies」這個欄位本身就是好幾項興趣的清單，「contact」底下又還有電話、地址
好幾個子欄位，這種「一層包一層」的資料，硬要塞進扁平的表格會很彆扭。

JSON（JavaScript
Object Notation）就是為了這種情況設計的純文字格式：
\index{JSON（JavaScript Object Notation）}
每個值都有名字標籤，而且可以巢狀，
dict 裡面可以包 list、list 裡面可以包 dict，層層疊起來都沒問題。
\index{巢狀（nested）}
跟 CSV 一樣，JSON 也是
純文字，不管換到哪一種程式語言、哪一個系統都讀得懂，這也是它常常被用來在不同程式之間
交換資料的原因。

Python 標準函式庫的 \texttt{json} 模組負責在「JSON 格式的字串」跟「Python 的值」之間
互相轉換。
\index{json 模組／json.loads()／json.dumps()}

\begin{code}{python}{JSON 字串與 Python 值互轉}
import json

data = json.loads(json_text)      # 把 JSON 格式的字串解析成 Python 的值
json_text2 = json.dumps(data)     # 把 Python 的值轉換成 JSON 格式的字串
\end{code}

\begin{example}{python}{(parse-student-profile.py) 讀出一筆巢狀的個人資料}
import json

text = '''
{
    "name": "Alice",
    "grade": 8,
    "hobbies": ["basketball", "coding", "music"],
    "contact": {"phone": "0912345678", "address": "São Paulo"}
}
'''

student = json.loads(text)

print(student)
print(student["name"])
print(student["hobbies"])
print(student["contact"]["phone"])
print(json.dumps(student, ensure_ascii=False))
\end{example}

\iosample{}{\{'name': 'Alice', 'grade': 8, 'hobbies': ['basketball', 'coding', 'music'], 'contact': \{'phone': '0912345678', 'address': 'São Paulo'\}\} \\{} Alice \\{} ['basketball', 'coding', 'music'] \\{} 0912345678 \\{} \{"name": "Alice", "grade": 8, "hobbies": ["basketball", "coding", "music"], "contact": \{"phone": "0912345678", "address": "São Paulo"\}\} \\}

\texttt{json.loads()} 把整段 JSON 文字讀進來之後，變成一個很普通的 Python dict。
\texttt{"hobbies"} 這個 key 對應到的值是一份 list，\texttt{"contact"} 對應到的值又是另一
個 dict，你可以像平常操作 dict、list 一樣，用 \texttt{student["contact"]["phone"]}
這樣一路串下去取值。跟 CSV 讀出來的東西不一樣的地方是：JSON 裡的數字（例如
\texttt{"grade": 8}）讀進來之後就是真正的整數，不是字串，因為 JSON 格式本身就有分
字串、數字、布林值等不同型別，不像 CSV 全部都是純文字。

\keypoint{\texttt{json.dumps()} 如果不加 \texttt{ensure\_ascii=False}，非 ASCII
字元（例如帶重音符號的字母 \texttt{ã}）會被轉成像 \texttt{\textbackslash u00e3} 這樣的
跳脫序列。這不是壞掉，Python 只是預設用任何系統都看得懂的純 ASCII 方式表示字串。但如果
你的輸出是要給人看的（例如存成檔案、印在畫面上），幾乎都會想加上
\texttt{ensure\_ascii=False}，讓這些字元照原樣顯示。}

\section{把 JSON 資料直接讀寫成檔案：json.dump() 與 json.load()}
\label{sec:json-dump-load}

剛剛的 \texttt{json.loads()} 和 \texttt{json.dumps()} 操作的對象都是字串：先準備好
一段 JSON 格式的文字，或是先把 Python 的值轉成一段 JSON 格式的文字。但實務上更常見的
情境是直接把資料存成一份 \texttt{.json} 檔案，或是直接從一份 \texttt{.json} 檔案讀出
資料，不需要自己先手動讀出字串、再轉換。\texttt{json} 模組另外提供了 \texttt{json.dump()}
和 \texttt{json.load()} 這一對函式，可以直接對著一個已經打開的檔案物件操作。
\index{json.dump()／json.load()}

\keypoint{\texttt{json.dump()}（沒有 s）跟 \texttt{json.dumps()}（有 s）只差一個字母，
功能卻不一樣：\texttt{json.dump(值, 檔案物件)} 直接把值寫進「檔案」；
\texttt{json.dumps(值)} 是把值轉成一段 JSON「字串」讓你自己處理。\texttt{load()} 跟
\texttt{loads()} 也是同樣的分工：\texttt{load(檔案物件)} 直接從檔案讀，
\texttt{loads(字串)} 是從一段已經準備好的字串讀。這兩組函式名稱只差一個字母，卻是完全
不同的用法，第一次用的時候務必看清楚自己到底要對「檔案」還是「字串」操作。}

\begin{code}{python}{直接對檔案讀寫 JSON}
import json

with open("filename.json", "w", encoding="utf-8") as f:
    json.dump(data, f, ensure_ascii=False, indent=2)

with open("filename.json", encoding="utf-8") as f:
    data = json.load(f)
\end{code}

跟第~\ref{sec:csv-encoding}~節學到的習慣一樣，開啟 \texttt{.json} 檔案時也建議明確加
上 \texttt{encoding="utf-8"}，理由完全相同：不要依賴電腦的預設編碼。用
第~\ref{sec:json-basics}~節的班級興趣資料，示範怎麼把一份巢狀資料直接寫成一份排版整齊
的 JSON 檔案：

\begin{example}{python}{(write-members-json.py) 把巢狀資料寫成 JSON 檔案}
import json

members = [
    {"name": "Alice", "grade": 8, "hobbies": ["basketball", "coding"]},
    {"name": "Bob", "grade": 7, "hobbies": ["music"]},
]

with open("members.json", "w", encoding="utf-8") as f:
    json.dump(members, f, ensure_ascii=False, indent=2)

print("Wrote members.json")
\end{example}

\iosample{}{Wrote members.json \\}

程式執行完之後，\texttt{members.json} 這個檔案的內容會是：

\begin{verbatim}
[
  {
    "name": "Alice",
    "grade": 8,
    "hobbies": [
      "basketball",
      "coding"
    ]
  },
  {
    "name": "Bob",
    "grade": 7,
    "hobbies": [
      "music"
    ]
  }
]
\end{verbatim}

留意 \texttt{indent=2} 這個參數：它讓 \texttt{json.dump()} 自動幫巢狀結構加上縮排、
換行，檔案打開來一眼就看得出層次，跟前面 \texttt{json.dumps()} 預設擠成一整行的樣子
差很多。\texttt{indent} 不是必要參數，純粹是為了讓存出來的檔案給人看的時候更好讀；如果
這份檔案只是要給另一支程式讀，不加 \texttt{indent} 完全沒問題，甚至檔案還會更小。

再用 \texttt{json.load()} 把這份檔案讀回來，確認資料真的原封不動：

\begin{example}{python}{(read-members-json.py) 從 JSON 檔案讀回巢狀資料}
import json

with open("members.json", encoding="utf-8") as f:
    data = json.load(f)

print(data)
print(data[0]["hobbies"])
\end{example}

\iosample{}{[\{'name': 'Alice', 'grade': 8, 'hobbies': ['basketball', 'coding']\}, \{'name': 'Bob', 'grade': 7, 'hobbies': ['music']\}] \\{} ['basketball', 'coding'] \\}

讀回來的 \texttt{data} 是一份 list，裡面每個元素都是 dict，跟寫入前的 \texttt{members}
結構一模一樣，\texttt{indent=2} 只影響檔案存起來給人看的排版，不會改變讀回 Python 之後
的資料內容。

\begin{drills}[label=drills:json-dump-load]
\item 填空：下面這行呼叫少了第二個引數，請把空格填上正確內容：
\begin{verbatim}
json.dump(data, ______)
\end{verbatim}
\item 概念題：\texttt{json.dump()} 和 \texttt{json.dumps()} 這兩個名稱只差一個字母
\texttt{s}，功能卻不一樣，差別是什麼？
\end{drills}

\section{從巢狀 JSON 資料裡安全取值}
\label{sec:json-nested}

實際拿到的 JSON 資料常常不是每一筆都長得一模一樣，有些成員填了「hobbies」欄位，有些人沒
填。如果直接用 \texttt{member["hobbies"]} 取值，遇到沒有這個欄位的成員就會出錯。這裡剛好
呼應第 \ref{ch:files-exceptions} 章的 EAFP 心智模型，也呼應第 \ref{ch:dict-set} 章學過的 \texttt{dict.get()}：與其一路寫
\texttt{if "hobbies" in member} 這種先檢查再取值的寫法，不如直接嘗試取值，用
\texttt{.get()} 或 \texttt{try/except KeyError} 接住「查不到」的情況。

\begin{example}{python}{(get-nested-hobbies.py) 安全地取出可能不存在的欄位}
import json

text = '''
[
    {"name": "Alice", "grade": 8, "hobbies": ["basketball", "coding"]},
    {"name": "Bob", "grade": 7},
    {"name": "Taylor", "grade": 9, "hobbies": ["music"]}
]
'''

members = json.loads(text)

for member in members:
    hobbies = member.get("hobbies", [])
    print(member["name"], "hobbies:", hobbies)

print("---")

for member in members:
    try:
        hobbies = member["hobbies"]
    except KeyError:
        hobbies = []
    print(member["name"], "hobbies:", hobbies)
\end{example}

\iosample{}{Alice hobbies: ['basketball', 'coding'] \\{} Bob hobbies: [] \\{} Taylor hobbies: ['music'] \\{} --- \\{} Alice hobbies: ['basketball', 'coding'] \\{} Bob hobbies: [] \\{} Taylor hobbies: ['music'] \\}

這份 JSON 資料裡，Bob 沒有「hobbies」這個欄位。第一段用 \texttt{member.get("hobbies", [])}，
查不到就直接拿到預設值 \texttt{[]}（一份空清單），完全不用寫任何 \texttt{if} 判斷；第
二段改用 \texttt{try/except KeyError}，做的事情一樣，只是換一種寫法。

兩段程式印出來的
結果完全相同，因為 \texttt{.get(key, 預設值)} 本質上就是「查得到就回傳值、查不到就回傳
預設值」這件事的一個現成寫法，不用自己寫 \texttt{try/except} 就能達到一樣的效果。這正
是 EAFP 精神的具體展現：不要先小題大作去確認資料存不存在，直接嘗試取值，讓語言本身的行
為替你簡化判斷。

巢狀資料底下的值，一樣可以用已經很熟悉的迴圈邏輯處理，不會因為多包了一層字典就變得特別
困難。假設每個成員的資料除了姓名，還多了一個「scores」欄位，本身又是一個 dict（科目對應
到分數）：

\begin{example}{python}{(average-nested-scores.py) 從巢狀資料算出平均分數}
import json

text = '''
[
    {"name": "Alice", "scores": {"Literature": 85, "Math": 92, "English": 78}},
    {"name": "Bob", "scores": {"Literature": 70, "Math": 65, "English": 80}}
]
'''

members = json.loads(text)

for member in members:
    name = member["name"]
    scores = member["scores"]

    total = 0
    for subject in scores:
        total += scores[subject]
    average = total / len(scores)

    print(f"{name}'s average score is {average:.1f}")
\end{example}

\iosample{}{Alice's average score is 85.0 \\{} Bob's average score is 71.7 \\}

拆開來看，外層的 \texttt{for member in members} 跟前面的寫法一樣，走訪清單裡每一個
成員；內層的 \texttt{scores = member["scores"]} 先把這個成員的成績字典整個取出來，接著
\texttt{for subject in scores} 走訪這個字典的每一個 key（也就是每個科目），把對應的
分數累加起來，這正是第~\ref{sec:csv-manual-average}~節學過的「累加器」樣板，只是這次
累加的數字是從巢狀字典裡一層層取出來的，不是直接從 CSV 欄位讀到的。

JSON 的巢狀結構乍看
複雜，但只要一層一層拆開來看，每一層用到的其實都是已經會的工具：走訪 list 用
\texttt{for}，走訪 dict 用 \texttt{for}，取值用中括號，全部都不是新東西。

到這裡，你已經認識了兩種最常見的純文字資料格式，也知道每一種都可以直接對著檔案讀寫，
不用先手動處理字串。該用哪一種，其實看資料本身的形狀：資料是一張表格、每一筆記錄的欄位
都固定的時候，用 CSV；資料有層次、欄位可能不是每筆都有的時候，用 JSON。兩者都是純文字，
這也是它們能在不同系統、不同程式語言之間暢行無阻的原因。

\begin{drills}[label=drills:json-nested]
\item 預測輸出：如果有一個 \texttt{member = \{"name": "Charlie"\}}（沒有 "grade" 欄位），
程式寫 \texttt{member.get("grade", "Unknown")}，印出來的結果會是什麼？
\item 填空：完成下面用 \texttt{try/except} 取代 \texttt{.get()} 的寫法：
\begin{verbatim}
try:
    grade = member["grade"]
except ______:
    grade = "Unknown"
\end{verbatim}
\item 抓蟲：如果直接寫 \texttt{member["grade"]}（沒有用 \texttt{.get()}，也沒有包
\texttt{try/except}），而 \texttt{member} 剛好沒有這個欄位，會發生什麼事？
\end{drills}

\section*{練習解答}
\addcontentsline{toc}{section}{練習解答}

\paragraph{第~\ref{sec:csv-manual-average}~節}
\begin{enumerate}
\item 總分是 150，平均分數是 30.0（\texttt{10+20+30+40+50=150}，\texttt{150/5=30.0}）。
\item 依序是 \texttt{0}、\texttt{+=}、\texttt{len(numbers)}。
\item 因為第一行標題「name,score」的 \texttt{fields[1]} 是字串 \texttt{"score"}，
\texttt{int("score")} 會丟出 \texttt{ValueError: invalid literal for int() with base
10: 'score'}，程式在還沒印出任何總分或平均之前就直接中斷。
\end{enumerate}

\paragraph{第~\ref{sec:csv-module-write}~節}
\begin{enumerate}
\item 不會報錯。\texttt{csv.writer} 不會檢查你給的值的數量跟標題列是否一致，它就是照你
給的清單原封不動寫成一行。所以檔案裡會出現一行只有兩個欄位、一行有四個欄位，跟標題的三
欄對不上，之後如果用 \texttt{csv.DictReader} 或依照固定欄位數取值就會出錯或錯位。
\item 因為這個坑的根源是「作業系統本身的換行慣例」跟「\texttt{csv} 模組自己輸出的換行
符號」重複翻譯造成的，macOS／Linux 的換行字元本來就只有 \texttt{\textbackslash n}，不
會被文字模式的換行轉換再多加一個 \texttt{\textbackslash r}；但 Windows 的換行慣例是
\texttt{\textbackslash r\textbackslash n}，同一段程式在 Windows 上執行仍然會出現每列
資料間多一行空白。因為程式不一定只在自己的電腦上執行，還是要養成永遠加上
\texttt{newline=""} 的習慣。
\item 把 \texttt{open(filename, "w", newline="")} 改成 \texttt{open(filename, "a",
newline="")}。\texttt{'a'} 是附加模式，不會清空舊資料，這是第 \ref{ch:files-exceptions} 章教過的
\texttt{'w'} 跟 \texttt{'a'} 差異，在 csv 檔案上一樣適用。
\end{enumerate}

\paragraph{第~\ref{sec:csv-no-header}~節}
\begin{enumerate}
\item 會印出 4 筆資料，最後一筆是
\texttt{\{'date': '2024-01-04', 'city': 'Springfield', 'temperature': '19.0'\}}。
\item \texttt{["name", "phone"]}。
\item \texttt{\{'a': 'name', 'b': 'score'\}}，也就是把標題列本身當成一筆普通資料，欄位
名稱變成自己指定的 \texttt{'a'}、\texttt{'b'}，跟檔案裡真正的欄位「name」「score」完全
對不上，但程式不會有任何錯誤訊息提醒你。
\end{enumerate}

\paragraph{第~\ref{sec:csv-encoding}~節}
\begin{enumerate}
\item 不會。讀取跟寫入用的是同一種編碼，位元組能正確還原成原本的文字，不會出現
\texttt{UnicodeDecodeError}。
\item \texttt{"utf-8"}。
\item 檔案開頭藏了一個看不見的 BOM（Excel 存成「CSV UTF-8」格式時會自動加上），黏在
\texttt{"name"} 這個欄位名稱前面。解法是把 \texttt{encoding} 從 \texttt{"utf-8"} 換成
\texttt{"utf-8-sig"}，讓 Python 自動辨認並拿掉這個記號。
\end{enumerate}

\paragraph{第~\ref{sec:json-dump-load}~節}
\begin{enumerate}
\item 已經用 \texttt{open()} 打開的檔案物件（也就是 \texttt{open()} 的回傳值）。
\item \texttt{json.dump()}（沒有 s）直接把資料寫進一個已經打開的檔案物件；
\texttt{json.dumps()}（有 s）是把資料轉成一段 JSON 字串，回傳給你自己處理，不會碰任何
檔案。
\end{enumerate}

\paragraph{第~\ref{sec:json-nested}~節}
\begin{enumerate}
\item Unknown，因為 \texttt{member} 沒有 \texttt{"grade"} 這個 key，\texttt{.get()} 會
回傳你給的預設值。
\item \texttt{KeyError}。
\item 會丟出 \texttt{KeyError: 'grade'}，程式直接中斷（除非外層有 \texttt{try/except}
接住這個錯誤）。
\end{enumerate}

\chapter{資料工具一瞥（numpy／pandas／matplotlib）}
\label{ch:data-tools}

到目前為止，處理資料靠的都是你自己一行一行寫迴圈：讀進一份 CSV、一筆一筆把數字挑出來、
自己維護一個總和變數慢慢加、自己數共有幾筆才能算平均。

上一章你就是這樣算出一整班的平均
分數的，先用 \texttt{open()} 開檔、用 \texttt{.strip().split(",")} 把每一行拆成欄位、
把字串轉成數字，再用一個 \texttt{total} 變數、一個計數器，一步一步把平均值湊出來。這套
辦法沒有錯，資料筆數不多的時候也夠用，但這一章想讓你先體會一次它的極限：如果那份 CSV
一口氣多到全校三十個班、一千個學生呢？如果不只要算平均，還要挑出某個班級、某個分數區間
的資料呢？

自己手寫的迴圈會愈寫愈長、愈寫愈容易出錯，你花的力氣其實大半耗在跟語法奮鬥，
離真正想解決的問題反而愈來愈遠。這一章要告訴你一件事：這種需求太常見了，早就有別人把
常用的做法寫成現成的工具，你可以直接拿來用。而「拿別人寫好的工具來用」這件事本身，也是
寫程式的人每天都在做的事，值得專門認識一下它是怎麼運作的。

所以這一章分成八段推進，比起單純學三個新套件，更像是帶你重新走一次上一章的旅程，只是
這次手上多了工具可以用：

\begin{itemize}
\item 先搞懂「套件」是什麼、\texttt{pip} 這個工具在做什麼、怎麼把別人
寫好的套件裝進你的電腦（\ref{sec:packages-and-pip}、\ref{sec:install-and-pip} 節）。
\item 接著從最小的痛點下手，用 numpy 解決「一次對一
整批數字做運算」這件事，不用再自己寫迴圈一個一個算，也順便看看「比較運算」套用在整批
數字上會發生什麼事（\ref{sec:numpy-vectorize}、\ref{sec:numpy-boolean} 節）。
\item 再來是這一章份量最重的一段，用 pandas 重做一次你
在上一章已經很熟悉的事，讀 CSV、挑欄位、篩資料、算平均，甚至進一步學會同時套用好幾個
篩選條件、簡單處理缺漏的資料（\ref{sec:pandas-read-select} 到 \ref{sec:pandas-multi-filter} 節），你會發現同一件事現在只要幾行就做完，
直接對照出「工具」到底幫你省了什麼。
\item 最後，用 matplotlib 把整理好的資料畫成長條圖和折線
圖各一張，作為這一章、也是這三個套件的一個小小成就感收束（\ref{sec:matplotlib-basic} 節）。
\end{itemize}

這一整章會延續同
一份資料，上一章那份「全班小考成績」的 CSV，讓你清楚感受到「同一份資料，換了工具再做
一次」的對照感。

先說在前面：這一章的目標不是把你訓練成資料科學家，這三個套件各自都還有
大片你這裡碰不到的功能，這一章要的只是「認識、不精熟」：知道這些工具存在、知道它們大概
能幫你做什麼、真的動手用過幾次，之後如果工作或生活中遇到類似的大量資料需求，你會想到
「這個好像可以用 pandas 處理」，這樣就達到目的了。

\section{你已經知道的極限，以及套件是什麼}
\label{sec:packages-and-pip}

先問你一個問題：如果上一章那份成績單一口氣多到一千筆，你打算怎麼辦？答案大概是，還是
可以用一樣的迴圈邏輯硬做，只是程式會愈跑愈久，而且只要你多想要一個功能（例如「只看某個
班級」「只看超過某個分數的人」），就要多寫一段新的迴圈跟新的變數。這種「每多一個需求就
多寫一段迴圈」的做法，遲早會把你自己繞暈。

到目前為止，你 \texttt{import} 過的 \texttt{csv}、\texttt{math}、\texttt{json} 這些，
都是 Python「內建」就有的東西，只要裝好 Python，它們就已經在你的電腦上了，不用另外
安裝，這種東西有個統稱，叫「標準函式庫」。

世界上還有更多好用的工具，是\textbf{別人
額外寫好、放到網路上讓大家下載}的東西，跟 Python 官方原廠附帶的內建工具是兩回事，這種
東西叫「套件」（package）。\index{套件}numpy、pandas、matplotlib 就是三個非常有名的套件，分別擅長
「對整批數字做運算」「把表格式資料讀進來、整理、篩選」「把資料畫成圖表」，剛好對應到你
在上一章手動處理資料時會卡關的地方。

要拿到這些套件，你需要 Python 內建的一個小工具，叫 \texttt{pip}。\index{pip}你可以把 \texttt{pip}
想成一個「套件的搬運工」：你告訴它套件的名字，它就去網路上把套件抓下來、安裝進你的電腦，
之後你才能在程式裡 \texttt{import} 它。這些套件實際上放在一個叫 PyPI（Python Package
Index）\index{PyPI}的公開套件倉庫裡，全世界任何人都可以把自己寫的套件上傳到那裡，也可以透過
\texttt{pip} 免費把別人上傳的套件下載下來用，numpy、pandas、matplotlib 都是放在 PyPI
上的套件。裝 Python 的時候，\texttt{pip} 通常已經跟著一起裝好了，不用你另外安裝這個
工具本身。下一節就會實際動手，把三個套件裝進你的電腦。

\texttt{pip} 除了「裝」套件之外，也能「拆」套件：如果哪天不想再用某個套件了，可以打
\texttt{pip uninstall 套件名稱} 把它移除；套件本身也會不時推出新版本，修掉舊版的
bug、加上新功能，這時候可以打 \texttt{pip install --upgrade 套件名稱} 把已經裝好的
套件升級到最新版。\index{pip uninstall／pip install --upgrade}這幾個指令用法都差不多，記得「\texttt{pip} 加一個動作、加套件名稱」
這個固定格式，之後真的需要拆套件、升級套件時，回頭查一下指令名稱就好，不用死背。

\section{安裝套件，第一次見到 \texttt{pip}}
\label{sec:install-and-pip}

\begin{code}{bash}{用 pip 安裝套件的固定寫法}
pip install package1 package2 package3
\end{code}

打開終端機（第 \ref{ch:introduction} 章已經用過的那個視窗），輸入以下指令，一次裝好這一章要用的三個套件：

\begin{verbatim}
$ pip install numpy pandas matplotlib
Collecting numpy
  Downloading numpy-2.5.2-...whl (5.3 MB)
Collecting pandas
  Downloading pandas-3.0.5-...whl (9.9 MB)
Collecting matplotlib
  Downloading matplotlib-3.11.1-...whl (9.3 MB)
（中間還會看到一串安裝相依套件的訊息捲過去，這裡省略）
Successfully installed matplotlib-3.11.1 numpy-2.5.2 pandas-3.0.5 ...
\end{verbatim}

實際畫面上看到的版本號、檔案大小、捲動的訊息內容，會因為你安裝的當下套件釋出到第幾版
而略有不同，不用擔心跟這裡不完全一樣，只要最後看到類似 \texttt{Successfully installed}
這樣的字樣，就代表三個套件都裝好了。

裝好之後，如果你在程式裡寫 \texttt{import numpy}，Python 就能正常找到它。但如果你剛剛
跳過了安裝這一步，直接寫下面這支程式會發生什麼事？

\begin{example}{python}{(try-numpy.py) 還沒安裝套件就 import}
import numpy
\end{example}

在一台還沒裝過 numpy 的電腦上執行，畫面上會出現：

\begin{verbatim}
Traceback (most recent call last):
  File "try-numpy.py", line 1, in <module>
    import numpy
ModuleNotFoundError: No module named 'numpy'
\end{verbatim}

\texttt{ModuleNotFoundError}\index{ModuleNotFoundError} 意思很直白，就是「找不到一個叫這個名字的模組」。Python
到你電腦上該找套件的地方翻了一圈，沒看到 numpy，因為你還沒有用 \texttt{pip} 把它搬過來。
修法也很直接：回頭執行 \texttt{pip install numpy}，裝好之後同一支程式就能正常執行。

\keypoint{你不需要在 \texttt{import} 之前先確認某個套件到底裝了沒，直接寫下去就好，
Python 會用 \texttt{ModuleNotFoundError} 這個錯誤訊息老實告訴你答案，缺什麼就補什麼。
這跟這本書前面提過的「先做、錯了再處理」是同一種精神，只是這次犯錯的地方換成了套件
安裝，程式邏輯本身反而沒有問題。}

如果哪天你想確認某個套件到底裝好了沒，或是想知道自己裝的是第幾版，可以在終端機打
\texttt{pip show 套件名稱}：

\begin{verbatim}
$ pip show numpy
Name: numpy
Version: 2.5.2
Summary: Fundamental package for array computing in Python
Home-page: https://numpy.org
...
\end{verbatim}

如果那個套件根本沒裝過，畫面上會出現一行 \texttt{WARNING: Package(s) not found:
numpy} 提醒你。這個指令平常不會天天用到，但排查問題、確認環境的時候很方便，知道有這麼
一個工具就夠了。

\section{numpy：一次對整批數字做運算}
\label{sec:numpy-vectorize}

先從最小的痛點下手。假設老師覺得這次小考題目偏難，決定幫全班的分數統一套用一個調分
公式：\texttt{調整後分數 = 原始分數 * 0.9 + 10}（分數愈高，乘以 0.9 扣的愈多；分數愈
低，加上 10 分墊高的比例愈高，是一種常見的拉近高低分距的調分方式）。

如果全班只有幾個
人，你當然可以像上一章那樣寫個迴圈，一個一個算；但如果要調的是全校好幾百個人的分數，
每次都要重新寫、重新想這段迴圈，未免太累。numpy 就是為了解決「對一整批數字做同一種
運算」這件事而生的套件。

\begin{code}{python}{numpy 建立 array 與四則運算的固定寫法}
import numpy as np

array_name = np.array(regular_list)
result = array_name * constant + constant
\end{code}

慣例上，大家都把 numpy 用 \texttt{as np} 縮寫成 \texttt{np}，之後每次呼叫 numpy 的
功能就少打幾個字。\texttt{np.array()}\index{numpy array} 可以把一般的 list 包成 numpy 專屬的「array」，
包好之後，對這個 array 做加、減、乘、除，會\textbf{自動對裡面每一個元素分別做同一次
運算}，不用你自己寫迴圈一個一個算。

先看看用你已經很熟悉的迴圈寫法，怎麼幫全班調分：

\begin{example}{python}{(adjust-scores-loop.py) 用迴圈幫全班調分}
scores = [88, 72, 95, 60, 78, 85, 91, 55, 68]

# 用迴圈一個一個調分：老師決定套用「乘以 0.9 再加 10」的公式
adjusted_scores = []
for score in scores:
    adjusted_scores.append(score * 0.9 + 10)

print(adjusted_scores)
\end{example}

\iosample{}{[89.2, 74.8, 95.5, 64.0, 80.2, 86.5, 91.9, 59.5, 71.2] \\}

再看看同一件事，換成 numpy 怎麼做：

\begin{example}{python}{(adjust-scores-numpy.py) 用 numpy 幫全班調分}
import numpy as np

scores = [88, 72, 95, 60, 78, 85, 91, 55, 68]
scores_array = np.array(scores)

# 同樣的調分公式：array 之間的運算會自動套用到每一個元素
adjusted_scores_array = scores_array * 0.9 + 10

print(adjusted_scores_array)
\end{example}

\iosample{}{[89.2 74.8 95.5 64.  80.2 86.5 91.9 59.5 71.2] \\}

兩段程式算出來的數字完全一樣，但這就是「向量化運算」\index{向量化運算}帶來的差別：迴圈版要開一個空
list、寫一行 \texttt{for}、呼叫一次 \texttt{.append()}，三個步驟才能湊出結果；numpy
版一行 \texttt{scores\_array * 0.9 + 10} 就做完，等於把三行邏輯濃縮成一行，行數少了
三分之二，好處不只是概念上比較方便，寫程式碼的份量也真的少了一大截。

你可能注意到
numpy 印出來的那個 \texttt{64.} 後面多了一個空格，這其實是 numpy 自動幫忙對齊小數點
位置的結果，不是排版出了錯，方便你一眼看出每個數字對得整不整齊。

這裡要特別提醒一件事：array 的 \texttt{*}、\texttt{+} 跟你在第 \ref{ch:lists} 章學過的 list 的
\texttt{*}（例如 \texttt{[0] * 5} 是把清單內容重複 5 次）、\texttt{+}（把兩個清單接
在一起）意思完全不同：array 的 \texttt{*}、\texttt{+} 是逐一元素做數學運算，list 的
\texttt{*}、\texttt{+} 是重複或接合整份清單。這正是 numpy array 跟一般 list 最大的
差別之一，看到同樣的符號，要記得先確認自己手上的是 array 還是 list。

numpy 能做的事情遠不只這裡示範的一維、逐一運算，它還有更多維度陣列（例如把資料排成
像棋盤一樣有列有欄的表格）、矩陣運算，甚至「broadcasting」（讓不同形狀的 array 也能
互相運算的規則）這些功能，但那些屬於更進階的數學／科學運算範疇，這本書不會用到，這裡
認識「把 list 包成 array、直接做四則運算和比較運算」這個小角落就夠用了，比較運算的
部分下一節馬上會看到。

\begin{drills}[label=drills:numpy-vectorize]
\item 這段程式執行後，會印出什麼？
\begin{verbatim}
import numpy as np

values = np.array([2, 4, 6])
print(values * 3)
\end{verbatim}
\item 這段程式執行後，會印出什麼？
\begin{verbatim}
import numpy as np

a = np.array([10, 20, 30])
b = np.array([1, 2, 3])
print(a - b)
\end{verbatim}
\end{drills}

\section{numpy：比較運算與布林陣列}
\label{sec:numpy-boolean}

上一節示範的是四則運算，其實 numpy 的向量化不只適用在加減乘除上，第 \ref{ch:conditionals} 章學過的比較
運算子（\texttt{==}、\texttt{>}、\texttt{>=} 這些）套用在 array 上，一樣會自動對每一
個元素分別做比較，得到一整個由 \texttt{True}／\texttt{False} 組成的新 array。

\begin{code}{python}{numpy 比較運算的固定寫法}
import numpy as np

array_name = np.array(regular_list)
result = array_name comparison_operator constant
\end{code}

還是拿全班的分數當例子：老師想知道這次考試哪些人拿到 80 分以上，一次看清楚全班每個人
「有沒有」達標，而不是一個一個分數自己比對。

\begin{example}{python}{(high-score-mask.py) 一次比較全班分數是否達標}
import numpy as np

scores = [88, 72, 95, 60, 78, 85, 91, 55, 68]
scores_array = np.array(scores)

# 比較運算子套用在 array 上，一樣會逐一元素做比較
is_high = scores_array >= 80
print(is_high)
\end{example}

\iosample{}{[ True False  True False False  True  True False False] \\}

印出來的結果是一整條由 \texttt{True}、\texttt{False} 組成的 array，每個位置對應原本
分數 array 同一個位置的人：第一個人 88 分，達標，對應 \texttt{True}；第二個人 72 分，
沒達標，對應 \texttt{False}，以此類推。這種由比較運算得出的布林 array，很適合拿來當
「這批資料裡有哪些符合條件」的一次性總覽。

如果你只想知道「總共有幾個人達標」，不需要自己一個一個數 \texttt{True} 有幾個，
可以直接呼叫 \texttt{.sum()}：

\begin{example}{python}{(count-high-scores.py) 算出達標人數}
import numpy as np

scores_array = np.array([88, 72, 95, 60, 78, 85, 91, 55, 68])
is_high = scores_array >= 80

count_high = is_high.sum()
print(count_high)
\end{example}

\iosample{}{4 \\}

\keypoint{在 Python 裡，\texttt{True} 其實可以被當成數字 \texttt{1} 使用，
\texttt{False} 可以被當成 \texttt{0}，這一點跟你在數學課學到的「布林值就是布林值、
數字就是數字」不太一樣。把一整條布林 array 丟進 \texttt{.sum()}，效果就等於把每個
\texttt{True} 當 1、每個 \texttt{False} 當 0 全部加起來，加總的結果剛好就是符合條件
的元素個數，不用另外寫迴圈自己數。}

你可能已經注意到，\texttt{.sum()} 這個方法名稱看起來有點眼熟。下一節認識 pandas 的
時候，你會看到 DataFrame 的某一欄也有一模一樣叫 \texttt{.sum()}、\texttt{.mean()}
的方法，這不是巧合，pandas 實際上是建立在 numpy 之上的套件，兩者共用很多類似的操作
方式，先在這裡認識過一次，等一下在 pandas 那邊看到同樣的名字會覺得眼熟，不是全新的
東西。

\begin{drills}[label=drills:numpy-boolean]
\item 這段程式執行後，會印出什麼？
\begin{verbatim}
import numpy as np

values = np.array([50, 85, 62, 90, 40])
print(values >= 60)
\end{verbatim}
\item 這段程式執行後，會印出什麼？
\begin{verbatim}
import numpy as np

passed = np.array([True, True, False, True, False])
print(passed.sum())
\end{verbatim}
\end{drills}

\section{pandas：把 CSV 讀成一張表}
\label{sec:pandas-read-select}

份量最重的一段來了。還記得上一章你怎麼把一份 CSV 讀進 Python 的嗎？要先 \texttt{open()}
開檔、用 \texttt{csv.reader()} 或 \texttt{csv.DictReader} 逐行讀、自己判斷每一欄代表
什麼、自己決定要不要把字串轉成數字。pandas 是專門用來處理「表格式資料」的套件，讀 CSV
這件事，它一行就能做完，而且連型別都幫你猜好了。

先看看這一章要用的資料檔長什麼樣子，沿用上一章那份「全班小考成績」，每一列記錄一個
學生的姓名、班級、分數：

\begin{example}{text}{(quiz\_scores.csv) 全班小考成績原始資料}
Name,Class,Score
Alice,501,88
Bob,501,72
Chloe,501,95
David,502,60
Emma,502,78
Frank,502,85
Grace,503,91
Henry,503,55
Ivy,503,68
\end{example}

\begin{code}{python}{pandas 讀取 CSV 與選欄的固定寫法}
import pandas as pd

df = pd.read_csv("filename.csv")
df["column_name"]
\end{code}

慣例上，pandas 用 \texttt{as pd} 縮寫成 \texttt{pd}。讀進來的東西叫做「DataFrame」，\index{DataFrame}
你可以把它想成一整張表格，用 \texttt{df["欄位名稱"]} 就能把某一欄整個選出來看。除了
讀 CSV 之外，也可以不透過檔案，直接把一個 dict 包給 \texttt{pd.DataFrame()} 建立一個
DataFrame，dict 的 key 會變成欄位名稱，value（一個 list）會變成那一欄的內容，例如
\texttt{pd.DataFrame(\{"Score": [80, 95]\})}。

\begin{example}{python}{(read-quiz-scores.py) 讀進小考成績、選出一欄}
import pandas as pd

df = pd.read_csv("quiz_scores.csv")

print(df)
print(df["Score"])
\end{example}

執行後會看到：

\begin{verbatim}
    Name  Class  Score
0  Alice    501     88
1    Bob    501     72
2  Chloe    501     95
3  David    502     60
4   Emma    502     78
5  Frank    502     85
6  Grace    503     91
7  Henry    503     55
8    Ivy    503     68
0    88
1    72
2    95
3    60
4    78
5    85
6    91
7    55
8    68
Name: Score, dtype: int64
\end{verbatim}

比較一下：上一章要開檔、逐行讀、自己判斷欄位順序、自己把字串轉成數字，這裡一行
\texttt{pd.read\_csv()} 全部做完。

\keypoint{印出 DataFrame 時，最左邊那一欄由 0 開始遞增的數字是 pandas 自動幫每一列加上
的「列索引」\index{列索引}，用來標記這一列在表格裡的位置，並不是你資料裡原本就有的內容。另外，
\texttt{Score} 這一欄印出來的最後一行寫著 \texttt{dtype: int64}，這是 pandas 自動猜
出來的型別（整數），代表它已經幫你把字串轉成數字了，不用像上一章那樣自己呼叫
\texttt{int()} 轉型。}

如果一份 CSV 有幾千、幾萬筆資料，你應該不會想把整張表格全部印出來，只是想很快看一眼
資料長什麼樣子。這時候可以用 \texttt{df.head()}\index{df.head()}，預設印出最前面 5 筆，也可以自己指定
筆數，例如 \texttt{df.head(3)} 只印前 3 筆：

\begin{example}{python}{(peek-quiz-scores.py) 只看前幾筆資料}
import pandas as pd

df = pd.read_csv("quiz_scores.csv")
print(df.head(3))
\end{example}

執行後會看到：

\begin{verbatim}
    Name  Class  Score
0  Alice    501     88
1    Bob    501     72
2  Chloe    501     95
\end{verbatim}

再來看看選欄的另一個變化：如果一次想選不只一欄，把好幾個欄位名稱裝進一個 list，再放進
中括號裡就好，例如 \texttt{df[["Name", "Score"]]}\index{df[[...]] 多欄選取}：

\begin{example}{python}{(select-two-columns.py) 一次選出兩欄}
import pandas as pd

df = pd.read_csv("quiz_scores.csv")
print(df[["Name", "Score"]])
\end{example}

\begin{verbatim}
    Name  Score
0  Alice     88
1    Bob     72
2  Chloe     95
3  David     60
4   Emma     78
5  Frank     85
6  Grace     91
7  Henry     55
8    Ivy     68
\end{verbatim}

\keypoint{選一欄用單層中括號（\texttt{df["Score"]}），選好幾欄要把欄位名稱先包成一個
list，再整份放進中括號（\texttt{df[["Name", "Score"]]}），兩層括號一層都不能少。這兩
種寫法拿到的東西也不一樣：選單一欄拿到的是 pandas 另一種資料結構，叫 Series，可以想成
表格裡「一整條」資料；選好幾欄拿到的還是一張表格（DataFrame），只是欄位變少了。這本書
不會深入 Series 本身的操作，這裡只要記得「一層中括號選一欄、兩層中括號選好幾欄」就夠
用了。}

\begin{drills}[label=drills:pandas-read-select]
\item 這段程式執行後，會印出什麼？
\begin{verbatim}
import pandas as pd

data = pd.DataFrame({"Subject": ["Chinese", "Math"], "Score": [80, 95]})
print(data["Score"])
\end{verbatim}
\item 承接本節範例的 \texttt{quiz\_scores.csv}：如果你執行 \texttt{print(df["Name"])}，
印出來的每一行前面除了姓名之外，還會多出一欄由 0 開始遞增的數字，這欄數字代表什麼？
\item 如果你想同時選出 \texttt{quiz\_scores.csv} 的「Name」和「Score」兩欄，不要
「Class」，程式碼該怎麼寫？
\end{drills}

\section{pandas：篩選與聚合，一行做完平均}
\label{sec:pandas-filter-agg}

讀進 DataFrame 之後，接下來要做的事跟上一章一模一樣：篩出某個班級的記錄、算出那個
班級的平均分數。上一章要自己寫 \texttt{if} 判斷符不符合條件、自己維護一個總和變數和
一個計數器變數才能算平均；pandas 把這整套動作濃縮成「篩選」加「聚合」兩步。

\begin{code}{python}{pandas 篩選與聚合的固定寫法}
df[df["column_name"] condition]
df["column_name"].mean()
\end{code}

\texttt{df[條件]} 會篩出符合條件的所有列，回傳一個新的、比較小的 DataFrame；\index{df[條件] 篩選}接在欄位
選取後面的 \texttt{.mean()}、\texttt{.sum()} 則是直接算出那一欄的平均值、總和，不用
自己開變數手動累加。\index{.mean()／.sum()}等一下的範例裡還會用到 \texttt{round()}，這是 Python 內建的函式，
跟 pandas 沒有關係：\texttt{round(數字, 小數位數)} 會把數字四捨五入到指定的小數位數，
例如 \texttt{round(85.333, 1)} 會得到 \texttt{85.3}。\index{round()}

\begin{example}{python}{(filter-quiz-scores.py) 篩出 501 班、算平均分數}
import pandas as pd

df = pd.read_csv("quiz_scores.csv")

# 篩出 501 班的所有記錄
class_501 = df[df["Class"] == 501]
print(class_501)

# 算出 501 班的平均分數
average_501 = class_501["Score"].mean()
print(round(average_501, 1))

# 篩出分數 80 分以上的記錄，並算出筆數
high_scores = df[df["Score"] >= 80]
print(len(high_scores))
\end{example}

執行後會看到：

\begin{verbatim}
    Name  Class  Score
0  Alice    501     88
1    Bob    501     72
2  Chloe    501     95
85.0
4
\end{verbatim}

留意一件事：篩出 80 分以上總共 4 筆，跟你在上一節用 numpy 布林 array 算出來的
\texttt{is\_high.sum()} 結果完全一樣，因為兩邊做的其實是同一件事，只是一個套用在
list 包出來的 array 上，一個套用在 DataFrame 的某一欄上。

\stylepoint{\texttt{df[df["Score"] >= 80]} 篩完直接接 \texttt{.mean()} 或用
\texttt{len()} 算筆數，可以串成一行；上一章的手寫版本要先開一個 \texttt{if} 判斷，
再開一個累加變數、一個計數器變數，跑完迴圈才能算出同樣的結果。實際比較一下行數：上
一章 \texttt{average-quiz-score.py} 光是讀檔案、拆欄位、轉型別、累加、算平均，就用掉
十一行程式邏輯；這裡讀檔、篩選、呼叫 \texttt{.mean()} 三行就做完同樣的事，行數少了
將近四分之三。同一件事，pandas 幫你把好幾行邏輯收進一行呼叫裡。}

\begin{drills}[label=drills:pandas-filter-agg]
\item 這段程式執行後，會印出什麼？
\begin{verbatim}
import pandas as pd

data = pd.DataFrame({"Class": [601, 601, 602], "Score": [70, 90, 100]})
print(data[data["Class"] == 601]["Score"].mean())
\end{verbatim}
\item 要篩出 \texttt{df} 裡 \texttt{"Score"} 欄位大於等於 90 分的所有記錄，程式碼該怎麼寫？
\item 這段程式執行後，會印出什麼？
\begin{verbatim}
import pandas as pd

data = pd.DataFrame({"Score": [55, 62, 78, 91]})
print(len(data[data["Score"] < 60]))
\end{verbatim}
\end{drills}

\section{pandas：多重條件篩選與簡單資料清理}
\label{sec:pandas-multi-filter}

上一節只用了一個條件去篩選，實務上常常需要「同時」符合好幾個條件，例如「501 班而且
分數在 80 分以上」。pandas 支援把好幾個條件串在一起，但寫法跟第 \ref{ch:conditionals} 章教過的
\texttt{and} 不太一樣，這裡要特別小心。

\begin{code}{python}{pandas 多重條件篩選的固定寫法}
df[(df["columnA"] conditionA) & (df["columnB"] conditionB)]
\end{code}

每個條件都要用小括號包起來，中間用 \texttt{\&} 串接，不是第 \ref{ch:conditionals} 章學過的 \texttt{and}：

\begin{example}{python}{(filter-multi-condition.py) 篩出 501 班且分數 80 分以上}
import pandas as pd

df = pd.read_csv("quiz_scores.csv")

result = df[(df["Class"] == 501) & (df["Score"] >= 80)]
print(result)
\end{example}

執行後會看到：

\begin{verbatim}
    Name  Class  Score
0  Alice    501     88
2  Chloe    501     95
\end{verbatim}

還記得第 \ref{ch:conditionals} 章教的 \texttt{and} 嗎？如果你在這裡把 \texttt{\&} 換成 \texttt{and}，一樣
會出錯，因為 \texttt{and} 期待兩邊各自是單一個 \texttt{True}／\texttt{False}，但
\texttt{df["Class"] == 501} 產生的其實是一整欄由好幾個 \texttt{True}／\texttt{False}
組成的資料，不是單一個值。pandas 因此改用 \texttt{\&} 這個逐一元素比對的運算子，效果
類似 \texttt{and}，但套用在整欄資料上。

如果你偷懶把括號省略掉，會發生什麼事？

\begin{example}{python}{(wrong-filter.py) 忘記加括號的多重條件篩選}
import pandas as pd

df = pd.read_csv("quiz_scores.csv")

result = df[df["Class"] == 501 & df["Score"] >= 80]
print(result)
\end{example}

執行後會看到：

\begin{verbatim}
Traceback (most recent call last):
  File "wrong-filter.py", line 5, in <module>
    result = df[df["Class"] == 501 & df["Score"] >= 80]
ValueError: The truth value of a Series is ambiguous. Use a.empty,
a.bool(), a.item(), a.any() or a.all().
\end{verbatim}

這是因為 \texttt{\&} 在 Python 裡的運算順序比 \texttt{==}、\texttt{>=} 這些比較符號
更優先，如果沒有括號，Python 會先算 \texttt{\&} 兩側，跟你原本想要的「先個別完成兩個
比較、再把兩個結果放在一起看」完全不同，才會冒出這個一時看起來莫名其妙的錯誤。

同樣的規則也適用在「符合其中一個條件就好」的情況，只是把 \texttt{\&} 換成
\texttt{|}（或），例如 \texttt{df[(df["Class"] == 501) | (df["Class"] == 502)]} 會篩出
501 班或 502 班的所有記錄，括號一樣不能省略。

\keypoint{pandas 篩選多重條件時，每個條件務必用小括號個別包起來，中間用 \texttt{\&}
（且）或 \texttt{|}（或）串接，不能直接寫 \texttt{and}／\texttt{or}，也不能省略括號。\index{pandas 條件組合（\&／\textbar{}）}
記不住細節沒關係，只要記得「多重條件、pandas、括號」這三件事會一起出現，遇到類似的
\texttt{ValueError} 就知道該回頭檢查括號了。}

除了篩選，資料收集回來的表格常常會有「缺考」「沒填答案」這種缺漏的情況，pandas 讀進來
的時候，這種缺漏值會顯示成 \texttt{NaN}（Not a Number 的縮寫，意思是「這裡沒有一個
正常的數值」）。這裡只做最簡單的認識，示範怎麼把整份表格裡有缺漏的那幾列直接濾掉：

\begin{code}{python}{pandas 濾掉缺漏資料的固定寫法}
df.dropna()
\end{code}

\begin{example}{python}{(dropna-example.py) 濾掉有缺考記錄的資料}
import pandas as pd

makeup_scores = pd.DataFrame({
    "Name": ["Alice", "Bob", "Chloe"],
    "Score": [88, None, 95],
})
print(makeup_scores)

cleaned = makeup_scores.dropna()
print(cleaned)
\end{example}

執行後會看到：

\begin{verbatim}
    Name  Score
0  Alice   88.0
1    Bob    NaN
2  Chloe   95.0
    Name  Score
0  Alice   88.0
2  Chloe   95.0
\end{verbatim}

Bob 那一列因為分數是 \texttt{None}（代表沒有值，第 \ref{ch:lists} 章學過的那個 \texttt{None}），
pandas 把它顯示成 \texttt{NaN}，呼叫 \texttt{dropna()} 之後那一列就整列被拿掉了。\index{NaN／dropna()}你
可能也注意到分數欄位變成 \texttt{88.0}、\texttt{95.0}，多了小數點，這是因為欄位裡只要
出現一個 \texttt{NaN}，pandas 為了讓整欄型別維持一致，會把原本的整數欄位一起轉成
帶小數點的型別，這是 pandas 內部的型別規則，不影響你後續用 \texttt{.mean()} 之類的
方法計算。

pandas 對缺漏值還有更細緻的處理方式，例如用平均值或某個固定值補上缺漏的格子，而不是
整列刪掉，這裡不深入，只要知道 \texttt{dropna()} 能幫你先把明顯不完整的資料濾掉就好。

pandas 還有更多沒碰到的進階功能，例如依類別分組聚合（\texttt{groupby}）、合併好幾張
表格（\texttt{merge}），這裡也不深入探討，這一章你已經學會篩選、算平均、算筆數、同時
套用多個條件、簡單濾掉缺漏資料，這些已經是資料整理最常用到的一批工具，之後真的遇到更
複雜的資料操作需求，你會知道 pandas 也做得到，想深入可以再自己查文件學。

\begin{drills}[label=drills:pandas-multi-filter]
\item 這段程式執行後，會印出什麼？
\begin{verbatim}
import pandas as pd

data = pd.DataFrame({"Class": [601, 601, 602], "Score": [70, 90, 55]})
print(data[(data["Class"] == 601) & (data["Score"] >= 80)])
\end{verbatim}
\item 如果把上一題的條件拿掉括號，寫成
\texttt{data[data["Class"] == 601 \& data["Score"] >= 80]}，會發生什麼事？
\item 這段程式執行後，會印出什麼？
\begin{verbatim}
import pandas as pd

data = pd.DataFrame({"Score": [80, None, 70]})
print(len(data.dropna()))
\end{verbatim}
\end{drills}

\section{matplotlib：畫兩張圖幫這章收尾}
\label{sec:matplotlib-basic}

最後這一段，要做的事很單純：把整理好的資料畫成看得懂的圖，當作這一章的成就感收束。你
已經會用 pandas 篩選、算平均了，matplotlib 要做的事只有一件，把算出來的數字變成一張
圖表。這裡會示範兩種最基本的圖：長條圖，適合比較不同類別之間的差異；折線圖，適合看
數值隨著時間或順序推進的變化趨勢。

\begin{code}{python}{matplotlib 畫長條圖的固定寫法}
import matplotlib.pyplot as plt

plt.bar(x_data, y_data)
plt.show()
\end{code}

慣例上，matplotlib 的繪圖功能用 \texttt{as plt} 縮寫成 \texttt{plt}。
\texttt{plt.bar()} 畫長條圖，把 x 軸、y 軸的資料傳進去；\texttt{plt.show()} 會跳出一個
視窗把圖顯示出來。\index{plt.bar()／plt.show()}

\begin{example}{python}{(class-average-bar-chart.py) 畫出各班平均分數長條圖}
import pandas as pd
import matplotlib.pyplot as plt

df = pd.read_csv("quiz_scores.csv")

# 沿用上一節的篩選＋平均邏輯，分別算出三個班級
avg_501 = df[df["Class"] == 501]["Score"].mean()
avg_502 = df[df["Class"] == 502]["Score"].mean()
avg_503 = df[df["Class"] == 503]["Score"].mean()

classes = ["501", "502", "503"]
averages = [avg_501, avg_502, avg_503]

plt.bar(classes, averages)
plt.xlabel("Class")
plt.ylabel("Average Score")
plt.title("Average Quiz Score by Class")
plt.show()
\end{example}

執行完你會看到一個彈出的視窗，裡面是三根長條，分別代表 501、502、503 班的平均分數
（依照這份資料算出來，大約是 85 分、74 分、71 分），長條愈高代表那個班的平均分數愈高，
一眼就能比較出三個班的差異，不用像前面那樣還得自己讀數字才能比較。

長條圖適合「比較不同類別」，例如剛剛的三個班級；如果你想看的是同一個對象的數值隨著
時間、或某種順序推進的變化趨勢，比較合適的圖表是折線圖。假設 501 班連續考了三次小考，
你想看看整體是進步還是退步：

\begin{code}{python}{matplotlib 畫折線圖的固定寫法}
import matplotlib.pyplot as plt

plt.plot(x_data, y_data)
plt.show()
\end{code}

\begin{example}{python}{(class-501-trend-line-chart.py) 畫出 501 班三次小考的平均分數趨勢}
import matplotlib.pyplot as plt

quiz_numbers = ["Quiz 1", "Quiz 2", "Quiz 3"]
class_501_averages = [78.5, 82.0, 85.0]

plt.plot(quiz_numbers, class_501_averages)
plt.xlabel("Quiz")
plt.ylabel("Average Score")
plt.title("Class 501 Average Score Over Three Quizzes")
plt.show()
\end{example}

執行完你會看到一個彈出的視窗，裡面是一條從左到右往上爬的折線，代表 501 班的平均分數從
第一次小考的 78.5 分，進步到第二次的 82.0 分，再進步到第三次的 85.0 分，整體呈現一路
上升的走勢，一眼就能看出這個趨勢，不用像列表那樣一個一個數字比對。\index{plt.plot()}

\stylepoint{長條圖跟折線圖都是「把一串數字畫成圖」，差別在於資料的性質：長條圖裡的
x 軸通常是彼此獨立、沒有順序關係的類別（像「501 班」「502 班」，換個順序排列也不影響
解讀）；折線圖裡的 x 軸通常本身就有順序或時間關係（像「第一次、第二次、第三次小考」，
順序亂掉整張圖就沒意義了）。選圖表類型之前，先想清楚你的資料是「比較類別」還是「看
變化趨勢」，就知道該用哪一種。}

這裡的座標軸標籤跟標題刻意用英文寫（\texttt{"Class"}、\texttt{"Average Score"}）而
不是中文。實際測試時如果把這幾個字串換成中文，圖上會看到一格一格的方框，那並不是真的
顯示不出中文，是 matplotlib 預設字型沒有包含中文字集的關係，跟程式碼寫得對不對無關。

要讓圖表正常顯示中文，需要額外設定一個支援中文的字型，已經超出這一章「畫一張看得懂的
圖」的範圍，這裡選擇直接用英文標籤，你之後如果真的需要中文圖表，可以再自行查怎麼設定
中文字型。

matplotlib 能客製化的東西非常多，顏色、線型、圖例位置、多張子圖排版都還沒碰到，這裡
只示範最基本的兩種畫圖方式，能滿足「我想很快看一眼資料長什麼樣子」的需求就夠了。

\begin{drills}[label=drills:matplotlib-basic]
\item 如果你想畫「長條圖」，該呼叫 \texttt{plt.} 底下的哪個函式？
\item 下面這行程式碼想畫的是哪一種圖：\texttt{plt.plot(months, sales)}？
\item 如果你想比較五個不同班級的平均分數，比較適合用長條圖還是折線圖？
\end{drills}

回頭看看這一章走過的路：一開始你手上只有第 \ref{ch:csv} 章那套手寫迴圈，讀 CSV、拆欄位、轉型別、
自己維護總和跟計數器，資料一多就顯得又長又容易出錯。

這一章帶你認識了三個現成的套件：
numpy 把「對整批數字做運算」濃縮成一行 array 運算；pandas 把「讀表格、篩資料、算平均」
濃縮成幾行 DataFrame 操作，甚至能同時套用好幾個條件、簡單濾掉缺漏值；matplotlib 則把
整理好的數字變成一眼就看得懂的圖表。這三個套件合起來，把上一章要十幾行才能做完的事，
壓縮成幾行就能做完，這正是「拿別人寫好的工具來用」的價值所在。

再次提醒：這一章的目標
是「認識、不精熟」，你還沒有真的成為資料科學家，這三個套件各自都還有大片沒碰到的功能
（numpy 的多維陣列、pandas 的 \texttt{groupby}／\texttt{merge}、matplotlib 的各種
客製化樣式），但下次工作或生活中遇到「一堆資料要整理」的場合，你會知道有這些工具存在、
大概能幫上什麼忙，這樣就足夠了。

\section*{練習解答}
\addcontentsline{toc}{section}{練習解答}

\paragraph{第~\ref{sec:numpy-vectorize}~節}
\begin{enumerate}
\item \texttt{[ 6 12 18]}：array 裡每個元素分別乘以 3。
\item \texttt{[ 9 18 27]}：兩個 array 對應位置的元素分別相減。
\end{enumerate}

\paragraph{第~\ref{sec:numpy-boolean}~節}
\begin{enumerate}
\item \texttt{[False  True  True  True False]}：分別比較 50、85、62、90、40 是否
大於等於 60。
\item \texttt{3}：把 \texttt{True} 當 1、\texttt{False} 當 0 加總，陣列裡有三個
\texttt{True}。
\end{enumerate}

\paragraph{第~\ref{sec:pandas-read-select}~節}
\begin{enumerate}
\item
\begin{verbatim}
0    80
1    95
Name: Score, dtype: int64
\end{verbatim}
\item 那是 pandas 自動加上的「列索引」（row index），不是資料本身的一部分，只是用來
標記每一列在 DataFrame 裡的位置，方便之後用來定位或參照某一列。
\item \texttt{df[["Name", "Score"]]}。
\end{enumerate}

\paragraph{第~\ref{sec:pandas-filter-agg}~節}
\begin{enumerate}
\item \texttt{80.0}：601 班的兩筆分數是 70 和 90，平均是 80.0。
\item \texttt{df[df["Score"] >= 90]}。
\item \texttt{1}：四筆分數裡只有 55 小於 60。
\end{enumerate}

\paragraph{第~\ref{sec:pandas-multi-filter}~節}
\begin{enumerate}
\item
\begin{verbatim}
   Class  Score
1    601     90
\end{verbatim}
601 班裡只有分數 90 那一列同時符合兩個條件。
\item 會出現 \texttt{ValueError: The truth value of a Series is
ambiguous.}，因為 \texttt{\&} 的運算順序比 \texttt{==}、\texttt{>=} 更優先，沒有
括號時 Python 不會照你原本想要的順序去比較。
\item \texttt{2}：三筆資料裡有一筆分數是 \texttt{None}，\texttt{dropna()} 把那一列
濾掉之後剩下兩筆。
\end{enumerate}

\paragraph{第~\ref{sec:matplotlib-basic}~節}
\begin{enumerate}
\item \texttt{plt.bar()}。
\item 折線圖：\texttt{plt.plot()} 是畫折線圖的函式。
\item 長條圖：五個班級是彼此獨立的類別，適合用長條圖比較差異；折線圖比較適合有順序或
時間關係的資料。
\end{enumerate}

\chapter{除錯與讀懂錯誤訊息}
\label{ch:debugging}

程式碼跑不動、螢幕跳出一片紅字，這件事你在第一章已經親身經歷過兩次：漏了一個引號的時候，
Python 丟出 \texttt{SyntaxError}；把 \texttt{print} 打成 \texttt{pint} 的時候，Python
丟出 \texttt{NameError}。那時候你學到的是一個求生技巧：紅字不用整段從頭讀到尾，先看最後
一行、掃到錯誤類型跟訊息，就抓到大半個線索了。

這一章要做的事，是把那個求生技巧升級成一套真正管用的排查方法，分四段推進：

\begin{enumerate}
\item 先把你的「錯誤字典」擴大，除了已經認識的 \texttt{SyntaxError}、\texttt{NameError}，
再加上五個你接下來寫程式一定會撞見的老朋友：\texttt{TypeError}、\texttt{AttributeError}、
\texttt{ValueError}、\texttt{IndexError}、\texttt{KeyError}，每一種都對應一種具體的犯錯
情境，看到訊息就能立刻猜到問題大概出在哪裡。
\item 回到「讀 traceback」這件事本身，這次的程式碼不再只有一兩行；你已經學過函式，程式碼
裡開始出現函式呼叫函式的情況，traceback 也會跟著變長，這段要教你怎麼在一長串
\texttt{File}、\texttt{line} 裡，順著呼叫鏈找到真正該修的地方。
\item 處理一個更麻煩的情況：程式沒有崩潰、沒有紅字，卻給了你錯的答案。這種「安靜的錯」
比丟例外的錯更難抓，因為沒有現成的錯誤訊息可以讀，你得自己主動出擊，一邊用「把懷疑範圍
砍半」的思路縮小搜尋範圍，一邊用 \texttt{print()} 一步步把程式內部到底發生了什麼事挖出來。
\item 介紹 \texttt{assert}，一個讓你先把自己的假設寫進程式碼裡的工具，如果哪天假設不成立，
程式會馬上告訴你，不用等到很多行程式碼之後才崩潰在一個完全看不出原因的地方；這段也會順便
講清楚一件容易搞混的事：\texttt{assert} 跟你在檔案與例外那一章學過的 \texttt{try/except}，
到底該用哪一個、分工在哪裡。
\end{enumerate}

讀完這一章，紅字對你來說不會再是「完蛋了」的訊號，會變成程式在告訴你「往這裡看」的路標。

在正式開始之前，先花一句話複習：\texttt{SyntaxError} 是 Python 連「讀」你的程式碼這一步
都卡住了，代表程式碼不符合語法規則；\texttt{NameError} 是 Python 找不到一個叫這個名字的
東西，通常是打錯字或忘了先定義。這兩個你已經很熟了，接下來要把你認識的錯誤類型，從兩種
擴充到七種。

\section{型別不對：TypeError}
\label{sec:type-error}

\texttt{TypeError} 是「型別不合」的意思，你叫 Python 做一件事，但它手上這幾個東西的種類
兜不起來，沒辦法照你的要求做。這種錯誤最常出現在兩種情境：把不同型別的東西直接湊在一起做
運算，或是呼叫函式時給的東西跟函式原本期待的對不上。

\begin{code}{python}{TypeError 常見寫法}
text + number        # 字串不能直接跟數字相加
function_name(argument)     # 引數的數量／型別要跟函式定義一致
\end{code}

先看第一種情境，把字串跟數字直接相加：

\begin{example}{python}{(type-error-concat.py) 印出分數報告}
name = "Alice"
score = 87
print(name + "'s score: " + score)
\end{example}

存檔執行後：

\begin{verbatim}
Traceback (most recent call last):
  File "type-error-concat.py", line 3, in <module>
    print(name + "'s score: " + score)
          ~~~~~~~~~~~~~~~~~~~~^~~~~~~
TypeError: can only concatenate str (not "int") to str
\end{verbatim}

\texttt{TypeError} 的意思是「型別不合」；\texttt{can only concatenate str (not "int")
to str} 更具體地說：「字串只能跟字串相接，你這裡想接的卻是一個 \texttt{int}（整數）」。
\texttt{score} 是數字，不是文字，Python 沒辦法把數字直接「接」在字串後面。修法是先用
\texttt{str()} 把數字轉成字串：\texttt{name + "'s score: " + str(score)}。

再看第二種情境，呼叫函式時，給的引數兜不起來：

\begin{example}{python}{(type-error-missing-arg.py) 計算商品總價}
def calc_total_price(unit_price, quantity):
    return unit_price * quantity

print(calc_total_price(35))
\end{example}

\begin{verbatim}
Traceback (most recent call last):
  File "type-error-missing-arg.py", line 4, in <module>
    print(calc_total_price(35))
          ~~~~~~~~~~~~~~~~^^^^
TypeError: calc_total_price() missing 1 required positional argument: 'quantity'
\end{verbatim}

這條訊息很老實：\texttt{calc\_total\_price()} 這個函式定義時要兩個引數，你呼叫的時候只
給了一個，少了 \texttt{quantity}。這類 \texttt{TypeError} 不一定是「數量」給錯，也可能是
「型別」給錯，例如同一個函式，如果呼叫時把折扣率誤打成一段文字（本來該是數字），Python
在函式內部真正要做乘法、減法運算的那一行才會爆出來，錯誤訊息會變成類似
\texttt{unsupported operand type(s) for -: 'int' and 'str'}，意思一樣：兩個型別不合，
沒辦法做這個運算。

\keypoint{看到 \texttt{TypeError}，先問自己兩個問題：「這裡的東西型別對嗎？」「函式要的
引數，我是不是少給了、或給錯種類了？」訊息裡通常會直接點名是哪兩種型別衝突，順著這個線索
去對照程式碼，八九不離十就能找到問題。}

\section{方法不存在：AttributeError}
\label{sec:attribute-error}

\texttt{AttributeError} 是「這個東西身上根本沒有你要呼叫的那個方法」的意思。它常常出現
在你以為一個變數是某種型別，但因為前面某一步，它其實已經變成另一種型別，而新的型別身上沒有
你接下來呼叫的那個方法。

最容易踩到的情況，是第 \ref{ch:lists} 章提過的一個小陷阱：很多會修改清單的方法
（像 \texttt{sort()}、\texttt{append()}、\texttt{insert()}、\texttt{remove()}）預設回傳
的是 \texttt{None}，不是改好的那份清單。如果你把回傳值重新指定給原本的變數，這個變數就會
被名字牌重新貼到 \texttt{None} 這個物件上，之後對它呼叫任何清單方法，都會撞上
\texttt{AttributeError}。

\begin{code}{python}{AttributeError 常見寫法}
variable = variable.sort()     # sort() 回傳 None，所以 variable 現在指向 None
variable.append(new_item)      # 對 None 呼叫 list 的方法會丟出 AttributeError
\end{code}

\begin{example}{python}{(attr-error-names.py) 排序後加入新成員}
names = ["Alice", "Bob", "Charlie"]
names = names.sort()
names.append("Taylor")
\end{example}

\begin{verbatim}
Traceback (most recent call last):
  File "attr-error-names.py", line 3, in <module>
    names.append("Taylor")
    ^^^^^^^^^^^^
AttributeError: 'NoneType' object has no attribute 'append'
\end{verbatim}

\texttt{AttributeError: 'NoneType' object has no attribute 'append'} 這句話拆開來看：
\texttt{'NoneType'} 是 \texttt{None} 這個特殊值的型別名稱，意思是「一個型別是
\texttt{NoneType} 的物件身上，沒有 \texttt{append} 這個屬性（方法）」。

回頭看程式碼：\texttt{names.sort()} 把 \texttt{names} 這份清單原地排序，但它的回傳值是
\texttt{None}；第二行 \texttt{names = names.sort()} 把 \texttt{names} 這張名字牌，從
原本指著清單的物件，改貼到了 \texttt{None} 身上。等到第三行呼叫
\texttt{names.append(...)}，\texttt{names} 早就不是清單了，是 \texttt{None}，自然找不到
\texttt{append} 這個方法。修法很直接：把 \texttt{names = names.sort()} 這一行改成單獨
一行 \texttt{names.sort()}，不要把回傳值重新指定給 \texttt{names}。

\texttt{AttributeError} 第二個常見來源，是方法名稱本身打錯字。這種情況型別完全沒有跑掉，
單純是手指打滑：

\begin{example}{python}{(attr-error-typo.py) 方法名稱打錯字}
scores = [82, 91, 76]
scores.apend(88)
\end{example}

\begin{verbatim}
Traceback (most recent call last):
  File "attr-error-typo.py", line 2, in <module>
    scores.apend(88)
    ^^^^^^^^^^^^
AttributeError: 'list' object has no attribute 'apend'. Did you mean: 'append'?
\end{verbatim}

\texttt{scores} 明明是一份 list，型別完全沒問題，問題出在 \texttt{apend} 這個名字，
list 身上根本沒有叫這個名字的方法，正確拼法是 \texttt{append}。比較新版的 Python 甚至會
貼心地在錯誤訊息最後加上 \texttt{Did you mean: 'append'?}，直接把最可能的正確拼法猜給你
看，這種提示值得養成習慣去讀，往往一眼就能抓到打錯字的地方。

\keypoint{看到 \texttt{AttributeError: 'X' object has no attribute 'Y'}，先確認這個
「\texttt{'X'} 型別」是不是你原本以為的型別：如果型別本身就不對，往回找是哪一行把這個
變數的名字牌改貼到別的物件上了，\texttt{sort()}、\texttt{append()} 這類回傳 \texttt{None}
的方法是最常見的肇事者；如果型別是對的，那就檢查 \texttt{Y}（方法名稱）是不是打錯字，
訊息後面如果帶著 \texttt{Did you mean} 的提示，通常就是答案。}

\stylepoint{懷疑 \texttt{AttributeError} 是型別跑掉造成的，可以在出錯的那一行前面插一行
\texttt{print(type(variable))}，把它目前實際的型別印出來看看，跟你以為的型別是否一致。這個
技巧在本章稍後談「用 print 主動除錯」時還會再用到。}

\section{值不合理：ValueError}
\label{sec:value-error}

\texttt{ValueError} 早在第 \ref{ch:variables} 章型別轉換時就出現過：型別轉換函式如果拿到一個內容不合理、
沒辦法真的轉換成目標型別的值，就會丟出 \texttt{ValueError}。跟 \texttt{TypeError} 不
一樣的地方在於：\texttt{TypeError} 是型別本身就對不上（例如硬要把清單當數字加），
\texttt{ValueError} 是型別本身沒問題（都是字串），只是這個字串裡裝的內容轉換不了。

\begin{code}{python}{ValueError 常見寫法}
int(text)      # 如果文字不是合法的整數格式就會出錯
float(text)    # 如果文字不是合法的數字格式就會出錯
\end{code}

\begin{example}{python}{(value-error-age.py) 讀取問卷年齡}
age_text = "twenty-five"
age = int(age_text)
print(age)
\end{example}

\begin{verbatim}
Traceback (most recent call last):
  File "value-error-age.py", line 2, in <module>
    age = int(age_text)
ValueError: invalid literal for int() with base 10: 'twenty-five'
\end{verbatim}

\texttt{invalid literal for int() with base 10: 'twenty-five'} 這句話在說：
「twenty-five」這個字串，用十進位的規則去解讀，不是一個合法的整數寫法。\texttt{int()}
只認得阿拉伯數字組成的字串（像 \texttt{"25"}），拼出來的英文數字詞「twenty-five」對它
來說只是一串它看不懂的符號。這種錯誤特別容易發生在處理使用者輸入或外部資料的時候，你以為
對方會乖乖打數字，實際上拿到的內容五花八門。

這正是 EAFP 精神該派上用場的地方：與其在轉換之前先寫一堆判斷式，檢查這段文字是不是「看起來
像數字」，不如直接呼叫 \texttt{int()}，用第 \ref{ch:files-exceptions} 章學過的 \texttt{try/except} 接住
\texttt{ValueError} 就好：

\begin{example}{python}{(value-error-try.py) 用 try/except 接住轉換失敗}
age_text = "twenty-five"
try:
    age = int(age_text)
except ValueError:
    print("Could not parse the age, defaulting to age 0")
    age = 0
print(age)
\end{example}

\iosample{}{Could not parse the age, defaulting to age 0 \\
0 \\}

\texttt{float()} 遇到轉換不了的內容，一樣會丟出 \texttt{ValueError}，只是訊息的措辭跟
\texttt{int()} 不太一樣：

\begin{verbatim}
Traceback (most recent call last):
  File "value-error-float.py", line 2, in <module>
    price = float(price_text)
ValueError: could not convert string to float: 'fifteen dollars'
\end{verbatim}

\texttt{int()} 的訊息習慣說 \texttt{invalid literal for int() with base 10}，
\texttt{float()} 則是 \texttt{could not convert string to float}，兩句話字面不同，但
意思一樣：手上這段字串，不是這個型別轉換函式看得懂的格式。習慣了 \texttt{int()} 的訊息
之後，看到 \texttt{float()} 的版本也不用重新學一遍，抓到「\texttt{ValueError} 加上
轉換函式的名字」這個共同的形狀就夠了。

\keypoint{\texttt{TypeError} 跟 \texttt{ValueError} 很容易搞混，分辨的訣竅是問：型別
本身對不對？如果型別本身就用錯了（例如拿一個清單當字典的鍵、字串硬要加數字），是
\texttt{TypeError}；如果型別沒問題，只是裡面的內容不合理、轉換不出你要的結果（例如
\texttt{int("twenty-five")}、\texttt{float("fifteen dollars")}），是 \texttt{ValueError}。}

\section{超出範圍：IndexError}
\label{sec:index-error}

\texttt{IndexError} 專屬於 list（還有字串）這類「有順序、可以用位置去拿東西」的資料，
你要的位置根本不存在，list 沒那麼長。這種錯誤特別容易發生在你「以為」清單長度是固定的，
但其實它是會變的（例如根據使用者輸入而長短不一）。

\begin{code}{python}{IndexError 常見寫法}
my_list[index]     # 索引必須小於清單目前的長度，否則會出錯
\end{code}

\begin{example}{python}{(index-error-survey.py) 讀取問卷回答}
survey_answers = ["Very satisfied", "Neutral"]
print("4th respondent's answer: " + survey_answers[3])
\end{example}

\begin{verbatim}
Traceback (most recent call last):
  File "index-error-survey.py", line 2, in <module>
    print("4th respondent's answer: " + survey_answers[3])
                                        ~~~~~~~~~~~~~~^^^
IndexError: list index out of range
\end{verbatim}

\texttt{survey\_answers} 只收到兩份回答，索引只有 \texttt{0} 跟 \texttt{1} 可以用，程式
裡卻寫死了 \texttt{[3]}，以為第四位受訪者一定存在。這是很典型的「你的程式碼假設了一個
固定長度，但實際資料沒那麼長」的情況，寫死的索引數字，正是最該懷疑的地方。

另一種常見的寫法是迴圈裡的位置算錯，例如故意讓迴圈多跑一輪：

\begin{verbatim}
weekly_expenses = [120, 340, 95]
for i in range(1, len(weekly_expenses) + 1):
    print(weekly_expenses[i])
\end{verbatim}

這段程式碼會先正常印出 \texttt{340}、\texttt{95}，跑到最後一輪 \texttt{i} 變成 \texttt{3}
時才爆出 \texttt{IndexError}。\texttt{weekly\_expenses} 只有三個元素，合法的索引是
\texttt{0}、\texttt{1}、\texttt{2}，並沒有 \texttt{3}。

\stylepoint{看到 \texttt{IndexError}，先檢查兩件事：這個清單的長度是固定的還是會變的？
如果是會變的，索引是不是被你寫死了；如果是用 \texttt{range()} 算出來的，範圍是不是多算
了一輪。}

\section{找不到 key：KeyError}
\label{sec:key-error}

dict 用「鍵（key）」去查「值（value）」，但如果你查的那個鍵根本沒被放進這個 dict 裡，
Python 不會幫你猜一個答案，會直接丟出 \texttt{KeyError}。

\begin{code}{python}{KeyError 常見寫法}
my_dict[key]     # 如果字典裡沒有這個鍵就會出錯
\end{code}

\begin{example}{python}{(key-error-word-count.py) 統計字出現的次數}
word_counts = {"cat": 3, "dog": 5}
print(word_counts["bird"])
\end{example}

\begin{verbatim}
Traceback (most recent call last):
  File "key-error-word-count.py", line 2, in <module>
    print(word_counts["bird"])
          ~~~~~~~~~~~^^^^^^^^
KeyError: 'bird'
\end{verbatim}

\texttt{KeyError: 'bird'} 直接把找不到的那個鍵印給你看，\texttt{word\_counts} 裡只登記了
\texttt{"cat"} 跟 \texttt{"dog"}，從來沒有出現過 \texttt{"bird"} 這個鍵，查詢自然失敗。

這種「查詢一個東西，但它可能根本不存在」的情境，正是第 \ref{ch:files-exceptions} 章 \texttt{try/except}／
\texttt{dict.get()} 派上用場的地方，如果你已經預期某個鍵可能不存在，直接用那時候學過的
方式接住，會比讓程式在這裡整個崩潰更好。這一章後半段會再回頭把這件事講得更完整。

這也呼應第 \ref{ch:files-exceptions} 章教的 EAFP 精神：不用動手前把條件檢查得滴水不漏，先讓程式跑，出錯了再用
這一章教的 traceback 抓出問題、用 \texttt{try/except} 接住就好。

\begin{drills}[label=drills:error-types]
\item 這段程式碼執行後，會是哪一種錯誤類型？
\begin{verbatim}
temperatures = [22, 25, 19, 21]
print("Day 5 temperature: " + str(temperatures[4]))
\end{verbatim}
\item 這段程式碼執行後，會是哪一種錯誤類型？
\begin{verbatim}
menu_prices = {"milk tea": 50, "black tea": 35}
print(menu_prices["green tea"])
\end{verbatim}
\item 這段程式碼執行後，會是哪一種錯誤類型？
\begin{verbatim}
total_score = 90
print(total_scroe)
\end{verbatim}
\item 這段程式碼執行後，會是哪一種錯誤類型？
\begin{verbatim}
inventory = {"apple": 10, "banana": 5}
inventory.add("cherry")
\end{verbatim}
\item 這段程式碼執行後，會是哪一種錯誤類型？
\begin{verbatim}
weight_text = "thirty-five"
weight = int(weight_text)
\end{verbatim}
\end{drills}

\section{讀懂變長的 traceback：多層函式呼叫}
\label{sec:call-stack}

前面幾個範例的程式都只有短短兩三行，traceback 裡自然只會出現一組
\texttt{File}／\texttt{line}。但你已經學過函式，程式碼開始出現函式呼叫函式的情況，一個
函式裡面又呼叫了另一個函式，這時候如果錯誤發生在最裡面那一層，traceback 會把整條呼叫鏈
都攤開給你看，不只一組 \texttt{File}／\texttt{line}，會有好幾組。

\begin{example}{python}{(call-stack-average.py) 印出成績報告}
def calculate_average(scores):
    return sum(scores) / len(scores)

def print_score_report(student_name, scores):
    avg = calculate_average(scores)
    print(f"{student_name}'s average score: {avg}")

print_score_report("Alice", [])
\end{example}

\begin{verbatim}
Traceback (most recent call last):
  File "call-stack-average.py", line 8, in <module>
    print_score_report("Alice", [])
    ~~~~~~~~~~~~~~~~~~^^^^^^^^^^^^^
  File "call-stack-average.py", line 5, in print_score_report
    avg = calculate_average(scores)
  File "call-stack-average.py", line 2, in calculate_average
    return sum(scores) / len(scores)
           ~~~~~~~~~~~~^~~~~~~~~~~~~
ZeroDivisionError: division by zero
\end{verbatim}

跟前面幾種錯誤一樣，先看最後一行：\texttt{ZeroDivisionError} 是除數為 0 時，Python
丟出的錯誤。\index{ZeroDivisionError}這一次 traceback 有三組 \texttt{File}：

\begin{itemize}
\item 最下面一組（\texttt{line 2}，在 \texttt{calculate\_average} 裡）才是真正出錯的
地方，\texttt{sum(scores) / len(scores)} 這一行，除數 \texttt{len(scores)} 剛好是
\texttt{0}，因為傳進來的 \texttt{scores} 是一份空清單。
\item 中間一組（\texttt{line 5}，在 \texttt{print\_score\_report} 裡）告訴你：
\texttt{calculate\_average} 這次是被 \texttt{print\_score\_report} 呼叫的，而且是在它
的第 5 行呼叫的。
\item 最上面一組（\texttt{line 8}，在 \texttt{<module>}，也就是程式最外層）告訴你這整條
呼叫鏈最初是從哪裡發動的：\texttt{print\_score\_report("Alice", [])} 這一行，傳進去的
成績清單根本是空的。
\end{itemize}

順著這條鏈子往回看，你會發現真正該修的不只是「除法會除以 0」這件事本身，還要往回問
「這個空清單是從哪裡來的？」，也許是 \texttt{print\_score\_report} 應該在呼叫
\texttt{calculate\_average} 之前，先確認清單不是空的。

\keypoint{traceback 由上到下，其實是呼叫發生的順序：最上面是最早發生的呼叫，最下面才是
真正出錯的那一行；讀的時候還是先看最下面的錯誤類型跟訊息抓大方向，再回頭用中間這幾組
\texttt{File}／\texttt{line} 定位是「哪一層」呼叫出的問題，兩件事不衝突，只是這次要多做
一步：確認是呼叫鏈裡的哪一層發動的。}

\begin{drills}[label=drills:call-stack]
\item 有一段程式，執行後出現以下 traceback：
\begin{verbatim}
Traceback (most recent call last):
  File "cart.py", line 8, in <module>
    print_shopping_summary(["apple", "banana"], [])
    ~~~~~~~~~~~~~~~~~~~~~~^^^^^^^^^^^^^^^^^^^^^^^^^
  File "cart.py", line 5, in print_shopping_summary
    avg_price = get_average_price(prices)
  File "cart.py", line 2, in get_average_price
    return sum(prices) / len(prices)
           ~~~~~~~~~~~~^~~~~~~~~~~~~
ZeroDivisionError: division by zero
\end{verbatim}
真正執行到出錯那一刻的是哪一個檔案、哪一行、哪個函式？這次呼叫又是從程式的哪裡發動的、
傳進去的是什麼？
\end{drills}

\section{程式很長、不知道從哪裡開始查：用二分法縮小範圍}
\label{sec:bisect-debug}

traceback 很好用的前提是程式真的丟出了例外。但接下來要面對的情況更常見：程式一路執行到
結束，沒有任何紅字，答案卻是錯的。這種時候你連 traceback 這個路標都沒有，只能自己動手找
問題。

當程式只有三五行，用眼睛掃過去可能就找到問題了；但如果整個函式有十幾行、甚至呼叫了
好幾個其他函式，從頭到尾一行一行看效率很低。這裡有一個好用的思路，叫二分法：

\begin{itemize}
\item 在程式執行過程的正中間找一個點，檢查那個時間點的狀態對不對。
\item 如果中間就已經不對了，問題出在前半段；如果中間是對的，問題出在後半段。
\item 每檢查一次，就能把懷疑範圍砍半，比從頭到尾每一行都檢查快得多。
\end{itemize}

\begin{example}{python}{(bisect-temps-buggy.py) 計算炎熱日的平均氣溫}
def analyze_temperatures(temps):
    hot_days = []
    for t in temps:
        if t > 30:
            hot_days.append(t)
    total = 0
    for t in hot_days:
        total = t
    average = total / len(hot_days)
    return average

temps = [28, 32, 35, 29, 31, 33]
print(analyze_temperatures(temps))
\end{example}

\iosample{}{8.25 \\}

六筆氣溫裡，超過 30 度的有 32、35、31、33 四筆，這四筆的平均應該是 32.75，但程式印出
8.25。這個函式做了兩件事：先過濾出炎熱日，再把它們加總平均，寫在兩段各自獨立的迴圈裡。如果
從頭到尾一行一行讀，要在兩段迴圈之間來回對照並不輕鬆，這時候正好可以用二分法，從中間切一刀：
在兩段迴圈之間插入 \texttt{print()}，檢查篩選出來的 \texttt{hot\_days} 對不對：

\begin{example}{python}{(bisect-temps-print.py) 從中間插入 print 檢查}
def analyze_temperatures(temps):
    hot_days = []
    for t in temps:
        if t > 30:
            hot_days.append(t)
    print(f"hot_days = {hot_days}")
    total = 0
    for t in hot_days:
        total = t
    print(f"total = {total}")
    average = total / len(hot_days)
    return average

temps = [28, 32, 35, 29, 31, 33]
print(analyze_temperatures(temps))
\end{example}

\iosample{}{hot\_days = [32, 35, 31, 33] \\
total = 33 \\
8.25 \\}

第一行印出來的 \texttt{hot\_days = [32, 35, 31, 33]} 完全正確，代表問題不在前半段的
篩選迴圈，範圍一口氣砍掉一半，只剩下加總的那段迴圈值得懷疑。再看第二行：
\texttt{total = 33}，這只是 \texttt{hot\_days} 裡最後一個元素，不是四個數字加起來的
總和，問題就在這裡。往回看程式碼，第二段迴圈裡寫的是 \texttt{total = t}，每一輪都把
\texttt{total} 整個換掉，跟第 \ref{ch:loops} 章學過的累加寫法對照一下就能看出來，這裡少了一個加號，
應該寫成 \texttt{total += t} 才對：

\begin{verbatim}
def analyze_temperatures(temps):
    hot_days = []
    for t in temps:
        if t > 30:
            hot_days.append(t)
    total = 0
    for t in hot_days:
        total += t
    average = total / len(hot_days)
    return average
\end{verbatim}

改完之後重新執行，會正確印出 \texttt{32.75}。

\keypoint{二分法除錯的精神不是每次都要剛好從程式正中間切一刀，是找一個「這裡的結果我可以
判斷對不對」的中繼點，插入 \texttt{print()} 確認狀態，再依照結果決定往前半段還是後半段
繼續縮小範圍。函式裡如果有好幾個各自獨立的步驟（像這裡的「先篩選、再加總」），步驟跟步驟
之間往往就是最自然的切點。}

\begin{drills}[label=drills:bisect-debug]
\item 一個函式做了三件事：先把清單裡的數字都乘以 2、再篩選出偶數、最後加總。執行後總和
不對，你會在哪兩個位置插入 \texttt{print()} 來縮小範圍？各自預期看到什麼？
\end{drills}

\section{程式沒有崩潰，但結果不對：用 print 主動除錯}
\label{sec:print-debug}

前面幾段的錯誤都會讓程式崩潰、直接把線索（traceback）遞到你手上。但很多時候程式跑得好好
的、沒有半個紅字，答案卻是錯的。這種「安靜的錯」比丟例外的錯更麻煩，因為沒有現成的錯誤
訊息可以讀，你得自己想辦法讓程式「說話」。

最直接的做法，就是在你懷疑出問題的地方插入 \texttt{print()}，把中間過程的變數值印出來，
對照「你以為它應該是什麼」跟「它實際上是什麼」。懷疑範圍如果很大，上一節教的二分法可以先
幫你把範圍砍到剩下一兩行，這一節接著處理縮小範圍之後，該怎麼真正用 \texttt{print()}
挖出問題。

\begin{code}{python}{插入 print 除錯}
print("debug:", variable_name)     # 在你懷疑有問題的地方印出目前的值
\end{code}

先看一段有問題的程式：它想算出一批成績裡有幾位及格，但輸出的數字明顯不對。

\begin{example}{python}{(debug-passing-count.py) 統計及格人數}
def count_passing_students(scores, passing_score):
    passing_count = 0
    for score in scores:
        passing_count = 0
        if score >= passing_score:
            passing_count += 1
    return passing_count

scores = [55, 78, 92, 60, 88]
print(count_passing_students(scores, 60))
\end{example}

\iosample{}{1 \\}

五筆成績裡有 78、92、60、88 四筆及格，答案應該是 \texttt{4}，但程式印出 \texttt{1}。
沒有任何紅字，程式「正常」執行完畢，只是答案是錯的，這正是最需要主動出擊的情況。把懷疑的
地方（迴圈每一輪 \texttt{passing\_count} 的變化）插上 \texttt{print()}：

\begin{example}{python}{(debug-passing-count-print.py) 加上 print 觀察每一輪}
def count_passing_students(scores, passing_score):
    passing_count = 0
    for score in scores:
        passing_count = 0
        print(f"Checking {score}, current passing_count = {passing_count}")
        if score >= passing_score:
            passing_count += 1
    return passing_count

scores = [55, 78, 92, 60, 88]
print(count_passing_students(scores, 60))
\end{example}

\iosample{}{Checking 55, current passing\_count = 0 \\
Checking 78, current passing\_count = 0 \\
Checking 92, current passing\_count = 0 \\
Checking 60, current passing\_count = 0 \\
Checking 88, current passing\_count = 0 \\
1 \\}

問題一眼就看出來了：每一輪 \texttt{passing\_count} 都在印出的當下是 \texttt{0}，它從來
沒有把上一輪累加的結果留住，因為 \texttt{passing\_count = 0} 這一行被放進了 \texttt{for}
迴圈裡面，每一輪都重新歸零一次。修法是把這一行搬到迴圈外面，只在一開始歸零一次：

\begin{verbatim}
def count_passing_students(scores, passing_score):
    passing_count = 0
    for score in scores:
        if score >= passing_score:
            passing_count += 1
    return passing_count
\end{verbatim}

改完之後重新執行，答案就會正確印出 \texttt{4}。

\subsection*{寫出有意義的 print 除錯訊息}

上面這個例子只印一個變數，一眼就能看懂哪個數字對應哪個東西。但除錯用的 \texttt{print()}
一多，如果每一行都只印裸的數字，很快就會分不清楚哪一個數字是哪個變數、是哪一輪印出來的。
比較一下同一段程式，兩種不同的除錯訊息寫法：

\begin{example}{python}{(print-msg-bad.py) 沒有標籤的除錯訊息}
def calculate_average_spending(receipts):
    total = 0
    count = 0
    for amount in receipts:
        print(amount)
        total += amount
        count += 1
    return total / count

receipts = [320, 150, 480]
print(calculate_average_spending(receipts))
\end{example}

\iosample{}{320 \\
150 \\
480 \\
316.6666666666667 \\}

這四行輸出裡，前三行是除錯用的 \texttt{print(amount)}，最後一行才是函式真正的回傳值，
但光看數字完全分不出哪個是哪個，如果程式裡同時有好幾個變數在印，更容易看到眼花。改成幫
每個值加上標籤，說清楚「這是什麼」跟「當下發生在哪個階段」：

\begin{example}{python}{(print-msg-good.py) 加上標籤的除錯訊息}
def calculate_average_spending(receipts):
    total = 0
    count = 0
    for amount in receipts:
        print(f"amount = {amount}, running total = {total}")
        total += amount
        count += 1
    return total / count

receipts = [320, 150, 480]
print(calculate_average_spending(receipts))
\end{example}

\iosample{}{amount = 320, running total = 0 \\
amount = 150, running total = 320 \\
amount = 480, running total = 470 \\
316.6666666666667 \\}

用 f-string 幫每個數字加上「這是什麼變數」「這是哪一輪」的標籤之後，就算同時印好幾個
變數、印好幾輪，也能一眼看出每個數字的來歷，不用回頭去對照程式碼第幾行才知道印出來的是
哪個東西。除錯訊息值不值得多打幾個字，差別往往就在這裡。

\stylepoint{懷疑某個變數的型別跑掉（像是前面 \texttt{AttributeError} 那一節提到的
情況），除了印值本身，也可以順便印 \texttt{type(variable)}，一次確認「值對不對」跟「型別
對不對」兩件事，不用等到真的撞上 \texttt{AttributeError} 才回頭查。}

\keypoint{\texttt{print()} 除錯法用完之後，記得把除錯用的那幾行刪掉或註解掉，不然它們會
跟正式的輸出混在一起，讓別人（或未來的你）看不懂哪些是「真正的結果」、哪些是「除錯的
碎碎念」。之後如果程式規模變大，會有更適合大量除錯訊息的專門工具，這本書用
\texttt{print()} 就夠了，它門檻最低，也是你從第一章就已經會用的工具。}

\begin{drills}[label=drills:print-debug]
\item 下面這段程式想算出購物車的總金額，但輸出結果不對：
\begin{verbatim}
def calculate_cart_total(items):
    total = 0
    for item_name, price, quantity in items:
        total = price * quantity
    return total

cart = [("apple", 20, 3), ("banana", 15, 2)]
print(calculate_cart_total(cart))
\end{verbatim}
不要真的動手執行，用紙上推演的方式回答：你會在哪幾個地方插入 \texttt{print()}？你預期
會看到什麼樣的輸出，藉此發現問題出在哪裡？
\item 下面這段程式印出來的除錯訊息完全沒有標籤，難以判讀：
\begin{verbatim}
def find_max_price(prices):
    max_price = 0
    for p in prices:
        print(p)
        if p > max_price:
            max_price = p
    return max_price
\end{verbatim}
把 \texttt{print(p)} 改寫成一行更有意義的除錯訊息，讓輸出清楚說明這個數字是目前這一輪
的 \texttt{p}，還是目前為止的 \texttt{max\_price}。
\end{drills}

\section{assert：把假設寫進程式碼裡}
\label{sec:assert}

寫程式的時候，你心裡常常會有一些「這裡應該一定會怎樣」的假設，例如一個算折扣後金額的
函式，你會假設傳進來的價格一定是正數。

如果這個假設有一天不成立，你會希望程式馬上停下來、清楚告訴你「你的假設錯了」，不用等到
一個荒謬的數字（例如負的金額）悄悄流到程式後面幾十行之後，才在一個完全看不出原因的地方
出狀況。\texttt{assert} 就是讓你把這種假設直接寫進程式碼裡的工具。\index{assert}

\begin{code}{python}{assert 語句}
assert condition, "message to show if the condition is False"
\end{code}

\begin{example}{python}{(assert-discount.py) 計算折扣後金額}
def apply_discount(price, discount_rate):
    assert price > 0, "price must be positive"
    return price * (1 - discount_rate)

print(apply_discount(-100, 0.2))
\end{example}

\begin{verbatim}
Traceback (most recent call last):
  File "assert-discount.py", line 5, in <module>
    print(apply_discount(-100, 0.2))
          ~~~~~~~~~~~~~~^^^^^^^^^^^
  File "assert-discount.py", line 2, in apply_discount
    assert price > 0, "price must be positive"
           ^^^^^^^^^
AssertionError: price must be positive
\end{verbatim}

\texttt{assert condition, message} 的意思是：「檢查這個條件，如果條件是 \texttt{True}，什麼事都
不會發生，程式繼續往下跑；如果條件是 \texttt{False}，馬上丟出
\texttt{AssertionError}，並把後面那段訊息顯示出來」。\index{AssertionError}

這裡呼叫 \texttt{apply\_discount(-100, 0.2)}，\texttt{price} 是 \texttt{-100}，不滿足
\texttt{price > 0}，assert 就當場攔了下來。如果改成呼叫
\texttt{apply\_discount(500, 0.2)}，因為 \texttt{500 > 0} 成立，assert 完全不會發出任何
聲音，程式照樣算出 \texttt{400.0} 繼續執行，你平常根本不會注意到 assert 的存在，只有
假設被打破的那一刻，它才會跳出來說話。

\stylepoint{assert 失敗的時候，該做的事是回頭檢查「我的假設哪裡想錯了」，不是把那行
assert 刪掉讓它閉嘴。assert 會失敗，通常代表程式在更早的地方就已經做錯事了，把它刪掉只
是把問題藏起來，不是解決問題。}

\keypoint{\texttt{assert} 跟你在第 \ref{ch:files-exceptions} 章學過的 \texttt{try/except}，聽起來都是在「處理
可能出錯的情況」，但分工其實很不一樣：
\begin{itemize}
\item \texttt{try/except}（也就是先做、錯了再處理的 EAFP 精神）適合用在「這件事本來就
可能因為外部因素失敗，而且失敗是你意料中、程式應該要能撐住繼續跑」的情境，例如檔案不
存在、使用者亂打、字典裡查不到那個鍵（就像第 \ref{sec:key-error}~節看到的
\texttt{KeyError}）。
\item \texttt{assert} 則適合用在「這裡假設一定要成立」的情境：如果不成立，代表程式本身
某個地方寫錯了，應該讓程式當場停下來，不該用 \texttt{try/except} 把
\texttt{AssertionError} 接住、假裝沒事繼續跑，因為那等於是把一個已經知道有錯的程式，
硬撐著繼續執行下去。
\end{itemize}}

\subsection*{assert 的限制：兩個容易踩的坑}

\texttt{assert} 好用，但有兩個地方特別容易讓新手誤用，需要先講清楚。

第一個是語法陷阱：把條件跟訊息用括號包在一起，變成一個 tuple。

\begin{code}{python}{assert 常見的坑：條件誤包成 tuple}
assert (condition, "message")   # 危險：這個括號會讓它變成 tuple，永遠是 True
\end{code}

\begin{example}{python}{(assert-tuple-trap.py) 誤把 assert 條件包成 tuple}
def apply_discount(price, discount_rate):
    assert (price > 0, "price must be positive")
    return price * (1 - discount_rate)

print(apply_discount(-100, 0.2))
\end{example}

\begin{verbatim}
assert-tuple-trap.py:2: SyntaxWarning: assertion is always true, perhaps remove parentheses?
  assert (price > 0, "price must be positive")
-80.0
\end{verbatim}

這裡完全沒有出現 \texttt{AssertionError}，價格是 \texttt{-100} 這個明顯不合理的數字，
照樣被拿去算出 \texttt{-80.0} 印出來，assert 形同虛設。原因藏在括號裡：
\texttt{(price > 0, "price must be positive")} 這一整包，因為外層多包了一層括號、中間又用逗號
隔開，被 Python 解讀成一個裝了兩個元素的 tuple（第 \ref{ch:tuples} 章學過的容器型別）。

根據第 \ref{ch:conditionals} 章教過的 truthy/falsy 規則，非空的容器一律當作 \texttt{True}，這個 tuple 裡有兩個
元素，不管 \texttt{price} 到底是多少，它永遠是非空的，所以這個 assert 不管 \texttt{price}
是多少，永遠通過檢查。Python 自己也發現了這個常見的坑，執行時印出的
\texttt{SyntaxWarning: assertion is always true, perhaps remove parentheses?} 就是在
提醒你。

正確寫法是拿掉外層括號，寫回原本 \texttt{assert condition, message} 的樣子，讓 Python
讀成兩個獨立的東西，不是一整包 tuple。

\keypoint{\texttt{assert condition, message} 中間那個逗號，是 assert 語法本身的一部分，不是用來
把兩者包進同一個 tuple。只要在條件外面多加一層括號、把訊息也包進去，就會不小心造出一個
恆真的 tuple，讓 assert 失去作用。如果執行程式時看到
\texttt{SyntaxWarning: assertion is always true}，幾乎可以直接斷定就是這個坑，把外層
括號拿掉即可。}

第二個限制更根本：assert 是可以被整個關掉的。Python 執行程式時如果加上 \texttt{-O} 這個
選項（例如 \texttt{python -O program.py}），所有的 assert 陳述句都會被直接跳過，連檢查
都不會檢查，不管條件是真是假。用本節最開始的 \texttt{assert-discount.py} 示範，平常直接
執行會照原本示範的方式丟出 \texttt{AssertionError}；但如果改成加上 \texttt{-O} 選項
執行：

\begin{verbatim}
$ python3 assert-discount.py
Traceback (most recent call last):
  File "assert-discount.py", line 5, in <module>
    print(apply_discount(-100, 0.2))
          ~~~~~~~~~~~~~~^^^^^^^^^^^
  File "assert-discount.py", line 2, in apply_discount
    assert price > 0, "price must be positive"
           ^^^^^^^^^
AssertionError: price must be positive

$ python3 -O assert-discount.py
-80.0
\end{verbatim}

同一段程式、同樣不合理的 \texttt{-100}，加上 \texttt{-O} 執行後，assert 完全沒有發揮
作用，一路算出 \texttt{-80.0}。這代表 assert 只適合拿來抓「開發階段就該抓到、抓到了是
在提醒自己犯了邏輯錯誤」的問題。

像驗證使用者輸入是否合法、確認金額或密碼格式正確這類攸關程式正確性與安全性、無論如何都
必須執行的檢查，該用第 \ref{ch:files-exceptions} 章學過的 \texttt{if} 搭配例外處理來處理，因為你沒辦法保證這段
程式將來執行時，會不會被人不小心（或刻意）加上 \texttt{-O} 執行，讓原本該擋下來的檢查
整個消失。

\keypoint{把 assert 想成寫給自己看的除錯用檢查點，不是程式正式運作時保證一定會執行的
保護機制。開發階段拿來快速抓自己的邏輯錯誤很好用，但任何「不管什麼情況都一定要擋下來」
的檢查（輸入驗證、權限確認），都應該用 \texttt{if}／\texttt{try except} 來寫，不要依賴
assert。}

\begin{drills}[label=drills:assert]
\item 幫下面這個函式加上一個合理的 \texttt{assert}，確保傳進來的邊長一定是正數：
\begin{verbatim}
def calculate_square_area(side_length):
    return side_length ** 2
\end{verbatim}
\item 你要寫一段程式，讀取使用者從鍵盤輸入的年齡（文字），再用 \texttt{int()} 轉成數字。
使用者有可能打了「二十」這種沒辦法轉換的內容，這種情況你會用 \texttt{assert} 還是
\texttt{try/except} 來處理？為什麼？
\item 你寫了一個「取清單中間那個元素」的函式，你自己在程式裡已經確保呼叫這個函式之前，
清單一定不會是空的。如果不小心真的傳了一份空清單進來，代表的是使用者輸入有問題，還是
你自己程式邏輯哪裡寫錯了？這種情況比較適合用 \texttt{assert} 還是 \texttt{try/except}？
\item 你在程式裡看到這一行 assert，執行後卻發現不管條件寫得多離譜，好像都沒有真的擋下
任何東西：
\begin{verbatim}
assert (age >= 0, "age cannot be negative")
\end{verbatim}
這裡的問題出在哪裡？該怎麼修正？
\end{drills}

紅字曾經是你最怕看到的東西，但走到這裡，你已經有一整套面對它的方法：

\begin{itemize}
\item 認得七種最常見的錯誤類型。
\item 知道怎麼順著變長的 traceback 找到真正該修的地方。
\item 遇到不報錯但答案不對的「安靜的錯」，知道怎麼用二分法縮小範圍、用有標籤的 print
主動挖出線索。
\item 學會用 assert 把自己的假設寫清楚，也知道它的限制在哪裡。
\end{itemize}

下次程式又跳出一片紅字，那不是「完蛋了」的訊號，會變成程式在告訴你「往這裡看」的提醒。

\section*{練習解答}
\addcontentsline{toc}{section}{練習解答}

\paragraph{第~\ref{sec:key-error}~節}
\begin{enumerate}
\item \texttt{IndexError}，\texttt{temperatures} 只有四筆資料，合法索引是
\texttt{0}～\texttt{3}，\texttt{temperatures[4]} 超出範圍。
\item \texttt{KeyError}，\texttt{menu\_prices} 裡沒有 \texttt{"green tea"} 這個鍵。
\item \texttt{NameError}，變數叫 \texttt{total\_score}，但查詢時打成了
\texttt{total\_scroe}，兩個拼法對不起來。
\item \texttt{AttributeError}，dict 沒有 \texttt{add()} 這個方法（\texttt{add()} 是
set 才有的方法），\texttt{inventory.add("cherry")} 會因為找不到這個方法而出錯。
\item \texttt{ValueError}，\texttt{"thirty-five"} 不是 \texttt{int()} 認得的阿拉伯數字格式，
轉換會失敗。
\end{enumerate}

\paragraph{第~\ref{sec:call-stack}~節}
\begin{enumerate}
\item 真正執行到出錯那一刻的是 \texttt{cart.py} 第 2 行，在 \texttt{get\_average\_price}
函式裡的 \texttt{sum(prices) / len(prices)}，因為 \texttt{prices} 是空清單、除數是
\texttt{0}。這次呼叫是從 \texttt{cart.py} 第 5 行、\texttt{print\_shopping\_summary}
函式裡的 \texttt{get\_average\_price(prices)} 發動的；而 \texttt{print\_shopping\_summary}
本身又是被最外層第 8 行的呼叫（傳進了一份空的價錢清單）觸發的，真正該回頭檢查的是
「這份價錢清單為什麼是空的」。
\end{enumerate}

\paragraph{第~\ref{sec:bisect-debug}~節}
\begin{enumerate}
\item 可以先在「乘以 2」跟「篩選偶數」之間插入 \texttt{print()}，檢查乘以 2 之後的清單
內容對不對；如果這裡是對的，再往後移，改成在「篩選偶數」跟「加總」之間插入
\texttt{print()}，檢查篩選出來的偶數清單對不對。哪一個中繼點先出現不合預期的內容，問題
就落在那一步驟；如果兩個中繼點看起來都對，問題就落在最後的加總那一段。
\end{enumerate}

\paragraph{第~\ref{sec:print-debug}~節}
\begin{enumerate}
\item 可以在 \texttt{for} 迴圈裡面、\texttt{total = price * quantity} 這一行前後插入
\texttt{print()}，把每一輪的 \texttt{item\_name}、\texttt{price}、\texttt{quantity}
跟當下的 \texttt{total} 印出來。預期會看到 \texttt{total} 每一輪都被整個換掉（例如第一輪
是 \texttt{60}，第二輪變成 \texttt{30}），不是越加越多，這樣就能發現問題出在
\texttt{total = price * quantity} 應該改成 \texttt{total += price * quantity}，不然
每一輪都會把前面算好的總額蓋掉，只留下最後一項商品的金額。
\item 可以把 \texttt{print(p)} 改成
\texttt{print(f"p = \{p\}，目前 max\_price = \{max\_price\}")}，這樣每一行印出來的
內容，會同時秀出「這一輪拿到的價格」跟「目前為止看過的最大值」，不用回頭對照程式碼就能
看懂每個數字代表什麼。
\end{enumerate}

\paragraph{第~\ref{sec:assert}~節}
\begin{enumerate}
\item 例如 \texttt{assert side\_length > 0, "side\_length must be positive"}，加在函式的第一行。
\item 用 \texttt{try/except}，使用者打了什麼是你控制不了的外部因素，這種「本來就可能會
失敗、失敗了程式該撐住繼續跑」的情境正是 EAFP 的主場，用
\texttt{try: ... except ValueError: ...} 接住轉換失敗的情況即可。
\item 代表你自己程式邏輯哪裡寫錯了（或呼叫的地方沒有真的做到你以為的保證），這種「照
理說一定成立」的假設被打破，適合用 \texttt{assert}，讓程式當場停下來提醒你，不要悄悄用
\texttt{try/except} 把錯誤吞掉、假裝沒事。
\item 問題出在 \texttt{(age >= 0, "age cannot be negative")} 外層多包了一層括號，被 Python
解讀成一個裝了兩個元素的 tuple，非空的 tuple 永遠是 \texttt{True}，這個 assert 不管
\texttt{age} 是多少都會通過檢查。修法是拿掉外層括號，改回
\texttt{assert age >= 0, "age cannot be negative"} 的寫法。
\end{enumerate}

\chapter{模組與標準函式庫}
\label{ch:modules}

到目前為止，你寫的每一支程式，用到的功能幾乎都是自己一行一行打出來的。但真實世界寫程式不
是這樣。你不會為了算一個平方根就自己重新推導數學公式，也不會為了抽一個隨機數字就自己
發明一套演算法，這些事情早就有人寫好、測試過，放在那裡等你直接拿來用。

這一章要帶你認識
Python 自己內建的「標準函式庫」，你可以把它想成一整套已經放在你書架上的工具書，電腦裡
本來就有，不用另外安裝（這一點跟第 \ref{ch:data-tools} 章不一樣：那裡教的 \texttt{pip} 是去外面的書店買
書；這裡的書，你打開電腦的那一刻就已經在書架上了），你只需要知道怎麼把對的那本書拿下來
用。

這一章分五段推進：

\begin{itemize}
\item 先搞懂「拿書」這個動作本身怎麼做：\texttt{import} 到底在做什麼，有哪幾種常見的
寫法（第~\ref{sec:import-syntax}~節）。
\item 接著實際去書架上翻幾本你以後會很常用到的書，處理檔案路徑和系統資訊的 \texttt{os}
和 \texttt{sys}（第~\ref{sec:os-and-sys}~節）。
\item 再翻抽隨機數的 \texttt{random} 和處理日期的 \texttt{datetime}（第~
\ref{sec:random-and-datetime}~節）。
\item 接著認識一個處理「整理一批資料」時特別好用的工具，\texttt{collections} 模組裡的
\texttt{Counter}，並簡短認識一下 \texttt{itertools}
（第~\ref{sec:collections-and-itertools}~節）。
\item 最後，我們要把角色反過來：你自己寫的 \texttt{.py} 檔案，其實也可以變成書架上的一
本書，讓別的檔案 \texttt{import} 它、引用裡面的函式（第~\ref{sec:write-your-own-module}
~節）。
\end{itemize}

我們不會把每一本書從頭到尾念給你聽，因為那既不可能，也沒必要，這一章真正的重點，是讓你
養成「不知道就去查官方文件」的習慣，這個習慣比背下每個函式的名字重要得多。

讀完這一章，你不會記得住標準函式
庫裡的每一個函式，沒有人記得住，但你會知道要用的時候去哪裡找、怎麼看懂文件在講什麼，這
個能力比死背 API 更耐用。

\section{先搞懂「拿書」這個動作：import 語法}
\label{sec:import-syntax}

在認識任何一個模組的功能之前，得先搞懂一件更根本的事：你要怎麼跟 Python 說「我要用某
個模組」。這個動作叫做 \texttt{import}，它有幾種常見的寫法，接下來每一節、每一個模組的
範例都會用到這裡教的語法，不會再重講一次，值得在這裡花點時間把差異搞清楚。這一節借用
\texttt{math}，Python 內建的數學模組，來示範，你不需要先知道 \texttt{math} 裡有什麼，
跟著範例做就會自然理解。
\index{import}

第一種寫法，是把整個模組匯入進來，之後每次要用模組裡的東西，都要在前面加上模組的名字：

\begin{code}{python}{Importing an entire module}
import module_name

module_name.member_name
\end{code}

\begin{example}{python}{(circle-area.py) Using import math to compute a circle's area}
import math

radius = 4
area = math.pi * radius ** 2
print(f"The area of a circle with radius {radius} is approximately {area:.2f}")
\end{example}

\iosample{}{The area of a circle with radius 4 is approximately 50.27 \\}

\texttt{import math} 之後，\texttt{math.pi}（圓周率）跟 \texttt{math.sqrt(...)} 之類的
東西都要透過 \texttt{math.} 這個首碼才能用。
\index{math 模組}
第二種寫法，是只精確匯入你要用的那一兩個名
字，之後就可以直接用，不用再加模組名字首碼：

\begin{code}{python}{Importing a single name precisely}
from module_name import member_name

member_name
\end{code}

\begin{example}{python}{(hypotenuse.py) Using from math import sqrt to compute the hypotenuse}
from math import sqrt

a = 3
b = 4
c = sqrt(a ** 2 + b ** 2)
print(f"The hypotenuse length is {c}")
\end{example}

\iosample{}{The hypotenuse length is 5.0 \\}

兩種寫法做的事情本質上一樣，差別只在「要不要打模組名字首碼」。什麼時候選哪一種？

\begin{itemize}
\item 如果你只會用到模組裡一兩個東西，而且那個名字不會跟你自己程式裡的變數撞名，用
\texttt{from...import} 比較精簡好讀。
\item 如果你會用到模組裡很多不同的功能，或者擔心名字撞
在一起搞混，整模組匯入（\texttt{import math}，之後都寫 \texttt{math.sqrt(...)}）會比較
清楚安全，一看到 \texttt{math.} 開頭，就知道這是從 \texttt{math} 模組來的東西，不會誤
以為是你自己定義的函式。
\end{itemize}

\keypoint{\texttt{import math} 這個動作，其實跟你已經很熟的 \texttt{x = 5} 是同一件事，
只是這次貼名字牌的對象換成了一整包別人寫好的程式碼，不再是一個數字。\texttt{import} 把
\texttt{math} 這個模組物件，貼上了一張叫 \texttt{math} 的名字牌，之後你打 \texttt{math}
這三個字，Python 就知道你在說哪一包程式碼。}

這兩種寫法最容易搞混的地方，是「你用哪種方式匯入，就只能用哪種方式呼叫」，混著用會直
接出錯。先看第一種混淆：明明只用 \texttt{import math} 整模組匯入，卻忘了打模組名字首
碼，直接呼叫 \texttt{math} 裡的東西：

\begin{example}{python}{(import-mistake-1.py) Using something without importing it first}
radius = 4
area = math.pi * radius ** 2
print(area)
\end{example}

存檔執行後，你會看到：

\begin{verbatim}
Traceback (most recent call last):
  File "import-mistake-1.py", line 2, in <module>
    area = math.pi * radius ** 2
           ^^^^
NameError: name 'math' is not defined. Did you forget to import 'math'?
\end{verbatim}

\texttt{NameError} 是說 Python 從頭到尾都沒看過 \texttt{math} 這個名字，這支程式從頭
到尾根本沒有寫 \texttt{import math} 這一行，Python 當然不知道 \texttt{math} 是什麼。
新版 Python 其實很貼心，直接在錯誤訊息裡猜了一句「你是不是忘了 import \texttt{math}」，
這次剛好被它猜中了。

再看反過來的混淆：用了 \texttt{from math import sqrt} 精確匯入，卻在呼叫的時候又多打
了一次模組名字首碼：

\begin{example}{python}{(import-mistake-2.py) Using the module prefix after a precise import}
from math import sqrt

result = math.sqrt(16)
print(result)
\end{example}

執行結果同樣是 \texttt{NameError}：

\begin{verbatim}
Traceback (most recent call last):
  File "import-mistake-2.py", line 3, in <module>
    result = math.sqrt(16)
             ^^^^
NameError: name 'math' is not defined. Did you forget to import 'math'?
\end{verbatim}

這裡容易讓人困惑的地方是：明明程式裡有一行 \texttt{import} 相關的敘述，怎麼還是
\texttt{NameError}？原因是 \texttt{from math import sqrt} 只把 \texttt{sqrt} 這一個
名字貼上了名字牌，\texttt{math} 這個名字本身完全沒有被貼過名字牌，這支程式應該直接寫
\texttt{sqrt(16)}，不用、也不能加 \texttt{math.} 首碼。

這兩種寫法之外，還有第三種常見的變化：幫模組或名字另外取一個比較短、比較好打，或是比較
不容易跟你自己變數撞名的名字，這個動作叫做「取別名」，寫法是在 \texttt{import} 敘述後
面加上 \texttt{as 別名}：
\index{import ... as ...（別名）}

\begin{code}{python}{Giving a module or name an alias}
import module_name as alias

alias.member_name

from module_name import member_name as alias

alias
\end{code}

\begin{example}{python}{(dice-roll.py) Using import random as rd for an alias}
import random as rd

roll = rd.randint(1, 6)
print(f"Rolled a {roll}")
\end{example}

因為骰出來的點數每次執行都會不一樣，下面顯示的是作者某一次執行的結果，你自己執行時骰出
的點數幾乎一定不一樣：

\iosample{}{Rolled a 2 \\}

\texttt{import random as rd} 之後，這支程式接下來都改用 \texttt{rd} 稱呼
\texttt{random} 模組，效果跟打完整的 \texttt{random} 一模一樣，只是打字比較少。
\texttt{as} 也可以接在 \texttt{from...import} 後面，把精確匯入的名字換一個名字：

\begin{example}{python}{(hypotenuse-alias.py) Using from math import sqrt as square\_root for an alias}
from math import sqrt as square_root

a = 6
b = 8
c = square_root(a ** 2 + b ** 2)
print(f"The hypotenuse length is {c}")
\end{example}

\iosample{}{The hypotenuse length is 10.0 \\}

\keypoint{\texttt{from...import} 和 \texttt{import...as} 常常被搞混，但它們解決的是兩
件不同的事。\texttt{from module\_name import member\_name} 決定的是「要不要打模組名字
首碼」：把模組裡的某個名字直接搬出來，之後可以省略首碼。\texttt{import module\_name as
alias} 決定的是「這個模組要叫什麼名字」：模組本身還是要透過首碼才能用它裡面的東西，只是
首碼換成你自訂的別名，通常是為了少打幾個字，或是避免跟你自己的變數名字撞在一起。兩者可
以合併使用，像 \texttt{from module\_name import member\_name as alias}，把「精確匯入」
和「取別名」一次做完。}

取了別名之後，如果不小心又用回原本的模組名字，一樣會出錯，道理跟前面兩個 NameError 一
模一樣：

\begin{example}{python}{(alias-mistake.py) Using the original module name after aliasing it}
import random as rd

number = random.randint(1, 6)
print(number)
\end{example}

\begin{verbatim}
Traceback (most recent call last):
  File "alias-mistake.py", line 3, in <module>
    number = random.randint(1, 6)
             ^^^^^^
NameError: name 'random' is not defined. Did you forget to import 'random'?
\end{verbatim}

\texttt{import random as rd} 只把名字牌 \texttt{rd} 貼到 random 模組上，\texttt{random}
這個名字本身，從頭到尾都沒有被貼過名字牌。取別名之後，程式碼裡就只認得那個別名，原本的
模組名字等於不存在。

\begin{drills}[label=drills:import-syntax]
\item 這段程式碼執行後會發生什麼事？為什麼？
\begin{verbatim}
import math
print(sqrt(25))
\end{verbatim}
\item 這段呢？
\begin{verbatim}
from math import pi
print(math.pi)
\end{verbatim}
\item 這段程式碼執行後會印出什麼？
\begin{verbatim}
from math import pow
print(pow(2, 10))
\end{verbatim}
\item 這段程式碼執行後會發生什麼事？為什麼？
\begin{verbatim}
import datetime as dt
print(datetime.date.today())
\end{verbatim}
\item 如果你想用 \texttt{from random import choice as pick} 這種寫法匯入，之後要怎麼
呼叫這個函式？
\end{drills}

\section{os 與 sys：跟你的電腦、跟執行方式對話}
\label{sec:os-and-sys}

書架上翻的頭兩本書，鎖定兩個你之後會很常打交道的模組：\texttt{os}（跟作業系統、跟檔案路
徑打交道）和 \texttt{sys}（跟「這支程式是怎麼被執行的」有關的資訊）。這兩個模組的名字都
是英文縮寫，\texttt{os} 是 operating system（作業系統），\texttt{sys} 是 system
（系統），光看名字大概就能猜到它們在管什麼。

\begin{code}{python}{Commonly used os module functions}
import os

os.getcwd()
os.listdir(folder_path)
os.path.join(path_part1, path_part2)
\end{code}

\texttt{os.getcwd()} 不用傳任何引數，回傳一個字串，代表「目前所在的資料夾路徑」（cwd
是 current working directory 的縮寫，意思就是「你現在站在檔案系統的哪裡」）；
\index{os.getcwd()}
\texttt{os.listdir(folder\_path)} 傳入一個資料夾路徑字串，回傳一個 list，裡面是那個資
料夾底下所有檔案和子資料夾的名字；
\index{os.listdir()}
\texttt{os.path.join(...)} 把好幾段路徑片段組合成一
個完整路徑字串，比你自己手動用斜線接起來安全，Windows 用反斜線、macOS 和 Linux 用斜
線，\texttt{os.path.join} 會自動幫你選對的符號，你不用自己判斷。
\index{os.path.join()}

\begin{example}{python}{(list-files.py) Listing files in the current folder}
import os

current_folder = os.getcwd()
print(f"Current folder: {current_folder}")

for filename in sorted(os.listdir(current_folder)):
    full_path = os.path.join(current_folder, filename)
    print(full_path)
\end{example}

下面是作者自己電腦上、在一個叫 \texttt{python-practice} 的資料夾（就是第 \ref{ch:introduction} 章請你建立
的那一個）裡執行的結果，你自己執行時，資料夾路徑和列出來的檔案清單，一定會跟這裡不一
樣，這正是這幾個函式該做的事：老實反映你電腦上這個資料夾現在的樣子。清單裡也包含了這支
程式檔案自己（\texttt{list-files.py}），因為它自己也放在這個資料夾裡：

\iosample{}{Current folder: /Users/nathan/python-practice \\
/Users/nathan/python-practice/add.py \\
/Users/nathan/python-practice/greet.py \\
/Users/nathan/python-practice/hello.py \\
/Users/nathan/python-practice/list-files.py \\
/Users/nathan/python-practice/notes.txt \\}

到這裡你可能會想問：\texttt{os} 模組裡還有其他函式嗎？當然有，而且遠比這三個多。這是本
章最重要的示範：帶你實際去官方文件查一次，把整份清單背下來完全不是重點。打開瀏覽器，到
\texttt{https://docs.python.org/3/library/os.html}（如果你在 Google 搜尋「python os
module documentation」，通常第一個結果就是它）。

頁面上每一個函式，都會先列出「函式簽
名」，
\index{函式簽名（signature）}
長得像 \texttt{os.listdir(path='.')} 這樣的一行，你可以把它讀成：「\texttt{listdir}
這個函式可以傳入一個叫 \texttt{path} 的引數，如果你沒傳，預設值是 \texttt{'.'}」（
\texttt{.} 這個符號在檔案路徑裡代表「目前所在的資料夾」，這也是為什麼前面的範例就算不特
別傳 \texttt{path} 也能用的原因）。

函式簽名下面通常會有一兩段文字說明這個函式在做什
麼、可能丟出什麼樣的例外，再往下常常會附一小段範例程式碼。你不需要一次讀懂整頁，先掃
函式簽名和第一段說明，看看是不是你要的東西，需要更細節再往下讀範例就好。

這一整套「打開
文件、看簽名、看說明、看範例」的流程，之後每次你遇到不熟的模組或函式都可以重複套用，
接下來這一章剩下的模組，我們不會再把這個流程重講一次，只會簡短提醒你「一樣的方法，自己
去查」。

現在就練習一次：\texttt{os.path} 底下還有 \texttt{basename()} 和 \texttt{dirname()}
這兩個函式，這裡完全不打算解釋它們在做什麼，換你自己去文件頁面上找。找到之後應該會發
現，兩者都是傳入一個完整路徑字串，\texttt{basename()} 只回傳路徑最後一段（通常是檔名本
身），\texttt{dirname()} 回傳去掉最後一段之後剩下的資料夾路徑。
\index{os.path.basename()／os.path.dirname()}
動手試一次：

\begin{code}{python}{os.path.basename() and os.path.dirname()}
os.path.basename(path)
os.path.dirname(path)
\end{code}

以 \texttt{path = "/Users/nathan/python-practice/notes.txt"} 為例，
\texttt{os.path.basename(path)} 會得到 \texttt{"notes.txt"}，
\texttt{os.path.dirname(path)} 會得到 \texttt{"/Users/nathan/python-practice"}。
這兩個函式常常搭配 \texttt{os.listdir()} 一起用，例如只想挑出某個副檔名的檔案時，會
需要先用 \texttt{os.path.basename()} 拿到檔名，再用字串方法檢查結尾。

\texttt{os} 管的是檔案跟資料夾，\texttt{sys} 管的是另一件事：這支程式本身是怎麼被執
行的。這裡只教一個你以後會很常用到的東西：\texttt{sys.argv}。

\begin{code}{python}{sys.argv: arguments typed after the filename at runtime}
import sys

sys.argv[0]
sys.argv[1]
\end{code}

\texttt{sys.argv} 是一個 list，裡面裝的是你在終端機打 \texttt{python3} 之後、那整行文
字被空白切開的每一段。
\index{sys 模組／sys.argv}

\keypoint{很多新手以為 \texttt{sys.argv[0]} 是自己打的第一個參數，其實不是。
\texttt{sys.argv[0]} 永遠是這支程式自己的檔名，你自己想額外傳進去的參數，要從
\texttt{sys.argv[1]} 才開始算。}

\begin{example}{python}{(greet-arg.py) Greeting using command-line arguments}
import sys

print(f"This program's filename is: {sys.argv[0]}")

if len(sys.argv) > 1:
    name = sys.argv[1]
    print(f"Hello, {name}!")
else:
    print("Hello, stranger! You can pass a name when you run this, like this:")
    print("python3 greet-arg.py Alice")
\end{example}

假設你在終端機打 \texttt{python3 greet-arg.py Alice}，執行結果會是：

\iosample{}{This program's filename is: greet-arg.py \\
Hello, Alice! \\}

如果不加任何名字，只打 \texttt{python3 greet-arg.py}，結果會變成：

\iosample{}{This program's filename is: greet-arg.py \\
Hello, stranger! You can pass a name when you run this, like this: \\
python3 greet-arg.py Alice \\}

\texttt{greet-arg.py} 裡那一行 \texttt{if len(sys.argv) > 1} 不是隨便寫寫的，拿掉它會
出事。看看少了這道檢查會發生什麼：

\begin{example}{python}{(greet-arg-mistake.py) Accessing an argument without checking its count}
import sys

name = sys.argv[1]
print(f"Hello, {name}!")
\end{example}

假設你直接執行 \texttt{python3 greet-arg-mistake.py}，沒有多打任何名字：

\begin{verbatim}
Traceback (most recent call last):
  File "greet-arg-mistake.py", line 3, in <module>
    name = sys.argv[1]
           ~~~~~~~~^^^
IndexError: list index out of range
\end{verbatim}

沒多打參數時，\texttt{sys.argv} 只有 \texttt{sys.argv[0]}（程式檔名）這一個元素，沒有
第二個位置可以取。這就是為什麼 \texttt{greet-arg.py} 動手存取 \texttt{sys.argv[1]} 之
前，要先用 \texttt{if len(sys.argv) > 1} 確認清單裡真的有這個位置，這跟第 \ref{ch:strings} 章提過的
差一錯誤是同一種提醒：存取索引之前，先想清楚清單到底有幾個元素。

\texttt{os} 模組底下還有一個函式，值得特別停下來講：\texttt{os.path.exists(path)}，
傳入一個路徑字串，回傳 \texttt{True} 或 \texttt{False}，告訴你這個路徑存不存在。乍看之
下，這好像是「打開檔案之前先檢查檔案在不在」的完美工具：

\begin{code}{python}{Checking whether a file exists: LBYL vs. EAFP}
# 先檢查再做 (LBYL)：可以運作，但不是 Python 慣用寫法
if os.path.exists(file_path):
    process_file(file_path)
else:
    handle_missing_file()

# 先做，出錯再處理 (EAFP)：更符合 Python 的慣用寫法
try:
    process_file(file_path)
except FileNotFoundError:
    handle_missing_file()
\end{code}

\stylepoint{\texttt{os.path.exists()} 可以先檢查檔案存不存在，但還記得第 \ref{ch:files-exceptions} 章教過的
嗎？從你檢查完到你真的打開檔案之間，這段時間檔案有可能被刪掉、被搬走，狀態已經變了，
「檢查過」不等於「安全」。優先用 \texttt{try/except} 接住 \texttt{FileNotFoundError}，
還是比先檢查再動手更保險，這跟第 \ref{ch:files-exceptions} 章教的 EAFP 是同一個道理。}

跟 \texttt{os.path.exists()} 同一家族的還有兩個，用來分辨一個路徑指向的是資料夾還是
檔案：\texttt{os.path.isdir(path)} 跟 \texttt{os.path.isfile(path)}，各自回傳
\texttt{True} 或 \texttt{False}。
\index{os.path.isdir()／os.path.isfile()}
\texttt{os.listdir()} 交給你的只是一串名字，光看名字分不出哪個是資料夾、哪個是檔案，
這兩個函式就是用來問這件事的。

\begin{example}{python}{(list-what-is-what.py) 分辨資料夾與檔案}
import os

for name in sorted(os.listdir(".")):
    full_path = os.path.join(".", name)
    if os.path.isdir(full_path):
        print(name, "是資料夾")
    else:
        print(name, "是檔案")
\end{example}

\iosample{}{notes\_folder 是資料夾 \\ readme.txt 是檔案 \\}

要注意的是，這兩個函式吃的是\textbf{完整路徑}，不是 \texttt{os.listdir()} 給你的那個
名字——所以上面先用 \texttt{os.path.join()} 把資料夾位置跟名字接起來，才拿去問。
直接把 \texttt{name} 丟進 \texttt{os.path.isdir()}，除非你剛好就站在那個資料夾裡，
否則會得到 \texttt{False}，而且不會報錯，是那種安靜的錯。

\section{random 與 datetime：抽個隨機數、算個日期}
\label{sec:random-and-datetime}

書架上再翻兩本書：\texttt{random}（抽隨機數）跟 \texttt{datetime}（處理日期）。這兩個
模組的示範方式跟前一節一樣，只教一兩個你以後最常用到的功能，剩下的，你已經知道怎麼自
己去查文件了。

\begin{code}{python}{Commonly used random module functions}
import random

random.randint(low, high)
random.choice(sequence)
\end{code}

\texttt{random.randint(low, high)} 回傳一個 \texttt{low} 到 \texttt{high} 之間的整
數；
\index{random.randint()}
\texttt{random.choice(sequence)} 從一個序列（list、tuple、字串都可以）裡隨機挑一
個元素回傳。
\index{random.choice()}

\keypoint{\texttt{random.randint(low, high)} 的範圍是「包含頭尾兩端」的，\texttt{low}
跟 \texttt{high} 都有可能被抽到。這跟你已經很熟的 \texttt{range(low, high)}「不包含尾
端」剛好不一樣，很容易搞混，用之前先想一下這次要的是哪一種。}

\begin{example}{python}{(guess-number.py) A number-guessing game}
import random

secret_number = random.randint(1, 10)
guess_count = 0

while True:
    guess = int(input("Guess a number between 1 and 10: "))
    guess_count += 1
    if guess < secret_number:
        print("Too low, try again!")
    elif guess > secret_number:
        print("Too high, try again!")
    else:
        print(f"Correct! The answer was {secret_number}. You guessed {guess_count} times.")
        break
\end{example}

因為 \texttt{secret\_number} 是每次執行都重新抽的，下面這組輸入輸出只是作者某一次執行
的結果，那次亂數剛好抽到 7，你自己執行時抽到的數字幾乎一定不一樣，猜的過程也會不一
樣，這是正常的：

\iosample{3 \\ 9 \\ 7 \\}{Guess a number between 1 and 10: Too low, try again! \\
Guess a number between 1 and 10: Too high, try again! \\
Guess a number between 1 and 10: Correct! The answer was 7. You guessed 3 times. \\}

\begin{example}{python}{(pick-helper.py) Randomly picking a helper for the day}
import random

students = ["Alice", "David", "Emma", "Bob"]
today_helper = random.choice(students)
print(f"Today's helper is: {today_helper}")
\end{example}

一樣，這裡顯示的是作者某一次執行的結果，你自己執行，抽到的名字幾乎一定不一樣：

\iosample{}{Today's helper is: Emma \\}

到這裡你已經看過一次「打開官方文件、看函式簽名、看說明」的完整示範（第~
\ref{sec:os-and-sys}~節示範了怎麼查 \texttt{os} 模組）。這一次換你自己動手：
\texttt{random} 模組裡除了 \texttt{randint} 和 \texttt{choice}，還有一個很好用的函式
叫 \texttt{sample}，這一章完全沒教過它，要自己去
\texttt{https://docs.python.org/3/library/random.html} 找。

頁面上往下捲，會找到一行
函式簽名長得像 \texttt{random.sample(population, k)}，說明寫著「從 population 這個序
列裡，不重複地隨機抽出 k 個元素，回傳一個 list」。
\index{random.sample()}
光看這一行，大概就能猜到它跟
\texttt{random.choice()} 的差別：\texttt{choice()} 每次只抽一個，同一批人裡可能被抽到
好幾次；\texttt{sample()} 一次抽好 k 個，保證彼此不重複，很適合抽獎、分組這類「不能抽
到同一個人兩次」的情境。

\begin{example}{python}{(lottery-pick.py) Using random.sample() from the docs for a lottery draw}
import random

numbers = list(range(1, 43))
winning_numbers = random.sample(numbers, 6)
print(f"This week's winning numbers: {sorted(winning_numbers)}")
\end{example}

跟前面的 random 範例一樣，下面顯示的是作者某一次執行的結果，你自己執行，抽到的號碼幾
乎一定不一樣：

\iosample{}{This week's winning numbers: [11, 25, 26, 30, 32, 39] \\}

\stylepoint{\texttt{random.sample()} 本身回傳的 list 沒有排序，抽出來的順序就是抽出的
順序。範例裡特地包了一層 \texttt{sorted(...)}，只是為了讓開獎號碼看起來由小到大比較好
讀，不影響 \texttt{sample()} 抽樣本身的公平性。}

接下來看 \texttt{datetime}，處理日期最常用的模組。

\begin{code}{python}{Commonly used datetime module functions}
import datetime

datetime.date.today()
datetime.date(year, month, day)
some_date + datetime.timedelta(days=n)
\end{code}

\texttt{datetime.date.today()} 回傳一個代表「今天」的日期物件，印出來會是
\texttt{YYYY-MM-DD} 這種格式；
\index{datetime.date.today()}
如果你要的日期換成某一天你自己指定的日期，不是今天，可以用
\texttt{datetime.date(year, month, day)}，例如 \texttt{datetime.date(2027, 1, 1)} 就是
「2027 年 1 月 1 日」這一天的日期物件，跟 \texttt{today()} 回傳的東西是同一種型別，可以
放在一起做運算；
\index{datetime.date(year, month, day)}
兩個日期物件相減，會得到一個 \texttt{timedelta} 物件，
裡面的 \texttt{.days} 屬性就是「相差幾天」的整數；一個日期物件加上
\texttt{datetime.timedelta(days=n)}，會得到 \texttt{n} 天後的另一個日期物件，減法自
然也可以，會得到 \texttt{n} 天前的日期。
\index{datetime.timedelta()}

\begin{example}{python}{(event-countdown.py) Counting the days until a date}
import datetime

today = datetime.date.today()
print(f"Today's date is: {today}")

new_year = datetime.date(2027, 1, 1)
days_left = (new_year - today).days
print(f"There are {days_left} days left until January 1, 2027")

reminder_date = today + datetime.timedelta(days=10)
print(f"10 days from today is: {reminder_date}")
\end{example}

「今天」跟「還有幾天」這兩行的數字，會隨著你實際執行的日期改變，下面顯示的是作者寫這
段文字那天（2026 年 9 月 1 日）執行的結果：

\iosample{}{Today's date is: 2026-09-01 \\
There are 122 days left until January 1, 2027 \\
10 days from today is: 2026-09-11 \\}

日期物件加減的時候，如果忘了用 \texttt{timedelta} 包起來，直接加一個整數，會出什麼事？

\begin{example}{python}{(date-mistake.py) Adding an integer directly instead of using timedelta}
import datetime

today = datetime.date.today()
next_week = today + 7
print(next_week)
\end{example}

執行結果：

\begin{verbatim}
Traceback (most recent call last):
  File "date-mistake.py", line 4, in <module>
    next_week = today + 7
                ~~~~~~^~~
TypeError: unsupported operand type(s) for +: 'datetime.date' and 'int'
\end{verbatim}

這裡的 \texttt{TypeError} 是在提醒你：日期物件不知道「加 7」該當成加 7 天、7 小時，還
是別的什麼單位，它只認得 \texttt{timedelta} 這種專門表示「一段時間長度」的物件。正確寫
法要把 7 包成 \texttt{datetime.timedelta(days=7)}，像本節稍早的
\texttt{event-countdown.py} 那樣。

\begin{drills}[label=drills:stdlib-modules]
\item \texttt{random.randint(1, 5)} 可能抽到的最小值、最大值各是多少？
\item 你的程式檔名是 \texttt{weather.py}，在終端機打
\texttt{python3 weather.py taipei}，\texttt{sys.argv} 這個 list 裡有幾個元素？
\texttt{sys.argv[0]}、\texttt{sys.argv[1]} 分別是什麼？
\item 小華覺得 \texttt{random.randint(1, 10)} 應該跟 \texttt{range(1, 10)} 一樣不包含
10，所以他很確定下面這段程式碼絕對不會印出「Got a 10!」：
\begin{verbatim}
number = random.randint(1, 10)
if number == 10:
    print("Got a 10!")
\end{verbatim}
小華的想法哪裡錯了？
\item \texttt{random.sample([10, 20, 30, 40], 2)} 會回傳一個裝有幾個元素的 list？裡面
的元素有沒有可能重複？
\item \texttt{path = "/Users/nathan/python-practice/notes.txt"}，
\texttt{os.path.basename(path)} 跟 \texttt{os.path.dirname(path)} 分別會得到什麼？
\end{drills}

\section{collections.Counter：專門用來計數的字典}
\label{sec:collections-and-itertools}

第 \ref{ch:dict-set} 章教過的 dict，很適合用來記錄「一個線索對應一個答案」的資料。有一種特別常見的任
務，剛好是 dict 拿手的另一種用法：給你一份清單，數清單裡每一種東西各出現幾次。用你現
有的工具就寫得出來：

\begin{example}{python}{(vote-count-manual.py) Manually counting with a dict}
votes = ["Alice", "David", "Alice", "Emma", "Alice", "David", "Alice", "Bob"]
vote_counts = {}

for name in votes:
    vote_counts[name] = vote_counts.get(name, 0) + 1

print(vote_counts)
\end{example}

\iosample{}{\{'Alice': 4, 'David': 2, 'Emma': 1, 'Bob': 1\} \\}

這裡用了第 \ref{ch:dict-set} 章教過的 \texttt{get(key, default)}：如果 \texttt{name} 這個 key 已經在
\texttt{vote\_counts} 裡，就把它原本的票數加 1；如果還沒出現過，\texttt{get} 會給一個
預設值 0，加 1 之後剛好就是這個人的第一票。這樣寫完全沒有問題，但「用 dict 計數」這個
模式實在太常見，Python 標準函式庫的 \texttt{collections} 模組，已經把它包成一個現成的
容器型別，叫 \texttt{Counter}：
\index{collections.Counter}

\begin{code}{python}{Common usage of collections.Counter}
from collections import Counter

counts = Counter(iterable)
counts.most_common(n)
\end{code}

把任何一個可疊代物件（list、字串都可以）直接餵給 \texttt{Counter(...)}，它會自動數出
裡面每個元素各出現幾次，回傳一個 \texttt{Counter} 物件，可以把它想成一種「加強版的
dict」：查詢、走訪的用法跟 dict 完全一樣，只是多了幾個計數專用的方法。

\texttt{counts.most\_common(n)} 依照出現次數由多到少排序，回傳前 \texttt{n} 名，每一名
是一個 \texttt{(元素, 次數)} 的 tuple，整體裝在一個 list 裡（呼應第 \ref{ch:tuples} 章教過的
unpacking，可以直接用 \texttt{for 元素, 次數 in counts.most\_common():} 走訪）；不傳
\texttt{n} 的話，預設回傳全部，一樣是由多到少排好序。
\index{most\_common()}

\begin{example}{python}{(vote-count.py) Rewriting the vote count with Counter}
from collections import Counter

votes = ["Alice", "David", "Alice", "Emma", "Alice", "David", "Alice", "Bob"]
vote_counts = Counter(votes)

print(vote_counts)

winner, winner_votes = vote_counts.most_common(1)[0]
print(f"{winner} is currently leading with {winner_votes} votes")

for name, count in vote_counts.most_common():
    print(f"{name}: {count} votes")
\end{example}

\iosample{}{Counter(\{'Alice': 4, 'David': 2, 'Emma': 1, 'Bob': 1\}) \\
Alice is currently leading with 4 votes \\
Alice: 4 votes \\
David: 2 votes \\
Emma: 1 votes \\
Bob: 1 votes \\}

跟 \texttt{vote-count-manual.py} 對照著看：兩支程式數出來的結果一模一樣，但
\texttt{Counter} 版本不用自己寫 \texttt{get(name, 0) + 1} 這個累加邏輯，還多送了
\texttt{most\_common()} 這個直接排好名次的方法。

\keypoint{查詢一個不存在的 key 時，\texttt{Counter} 跟一般 dict 的行為不一樣。一般
dict 用中括號查一個不存在的 key，會丟出第 \ref{ch:dict-set} 章教過的 \texttt{KeyError}；
\texttt{Counter} 遇到同樣的情況，會直接回傳 0，不會報錯。這其實也是一種對讀者友善的
EAFP 精神：計數這件事，「查無此人」本來就等於「目前是 0 票」，\texttt{Counter} 幫你把
這個常見情境內建好了，不用自己再寫 \texttt{try/except} 或 \texttt{get()} 去接。}

\begin{example}{python}{(counter-missing-key.py) Looking up a name that never appeared}
from collections import Counter

votes = ["Alice", "David", "Alice", "Emma"]
vote_counts = Counter(votes)

print(vote_counts["Bob"])
\end{example}

\iosample{}{0 \\}

\texttt{Counter} 是 \texttt{collections} 模組裡的名字，不是 Python 內建、隨時都能用的
東西，忘記 \texttt{import} 一樣會出現本章一開始就看過的 \texttt{NameError}：

\begin{example}{python}{(counter-import-mistake.py) Forgetting to import Counter}
votes = ["Alice", "David", "Alice"]
vote_counts = Counter(votes)
print(vote_counts)
\end{example}

\begin{verbatim}
Traceback (most recent call last):
  File "counter-import-mistake.py", line 2, in <module>
    vote_counts = Counter(votes)
                  ^^^^^^^
NameError: name 'Counter' is not defined
\end{verbatim}

這支程式缺了開頭那一行 \texttt{from collections import Counter}，跟第~
\ref{sec:import-syntax}~節示範過的道理完全一樣：不管一個名字來自標準函式庫裡的哪個模
組，用之前一定要先 \texttt{import}，Python 不會自動幫你找。

標準函式庫裡跟 \texttt{collections} 放在一起的，還有一個叫 \texttt{itertools} 的模
組，專門處理「疊代」相關的各種小技巧，例如把好幾個序列的元素兩兩配對出所有組合。這裡只
點到為止，讓你知道遇到「窮舉所有組合」這類需求時，可以想到有這個模組存在，值得去官方文
件看看，不會要求你這裡就精熟它：
\index{itertools（product()／combinations()）}

\begin{example}{python}{(itertools-peek.py) Using itertools.product to list all combinations}
import itertools

sizes = ["S", "M", "L"]
colors = ["Red", "Blue"]

for size, color in itertools.product(sizes, colors):
    print(f"Size {size}, Color {color}")
\end{example}

\iosample{}{Size S, Color Red \\
Size S, Color Blue \\
Size M, Color Red \\
Size M, Color Blue \\
Size L, Color Red \\
Size L, Color Blue \\}

如果自己寫，這段程式碼要用第 \ref{ch:loops} 章教過的巢狀迴圈（一個 \texttt{for} 迴圈裡面再包一個
\texttt{for} 迴圈）才寫得出同樣的效果，\texttt{itertools.product(...)} 把這個常見模式
包成一行。\texttt{itertools} 模組裡類似的函式還有不少，例如
\texttt{itertools.combinations(sequence, r)}，列出「從一批東西裡挑 \texttt{r} 個、不
考慮挑選順序」的所有搭配（適合「兩兩配對但不重複算兩次」這種情境，跟 \texttt{product}
「每個位置各自獨立配對」不太一樣），用法都大同小異，需要的時候查文件就好。

標準函式庫裡
類似 \texttt{itertools} 這種「會用就好、不強求精熟」的模組還有很
多，之後遇到「這個任務感覺 Python 應該已經有現成工具」的情況，養成先去官方文件搜尋一
下的習慣，往往比自己從頭寫更快。

\begin{drills}[label=drills:collections-and-itertools]
\item 下面這段程式碼會印出什麼？
\begin{verbatim}
from collections import Counter

grades = ["A", "B", "A", "A", "C", "B"]
grade_counts = Counter(grades)
print(grade_counts["A"])
print(grade_counts["F"])
\end{verbatim}
\item \texttt{Counter(["x", "y", "x"]).most\_common(1)} 這一行會回傳什麼？
\item \texttt{itertools.product(["a", "b"], [1, 2])} 這一行會依序配對出幾組結果？
\end{drills}

\section{寫自己的模組}
\label{sec:write-your-own-module}

前面幾節，你都是在拿別人寫好的書來用：\texttt{math}、\texttt{os}、\texttt{sys}、
\texttt{random}、\texttt{datetime}、\texttt{collections}，這些都是別人寫好、測試過，
放在書架上的模組。這一節要把角色反過來：你自己寫的 \texttt{.py} 檔案，本質上跟這些模組
是同一種東西，一樣可以被 \texttt{import}，一樣可以被別的檔案拿去用。
\index{自訂模組}

\begin{code}{python}{Saving a function as a standalone module}
# my_module.py
def function_name(argument):
    ...

# main.py
import my_module

my_module.function_name(value)
\end{code}

具體做法是：把你想重複使用的函式，存進一個獨立的 \texttt{.py} 檔案，檔名就是之後
\texttt{import} 要打的名字（不含 \texttt{.py} 這個副檔名）。下面先準備一個工具檔案，裡
面放兩個小函式：

\begin{example}{python}{(my\_tools.py) Saving two small utility functions as a standalone file}
def count_vowels(text):
    """Count the number of vowel letters in a string (a, e, i, o, u, case-insensitive)."""
    vowels = "aeiouAEIOU"
    count = 0
    for char in text:
        if char in vowels:
            count += 1
    return count


def celsius_to_fahrenheit(celsius):
    """Convert Celsius temperature to Fahrenheit."""
    return celsius * 9 / 5 + 32
\end{example}

你可能注意到，這兩個函式主體的第一行都是一個三引號字串，例如
\texttt{"""Count the number of vowel letters..."""}，這叫做 docstring（說明文件字串），用來說
明這個函式在做什麼，不是必要的，但是個好習慣，之後別人（或你自己）用 \texttt{help()} 或
編輯器查詢這個函式時，就會看到這段說明。
\index{docstring（說明文件字串）}
實際用 \texttt{help()} 查一次看看：

\begin{example}{python}{(check-docstring.py) Using help() to view a function's docstring}
import my_tools

help(my_tools.count_vowels)
\end{example}

執行後，終端機會顯示：

\begin{verbatim}
Help on function count_vowels in module my_tools:

count_vowels(text)
    Count the number of vowel letters in a string (a, e, i, o, u, case-insensitive).
\end{verbatim}

\texttt{help()} 不只可以查你自己寫的函式，前面幾節提到的 \texttt{os.listdir}、
\texttt{random.randint}，只要模組已經 \texttt{import} 過，都可以直接
\texttt{help(os.listdir)} 這樣查。這跟本章稍早示範的「打開官方文件查函式簽名」是同一
件事的另一種做法，差別在 \texttt{help()} 直接印在終端機裡，不用開瀏覽器；缺點是網頁文
件通常還會有額外的範例和更完整的說明，兩種方式互相搭配使用即可。

這支檔案本身執行不會印出任何東西，它只是定義了兩個函式，等別的檔案來呼叫，所以這裡沒
有對應的輸出範例。把它存成 \texttt{my\_tools.py}，跟下面這支檔案放在同一個資料夾裡：

\begin{example}{python}{(use-my-tools.py) Calling your own function with import my\_tools}
import my_tools

sentence = "Hello, World!"
vowel_count = my_tools.count_vowels(sentence)
print(f"'{sentence}' has {vowel_count} vowels")

temperature = my_tools.celsius_to_fahrenheit(30)
print(f"30 Celsius is {temperature} Fahrenheit")
\end{example}

\iosample{}{'Hello, World!' has 3 vowels \\
30 Celsius is 86.0 Fahrenheit \\}

你剛剛做的事，跟第~\ref{sec:import-syntax}~節的 \texttt{import math} 是一模一樣的動
作，只是這次書架上的書換成了你自己寫的 \texttt{my\_tools.py}。既然是一模一樣的動作，
第~\ref{sec:import-syntax}~節教過的幾種寫法，在這裡當然也完全適用：

\begin{example}{python}{(use-my-tools-2.py) Calling your own function with a precise import}
from my_tools import celsius_to_fahrenheit

temperature = celsius_to_fahrenheit(0)
print(f"0 Celsius is {temperature} Fahrenheit")
\end{example}

\iosample{}{0 Celsius is 32.0 Fahrenheit \\}

\stylepoint{你自己的檔案被 \texttt{import} 之後，也是被貼上一張名字牌（模組名
\texttt{my\_tools}），跟 \texttt{os}、\texttt{random} 這些別人寫的模組沒有兩樣，
Python 並不區分「標準函式庫的模組」跟「你自己寫的模組」，兩者是同一套機制。}

自己寫的模組要能被 \texttt{import}，有一個容易忽略的前提：檔案要跟呼叫它的那支程式放
在同一個資料夾裡（至少目前這個階段，你只需要記得這一點，更複雜的資料夾結構是之後才會
遇到的事）。如果不是，會發生什麼事？

\begin{example}{python}{(missing-module.py) Importing when the file isn't in the same folder}
import my_tools

print(my_tools.count_vowels("test"))
\end{example}

假設 \texttt{missing-module.py} 跟 \texttt{my\_tools.py} 沒有放在同一個資料夾，執行時
會看到：

\begin{verbatim}
Traceback (most recent call last):
  File "missing-module.py", line 1, in <module>
    import my_tools
ModuleNotFoundError: No module named 'my_tools'
\end{verbatim}

還記得第 \ref{ch:data-tools} 章的 \texttt{ModuleNotFoundError} 嗎？那裡是套件根本沒有安裝；這裡的情況不
一樣，\texttt{my\_tools.py} 確實存在，只是不在 Python 找模組的地方。Python 找你自己寫
的模組，預設只會找「目前執行的這支程式所在的資料夾」，不會自動去別的資料夾翻找，這也是
為什麼前面一直提醒你，兩個檔案要放在同一個資料夾裡。

\begin{drills}[label=drills:write-your-own-module]
\item 你把一個函式 \texttt{def add(a, b): return a + b} 存成一個檔案，想在另一個檔案裡
用 \texttt{import my\_math} 的方式使用它，這個檔案該存成什麼檔名？
\item 如果你想省略模組名字首碼，直接用 \texttt{add(3, 5)} 呼叫這個函式，
\texttt{import} 那一行該怎麼寫？
\item 小美把 \texttt{my\_tools.py} 存在桌面的一個資料夾，把 \texttt{use-my-tools.py}
存在另一個資料夾，執行後出現 \texttt{ModuleNotFoundError}，該怎麼修正？
\end{drills}

\section*{練習解答}
\addcontentsline{toc}{section}{練習解答}

\paragraph{第~\ref{sec:import-syntax}~節}
\begin{enumerate}
\item 會出現 \texttt{NameError: name 'sqrt' is not defined}，這裡只用
\texttt{import math} 整模組匯入，\texttt{sqrt} 這個名字沒有被單獨貼上名字牌，只有
\texttt{math} 有；要嘛改成 \texttt{math.sqrt(25)}，要嘛開頭改成
\texttt{from math import sqrt}。
\item 會出現 \texttt{NameError: name 'math' is not defined}，這裡只用
\texttt{from math import pi} 精確匯入了 \texttt{pi}，\texttt{math} 這個名字本身沒有被
貼過名字牌，應該直接寫 \texttt{print(pi)}。
\item 會印出 \texttt{1024.0}，\texttt{from math import pow} 已經把 \texttt{pow} 精確
匯入，可以直接呼叫，\texttt{2} 的 \texttt{10} 次方是 1024，\texttt{math.pow} 回傳的是
浮點數。
\item 會出現 \texttt{NameError: name 'datetime' is not defined}，\texttt{import
datetime as dt} 只把名字牌 \texttt{dt} 貼到 \texttt{datetime} 模組上，\texttt{datetime}
這個名字本身沒有被貼過名字牌，應該寫成 \texttt{dt.date.today()}。
\item \texttt{pick(sequence)}，取別名之後，程式碼裡只認得別名 \texttt{pick}，原本的函
式名字 \texttt{choice} 已經不能再拿來呼叫。
\end{enumerate}

\paragraph{第~\ref{sec:random-and-datetime}~節}
\begin{enumerate}
\item 最小值是 1，最大值是 5，\texttt{random.randint(low, high)} 的範圍包含頭尾兩端。
\item \texttt{sys.argv} 有 2 個元素：\texttt{sys.argv[0]} 是 \texttt{"weather.py"}，
\texttt{sys.argv[1]} 是 \texttt{"taipei"}。
\item 小華把 \texttt{random.randint()} 跟 \texttt{range()} 的邊界行為搞混了，
\texttt{random.randint(1, 10)} 是「包含頭尾兩端」的閉區間，10 是有可能被抽到的，
所以這段程式碼有可能印出「Got a 10!」，跟不包含尾端的 \texttt{range(1, 10)} 不一樣。
\item 會回傳一個裝有 2 個元素的 list，元素之間不會重複，因為 \texttt{random.sample()}
是不重複抽樣，跟每次都可能抽到同一個的 \texttt{random.choice()} 不一樣。
\item \texttt{os.path.basename(path)} 會得到 \texttt{"notes.txt"}（路徑最後一段）；
\texttt{os.path.dirname(path)} 會得到 \texttt{"/Users/nathan/python-practice"}（去掉
最後一段之後剩下的資料夾路徑）。
\end{enumerate}

\paragraph{第~\ref{sec:collections-and-itertools}~節}
\begin{enumerate}
\item 依序印出 \texttt{3} 和 \texttt{0}：\texttt{"A"} 在清單裡出現 3 次；\texttt{"F"}
從沒出現過，\texttt{Counter} 對不存在的 key 直接回傳 0，不會丟出 \texttt{KeyError}。
\item 會回傳 \texttt{[('x', 2)]}，一個裝著一個 tuple 的 list，\texttt{"x"} 出現 2 次，
是清單裡出現次數最多的元素，\texttt{most\_common(1)} 只取排名第一名。
\item 會依序配對出 4 組：\texttt{("a", 1)}、\texttt{("a", 2)}、\texttt{("b", 1)}、
\texttt{("b", 2)}，第一個序列有 2 個元素、第二個序列有 2 個元素，總共 $2 \times 2 = 4$
種組合。
\end{enumerate}

\paragraph{第~\ref{sec:write-your-own-module}~節}
\begin{enumerate}
\item \texttt{my\_math.py}，檔名（不含 \texttt{.py}）要跟 \texttt{import} 後面打的
名字完全一致。
\item \texttt{from my\_math import add}。
\item 把 \texttt{my\_tools.py} 搬到跟 \texttt{use-my-tools.py} 同一個資料夾裡。
Python 找自己寫的模組，預設只會找「目前執行的這支程式所在的資料夾」，兩個檔案不在同一
個資料夾就會找不到。
\end{enumerate}

\chapter{函式進階：把函式當成資料來用}
\label{ch:advanced-functions}

請人幫忙做事的時候，交出去的不一定只有「東西」，也可能是「怎麼做」這條規則本身。你把一疊
表單推給別人說「幫我排好」，對方一定會反問一句：照什麼排？照日期、照名字，還是照金額？
「整理」這個動作是固定的，會變的是那條比大小的規則。你把規則講出口的那一刻，其實就把
「怎麼比」這件事整個交出去了——對方負責跑流程，你負責決定標準。

第 \ref{ch:functions} 章的最後幾節，你已經摸到這件事的邊了：第 \ref{sec:function-first-class} 節說過，
用 \texttt{def} 定義出來的函式本身也是一個物件，函式名稱只是貼在它身上的一張名字牌，所以
函式可以被指定給別的變數（\texttt{operation = square}）、可以被當作引數傳給另一個函式
（\texttt{apply\_twice(add\_ten, 5)}）。不過那一節裡收下函式的 \texttt{apply\_twice}，是
我們自己寫的。真實世界的 Python 程式碼裡，情況通常反過來——\textbf{是別人早就寫好的函式，
張著手等你把你的規則交給它}。這一章就是從這個轉向開始。

這一章分四段推進：

\begin{itemize}
\item 第一段是「把規則交出去」：先學一個更短的寫法來寫那些用完就丟的小函式，再看它最常
出現的三個位置——排序時的比較依據、找最大最小值的依據、以及一次處理一整串資料的處理方式
（第 \ref{sec:lambda}～\ref{sec:map-filter} 節）。
\item 第二段整個翻面，談「把規則帶回來」：函式不只可以被傳進去，也可以被別的函式生出來、
再交還給你，而且它會記得自己出生時看到的東西（第 \ref{sec:closure} 節）。
\item 第三段是這一章的最高點，也是前兩段的合體：在一個已經能用的函式外面「包一層」，讓它
多做點事，但完全不動它原本的程式碼；先手動包一次，再揭曉 Python 替這件事準備的簡寫
（第 \ref{sec:wrap-by-hand}～\ref{sec:decorator} 節）。
\item 第四段回頭談一件完全不同的事：怎麼在 \texttt{def} 那一行就寫清楚「這個函式吃什麼、
吐什麼」，以及這種標註一個非常反直覺的性質（第 \ref{sec:type-hints} 節）。
\end{itemize}

\section{不用取名字的小函式：\texttt{lambda}}
\label{sec:lambda}

回頭看第 \ref{sec:function-first-class} 節那個 \texttt{apply\_twice(add\_ten, 5)} 的例子。
為了把「加十」這條規則交給 \texttt{apply\_twice}，我們先花了三行寫一個 \texttt{add\_ten}
函式，替它想了一個名字，然後只用了一次就再也沒提過它。函式本體其實只有 \texttt{n + 10}
這麼一點點東西，包裝卻比內容還大。

Python 對這種「只有一個算式、用完就丟」的小函式，準備了一個更短的寫法：\texttt{lambda}。
它做的事跟 \texttt{def} 完全一樣——造出一個函式物件——差別只在它不需要名字，可以直接寫在
需要它的地方。
\index{lambda}

\begin{code}{python}{用 lambda 寫一個不具名的小函式}
lambda parameter: expression
\end{code}

\texttt{lambda} 這個字後面接參數名稱（要好幾個就用逗號隔開），一個冒號，然後是一個
\textbf{算得出值的式子}。這個式子算出來的結果，就是這個函式的回傳值——不用寫
\texttt{return}，也不能寫。

先看看它跟 \texttt{def} 版本擺在一起是什麼樣子。第 \ref{ch:functions} 章的寫法是：

\begin{code}{python}{第 \ref{ch:functions} 章原本的 def 寫法}
def add_ten(n):
    return n + 10


print(apply_twice(add_ten, 5))
\end{code}

換成 \texttt{lambda} 之後，三行的定義直接縮成呼叫裡的一小段：

\begin{example}{python}{(apply-twice-lambda.py) 用 lambda 取代只用一次的小函式}
def apply_twice(func, value):
    return func(func(value))


print(apply_twice(lambda n: n + 10, 5))
\end{example}

\iosample{}{25 \\}

輸出跟第 \ref{sec:function-first-class} 節那個範例一模一樣，也還是 \texttt{25}：換了寫法，
沒換行為。\texttt{lambda n: n + 10} 整串本身就是一個函式物件，它被直接當成引數塞進
\texttt{apply\_twice} 的 \texttt{func} 參數裡，跟原本傳 \texttt{add\_ten} 沒有任何差別。
唯一的差別是，這個函式物件從頭到尾沒有被貼上任何一張名字牌。

\keypoint{\texttt{lambda ...} 整串寫出來的東西，就是一個函式物件，跟 \texttt{def} 造出來的
是同一種東西——差別只在它沒有被貼上名字牌。這也解釋了為什麼你可以把它直接寫在括號裡：既然
它本身就是那個「函式物件」，當然不用先存進變數再傳過去。要判斷一段程式碼給出去的是「函式
本身」還是「呼叫函式的結果」，判準跟第 \ref{sec:function-first-class} 節一樣，永遠是看有沒有加括號。}

既然 \texttt{lambda} 造出來的也是函式物件，那它當然也可以被存進變數：

\begin{example}{python}{(say-hi-lambda.py) 把 lambda 存進變數再呼叫}
say_hi = lambda name: f"Hi, {name}!"

print(say_hi("Amy"))
print(say_hi("Ben"))
\end{example}

\iosample{}{Hi, Amy! \\ Hi, Ben! \\}

這樣寫可以動，而且看得懂——\texttt{say\_hi} 這張名字牌被貼到 \texttt{lambda} 造出來的函式
物件上，之後 \texttt{say\_hi("Amy")} 就是在呼叫它。但這種寫法在 Python 圈子裡評價不高，
原因很單純：繞了一圈，做的其實就是 \texttt{def} 的事。

\stylepoint{如果你打算替一個 \texttt{lambda} 取名字（也就是寫成
\texttt{名字 = lambda ...}），那就直接用 \texttt{def} 寫。兩種寫法效果一樣，但 \texttt{def}
版本更好讀，出錯時的錯誤訊息也會告訴你是哪個函式出問題。\texttt{lambda} 真正的舞台，是
「當引數直接塞進別的函式裡」的場合——也就是下一節要看的東西。}

順帶一提：既然函式是物件，它就跟數字、字串一樣，也可以被放進 list 或 dict 裡當成資料收著。
這件事在寫比較大的程式時偶爾會派上用場，這裡先知道「可以」就好。

\subsection*{常見錯誤：在 lambda 裡面寫 return}

\texttt{lambda} 的冒號後面只能放一個「算得出值的式子」。很多人第一次寫的時候會照著
\texttt{def} 的習慣補上 \texttt{return}：

\begin{example}{python}{(lambda-return-broken.py) 在 lambda 主體裡寫 return}
add_one = lambda n: return n + 1

print(add_one(5))
\end{example}

\begin{verbatim}
  File "lambda-return-broken.py", line 1
    add_one = lambda n: return n + 1
                        ^^^^^^
SyntaxError: invalid syntax
\end{verbatim}

這次連程式都沒開始跑就出錯了——\texttt{SyntaxError} 是「你寫的東西不符合 Python 的文法」，
Python 連讀都讀不下去。原因是 \texttt{return} 是一個\textbf{敘述}（statement），是一個
「動作」；而 \texttt{lambda} 的冒號後面要的是一個\textbf{式子}（expression），是一個「算得
出值的東西」。同樣的理由，\texttt{if} 的整段結構、\texttt{for} 迴圈、\texttt{print} 之外的
賦值，通通不能塞進 \texttt{lambda}。正確的寫法是把 \texttt{return} 拿掉，只留
\texttt{lambda n: n + 1}——那個式子算出來的值自動就是回傳值。

這條限制看起來很嚴格，其實正好說明了 \texttt{lambda} 的定位：它只適合真的很小的函式。
只要邏輯稍微複雜一點，你自然就得寫回 \texttt{def}。

\begin{drills}[label=drills:lambda]
\item 有 \texttt{f = lambda a, b: a * b}，那麼 \texttt{print(f(3, 4))} 會印出什麼？
\item 把 \texttt{lambda n: n * n} 改寫成一個叫 \texttt{square} 的 \texttt{def} 函式。
\end{drills}

\section{把「照什麼標準」交出去：\texttt{sorted(key=)}、\texttt{max}、\texttt{min}}
\label{sec:key-function}

開場那個「照什麼排」的問題，在這一節兌現。這件事你其實已經做過一次了：第
\ref{sec:sort-key} 節教過 \texttt{sort()}／\texttt{sorted()} 的 \texttt{key} 參數，
你在那裡寫過 \texttt{key=len}，把「量長度」這條規則交給排序函式，讓它照字數排而不是照
字母順序。當時沒有特別點破的是：\textbf{那一行就是把一個函式當成引數傳給另一個函式}——只是那時候
你還沒有「函式是一種值」這個說法可以用。

現在你有了，而且可以往前走一步。\texttt{key=len} 之所以方便，是因為「量長度」剛好有一個
現成的內建函式。但真實的資料很少這麼配合：如果你手上是一串商品，每一項都是一個 dict，
要照價格排，Python 沒有一個現成的「取出價格」函式可以給你。這種時候，你要自己生一個——
而且它短到不值得取名字，正是上一節 \texttt{lambda} 的舞台。

先把規則講完整：\texttt{key} 收的是一個函式，排序時它會拿這個函式對每一項各跑一次，用算出來
的結果當比大小的依據。\texttt{sorted} 負責排，你負責給「照什麼排」。\texttt{max} 跟
\texttt{min} 也是同一套。

\begin{code}{python}{用 key 指定比較的依據}
sorted(iterable, key=function)
max(iterable, key=function)
min(iterable, key=function)
\end{code}

有一件事第 \ref{sec:sort-key} 節已經強調過，這裡再提醒一次，因為接下來每一個範例都靠它：
\texttt{key=} 後面交出去的是函式\textbf{本身}，名字後面\textbf{不加括號}。加了括號就變成
「現在就呼叫它、把結果交出去」，那是完全不同的一件事——第 \ref{sec:sort-key} 節那個
\texttt{key=len()} 的 \texttt{TypeError}，跟第 \ref{sec:function-first-class} 節那個
\texttt{operation = square()}，錯的都是同一件事。

接下來是這一節真正的重點，也是你以後最常真的會寫的形狀。第 \ref{ch:dict-set} 章那種
「一串 dict」是實務上最常見的資料長相：

\begin{example}{python}{(sort-products.py) 把一串商品資料照價格排序}
products = [
    {"name": "notebook", "price": 120},
    {"name": "pen", "price": 25},
    {"name": "folder", "price": 60},
]

by_price = sorted(products, key=lambda item: item["price"])

for product in by_price:
    print(product["name"], product["price"])

print(products[0]["name"])
\end{example}

\iosample{}{pen 25 \\ folder 60 \\ notebook 120 \\ notebook \\}

\texttt{lambda item: item["price"]} 這一段翻成白話就是：「給我一項商品，我回你它的價格。」
\texttt{sorted} 拿到這條規則之後，會對清單裡的每一個 dict 各呼叫它一次，拿回三個數字
（120、25、60），然後照這三個數字的大小去排原本那三個 dict。最後一行印出的還是
\texttt{notebook}，再次確認 \texttt{products} 自己完全沒被動過。

同一份資料，如果要的不是整串排序，而只是「最貴的那一項」，\texttt{max} 收的
\texttt{key} 一模一樣：

\begin{example}{python}{(most-expensive.py) 找出最貴與最便宜的商品}
products = [
    {"name": "notebook", "price": 120},
    {"name": "pen", "price": 25},
    {"name": "folder", "price": 60},
]

top = max(products, key=lambda item: item["price"])
cheapest = min(products, key=lambda item: item["price"])

print(top["name"], top["price"])
print(cheapest["name"], cheapest["price"])
\end{example}

\iosample{}{notebook 120 \\ pen 25 \\}

注意 \texttt{max} 回傳的是\textbf{那一整項商品}（整個 dict），不是那個價格數字——
\texttt{key} 只影響「怎麼比」，不影響「回傳什麼」。\texttt{min} 是同一回事，只是挑的是
算出來最小的那一項。

\keypoint{\texttt{key=} 要的是一個「吃一項、吐出比較依據」的函式。\texttt{sorted}／
\texttt{max}／\texttt{min} 會自己對每一項呼叫它一次，你不用寫迴圈；而且它們只拿算出來的
結果當比較依據，原本的資料一個字都不會被改動，回傳的也還是原本那些東西。}

\begin{drills}[label=drills:key-function]
\item \texttt{print(sorted(["kiwi", "fig", "banana"], key=len))} 會印出什麼？
\item 有 \texttt{people = [\{"name": "Amy", "age": 30\}, \{"name": "Ben", "age": 24\}]}，
那麼 \texttt{print(sorted(people, key=lambda p: p["age"])[0]["name"])} 會印出什麼？
\item 有人寫了 \texttt{sorted(products, key=lambda item: item["price"]())}，執行時出錯了。
問題出在哪裡？
\end{drills}

\section{一次處理一整串：\texttt{map} 與 \texttt{filter}}
\label{sec:map-filter}

\texttt{sorted(key=)} 是「把比較的依據交出去」。同樣是交一條規則出去，還有兩個內建函式收的
是別種規則：\texttt{map} 收的是「每一項要怎麼處理」，\texttt{filter} 收的是「哪些項目要
留下來」。
\index{map}
\index{filter}

\begin{code}{python}{把處理方式或篩選條件交給 map／filter}
map(function, iterable)
filter(function, iterable)
\end{code}

先看 \texttt{map}，而且情境是你已經吃過苦頭的：第 \ref{ch:files-exceptions} 章跟
第 \ref{ch:csv} 章都提過，從檔案裡讀進來的東西一律是字串，就算長得像數字也一樣，要算數之前
得先一個一個轉成整數。這件事交給 \texttt{map} 一行就完成，而且要交出去的規則同樣是一個
你早就認識的函式——\texttt{int}。

\begin{example}{python}{(convert-prices.py) 把一串數字字串一次轉成整數}
prices = ["12", "5", "38"]

numbers = list(map(int, prices))
print(numbers)
print(sum(numbers))
\end{example}

\iosample{}{[12, 5, 38] \\{} 55 \\}

\texttt{map(int, prices)} 的意思是「拿 \texttt{int} 這條規則，套到 \texttt{prices} 的每一
項上」。跟 \texttt{key=len} 一樣，\texttt{int} 後面沒有括號——交出去的是規則本身。外面那層
\texttt{list()} 等一下就會解釋為什麼需要。

\texttt{filter} 的形狀幾乎一模一樣，差別只在它交出去的函式要回答一個是非題：這一項要不要
留下來？回傳 \texttt{True} 就留，\texttt{False} 就丟掉。這種「一句話就講完」的條件，正是
\texttt{lambda} 的主場。

\begin{example}{python}{(passing-scores.py) 挑出及格的分數}
scores = [92, 45, 78, 59, 60]

passed = list(filter(lambda score: score >= 60, scores))
print(passed)
\end{example}

\iosample{}{[92, 78, 60] \\}

\texttt{lambda score: score >= 60} 對每個分數各跑一次，算出 \texttt{True} 或
\texttt{False}，\texttt{filter} 只把答 \texttt{True} 的那幾項留下來。

\subsection*{常見錯誤：忘記用 \texttt{list()} 包起來}

上面兩個範例外面那層 \texttt{list()} 不是裝飾用的。如果直接把 \texttt{map} 的結果印出來：

\begin{example}{python}{(convert-prices-broken.py) 直接印出 map 的結果}
prices = ["12", "5", "38"]

print(map(int, prices))
\end{example}

\iosample{}{<map object at 0x102bbfee0> \\}

（後面那串十六進位數字每次執行都不一樣，那是這個東西在記憶體裡的位置，不用管它。）

畫面上沒有你要的 \texttt{[12, 5, 38]}，而是一個看不出內容的東西。原因是 \texttt{map} 並
\textbf{不會馬上把結果全部算好交給你}，它回傳的是一個「等你來拿才現算」的東西。要真的把
每一項都算出來、擺成一份清單，就得用 \texttt{list()} 把它整個攤開。\texttt{filter} 也是
一樣的性質。

這種「先不急著算」的設計在 Python 裡其實到處都是，它有一整套背後的道理，完整的討論留到
第 \ref{ch:iterators-generators} 章。這裡你只要記住一件事：\texttt{map} 或 \texttt{filter} 的結果，
想看內容就用 \texttt{list()} 包起來。

\subsection*{同一件事的另一種寫法}

第 \ref{ch:lists} 章你已經看過 comprehension 的初探——那個把「造一份空清單、跑迴圈、逐一
append」濃縮成一行的寫法。上面挑及格分數那件事，用 comprehension 寫是這樣：

\begin{code}{python}{同一件事的兩種寫法}
passed = list(filter(lambda score: score >= 60, scores))

passed = [score for score in scores if score >= 60]
\end{code}

兩行做的事完全一樣，跑出來的結果也一樣（都是 \texttt{[92, 78, 60]}）。

\stylepoint{\texttt{map} 與 \texttt{filter} 的寫法沒有錯，但 Python 圈子普遍覺得
comprehension 比較好讀，實務上你在別人的 Python 程式碼裡看到 comprehension 的機會，也遠
多於看到 \texttt{map}／\texttt{filter}。所以這一節的定位是「看得懂就好」：讀到別人這樣寫
你要知道它在做什麼，但輪到你自己寫的時候，優先考慮 comprehension。它完整的寫法（包含為
什麼比較好讀）留到第 \ref{ch:comprehensions} 章。}

\begin{drills}[label=drills:map-filter]
\item \texttt{print(list(map(lambda n: n * 2, [3, 1, 2])))} 會印出什麼？
\item 有 \texttt{words = ["kiwi", "fig", "banana"]}，把
\texttt{list(filter(lambda w: len(w) > 3, words))} 改寫成 comprehension 的形式。
\end{drills}

\section{函式可以被生出來，而且會記得出生時看到的東西}
\label{sec:closure}

到目前為止，三節談的都是同一個方向：你手上有一條規則，把它交給別人寫好的函式去跑。這一節
整個翻面。

回到開場那疊表單。除了「把規則交出去」之外，還有另一種常見的情況：你請一位很懂行的同事
\textbf{幫你訂一條規則}。你跟他說「這次要照日期排」，他回頭給你的不是一份排好的表單，而是
一張寫著「照日期比大小」的卡片——一條你之後可以隨時拿去用、想用幾次就用幾次的規則。而且
關鍵是，這張卡片會記得他當初訂它的時候你說的那句「照日期」。下次你再請他訂一條「照金額」
的，他會給你另一張卡片，兩張互不干擾。

在 Python 裡，這件事寫出來就是：一個函式裡面再定義一個函式，然後把裡面那個函式
\textbf{回傳出去}。

\begin{code}{python}{在函式裡造一個函式，再把它回傳出去}
def outer_function(value):
    def inner_function(parameter):
        code using value and parameter

    return inner_function
\end{code}

最後那行 \texttt{return inner\_function} 沒有括號——這是第 \ref{sec:function-first-class} 節那條
判準的又一次應用：這裡要交回去的是函式物件本身，不是呼叫它的結果。

\begin{example}{python}{(make-greeter.py) 訂做兩句不同的招呼語}
def make_greeter(greeting):
    def greet(name):
        print(f"{greeting}, {name}!")

    return greet


morning = make_greeter("Good morning")
evening = make_greeter("Good evening")

morning("Amy")
evening("Amy")
morning("Ben")
\end{example}

\iosample{}{Good morning, Amy! \\ Good evening, Amy! \\ Good morning, Ben! \\}

一行一行看清楚發生了什麼事。\texttt{make\_greeter("Good morning")} 這一行做了三件事：
Python 進到 \texttt{make\_greeter} 裡面，先把 \texttt{greeting} 這張名字牌貼到字串
\texttt{"Good morning"} 上，接著執行 \texttt{def greet(name):} 造出一個新的函式物件，最後
把這個函式物件回傳出來，被外面的 \texttt{morning} 接住。\texttt{make\_greeter} 到這裡就
執行完畢、結束了。

然後我們又呼叫了一次 \texttt{make\_greeter("Good evening")}。這一次是全新的一輪呼叫，
Python 會重新執行一次那個 \texttt{def greet(name):}，造出\textbf{另一個}函式物件，交給
\texttt{evening}。所以 \texttt{morning} 跟 \texttt{evening} 是兩張貼在\textbf{兩個不同函式
物件}上的名字牌，儘管它們的程式碼長得一模一樣。

真正有意思的是輸出：\texttt{morning("Amy")} 印的是 Good morning，\texttt{evening("Amy")}
印的是 Good evening。兩個函式各自記得了自己出生時看到的那句招呼語，而且互不干擾。這種
「記得出生時周遭環境」的函式，有一個正式名字叫 \textbf{closure}（閉包）。
\index{closure（閉包）}

\subsection*{它為什麼記得？名字牌最後一次派上用場}

這裡有一件事跟直覺正面衝突，值得慢慢講。

第 \ref{sec:function-scope} 節說得很清楚：函式裡面的區域變數，只在那一次呼叫執行期間存在，
函式執行完就消失，外面完全存取不到。而 \texttt{greeting} 就是 \texttt{make\_greeter} 的
一個區域變數（它是參數，參數也是區域變數）。按照這條規則，\texttt{make\_greeter} 早在
\texttt{morning} 接住回傳值的那一刻就已經執行完畢、\texttt{greeting} 應該早就不見了才對——那 \texttt{morning("Amy")}
執行到 \texttt{f"\{greeting\}, \{name\}!"} 的時候，是從哪裡撈出那句 Good morning 的？

答案要用「名字牌」來講，才會完全講得通。這是我們從第 \ref{ch:variables} 章一路用到現在的那個
比喻——變數不是裝東西的盒子，是貼在物件上的名字牌——而這裡是它在這本書裡\textbf{最後一次}
派上用場。

真正消失的，是 \texttt{make\_greeter} 那一輪呼叫所使用的那個區域名字空間，不是名字牌指著
的\textbf{物件}。\texttt{greet} 這個內層函式在被造出來的時候，帶走的不是
\texttt{"Good morning"} 這個字串的一份副本，而是 \texttt{greeting} 那張名字牌本身。只要
還有函式握著那張名字牌，名字牌指著的那個字串物件就不會不見——Python 不會把還有人握著的
東西丟掉。

所以第二次呼叫 \texttt{make\_greeter} 的時候，發生的是一整套全新的事：新的一輪呼叫、新的
一張 \texttt{greeting} 名字牌、貼在新的字串物件 \texttt{"Good evening"} 上、新造出來的
\texttt{greet} 函式物件握著這張新名字牌。兩輪呼叫從頭到尾沒有共用任何東西，這就是為什麼
\texttt{morning} 跟 \texttt{evening} 各說各話，一點也不會混。

\keypoint{closure 記住的是\textbf{名字牌}，不是值的副本。外層函式執行完畢之後，它的區域
名字空間確實消失了，但只要內層函式還握著其中某張名字牌，那張名字牌指著的物件就會留著。
這也是為什麼每呼叫一次外層函式，就會得到一個記得不同東西的內層函式——每一輪呼叫都有自己
全新的一組名字牌。}

\subsection*{只讀沒問題，要改就出事}

\texttt{make-greeter.py} 裡的內層函式只是\textbf{讀}外層的 \texttt{greeting}，從頭到尾
沒有想改它。如果內層函式想\textbf{改}外層的東西呢？來寫一個計數器試試看：每呼叫一次就
往上加一。

\begin{example}{python}{(make-counter-broken.py) 內層函式想修改外層的變數}
def make_counter():
    count = 0

    def increase():
        count = count + 1
        print(count)

    return increase


counter = make_counter()
counter()
\end{example}

\begin{verbatim}
Traceback (most recent call last):
  File "make-counter-broken.py", line 12, in <module>
    counter()
    ~~~~~~~^^
  File "make-counter-broken.py", line 5, in increase
    count = count + 1
            ^^^^^
UnboundLocalError: cannot access local variable 'count' where it is not associated with a value
\end{verbatim}

這位是老朋友。第 \ref{sec:function-scope} 節談 \texttt{global} 的時候，你看過一模一樣的
\texttt{UnboundLocalError}，原因也完全一樣：\textbf{只要函式主體裡有任何一行對某個名字
賦值，Python 就把這個名字在整段主體裡都當成區域變數}。\texttt{count = count + 1} 的等號
左邊有 \texttt{count}，所以在 \texttt{increase} 眼中，\texttt{count} 是一個屬於它自己的
新名字——跟外層那個完全無關。可是等號右邊又要先讀取 \texttt{count} 才能算，而這個「內層
版本」的 \texttt{count} 到那一刻根本還沒被賦值過，於是當場報錯。

修法也跟 \texttt{global} 是同一個路數，只是換一個關鍵字：\texttt{nonlocal}。
\index{nonlocal}

\begin{example}{python}{(make-counter.py) 用 nonlocal 修改外層函式的變數}
def make_counter():
    count = 0

    def increase():
        nonlocal count
        count = count + 1
        print(count)

    return increase


counter = make_counter()
counter()
counter()

another = make_counter()
another()
\end{example}

\iosample{}{1 \\ 2 \\ 1 \\}

有了 \texttt{nonlocal count} 這一行，\texttt{increase} 裡的 \texttt{count} 不再是一張新的
內層名字牌，而是明確指向\textbf{外層那一輪呼叫}的那張 \texttt{count}。每次呼叫
\texttt{counter()}，改動的都是同一張名字牌，所以數字才會一路累加上去。

最後三行是關鍵的對照：\texttt{another = make\_counter()} 是全新的一輪呼叫，配一組全新的
名字牌，所以 \texttt{another()} 從 \texttt{1} 開始數，跟已經數到 \texttt{2} 的
\texttt{counter} 毫無關係。

\texttt{global} 跟 \texttt{nonlocal} 是同一件事的兩個層次：\texttt{global} 說的是「我要動
的是\textbf{最外層那個全域}的名字牌」，\texttt{nonlocal} 說的是「我要動的是\textbf{外面那
層函式}的名字牌」。也因為這樣，\texttt{nonlocal} 只有在「函式裡面又定義了一個函式」的時候
才有意義——這正是第 \ref{ch:functions} 章沒有教它的原因，那時候還沒有這種場景。

而範例 \texttt{make-greeter.py} 之所以完全不需要 \texttt{nonlocal}，就是因為它只
\textbf{讀}外層的 \texttt{greeting}，沒有對它賦值。這條界線很簡單，也就是 \texttt{nonlocal}
存在的全部理由：\textbf{只讀不用宣告，要改才要宣告}。

\begin{drills}[label=drills:closure]
\item 有以下程式，\texttt{print(double(10), triple(10))} 會印出什麼？
\begin{quote}
\texttt{def make\_multiplier(factor):} \\
\texttt{~~~~def multiply(n):} \\
\texttt{~~~~~~~~return n * factor} \\
\texttt{~~~~return multiply} \\
\texttt{double = make\_multiplier(2)} \\
\texttt{triple = make\_multiplier(3)}
\end{quote}
\item 承 \texttt{make-counter.py}：如果連續呼叫三次 \texttt{counter()}，第三次會印出什麼？
如果這時候才建立 \texttt{another = make\_counter()} 並呼叫一次，會印出什麼？
\item 有人照著 \texttt{make-counter.py} 寫，但把 \texttt{nonlocal count} 那一行漏掉了。
執行時會出什麼錯？為什麼？
\end{drills}

\section{自己動手替函式包一層}
\label{sec:wrap-by-hand}

現在把前面兩件事接起來：函式可以被當引數收下（第 \ref{sec:lambda}～\ref{sec:map-filter} 節），
函式也可以被造出來、被回傳（第 \ref{sec:closure} 節）。兩者合在一起，就能做到一件很有用
的事——在一個已經能用的函式外面「包一層」。

情境是這樣的：你手上有一個寫好、也測試過、正常運作的函式。現在你想在它執行前後各印一行
記錄，方便追蹤程式跑到哪裡了。但你\textbf{不想動它原本的程式碼}——可能是那段邏輯很敏感、
可能是別人寫的、也可能是你想在很多函式上都做同一件事，不想每個都改一遍。

做法就是開場類比的第三個方向：原本的做法不動，只在前後各補一道手續，交出去的變成「包過的
那一份」。寫成 Python，就是寫一個函式，它收下原本那個函式，造一個新的函式（新的裡面前後
各印一行、中間呼叫原本那個），再把新的交回來。這正是上一節 closure 的形狀，只是這次被
記住的東西不是一句招呼語，而是\textbf{一個函式}。

\begin{example}{python}{(wrap-by-hand.py) 手動替函式包上一層記錄}
def add_log(func):
    def wrapper(name):
        print("[log] start")
        result = func(name)
        print("[log] done")
        return result

    return wrapper


def greet(name):
    return f"Hello, {name}!"


greet = add_log(greet)

print(greet("Amy"))
\end{example}

\iosample{}{[log] start \\{} [log] done \\ Hello, Amy! \\}

慢慢拆。\texttt{add\_log} 是一個收下函式、回傳函式的函式：它的參數 \texttt{func} 接住被傳
進來的那個函式物件，然後定義一個叫 \texttt{wrapper} 的內層函式。\texttt{wrapper} 做三件事
——先印一行「開始」，接著呼叫 \texttt{func(name)} 做原本該做的事、把結果接住，再印一行
「完成」，最後把結果原封不動 \texttt{return} 出去。因為 \texttt{wrapper} 是一個 closure，
它記得住 \texttt{func} 這張名字牌指著的是哪個函式。最後 \texttt{add\_log} 回傳
\texttt{wrapper} 本身（沒有括號）。

輸出的順序值得確認一次：兩行記錄先出現，\texttt{Hello, Amy!} 最後才印。這是因為
\texttt{greet("Amy")} 整個跑完（包含兩行記錄）之後才把字串交出來，最外面那個
\texttt{print} 這時候才有東西可印。

\subsection*{整章的樞紐：\texttt{greet = add\_log(greet)}}

上面那支程式裡，最需要盯著看的是這一行：

\begin{code}{python}{把包好的那一個重新貼回原本的名字牌}
greet = add_log(greet)
\end{code}

等號右邊 \texttt{add\_log(greet)} 先跑：把原本的 \texttt{greet} 函式物件交給
\texttt{add\_log}，拿回一個包好的 \texttt{wrapper} 函式物件。然後等號把 \texttt{greet} 這張
名字牌\textbf{撕下來、改貼到那個新的函式物件上}。

從這一行之後，\texttt{greet} 指的已經不是原本那個函式了，而是包過的那一個。但原本那個函式
並沒有消失——它還被 \texttt{wrapper} 握在手上（就是 \texttt{func} 那張名字牌），只是外面
的人喊不到它了。

這一行的巧妙之處在於：\textbf{呼叫端完全不用改}。程式其他地方原本怎麼寫 \texttt{greet(...)}，
現在還是那樣寫，只是行為多了兩行記錄。換一張名字牌這件事，你在第 \ref{sec:function-first-class} 節
就學過了（\texttt{operation = square}），只是當時大概沒想過可以這樣用。

\keypoint{這裡的 \texttt{wrapper} 為了看得清楚，固定只收一個參數 \texttt{name}，所以
\texttt{add\_log} 只能拿來包「剛好收一個引數」的函式。真正在專案裡通用的寫法，內層會用
第 \ref{sec:function-args} 節的 \texttt{*args}／\texttt{**kwargs} 收下任意引數，這樣不管被包的
函式長什麼樣子都應付得了。你以後在別人的程式碼裡看到滿滿的星號，不用意外——那就是這裡的
\texttt{name} 換成了「什麼都收」的寫法而已。}

\subsection*{常見錯誤：忘記 \texttt{return wrapper}}

\texttt{add\_log} 最後那一行 \texttt{return wrapper} 很容易漏掉，畢竟前面的縮排區塊看起來
已經「寫完了」。漏掉的話會這樣：

\begin{example}{python}{(wrap-by-hand-broken.py) add\_log 忘記把 wrapper 回傳出去}
def add_log(func):
    def wrapper(name):
        print("[log] start")
        result = func(name)
        print("[log] done")
        return result


def greet(name):
    return f"Hello, {name}!"


greet = add_log(greet)

print(greet("Amy"))
\end{example}

\begin{verbatim}
Traceback (most recent call last):
  File "wrap-by-hand-broken.py", line 15, in <module>
    print(greet("Amy"))
          ~~~~~^^^^^^^
TypeError: 'NoneType' object is not callable
\end{verbatim}

這個錯誤訊息值得學會怎麼讀，因為它指的地方不是真正出錯的地方。整句話翻成中文是「\texttt{None}
這種東西不能被呼叫」。從哪裡冒出一個 \texttt{None}？回想第 \ref{ch:functions} 章講過的規則：
函式如果沒有寫 \texttt{return}，呼叫它得到的值就是 \texttt{None}。\texttt{add\_log} 現在正是
這種函式，所以 \texttt{add\_log(greet)} 交回來的是 \texttt{None}，接著
\texttt{greet = add\_log(greet)} 把 \texttt{greet} 這張名字牌貼到了 \texttt{None} 身上。等到
下一行寫 \texttt{greet("Amy")}，等於是在呼叫 \texttt{None}——當然不行。

注意裡面那個 \texttt{def wrapper} 本身完全沒問題，它確實被造出來了，只是造完就被丟在
\texttt{add\_log} 裡面沒交出來。看到 \texttt{'NoneType' object is not callable}，第一個要
檢查的就是：那個應該回傳函式的地方，是不是漏了 \texttt{return}？

\section{\texttt{@}：包一層的簡寫}
\label{sec:decorator}

上一節那支程式完全能動，但寫起來有點囉嗦：函式定義完之後，還要記得在下面補一行
\texttt{greet = add\_log(greet)}。而且這一行離 \texttt{def greet} 可能隔了一段距離，讀的人
不一定會注意到「這個函式其實被包過了」。

Python 為這個固定的動作準備了一個專用的簡寫。把上一節的最後兩步：

\begin{code}{python}{手動的寫法}
def greet(name):
    return f"Hello, {name}!"


greet = add_log(greet)
\end{code}

換成這樣：

\begin{code}{python}{用 @ 的簡寫}
@add_log
def greet(name):
    return f"Hello, {name}!"
\end{code}

在 \texttt{def} 上面那一行寫一個 \texttt{@} 加上函式名稱，意思就是：「這個函式定義完之後，
立刻做一次 \texttt{greet = add\_log(greet)}」。就這樣，沒有別的了。

\begin{example}{python}{(decorator-at.py) 用 @ 改寫上一節那支程式}
def add_log(func):
    def wrapper(name):
        print("[log] start")
        result = func(name)
        print("[log] done")
        return result

    return wrapper


@add_log
def greet(name):
    return f"Hello, {name}!"


print(greet("Amy"))
\end{example}

\iosample{}{[log] start \\{} [log] done \\ Hello, Amy! \\}

輸出跟 \texttt{wrap-by-hand.py} 一模一樣，因為做的事情本來就完全一樣。\texttt{add\_log} 這個
函式一個字都沒改，改的只有「怎麼把它套上去」的寫法。

這種寫在 \texttt{def} 上面、用來替函式包一層的東西，正式名稱叫 \textbf{decorator}
（裝飾器）。它在真實的 Python 專案裡幾乎無所不在——網頁框架用它來標記「這個函式負責處理
哪個網址」，測試工具用它來標記「這是一個要跑的測試」，權限檢查、快取、記錄，全都靠它。
以後你打開任何一個像樣的 Python 專案，看到滿滿的 \texttt{@} 開頭的行，現在你知道那是什麼
了：每一個都只是「定義完之後，把這個函式交給某個函式包一層，再貼回原本的名字」。
\index{decorator（裝飾器）}

\keypoint{\texttt{@add\_log} 不是一種新的語言機制，它只是
\texttt{greet = add\_log(greet)} 的簡寫——這種「只是為了寫起來順手，沒有新增任何能力」的
語法，一般叫做語法糖（syntactic sugar）。所以看到任何 \texttt{@} 開頭的東西，你都可以在
腦子裡翻譯回第 \ref{sec:wrap-by-hand} 節那個手動的兩步形狀，再回頭想它到底做了什麼。}
\index{語法糖（syntactic sugar）}

要說明的是，decorator 這個主題還可以做得更花——讓 decorator 自己也帶參數、一個函式上疊好
幾個 decorator、用它來包整個 class。這些寫法的地基都是這一節的那個等式，等你真的需要的
時候再學不遲。這本書要你精熟的只有一件事：\texttt{@deco} 就是 \texttt{f = deco(f)}，
兩邊可以互相翻譯。

\begin{drills}[label=drills:decorator]
\item 把下面這段用 \texttt{@} 的寫法，改寫成不用 \texttt{@} 的兩步寫法：
\begin{quote}
\texttt{@shout} \\
\texttt{def say(word):} \\
\texttt{~~~~return f"hello \{word\}"}
\end{quote}
\item 承上題，如果 \texttt{shout} 的定義是下面這樣，\texttt{print(say("amy"))} 會印出什麼？
\begin{quote}
\texttt{def shout(func):} \\
\texttt{~~~~def wrapper(word):} \\
\texttt{~~~~~~~~return func(word).upper()} \\
\texttt{~~~~return wrapper}
\end{quote}
\item 有一個 decorator 叫 \texttt{add\_note}，它的 \texttt{wrapper} 裡寫的是
\texttt{return func(text) + " (checked)"}。現在有一個函式：
\begin{quote}
\texttt{def clean(text):} \\
\texttt{~~~~return text.strip()}
\end{quote}
如果用手動的方式套上去（也就是補一行 \texttt{clean = add\_note(clean)}），那麼
\texttt{print(clean("~~hi~~"))} 會印出什麼？
\end{drills}

\section{幫函式的參數與回傳值貼上標籤：type hints}
\label{sec:type-hints}

前面六節談的都是「拿函式來用」。這一節換個方向，回頭談「怎麼描述一個函式」。

第 \ref{sec:function-docstring} 節教過 docstring：在函式主體第一行寫一段說明文字，讓下一個
讀到這段程式碼的人知道它在做什麼。那是給人看的散文。Python 還準備了一種更短、更結構化的
寫法，可以直接寫在 \texttt{def} 那一行上，讓人一眼看出「這個函式吃什麼、吐什麼」——這種
標註叫做 type hint（型別標註）。
\index{type hint（型別標註）}

\begin{code}{python}{在 def 那一行標註參數與回傳值的型別}
def function_name(parameter: type) -> return_type:
    code
\end{code}

參數名稱後面加一個冒號、接型別名稱；整串參數的右括號後面用 \texttt{->} 接回傳值的型別，
再接原本那個冒號。

\begin{example}{python}{(total-price-hint.py) 替兩個小函式加上標註}
def total_price(prices: list[float]) -> float:
    return sum(prices)


def make_label(name: str, count: int) -> str:
    return f"{name} x {count}"


print(total_price([12.5, 8.0, 30.0]))
print(make_label("notebook", 3))
\end{example}

\iosample{}{50.5 \\ notebook x 3 \\}

\texttt{def total\_price(prices: list[float]) -> float:} 這一行，不用看函式主體就讀得懂：
它收一份裝著小數的清單，吐回一個小數。可以用的型別就是你已經認識的那些——\texttt{int}、
\texttt{float}、\texttt{str}、\texttt{bool}，容器則寫成 \texttt{list[str]}、
\texttt{list[float]} 這種「清單，裡面裝的是字串／小數」的形式。

順帶一提，如果一個函式只負責印東西、沒有 \texttt{return} 任何值（第 \ref{ch:functions} 章
說過，這種函式呼叫後拿到的是 \texttt{None}），標註就寫成 \texttt{-> None}。這在別人的程式
碼裡很常見，看到不用緊張。

\subsection*{標錯了會怎樣？答案是什麼事都不會發生}

看到 \texttt{name: str} 這種寫法，直覺上會覺得：既然都寫出來了，那傳錯東西進去 Python 應該
會擋下來吧？來實測一次。

\begin{example}{python}{(hint-not-checked.py) 故意傳跟標註完全不符的東西}
def make_label(name: str, count: int) -> str:
    return f"{name} x {count}"


print(make_label("notebook", 3))
print(make_label(99, "three"))
\end{example}

\iosample{}{notebook x 3 \\ 99 x three \\}

第二次呼叫，\texttt{name} 收到的是數字 \texttt{99}、\texttt{count} 收到的是字串
\texttt{"three"}，兩個都跟標註完全相反。程式照跑不誤，連一句警告都沒有。

\keypoint{Python \textbf{完全不會}因為你標錯型別就報錯。type hint 對程式執行時的行為
零影響——Python 執行的時候根本不看它們。這件事跟直覺衝突（「寫了就會檢查吧？」），但務必
記牢：它們是標籤，不是門禁。程式會不會出事，取決於函式主體裡那些運算撐不撐得住你傳進去
的東西，跟標註無關。}

\subsection*{那它到底有什麼用}

既然不檢查，為什麼還要寫？因為它服務的不是 Python，是\textbf{人跟工具}。

第一，是給讀程式碼的人看的。看到 \texttt{def total\_price(prices: list[float]) -> float:}，
你不用跳進函式裡讀主體，也不用猜參數名稱的言外之意，一眼就知道該餵什麼、會拿回什麼。

第二，是給編輯器看的。現代的程式編輯器會讀這些標註，你打字打到一半，它就能提示這個變數
有哪些方法可以用；如果你把字串傳給一個標了 \texttt{int} 的參數，它也會先在畫面上畫一條
波浪線提醒你——注意，這是編輯器在提醒，不是 Python 在報錯。

第三，是團隊的事。在一個好幾個人一起維護、跑好幾年的專案裡，這種標註等於是
\textbf{用程式碼本身寫下來的約定}：這個函式承諾收什麼、承諾給什麼。約定寫在註解裡容易
過期，寫在 \texttt{def} 那一行上就跟著程式碼一起被讀、一起被改。也因為這樣，許多公司的
專案會要求新寫的函式一律加上標註。

\subsection*{常見錯誤：冒號打成等號}

有一個手滑很容易看漏，因為兩種寫法都不會報錯，但意思完全不同：

\begin{example}{python}{(hint-typo.py) 把冒號打成等號}
def show(count=int):
    print(count)


show()
\end{example}

\iosample{}{<class 'int'> \\}

\texttt{def show(count=int):} 裡的等號不是標註，是\textbf{預設值}——它把 \texttt{count} 的
預設值設成 \texttt{int} 這個函式物件本身（第 \ref{ch:variables} 章看過，\texttt{print(type(x))}
印出來的 \texttt{<class 'int'>} 就是這個東西）。所以不給引數呼叫 \texttt{show()}，印出來的
就是這個奇怪的東西。真正的標註寫法是冒號：\texttt{def show(count: int):}。

\begin{drills}[label=drills:type-hints]
\item 有 \texttt{def double(n: int) -> int: return n * 2}。呼叫
\texttt{print(double("ab"))} 會不會報錯？會印出什麼？
\item 看到 \texttt{def find\_names(records: list[str], keyword: str) -> list[str]:}
這一行，在完全不看函式主體的情況下，你知道這個函式要收什麼、會給你什麼？
\end{drills}

\section*{練習解答}
\addcontentsline{toc}{section}{練習解答}

\paragraph{第~\ref{sec:lambda}~節}
\begin{enumerate}
\item 印出 \texttt{12}。\texttt{lambda} 也可以有好幾個參數，用逗號隔開就行。
\item \texttt{def square(n): return n * n}（寫成兩行：\texttt{def square(n):} 換行縮排
\texttt{return n * n}）。
\end{enumerate}

\paragraph{第~\ref{sec:key-function}~節}
\begin{enumerate}
\item 印出 \texttt{['fig', 'kiwi', 'banana']}：三個字的長度分別是 6、3、4，照長度由小到大
排就是 fig、kiwi、banana。
\item 印出 \texttt{Ben}：照 \texttt{age} 排完之後排第一個的是 24 歲的 Ben，再取他的
\texttt{name}。
\item 多打了一對括號。\texttt{item["price"]} 取出來的是一個數字，後面加上 \texttt{()}
等於是在呼叫這個數字，會丟出 \texttt{TypeError}。把括號拿掉就對了。
\end{enumerate}

\paragraph{第~\ref{sec:map-filter}~節}
\begin{enumerate}
\item 印出 \texttt{[6, 2, 4]}：每一項各乘 2，順序不變（\texttt{map} 不排序）。
\item \texttt{[w for w in words if len(w) > 3]}，結果是
\texttt{['kiwi', 'banana']}。
\end{enumerate}

\paragraph{第~\ref{sec:closure}~節}
\begin{enumerate}
\item 印出 \texttt{20 30}。\texttt{double} 跟 \texttt{triple} 是兩次不同呼叫造出來的兩個
函式物件，各自握著自己那張 \texttt{factor} 名字牌（分別指著 2 跟 3）。
\item 第三次呼叫 \texttt{counter()} 印出 \texttt{3}；\texttt{another()} 印出 \texttt{1}，
因為 \texttt{another} 來自全新的一輪 \texttt{make\_counter()} 呼叫，配的是一組全新的名字
牌，跟 \texttt{counter} 完全無關。
\item 會丟出 \texttt{UnboundLocalError}。少了 \texttt{nonlocal}，\texttt{count = count + 1}
這一行的賦值會讓 Python 把 \texttt{count} 當成 \texttt{increase} 自己的區域變數，可是等號
右邊又要先讀它，而它還沒被賦值過。
\end{enumerate}

\paragraph{第~\ref{sec:decorator}~節}
\begin{enumerate}
\item 先照常定義 \texttt{def say(word): return f"hello \{word\}"}，然後在下面補一行
\texttt{say = shout(say)}。
\item 印出 \texttt{HELLO AMY}：\texttt{wrapper} 先呼叫原本的 \texttt{say} 拿到
\texttt{"hello amy"}，再對它做 \texttt{.upper()}。
\item 印出 \texttt{hi (checked)}：原本的 \texttt{clean} 先把前後空白去掉得到
\texttt{"hi"}，包在外面那層再接上 \texttt{" (checked)"}。
\end{enumerate}

\paragraph{第~\ref{sec:type-hints}~節}
\begin{enumerate}
\item 不會報錯，會印出 \texttt{abab}。標註 \texttt{int} 對執行完全沒有影響，而
\texttt{"ab" * 2} 剛好是合法的字串重複運算，所以程式順利跑完——即使它回傳的其實是字串，
跟標註寫的 \texttt{int} 不符。
\item 它收兩個東西：一份裝著字串的清單 \texttt{records}，跟一個字串 \texttt{keyword}；
回傳的是另一份裝著字串的清單。
\end{enumerate}

\chapter{Comprehension 與內建工具函式}
\label{ch:comprehensions}

請人幫忙的時候，你講的是「我要什麼」，不是「你的手要怎麼動」。把一疊報名表推過去，你會說
「幫我把及格的挑出來」——一句話，講完了。你不會說：「先拿一張空白紙放旁邊，然後從第一張
開始，一張一張看，看到及格的就抄到那張紙上，抄完再翻下一張。」後面這一串話每個步驟都沒
錯，做出來的結果也一模一樣，只是聽的人得撐到最後一句，才知道你到底要什麼。

程式碼也有這兩種講法，而你已經寫過後面那一種了。回頭看第 \ref{sec:csv-dict} 節那支篩選
及格名單的 \texttt{filter-and-write-passed.py}：為了從一份成績檔案裡挑出及格的人，中間
這四行是這樣寫的：

\begin{code}{python}{你在成績檔案那一章親手打過的四行}
passed_students = []
for row in reader:
    if row["passed"] == "Yes":
        passed_students.append(row)
\end{code}

這四行到底在說什麼？答案是「我要及格的那些」。但這句話並沒有寫在任何一行裡——它散在四行
之間，你得先看到有人造了一份空清單，再看到有個迴圈在走，再看到一個判斷，最後看到
\texttt{append}，把這四件事在腦子裡合起來，才拼得出那句話。

同一件事，還可以這樣寫：

\begin{code}{python}{同一件事的一行版}
passed_students = [row for row in reader if row["passed"] == "Yes"]
\end{code}

這一行從左讀到右，就是「我要每一個 \texttt{row}，從 \texttt{reader} 裡面來，條件是它的
\texttt{passed} 欄位等於 Yes」。方括號本身就先告訴你「結果是一份清單」，剩下的只是填空。
這個寫法你在第 \ref{sec:list-comprehension} 節見過最陽春的版本，當時說「這裡只教最簡單的
形式，不帶條件式、也不是雙層迴圈，那些之後有需要再說」——現在就是之後。

這一章分四段推進：

\begin{itemize}
\item 第一段把 comprehension 本身補到完整：在原本的骨架上依序裝三個零件——後面的
\texttt{if}（挑掉一些）、前面的 \texttt{if}／\texttt{else}（每一項換個值）、第二個
\texttt{for}（多走一層）（第 \ref{sec:comp-filter}～\ref{sec:comp-nested} 節）。
\item 第二段把鏡頭拉遠，看「comprehension 這個形狀本身」：同樣的內容換一副括號，拿到的
東西就不一樣了——大括號給你 dict 或 set，小括號給你一種你還沒見過的東西
（第 \ref{sec:comp-dict-set}～\ref{sec:genexp-glimpse} 節）。
\item 第三段反過來問：如果你根本不想要那個新容器呢？有一組內建函式吃進一整串東西、只吐
一個答案，另一組則是改造「你要走訪的那一串」本身
（第 \ref{sec:reducing-builtins}～\ref{sec:zip-walk} 節）。
\item 第四段是煞車。學會一個好用的寫法之後，同樣重要的是知道什麼時候不要用它
（第 \ref{sec:comp-limits} 節）。
\end{itemize}

\section{加上條件：只留下要的那些}
\label{sec:comp-filter}

先把第 \ref{sec:list-comprehension} 節那個骨架喚回來：\texttt{[expression for item in
iterable]}，意思是「把 \texttt{iterable} 裡的每一項各拿去算一次 \texttt{expression}，
算出來的結果收成一份新清單」。那一節的版本有一個限制——來幾項就收幾項，一項都不能少。

現在把第一個零件裝上去：在最後面補一個 \texttt{if}。

\begin{code}{python}{帶條件的 list comprehension}
[expression for item in iterable if condition]
\end{code}
\index{帶條件的 comprehension}

多出來的 \texttt{if condition} 是一道關卡：每一項都要先通過它，通過的才會被收下來，沒通過
的直接跳過。所以這種寫法造出來的清單，長度通常會比原本那份短。

\begin{example}{python}{(passed-scores.py) 挑出及格的分數：兩種寫法並排}
scores = [88, 76, 92, 65, 79]

# 傳統寫法：造一份空清單、跑迴圈、符合條件才 append
passed = []
for score in scores:
    if score >= 70:
        passed.append(score)
print(passed)

# comprehension 寫法：同一件事，一行
passed_v2 = [score for score in scores if score >= 70]
print(passed_v2)
\end{example}

\iosample{}{[88, 76, 92, 79] \\{} [88, 76, 92, 79] \\}

這份分數就是第 \ref{ch:csv} 章那份 \texttt{quiz\_scores.csv} 裡的五個人（Alice、Bob、
Taylor、Charlie、David），及格線也沿用第 \ref{sec:csv-module-write} 節的 70 分，所以
Charlie 的 65 分被擋在外面。兩段程式的輸出一模一樣，而下面那一行跟上面那五行的對應關係
很整齊：\texttt{for score in scores} 就是原本的迴圈標頭，\texttt{if score >= 70} 就是原本
那個判斷，最前面那個 \texttt{score} 就是原本 \texttt{append} 進去的東西。位置換了，內容
一個也沒少。

真正好用的是——這兩格可以同時派上用場，而且互不干涉：\textbf{最前面那一格負責「變成
什麼」，最後面那一格負責「留不留」}。從檔案讀進來的資料通常又髒又亂，正好兩件事都要做。

\begin{example}{python}{(clean-scores.py) 一邊過濾髒資料，一邊轉成整數}
raw = ["88", "76", "", "92", "n/a", "65", "79"]

numbers = [int(s) for s in raw if s.isdigit()]

print(numbers)
print(len(raw), "->", len(numbers))
\end{example}

\iosample{}{[88, 76, 92, 65, 79] \\{} 7 -> 5 \\}

\texttt{raw} 裡有七項，其中一項是空字串、一項是 \texttt{"n/a"}，這兩個丟給
\texttt{int()} 都會出事。\texttt{if s.isdigit()} 這道關卡先把它們擋下來（\texttt{isdigit()}
是字串的方法，回答「這串字是不是每個字元都是數字」），通過的五項才輪到最前面的
\texttt{int(s)} 做轉換。順序上，Python 一定是先過關卡、再做轉換，所以那兩個壞資料根本沒有
機會被 \texttt{int()} 碰到。

\subsection*{常見錯誤：把 \texttt{if} 搬到前面去}

看多了之後，很多人會產生一種印象：\texttt{if} 好像放哪裡都行。試試看把它挪到最前面：

\begin{example}{python}{(filter-oops.py) 把 if 挪到 for 前面}
raw = ["88", "", "76"]

numbers = [s if s.isdigit() for s in raw]
print(numbers)
\end{example}

\begin{verbatim}
  File "filter-oops.py", line 3
    numbers = [s if s.isdigit() for s in raw]
               ^^^^^^^^^^^^^^^^
SyntaxError: expected 'else' after 'if' expression
\end{verbatim}

程式連跑都沒開始跑就被擋下來了。這次的錯誤訊息很不客氣，但也很好懂：「\texttt{if} 這個
式子後面應該要有一個 \texttt{else}」。這句話其實透露了一件重要的事——\texttt{if} 放到前面
去，Python 讀到的\textbf{不是}原本那個過濾用的 \texttt{if}，而是另一種完全不同的寫法，
而且那種寫法一定要配 \texttt{else}。下一節就講它。

\begin{drills}[label=drills:comp-filter]
\item 預測輸出：\texttt{nums = [3, 8, 1, 10, 6]}，那麼
\texttt{print([n for n in nums if n > 5])} 會印出什麼？
\item 改寫題：把下面這段迴圈改寫成一行 comprehension。
\begin{verbatim}
words = ["apple", "fig", "banana", "kiwi"]
long_words = []
for w in words:
    if len(w) > 4:
        long_words.append(w.upper())
\end{verbatim}
\item 抓蟲：有人想從 \texttt{words} 裡挑出長度大於 3 的字，寫成
\texttt{[w if len(w) > 3 for w in words]}，執行時直接報錯。錯在哪裡？該怎麼改？
\end{drills}

\section{前面的 \texttt{if}：兩個結果選一個}
\label{sec:conditional-expression}

上一節最後那個 \texttt{SyntaxError} 說「\texttt{if} 這個式子後面應該要有一個
\texttt{else}」。先把 comprehension 整個放到一邊，單獨把那個東西看清楚。

你早就會寫這種程式：依照一個條件，決定某個變數要拿到兩個值裡的哪一個。

\begin{code}{python}{你熟悉的四行寫法}
if score >= 70:
    label = "Pass"
else:
    label = "Fail"
\end{code}

這四行做的事其實只有一件：\textbf{算出一個值，指定給 \texttt{label}}。Python 準備了一個
專門寫這種情況的簡短寫法，叫做\textbf{條件運算式}。
\index{條件運算式（conditional expression）}

\begin{code}{python}{條件運算式：兩個值裡選一個}
value_if_true if condition else value_if_false
\end{code}

順序看起來有點怪，但照著字面朗讀一次就順了：「\texttt{label} 拿到 \texttt{"Pass"}，
如果分數大於等於 70；否則拿到 \texttt{"Fail"}。」中文講「如果……就……」，這個寫法講的是
「就……如果……否則……」，把結果擺在最前面。

\begin{example}{python}{(pass-label.py) 四行與一行，做的是同一件事}
score = 88

if score >= 70:
    label = "Pass"
else:
    label = "Fail"
print(label)

label = "Pass" if score >= 70 else "Fail"
print(label)
\end{example}

\iosample{}{Pass \\ Pass \\}

這個寫法你其實見過一次。第 \ref{sec:csv-module-write} 節那支
\texttt{write-quiz-results.py} 裡，替每個人算出「及格與否」那一欄的，就是
\texttt{passed = "Yes" if score >= 70 else "No"} 這一行——當時沒有特別停下來解釋，現在
你知道它是什麼了。

有一條界線要先講清楚：條件運算式的兩邊都只能是「算得出一個值的東西」。只要有一邊要做別的
事——印東西、把資料 \texttt{append} 進某份清單、開檔案——就得乖乖寫回四行的
\texttt{if}／\texttt{else}。這條限制跟第 \ref{sec:lambda} 節的 \texttt{lambda} 是同一個
道理：\textbf{式子算值，敘述做事，兩者不能混著用}。

現在把它裝進 comprehension 的第一格：

\begin{example}{python}{(score-labels.py) 兩種 if 放在一起比較}
scores = [88, 76, 92, 65, 79]

labels = [("Pass" if score >= 70 else "Fail") for score in scores]
passed_only = [score for score in scores if score >= 70]

print(labels)
print(passed_only)
print(len(scores), len(labels), len(passed_only))
\end{example}

\iosample{}{['Pass', 'Pass', 'Pass', 'Fail', 'Pass'] \\{} [88, 76, 92, 79] \\{} 5 5 4 \\}

\texttt{("Pass" if score >= 70 else "Fail")} 外面那對括號不是必要的，拿掉照樣能跑；加上去
只是為了讓「這一整串是第一格的內容」看得更清楚。剛開始寫的時候建議加，等眼睛習慣了再說。

最後那一行是這一節的重點：原本五項，貼標籤的版本還是五項，過濾的版本剩四項。

\keypoint{同一個 \texttt{if}，放的位置不同，做的事完全不同。
\texttt{[x for x in items if condition]}——\texttt{if} 在\textbf{後面}，是\textbf{過濾}，
每一項要嘛原封不動被收下、要嘛整個消失，所以結果\textbf{會變短}。
\texttt{[A if condition else B for x in items]}——\texttt{if} 在\textbf{前面}，是
\textbf{選值}，一項都不會少，結果的長度\textbf{跟原本一樣}，只是每一項變成 A 或 B。
判斷的方法很簡單：看 \texttt{for} 在哪裡。\texttt{for} 前面的東西是「收下什麼」，
\texttt{for} 後面的 \texttt{if} 才是「收不收」。}

\subsection*{常見錯誤：兩種 \texttt{if} 混在一起}

這兩種寫法各有各的規矩，混起來用兩個方向都會被擋下來。想過濾卻把 \texttt{if} 寫到前面，
就是上一節那個錯誤，現在可以正式解釋了：

\begin{example}{python}{(label-oops.py) 想過濾，卻把 if 寫到前面}
scores = [88, 65, 92]

labels = [score if score >= 70 for score in scores]
print(labels)
\end{example}

\begin{verbatim}
  File "label-oops.py", line 3
    labels = [score if score >= 70 for score in scores]
              ^^^^^^^^^^^^^^^^^^^^
SyntaxError: expected 'else' after 'if' expression
\end{verbatim}

Python 看到第一格裡有 \texttt{if}，就認定你要寫的是條件運算式，於是等著看 \texttt{else}
——結果等到的是 \texttt{for}。想過濾的話，這個 \texttt{if} 應該搬到最後面去。

反過來，想選值卻把 \texttt{else} 掛在後面那個 \texttt{if} 上：

\begin{example}{python}{(filter-oops2.py) 想選值，卻把 else 掛在後面的 if 上}
scores = [88, 65, 92]

passed = [score for score in scores if score >= 70 else 0]
print(passed)
\end{example}

\begin{verbatim}
  File "filter-oops2.py", line 3
    passed = [score for score in scores if score >= 70 else 0]
                                                       ^^^^
SyntaxError: invalid syntax
\end{verbatim}

這次的訊息只說「文法不對」，比較不友善，但箭頭指得很準——問題就在那個 \texttt{else}。
後面那個 \texttt{if} 是關卡，關卡只回答「過或不過」，沒有第二條路可以走，所以它後面不接
\texttt{else}。要讓不及格的變成 0，得把整串搬到最前面去：
\texttt{[score if score >= 70 else 0 for score in scores]}。

\begin{drills}[label=drills:conditional-expression]
\item 有 \texttt{nums = [4, 9, 2, 7]}。下面兩行各自印出什麼？印出來的清單各有幾項？
\begin{verbatim}
print([n for n in nums if n > 5])
print([n if n > 5 else 0 for n in nums])
\end{verbatim}
\item 填空：有 \texttt{temps = [18, 25, 30]}，想得到
\texttt{['cool', 'warm', 'warm']}（25 度以上算 warm），下面這行該怎麼填？
\begin{verbatim}
labels = [______ for t in temps]
\end{verbatim}
\end{drills}

\section{兩層 \texttt{for}：走訪有層次的資料}
\label{sec:comp-nested}

第三個零件是「多走一層」。動手之前要先分清楚兩種長得很像、但不是同一回事的東西。

第 \ref{sec:list-comprehension} 節那支 \texttt{comprehension-grid.py} 寫的是
\texttt{[[0 for \_ in range(3)] for \_ in range(3)]}——注意裡面那一整串 comprehension
是被包在\textbf{最前面那一格}裡的，它是「每繞一圈就重新造一份新的內層清單」，做出來的
結果是一份清單中的清單。這一節要講的是另一種：\textbf{兩個 \texttt{for} 排在同一組
方括號裡}，做出來的結果是一份平的清單。前者是往下長，後者是往外攤。

\begin{code}{python}{兩層 for 的 comprehension}
[expression for outer in outers for inner in outer]
\end{code}

最常見的用途就是「攤平」：手上有一份第 \ref{sec:nested-list} 節那種清單中的清單，你想把
所有最裡層的東西倒進同一份清單裡。

\begin{example}{python}{(flatten-scores.py) 把每班一份的成績攤成一份}
class_scores = [[88, 92], [76], [95, 60, 71]]

# 兩層迴圈版
all_scores = []
for one_class in class_scores:
    for score in one_class:
        all_scores.append(score)
print(all_scores)

# comprehension 版：兩個 for 排在同一組方括號裡
all_scores_v2 = [score for one_class in class_scores for score in one_class]
print(all_scores_v2)
\end{example}

\iosample{}{[88, 92, 76, 95, 60, 71] \\{} [88, 92, 76, 95, 60, 71] \\}

兩份輸出一模一樣，連順序都一樣。把兩種寫法對著看，會發現一件事：\textbf{comprehension 裡
兩個 \texttt{for} 的先後順序，跟寫成巢狀迴圈時由外到內的順序完全一致}——外層的
\texttt{for one\_class in class\_scores} 在前，內層的 \texttt{for score in one\_class}
在後，一個字都沒調換。

那為什麼這種寫法常被說「順序很反直覺」？反直覺的不是兩個 \texttt{for}，是被搬到最前面
那一格的 \texttt{score}。它在迴圈版裡明明是最後才出現的（在最裡層的 \texttt{append}
裡），在 comprehension 裡卻跑到了整串的最前面。

\keypoint{讀一段兩層的 comprehension，眼睛不要從最左邊開始讀。先跳過第一格，從第一個
\texttt{for} 開始由左往右讀完，最後再回頭看最前面那一格。
\texttt{[score for one\_class in class\_scores for score in one\_class]} 要讀成：
「for one\_class in class\_scores $\to$ for score in one\_class $\to$ 收下 score」。
先走外層，再走內層，最後才是「收下什麼」——這個順序，就是你寫巢狀迴圈時本來的順序。}

順序不能亂調換，這件事與其抽象地講，不如直接跑一次給你看。把兩個 \texttt{for} 對調：

\begin{example}{python}{(flatten-swapped.py) 故意把兩個 for 對調}
class_scores = [[88, 92], [76], [95, 60, 71]]

all_scores = [score for score in one_class for one_class in class_scores]
print(all_scores)
\end{example}

\begin{verbatim}
Traceback (most recent call last):
  File "flatten-swapped.py", line 3, in <module>
    all_scores = [score for score in one_class for one_class in class_scores]
                                     ^^^^^^^^^
NameError: name 'one_class' is not defined
\end{verbatim}

\texttt{NameError: name 'one\_class' is not defined} 直譯是「\texttt{one\_class} 這個名字
沒有被定義過」。原因很單純：照剛才那個讀法由左往右讀，第一個 \texttt{for} 就要拿
\texttt{one\_class} 來走訪，可是負責產生 \texttt{one\_class} 的那個 \texttt{for} 排在
它後面，還沒輪到。這跟你寫巢狀迴圈時不可能把內層迴圈寫在外層前面，是同一件事。

最後說清楚這一節的定位：\textbf{兩層 comprehension 你要讀得懂，但不必急著自己寫}。看到
別人的兩層版本，能在腦子裡翻譯回巢狀迴圈就算過關；輪到自己寫的時候，寫成兩層縮排的
\texttt{for} 迴圈一點也不丟臉，很多時候還更好讀——第 \ref{sec:comp-limits} 節會把這件事
講成一條明確的規則。另外，這一章只教到兩層；真的需要三層以上，那就是「該寫回迴圈」的
訊號，不是「該再多塞一個 \texttt{for}」的訊號。

\begin{drills}[label=drills:comp-nested]
\item 翻譯題：把 \texttt{[c for word in words for c in word]} 改寫成兩層縮排的
\texttt{for} 迴圈（含 \texttt{append}）。如果 \texttt{words = ["ab", "cd"]}，結果會是
什麼？
\item 預測輸出：\texttt{rows = [[1, 2], [3, 4]]}，那麼
\texttt{print([n * 10 for row in rows for n in row])} 會印出什麼？特別注意順序。
\end{drills}

\section{換成大括號：dict 與 set comprehension}
\label{sec:comp-dict-set}

前三節都在拆解「方括號裡面」有哪些零件。這一節換個方向：零件不動，只換掉外面那副括號。

第 \ref{sec:comprehension} 節已經給過語法了，一句話喚回來就好：把方括號換成大括號，
\textbf{沒有冒號}就得到一個 set，\textbf{有冒號}把 key 跟 value 分開就得到一個 dict。
那一節說「你現在看得懂、知道長什麼樣子就好」，這一節要把它升級成「寫得出來」——加上剛學會
的條件，換一個你真的會遇到的情境。

\begin{code}{python}{dict comprehension 與 set comprehension（都可以帶條件）}
{key_expr: value_expr for item in iterable if condition}
{expression for item in iterable if condition}
\end{code}

dict comprehension 在真實程式碼裡最常見的用途只有一個，而且你在第 \ref{ch:csv} 章已經
吃過對應的苦頭：用 \texttt{DictReader} 讀進來的資料是「一列一列」的，每一列是一個 dict。
這種形狀很適合從頭走到尾，但只要你想問「Taylor 考幾分」，就得重跑一次迴圈去找。
一行 dict comprehension 可以把它翻成一張\textbf{查找表}。

\begin{example}{python}{(score-lookup.py) 把一列一列的資料變成一張查找表}
rows = [
    {"name": "Alice", "score": "88"},
    {"name": "Bob", "score": "76"},
    {"name": "Taylor", "score": "92"},
    {"name": "Charlie", "score": "65"},
    {"name": "David", "score": "79"},
]

score_of = {row["name"]: int(row["score"]) for row in rows}

print(score_of)
print(score_of["Taylor"])
print(score_of.get("Ivy", 0))
\end{example}

\iosample{}{\{'Alice': 88, 'Bob': 76, 'Taylor': 92, 'Charlie': 65, 'David': 79\} \\ 92 \\ 0 \\}

冒號左邊的 \texttt{row["name"]} 決定這一筆的 key，右邊的 \texttt{int(row["score"])} 決定
value——順便把字串轉成整數，這樣之後拿來計算就不用再轉一次。造好之後，
\texttt{score\_of["Taylor"]} 直接就是 92，不用再跑迴圈找人。最後一行用的是第
\ref{sec:dict-get} 節的 \texttt{.get(name, 0)}：名單裡沒有 Ivy 這個人，但這樣寫不會炸，
會拿到你指定的預設值 0。

順帶提醒一件第 \ref{sec:comprehension} 節講過的事：如果 \texttt{rows} 裡有兩筆同名的
資料，後面那筆會安靜地蓋掉前面那筆，不會報錯。查找表造出來的筆數比原始資料少，通常就是
這個原因。

set comprehension 則適合「去重複順便整理」。第 \ref{sec:set-ops} 節那個
\texttt{set(signups)} 只能原封不動去重複，如果資料本身大小寫不一致，就得先跑一趟迴圈
整理好、再丟進 \texttt{set()}。用 comprehension 一步到位：

\begin{example}{python}{(clean-tags.py) 標籤去重複，順手統一成小寫}
tags = ["Python", "csv", "CSV", "python", "JSON", "json", "csv"]

unique_tags = {tag.lower() for tag in tags}

print(len(tags), "->", len(unique_tags))
print(sorted(unique_tags))
\end{example}

\iosample{}{7 -> 3 \\{} ['csv', 'json', 'python'] \\}

七個標籤變成三個。這裡最後印的是 \texttt{sorted(unique\_tags)} 而不是 set 本身，因為第
\ref{ch:dict-set} 章說過 set 沒有順序，直接印出來每次執行的排列可能都不一樣，排過再印
才看得清楚內容。

\subsection*{常見錯誤：忘了冒號，拿到的是 set}

這一節唯一的坑，也是整章最陰險的一個——因為它不會給你任何錯誤訊息，只會給你錯的東西。

\begin{example}{python}{(dict-or-set.py) 有冒號跟沒冒號，差在哪裡}
rows = [
    {"name": "Alice", "score": "88"},
    {"name": "Bob", "score": "76"},
]

no_colon = {row["name"] for row in rows}
with_colon = {row["name"]: int(row["score"]) for row in rows}

print(type(no_colon))
print(type(with_colon))
print(sorted(no_colon))
print(with_colon["Alice"])
\end{example}

\iosample{}{<class 'set'> \\ <class 'dict'> \\{} ['Alice', 'Bob'] \\ 88 \\}

\texttt{\{row["name"] for row in rows\}} 一樣用大括號、一樣能跑、一樣印得出東西，但
\texttt{type()} 老實告訴你那是一個 \texttt{set}——一堆名字，不是查找表。它要到你寫
\texttt{no\_colon["Alice"]} 的時候才會爆炸，而那可能已經是好幾十行以後的事了。
\textbf{大括號裡有沒有那個冒號，決定你拿到的是兩種完全不同的東西}，寫完順手用
\texttt{type()} 確認一次，比事後回頭找便宜得多。

\begin{drills}[label=drills:comp-dict-set]
\item 預測輸出：\texttt{words = ["fig", "kiwi", "fig"]}，那麼
\texttt{print(\{w: len(w) for w in words\})} 會印出什麼？（注意有幾筆。）
\item 判斷題：下面這一行執行完，\texttt{result} 是一個 dict 還是一個 set？為什麼？
\begin{verbatim}
result = {row["name"] for row in rows}
\end{verbatim}
\item 造表題：有
\texttt{rows = [\{"name": "Amy", "score": "91"\}, \{"name": "Ben", "score": "58"\}]}，
寫一行 dict comprehension 造出「名字對應到整數分數」的查找表，並只收下 70 分以上的人。
\end{drills}

\section{換成小括號：你會拿到別的東西}
\label{sec:genexp-glimpse}

方括號給你 list，大括號給你 dict 或 set。那小括號呢？照這個規律推下去，答案應該是
tuple——第 \ref{ch:tuples} 章那種「定死了不能改」的容器。合理的猜測，可惜猜錯了。

\begin{example}{python}{(genexp-glimpse.py) 同樣的內容，換成小括號}
numbers = [1, 2, 3, 4]

squares_list = [n ** 2 for n in numbers]
squares = (n ** 2 for n in numbers)

print(squares_list)
print(squares)

for value in squares:
    print(value)

print(list(n ** 2 for n in numbers))
print(sum(n ** 2 for n in numbers))
\end{example}

\iosample{}{[1, 4, 9, 16] \\ <generator object <genexpr> at 0x104e769b0> \\{} 1 \\ 4 \\ 9 \\ 16 \\{} [1, 4, 9, 16] \\ 30 \\}

（那串十六進位數字是它在記憶體裡的位置，每次執行都不一樣，不用管它。）

方括號版本印出來就是你要的清單；小括號版本印出來卻是一個看不出內容的東西，叫做
\textbf{generator}（產生器）。這個畫面你在第 \ref{sec:map-filter} 節見過類似的——那時候
\texttt{print(map(int, prices))} 印出來的是 \texttt{<map object at ...>}，同樣是一個
「看不到內容」的東西。
\index{generator（產生器）}

關於它，這一章只講三件事：

\begin{enumerate}
\item \textbf{Python 沒有 tuple comprehension。} 小括號的位置被 generator 佔走了。真的
想要一個 tuple，就用 \texttt{tuple(...)} 把整串包起來。
\item \textbf{它不會馬上把結果算好給你。} 它先不動手，等你真的要用的時候才一個一個吐出
來。所以你可以把它直接丟進 \texttt{for} 走一次（上面那四行 1、4、9、16 就是這樣印出來
的），也可以用 \texttt{list(...)} 把它整個攤開成清單——跟 \texttt{map} 要用
\texttt{list()} 包起來才看得到內容，是同一個道理。
\item \textbf{當它是某個函式的唯一引數時，外層那對小括號可以省略。} 上面最後一行
\texttt{sum(n ** 2 for n in numbers)} 就是這樣寫的——本來應該是
\texttt{sum((n ** 2 for n in numbers))}，兩對括號疊在一起太醜，Python 允許你少寫一對。
\end{enumerate}

第三件事看起來只是省括號的小事，卻是下一節的入場券：它讓你可以寫出
\texttt{sum(...\ for ...\ in ...)} 這種「不先造出清單，直接要答案」的形狀。

至於 generator 為什麼要設計成「先不動手」、這樣做換到了什麼好處、用它的時候要注意什麼，
第 \ref{ch:iterators-generators} 章會完整回答。這一節你只要認得出來就好：看到小括號版本，
知道「這不是 tuple，也不是一份算好的清單」。

\section{一整串值，換一個答案：\texttt{sum}／\texttt{any}／\texttt{all}}
\label{sec:reducing-builtins}

前面幾節做的都是同一件事：造一個新的東西出來。這一節問反過來的問題——如果你要的根本不是
容器，是一個答案呢？

Python 有一組內建函式專門做這件事：吃進一整串東西，吐出一個答案。你手上已經有一個了。

\begin{code}{python}{吃一整串，吐一個答案}
sum(iterable)
any(iterable)
all(iterable)
\end{code}

\texttt{sum} 你在第 \ref{sec:list-tools} 節見過，當時跟 \texttt{min()}、\texttt{max()}
擺在一起，一句話帶過「拿總和」。現在要幫它升級：它不只能吃一份現成的清單，還可以直接吃
上一節那個「還沒算好的東西」。這一升級的威力，拿第 \ref{sec:csv-manual-average} 節那支
\texttt{average-quiz-score.py} 來比就很清楚——那裡為了算出總分，先用一個迴圈把分數收成
清單，再用一個累加器迴圈把它們加起來，前後六行。

\begin{example}{python}{(total-with-sum.py) 把兩個迴圈換成一行}
rows = [
    {"name": "Alice", "score": "88"},
    {"name": "Bob", "score": "76"},
    {"name": "Taylor", "score": "92"},
    {"name": "Charlie", "score": "65"},
    {"name": "David", "score": "79"},
]

total = sum(int(row["score"]) for row in rows)
average = total / len(rows)

print("Total:", total)
print("Average score:", average)
\end{example}

\iosample{}{Total: 400 \\{} Average score: 80.0 \\}

輸出跟第 \ref{sec:csv-manual-average} 節那支程式一模一樣，只是中間那六行變成一行。
\texttt{int(row["score"]) for row in rows} 是上一節那個「先不動手」的東西，
\texttt{sum} 收下它、一項一項跟它要值、邊要邊加，全程沒有任何一份中間清單被造出來。

不過第 \ref{sec:csv-manual-average} 節那個「\texttt{total = 0} $\to$ 迴圈裡 \texttt{+=} $\to$
最後除以數量」的樣板不是白學的：\texttt{sum} 只會做加法，一旦你要邊走邊做別的判斷、
或是同時累積好幾個東西，那個手寫樣板還是你的主力。

另外兩個是新面孔，它們回答的不是「總共多少」，而是是非題。\texttt{any} 問「\textbf{有沒有
任何一個}成立」，\texttt{all} 問「\textbf{是不是每一個都}成立」。這兩個問題你在第
\ref{ch:csv} 章一定遇過：讀進來的資料到底有沒有缺值？是不是每一筆都真的是數字？

\begin{example}{python}{(check-data.py) 這批資料有沒有問題}
rows = [
    {"name": "Alice", "score": "88"},
    {"name": "Bob", "score": ""},
    {"name": "Taylor", "score": "92"},
]

# 傳統寫法：有沒有哪一筆缺分數？
has_blank = False
for row in rows:
    if row["score"] == "":
        has_blank = True
        break
print(has_blank)

# 一行版
print(any(row["score"] == "" for row in rows))
print(all(row["score"].isdigit() for row in rows))

print(all([]))
print(any([]))
\end{example}

\iosample{}{True \\ True \\ False \\ True \\ False \\}

前面那五行迴圈跟後面那一行 \texttt{any(...)} 給的答案一樣，都是 \texttt{True}——Bob 那筆
的分數是空字串。\texttt{all(row["score"].isdigit() for row in rows)} 是
\texttt{False}，因為並不是每一筆都是純數字。順帶一提，\texttt{any} 跟迴圈版一樣聰明：
找到第一個成立的就停手，不會把整串走完。

\texttt{any} 跟 \texttt{all} 判斷的是第 \ref{ch:conditionals} 章那套 truthy／falsy 規則
——不是只有 \texttt{True}／\texttt{False} 才算數，空字串、0、空清單在它們眼裡也都是
「不成立」。所以 \texttt{any(row["score"] for row in rows)} 其實也能問出「有沒有哪一筆
是有東西的」，不用另外寫比較。

最後兩行是這一節的重頭戲，也是最容易答錯的一題。

\keypoint{對一個\textbf{空的}東西，\texttt{all} 回 \texttt{True}，\texttt{any} 回
\texttt{False}。「一個都沒有，怎麼會是『全部都符合』？」——關鍵在於 \texttt{all} 真正問的
是「\textbf{有沒有任何一個不符合}」。裡面一個東西都沒有，當然找不出不符合的，所以答
\texttt{True}。\texttt{any} 反過來，問的是「有沒有找到任何一個符合的」，空的東西裡什麼
都找不到，所以答 \texttt{False}。這個邊界在處理資料時很容易踩到：檔案剛好是空的、篩選
完剛好一筆都不剩，\texttt{all(...)} 就會回一個看起來「全部都通過檢查」的 \texttt{True}。}

\begin{drills}[label=drills:reducing-builtins]
\item 預測輸出：\texttt{nums = [2, 4, 6]}，下面四行各印出什麼？
\begin{verbatim}
print(sum(n for n in nums if n > 3))
print(all(n % 2 == 0 for n in nums))
print(any(n > 10 for n in nums))
print(all(n > 100 for n in []))
\end{verbatim}
\item 改寫題：把下面這段改寫成一行。
\begin{verbatim}
total = 0
for row in rows:
    total += int(row["score"])
\end{verbatim}
\end{drills}

\section{換一種走訪方式：\texttt{zip}}
\label{sec:zip-walk}

上一節那組工具，是把一整串壓成一個答案。這一節是另一類工具，它們\textbf{不改變你要做的
事，只改變「你走訪的那一串東西」長什麼樣子}。

這個定位一講出來，有一個工具就自動歸位了——\texttt{enumerate()}。第 \ref{sec:enumerate}
節教過它：你要位置又要內容的時候，不要自己造一個計數器手動累加，讓 \texttt{enumerate()}
一次給你兩個。它做的事正是「改造你走訪的那一串」：本來一項一項給你內容，包上
\texttt{enumerate()} 之後，變成一組一組給你「位置跟內容」。

那時候你只在 \texttt{for} 迴圈裡用過它。既然 comprehension 的第二格就是一個
\texttt{for}，它當然也能放進去，而且很好用——例如造一份帶編號的清單：

\begin{example}{python}{(numbered-list.py) 用 enumerate 造一份帶編號的清單}
names = ["Alice", "Bob", "Taylor"]

lines = [f"{i + 1}. {name}" for i, name in enumerate(names)]

for line in lines:
    print(line)
\end{example}

\iosample{}{1. Alice \\ 2. Bob \\ 3. Taylor \\}

\texttt{for i, name in enumerate(names)} 這一段跟你在 \texttt{for} 迴圈裡寫的完全一樣，
一次接住兩個名字；差別只在它現在被塞進了方括號裡，而最前面那一格用 f-string 把編號跟
名字組成一行字。編號寫 \texttt{i + 1} 是因為 \texttt{enumerate()} 從 0 開始數，這也是
第 \ref{sec:enumerate} 節的老招。

\subsection*{兩份資料要一起走：\texttt{zip}}

\texttt{enumerate} 解決的是「內容之外還要位置」。另一個很常見的需求是：你手上有\textbf{兩
份平行的資料}，第一份的第一項對應第二份的第一項，以此類推。從檔案裡分欄讀資料很容易造出
這種形狀——名字收成一份、分數收成一份。

\begin{code}{python}{用 zip 一次走兩份資料}
zip(iterable_a, iterable_b)
\end{code}
\index{zip()}

\begin{example}{python}{(walk-two-lists.py) 同時走訪名字與分數}
names = ["Alice", "Bob", "Taylor", "Charlie", "David"]
scores = [88, 76, 92, 65, 79]

for name, score in zip(names, scores):
    print(name, score)

print(zip(names, scores))
\end{example}

\iosample{}{Alice 88 \\ Bob 76 \\ Taylor 92 \\ Charlie 65 \\ David 79 \\ <zip object at 0x1051a27c0> \\}

\texttt{zip(names, scores)} 把兩份清單按位置配成一組一組，\texttt{for name, score in ...}
一次接住兩個名字——跟 \texttt{enumerate} 完全同一個手法，只是這次的第一個值不是位置編號，
而是另一份資料裡對應的那一項。

沒有 \texttt{zip} 的話，你只能寫 \texttt{for i in range(len(names))}，再用
\texttt{names[i]} 跟 \texttt{scores[i]} 分別去兩份清單裡取值。那樣寫能動，但你得自己管
索引，而且一不小心就會寫錯一份的名字。

最後那一行印出 \texttt{<zip object at ...>}，看不到內容——又是同一類的東西。跟上一節那個
generator 一樣，它也不會馬上把結果算好給你，想看內容就用 \texttt{list()} 包起來。

\texttt{zip} 最順手的一個用法，是一行造出查找表：

\begin{example}{python}{(zip-lookup.py) 一行造查找表，順便看看長度不一樣會怎樣}
names = ["Alice", "Bob", "Taylor"]
scores = [88, 76, 92]

score_of = dict(zip(names, scores))
print(score_of)
print(score_of["Bob"])

short_names = ["Alice", "Bob"]
pairs = [f"{name}={score}" for name, score in zip(short_names, scores)]
print(pairs)
\end{example}

\iosample{}{\{'Alice': 88, 'Bob': 76, 'Taylor': 92\} \\ 76 \\{} ['Alice=88', 'Bob=76'] \\}

前三行跟第 \ref{sec:comp-dict-set} 節那個 \texttt{score-lookup.py} 得到的是同一種東西
——一張用名字查分數的表——但走的是完全不同的路：那邊的原始資料是「一列一列的 dict」，
所以用 dict comprehension；這邊的原始資料是「兩份平行的清單」，所以用
\texttt{dict(zip(...))}。\textbf{資料原本長什麼形狀，決定你該走哪一條路。}

後面那三行藏著一個坑。\texttt{short\_names} 只有兩項，\texttt{scores} 有三項，
\texttt{zip} 走到\textbf{短的那一份用完就停}，Taylor 的 92 分就這樣安靜地消失了——沒有
報錯，沒有警告，只有結果少一筆。這種問題比報錯難抓得多，因為程式看起來完全正常。兩份
資料如果理論上該一樣長，配對之前先確認一次
（\texttt{len(names) == len(scores)}）不算多此一舉。

\stylepoint{要位置就用 \texttt{enumerate}，要平行走兩份就用 \texttt{zip}，這兩種情況都
不要寫 \texttt{for i in range(len(items))} 再自己去索引取值。第 \ref{sec:enumerate} 節
對 \texttt{enumerate} 已經講過半句，這裡把 \texttt{zip} 補上：\texttt{range(len(...))}
在 Python 裡幾乎永遠有更好的替代寫法，看到它就先想想是不是這兩個之一。}

\begin{drills}[label=drills:zip-walk]
\item 有 \texttt{names = ["a", "b", "c"]} 跟 \texttt{scores = [1, 2, 3, 4, 5]}。下面這個
迴圈會印出幾行？沒被印到的那幾項會怎麼樣？
\begin{verbatim}
for n, s in zip(names, scores):
    print(n, s)
\end{verbatim}
\item 預測輸出：\texttt{print(dict(zip(["a", "b"], [1, 2, 3])))} 會印出什麼？
\end{drills}

\section{什麼時候該用、什麼時候不要用}
\label{sec:comp-limits}

前面七節都在教你怎麼寫。這一節不教任何新語法，只做一件事：幫你建立「該不該用」的判斷力。
前半支持，後半反對，兩半合起來才完整。

\subsection*{前半：為什麼 comprehension 通常比較好讀}

第 \ref{sec:map-filter} 節結尾留了一個伏筆：那裡說「Python 圈子普遍覺得 comprehension
比較好讀」，但只給了立場，沒給理由。現在把理由補上，一共三條。三條都不需要你重新學
\texttt{map}／\texttt{filter} 怎麼用，只要把兩段程式碼並排看就好。

\textbf{第一條，意圖明確。} \texttt{for} 迴圈是萬用工具：它可能在數東西、可能在算總和、
可能在印東西、可能在改一份既有的清單。你得把整個迴圈讀完，才知道它到底想做什麼。
comprehension 只做一件事——\textbf{造一個新的容器出來}。看到方括號，你在讀懂內容之前就
已經知道結果是一份清單了；看到大括號有冒號，就知道是一張表。這正是這一章開場那個對比：
「說你要什麼」跟「一步一步交代怎麼做」的差別。

\textbf{第二條，讀的方向。} 兩段做同一件事的程式碼：

\begin{code}{python}{同一件事，兩種讀法}
numbers = list(map(int, filter(lambda s: s.isdigit(), raw)))

numbers = [int(s) for s in raw if s.isdigit()]
\end{code}

上面那行要由內往外讀：先找到最裡面的 \texttt{filter}（先過濾），再看外面的
\texttt{map}（再轉換），再看最外面的 \texttt{list}（最後收成清單）。眼睛得先跳到中間、
再往兩邊擴散。下面那行由左往右讀一遍就完了。這裡不多解釋，兩行並排放著，你自己看得出來
差別。

\textbf{第三條，少一個東西要學。} 注意下面那行裡 \texttt{lambda} 整個消失了。條件直接
寫成 \texttt{s.isdigit()}，不需要先包成一個匿名函式再交出去。

\stylepoint{這三條理由不代表 \texttt{map}／\texttt{filter} 是錯的或過時的。它們是好好的
內建函式，你在別人的程式碼裡還是會看到，第 \ref{sec:map-filter} 節的定位維持不變：
\textbf{看得懂就好}。這一節只是把「為什麼輪到自己寫的時候優先選 comprehension」這個理由
補齊，不是要你去改別人的程式碼。}

\subsection*{後半：三條煞車}

好用的東西都有被用過頭的風險，comprehension 尤其容易——因為「把好幾行縮成一行」這件事
本身很有成就感。以下三條，只要踩到任何一條，就該寫回 \texttt{for} 迴圈。

\textbf{第一條：超過兩層就拆開。} 第 \ref{sec:comp-nested} 節說過這一章只教到兩層，這裡
把它收成規則。一個更好記的門檻是：\textbf{一段 comprehension 如果寫出來超過兩行，就該
考慮改回迴圈}。一行讀完是它最大的優點，一旦要折行、要往右捲，這個優點就沒了。

\textbf{第二條：你不要那份結果清單的時候，就不該用它。} comprehension 存在的理由是
「造一個新容器」。如果你只是想「做某件事」，那就寫迴圈。這個誤用很常見，而且它不會報錯：

\begin{example}{python}{(print-comprehension.py) 用 comprehension 來印東西}
names = ["Alice", "Bob", "Taylor"]

result = [print(name) for name in names]

print(result)
\end{example}

\iosample{}{Alice \\ Bob \\ Taylor \\{} [None, None, None] \\}

三個名字確實印出來了，看起來一切正常。但最後一行揭穿了它：那份 comprehension 老老實實
造了一份清單交給 \texttt{result}，裡面裝著三個 \texttt{None}——因為 \texttt{print()} 這個
函式沒有回傳值，呼叫它拿到的就是 \texttt{None}（第 \ref{ch:functions} 章講過）。你完全
不需要這份清單，它卻被造了出來、佔著記憶體，而且沒有任何人會去看它。

判準只有一句話：\textbf{你不要那個結果清單的時候，就不該用這個語法}。這裡該寫的是
\texttt{for name in names: print(name)}，兩行，而且誠實。

\textbf{第三條：comprehension 裡放不進 \texttt{try}／\texttt{except}。} 這一條會讓你有點
不舒服，因為它讓你剛學會的兩個好習慣正面衝突。第 \ref{sec:lbyl-vs-eafp} 節教過：可能失敗
的操作，Python 的作風是先做、錯了再用 \texttt{try}／\texttt{except} 接住。而從檔案讀進來
的東西一律是字串，\texttt{int()} 隨時可能踩到空字串或髒資料。

問題是，comprehension 裡\textbf{沒有地方能放 \texttt{try}}。它整串是一個式子，而
\texttt{try}／\texttt{except} 是敘述——這正是第 \ref{sec:conditional-expression} 節那條
界線的又一次應用。所以下面這行遇到髒資料時，你只能眼睜睜看它整個倒下：

\begin{code}{python}{沒有地方能接住錯誤的一行}
scores = {row["name"]: int(row["score"]) for row in raw_rows}
\end{code}

\begin{verbatim}
Traceback (most recent call last):
  File "dirty-comprehension.py", line 6, in <module>
    scores = {row["name"]: int(row["score"]) for row in raw_rows}
                           ~~~^^^^^^^^^^^^^^
ValueError: invalid literal for int() with base 10: ''
\end{verbatim}

整串倒下，前面已經處理好的資料也一起沒了。這種時候答案很乾淨：\textbf{寫回 \texttt{for}
迴圈}。

\begin{example}{python}{(clean-dirty-scores.py) 需要接住錯誤，就寫回迴圈}
raw_rows = [
    {"name": "Alice", "score": "88"},
    {"name": "Bob", "score": ""},
    {"name": "Taylor", "score": "92分"},
    {"name": "David", "score": "79"},
]

scores = {}
for row in raw_rows:
    try:
        scores[row["name"]] = int(row["score"])
    except ValueError:
        print("skipped", row["name"])

print(scores)
\end{example}

\iosample{}{skipped Bob \\ skipped Taylor \\ \{'Alice': 88, 'David': 79\} \\}

四筆資料，兩筆有問題，程式沒有倒，好的兩筆照樣收進表裡，壞的兩筆還留下了紀錄。這段程式
比一行版長了很多，但它是對的，而一行版不是。

\stylepoint{三條煞車一起記：\textbf{超過兩層（或超過兩行）就拆開}；\textbf{只是想做某件
事、不要那份結果清單的時候不要用}；\textbf{需要接住錯誤的時候不要用}。這三條的共同精神
跟第 \ref{sec:lbyl-vs-eafp} 節那條提醒是一樣的——一個寫法再怎麼慣用，也不是每個地方都該
用。判斷力比語法值錢。}

\begin{drills}[label=drills:comp-limits]
\item 判斷題：有人寫了 \texttt{[send\_email(user) for user in users]} 來寄一輪通知信，
\texttt{send\_email()} 只負責寄信、沒有回傳值。這樣寫該不該？為什麼？
\item 判斷題：一份從 CSV 讀進來的資料，分數欄位偶爾會是空字串、或夾著文字（例如
\texttt{"92分"}），你要把它們轉成整數收成一份清單。這個工作該不該寫成 comprehension？
為什麼？
\end{drills}

\section*{練習解答}
\addcontentsline{toc}{section}{練習解答}

\paragraph{第~\ref{sec:comp-filter}~節}
\begin{enumerate}
\item 印出 \texttt{[8, 10, 6]}：只有大於 5 的三項通過關卡，順序維持原本的順序。
\item \texttt{long\_words = [w.upper() for w in words if len(w) > 4]}，結果是
\texttt{['APPLE', 'BANANA']}。最前面那格負責「變成什麼」（轉大寫），最後面那格負責
「留不留」。
\item \texttt{if} 是過濾用的，必須放在最後面：\texttt{[w for w in words if len(w) > 3]}。
放在前面 Python 會以為你要寫條件運算式，於是要求你補一個 \texttt{else}，丟出
\texttt{SyntaxError: expected 'else' after 'if' expression}。
\end{enumerate}

\paragraph{第~\ref{sec:conditional-expression}~節}
\begin{enumerate}
\item 第一行印出 \texttt{[9, 7]}，兩項；第二行印出 \texttt{[0, 9, 0, 7]}，四項。
\texttt{if} 在 \texttt{for} 後面是過濾，會變短；在 \texttt{for} 前面是選值，長度跟原本
一樣。
\item \texttt{labels = ["warm" if t >= 25 else "cool" for t in temps]}。
\end{enumerate}

\paragraph{第~\ref{sec:comp-nested}~節}
\begin{enumerate}
\item 翻譯成：\texttt{result = []} $\to$ \texttt{for word in words:} $\to$ 縮排一層
\texttt{for c in word:} $\to$ 再縮排一層 \texttt{result.append(c)}。兩個 \texttt{for} 的
順序跟 comprehension 裡完全一樣。結果是 \texttt{['a', 'b', 'c', 'd']}。
\item 印出 \texttt{[10, 20, 30, 40]}。先走完第一列 \texttt{[1, 2]} 才換到第二列，跟兩層
縮排迴圈的順序一致。
\end{enumerate}

\paragraph{第~\ref{sec:comp-dict-set}~節}
\begin{enumerate}
\item 印出 \texttt{\{'fig': 3, 'kiwi': 4\}}，只有兩筆——\texttt{"fig"} 出現兩次，但同一個
key 後面那次會蓋掉前面那次。
\item 是 set。大括號裡\textbf{沒有冒號}，Python 只看到「一個式子」，就造 set 給你；要
dict 必須寫成 \texttt{\{row["name"]: 某個值 for row in rows\}}。這個錯不會報錯，要用
\texttt{type()} 或後面存取出錯時才發現得了。
\item \texttt{\{row["name"]: int(row["score"]) for row in rows if int(row["score"]) >= 70\}}，
結果是 \texttt{\{'Amy': 91\}}。（Ben 只有 58 分，被關卡擋下來。）
\end{enumerate}

\paragraph{第~\ref{sec:reducing-builtins}~節}
\begin{enumerate}
\item 依序印出 \texttt{10}（只加 4 跟 6）、\texttt{True}（三個都是偶數）、
\texttt{False}（沒有任何一個大於 10）、\texttt{True}（空的東西丟給 \texttt{all} 一律回
\texttt{True}）。
\item \texttt{total = sum(int(row["score"]) for row in rows)}。
\end{enumerate}

\paragraph{第~\ref{sec:zip-walk}~節}
\begin{enumerate}
\item 跑 3 輪、印 3 行。\texttt{zip} 走到短的那一份用完就停，\texttt{scores} 多出來的
\texttt{4} 跟 \texttt{5} 會被安靜地丟掉，不會報錯。
\item 印出 \texttt{\{'a': 1, 'b': 2\}}。理由同上：只配到兩組，\texttt{3} 被丟掉了。
\end{enumerate}

\paragraph{第~\ref{sec:comp-limits}~節}
\begin{enumerate}
\item 不該。這裡的目的是「寄信」這件事本身，不是要一份結果清單，而
\texttt{send\_email()} 沒有回傳值，所以會白白造出一份裝滿 \texttt{None} 的清單。應該寫成
\texttt{for user in users: send\_email(user)}。
\item 不該。髒資料表示 \texttt{int()} 隨時可能丟出 \texttt{ValueError}，而 comprehension
裡沒有地方能放 \texttt{try}／\texttt{except}，一筆壞資料就會讓整串倒下、連前面處理好的
也一起沒了。要能接住錯誤、跳過壞資料繼續做，就得寫回 \texttt{for} 迴圈。
\end{enumerate}

\chapter{Iterator 與 Generator}
\label{ch:iterators-generators}

這本書欠你五個解釋。不是那種「以後有機會再說」的客套話，是白紙黑字寫在前面某一章裡、
說了「這個之後會回頭講」然後真的先跳過去的地方。這一章要一次還清。

先看三段你一定認得的程式碼。第一段是第 \ref{ch:loops} 章走訪字串的寫法：

\begin{code}{python}{第 \ref{ch:loops} 章：走訪一句話裡的每個字元}
for char in sentence:
    if char == " ":
        space_count += 1
\end{code}

第二段是第 \ref{ch:files-exceptions} 章逐行讀檔案：

\begin{code}{python}{第 \ref{ch:files-exceptions} 章：一行一行讀進來}
with open("names.txt") as names_file:
    for line in names_file:
        print(line.rstrip())
\end{code}

第三段是第 \ref{ch:csv} 章跳過標題列那一行：

\begin{code}{python}{第 \ref{ch:csv} 章：把標題列先讀掉}
next(reader)
\end{code}

這三段做的其實是同一件事，而這本書從頭到現在都沒告訴你那件事是什麼。第
\ref{ch:loops} 章第 \ref{sec:for-vs-while} 節那個學習要點是這樣說的：\texttt{range()}
和直接走訪序列「\textbf{都會替你處理好}『這一輪走到哪、還有沒有下一輪』」——只說了有人
替你處理，沒說那個人是誰。

這一章就是回答這個問題：\textbf{替你處理的那個「誰」，到底是什麼？} 然後把答案變成你
自己也做得出來的東西。

這一章分四段推進：

\begin{itemize}
\item 第一段掀開 \texttt{for} 的蓋子。先把一個 \texttt{for} 迴圈手動拆開來跑一次，再
把「一批東西」跟「真正一個一個交出來的那個東西」分成兩種角色
（第 \ref{sec:what-for-does}～\ref{sec:iterable-vs-iterator} 節）。
\item 第二段換你動手做一台。\texttt{yield} 這個關鍵字會把一個普通函式變成完全不同的
東西，而第 \ref{ch:comprehensions} 章那個換成小括號拿到的神祕物件，其實就是它的一行版
（第 \ref{sec:yield-basics}～\ref{sec:generator-expression} 節）。
\item 第三段是這一章的重頭戲，一好一壞剛好是同一件事的正反兩面：不先算好省下了什麼，
以及不先算好要付出什麼代價（第 \ref{sec:lazy-memory}～\ref{sec:one-shot} 節）。
\item 第四段收尾。有一個標準函式庫模組整個就是在做這件事，你在第 \ref{ch:modules} 章
已經跟它打過一次照面（第 \ref{sec:itertools-peek} 節）。
\end{itemize}

\section{\texttt{for} 底下到底發生什麼事}
\label{sec:what-for-does}

先講一個大家都排過的隊。去醫院、去郵局、去監理所，門口通常有一台機器，按一下吐一張
號碼牌。這台機器有幾個特性，等一下每一個都會對應到程式裡的某個東西：

\begin{itemize}
\item 今天要辦事的那一批人，跟門口那台機器，是\textbf{兩個不同的東西}。
\item 你按一下，它才吐一張；不按，它就靜靜站在那裡。
\item 它\textbf{不會}一大早就把今天所有號碼牌都印好堆在櫃台上。
\item 號碼發完了，它會明確告訴你「今天發完了」。
\item 撕下來的號碼牌不會自己飛回機器裡。想從頭再來一輪，得換一台新的機器。
\end{itemize}

Python 的 \texttt{for} 迴圈底下，就藏著一台這樣的機器。它有三個零件，而這三個零件用到
的語法你全部都學過了——真正新的只有三個名字。

\subsection*{把一個 \texttt{for} 迴圈拆開來跑一次}

下面這支程式跑了兩次同樣的事情。上半段是你平常會寫的；下半段一個 \texttt{for} 都沒有，
全部手工完成。

\begin{example}{python}{(for-by-hand.py) 同一件事，一次用 for，一次全手工}
names = ["Alice", "Bob", "Taylor"]

# 你平常寫的版本
for name in names:
    print(name)

print("---")

# 完全等價的手動版本
it = iter(names)
while True:
    try:
        name = next(it)
    except StopIteration:
        break
    print(name)
\end{example}

\iosample{}{Alice \\ Bob \\ Taylor \\ --- \\ Alice \\ Bob \\ Taylor \\}

兩段輸出一模一樣，一個字都不差。這就是這一節唯一要你相信的事：\textbf{你寫的
\texttt{for}，Python 在底下做的就是下半段那件事，不多也不少。}

下半段裡新出現的東西只有三個名字：

\begin{itemize}
\item \texttt{iter(names)}——跟那批東西要一台機器。
\index{iter()}
\item \texttt{next(it)}——按一下，拿走下一個。
\item \texttt{StopIteration}——機器說「發完了」。
\index{StopIteration}
\end{itemize}

其他零件你全都認識：\texttt{while True} 是第 \ref{ch:loops} 章第
\ref{sec:while-basics} 節的無窮迴圈，\texttt{break} 是第 \ref{sec:break-search} 節的
提早結束，\texttt{try}／\texttt{except} 是第 \ref{ch:files-exceptions} 章第
\ref{sec:lbyl-vs-eafp} 節的先做再處理。手動版之所以比較長，不是因為它比較難，是因為
\texttt{for} 把這五行濃縮成了一行。

那台從 \texttt{names} 要來的機器，有個正式名字叫 \textbf{iterator}——它就是「真正負責
一個一個把東西交出來」的那個東西。
\index{iterator}

\subsection*{一直按下去會怎樣}

機器發完了會怎麼樣？直接按給你看。這次換一個更小的素材：字串 \texttt{"cat"} 只有三個
字元，所以第四次一定會出事。

\begin{example}{python}{(press-until-empty.py) 連按四次，第四次沒東西了}
it = iter("cat")

print(next(it))
print(next(it))
print(next(it))
print(next(it))
\end{example}

\iosample{}{c \\ a \\ t \\}

前三次乖乖吐出 \texttt{'c'}、\texttt{'a'}、\texttt{'t'}，第四次程式停了下來，畫面上出現：

\begin{verbatim}
Traceback (most recent call last):
  File "press-until-empty.py", line 6, in <module>
    print(next(it))
          ~~~~^^^^
StopIteration
\end{verbatim}

看到 \texttt{Traceback} 先不要緊張。\textbf{這一次不是你打錯字}，程式碼一個字都沒錯，
這就是機器在說「今天的號碼發完了」。\texttt{StopIteration} 直譯就是「停止走訪」，它是
Python 用來表達「沒有下一個了」的方式。

那為什麼你平常寫 \texttt{for} 從來沒看過它？因為 \texttt{for} 每一輪都在替你接住它，
接到了就默默結束迴圈——就是上一支程式那個 \texttt{except StopIteration: break}。

\keypoint{一台 iterator 沒有「還有沒有下一個」可以問。Python \textbf{刻意不提供}那種
先問一句再決定要不要拿的做法，它只給你一個動作：按下去。按到沒有了，它就丟
\texttt{StopIteration} 告訴你。這正是第 \ref{ch:files-exceptions} 章第
\ref{sec:lbyl-vs-eafp} 節那個 EAFP——先做，錯了再處理。那一節教的是「你寫程式時該用這
種風格」，這裡能給你一句更強的：\textbf{連 Python 自己的 \texttt{for} 迴圈都是這樣做的}，
它不先檢查，它先按，然後準備好接住那個「沒了」。}

\subsection*{還一張欠條：\texttt{next(reader)} 到底是什麼}

第 \ref{ch:csv} 章第 \ref{sec:csv-module-read} 節有一行程式，當時被稱作「一個小技巧」：
\texttt{next(reader)} 放在 \texttt{for} 迴圈前面，就能把標題列先讀掉。你照著打了，但沒人
告訴你它是什麼。

現在可以說了：\texttt{reader} 本身就是一台這樣的機器，\texttt{next(reader)} 就是按它一下。
第一張號碼牌（標題列）被你撕走了，剩下的才輪到 \texttt{for}。

\begin{example}{python}{(skip-header.py) 現在看得懂這一行了}
import csv

with open("quiz_scores.csv") as f:
    reader = csv.reader(f)
    header = next(reader)
    print("撕掉的第一張：", header)
    for row in reader:
        print(row)
\end{example}

\iosample{}{撕掉的第一張： ['name', 'score'] \\{} ['Alice', '88'] \\{} ['Bob', '76'] \\{} ['Taylor', '92'] \\{} ['Charlie', '65'] \\{} ['David', '79'] \\}

用的還是第 \ref{ch:csv} 章那份 \texttt{quiz\_scores.csv}。這裡多做了一件事：把撕下來的
那一張接住印出來，證明它真的沒有消失，只是不在 \texttt{for} 拿得到的地方了。

最後把這一節的定位講清楚：\textbf{手動版你要看得懂，但不必自己寫}。日常寫 Python 沒有
人會去手動 \texttt{iter} 加 \texttt{while True}，因為 \texttt{for} 已經幫你寫好了。手動版
存在的唯一理由，是讓你看清楚那台機器長什麼形狀——看過一次之後，以後在別人的程式碼裡撞見
\texttt{next(...)}，你就知道那是在按什麼。

\begin{drills}[label=drills:what-for-does]
\item 翻譯題：把下面這個迴圈改寫成手動版（用 \texttt{iter}、\texttt{next}、
\texttt{while True} 與 \texttt{try}／\texttt{except StopIteration}）。
\begin{verbatim}
for score in [88, 76, 92]:
    print(score * 2)
\end{verbatim}
\item 預測輸出：下面這段程式印出什麼？跑到哪一行會停下來、為什麼？
\begin{verbatim}
it = iter([10, 20])
print(next(it))
print(next(it))
print(next(it))
\end{verbatim}
\end{drills}

\section{兩種角色：一批東西，和一台發號機}
\label{sec:iterable-vs-iterator}

第 \ref{ch:lists} 章第 \ref{sec:list-comprehension} 節給過一句話定義：「\texttt{iterable}
就是你可以放進 \texttt{for} 迴圈裡一個一個處理的東西，像 list、字串、\texttt{range()}
都是。」第 \ref{ch:modules} 章也用過「可疊代物件」這個中文說法。
\index{iterable（可疊代物件）}

那句定義沒有錯，但它把兩個東西講成了一個。上一節你已經看到它們同時出現：\texttt{names}
是一個，\texttt{iter(names)} 拿到的是另一個。這一節把它們正式分開。

\begin{code}{python}{兩種角色，兩個內建函式}
iterator = iter(iterable)
item = next(iterator)
\end{code}

\textbf{iterable（可疊代物件）是那一批東西}——list、字串、\texttt{range()}、dict、set，
你能放進 \texttt{for} 的都算。\textbf{iterator 是那台機器}，你只能對它做一件事：按一下。
兩者的關係就是 \texttt{iter()} 這個函式：拿一批東西去要一台機器。

\subsection*{對一份 list 直接按下去}

分不清楚的話，Python 會直接告訴你。

\begin{example}{python}{(next-on-a-list.py) 把一份 list 當成機器來按}
names = ["Alice", "Bob"]

print(next(names))
\end{example}

\begin{verbatim}
Traceback (most recent call last):
  File "next-on-a-list.py", line 3, in <module>
    print(next(names))
          ~~~~^^^^^^^
TypeError: 'list' object is not an iterator
\end{verbatim}

這句錯誤訊息本身就是這一節最好的教材：\texttt{'list' object is not an iterator}——
「list 不是一台 iterator」。它是一批東西，不是負責發號的那台機器。要按之前，得先用
\texttt{iter()} 跟它要一台。

同一段話對 \texttt{range} 也成立：\texttt{next(range(5))} 一樣會得到
\texttt{TypeError: 'range' object is not an iterator}。\texttt{range} 從第
\ref{ch:loops} 章用到現在，你大概一直把它當成「一串數字」，這裡先記住它是「一批東西」
這個身分，第 \ref{sec:lazy-memory} 節還會回頭找它算一次帳。

\begin{example}{python}{(list-vs-iterator.py) 同一份資料連走兩趟，兩種角色差在哪}
names = ["Alice", "Bob"]

# 一批東西：想走幾次就走幾次
for name in names:
    print("第一趟", name)
for name in names:
    print("第二趟", name)

print("---")

# 一台發號機：只能用完一次
it = iter(names)
for name in it:
    print("第一趟", name)
for name in it:
    print("第二趟", name)

print("結束")
\end{example}

\iosample{}{第一趟 Alice \\ 第一趟 Bob \\ 第二趟 Alice \\ 第二趟 Bob \\ --- \\ 第一趟 Alice \\ 第一趟 Bob \\ 結束 \\}

上半段兩趟都有輸出，下半段第二趟一行都沒有——中間那兩行「第二趟」直接不見了，程式卻
一聲不吭地跑完了。

\keypoint{一批東西可以被走很多次，一台發號機只能用完一次。同一份 list 你想 \texttt{for}
幾次就幾次，是因為\textbf{每次 \texttt{for} 開跑，都會先替你跟它要一台全新的機器}——
第一趟用掉的那台，跟第二趟用的不是同一台。你自己用 \texttt{iter()} 要來的那台則不會自動
換新，用完就是用完了。這一條是第 \ref{sec:one-shot} 節那個坑的全部原因，這裡先講原理，
那一節再演給你看。}

\subsection*{你早就撞見過三台機器了}

把這個分辨學會之後，前面幾章有三個「當時只能照著打」的地方會自動歸位。

\begin{example}{python}{(three-machines.py) 檔案物件、csv.reader、map 的結果}
import csv

# 一、檔案物件本身就是一台發號機
with open("quiz_scores.csv") as f:
    print(next(f).rstrip())
    print(next(f).rstrip())

# 二、csv.reader 的結果也是
with open("quiz_scores.csv") as f:
    reader = csv.reader(f)
    print(next(reader))

# 三、map 的結果也是
numbers = map(int, ["12", "5", "38"])
print(next(numbers))
print(list(numbers))
\end{example}

\iosample{}{name,score \\ Alice,88 \\{} ['name', 'score'] \\ 12 \\{} [5, 38] \\}

三件事各講一句：

\textbf{檔案物件}。第 \ref{ch:files-exceptions} 章第 \ref{sec:reading-lines} 節說「檔案
物件本身就是可以放進 \texttt{for} 迴圈裡的東西」，機制就在這裡——它本身就是一台機器，
所以你可以直接按它。

\textbf{\texttt{csv.reader} 的結果}。所以上一節那個 \texttt{next(reader)} 才拿得走一列。

\textbf{\texttt{map} 的結果}。第 \ref{ch:advanced-functions} 章第 \ref{sec:map-filter}
節那個 \texttt{<map object at ...>}，就是一台你還沒去按的機器。那時候書上叫你用
\texttt{list()} 包起來才看得到內容，並且明說「這種『先不急著算』的設計……完整的討論留到
第 \ref{ch:iterators-generators} 章」——現在兌現：\texttt{list()} 做的事就是\textbf{一口氣
把它按到底，把吐出來的東西全部收好}。上面最後兩行正是這件事拆成兩半：先手動按一下拿到
\texttt{12}，剩下的才被 \texttt{list()} 掃光，所以拿到的是 \texttt{[5, 38]} 而不是三項。
（\texttt{filter} 的結果、第 \ref{sec:zip-walk} 節那個 \texttt{<zip object at ...>}，
全都是同一類東西。）

\begin{drills}[label=drills:iterable-vs-iterator]
\item 判斷題：下面兩行哪一行會報錯？報什麼錯？另一行印出什麼？
\begin{verbatim}
print(next("cat"))
print(next(iter("cat")))
\end{verbatim}
\item 預測輸出：下面這段印出幾行？分別是什麼？
\begin{verbatim}
words = ["fig", "kiwi"]
it = iter(words)
for w in it:
    print("A", w)
for w in it:
    print("B", w)
\end{verbatim}
\item 歸類題：下面五樣東西，哪些是「一批東西」（iterable）、哪些是「一台發號機」
（iterator）？\texttt{[1, 2, 3]}、\texttt{"hello"}、用 \texttt{open()} 拿到的檔案物件、
\texttt{csv.reader(f)} 的結果、\texttt{range(10)}。
\end{drills}

\section{自己做一台：\texttt{yield}}
\label{sec:yield-basics}

前兩節都在認識別人做好的機器。這一節換你做一台。

做法出乎意料地簡單：\textbf{在一個函式裡寫一個 \texttt{yield}}。只要有這個字，這個函式
就不再是「算完給你答案」的函式了——你呼叫它，拿到的不是答案，是一台機器。這種函式叫
\textbf{generator function}，它產出的那台機器叫 \textbf{generator}（產生器）。
\index{yield}
\index{generator function（generator 函式）}

\begin{code}{python}{generator function 的形狀}
def generator_function(source):
    for item in source:
        yield transformed_item
\end{code}

跟你熟悉的函式比，形狀幾乎沒變，只有最後一行的 \texttt{return} 換成了 \texttt{yield}。
但這兩個字的行為差很多：\textbf{\texttt{return} 一講完函式就結束了；\texttt{yield} 只是
先停在這裡}，把手上這個值遞出去，等你下次按，它從停住的地方繼續往下跑。

\subsection*{它到底什麼時候才真的執行}

這件事光用講的沒有說服力，直接在函式裡塞幾個 \texttt{print}，看它們什麼時候才出現。

\begin{example}{python}{(generator-steps.py) 在函式裡埋三個 print，看它什麼時候跑}
def steps():
    print("開始")
    yield "第一個值"
    print("中間")
    yield "第二個值"
    print("結束")


print("--- 呼叫函式之前 ---")
g = steps()
print("--- 呼叫函式之後 ---")
print(g)
print("--- 開始走訪 ---")
for value in g:
    print(value)
print("--- 走訪完了 ---")
\end{example}

\iosample{}{--- 呼叫函式之前 --- \\ --- 呼叫函式之後 --- \\ <generator object steps at 0x1056619c0> \\ --- 開始走訪 --- \\ 開始 \\ 第一個值 \\ 中間 \\ 第二個值 \\ 結束 \\ --- 走訪完了 --- \\}

（那串十六進位數字是它在記憶體裡的位置，你跑出來的會跟這裡不一樣，不用管它。）

這段輸出有三件事值得看：

\textbf{第一，「呼叫函式之前」跟「呼叫函式之後」中間什麼都沒有。} \texttt{g = steps()}
這一行明明呼叫了函式，函式裡第一行就是 \texttt{print("開始")}，畫面上卻連一個字都沒出現。
因為呼叫 generator function 根本不會執行函式的內容，它只是\textbf{做一台機器交給你}。

\textbf{第二，\texttt{print(g)} 印出 \texttt{<generator object steps at 0x...>}。} 這個
畫面眼熟嗎？第 \ref{ch:advanced-functions} 章那個 \texttt{<map object at 0x...>}、第
\ref{ch:comprehensions} 章那個 \texttt{<generator object <genexpr> at 0x...>}，一模一樣的
形狀，一模一樣的意思：一台還沒被按過的機器。

\textbf{第三，走訪的時候輸出是交錯的。} 「開始」印完就吐出第一個值，然後停住等你；你再
按，「中間」才印出來。函式裡的程式碼不是一口氣跑完的，它\textbf{停在每個 \texttt{yield}
那裡等你}。

\subsection*{寫一個真的有用的}

看懂機制之後，寫一個實際會用到的。第 \ref{ch:files-exceptions} 章你寫過「逐行讀檔案、
每行 \texttt{strip()} 掉多餘空白」的迴圈。把那個迴圈的 \texttt{print} 換成
\texttt{yield}，它就變成一台可以到處插進 \texttt{for} 裡的機器。

假設有一份報名名單 \texttt{signups.txt}，是別人隨手貼進來的，前後有空白、中間還夾了一行空的：

\begin{verbatim}
  Alice
Bob

   Taylor
\end{verbatim}

\begin{example}{python}{(clean-lines.py) 一台專門吐乾淨資料的機器}
def clean_lines(filename):
    with open(filename) as f:
        for line in f:
            cleaned = line.strip()
            if cleaned:
                yield cleaned


for name in clean_lines("signups.txt"):
    print(f"[{name}]")
\end{example}

\iosample{}{[Alice] \\{} [Bob] \\{} [Taylor] \\}

方括號是印出來確認前後真的沒有空白用的。這支程式的價值在於\textbf{它幾乎沒變}——函式
裡那三行迴圈跟你第 \ref{ch:files-exceptions} 章寫過的一模一樣，唯一的改動是把
\texttt{print(cleaned)} 換成 \texttt{yield cleaned}。而換完之後，「怎麼清資料」跟「清完
要拿去做什麼」就分成了兩件事：以後你想印、想存成 list、想丟給 \texttt{sum} 算總和，都不
必再動 \texttt{clean\_lines} 一根寒毛。

用詞上順帶講一件事，以後讀文件會用得到：普通函式 \textbf{return（回傳）}一個值，
generator 是 \textbf{yield（吐出）}值。這本書從這裡開始會一直照這個分別用詞。

\subsection*{常見錯誤：呼叫了，卻沒有走訪}

這是新手寫完第一個 generator function 之後最常撞的牆：程式跑完了，什麼事都沒發生，
\textbf{而且不會報錯}。

\begin{example}{python}{(forgot-to-loop.py) 呼叫了卻沒走訪，等於什麼都沒做}
def shout(names):
    for name in names:
        print("處理中：", name)
        yield name.upper()


names = ["Alice", "Bob"]

shout(names)
print("跑完了，可是上面什麼都沒印")

print(list(shout(names)))
\end{example}

\iosample{}{跑完了，可是上面什麼都沒印 \\ 處理中： Alice \\ 處理中： Bob \\{} ['ALICE', 'BOB'] \\}

\texttt{shout(names)} 那一行單獨寫出來，就只是造了一台機器，然後把它丟了——函式裡的
\texttt{print("處理中：", ...)} 一次都沒跑。要它真的動起來，得有人去按：放進 \texttt{for}、
用 \texttt{list()} 收起來、丟給 \texttt{sum}，都算。最後那行 \texttt{list(shout(names))}
一按到底，兩行「處理中」才姍姍來遲地出現。

還有一個相關的手誤：\textbf{\texttt{yield} 打成 \texttt{return}}。那樣的話函式就變回
普通函式了，跑到那一行吐出一個值就整個結束，不會有第二個值，也不會拿到任何機器。

\begin{drills}[label=drills:yield-basics]
\item 預測輸出：下面這段程式會印出哪幾行？順序是什麼？
\begin{verbatim}
def numbers():
    print("a")
    yield 1
    print("b")
    yield 2
    print("c")

for n in numbers():
    print(n)
\end{verbatim}
\item 改寫題：把下面這個函式改寫成 generator function（去掉那份空清單與
\texttt{append}，改用 \texttt{yield}）。改寫後，呼叫端要怎麼寫才拿得到
\texttt{[2, 4, 6]}？
\begin{verbatim}
def doubled(items):
    result = []
    for n in items:
        result.append(n * 2)
    return result
\end{verbatim}
\item 抓蟲：有人寫了下面這段，執行後畫面上一個字都沒有，也沒有任何錯誤訊息。為什麼？
怎麼修？
\begin{verbatim}
def greet(names):
    for name in names:
        yield f"Hello, {name}!"

greet(["Alice", "Bob"])
\end{verbatim}
\end{drills}

\section{接住小括號：generator expression}
\label{sec:generator-expression}

第 \ref{ch:comprehensions} 章第 \ref{sec:genexp-glimpse} 節結尾留了一句話：「至於
generator 為什麼要設計成『先不動手』、這樣做換到了什麼好處、用它的時候要注意什麼，第
\ref{ch:iterators-generators} 章會完整回答。」這一節開始認領。

先把最直接的一句話講掉：\textbf{你在那一節換成小括號拿到的東西，就是上一節那台機器}——
只是換成一行寫法，\texttt{yield} 被藏起來了。
\index{generator expression}

\subsection*{同一件事，兩種寫法}

不用抽象比較，直接把兩種寫法擺在同一支程式裡。

\begin{example}{python}{(two-ways.py) 有 yield 的四行版，跟一行版}
scores = [88, 76, 92, 65, 79]


def passed_scores(all_scores):
    for score in all_scores:
        if score >= 70:
            yield score


from_function = passed_scores(scores)
from_genexp = (score for score in scores if score >= 70)

print(from_function)
print(from_genexp)
print(list(from_function))
print(list(from_genexp))
\end{example}

\iosample{}{<generator object passed\_scores at 0x1036f69b0> \\ <generator object <genexpr> at 0x103781150> \\{} [88, 76, 92, 79] \\{} [88, 76, 92, 79] \\}

（兩串位置數字你跑出來都會不一樣。）

兩邊印出來都是 \texttt{generator object}，收成 list 之後結果一模一樣。差別只有一個：
上面那台機器有名字（\texttt{passed\_scores}，因為它是你定義的函式），下面那台叫
\texttt{<genexpr>}，因為它根本沒有名字，是當場寫出來的。

\textbf{一行版跟四行版做的是同一件事，形狀不同而已。} 這句話就是這一節存在的全部理由。

順帶把第 \ref{ch:comprehensions} 章的定位升級一下：那一章對 generator expression 的要求
是「你只要認得出來就好」。從這一節開始，標準提高了——\textbf{你要寫得出來}。

\subsection*{什麼時候用哪一種}

判準很簡單：\textbf{一行寫得完，就用小括號；需要多寫幾行邏輯，就寫成有 \texttt{yield}
的函式。}

「多寫幾行邏輯」最常見的情況，第 \ref{ch:comprehensions} 章第 \ref{sec:comp-limits} 節
剛講過一次：\textbf{要接住錯誤的時候}。那一節的結論是「comprehension 裡放不進
\texttt{try}／\texttt{except}，所以寫回 \texttt{for} 迴圈」。genexp 也是一樣的式子，一樣
放不進去——但現在你多了第三條路。

\begin{example}{python}{(clean-scores-yield.py) 一行塞不下的時候，寫成 yield 版}
raw = ["88", "", "92分", "79"]


def clean_scores(texts):
    for text in texts:
        try:
            yield int(text)
        except ValueError:
            print("跳過壞資料:", repr(text))


print(sum(clean_scores(raw)))
\end{example}

\iosample{}{跳過壞資料: '' \\ 跳過壞資料: '92分' \\ 167 \\}

四筆資料裡兩筆是壞的，程式沒有倒，好的兩筆加起來是 167。這段程式碼有兩個好處疊在一起：
它像 \texttt{try}／\texttt{except} 版一樣擋得住髒資料，又像 genexp 一樣可以直接餵給
\texttt{sum}，中間不用先造出一份清單。第 \ref{sec:comp-limits} 節那條煞車在這裡有了
下文——不是只能寫回迴圈，還可以把迴圈包成一台機器。

\stylepoint{中括號還是小括號，取決於一件事：\textbf{你要不要留著那份結果}。只是要餵給
下一步——丟給 \texttt{sum}、丟給 \texttt{for}、丟給 \texttt{max}——用小括號；要留著、
要走第二次、要用索引拿第 3 個，用中括號。這條判準在第 \ref{sec:one-shot} 節會被補上
另外半句，變成這一章給你的最終規則。}

\begin{drills}[label=drills:generator-expression]
\item 判斷題：有 \texttt{nums = [1, 2, 3]}。下面三行各自拿到什麼型別的東西？
\begin{verbatim}
a = [n * 2 for n in nums]
b = (n * 2 for n in nums)
c = sum(n * 2 for n in nums)
\end{verbatim}
\item 判斷題：你要從一份很長的紀錄檔裡，把每一行的日期欄位轉成整數，而那個欄位偶爾會是
空的、轉型會失敗。這個工作該寫成一行小括號版，還是寫成有 \texttt{yield} 的函式？為什麼？
\end{drills}

\section{好處：不先算，就不用先佔位子}
\label{sec:lazy-memory}

到目前為止，這一章一直在講「它不先算」，但沒說過這樣做到底換到了什麼。這一節回答這件事。
\index{惰性求值（lazy evaluation）}

第 \ref{ch:files-exceptions} 章第 \ref{sec:reading-lines} 節已經給過你結論了：逐行讀比
\texttt{readlines()} 省記憶體，所以「很多 Python 使用者的預設習慣是逐行處理」。那一段
論證這裡不重講，只補上當時沒說的那一半——\textbf{檔案物件憑什麼辦得到}。答案上一節就
出現了：它是一台機器，它從頭到尾手上只有一行。

\subsection*{把你自己寫過的那支程式改一次}

第 \ref{ch:csv} 章第 \ref{sec:csv-manual-average} 節那支 \texttt{average-quiz-score.py}
是你親手打過的。它的第一個動作是 \texttt{lines = f.readlines()}，接著造一份空的
\texttt{scores}、跑迴圈 \texttt{append}、再跑一個累加器迴圈把它們加起來。

同一件事，現在可以寫成這樣：

\begin{example}{python}{(total-lazy.py) 算總分，中間一份清單都不造}
import csv

with open("quiz_scores.csv") as f:
    reader = csv.DictReader(f)
    total = sum(int(row["score"]) for row in reader)

print("Total:", total)
\end{example}

\iosample{}{Total: 400 \\}

答案跟第 \ref{sec:csv-manual-average} 節那支程式一樣是 400，但兩者在記憶體裡發生的事完全
不同：

\begin{itemize}
\item \textbf{舊寫法}：整份檔案讀進來 $\to$ 造出一份完整的字串清單 $\to$ 造出一份完整的
數字清單 $\to$ 才開始算。\textbf{同一時間記憶體裡有三份東西。}
\item \textbf{新寫法}：讀一行 $\to$ 轉成一個數字 $\to$ 加進總和 $\to$ 這一行就丟掉了，
換下一行。\textbf{同一時間記憶體裡只有一行。}
\end{itemize}

老實說一句：這份成績檔只有五行，兩種寫法你感覺不出任何差別，跑起來也一樣快。

真正的差別要靠放大數字才看得見。\textbf{如果這不是 5 行，是一份 500 萬行的紀錄檔呢？}
舊寫法必須\textbf{先}把 500 萬個字串塞進記憶體，才能開始算第一筆；新寫法從頭到尾手上
只有一行，500 萬行跟 5 行對它來說是同一回事。這就是那個結論的完整版本。

（順帶一提，這支程式只算了總分。要平均的話還得知道有幾筆，而那件事在這種寫法下沒那麼
順手——原因是下一節的主題。）

\subsection*{你早就在用惰性了：\texttt{range}}

第 \ref{sec:iterable-vs-iterator} 節埋了一句「\texttt{range} 是一批東西，不是機器」，
現在把另一半補上：\texttt{range} 也是惰性的。

\begin{example}{python}{(range-is-lazy.py) 一百萬個數字，其實一個都沒造出來}
numbers = range(1000000)

print(numbers)
print(len(numbers))
print(numbers[500000])

for n in range(1000000000):
    print(n)
    break
\end{example}

\iosample{}{range(0, 1000000) \\ 1000000 \\ 500000 \\ 0 \\}

\texttt{range(1000000)} 印出來還是 \texttt{range(0, 1000000)}，因為它\textbf{根本沒有造
出一百萬個數字}。它手上只記著三件事：從哪開始、到哪結束、每次跳多少。你問它有幾個，它用
減法算給你；你問它第 500000 個是什麼，它用加法算給你。

最後那個十億次的迴圈瞬間就印出 \texttt{0} 然後 \texttt{break} 掉——十億這個數字對它毫無
壓力，因為它一次只需要生一個。但如果你把同一個 \texttt{range} 寫成
\texttt{list(range(1000000000))}，那就是另一回事了：那一行是在要求 Python 真的把十億個
數字通通造出來擺好，你的電腦會很不開心。\textbf{這一行不要真的執行。}

\keypoint{惰性不是為了跑得比較快，是為了\textbf{不用先把東西全部生出來}。如果你本來就
要把每一項都走過一遍，兩種寫法的總工作量其實差不多，generator 省下的是
\textbf{同一時間佔的位子}，不是總時間。所以不要把它當成萬用加速器——它解決的是「東西太
多，記憶體放不下」這個問題，不是「跑太慢」那個問題。}

也不是所有東西都惰性。最好記的反例是 \texttt{sorted()}：它非得把全部看過一遍不可，不然
排不出順序，所以它一定會先把資料整份收下來。

\begin{drills}[label=drills:lazy-memory]
\item 問答題：下面兩種寫法，程式跑到一半時，同一時間記憶體裡各有幾份東西？
\begin{verbatim}
scores = [int(line.split(",")[1]) for line in lines]
total = sum(scores)

total = sum(int(line.split(",")[1]) for line in lines)
\end{verbatim}
\item 判斷題：你有一份 300 萬行的伺服器紀錄檔，要從裡面找出所有包含 \texttt{"ERROR"}
的行，把它們寫進另一個檔案。你會用 \texttt{readlines()} 一次讀進來，還是逐行處理？為什麼？
\end{drills}

\section{代價：generator 只能走一次}
\label{sec:one-shot}

上一節把 generator 講得很好。這一節翻面。

不先算好的東西，有一個躲不掉的代價：\textbf{算過的就沒了}。第
\ref{sec:iterable-vs-iterator} 節的學習要點已經講過原理——一台發號機只能用完一次——
這一節要演給你看，而且重點在演出它有多安靜。

\subsection*{第二個迴圈，什麼都沒有}

\begin{example}{python}{(walk-twice.py) 同一台機器，連跑兩個 for}
scores = [88, 76, 92, 65, 79]
passed = (score for score in scores if score >= 70)

print("第一趟：")
for score in passed:
    print(score)

print("第二趟：")
for score in passed:
    print(score)

print("（第二趟到此結束）")
\end{example}

\iosample{}{第一趟： \\ 88 \\ 76 \\ 92 \\ 79 \\ 第二趟： \\ （第二趟到此結束） \\}

「第二趟：」跟「（第二趟到此結束）」之間那一片空白，就是這一節的全部重點。

注意它做了什麼、沒做什麼。第二個 \texttt{for} 迴圈\textbf{確實執行了}，Python 老老實實
去按了那台機器，機器說沒有了，於是迴圈跑了零輪就結束。\textbf{沒有 traceback、沒有紅字、
沒有任何警告}，只有一個安安靜靜跑完、什麼都沒做的迴圈。

這種坑跟第 \ref{ch:comprehensions} 章第 \ref{sec:comp-dict-set} 節那個「大括號忘了冒號
拿到 set」是同一類：\textbf{不報錯，只給你錯的東西}。而不報錯的錯，永遠是最難抓的。

\subsection*{更陰險的版本：走到一半也算走過}

上面那個至少還算好認——第二趟整個是空的。真正難抓的是這種：

\begin{example}{python}{(half-eaten.py) 只問了一個問題，資料就少了一截}
scores = (score for score in [65, 70, 88, 76, 92])

print("有沒有人考 80 分以上？", any(score >= 80 for score in scores))
print("剩下的資料：", list(scores))
\end{example}

\iosample{}{有沒有人考 80 分以上？ True \\ 剩下的資料： [76, 92] \\}

原本五筆資料，走完一個 \texttt{any(...)} 之後只剩兩筆。為什麼？因為第
\ref{ch:comprehensions} 章第 \ref{sec:reducing-builtins} 節教過，\texttt{any} 找到第一個
成立的就停手——它按到 88 就得到答案了，前面的 65、70 已經被按掉，88 也被拿走了，機器就
停在那裡。接下來 \texttt{list(scores)} 只能收到剩下的 76 跟 92。

拿到的不是空的，是\textbf{少了前面幾筆的}。這種結果放進報表裡，你很可能一整天都不會發現。

\subsection*{三件你不能對 generator 做的事}

前兩個坑都不報錯。有兩件事倒是會直接罵你，正好可以擺在一起看。

\begin{example}{python}{(no-len.py) 問它有幾項}
scores = [88, 76, 92, 65, 79]
passed = (score for score in scores if score >= 70)

print(len(passed))
\end{example}

\begin{verbatim}
Traceback (most recent call last):
  File "no-len.py", line 4, in <module>
    print(len(passed))
          ~~~^^^^^^^^
TypeError: object of type 'generator' has no len()
\end{verbatim}

\texttt{has no len()}——它自己也不知道有幾項，因為它根本還沒算呢。要知道有幾項，就得先
全部算完，那不就違背了它存在的意義？

換成用索引拿第一項，一樣不行：

\begin{code}{python}{想拿第 0 項}
print(passed[0])
\end{code}

\begin{verbatim}
Traceback (most recent call last):
  File "no-index.py", line 4, in <module>
    print(passed[0])
          ~~~~~~^^^
TypeError: 'generator' object is not subscriptable
\end{verbatim}

\texttt{not subscriptable} 就是「不能用中括號取值」。它只會照順序往前吐，沒辦法跳到第 3 個。

第三件事就是前面那個：再走一次 \texttt{for}。而它\textbf{不會}罵你。

\keypoint{generator 走完之後不會壞掉、不會報錯，它只是變成空的。你拿到的不是錯誤訊息，
是一份空清單、一個 0、一個什麼都沒印的迴圈——\textbf{看起來像「資料沒讀到」，其實是
「已經被讀掉了」}。（空清單在 \texttt{if} 裡還是 falsy，所以連 \texttt{if data:} 這種
檢查都會安靜地跳過去。）下次你發現一段程式莫名其妙拿到空結果，先想想這份資料是不是已經
被走過一次了。順帶釐清一件事：不是 \texttt{passed} 這個變數被清空了，是 \texttt{passed}
指的那台機器已經走到底了——機器還在，只是裡面沒東西了。}

\subsection*{怎麼救：兩條路}

\textbf{第一條：重新拿一台。} generator function 再呼叫一次、genexp 再寫一次，就是全新的
一台。這也正是 list 不需要你操心的原因——list 是一批東西，每次 \texttt{for} 開跑都會自己
去要一台新的機器。

\textbf{第二條：要走很多次就收成 list。}

\begin{example}{python}{(rescue-with-list.py) 收成 list 之後，什麼都能做了}
scores = [88, 76, 92, 65, 79]

passed = list(score for score in scores if score >= 70)

print(len(passed))
print(passed[0])
for score in passed:
    print(score)
for score in passed:
    print(score)
\end{example}

\iosample{}{4 \\ 88 \\ 88 \\ 76 \\ 92 \\ 79 \\ 88 \\ 76 \\ 92 \\ 79 \\}

外面套一層 \texttt{list()}，剛才不能做的三件事全都可以了：有 \texttt{len()}、可以索引、
兩個 \texttt{for} 都有輸出。代價就是上一節講的那個——這四項現在真的躺在記憶體裡了。

\stylepoint{這一章的最終判準，兩句話：\textbf{只走一次、資料可能很大、或根本不知道有多長，
就用 generator；要走第二次、要 \texttt{len()}、要用索引，就老實用 list。} 這是第
\ref{sec:generator-expression} 節那條「要不要留著結果」的完整版——那裡問的是你要不要，
這裡告訴你不要的話會怎樣。}

\begin{drills}[label=drills:one-shot]
\item 預測輸出：下面這段印出什麼？
\begin{verbatim}
g = (n for n in [1, 2, 3])
print(list(g))
print(list(g))
\end{verbatim}
\item 抓蟲：下面這段程式，總分印得好好的，明細卻一行都沒有。為什麼？要怎麼修？
\begin{verbatim}
scores = (int(row["score"]) for row in rows)

print("總分：", sum(scores))
print("明細：")
for score in scores:
    print(score)
\end{verbatim}
\item 判斷題：你要從一份很大的檔案裡篩出符合條件的行，然後做兩件事：算出總共幾行，以及
把它們全部印出來。這個需求該不該先收成 list？為什麼？
\end{drills}

\section{別人已經做好的機器：\texttt{itertools} 一瞥}
\label{sec:itertools-peek}

你在第 \ref{ch:modules} 章第 \ref{sec:collections-and-itertools} 節見過
\texttt{itertools} 這個模組，當時挑的是 \texttt{product}——把兩批東西配對出所有組合。
那一節說它是「專門處理『疊代』相關的各種小技巧」的模組，但那時候你還不知道「疊代」是
什麼，所以只看得懂那一個例子。
\index{itertools（count()／islice()／chain()）}

現在你知道了，於是這個模組裡另外一批工具才真正講得通：\textbf{裡面每一個函式，吐出來的
都是一台機器。} 這一節挑三個，每一個都能直接打掉你前面章節寫過的某段程式碼。

定位先講死，跟第 \ref{ch:modules} 章完全一樣：\textbf{這一節不要求你記住任何一個函式的
名字}，只要求你知道「遇到這類需求時，這個模組裡大概有現成的」，需要時查文件就好。前面
六節是要精熟的，這一節不是。

\subsection*{\texttt{islice}：在機器上直接切一段}

第 \ref{ch:csv} 章那支 \texttt{average-quiz-score.py} 為了跳過標題列，寫的是
\texttt{lines[1:]}——而要能寫 \texttt{[1:]}，就得先 \texttt{readlines()} 拿到一份完整的
清單。上一節剛說過那樣要把整份檔案吃進記憶體。

\texttt{islice} 可以直接在一台機器上切一段，不用先收成 list。

\begin{example}{python}{(skip-header-islice.py) 不用 readlines 也能跳過標題列}
from itertools import islice

with open("quiz_scores.csv") as f:
    for line in islice(f, 1, None):
        print(line.strip())

print("---")

with open("quiz_scores.csv") as f:
    for line in islice(f, 3):
        print(line.strip())
\end{example}

\iosample{}{Alice,88 \\ Bob,76 \\ Taylor,92 \\ Charlie,65 \\ David,79 \\ --- \\ name,score \\ Alice,88 \\ Bob,76 \\}

\texttt{islice(f, 1, None)} 是「從第 1 項開始，一直到沒有為止」，也就是跳過第 0 項；
\texttt{islice(f, 3)} 只給一個數字，意思是「只要前 3 項」。切法跟你熟悉的切片幾乎一樣，
差別在\textbf{它不需要先有一份清單}，直接架在檔案這台機器上就能用。

\subsection*{\texttt{chain}：兩批東西接成一串走完}

手上有兩份資料，想用一個 \texttt{for} 走完。你原本會寫 \texttt{class\_a + class\_b}——
那會造出一份新的、把兩邊都裝進去的清單。\texttt{chain} 不造那份清單，它只是依序把兩邊
的東西吐出來。

\begin{example}{python}{(chain-two-classes.py) 兩班成績，一個迴圈走完}
from itertools import chain

class_a = [88, 76]
class_b = [92, 65, 79]

for score in chain(class_a, class_b):
    print(score)

print("Total:", sum(chain(class_a, class_b)))
\end{example}

\iosample{}{88 \\ 76 \\ 92 \\ 65 \\ 79 \\ Total: 400 \\}

兩份清單只有五個數字時，這跟 \texttt{class\_a + class\_b} 沒什麼差。但如果接的是兩份很大
的檔案，差別就是上一節那件事了：\texttt{+} 要先造出一份合併好的大清單，\texttt{chain}
從頭到尾只在轉手。

\subsection*{\texttt{count}：一台永遠不會發完的機器}

最後一個純粹是為了把「惰性」推到極端。\texttt{count} 從某個數字開始一直往上數，
\textbf{永遠不會停}。

聽起來很危險，但一點問題也沒有——因為它本來就不打算把東西全部生出來。你只要記得替它封口。

\begin{example}{python}{(endless-count.py) 用 islice 給無限的機器封口}
from itertools import count, islice

for n in islice(count(1), 5):
    print(n)

names = ["Alice", "Bob", "Taylor"]
for number, name in zip(count(1), names):
    print(number, name)
\end{example}

\iosample{}{1 \\ 2 \\ 3 \\ 4 \\ 5 \\ 1 Alice \\ 2 Bob \\ 3 Taylor \\}

上半段用剛學會的 \texttt{islice} 只取前 5 個。下半段更實用：\texttt{zip} 走到短的那一份
用完就停（第 \ref{ch:comprehensions} 章第 \ref{sec:zip-walk} 節那個坑），所以拿一台無限
的機器跟一份三人的名單配對，剛好得到從 1 開始的編號——這是第 \ref{sec:enumerate} 節
\texttt{enumerate} 的另一種寫法。

\textbf{一個警告}：\texttt{list(count())} 會試著把無限多個數字全部造出來，直到把記憶體
吃光為止。\textbf{這一行絕對不要真的執行。} 這正是上一節那條判準的極端案例——有些東西
根本不知道有多長，那就只能惰性，沒有第二個選項。

\subsection*{五張欠條，還完了}

回到這一章的開場。當時說這本書欠你五個解釋，現在逐一對一下：

\begin{enumerate}
\item \texttt{for} 底下「這一輪走到哪、還有沒有下一輪」是誰在管？——是 \texttt{iter()}
要來的那台 iterator，靠 \texttt{next()} 跟 \texttt{StopIteration} 運作。
\item 為什麼檔案物件可以直接放進 \texttt{for}，而且逐行讀比 \texttt{readlines()} 省
記憶體？——因為它本身就是一台 iterator，手上永遠只有一行。
\item \texttt{next(reader)} 是什麼？——按一下那台機器，把第一列拿走。
\item \texttt{<map object at 0x...>} 是什麼？為什麼要 \texttt{list()} 包起來？——一台
你還沒按的機器；\texttt{list()} 是一口氣按到底。
\item comprehension 換成小括號拿到的是什麼？為什麼要小心？——是一台 generator，
小心的理由是它只能走一次。
\end{enumerate}

還有一件事這一章沒有講完，但已經露出形狀：什麼時候該惰性、什麼時候該老實。第
\ref{ch:pythonic} 章會把「寫得像 Python」這件事整個收起來，這一章學到的判斷是其中一塊。

\begin{drills}[label=drills:itertools-peek]
\item 預測輸出：下面兩行各印出什麼？
\begin{verbatim}
from itertools import islice
print(list(islice([10, 20, 30, 40, 50], 1, 4)))
print(list(islice("abcdef", 2)))
\end{verbatim}
\item 判斷題：你要跳過一份 200 萬行檔案的前 3 行，剩下的逐行處理。用
\texttt{f.readlines()[3:]} 還是用 \texttt{islice(f, 3, None)}？為什麼？
\end{drills}

\section*{練習解答}
\addcontentsline{toc}{section}{練習解答}

\paragraph{第~\ref{sec:what-for-does}~節}
\begin{enumerate}
\item 手動版是：\texttt{it = iter([88, 76, 92])} $\to$ \texttt{while True:} $\to$ 縮排一層
\texttt{try: score = next(it)} $\to$ \texttt{except StopIteration: break} $\to$ 最後
\texttt{print(score * 2)}。輸出跟原本的 \texttt{for} 一模一樣：\texttt{176}、
\texttt{152}、\texttt{184}。
\item 先印出 \texttt{10}、再印出 \texttt{20}，第三個 \texttt{next(it)} 停下來，丟出
\texttt{StopIteration}。因為那台機器只有兩個值可以發，第三次按下去就是「發完了」。
\end{enumerate}

\paragraph{第~\ref{sec:iterable-vs-iterator}~節}
\begin{enumerate}
\item 第一行報錯：\texttt{TypeError: 'str' object is not an iterator}——字串是一批東西，
不是機器。第二行印出 \texttt{c}，因為 \texttt{iter("cat")} 先要來一台機器，才有得按。
\item 印出兩行：\texttt{A fig} 跟 \texttt{A kiwi}。第二個迴圈一行都沒有，因為
\texttt{it} 這台機器在第一個迴圈就被用完了，而且不會報錯。
\item 「一批東西」：\texttt{[1, 2, 3]}、\texttt{"hello"}、\texttt{range(10)}。
「一台發號機」：檔案物件、\texttt{csv.reader(f)} 的結果。（檔案物件比較特別，它同時
兩種身分都算得上，但關鍵在於它\textbf{可以直接被 \texttt{next()} 按}，而且用完就沒了。）
\end{enumerate}

\paragraph{第~\ref{sec:yield-basics}~節}
\begin{enumerate}
\item 依序印出：\texttt{a}、\texttt{1}、\texttt{b}、\texttt{2}、\texttt{c}。函式裡的
\texttt{print} 跟迴圈裡的 \texttt{print} 是交錯出現的，因為函式每次都停在 \texttt{yield}
那裡等你按下一次。
\item 改寫成：\texttt{def doubled(items):} $\to$ \texttt{for n in items:} $\to$
\texttt{yield n * 2}。呼叫端要寫 \texttt{list(doubled([1, 2, 3]))} 才會拿到
\texttt{[2, 4, 6]}；直接 \texttt{print(doubled([1, 2, 3]))} 只會印出一個
\texttt{<generator object ...>}。
\item 因為 \texttt{greet([...])} 只是造了一台機器就把它丟了，函式裡的程式碼一行都沒跑。
要有人去按它才會動，例如 \texttt{for line in greet(["Alice", "Bob"]): print(line)}，
或是 \texttt{print(list(greet(["Alice", "Bob"])))}。
\end{enumerate}

\paragraph{第~\ref{sec:generator-expression}~節}
\begin{enumerate}
\item \texttt{a} 是 list（中括號，當場算好），\texttt{b} 是 generator（小括號，還沒算），
\texttt{c} 是整數 \texttt{12}——因為 \texttt{sum} 已經把那台機器按到底、算出答案了。
\item 寫成有 \texttt{yield} 的函式。轉型會失敗表示你需要 \texttt{try}／\texttt{except}，
而一行小括號版跟 comprehension 一樣是式子，裡面放不進 \texttt{try}（第
\ref{sec:comp-limits} 節那條煞車）。寫成函式就能一邊接住錯誤、一邊照樣 \texttt{yield}
好的那幾筆。
\end{enumerate}

\paragraph{第~\ref{sec:lazy-memory}~節}
\begin{enumerate}
\item 第一種寫法在 \texttt{sum} 開跑的那一刻，記憶體裡有 \texttt{lines} 跟
\texttt{scores} 兩份完整的資料；第二種寫法從頭到尾只有「目前這一行」跟一個累加中的
\texttt{total}。
\item 逐行處理。300 萬行用 \texttt{readlines()} 要先整份塞進記憶體才能開始篩，而這個
任務根本不需要「同時看到所有行」——看一行、判斷一行、寫一行就夠了。
\end{enumerate}

\paragraph{第~\ref{sec:one-shot}~節}
\begin{enumerate}
\item 第一行印出 \texttt{[1, 2, 3]}，第二行印出 \texttt{[]}。第一個 \texttt{list(g)}
已經把那台機器按到底了，第二次去按只會拿到空的——而且不會報錯。
\item 因為 \texttt{sum(scores)} 已經把 \texttt{scores} 整台按到底了，輪到 \texttt{for}
的時候一項都不剩。修法有兩條：先收成 list（\texttt{scores = [int(row["score"]) for row
in rows]}），之後 \texttt{sum} 跟 \texttt{for} 都對著同一份 list；或是需要的時候各造一台
新的 generator。
\item 該。你要對同一份結果做兩件事——算數量跟印出來——那就是「要走第二次」，符合先收成
list 的條件。除非那份結果大到裝不下，那就得改成一趟走完同時做兩件事（一邊印一邊數）。
\end{enumerate}

\paragraph{第~\ref{sec:itertools-peek}~節}
\begin{enumerate}
\item 第一行印出 \texttt{[20, 30, 40]}（從第 1 項切到第 4 項之前，跟切片一樣不含結尾）；
第二行印出 \texttt{['a', 'b']}（只給一個數字就是「只要前 2 項」）。
\item 用 \texttt{islice(f, 3, None)}。\texttt{readlines()[3:]} 要先把 200 萬行全部讀進
記憶體造成一份清單，才能切掉前 3 行；\texttt{islice} 直接架在檔案這台機器上，讀一行處理
一行，記憶體裡永遠只有一行。
\end{enumerate}

\chapter{遞迴}
\label{ch:recursion}

第 \ref{ch:functions} 章第 \ref{sec:function-scope} 節的結尾留了一筆帳。那裡說：函式主體
裡除了可以呼叫別的函式，也可以呼叫函式自己，這種寫法叫做遞迴；至於完整的討論——怎麼讓它
停下來、什麼時候值得用它、什麼時候用迴圈更好——留到這一章。

這一章就是來還這筆帳的。三個問題裡，看起來最不像問題的是第三個，但它其實才是這一章真正
要交給你的東西。遞迴不是比較厲害的迴圈，也不是寫出來比較好看的迴圈；它是一種在特定形狀的
問題上非常好用、在其他形狀的問題上只會把事情弄複雜的工具。這一章的目標，是讓你會寫它，
更重要的是讓你看到一個問題的時候，知道這次該不該用它。

這一章分三段推進：

\begin{itemize}
\item 第一段把遞迴看清楚：一個函式呼叫自己，它到底怎麼運作、又怎麼停下來，還有怎麼讓
程式自己把「一層一層疊上去、又一層一層收回來」的過程印出來給你看
（第 \ref{sec:recursion-basics}～\ref{sec:recursion-trace} 節）。
\item 第二段用兩個經典的小例子練形狀：階乘與費氏數列。這兩個例子有個共同點——它們其實
都用迴圈寫更好。它們留在這一章的理由只有一個：夠小，小到你可以把整個過程從頭追到尾
（第 \ref{sec:recursion-return}～\ref{sec:fibonacci-cost} 節）。
\item 第三段才是遞迴真正划算的場合：走訪巢狀資料、走訪資料夾樹，順帶看一眼標準函式庫
替遞迴準備的一個小工具，最後收在一份「什麼時候該用、什麼時候不該用」的判準上
（第 \ref{sec:nested-data-walk}～\ref{sec:recursion-or-loop} 節）。
\end{itemize}

\section{遞迴初探：函式呼叫自己}
\label{sec:recursion-basics}

先從一件跟程式完全無關的事說起。你決定把一個塞滿雜物的抽屜整理乾淨，做法很單純：東西一件
一件拿出來歸位，筆放筆筒、單據丟掉、鑰匙掛起來。做到一半，你拿出一個鐵盒；打開鐵盒，裡面
又是一堆沒整理的東西，而且裡面還有一個更小的盒子。

這時候你不會為了「整理鐵盒」另外發明一套辦法。你會做的事跟整理抽屜完全一樣：東西一件一件
拿出來歸位，碰到盒子就打開它，用同一套做法處理裡面的東西。而且你一點也不擔心整理不完——
每打開一層，裡面的東西一定比外面少，總會有一個盒子，打開之後裡面沒有盒子了。

「做著做著又碰到同一件事，於是用同一套做法再做一次」——這個形狀，就是這一章要講的東西。

函式主體裡除了可以呼叫別的函式，其實也可以呼叫自己，這種寫法叫做遞迴（recursion）。第一
次聽到「函式呼叫自己」通常會覺得這樣寫會不會永遠繞不出來，這一節就用一個小例子，看看遞迴
實際上是怎麼運作、又怎麼停下來的。
\index{遞迴}

\begin{code}{python}{遞迴函式的基本形狀}
def function_name(parameter):
    if base_case_reached:
        code (does not call itself again)
        return
    code
    function_name(parameter closer to the base case)
\end{code}

\begin{example}{python}{(countdown.py) 倒數計時}
def countdown(n):
    if n <= 0:
        print("Liftoff!")
        return
    print(n)
    countdown(n - 1)


countdown(3)
\end{example}

\iosample{}{3 \\ 2 \\ 1 \\ Liftoff! \\}

追蹤這支程式實際發生的事：呼叫 \texttt{countdown(3)}，因為 \texttt{3 <= 0} 不成立，先
印出 \texttt{3}，再呼叫 \texttt{countdown(2)}；\texttt{countdown(2)} 又印出
\texttt{2}，再呼叫 \texttt{countdown(1)}；一路下去，直到呼叫 \texttt{countdown(0)}，
這次 \texttt{0 <= 0} 成立，印出「Liftoff!」之後就 \texttt{return}，不再呼叫自己。整個過程
總共呼叫了 \texttt{countdown} 四次（\texttt{n} 分別是 3、2、1、0），每次呼叫都有自己
獨立的一份 \texttt{n}（呼應第 \ref{ch:functions} 章的作用域規則：每次呼叫都有自己的一份
區域變數）。

讓函式停止繼續呼叫自己的那個判斷（這裡是 \texttt{n <= 0}），叫做終止條件。
\index{終止條件}
每一個遞迴
函式都一定要有終止條件，而且每次呼叫自己的時候，傳進去的參數都要比上一次更靠近終止條件
（這裡是 \texttt{n} 每次都減 1，遲早會碰到 \texttt{n <= 0}）。如果忘了寫終止條件，會發生
什麼事？

\begin{example}{python}{(countdown-broken.py) 忘記寫終止條件}
def countdown(n):
    print(n)
    countdown(n - 1)


countdown(3)
\end{example}

\begin{verbatim}
3
2
1
0
-1
-2
...
  File "countdown-broken.py", line 3, in countdown
    countdown(n - 1)
  File "countdown-broken.py", line 3, in countdown
    countdown(n - 1)
  [Previous line repeated 996 more times]
RecursionError: maximum recursion depth exceeded
\end{verbatim}

這支程式會一路印出 3、2、1、0、-1、-2……完全停不下來，因為沒有任何一個判斷會讓它
\texttt{return} 而不再呼叫自己。Python 沒辦法讓函式無止盡地一直呼叫下去，每多呼叫一層
就會多佔用一點電腦的資源，超過某個深度上限，Python 會直接放棄並丟出
\texttt{RecursionError}，
\index{RecursionError}
訊息是「超過最大遞迴深度」，畫面上還會看到一長串重複的呼叫紀錄
（上面的 \texttt{[Previous line repeated 996 more times]} 就是 Python 自己在說「這段
重複了 996 次，不再一行一行列出來」）。

\keypoint{寫遞迴函式一定要有終止條件，而且每次呼叫自己時，參數都要比上一次更靠近那個終止
條件，不然函式會沒完沒了地呼叫下去，最後被 Python 自己的深度限制擋下來，丟出
\texttt{RecursionError}。第 \ref{ch:loops} 章學過的 \texttt{for}／\texttt{while} 迴圈通常
也能做到遞迴能做的事，遞迴只是多一種思考、拆解問題的方式，不是每個情境都非用不可。}

\begin{drills}[label=drills:recursion-basics]
\item 如果把 \texttt{countdown(3)} 改成呼叫 \texttt{countdown(0)}，會印出什麼？
\item 有一個函式：
\begin{verbatim}
def count_up(n, limit):
    if n > limit:
        return
    print(n)
    count_up(n + 1, limit)
\end{verbatim}
如果呼叫 \texttt{count\_up(1, 4)}，會依序印出哪些數字？
\item 這裡的 \texttt{count\_up} 函式，改用第 \ref{ch:loops} 章教過的 \texttt{for} 迴圈搭配
\texttt{range()} 可以達到一樣的效果嗎？你會怎麼寫？
\end{drills}

\section{看見它一層一層疊上去，又一層一層收回來}
\label{sec:recursion-trace}

上一節那段「追蹤這支程式實際發生的事」，是這本書講給你聽的。講得再清楚，也還是別人告訴你
的；真正能讓人放下心裡那個「它到底會不會停」的疑慮的，是自己親眼看到過程。

這一節要做的事很單純：讓程式自己把過程印出來。方法只有一個小技巧——在遞迴函式裡多帶一個
參數，記住「現在是第幾層」，用它決定要縮排幾格；進入函式的時候印一行，要離開之前再印
一行。

\begin{code}{python}{幫遞迴函式加上縮排追蹤}
def function_name(parameter, depth=0):
    indent = "    " * depth
    print(indent + "進入這一層")
    if base_case_reached:
        print(indent + "離開這一層")
        return
    function_name(next_parameter, depth + 1)
    print(indent + "離開這一層")
\end{code}

\texttt{depth} 的預設值是 \texttt{0}，所以最外面那次呼叫可以照舊寫
\texttt{countdown(3)}，不用自己傳；只有函式呼叫自己的時候才會把 \texttt{depth + 1} 傳
下去，於是每深一層就多縮四格。拿上一節那支 \texttt{countdown} 直接改（題目不換，這樣你
只需要注意「多印了什麼」）：

\begin{example}{python}{(countdown-trace.py) 把倒數的每一層印出來}
def countdown(n, depth=0):
    indent = "    " * depth
    print(f"{indent}進入 countdown({n})")
    if n <= 0:
        print(f"{indent}Liftoff!")
        print(f"{indent}離開 countdown({n})")
        return
    print(f"{indent}{n}")
    countdown(n - 1, depth + 1)
    print(f"{indent}離開 countdown({n})")


countdown(3)
\end{example}

\iosample{}{進入 countdown(3) \\ 3 \\ ~~~~進入 countdown(2) \\ ~~~~2 \\
~~~~~~~~進入 countdown(1) \\ ~~~~~~~~1 \\ ~~~~~~~~~~~~進入 countdown(0) \\
~~~~~~~~~~~~Liftoff! \\ ~~~~~~~~~~~~離開 countdown(0) \\ ~~~~~~~~離開 countdown(1) \\
~~~~離開 countdown(2) \\ 離開 countdown(3) \\}

這個形狀就是這一節的重點，值得盯著看一會兒。輸出先一路往右縮排，縮到最深的那一層
（\texttt{countdown(0)}）之後轉向，再一層一層往左退回來。往右走的那一半是「疊上去」：
每一次呼叫自己，都是在上一層還沒做完的情況下，又開了新的一層。往左退的那一半是「收回
來」：最裡面那一層 \texttt{return} 之後，控制權交回給呼叫它的那一層，那一層才把自己
剩下的 \texttt{print} 做完，然後再交回給更外面一層。

注意「離開 countdown(3)」是最後印出來的。最早開始的那一層，是最晚結束的——因為它一直
在等裡面那一層做完。這件事講出來有點繞，但看著這個輸出就一目了然。

還有一個一直被提到、但到現在才真正看得見的東西：\textbf{每一層的 \texttt{n} 是各自
獨立的}。畫面上同時存在 \texttt{countdown(3)}、\texttt{countdown(2)}、
\texttt{countdown(1)}、\texttt{countdown(0)} 四層，每一層印出來的 \texttt{n} 都不一樣，
而且退回來的時候，每一層都還記得自己原本那個 \texttt{n}——「離開 countdown(2)」印的
還是 \texttt{2}。用第 \ref{ch:functions} 章第 \ref{sec:function-nametag} 節的說法：這
四個 \texttt{n} 是四組不同的名字牌，貼在四個不同的數字上；它們同名，但從來不是同一個
東西。

這種只在「想搞清楚程式在幹嘛」的時候才加、弄懂之後就會拿掉的程式碼，有個很貼切的稱呼：
鷹架。
\index{鷹架（scaffolding）}
蓋房子的時候要搭，房子蓋好就拆掉。花五分鐘替一支遞迴函式搭一次鷹架，常常比盯著
程式碼想半小時還快。

\subsection*{traceback 也是一張堆疊的照片}

回頭看第 \ref{sec:recursion-basics} 節那支 \texttt{countdown-broken.py} 丟出來的
\texttt{RecursionError}。第 \ref{ch:debugging} 章第 \ref{sec:call-stack} 節已經教過怎麼
讀這種有好幾組 \texttt{File}／\texttt{line} 的 traceback：由上到下就是呼叫發生的順序，
最下面才是真正出錯的那一行。那一節的讀法在這裡完全適用，不用重學。

這次唯一特別的地方是：那幾組 \texttt{File}／\texttt{line} 全部是同一個函式、同一行。
換句話說，這串 traceback 印的不是「A 呼叫 B 呼叫 C」，是「\texttt{countdown} 呼叫
\texttt{countdown} 呼叫 \texttt{countdown}……」。而
\texttt{[Previous line repeated 996 more times]} 這一行，就是 Python 直接告訴你它疊了
幾層——這是上面那個縮排輸出的另一種呈現，只是這一次是 Python 免費附贈的。

順帶一提，Python 有辦法把那個深度上限調高（標準函式庫裡有個
\texttt{sys.setrecursionlimit}），但幾乎沒有正當理由去動它。會撞到上限，絕大多數時候
代表終止條件寫錯了或走不到，不是上限太小。

\begin{drills}[label=drills:recursion-trace]
\item 把 \texttt{countdown-trace.py} 最後一行改成 \texttt{countdown(2)}，輸出會長什麼
樣子？縮排最深的是哪一行？
\item 有人寫了下面這支程式，想從 10 開始每次減 2，數到 1 為止：
\begin{verbatim}
def shrink(n):
    print(n)
    if n == 1:
        return
    shrink(n - 2)


shrink(10)
\end{verbatim}
執行之後畫面上先印出 \texttt{10}、\texttt{8}、\texttt{6}、\texttt{4}、\texttt{2}、
\texttt{0}、\texttt{-2}……一路往負數跑，最後結束在：
\begin{verbatim}
  File "shrink.py", line 5, in shrink
    shrink(n - 2)
    ~~~~~~^^^^^^^
  [Previous line repeated 996 more times]
RecursionError: maximum recursion depth exceeded
\end{verbatim}
這個函式明明有終止條件，為什麼還是停不下來？從 traceback 看得出它大約疊了幾層嗎？
\item 如果把 \texttt{countdown-trace.py} 裡「離開」那兩行 \texttt{print} 都刪掉，只留
「進入」的那一行，輸出會少掉什麼？剩下的部分還看得出「一層一層收回來」嗎？
\end{drills}

\section{讓每一層把答案交回上一層：階乘}
\label{sec:recursion-return}

到目前為止的 \texttt{countdown} 只做一件事：印東西。印完就沒了，沒有人需要拿到任何結果。
接下來這一步是新的：\textbf{每一層都要算出一個值，交回給呼叫它的那一層。}

生活裡有對應的畫面。你想知道一件事，問了旁邊的人，他不知道，於是他去問他認識的人；那個人
也不知道，再去問下一個……直到有人知道答案，答案才一個一個往回傳，最後傳到你手上。
下去的是問題，上來的是答案。遞迴函式的回傳值就是那個「上來的答案」。

拿階乘當例子。4 的階乘是 $4 \times 3 \times 2 \times 1$，也就是 24。用遞迴的
角度看它，只需要一句話：\texttt{n} 的階乘，就是 \texttt{n} 乘上「\texttt{n - 1} 的
階乘」。而 \texttt{0} 的階乘規定是 1，那就是終止條件。

\begin{code}{python}{有回傳值的遞迴形狀}
def function_name(parameter):
    if base_case_reached:
        return simple_answer
    return combine(parameter, function_name(smaller_parameter))
\end{code}

跟第 \ref{sec:recursion-basics} 節那個片語比起來，形狀沒變，兩塊還是那兩塊：終止條件
那一行不再呼叫自己，直接給出一個現成的答案；遞迴步驟那一行呼叫自己，然後拿它交回來的
東西再加工一次。

\begin{example}{python}{(factorial.py) 用遞迴算階乘}
def factorial(n):
    if n == 0:
        return 1
    return n * factorial(n - 1)


print(factorial(4))
print(factorial(0))
\end{example}

\iosample{}{24 \\ 1 \\}

\texttt{if n == 0: return 1} 是終止條件，\texttt{return n * factorial(n - 1)} 是遞迴
步驟。參數每次減 1，一定會走到 \texttt{0}。

光看程式碼還是有點抽象，所以把上一節的鷹架搭上去——這次鷹架不是新東西，是拿來當工具用：

\begin{example}{python}{(factorial-trace.py) 看清楚答案是怎麼一層一層交回來的}
def factorial(n, depth=0):
    indent = "    " * depth
    print(f"{indent}進入 factorial({n})")
    if n == 0:
        print(f"{indent}回傳 1")
        return 1
    result = n * factorial(n - 1, depth + 1)
    print(f"{indent}回傳 {result}")
    return result


print(factorial(4))
\end{example}

\iosample{}{進入 factorial(4) \\ ~~~~進入 factorial(3) \\ ~~~~~~~~進入 factorial(2) \\
~~~~~~~~~~~~進入 factorial(1) \\ ~~~~~~~~~~~~~~~~進入 factorial(0) \\
~~~~~~~~~~~~~~~~回傳 1 \\ ~~~~~~~~~~~~回傳 1 \\ ~~~~~~~~回傳 2 \\ ~~~~回傳 6 \\
回傳 24 \\ 24 \\}

跟 \texttt{countdown} 的追蹤比起來，往右的那一半幾乎一模一樣：五層呼叫一層一層疊上去，
一路疊到 \texttt{factorial(0)}。差別全在往左退的那一半——這次每一層退回來的時候，手上
都帶著一個數字：\texttt{1}、\texttt{1}、\texttt{2}、\texttt{6}、\texttt{24}。
每一層拿到下一層交回來的答案，乘上自己的 \texttt{n}，再交給上一層。

還有一件事值得注意：\texttt{factorial(4)} 這一層在寫下 \texttt{n * factorial(n - 1)}
的時候，並不需要知道 \texttt{factorial(3)} 是怎麼算出來的，只要相信它會回傳正確的值就
好。這是遞迴最省力的地方——你每次只需要想清楚「這一層要做什麼」，剩下的交給下一層。

\subsection*{有終止條件，不代表走得到}

第 \ref{sec:recursion-basics} 節的學習要點說了兩件事：要有終止條件，而且每次都要更靠近
它。前面幾個例子都只示範了第一件事出問題的樣子。第二件事出問題，長得比較陰險——程式看
起來完全正常，也真的有一行 \texttt{if} 寫著停，但就是永遠碰不到。

\begin{example}{python}{(factorial-negative.py) 傳進一個永遠碰不到 0 的數}
def factorial(n):
    if n == 0:
        return 1
    return n * factorial(n - 1)


print(factorial(-1))
\end{example}

\begin{verbatim}
Traceback (most recent call last):
  File "factorial-negative.py", line 7, in <module>
    print(factorial(-1))
          ~~~~~~~~~^^^^
  File "factorial-negative.py", line 4, in factorial
    return n * factorial(n - 1)
               ~~~~~~~~~^^^^^^^
  File "factorial-negative.py", line 4, in factorial
    return n * factorial(n - 1)
               ~~~~~~~~~^^^^^^^
  [Previous line repeated 996 more times]
RecursionError: maximum recursion depth exceeded
\end{verbatim}

\texttt{-1} 減 1 是 \texttt{-2}，再減是 \texttt{-3}，永遠不會等於 \texttt{0}。終止條件
明明就寫在那裡，卻是個擦身而過的目標。所以碰到 \texttt{RecursionError} 的時候，有一組
固定的兩個問題可以問：有沒有終止條件？有的話，傳進去的參數保證走得到它嗎？把
\texttt{if n == 0} 改成 \texttt{if n <= 0}，或是在函式開頭先擋掉負數，兩種修法都行。

\subsection*{老實說：這個例子用迴圈寫更好}

寫完了，該講實話。階乘用迴圈寫比較好：

\begin{example}{python}{(factorial-loop.py) 同一件事的迴圈版}
def factorial(n):
    result = 1
    for i in range(2, n + 1):
        result = result * i
    return result


print(factorial(4))
\end{example}

\iosample{}{24 \\}

迴圈版跟遞迴版一樣長、一樣好讀，但它不會疊層、不會有深度上限的問題，傳進去 \texttt{-1}
也只是回傳 \texttt{1}，不會炸。如果哪天你在真正的專案裡要算階乘，寫這一版。

那階乘為什麼還留在這一章？因為它夠小。小到你可以把整個過程從第一層追到最後一層，一行
不漏，而這件事對「相信遞迴會停」非常有幫助。它是拿來看清楚形狀用的，不是拿來當範本用
的。第 \ref{sec:nested-data-walk} 節與第 \ref{sec:folder-tree} 節會給你兩個反過來的
例子：用迴圈寫會很痛苦，用遞迴寫剛剛好。

\begin{drills}[label=drills:recursion-return]
\item 把 \texttt{factorial-trace.py} 最後一行改成 \texttt{print(factorial(3))}，輸出會
長什麼樣子？往左退回來的那幾行，數字依序是什麼？
\item 有人把終止條件寫成 \texttt{if n == 1: return 1}（想說「1 的階乘就是 1」）。
\texttt{factorial(4)} 還是算得出 24 嗎？那 \texttt{factorial(0)} 呢？會發生什麼事？
\item 下面這個函式用遞迴把一份清單加總：
\begin{verbatim}
def total(numbers):
    if len(numbers) == 0:
        return 0
    return numbers[0] + total(numbers[1:])
\end{verbatim}
\texttt{total([3, 1, 4, 1, 5])} 的結果是什麼？請把它改寫成 \texttt{for} 迴圈版。
\end{drills}

\section{費氏數列：形狀很漂亮，帳單很可怕}
\label{sec:fibonacci-cost}

費氏數列是這樣一串數字：0、1、1、2、3、5、8、13、21……規則是「每一項等於前兩項相加」，
前兩項規定是 0 跟 1。這個規則翻成遞迴幾乎不用轉換，兩個終止條件、一個遞迴步驟，三行
寫完：

\begin{example}{python}{(fibonacci.py) 用遞迴算費氏數列}
def fibonacci(n):
    if n <= 1:
        return n
    return fibonacci(n - 1) + fibonacci(n - 2)


for i in range(10):
    print(i, fibonacci(i))
\end{example}

\iosample{}{0 0 \\ 1 1 \\ 2 1 \\ 3 2 \\ 4 3 \\ 5 5 \\ 6 8 \\ 7 13 \\ 8 21 \\ 9 34 \\}

到這裡都很順，程式短、規則清楚、答案正確。問題藏在一個前面每個例子都沒有的地方：這個
函式呼叫自己\textbf{兩次}。

\subsection*{數數看它到底做了多少事}

不用談任何理論，直接數。加一個計數器，每進入函式一次就加 1（\texttt{global} 是第
\ref{ch:functions} 章教過的，用來表示「我要改的是函式外面那個變數」）：

\begin{example}{python}{(fibonacci-count.py) 數它一共呼叫了自己幾次}
call_count = 0


def fibonacci(n):
    global call_count
    call_count = call_count + 1
    if n <= 1:
        return n
    return fibonacci(n - 1) + fibonacci(n - 2)


for target in [5, 10, 20, 30]:
    call_count = 0
    answer = fibonacci(target)
    print(f"fibonacci({target}) = {answer}，一共呼叫了 {call_count} 次")
\end{example}

\iosample{}{fibonacci(5) = 5，一共呼叫了 15 次 \\ fibonacci(10) = 55，一共呼叫了 177 次 \\
fibonacci(20) = 6765，一共呼叫了 21891 次 \\
fibonacci(30) = 832040，一共呼叫了 2692537 次 \\}

\texttt{n} 從 5 加到 30，只多了 25，呼叫次數卻從 15 次變成兩百六十九萬次。在作者的電腦
上，\texttt{fibonacci(35)} 要跑將近一秒才會出來；再往上加幾，就要開始等了。

為什麼會這樣？算第 5 項要先算第 4 項跟第 3 項；而算第 4 項的時候，又得再算一次第 3 項
跟第 2 項。第 3 項被算了兩次，而且這兩次算出來的東西一模一樣。再往下，第 2 項被算了
三次、第 1 項被算了五次……每往下一層，重複的份數就再翻一次，而且沒有任何人記得上一次
算過。

請注意問題出在哪裡：\textbf{不是遞迴本身有問題，是這個問題的形狀有問題。} 它每一步會
岔成兩條路，而兩條路又長出一模一樣的子問題。\texttt{countdown} 跟 \texttt{factorial}
都沒有這個毛病，因為它們每次只呼叫自己一次，一路往下就是一條線，四層就是四次呼叫。

\keypoint{一個遞迴函式如果每次呼叫自己不只一次，要特別小心：重複的工作量不是用加的，
是用乘的長大。看到一個遞迴函式的主體裡出現兩個以上的自我呼叫，先停下來想一想，同一件
事會不會被算很多遍。}

\subsection*{老實說：這個例子也用迴圈寫更好}

跟階乘一樣，費氏數列在真正的程式裡該用迴圈寫。而且這次不只是「差不多好」，是明顯更好：

\begin{example}{python}{(fibonacci-loop.py) 同一件事的迴圈版}
def fibonacci(n):
    previous = 0
    current = 1
    for i in range(n):
        previous, current = current, previous + current
    return previous


for i in range(10):
    print(i, fibonacci(i))
\end{example}

\iosample{}{0 0 \\ 1 1 \\ 2 1 \\ 3 2 \\ 4 3 \\ 5 5 \\ 6 8 \\ 7 13 \\ 8 21 \\ 9 34 \\}

迴圈版從頭往上算一遍，每個值只算一次，\texttt{fibonacci(30)} 只跑 30 圈就結束，不是
兩百多萬次。它甚至可以輕鬆算 \texttt{fibonacci(1000)}，遞迴版連試都不用試。

那費氏數列為什麼還留在這一章？跟階乘一樣：它把「遞迴的形狀」照得非常清楚——規則怎麼寫，
程式就怎麼寫，幾乎不用翻譯。它同時也是這一章最好的警告：程式寫得像規則本身，不代表
它跑起來划算。

不過，「同樣的事被重算了幾千次」這件事，其實還有另一種處理方式——不是改寫成迴圈，而是
讓函式記得自己算過什麼。這個坑先挖在這裡，第 \ref{sec:lru-cache} 節回來填。

\begin{drills}[label=drills:fibonacci-cost]
\item 照 \texttt{fibonacci-count.py} 的做法，\texttt{fibonacci(5)} 一共會呼叫幾次
\texttt{fibonacci}？（提示：輸出裡已經有答案了，但請試著自己在紙上把
\texttt{fibonacci(5)} 岔出去的兩條路寫出來，數數看。）
\item 為什麼 \texttt{countdown(30)} 完全沒有這個問題？它跟 \texttt{fibonacci(30)} 差在
哪一個字？
\item 不看課文，自己把費氏數列的迴圈版寫一次。寫完之後跟 \texttt{fibonacci-loop.py}
對照，看看你的版本 \texttt{fibonacci(0)} 跟 \texttt{fibonacci(1)} 是不是也對。
\end{drills}

\section{遞迴真正划算的地方（一）：走訪巢狀資料}
\label{sec:nested-data-walk}

前面四節的遞迴函式，每一個都可以用迴圈寫得一樣好或更好。這一節開始不一樣了。

先回想一下你已經做過的事。第 \ref{ch:csv} 章第 \ref{sec:json-basics} 節讀進一份 JSON
之後，是這樣取值的：\texttt{student["contact"]["phone"]}；第 \ref{ch:api} 章第
\ref{sec:parsing-json} 節從天氣 API 的回應裡挖資料，寫的是
\texttt{data["current\_weather"]["temperature"]}。這種一路用中括號往下挖的寫法非常好用，
但它有一個沒說出口的前提：\textbf{你事先就知道有幾層、每一層叫什麼名字。}

那如果不知道呢？你拿到一份沒看過的 JSON，想先看看裡面到底有些什麼；或者資料本身就長得
不整齊，有的成員多一層社團名單、有的成員沒有。這種時候你沒辦法先決定要寫幾層中括號，
也沒辦法先決定要寫幾層 \texttt{for} 迴圈——因為層數是資料說了算，不是你說了算。

這就是這一章開頭那個抽屜：你不知道盒子裡還有幾個盒子。

\subsection*{一個新的形狀：迴圈包著遞迴}

前面四個遞迴函式都是「一路往下一條線」：一層只往下開一個新的層。從這一節開始，形狀會變
成另一種——\textbf{一層裡面有很多個東西}，要用迴圈跑過每一個，碰到還能再拆開的就遞迴
下去：

\begin{code}{python}{走訪「裡面還有同一種東西」的資料}
def function_name(value):
    for item in value:
        if item is a container:
            function_name(item)
        else:
            handle item directly
\end{code}

這個片語值得記起來，因為下一節的資料夾樹用的是\textbf{完全同一個骨架}，只是把「是不是
容器」換成「是不是資料夾」而已。

要問「這個東西是不是 dict」「是不是 list」，用內建函式
\texttt{isinstance()}：\texttt{isinstance(x, dict)} 問的是「\texttt{x} 是不是一個
dict」，回傳 \texttt{True} 或 \texttt{False}，\texttt{isinstance(x, list)} 同理。
\index{isinstance()}

這裡先擋一個可能的困惑。第 \ref{ch:files-exceptions} 章第 \ref{sec:lbyl-vs-eafp} 節說
Python 偏好「先做，錯了再處理」，不要一路寫 \texttt{if} 檢查；這一節卻在用
\texttt{isinstance()} 先檢查型別再決定怎麼做，看起來剛好相反。這兩件事不衝突——那一節
講的是\textbf{本來就可能因為外部因素失敗}的操作（開檔案、連網路、轉換使用者輸入），
那種時候先檢查沒有意義，因為狀態隨時會變。而「這個值是 dict 還是別的」是完全可預期、
跟外界無關的邏輯分岔，本來就該用 \texttt{if}，跟「這個數是不是偶數」是同一類的判斷。

\subsection*{把每一個值連同它的位置印出來}

任務是：不管資料有幾層，把最底層的每一個值都印出來，而且順便告訴你它是從哪裡挖到的。

\begin{example}{python}{(walk-nested-json.py) 印出巢狀資料裡的每一個值}
import json

text = '''
{
    "class": "8A",
    "teacher": {
        "name": "Wang",
        "contact": {"phone": "0912345678", "email": "wang@example.com"}
    },
    "students": [
        {"name": "Alice", "scores": {"math": 92, "english": 78}},
        {"name": "Bob", "clubs": ["basketball", "coding"]}
    ]
}
'''

data = json.loads(text)


def show_all(value, path):
    if isinstance(value, dict):
        for key in value:
            show_all(value[key], path + "." + key)
    elif isinstance(value, list):
        for index, item in enumerate(value):
            show_all(item, path + "[" + str(index) + "]")
    else:
        print(path + ": " + str(value))


show_all(data, "data")
\end{example}

\iosample{}{data.class: 8A \\ data.teacher.name: Wang \\ data.teacher.contact.phone: 0912345678 \\
data.teacher.contact.email: wang@example.com \\ data.students[0].name: Alice \\
data.students[0].scores.math: 92 \\ data.students[0].scores.english: 78 \\
data.students[1].name: Bob \\ data.students[1].clubs[0]: basketball \\
data.students[1].clubs[1]: coding \\}

這十行輸出，直接就是那份資料的地圖。左邊那串路徑還可以照著翻成你在第 \ref{ch:csv} 章
寫慣的中括號形式，例如最後一行：
\begin{verbatim}
data.students[1].clubs[1]            <- 程式印出來的路徑
data["students"][1]["clubs"][1]      <- 對應的取值寫法
\end{verbatim}
也就是說，這支程式順便替你把「這份資料到底能怎麼取值」整份列了出來。

還有一件事值得注意：兩位學生的深度根本不一樣（Alice 有 \texttt{scores}，Bob 有
\texttt{clubs}），而這支程式從頭到尾不在意這件事。

拆開來看，函式只有三個分支：

\begin{itemize}
\item 是 dict：對每一個 key，把對應的值再交給自己處理，路徑後面接上 \texttt{.key}。
\item 是 list：對每一個項目，把項目再交給自己處理，路徑後面接上 \texttt{[索引]}。
\item 都不是：那就是最底層的值了，印出來。
\end{itemize}

\textbf{終止條件在哪裡？} 這是這一節跟前面四節最不一樣的地方，一定要看清楚：這支程式
裡沒有任何一行 \texttt{if} 寫著「停」。終止條件藏在最後那個 \texttt{else} 裡——碰到不是
dict 也不是 list 的東西（一個字串、一個數字），就不再往下呼叫自己，直接印出來收工。

那「每次要更靠近終止條件」呢？這裡變少的不是某個數字，是\textbf{還沒拆開的層數}。每
遞迴一次，處理的就是更裡面一層的東西，而任何一份資料的層數都是有限的，所以一定會走到
底。回到抽屜的比喻：每打開一個盒子，裡面的東西一定比外面少，總會有一個盒子裡沒有盒子。

\subsection*{常見錯誤：想把結果收集起來，卻讓上一次的結果留了下來}

上面的函式是邊走邊印。更常見的需求是「收集起來回傳」，例如把資料裡所有的數字挑出來。
這時候有一種寫法會非常自然地浮現在腦中——用一個預設參數當收集用的袋子：

\begin{example}{python}{(collect-numbers-bug.py) 用可變物件當預設參數的下場}
def collect_numbers(value, found=[]):
    if isinstance(value, dict):
        for key in value:
            collect_numbers(value[key], found)
    elif isinstance(value, list):
        for item in value:
            collect_numbers(item, found)
    elif isinstance(value, int):
        found.append(value)
    return found


print(collect_numbers({"math": 92, "english": 78}))
print(collect_numbers({"math": 65}))
\end{example}

\iosample{}{[92, 78] \\{} [92, 78, 65] \\}

第二次呼叫傳的是一份全新的資料，卻拿回了上一次的 \texttt{92} 跟 \texttt{78}。這正是
第 \ref{ch:functions} 章第 \ref{sec:mutable-default} 節那個購物車陷阱：\texttt{found=[]}
只在 \texttt{def} 那一行執行的當下造了一次 list，之後每次沒自己傳 \texttt{found} 的
呼叫，用的都是同一個物件。

在遞迴裡踩到這個坑，比在購物車那個例子更難查——因為這支函式\textbf{本來就會不停地呼叫
自己並共用同一個袋子}，那是設計上就要的行為。看到結果不對的時候，很容易懷疑是遞迴寫錯
了，其實遞迴那部分完全正確，錯的是袋子從一開始就不是新的。

修法有兩種。一種是照第 \ref{ch:functions} 章的做法把預設值換成 \texttt{None}，函式裡
再判斷。另一種更適合遞迴：\textbf{乾脆不要共用袋子，讓每一層都回傳自己的一份清單，由
上一層合併起來。}

\begin{example}{python}{(collect-numbers-fixed.py) 每一層都給自己一個新袋子}
def collect_numbers(value):
    found = []
    if isinstance(value, dict):
        for key in value:
            found.extend(collect_numbers(value[key]))
    elif isinstance(value, list):
        for item in value:
            found.extend(collect_numbers(item))
    elif isinstance(value, int):
        found.append(value)
    return found


print(collect_numbers({"math": 92, "english": 78}))
print(collect_numbers({"math": 65}))
print(collect_numbers({"a": [1, {"b": 2}], "c": {"d": [3, 4]}}))
\end{example}

\iosample{}{[92, 78] \\{} [65] \\{} [1, 2, 3, 4] \\}

\texttt{found = []} 現在寫在函式主體裡，每一次呼叫都會重新執行一次，所以每一層都有自己
乾淨的一份；\texttt{extend()} 則是第 \ref{ch:lists} 章教過的，把下一層交回來的那份清單
拆開來一項一項加進自己這份裡。這個版本沒有任何跨呼叫的共用狀態，怎麼呼叫都不會互相
污染。

\begin{drills}[label=drills:nested-data-walk]
\item 用 \texttt{show\_all} 走訪下面這份資料，會依序印出哪幾行？
\begin{verbatim}
{"title": "note", "tags": ["work", "urgent"], "meta": {"author": "Amy"}}
\end{verbatim}
（呼叫方式是 \texttt{show\_all(data, "data")}。）
\item 抓蟲：有人把 \texttt{show\_all} 裡處理 list 的那個 \texttt{elif} 整段刪掉了，
只留下 dict 跟 \texttt{else} 兩個分支。用它走訪
\texttt{\{"name": "Bob", "clubs": [\{"title": "coding"\}]\}} 會印出什麼？漏掉了什麼？
\item 改寫：把 \texttt{show\_all} 改成只印出「值是數字」的那幾行，字串的部分跳過不印。
（提示：最後那個 \texttt{else} 裡再加一層判斷。）
\end{drills}

\section{遞迴真正划算的地方（二）：走訪資料夾樹}
\label{sec:folder-tree}

上一節的骨架，這一節原封不動再用一次。差別只有兩個字：把「這個東西是不是容器」換成
「這個東西是不是資料夾」。

第 \ref{ch:modules} 章第 \ref{sec:os-and-sys} 節教過用
\texttt{for name in sorted(os.listdir(folder))} 把一個資料夾裡的東西列出來，也教過
\texttt{os.path.join()} 把資料夾位置跟名字接成完整路徑，還有
\texttt{os.path.isdir()}／\texttt{os.path.isfile()} 用來問「這個路徑是資料夾還是檔案」。
那一節特別提醒過：這兩個函式吃的是完整路徑，不是 \texttt{os.listdir()} 給你的那個名字，
直接把名字丟進去會得到 \texttt{False}，而且不會報錯。

那個列資料夾的程式有個限制，當時沒有說破：\textbf{它只走一層。} 如果資料夾裡面還有
資料夾，它只會告訴你「這裡有一個叫 notes 的資料夾」，不會告訴你 notes 裡面有什麼。
現在你有工具把那一半補上了。

\begin{example}{python}{(walk-folder-tree.py) 把整棵資料夾樹印出來}
import os


def show_tree(folder, depth=0):
    indent = "    " * depth
    for name in sorted(os.listdir(folder)):
        full_path = os.path.join(folder, name)
        if os.path.isdir(full_path):
            print(indent + name + "/")
            show_tree(full_path, depth + 1)
        else:
            print(indent + name)


show_tree("project")
\end{example}

假設 \texttt{project} 資料夾裡有一個 \texttt{README.txt}、一個 \texttt{notes} 資料夾
（裡面有 \texttt{2026-01.txt} 跟一個 \texttt{meeting} 子資料夾）、還有一個
\texttt{photos} 資料夾（裡面有 \texttt{cover.jpg} 跟一個 \texttt{trip} 子資料夾）：

\iosample{}{README.txt \\ notes/ \\ ~~~~2026-01.txt \\ ~~~~meeting/ \\
~~~~~~~~agenda.txt \\ photos/ \\ ~~~~cover.jpg \\ ~~~~trip/ \\ ~~~~~~~~beach.jpg \\}

這支程式你其實早就會寫了，因為它就是上一節那一招。三個分支換成兩個，判斷式換一個函式，
其他完全一樣：用迴圈跑過這一層的每一個東西，是資料夾就遞迴下去，不是就直接處理。
\texttt{depth} 參數也還是第 \ref{sec:recursion-trace} 節那個鷹架——只是這一次它不是
除錯用的臨時程式碼，縮排本身就是輸出的一部分，直接把樹的形狀畫了出來。

順帶一提，\texttt{sorted()} 不是可有可無的裝飾。\texttt{os.listdir()} 交出來的順序不
保證固定，不同的電腦、不同的檔案系統可能給你不一樣的排序；加上 \texttt{sorted()}，你
在自己電腦上看到的結果才是穩定、可重現的。

\subsection*{老實說：標準函式庫已經有這個功能了}

Python 的 \texttt{os} 模組裡有一個叫 \texttt{os.walk()} 的函式，它已經幫你把整棵樹走完
了，實務上多數人要處理整棵資料夾樹的時候，會直接用它，不會自己寫遞迴。

那這一節在幹嘛？不是要取代 \texttt{os.walk()}，是要讓你看懂「走一棵樹」這件事到底在做
什麼。而且更重要的是：下一次你碰到的樹沒有現成工具——例如上一節那種誰也不知道長幾層的
JSON、或是一份「每個人底下可能還有下屬」的組織資料——你自己寫得出來。這一節跟上一節
共用的那個骨架，才是要帶走的東西；資料夾只是它最好認的一個例子。

\subsection*{只看，不動手}

這一節的程式從頭到尾只印路徑，沒有任何刪除、搬移、改名的動作，這是刻意的。走訪整棵樹的
程式一旦寫錯——例如遞迴的方向錯了、判斷式寫反了——它會錯得又快又廣。動手改動檔案之前，
先讓程式把「它打算動哪些檔案」印出來給你看一遍，確認清單沒問題再說。

還有一件事這一節沒做：真實世界的資料夾裡可能有你沒有權限打開的角落，
\texttt{os.listdir()} 碰到那種資料夾會丟出例外，整支程式就停在那裡。那是第
\ref{ch:files-exceptions} 章 \texttt{try}／\texttt{except} 的工作，這裡不展開，但你要
知道拿去跑真實資料之前有這一關。

\begin{drills}[label=drills:folder-tree]
\item 有一個 \texttt{photos} 資料夾，裡面是：一個 \texttt{2025} 資料夾（裡面有
\texttt{note.txt} 跟一個 \texttt{summer} 資料夾，\texttt{summer} 裡面只有
\texttt{sea.jpg}），以及一個檔案 \texttt{cover.png}。執行
\texttt{show\_tree("photos")} 會依序印出哪幾行？縮排各幾格？
\item 抓蟲：有人把 \texttt{os.path.isdir(full\_path)} 寫成
\texttt{os.path.isdir(name)}，遞迴呼叫也從 \texttt{show\_tree(full\_path, depth + 1)}
改成 \texttt{show\_tree(name, depth + 1)}。這支程式會報錯嗎？輸出會變成什麼樣子？
\item 改寫：把 \texttt{show\_tree} 改成只印出 \texttt{.txt} 檔案（資料夾的名字還是要
印，不然看不出結構）。
\end{drills}

\section{一瞥：\texttt{functools.lru\_cache}}
\label{sec:lru-cache}

回頭填第 \ref{sec:fibonacci-cost} 節挖的坑。當時的結論是：遞迴版的費氏數列慢，是因為
同樣的子問題被重複算了幾百萬次。既然如此，除了改寫成迴圈，還有另一條路——讓函式記得
自己算過的答案。

標準函式庫已經替你寫好了。從 \texttt{functools} 模組匯入 \texttt{lru\_cache}，在函式
上面加一行：

\begin{example}{python}{(fibonacci-cached.py) 加一行，讓函式記得算過的答案}
from functools import lru_cache

call_count = 0


@lru_cache
def fibonacci(n):
    global call_count
    call_count = call_count + 1
    if n <= 1:
        return n
    return fibonacci(n - 1) + fibonacci(n - 2)


answer = fibonacci(30)
print(f"fibonacci(30) = {answer}，一共呼叫了 {call_count} 次")
\end{example}

\iosample{}{fibonacci(30) = 832040，一共呼叫了 31 次 \\}

函式主體一個字都沒改，呼叫次數從兩百六十九萬次掉到 31 次。
\index{functools.lru\_cache}

\texttt{lru\_cache} 做的事一句話講得完：\textbf{它幫函式記住「這個參數算過的答案」，
下次再拿同樣的參數來問，直接把答案拿出來，不再重算一遍。} 費氏數列之所以受益這麼大，
正是因為它重複問的問題特別多。

那個 \texttt{@} 符號你在第 \ref{ch:advanced-functions} 章第 \ref{sec:decorator} 節看過，
就是那個「把函式包一層」的寫法。差別只在這次的裝飾器不是你自己寫的，是標準函式庫替你
準備好的，直接拿來用就好。

\subsection*{老實說：這沒有讓費氏該用遞迴變成對的}

\texttt{lru\_cache} 讓一支慢的遞迴變快了，但第 \ref{sec:fibonacci-cost} 節那個迴圈版
本來就又快又短，加了快取的遞迴版並沒有比它好。這一行不是「遞迴的萬用解藥」，加上去
也不會讓一個本來就不該用遞迴的問題變得該用遞迴。

它真正的用處是在另一種情境：\textbf{這個函式本來就該用遞迴，而且真的會重複算到同樣的
東西}。那時候加這一行，才是拿它來解決問題，不是拿它來遮蓋問題。

這一節到此為止，只到「知道有這個東西、看到別人用的時候看得懂那一行在做什麼」的程度。
它還有一些可以調的選項，等你哪天真的需要了再去查——第 \ref{ch:modules} 章那套「打開
文件、看簽名、看說明」的流程對它一樣管用。

\section{所以，什麼時候該用遞迴？}
\label{sec:recursion-or-loop}

這一章走到這裡，五個遞迴函式都寫過了，該把判準收起來。判準只有兩句話：

\begin{itemize}
\item \textbf{資料是平的、次數是知道的} $\to$ 用迴圈。
\item \textbf{東西裡面還有同一種東西，而且不知道有幾層} $\to$ 用遞迴。
\end{itemize}

拿這兩句話把這一章的五個例子重新排一次：

\begin{center}
\begin{tabular}{@{}lll@{}}
\toprule
例子 & 該用遞迴嗎 & 為什麼 \\
\midrule
\texttt{countdown} & 不必 & 要跑幾次是知道的，\texttt{for} 就夠了 \\
階乘 & 不必 & 同上；它的價值是形狀清楚，適合拿來看懂遞迴 \\
費氏數列 & 不該 & 同上，而且遞迴版還會把同一件事重算幾百萬次 \\
巢狀 JSON & 該用 & 不知道有幾層，用迴圈得先知道要寫幾層 \texttt{for} \\
資料夾樹 & 該用 & 同上（不過實務上有現成的 \texttt{os.walk()}） \\
\bottomrule
\end{tabular}
\end{center}

這張表最值得注意的地方是：前三個「不必／不該」的例子，正好是所有教遞迴的地方最常用的
入門例子。它們不是壞例子，它們是好的\textbf{教材}，但不是好的\textbf{範本}。看懂形狀
之後，請把它們留在教室裡。

除了上面那條判準，還有三件實務上的事值得記住：

\begin{itemize}
\item \textbf{很深的資料不適合遞迴。} 每疊一層都要花資源，而且 Python 有深度上限
（第 \ref{sec:recursion-trace} 節提過）。如果你的資料可能有幾萬層，遞迴會直接撞牆，
那種時候得換寫法。
\item \textbf{遞迴比較難用 \texttt{print} 追蹤}，因為同一行 \texttt{print} 會被跑很多
次、而且分不出是哪一層印的。這就是第 \ref{sec:recursion-trace} 節那個縮排鷹架值得學
起來的原因——它是遞迴專用的除錯手法。
\item \textbf{寫不出來就不要硬寫。} 一個你自己都追蹤不完的遞迴函式，三個月後的你也讀
不懂，而且到時候要改它的人多半就是你。能用迴圈清楚表達的事情，就用迴圈，沒有人會覺得
你不夠厲害。
\end{itemize}

\stylepoint{Python 社群對遞迴的態度偏保守：能用迴圈或現成工具寫清楚的事，多數人不會
特地寫成遞迴——這跟某些語言把遞迴當成預設寫法的文化不一樣。真正該用遞迴的訊號其實只有
一個：\textbf{你要處理的東西，裡面還有一模一樣的東西，而且你不知道有幾層。} 看到這個
訊號就放心用，沒看到就先寫迴圈。}

這一章的主題說到底是一種節制：明明有一個很聰明的工具，卻要判斷什麼時候不該拿出來用。
下一章談的是同一種節制的另一面——什麼樣的程式碼叫「寫得像 Python」，以及什麼時候
「太 Pythonic」反而讓人讀不懂（第 \ref{ch:pythonic} 章）。

\begin{drills}[label=drills:recursion-or-loop]
\item 你要把一份成績 list 全部加總起來，例如 \texttt{[88, 76, 92, 65]}。用遞迴還是
迴圈？為什麼？
\item 你有一份公司的組織資料，每個人底下可能還有下屬，下屬底下可能又還有下屬，沒有人
說得準最多幾層。你要把所有人的名字列出來。用遞迴還是迴圈？為什麼？
\item 你要印出 1 到 100。用遞迴還是迴圈？如果硬要用遞迴，會有什麼問題？
\end{drills}

\section*{練習解答}
\addcontentsline{toc}{section}{練習解答}

\paragraph{第~\ref{sec:recursion-basics}~節}
\begin{enumerate}
\item 只印出 \texttt{Liftoff!}：\texttt{n} 一開始就是 \texttt{0}，第一次呼叫就滿足終止條件
\texttt{n <= 0}。
\item 依序印出 \texttt{1}、\texttt{2}、\texttt{3}、\texttt{4}：每次呼叫 \texttt{n}
加 1，直到 \texttt{n} 超過 \texttt{limit}（也就是 4）才停止。
\item 可以，例如：
\begin{verbatim}
def count_up(n, limit):
    for i in range(n, limit + 1):
        print(i)
\end{verbatim}
用 \texttt{range(n, limit + 1)} 就能取代遞迴，一路數到 \texttt{limit}（含）為止。
\end{enumerate}

\paragraph{第~\ref{sec:recursion-trace}~節}
\begin{enumerate}
\item 少一層，形狀一樣：\texttt{進入 countdown(2)} $\to$ \texttt{2} $\to$ 縮四格的
\texttt{進入 countdown(1)} $\to$ \texttt{1} $\to$ 縮八格的 \texttt{進入 countdown(0)}、
\texttt{Liftoff!}、\texttt{離開 countdown(0)} $\to$ 縮四格的 \texttt{離開 countdown(1)}
$\to$ 不縮排的 \texttt{離開 countdown(2)}。縮排最深的是
\texttt{countdown(0)} 那三行（八格）。
\item 因為 \texttt{n} 從 10 開始每次減 2，走的是 10、8、6、4、2、0、-2、……全部是偶數，
永遠不會剛好等於 \texttt{1}，終止條件寫在那裡卻碰不到。這正是「有終止條件」跟「保證走
得到終止條件」的差別。把 \texttt{if n == 1} 改成 \texttt{if n <= 1} 就會停。從
traceback 看：\texttt{[Previous line repeated 996 more times]} 加上前面明白列出來的那
幾行，大約疊了一千層——那就是 Python 的深度上限。
\item 會少掉「一層一層收回來」的那一半，只剩往右縮排一路變深的前半段。看不出每一層是
什麼時候結束的，也看不出「最先開始的那一層最後才結束」這件事——而那正是遞迴最容易誤解
的地方。
\end{enumerate}

\paragraph{第~\ref{sec:recursion-return}~節}
\begin{enumerate}
\item 四層疊上去（\texttt{factorial(3)}、\texttt{factorial(2)}、\texttt{factorial(1)}、
\texttt{factorial(0)}），再四層退回來，退回來的數字依序是 \texttt{1}、\texttt{1}、
\texttt{2}、\texttt{6}，最後 \texttt{print} 再印一次 \texttt{6}。
\item \texttt{factorial(4)} 還是 24：4、3、2 一路乘下去，走到 \texttt{n == 1} 就停，
少乘的那個 1 不影響答案。但 \texttt{factorial(0)} 會出事——\texttt{0} 不等於 \texttt{1}，
於是往下走 \texttt{-1}、\texttt{-2}、\texttt{-3}……永遠碰不到 \texttt{1}，最後丟出
\texttt{RecursionError}。跟 \texttt{factorial(-1)} 是同一個病。
\item 結果是 \texttt{14}。迴圈版例如：
\begin{verbatim}
def total(numbers):
    result = 0
    for number in numbers:
        result = result + number
    return result
\end{verbatim}
（其實直接寫 \texttt{sum(numbers)} 就好，這一題只是要你看出遞迴版在做的就是這件事。）
\end{enumerate}

\paragraph{第~\ref{sec:fibonacci-cost}~節}
\begin{enumerate}
\item 15 次。\texttt{fibonacci(5)} 岔成 \texttt{fibonacci(4)} 跟 \texttt{fibonacci(3)}，
兩邊各自再岔下去，數到最後一共 15 次呼叫，其中 \texttt{fibonacci(1)} 被算了 5 次、
\texttt{fibonacci(0)} 被算了 3 次。
\item 因為 \texttt{countdown} 每次\textbf{只呼叫自己一次}，一路往下就是一條線，
\texttt{countdown(30)} 就是 31 次呼叫。\texttt{fibonacci} 每次呼叫自己\textbf{兩次}，
一條線變成一棵不斷分岔的樹，重複的工作用乘的長大。差別就在自我呼叫的次數。
\item 對照 \texttt{fibonacci-loop.py}。最容易寫錯的是頭兩項：第 0 項要是 \texttt{0}、
第 1 項要是 \texttt{1}。如果你的版本迴圈跑 \texttt{range(n)} 之後回傳
\texttt{current} 而不是 \texttt{previous}，整串會往前偏一格。
\end{enumerate}

\paragraph{第~\ref{sec:nested-data-walk}~節}
\begin{enumerate}
\item 依序四行：
\begin{verbatim}
data.title: note
data.tags[0]: work
data.tags[1]: urgent
data.meta.author: Amy
\end{verbatim}
\item 不會報錯，會印出兩行：\texttt{data.name: Bob} 跟
\texttt{data.clubs: [\{'title': 'coding'\}]}。因為少了 list 的分支，那份社團清單被當成
「最底層的值」整個印出來了，裡面那個 dict 根本沒被拆開——這是最典型的「安靜的錯」，
程式跑完了、看起來也有輸出，但少走了一整片資料。
\item 把最後的 \texttt{else} 改成：
\begin{verbatim}
    else:
        if isinstance(value, int):
            print(path + ": " + str(value))
\end{verbatim}
（也可以直接寫成 \texttt{elif isinstance(value, int):}，效果一樣。）走訪課文那份資料的
話，十行輸出只會剩下 Alice 那兩個分數：
\begin{verbatim}
data.students[0].scores.math: 92
data.students[0].scores.english: 78
\end{verbatim}
\end{enumerate}

\paragraph{第~\ref{sec:folder-tree}~節}
\begin{enumerate}
\item 依序是：\texttt{2025/}（不縮排）、\texttt{note.txt}（縮四格）、
\texttt{summer/}（縮四格）、\texttt{sea.jpg}（縮八格）、\texttt{cover.png}（不縮排）。
\texttt{2025} 排在 \texttt{cover.png} 前面，是因為 \texttt{sorted()} 按字串排序，
數字開頭的名字排在字母前面。
\item 不會報錯，這才是麻煩的地方。\texttt{os.path.isdir("notes")} 問的是「目前所在的
資料夾底下有沒有一個叫 notes 的資料夾」，而 \texttt{notes} 其實在 \texttt{project}
裡面，所以得到 \texttt{False}——每一個子資料夾都會被當成檔案印出來，程式一層都沒有
遞迴下去，最後只印出最上層那三個名字。這就是第 \ref{ch:modules} 章提醒過的：
\texttt{os.path.isdir()} 吃的是完整路徑，不是 \texttt{os.listdir()} 給你的名字。
\item 把 \texttt{else} 那個分支加上副檔名判斷：
\begin{verbatim}
        elif name.endswith(".txt"):
            print(indent + name)
\end{verbatim}
（也就是把原本的 \texttt{else:} 換成這一行。）資料夾的分支不動，樹的結構才留得住；
沒有 \texttt{.txt} 的資料夾會印出名字但底下是空的。
\end{enumerate}

\paragraph{第~\ref{sec:recursion-or-loop}~節}
\begin{enumerate}
\item 迴圈。資料是平的（一層 list，裡面就是數字），要跑幾次也很明確（就是清單的長度）。
一個 \texttt{for} 迴圈，或直接 \texttt{sum()}。用遞迴不但沒有好處，還多了撞到深度上限
的風險。
\item 遞迴。「每個人底下可能還有下屬」就是「東西裡面還有同一種東西」，而且沒有人說得
準最多幾層——這正是遞迴的招牌情境。用迴圈的話，你得事先知道要寫幾層巢狀 \texttt{for}，
而這件事根本不知道。骨架就是第 \ref{sec:nested-data-walk} 節那個片語。
\item 迴圈，\texttt{for i in range(1, 101)}。次數是知道的，資料是平的。硬要用遞迴的話
它會疊 100 層，這次還撐得住，但同樣的寫法拿去印 1 到 10000 就會丟
\texttt{RecursionError}——而迴圈版改個數字就沒事。
\end{enumerate}

\chapter{Pythonic 慣用寫法與風格}
\label{ch:pythonic}

你大概讀過那種句子：每個字都認得，文法也挑不出毛病，但讀起來就是怪，一看就知道是從別的
語言硬翻過來的。反過來也一樣——聽一個學中文沒多久的朋友講話，每個字都用對了，組起來就是
不像我們平常講話的樣子。他沒有講錯，只是帶著口音。

程式碼也有口音。從第 \ref{ch:variables} 章到第 \ref{ch:recursion} 章，語法該教的都教完
了：你會寫變數、條件、迴圈、函式、檔案、例外、comprehension、generator、遞迴。同一件事，
用 Python 寫得出來的方法常常不只一種，而且它們\textbf{全都能動}——這一章要談的不是哪一種寫法
對、哪一種錯，是\textbf{讀的人要不要多花力氣}。

上一章的結尾，還有第 \ref{ch:iterators-generators} 章的結尾，都留了同一句話：「第
\ref{ch:pythonic} 章會把『寫得像 Python』這件事整個收起來。」就是現在。

先把一件事講在前面，免得整章讀起來像在推銷：這一章的最後一節會反過來談「什麼時候不要太
Pythonic」。腔調本身不是目的，把話講清楚才是。

這一章分三段推進：

\begin{itemize}
\item 第一段是低風險的日常慣例：程式碼該長什麼樣子、條件怎麼寫成一句話、值怎麼搬運。
這三節的共同論證只有一條——程式碼被讀的次數，遠比被寫的次數多
（第 \ref{sec:pep8}～\ref{sec:value-idioms} 節）。
\item 第二段是這一章存在的理由。名字牌與 EAFP 這兩個從第 \ref{ch:variables} 章跟第
\ref{ch:files-exceptions} 章一路跟著你走到這裡的想法，會在這裡走完一圈：從「解釋為什麼
會發生怪事的圖像」，變成「你打字的當下會用到的判準」。中間補上 \texttt{with} 這塊一直
缺著的拼圖（第 \ref{sec:is-vs-equal}～\ref{sec:eafp-closure} 節）。
\item 第三段是煞車。學會一個寫法之後，還要學會什麼時候不用它
（第 \ref{sec:not-too-pythonic} 節）。
\end{itemize}

\section{大家講好的樣子：PEP 8 與命名慣例}
\label{sec:pep8}

先看一段程式。它能跑，答案也對。

\begin{example}{python}{(tidy-todo-before.py) 能動，但看得很累}
maxitems=5
l=["Milk","Eggs","Toast"]
def addItem(item,items):
    if len(items)<maxitems:
        items.append(item)
    return items
print(addItem("Apple",l))
\end{example}

\iosample{}{['Milk', 'Eggs', 'Toast', 'Apple'] \\}

沒有一行寫錯。但你讀它的時候，眼睛要多做好幾件事：\texttt{maxitems=5} 這一串黏在一起，
你得自己在心裡把等號兩邊切開；\texttt{l} 是什麼？是 list 嗎？還是 1？\texttt{addItem}
跟 \texttt{maxitems} 一個用大寫駝峰、一個全小寫，看不出規律；\texttt{def} 前面沒有空行，
整段擠成一坨。

同樣的程式，換一種排法：

\begin{example}{python}{(tidy-todo-after.py) 一樣的程式，照慣例排過}
MAX_ITEMS = 5
shopping_list = ["Milk", "Eggs", "Toast"]


def add_item(item, items):
    if len(items) < MAX_ITEMS:
        items.append(item)
    return items


print(add_item("Apple", shopping_list))
\end{example}

\iosample{}{['Milk', 'Eggs', 'Toast', 'Apple'] \\}

輸出一模一樣。Python 完全不在乎你選哪一種——但下一個讀它的人在乎，而下一個讀它的人，
第一個是三天後的你自己。第 \ref{ch:variables} 章第 \ref{sec:naming-practices} 節就提醒過
一次：幾天後回頭看，你會完全忘記某個短名字到底是什麼。這一節要把那句話擴大到整份檔案的
外觀上。

\subsection*{PEP 8 是什麼}

Python 社群對「程式碼該長什麼樣子」有一份講好的文件，叫做 \textbf{PEP 8}。
\index{PEP 8}
PEP 是 Python Enhancement Proposal 的縮寫，直譯是「Python 改進提案」——任何人想對 Python
提一件事，都可以寫成一份 PEP 給整個社群看，通過的就變成大家的共識。PEP 8 是最老的幾份
之一，它談的不是語言功能，而是排版與命名。

有一件事要立刻講清楚，因為這一章最後一節整節都建立在它上面：

\keypoint{\textbf{PEP 8 不是語法規則。} Python 不會因為你不照它做就報錯，也不會跑得比較
慢。它是一群人為了讓彼此的程式碼看起來像同一個世界的東西，自願講好的約定。約定跟規則
不一樣的地方在於：規則由編譯器執行，約定只能靠人守。}

\subsection*{排版：四條真的有差的}

PEP 8 很長，但對現在的你真正有差的只有四條，而且每一條都有理由，不是「規定就是這樣」。

\begin{itemize}
\item \textbf{縮排用 4 個空白，而且不要跟 tab 混用。} 4 個空白是視覺上剛好看得出層次、
又還留得下空間給多層縮排的寬度。至於 tab 跟空白混用會發生什麼事，第 \ref{ch:introduction}
章第 \ref{sec:editor-pitfalls} 節示範過那個錯誤訊息了——那次的教訓在這裡變成一條慣例。
\item \textbf{一行不要太長。} PEP 8 給的數字是 79 個字元。這個數字有歷史原因，但今天更
實際的理由是：真正在工作的人常常同時開好幾個檔案並排在同一個螢幕上，一行太長，並排就
放不下了。
\item \textbf{用空行做視覺分組，但不要過量。} 兩個函式之間空兩行、函式裡面的邏輯段落之間
空一行，讀的人一眼就知道哪幾行是一組。全部黏在一起看不出結構，每兩行就空一行也一樣看不
出結構。
\item \textbf{運算子前後各留一個空白。} \texttt{x = 5} 讀起來比 \texttt{x=5} 舒服；
式子一長，差別更明顯：\texttt{len(items) < MAX\_ITEMS} 一眼看得出是在比大小，
擠成 \texttt{len(items)<MAX\_ITEMS} 就得自己找斷點在哪。理由跟中文標點一樣，
空白是在幫讀的人斷句。
\end{itemize}

\subsection*{命名：四條}

\begin{itemize}
\item \textbf{變數與函式用 \texttt{snake\_case}}：全部小寫，單字之間用底線連接。
\texttt{shopping\_list}、\texttt{add\_item} 都是這種寫法。第 \ref{ch:functions} 章教
函式命名時已經提過這個詞，這裡只是把它一起放進慣例清單。
\item \textbf{常數用全大寫加底線}，例如 \texttt{MAX\_ITEMS}、\texttt{CONFIG\_PATH}。
\index{全大寫命名慣例}
這裡要特別說清楚一件事：\textbf{Python 沒有「常數」這種東西。} 你寫成
\texttt{MAX\_ITEMS = 5}，之後照樣可以寫 \texttt{MAX\_ITEMS = 99}，Python 不會攔你。
全大寫只是一個講好的暗號，意思是「這個值是設定，不打算在程式跑的過程中被改動」。它防
不了任何事，它只是在跟讀的人講話。
\item \textbf{class 用大寫開頭的駝峰式}，例如 \texttt{ShoppingCart}。class 是第
\ref{ch:objects} 章的主題，這裡先知道有這條慣例就好。
\item \textbf{布林值用 \texttt{is\_}／\texttt{has\_} 開頭}，例如
\texttt{is\_passing}、\texttt{has\_ticket}。這一條第 \ref{ch:variables} 章第
\ref{sec:naming-practices} 節已經完整講過，不重講。
\end{itemize}

\stylepoint{這些命名慣例本身沒有一條是神聖的——真正有價值的是\textbf{整份專案裡保持
一致}。一份檔案裡一半 \texttt{snake\_case}、一半駝峰式，比全部用駝峰式還難讀，因為讀的人
會一直懷疑這個差別是不是有意義的。}

\subsection*{有工具會替你檢查這件事}

風格慣例不需要用背的。有人寫了程式來幫你抓這類問題：\texttt{flake8} 會把不合慣例的地方
一條一條列給你看，\texttt{black} 更直接，它會直接把你的檔案照慣例重排一遍。這種工具的來源
跟第 \ref{ch:data-tools} 章裝過的那些套件一樣，用 \texttt{pip} 安裝。這裡點名它們，只是
要你知道「不必用背的」，實際怎麼裝、怎麼設定，等你真的需要的時候再查就好。

\subsection*{但是}

給完規則之後，馬上要補一句「但是」——這件事很重要，重要到 PEP 8 自己就有一節在講。那一節
的標題叫「A Foolish Consistency Is the Hobgoblin of Little Minds」（愚蠢的一致性是小心眼
的妖怪），裡面寫著：

\begin{quote}
與風格指南保持一致很重要。但最重要的是：知道什麼時候該不一致——有時候風格指南就是不適用。
有疑慮的時候，用你自己最好的判斷。
\end{quote}

也就是說，這份講好程式碼該長什麼樣子的文件，自己就留了一道門，讓你在有理由的時候走出去。
什麼樣的理由算數，第 \ref{sec:not-too-pythonic} 節會整節談。

最後補一句你現在用不到、但很快會用到的事：一群人一起改同一個專案的時候，會有一道叫
code review 的關卡——你的程式碼要先被別人讀過才進得去。風格一致在那裡不是美學問題，是
直接變成別人讀你程式碼要花多少力氣。那套流程長什麼樣子，第 \ref{ch:industry} 章會講。

\begin{drills}[label=drills:pep8]
\item 把下面這段程式照這一節的慣例改寫（命名與排版都要改）：
\begin{verbatim}
TotalItems=0
def CountThem(TheList):
    global TotalItems
    for x in TheList:
        TotalItems=TotalItems+1
    return TotalItems
\end{verbatim}
\item 你寫了 \texttt{MAX\_RETRY = 3}，後來在程式的另一個地方寫了
\texttt{MAX\_RETRY = 10}。Python 會阻止你嗎？如果不會，那全大寫這個慣例到底在幫誰？
\item PEP 8 說一行不要超過 79 個字元，但你的團隊決定改用 99 個字元當上限，而且整個專案
都照這個數字。這樣算不算「錯」？
\end{drills}

\section{把條件寫成一句話：\texttt{not}、\texttt{in}、連續比較}
\label{sec:condition-idioms}

上一節談的是程式碼的外觀，這一節往裡面走一層：同一個判斷，怎麼寫會讓下一個人一眼看懂。
這一節有三個小技巧，其中兩個你手上已經有材料了，只有最後一個是全新的語法。

\subsection*{\texttt{if not items:}}

第 \ref{ch:conditionals} 章第 \ref{sec:truthy-falsy} 節教過 Python 的一個脾氣：空的東西
在 \texttt{if} 裡會被當成 \texttt{False}。那時候教的是\textbf{正面}的形式——
\texttt{if name:} 就是「name 有內容的話」。第 \ref{ch:files-exceptions} 章第
\ref{sec:eafp-recap} 節也用容器示範過一次。

現在把否定的形式補上，順便第一次把完整的清單集中列出來。下面這些值放進 \texttt{if} 裡，
一律被當成 \texttt{False}：

\begin{center}
\begin{tabular}{@{}ll@{}}
\toprule
值 & 是什麼 \\
\midrule
\texttt{0}、\texttt{0.0} & 數字的零 \\
\texttt{""} & 空字串 \\
\texttt{[]} & 空 list \\
\texttt{()} & 空 tuple \\
\texttt{\{\}} & 空 dict \\
\texttt{set()} & 空 set \\
\texttt{None} & 「沒有值」這個值 \\
\bottomrule
\end{tabular}
\end{center}

除了這些，其他一律當成 \texttt{True}。有了這張表，「這個清單是不是空的」就有三種寫法，
而且三種都能動：

\begin{example}{python}{(empty-check-three-ways.py) 三種問法，同一個答案}
items = []

if len(items) == 0:
    print("A: nothing to do")

if items == []:
    print("B: nothing to do")

if not items:
    print("C: nothing to do")
\end{example}

\iosample{}{A: nothing to do \\ B: nothing to do \\ C: nothing to do \\}

三行都印出來了，所以問題不是對錯，是哪一種讀起來比較直接。

A 版繞了一圈：先算長度、再拿長度跟 0 比。你要問的其實不是「長度是幾」，是「有沒有東西」。
B 版更繞，它造了一個空 list 出來，只為了拿去跟另一個 list 比。而且 B 版換一種容器就得
重寫——換成字串要改成 \texttt{== ""}，換成 dict 要改成 \texttt{== \{\}}。

C 版直接問「這裡面沒東西吧？」，而且它對上面那張表裡的每一種值都通用：空 list、空字串、
空 dict、空 set，寫法完全一樣。

\stylepoint{要問一個容器是不是空的，寫 \texttt{if not items:}，不要寫
\texttt{if len(items) == 0:} 或 \texttt{if items == []:}。三種都能動，但第一種不用繞道
算長度、也不用造一個空容器出來比對，而且換成任何一種容器都不必改寫。正面的形式同理：
\texttt{if items:} 就是「裡面有東西的話」。}

\subsection*{一個誠實的例外}

\texttt{if not items:} 有一個它問不出來的問題：它分不出「空的」跟「根本沒有」。

\begin{example}{python}{(empty-or-missing.py) 兩種完全不同的情況，同一個答案}
def report(items):
    if not items:
        print("nothing to do")
    else:
        print(len(items), "items")


report([])
report(None)
\end{example}

\iosample{}{nothing to do \\ nothing to do \\}

兩次都印一樣的話。但這兩件事在現實裡差很多：\texttt{[]} 是「我確實去查了，結果是一份空
清單」，\texttt{None} 是「根本沒有這份清單」——可能是檔案沒讀到，可能是函式忘了回傳東西。
一個是「今天沒有待辦事項」，一個是「待辦清單不見了」。

要分辨這兩件事，得換一個問法，而那個問法就是下一節整節在講的東西。這裡先記住這個缺口。

\subsection*{\texttt{in} 取代一長串 \texttt{or}}

第二個技巧更小。你想問「今天是不是週末」：

\begin{example}{python}{(weekend-check.py) 兩種問法}
day = "Sun"

if day == "Sat" or day == "Sun":
    print("Weekend (or version)")

if day in ("Sat", "Sun"):
    print("Weekend (in version)")
\end{example}

\iosample{}{Weekend (or version) \\ Weekend (in version) \\}

上面那一行，\texttt{day} 這個名字出現了兩次，等號出現了兩次；選項如果有五個，就會出現
五次。下面那一行，\texttt{day} 只出現一次，而且「這幾個選項」被放在一起、看起來就像一份
清單。

\texttt{in} 你早就會了——第 \ref{ch:loops} 章第 \ref{sec:break-search} 節說過「\texttt{in}
其實就是幫你把搜尋迴圈包好了」。這裡只是把同一個工具用在另一個地方。順帶一提，如果選項
很多、而且程式要反覆檢查很多次，第 \ref{ch:dict-set} 章第 \ref{sec:set-membership} 節
講過改用 set 會快很多，寫法一樣是 \texttt{in}。

\subsection*{連續比較}

這是這一章唯一一個全新的語法。你想確認一個分數在 0 到 100 之間：

\begin{example}{python}{(score-in-range.py) 夾在中間的寫法}
score = 87

if score >= 0 and score <= 100:
    print("valid (and version)")

if 0 <= score <= 100:
    print("valid (chained version)")
\end{example}

\iosample{}{valid (and version) \\ valid (chained version) \\}

Python 允許把比較運算子串起來寫，這叫\textbf{連續比較}。
\index{連續比較}
\texttt{0 <= score <= 100} 的意思就是「score 大於等於 0，而且 score 小於等於 100」，
跟上面那一行完全等價——差別是它讀起來就是數學課上寫的那個樣子，而且 \texttt{score} 這個
名字只出現一次。名字只出現一次是有實際好處的：如果你要把 \texttt{score} 改成
\texttt{final\_score}，\texttt{and} 版本有兩個地方要改，改漏一個就是一個安靜的錯。

這個寫法有一條界線要守：\textbf{只有在「同一個東西被夾在中間」的時候才用它}。

\begin{code}{python}{連續比較該用與不該用的形狀}
0 <= score <= 100          # 好：score 被夾在 0 跟 100 中間
lower <= value < upper     # 好：同一個 value 被夾住

a < b == c                 # 不要：三個不同的東西串在一起，讀的人得停下來想
\end{code}

最後一行是合法的 Python，但沒有人這樣寫。連續比較的價值來自它剛好對應到「夾在中間」這個
大家都懂的形狀，一旦串的是三個不相干的比較，那個形狀就沒了，讀的人反而要多想一次。

這一節的三個技巧看起來各做各的，其實它們在做同一件事：\textbf{把一句判斷，還原成一句
人話}。「這裡沒東西」「今天是週末之一」「分數在 0 到 100 之間」——這三句話你講給人聽的
時候本來就是這樣講的，程式碼只是照著寫下來。第 \ref{ch:conditionals} 章第
\ref{sec:readable-conditions} 節那條線（把複雜的條件拆成有名字的布林變數）講的也是同一
件事，只是切入的角度不同。

\begin{drills}[label=drills:condition-idioms]
\item 把下面這段改寫成用 \texttt{in}：
\begin{verbatim}
if ext == ".jpg" or ext == ".png" or ext == ".gif":
    print("image")
\end{verbatim}
\item 把下面這段改寫成連續比較：
\begin{verbatim}
if temperature >= 18 and temperature <= 26:
    print("comfortable")
\end{verbatim}
\item 下面四行分別會印出 \texttt{True} 還是 \texttt{False}？
\begin{verbatim}
items = []
print(not items)
print(items is None)
items = None
print(not items)
print(items is None)
\end{verbatim}
\item 什麼情況下 \texttt{if not items:} 不夠用？舉一個具體的例子。
\end{drills}

\section{值的搬運：unpacking 與 f-string}
\label{sec:value-idioms}

這一節不教任何新語法。一行都不教。

第 \ref{ch:tuples} 章第 \ref{sec:tuple-unpacking} 節已經把 unpacking 完整教完了——基本
寫法、\texttt{*rest}、用底線忽略不要的值、\texttt{a, b = b, a} 交換兩個變數，連風格要點
都給過了。第 \ref{ch:strings} 章三節也把 f-string 教完了。你手上什麼都有。

那這一節在做什麼？書上教過\textbf{怎麼寫}，沒教過\textbf{什麼時候該想起它}。這是兩件不同
的事：學會一個工具跟看到情況會反射性拿出那個工具，中間差了一段路。這一節就是那段路——
一張「看到左邊，想到右邊」的對照表。

\begin{center}
\begin{tabular}{@{}p{5.4cm}p{5.4cm}p{3.8cm}@{}}
\toprule
看到自己寫成這樣 & 該想到 & 教在哪 \\
\midrule
\texttt{temp = a} 之後 \texttt{a = b}、\texttt{b = temp} &
\texttt{a, b = b, a} &
第 \ref{ch:tuples} 章第 \ref{sec:tuple-unpacking} 節 \\
\addlinespace
\texttt{for i in range(len(items)):} 再用 \texttt{items[i]} &
\texttt{for index, item in enumerate(items):} &
第 \ref{ch:loops} 章第 \ref{sec:enumerate} 節 \\
\addlinespace
\texttt{first = row[0]} 加上 \texttt{rest = row[1:]} &
\texttt{first, *rest = row} &
第 \ref{ch:tuples} 章 \\
\addlinespace
迴圈裡拿了 key，再回頭用 \texttt{d[key]} 取值 &
\texttt{for key, value in d.items():} &
第 \ref{ch:tuples} 章 \\
\addlinespace
用不到的那個值，也認真替它取了名字 &
用底線 \texttt{\_} 接住它 &
第 \ref{ch:tuples} 章 \\
\addlinespace
\texttt{"Hello, " + name + "! You are " + str(age)} &
f-string &
第 \ref{ch:strings} 章第 \ref{sec:f-string-basics} 節 \\
\bottomrule
\end{tabular}
\end{center}

左邊那一欄，每一條都能動，而且每一條都是從別的地方帶來的口音——尤其是
\texttt{for i in range(len(items))}。很多語言的迴圈就是這樣寫的：先要一個編號，再拿編號
去取東西。Python 也允許你這樣寫，但它是全書最典型的一種腔調，因為你真正想要的是「東西
本身」，編號只是取到東西的手段。

把左右兩邊並排看一次：

\begin{code}{python}{帶著口音的版本}
temp = a
a = b
b = temp

for i in range(len(items)):
    print(i, items[i])

first = row[0]
rest = row[1:]

for key in scores:
    print(key, scores[key])
\end{code}

\begin{code}{python}{同樣的四件事，本地話}
a, b = b, a

for index, item in enumerate(items):
    print(index, item)

first, *rest = row

for key, value in scores.items():
    print(key, value)
\end{code}

實際跑一次那個迴圈的例子，確認兩種寫法真的一模一樣：

\begin{example}{python}{(list-with-index.py) 同一份清單，兩種走法}
items = ["Milk", "Eggs", "Toast"]

for i in range(len(items)):
    print(i, items[i])

print("---")

for index, item in enumerate(items):
    print(index, item)
\end{example}

\iosample{}{0 Milk \\ 1 Eggs \\ 2 Toast \\ --- \\ 0 Milk \\ 1 Eggs \\ 2 Toast \\}

\subsection*{f-string 那一條，值得單獨講}

第 \ref{ch:strings} 章介紹 f-string 的時候只說了一句「變數一多就難維護」，沒有真的把
另一種寫法擺出來給你看。現在擺出來。

用 \texttt{+} 把一句話串起來，會踩到兩個實際的痛。第一個痛會直接讓程式當掉：

\begin{example}{python}{(greeting-plus.py) 用加號串一句話，忘了轉型}
name = "Alice"
age = 30

print("Hello, " + name + "! You are " + age + " years old.")
\end{example}

\begin{verbatim}
Traceback (most recent call last):
  File "greeting-plus.py", line 4, in <module>
    print("Hello, " + name + "! You are " + age + " years old.")
          ~~~~~~~~~~~~~~~~~~~~~~~~~~~~~~~~^~~~~
TypeError: can only concatenate str (not "int") to str
\end{verbatim}

\texttt{TypeError: can only concatenate str (not "int") to str}，意思是「字串只能跟字串
相加，不能跟 int 相加」。\texttt{age} 是數字，得先用 \texttt{str()} 包起來才行。這個錯不
難修，但它會反覆發生——每次你在句子裡多塞一個數字，就得記得包一次。

第二個痛不會報錯，所以更煩：標點跟空白全部要自己數。\texttt{"Hello, "} 後面那個空白、
\texttt{"! You are "} 前後那兩個空白，少一個就是輸出擠在一起，多一個就是中間多一格，而且
它們散落在一堆引號跟加號之間，很難一眼看出來。

\begin{example}{python}{(greeting-fstring.py) 兩種寫法，同一句話}
name = "Alice"
age = 30

print("Hello, " + name + "! You are " + str(age) + " years old.")
print(f"Hello, {name}! You are {age} years old.")
\end{example}

\iosample{}{Hello, Alice! You are 30 years old. \\ Hello, Alice! You are 30 years old. \\}

下面那一行裡，那句話本來的樣子是完整的——空白跟標點就長在它們該在的位置上，你不用在腦中
把五個片段拼起來，而且數字不必轉型。

\stylepoint{這張對照表裡的寫法，沒有一個是必要的——左邊那些全部能動。它們的共同好處只有
一個：\textbf{讓「你在做什麼」直接寫在程式碼的表面上}，不用讀的人自己拼。
\texttt{a, b = b, a} 表面上就寫著「交換」，\texttt{enumerate} 表面上就寫著「我要編號跟
東西」，f-string 表面上就是那句話本來的樣子。}

\begin{drills}[label=drills:value-idioms]
\item 把下面這段改寫成用 \texttt{enumerate}：
\begin{verbatim}
names = ["Alice", "Bob", "Cathy"]
for i in range(len(names)):
    print(str(i + 1) + ". " + names[i])
\end{verbatim}
\item 把下面這段改寫成 f-string：
\begin{verbatim}
city = "Taipei"
temp = 26
print("Today in " + city + " it is " + str(temp) + " degrees.")
\end{verbatim}
\item 下面這段會出什麼錯？錯誤訊息是什麼？怎麼修？
\begin{verbatim}
count = 3
print("You have " + count + " new messages.")
\end{verbatim}
\end{drills}

\section{\texttt{is} 還是 \texttt{==}：名字牌模型的收束}
\label{sec:is-vs-equal}

上一節結尾留了一個問題：要分辨「空的」跟「根本沒有」，該怎麼寫？答案很短，但要拿到它
之前，得先把一件跟了你十七章的東西處理掉。

\subsection*{一張走過的路}

名字牌這個比喻——變數不是裝東西的盒子，是貼在物件上的名字牌——在這本書裡出場過五次：

\begin{center}
\begin{tabular}{@{}p{5.2cm}p{9.2cm}@{}}
\toprule
在哪裡 & 那一次它解釋的是什麼 \\
\midrule
第 \ref{ch:variables} 章第 \ref{sec:name-tags} 節 &
為什麼 \texttt{x = 5} 之後 \texttt{x = "hi"} 完全合法 \\
\addlinespace
第 \ref{ch:lists} 章第 \ref{sec:aliasing} 節 &
為什麼改了 \texttt{b}，\texttt{a} 也跟著變了 \\
\addlinespace
第 \ref{ch:functions} 章第 \ref{sec:function-nametag}、\ref{sec:mutable-default} 節 &
為什麼函式裡的 \texttt{append} 會影響外面，以及 \texttt{cart=[]} 那個坑 \\
\addlinespace
第 \ref{ch:tuples} 章第 \ref{sec:tuple-vs-list} 節 &
為什麼「不能改」的 tuple 裡面的 list 還是改得動 \\
\addlinespace
第 \ref{ch:advanced-functions} 章第 \ref{sec:closure} 節 &
為什麼函式回傳之後，裡面那個計數器還記得上次的數字 \\
\bottomrule
\end{tabular}
\end{center}

這張表不是要你回去重讀那五節。它是要你注意一件你大概沒發現的事：\textbf{這五次，名字牌都是被
拿來解釋「一件看起來很怪的事，為什麼會發生」}。它一直是一張事後的圖，用來收拾直覺跟
事實對不上的場面。

這一節要做的，是換掉它的身分。

\subsection*{兩個問法}

\texttt{==} 跟 \texttt{is} 問的是兩個不同的問題：

\begin{itemize}
\item \texttt{a == b} 問的是：\textbf{這兩張名字牌指的物件，裝的東西一不一樣？}
\item \texttt{a is b} 問的是：\textbf{這兩張名字牌，是不是貼在同一個物件上？}
\end{itemize}

第 \ref{ch:lists} 章第 \ref{sec:aliasing} 節那個範例最後印出 True、True、True、False，
就是這兩個問題各問兩次的結果。這裡不重貼那段程式，只用一個更短的版本確認一次：

\begin{example}{python}{(compare-carts.py) 兩台內容一樣的購物車}
my_cart = ["Milk", "Eggs"]
your_cart = ["Milk", "Eggs"]

print(my_cart == your_cart)
print(my_cart is your_cart)
\end{example}

\iosample{}{True \\ False \\}

裝的東西一樣（\texttt{==} 是 True），但它們是兩個各自造出來的 list，不是同一個
（\texttt{is} 是 False）。

\subsection*{那我日常到底什麼時候該用 \texttt{is}？}

這個問題從第 \ref{ch:lists} 章就欠著了。那裡說 \texttt{is} 用在「比較少見的情境」，卻
沒告訴你少見到什麼程度。現在給答案，而且答案短得有點意外：

\keypoint{日常寫程式，你要的幾乎永遠是 \texttt{==}。\texttt{is} 只有一個常見用途：
\textbf{跟 \texttt{None} 比}。不確定的時候，用 \texttt{==}——它通常就是你要的，而且對
\texttt{None} 也管用。}

為什麼 \texttt{None} 特別？因為 \texttt{None} 全世界只有一個。這句話第 \ref{ch:functions}
章第 \ref{sec:mutable-default} 節講過了：那一節教的安全寫法 \texttt{if cart is None:}
正是這條規則最好的既有例子。整個程式裡不管哪裡出現 \texttt{None}，都是同一個物件，所以
問「是不是同一個」跟問「值一不一樣」會得到一樣的答案，而前者更直接。

\subsection*{兩個反例}

第一個反例是 \texttt{if x == None:}。它能動，答案也對，但沒有人這樣寫。既然全世界只有
一個 \texttt{None}，「是不是同一個」才是你真正想問的問題；繞去問「值一不一樣」，是拿一
把更大的工具做一件更小的事。

\stylepoint{跟 \texttt{None} 比，寫 \texttt{if x is None:}／\texttt{if x is not None:}，
不要寫 \texttt{if x == None:}／\texttt{if x != None:}。兩種都能動，但前者是所有 Python
程式碼裡的統一寫法，讀的人掃過去不用停。}

第二個反例是真的會害人的。看到上面那條「跟 \texttt{None} 比用 \texttt{is}」之後，很容易
順手把 \texttt{is} 用到別的地方去：

\begin{example}{python}{(is-zero.py) 拿 is 去比一個數字}
count = 0

if count is 0:
    print("empty")
\end{example}

\begin{verbatim}
is-zero.py:3: SyntaxWarning: "is" with 'int' literal. Did you mean "=="?
  if count is 0:
empty
\end{verbatim}

注意兩件事。第一，Python 發出的是 \texttt{SyntaxWarning}（警告），不是錯誤——
\index{SyntaxWarning}
程式沒有停下來，它照樣印出了 \texttt{empty}。第二，訊息裡 Python 直接問你「你是不是想
寫 \texttt{==}？」，它已經看出你大概搞錯了。

這一類的錯屬於第 \ref{ch:strings} 章第 \ref{sec:f-string-basics} 節說過的「安靜的錯」：
程式跑完了、有輸出、看起來一切正常，但它成立的理由不是你以為的那個。\texttt{count is 0}
這次剛好是 \texttt{True}，是因為 Python 內部對很小的整數做了一些最佳化，讓它們共用同一
個物件——換一個大一點的數字、或換一個字串，同樣的寫法就可能變成 \texttt{False}，而且一樣
不會報錯。

這就是為什麼那條規則要寫得那麼死：\textbf{跟 \texttt{None} 比用 \texttt{is}，其他一律
\texttt{==}}。不是因為別的用法一定會錯，是因為它什麼時候會錯、什麼時候不會錯，取決於你
不需要知道、也不該依賴的細節。

\subsection*{回到上一節那個問題}

現在可以回答「空的」跟「根本沒有」怎麼分了：

\begin{example}{python}{(missing-vs-empty.py) 三種情況，三種回答}
def report(items):
    if items is None:
        print("no list at all")
    elif not items:
        print("list is empty")
    else:
        print(len(items), "items")


report(None)
report([])
report(["Milk", "Eggs"])
\end{example}

\iosample{}{no list at all \\ list is empty \\ 2 items \\}

順序很重要：先問 \texttt{is None}，再問 \texttt{not items}。反過來的話，
\texttt{not items} 會把 \texttt{None} 一起吞掉，因為 \texttt{None} 本身也是 falsy。

\subsection*{名字牌，從這裡開始換一個身分}

名字牌從第 \ref{ch:variables} 章走到這裡，一直在替你解釋「為什麼會發生怪事」。每一次它
出現，都是因為某件事跟你的直覺對不上，需要一張圖來收拾。

現在它變成一件更小、也更常用的東西：\textbf{一條你打字的時候會用到的規則。}

跟 \texttt{None} 比用 \texttt{is}，其他一律 \texttt{==}。就這樣。你不需要每次打字都想起
那張圖，也不需要在心裡重演一次名字牌貼上去的過程——這條規則已經把那張圖要講的事收在裡面
了。而如果哪一天你忘了它為什麼是這樣，那張圖還在第 \ref{ch:variables} 章那裡等你回去看。

\begin{drills}[label=drills:is-vs-equal]
\item 下面這段會印出什麼？
\begin{verbatim}
a = [1, 2, 3]
b = [1, 2, 3]
c = a
print(a == b)
print(a is b)
print(a is c)
\end{verbatim}
\item 把 \texttt{if result == None:} 改成慣用寫法。如果要問的是「result 不是 None」呢？
\item \texttt{if status is "done":} 為什麼危險？它一定會出錯嗎？
\item 你要判斷兩台購物車裝的東西一不一樣（不在乎它們是不是同一台車）。該用
\texttt{==} 還是 \texttt{is}？
\end{drills}

\section{\texttt{with} 到底在解決什麼問題}
\label{sec:with-explained}

第 \ref{ch:files-exceptions} 章第 \ref{sec:reading-files} 節第一次教開檔案的時候，除了
把 \texttt{with} 的作用講成一句話（「這段縮排區塊做完之後，自動幫我把檔案關掉」）跟一個
比喻（用完自動幫你把號碼牌繳回去、置物櫃自動關好），還留了一句話：

\begin{quote}
你之後在很多情境都會看到 \texttt{with} 這個寫法，不只用在檔案上。
\end{quote}

這一節來把那句話講完。

順帶把當時講不完的另外半句也補上：那一章第 \ref{sec:finally} 節說過「\texttt{with} 背後
的機制，精神上跟 \texttt{finally} 很像」。當時只能講到「精神上」，因為 \texttt{finally}
是那一節才剛教的。現在你兩個都有了，可以把那句話從比喻換成一段真的程式碼。

\subsection*{不用 \texttt{with} 的話，要做哪些事？}

這個問題你答得出來：開檔、讀、關檔。

\begin{example}{python}{(open-close-by-hand.py) 自己開、自己關}
file_object = open("shopping_list.txt")
content = file_object.read()
file_object.close()

print(content)
\end{example}

\iosample{}{Milk \\ Eggs \\ Toast \\ Apple \\}

三行，看起來也還好。問題是這三行有三種壞掉的方式。

\textbf{第一種：忘了寫 \texttt{close()}。} 每個開著的檔案都會佔掉一點資源，而且大部分
作業系統對「同時能開幾個檔案」有上限。開一兩個沒事，一支程式在迴圈裡開了幾千個而且都沒關，
就會出事。更常見的後果是寫檔：沒有正確關閉的檔案，最後那一段內容可能根本還沒真的寫進去。

\textbf{第二種：\texttt{close()} 放太早。} 關掉之後才發現還要再讀一次，於是對著一個已經
關掉的檔案動手。這種錯的訊息等一下就會看到。

\textbf{第三種，也是最要命的：中間那一行炸了。} 如果 \texttt{read()} 那一行丟出例外，
程式會直接跳走，\texttt{close()} 那一行根本輪不到執行。前兩種是你會忘記，第三種是你就算
沒忘記也躲不掉——因為那一行寫在那裡，只是永遠沒被跑到。

\subsection*{把 \texttt{with} 拆回 \texttt{try/finally}}

第三種問題，你其實已經有工具可以解決了。第 \ref{ch:files-exceptions} 章第
\ref{sec:finally} 節教過 \texttt{finally}：不管 \texttt{try} 裡發生什麼事，
\texttt{finally} 裡的事一定會做。所以自己開檔而想寫得安全，要寫成這樣：

\begin{example}{python}{(open-with-finally.py) 不用 with，但要寫得安全}
file_object = open("shopping_list.txt")
try:
    content = file_object.read()
finally:
    file_object.close()

print(content)
\end{example}

\iosample{}{Milk \\ Eggs \\ Toast \\ Apple \\}

現在把 \texttt{with} 版本擺在旁邊：

\begin{code}{python}{同一件事的兩種寫法}
# 手寫版
file_object = open(filename)
try:
    content = file_object.read()
finally:
    file_object.close()

# with 版
with open(filename) as file_object:
    content = file_object.read()
\end{code}

\keypoint{\texttt{with} 就是上面那五行的簡寫。它替你保證的事只有一件，但那件事很重要：
\textbf{不管縮排區塊裡發生什麼——正常跑完、中途 \texttt{return}、還是丟出例外——收尾那
一步一定會執行。} 這正是第 \ref{sec:finally} 節那個 \texttt{finally} 在做的事，只是
\texttt{with} 把它縮成一行，而且連 \texttt{close()} 都不用你自己寫。}

到這裡，「\texttt{with} 是什麼」就有一個具體的答案了：它不是魔法，它是一段 \texttt{try/
finally} 的固定寫法，被包成一個關鍵字。

\subsection*{離開 \texttt{with} 之後，那個名字還在}

有一件事值得親眼看一次，因為它常被誤會：

\begin{example}{python}{(file-after-with.py) 離開之後再去碰那個檔案}
with open("shopping_list.txt") as file_object:
    content = file_object.read()

print(file_object)
print(file_object.closed)
print(file_object.read())
\end{example}

\begin{verbatim}
<_io.TextIOWrapper name='shopping_list.txt' mode='r' encoding='UTF-8'>
True
Traceback (most recent call last):
  File "file-after-with.py", line 6, in <module>
    print(file_object.read())
          ~~~~~~~~~~~~~~~~^^
ValueError: I/O operation on closed file.
\end{verbatim}

三行輸出各講一件事。第一行：離開 \texttt{with} 之後，\texttt{file\_object} 這個名字還在，
\texttt{print} 印得出東西——\texttt{with} 不像函式那樣會建立一個獨立的空間，區塊裡的名字
出來之後照樣用得到。第二行：那個檔案物件的 \texttt{closed} 屬性是 \texttt{True}，它確實
被關掉了。第三行：真的想再讀一次，就會拿到
\texttt{ValueError: I/O operation on closed file.}（對一個已經關掉的檔案做輸入輸出）。

用上一節剛收好的那條規則的原話來說：名字牌沒有被撕掉，被改變的是它指著的那個物件的狀態。

\subsection*{「不只用在檔案上」是什麼意思}

回到那張欠條。\texttt{with} 建立的是一段\textbf{有開頭、有結尾的區間}：進去的時候做一件
事，出來的時候一定做另一件事——不管你是正常走出來，還是中途被例外叫走。

開檔跟關檔只是這種「一對」裡最常見的一組。其他常見的組合，你以後遇到的時候就會認得：

\begin{itemize}
\item 連上資料庫，用完斷開連線。
\item 打開一個壓縮檔往裡面塞東西，塞完把它封好。第 \ref{ch:automating} 章做備份的時候
就會再用到一次同樣的寫法，那時候你會發現它跟開檔案長得一模一樣。
\end{itemize}

至於「怎麼自己做一個 \texttt{with} 用得上的東西」，那需要先學會 class，超出這本書的範圍。
好消息是你幾乎不需要——標準函式庫跟常見的套件早就替你做好了，你要做的是\textbf{看到
\texttt{with} 的時候知道它在保證什麼}，以及看到一段「開了就一定要關」的程式碼時，想起有
這個東西。

寄物櫃那個比喻現在可以說得更準一點：不管你是正常辦完事走出來，還是中途被叫走，那個櫃子
都會被關好，號碼牌都會被繳回去。

\begin{drills}[label=drills:with-explained]
\item 下面這段會印出什麼？（提示：不只一行）
\begin{verbatim}
with open("notes.txt") as f:
    text = f.read()

print(f.closed)
print(f.read())
\end{verbatim}
\item 把下面這段改寫成用 \texttt{with}：
\begin{verbatim}
log_file = open("log.txt", "a")
log_file.write("done\n")
log_file.close()
\end{verbatim}
\item 如果 \texttt{with} 區塊裡的程式丟出例外，檔案還會被關上嗎？為什麼？
\end{drills}

\section{先做再說：EAFP 是一整套，不是一個語法}
\label{sec:eafp-closure}

第 \ref{ch:files-exceptions} 章第 \ref{sec:lbyl-vs-eafp} 節教過一個縮寫：EAFP，
easier to ask forgiveness than permission——與其先申請許可，不如先做了再道歉。那一章還
用第 \ref{sec:eafp-recap} 節整節回頭指出，這套邏輯遠比「檔案讀寫」這一件事大。

這一節要再往上一層。往上一層具體是三件事。

\subsection*{一、它不是一個語法，是一整個家族}

第 \ref{ch:files-exceptions} 章講的是「面對可能失敗的事，Python 偏好先做」。但這個偏好
影響的不只是你怎麼寫 \texttt{try}——它長在整個 Python 的設計裡。下面這張表裡的每一項，
你都已經在用了，只是從來沒有人告訴你它們是同一家人：

\begin{center}
\begin{tabular}{@{}p{3.6cm}p{3.6cm}p{7.2cm}@{}}
\toprule
你已經在用的東西 & 教在哪 & 它是這個家族的哪一種形狀 \\
\midrule
\texttt{try} / \texttt{except} &
第 \ref{ch:files-exceptions} 章第 \ref{sec:lbyl-vs-eafp} 節 &
本體：直接做，錯了接住 \\
\addlinespace
\texttt{dict.get(key, default)} &
第 \ref{ch:dict-set} 章第 \ref{sec:dict-get} 節 &
Python 替你把「查不到就用預設值」整包做好了 \\
\addlinespace
\texttt{set.discard()} &
第 \ref{ch:dict-set} 章第 \ref{sec:set-membership} 節 &
同上：不在裡面就當沒事，不必先問 \\
\addlinespace
\texttt{if items:} / \texttt{if not items:} &
第 \ref{ch:conditionals} 章、第 \ref{sec:condition-idioms} 節 &
語言自己把「有沒有內容」換成答案，不必你先量長度 \\
\addlinespace
\texttt{with} &
第 \ref{sec:with-explained} 節 &
\texttt{try/finally} 的簡寫——負責收尾的那一半 \\
\addlinespace
讀得懂的錯誤訊息 &
第 \ref{ch:debugging} 章 &
既然失敗是預期會發生的事，訊息當然要寫得能讀 \\
\bottomrule
\end{tabular}
\end{center}

前面三格用程式碼寫出來，你會發現它們短得幾乎不像在處理失敗：

\begin{code}{python}{同一個家族的三種說法}
value = scores.get(name, 0)    # 查不到？給 0，不必先問在不在
tags.discard("draft")          # 不在裡面？當沒事，不必先問在不在
if not items:                  # 空的？直接問，不必先量長度
    print("nothing to do")
\end{code}

這張表裡最值得停一下的是 \texttt{with} 那一格。上一節剛把 \texttt{with} 拆回
\texttt{try/finally}，所以現在你看得懂它為什麼在這張表裡：處理失敗有兩半，一半是「接住
它」（\texttt{except}），一半是「不管有沒有失敗都要收尾」（\texttt{finally}），而
\texttt{with} 就是後面那一半的固定寫法。第 \ref{ch:files-exceptions} 章沒辦法把這一格
填進來，因為那時候 \texttt{with} 還只是一個「照著寫就對了」的關鍵字。

\subsection*{二、這個偏好是從哪來的}

不是因為 Python 比較潮，也不只是因為第 \ref{ch:files-exceptions} 章講過的那個理由（檢查
完到真的動手之間，狀態可能已經變了）。往上一層的說法是這樣的：

\keypoint{Python 把「失敗」當成一件\textbf{會正常發生的事}，而不是一件\textbf{該被預防的
意外}。所以它給你的是接住失敗的工具，而不是要求你事先把所有可能出錯的情況都想過一遍。
查不到、開不起來、轉不了型——這些事本來就會發生，程式的工作是處理它們，不是假裝它們不
存在。}

這句話解釋了很多你早就見過、卻沒想過為什麼的事。為什麼 \texttt{dict} 有一個
\texttt{get()}？為什麼 \texttt{set} 除了 \texttt{remove()} 還要有一個 \texttt{discard()}？
為什麼 Python 的錯誤訊息會長到用一整段話告訴你哪裡不對，還畫箭頭指出是哪個位置——如果
失敗被當成不該發生的意外，訊息只要寫「失敗」兩個字就夠了。

\subsection*{三、兩條線在這裡會合}

這本書從頭到尾只有兩個心智模型：名字牌，還有 EAFP。它們一個在第 \ref{ch:variables} 章
起頭，一個在第 \ref{ch:files-exceptions} 章起頭，一路上看起來完全沒有關係——
一個在講變數怎麼運作，一個在講錯誤怎麼處理。

把它們各自要修正的那個直覺並排看，關係就出來了：

\begin{center}
\begin{tabular}{@{}p{3.8cm}p{4.4cm}p{6.2cm}@{}}
\toprule
你原本的直覺 & Python 實際做的事 & 這個直覺假設了什麼 \\
\midrule
\texttt{b = a} 會複製一份 &
只是多貼一張名字牌 &
假設 Python 會替我多做一步（複製） \\
\addlinespace
動手之前要先檢查 &
直接做，失敗了告訴你 &
假設 Python 應該替我多做一步（防呆） \\
\bottomrule
\end{tabular}
\end{center}

兩個直覺其實是同一個直覺：\textbf{以為 Python 會在你沒開口的時候，替你多做一步。}

它不會。Python 的預設是「不多做」——不替你複製，也不替你預檢。這件事的代價是你得知道它
預設不做什麼，這也是為什麼這兩件事要花十幾章反覆提醒。好處是另一面：\textbf{沒有任何一
步是在你背後偷偷發生的。} 你想要複製，就自己寫 \texttt{.copy()}；你想要檢查，就自己寫
\texttt{if}。你寫幾步，Python 就做幾步。

這兩個模型從第 \ref{ch:variables} 章跟第 \ref{ch:files-exceptions} 章分頭走了十幾章，走
到這裡才會合。

\begin{drills}[label=drills:eafp-closure]
\item 下面四段程式，哪幾段屬於這一節那張表裡的家族？
\begin{verbatim}
# (1)
name = people.get(user_id, "unknown")

# (2)
if x % 2 == 0:
    print("even")

# (3)
with open("data.txt") as f:
    text = f.read()

# (4)
if os.path.exists(path):
    f = open(path)
\end{verbatim}
\item \texttt{with} 為什麼算 EAFP 家族的一員？用一句話說明。
\item 「這個數是不是偶數」這種判斷，該不該包成 \texttt{try/except}？為什麼？
\end{drills}

\section{什麼時候不要太 Pythonic}
\label{sec:not-too-pythonic}

從第 \ref{ch:advanced-functions} 章到現在，你連續學了五章新寫法：\texttt{lambda}、
\texttt{sorted(key=)}、decorator、comprehension、generator、遞迴，再加上這一章的六節
慣用寫法。這一節是煞車。

\subsection*{同一件事，三個版本}

給你一份成績資料（假設是從檔案讀進來的），要做三件事：篩掉不及格的、換算成等第、按分數
由高到低排序。

A 版，全部擠成一行：

\begin{example}{python}{(grades-one-liner.py) A 版：一行寫完}
rows = [
    {"name": "Alice", "score": 88},
    {"name": "Bob", "score": 54},
    {"name": "Cathy", "score": 92},
    {"name": "David", "score": 71},
]

result = [(r["name"], "A" if r["score"] >= 90 else "B" if r["score"] >= 80 else "C") for r in sorted([r for r in rows if r["score"] >= 60], key=lambda r: r["score"], reverse=True)]

print(result)
\end{example}

\iosample{}{[('Cathy', 'A'), ('Alice', 'B'), ('David', 'C')] \\}

B 版，拆成幾個有名字的步驟：

\begin{example}{python}{(grades-steps.py) B 版：一步一個名字}
rows = [
    {"name": "Alice", "score": 88},
    {"name": "Bob", "score": 54},
    {"name": "Cathy", "score": 92},
    {"name": "David", "score": 71},
]


def grade_of(score):
    if score >= 90:
        return "A"
    if score >= 80:
        return "B"
    return "C"


passed = [row for row in rows if row["score"] >= 60]
ranked = sorted(passed, key=lambda row: row["score"], reverse=True)
result = [(row["name"], grade_of(row["score"])) for row in ranked]

print(result)
\end{example}

\iosample{}{[('Cathy', 'A'), ('Alice', 'B'), ('David', 'C')] \\}

C 版，完全不用這五章教過的任何東西，全部手寫迴圈：

\begin{example}{python}{(grades-plain-loop.py) C 版：全部手寫}
rows = [
    {"name": "Alice", "score": 88},
    {"name": "Bob", "score": 54},
    {"name": "Cathy", "score": 92},
    {"name": "David", "score": 71},
]

passed = []
for row in rows:
    if row["score"] >= 60:
        passed.append((row["score"], row["name"]))

passed.sort(reverse=True)

result = []
for score, name in passed:
    if score >= 90:
        grade = "A"
    elif score >= 80:
        grade = "B"
    else:
        grade = "C"
    result.append((name, grade))

print(result)
\end{example}

\iosample{}{[('Cathy', 'A'), ('Alice', 'B'), ('David', 'C')] \\}

三個版本輸出一模一樣。哪一個比較好？

先不要急著回答。

\subsection*{三條判準，而且判準的對象是人}

\begin{itemize}
\item \textbf{讀的人是誰？} 只有你自己看、跑一次就丟的小腳本，A 版沒問題——它確實比較快
寫完。要交給別人維護，或者三個月後你自己還要回來改，那就是 B 版。這個問題沒有一般答案，
因為它問的不是程式碼，是情況。
\item \textbf{一行讀不完，就沒有優點了。} 慣用寫法的價值來自「一眼看懂」，一旦要折行、
要往右捲、要從中間往兩邊讀，那個價值就不見了，剩下的只有難度。這條就是第
\ref{ch:comprehensions} 章第 \ref{sec:comp-limits} 節那條「超過兩層或兩行就拆開」的
一般化——那裡講的是 comprehension，這裡對這五章教過的每一種工具都成立。
\item \textbf{寫的時候用光了聰明，除錯的時候就沒有聰明可以用了。} 這句話改寫自
Kernighan 的一句老話：「每個人都知道，除錯比一開始寫程式難兩倍。所以，
如果你寫的時候就用盡了自己所有的聰明，你要拿什麼去除錯它？」第 \ref{ch:debugging} 章教的每一種除錯手法——讀 traceback、插
\texttt{print}、一步一步縮小範圍——全部都有一個前提：你得看得懂那段程式碼在做什麼。
\end{itemize}

\stylepoint{要不要用一個比較精簡的寫法，問三句話就夠了：\textbf{誰會讀它？一行讀不讀得
完？出錯的時候我還追得動嗎？} 三句話問的都不是程式碼本身，是人——寫程式碼的你、讀程式碼
的人、還有半年後回來修它的那個人。}

\subsection*{回答那個問題}

B 版。

理由不是「B 比較 Pythonic」——這種說法沒有意義，因為 A 版用到的 Python 特性其實比 B 版
還多。理由是很具體的一件事：A 版把篩選、排序、換算三件事壓在同一行裡，所以當結果不對的
時候，你分不出是哪一件出的錯。你沒辦法在中間插一個 \texttt{print} 看看篩完剩下幾筆，因為
「篩完」這個中間狀態根本沒有名字。B 版有 \texttt{passed}、有 \texttt{ranked}，每一步都
停得下來、看得到。

順帶一提，A 版還有一個實務上的麻煩：它有 152 個字元，在一般的螢幕寬度下讀不完。

\subsection*{反過來的錯誤，也是錯的}

讀到這裡有一個很自然的反應：那我乾脆全部寫 C 版好了，反正也能動，而且絕對不會太花俏。

這也是錯的，而且錯得沒有比較少。有兩件事要講清楚。

\textbf{第一，讀跟寫的標準不一樣。} 別人的程式碼裡一定會出現 comprehension、
\texttt{with}、\texttt{if not items:}、\texttt{a, b = b, a}——這些東西用得非常多，多到
你不可能繞開。\textbf{看得懂它們是不能談判的}；輪到自己寫的時候，你才有選擇權。第
\ref{ch:advanced-functions} 章第 \ref{sec:map-filter} 節對 \texttt{map} 跟 \texttt{filter}
說過「看得懂就好」，那句話的一般化就是這條：能讀是底線，怎麼寫是選擇。

\textbf{第二，回到這一章開頭那個比喻。} 怕講得有腔調，就乾脆不講話——那不叫沒有口音，
那叫沒有在講話。C 版不是「比較保守的選擇」，它是把一件三句話講得完的事講成十五句。

所以這一節不是在勸你少用工具，是在勸你為每一次使用給得出理由。「因為看起來比較厲害」
不是理由，「因為這樣一眼看得出我在篩什麼」是。

\subsection*{最後一句}

這一章從頭到尾沒有教你任何新東西——嚴格說起來只有一個連續比較，而那也只是把
\texttt{and} 少寫一次。它做的只有一件事：把「這樣寫可以」變成「這樣寫，下一個人讀起來
會比較輕鬆」。

誰是下一個人，只有你知道。

\begin{drills}[label=drills:not-too-pythonic]
\item 下面這一行出現在一支你要交給同事接手的程式裡。該不該拆開？
\begin{verbatim}
out = "\n".join(f"{k}: {sum(v)/len(v):.1f}" for k, v in data.items())
\end{verbatim}
\item 同事的程式碼裡有一段你看不懂的 comprehension。你該叫他改掉嗎？
\item 你寫了一支只有自己會用、跑一次就刪掉的小腳本。需要照 PEP 8 排版嗎？
\end{drills}

下一章換一個方向。到目前為止，你的資料一直是散的：一份 list 裝分數、另一份 dict 裝名字，
要用的時候自己記得它們是一組。第 \ref{ch:objects} 章要教一種把資料跟操作它的動作綁在
一起的做法，叫做 class。

\section*{練習解答}
\addcontentsline{toc}{section}{練習解答}

\paragraph{第~\ref{sec:pep8}~節}
\begin{enumerate}
\item 一種寫法：
\begin{verbatim}
total_items = 0


def count_them(the_list):
    global total_items
    for item in the_list:
        total_items = total_items + 1
    return total_items
\end{verbatim}
改了四件事：變數與函式名改成 \texttt{snake\_case}；\texttt{x} 換成看得出意思的
\texttt{item}；等號前後補空白；\texttt{def} 前面空兩行。（\texttt{TotalItems} 不是常數，
它會被改動，所以不該用全大寫。）
\item 不會阻止你，第二次指派會直接生效。全大寫幫的是\textbf{讀的人}——它在說「這個值是
設定，不打算被改」。Python 沒有常數這種東西，這個慣例防不了任何事，它只是一個講好的暗號。
\item 不算錯。PEP 8 自己就寫著「知道什麼時候該不一致」，而且行長那條特別常被團隊改。真正
會出問題的是團隊裡一半的人用 79、一半的人用 99——不一致的成本比數字本身高得多。
\end{enumerate}

\paragraph{第~\ref{sec:condition-idioms}~節}
\begin{enumerate}
\item \texttt{if ext in (".jpg", ".png", ".gif"):}
\item \texttt{if 18 <= temperature <= 26:}
\item 依序是 \texttt{True}、\texttt{False}、\texttt{True}、\texttt{True}。前兩行：空 list
是 falsy，所以 \texttt{not items} 是 \texttt{True}，但它不是 \texttt{None}。後兩行：
\texttt{None} 也是 falsy，所以 \texttt{not items} 一樣是 \texttt{True}，而這次
\texttt{is None} 也是 \texttt{True}。關鍵是第一行跟第三行答案一樣——這正是
\texttt{if not items:} 分不出「空的」跟「根本沒有」的地方。
\item 當你需要分辨「空的」跟「根本沒有」的時候。例如一個函式回傳「搜尋結果」，
\texttt{[]} 代表「找過了，沒有符合的」，\texttt{None} 代表「根本沒搜尋成功」——這兩件事
要給使用者的訊息完全不同，用 \texttt{if not items:} 會把它們混成一種。這種時候要先問
\texttt{if items is None:}。
\end{enumerate}

\paragraph{第~\ref{sec:value-idioms}~節}
\begin{enumerate}
\item
\begin{verbatim}
names = ["Alice", "Bob", "Cathy"]
for index, name in enumerate(names, start=1):
    print(f"{index}. {name}")
\end{verbatim}
（不用 \texttt{start=1} 也可以，改寫成 \texttt{f"\{index + 1\}. \{name\}"}。順帶把
\texttt{+} 串接也一起換掉了。）
\item \texttt{print(f"Today in \{city\} it is \{temp\} degrees.")}
\item 錯誤是
\begin{verbatim}
TypeError: can only concatenate str (not "int") to str
\end{verbatim}
因為 \texttt{count} 是整數，不能直接用 \texttt{+} 接在字串後面。修法有兩種：包成
\texttt{str(count)}，或直接改寫成
\texttt{print(f"You have \{count\} new messages.")}——後者不必轉型，也少數幾個空白。
\end{enumerate}

\paragraph{第~\ref{sec:is-vs-equal}~節}
\begin{enumerate}
\item 依序是 \texttt{True}、\texttt{False}、\texttt{True}。\texttt{a} 跟 \texttt{b}
內容一樣但是兩個各自造出來的 list，所以 \texttt{==} 是 \texttt{True}、\texttt{is} 是
\texttt{False}；\texttt{c = a} 只是多貼一張名字牌，所以 \texttt{a is c} 是 \texttt{True}。
\item 改成 \texttt{if result is None:}。要問「不是 \texttt{None}」的話，
慣用寫法是 \texttt{if result is not None:}，不是
\texttt{if not result is None:}——後者雖然也能動，但讀起來要多繞一次。
\item 因為它問的是「這兩個是不是同一個物件」，而你想問的其實是「內容一不一樣」。
它不一定會出錯——如果 \texttt{status} 是程式裡直接寫死的 \texttt{"done"}，Python 可能
讓兩邊共用同一個字串物件，結果剛好是 \texttt{True}；但如果 \texttt{status} 是拼出來的
（例如從檔案讀進來、或用 \texttt{join()} 組出來的），同樣的比較就會變成 \texttt{False}，
而且不會報錯。「有時候對、有時候錯，還不告訴你」比「一定錯」更難抓。正確寫法是
\texttt{if status == "done":}。
\item \texttt{==}。你要問的是內容，不是身分。用 \texttt{is} 的話，兩台裝著一樣東西的
購物車會得到 \texttt{False}——除非它們本來就是同一台。
\end{enumerate}

\paragraph{第~\ref{sec:with-explained}~節}
\begin{enumerate}
\item 先印出 \texttt{True}（檔案已經被 \texttt{with} 關掉了），然後第二個 \texttt{print}
會丟出 \texttt{ValueError: I/O operation on closed file.}，程式停在那裡。重點是
\texttt{f} 這個名字本身還在、還用得到，被關掉的是它指著的那個檔案物件。
\item
\begin{verbatim}
with open("log.txt", "a") as log_file:
    log_file.write("done\n")
\end{verbatim}
\item 會。這正是 \texttt{with} 存在的理由。把它拆回 \texttt{try/finally} 就看得很清楚：
關檔那一步寫在 \texttt{finally} 裡，而 \texttt{finally} 不管 \texttt{try} 裡發生什麼事
都會執行。如果自己寫 \texttt{open()} 加 \texttt{close()} 而中間丟了例外，\texttt{close()}
那一行根本輪不到跑。
\end{enumerate}

\paragraph{第~\ref{sec:eafp-closure}~節}
\begin{enumerate}
\item (1)、(3) 是。(1) 是 \texttt{dict.get()}，把「查不到就用預設值」包好了；(3) 是
\texttt{with}，負責收尾的那一半。(2) 不是——它是一個完全可預期、跟外部無關的條件判斷，
本來就該用 \texttt{if}，跟「可能失敗」無關。(4) 剛好是相反的那一種：它先檢查再動手，
是典型的 LBYL，而且檢查完到真的 \texttt{open()} 之間，檔案還可能被刪掉。
\item 因為它是 \texttt{try/finally} 的簡寫，保證「不管中間成不成功，收尾一定會做」——
而「假設事情可能失敗，並且先把善後安排好」正是這個家族的共同精神。
\item 不該。這是一個完全可預期、只跟手上這個數字有關的判斷，不是「可能因為外部因素失敗」
的操作，用 \texttt{if x \% 2 == 0:} 就好。EAFP 不是「所有地方都改成 \texttt{try}」，
它針對的是檔案、網路、使用者輸入、字典查詢這一類本來就可能落空的事。
\end{enumerate}

\paragraph{第~\ref{sec:not-too-pythonic}~節}
\begin{enumerate}
\item 該拆。它同時做了三件事：走過每一組資料、算平均並格式化到小數點後一位、把結果
組成多行字串，而且一行讀不完。至少拆成兩步（先算好一份 \texttt{averages}，再組字串），
出錯的時候才有地方插 \texttt{print} 看中間結果。如果這是你自己跑一次就丟的腳本，
留著也無妨——但題目說了要交給同事接手。
\item 不該叫他改掉，先去讀懂它。這是「讀跟寫的標準不一樣」那條界線：看得懂別人的慣用寫法
是不能談判的，那是你的功課不是他的。讀懂之後如果發現它真的違反了上面那三條判準（一行讀
不完、三件事壓在一起），那時候再提，而且要講的是具體理由，不是「我看不懂」。
\item 不需要，但排一下對你自己有好處。PEP 8 自己就說了「知道什麼時候該不一致」，一支馬上
要刪掉的腳本沒有下一個讀者。不過有兩件事值得留：變數名取得看得懂，還有等號前後留空白——
因為「跑一次就刪掉」的腳本，往往不會真的被刪掉。
\end{enumerate}

\chapter{物件與類別基礎}
\label{ch:objects}

前面十四章，你已經很會用 Python 內建的資料結構（list、dict、tuple）把資料存起來、
一個一個處理。但如果要記錄的是「一個學生」這種東西，這個學生有姓名、有分數，那個學生
也有姓名、有分數，一次要記好幾份，彼此獨立，還不能搞混，用一堆各自獨立的變數，或是一層
疊一層的 dict，程式會越寫越亂，稍微一個地方漏改一個欄位名稱，就整個對不上。

這一章要教的「類別」（class）跟「物件」（object），就是 Python 用來解決這個問題的辦法：
你先寫一份「模板」，說清楚「這種東西應該長什麼樣子、會做什麼事」，之後每次要用就照著模板
造一個出來，各自獨立，互不影響。

這一章分成幾段走：

\begin{itemize}
\item 先從一個你其實已經很熟的東西重新認識一次「物件」這個詞：你前面章節一路用到現在的
list、字串、dict，全部都是物件。
\item 接著動手寫你自己的第一個類別，搞懂 \texttt{class}、\texttt{\_\_init\_\_}、
\texttt{self} 這幾個新語法到底在做什麼，還有方法怎麼定義、怎麼呼叫，甚至怎麼在同一個
物件裡互相借用彼此的結果。
\item 再來會處理兩個物件互動的情境：一個方法怎麼接住另一個物件當參數，直接拿它身上的
屬性跟自己比較。
\item 最後把範圍拉大一點，處理一個很多初學者會卡住的地方，「每個物件自己的屬性」跟
「整個類別共用的屬性」是兩回事，還有為什麼值得花力氣把資料跟操作綁在一起、寫這些額外的
語法。
\end{itemize}

這一章只處理「基礎」，類別怎麼定義、物件怎麼建立、屬性怎麼分兩種；「一個類別可以繼承
另一個類別的能力」這件事，留到下一章。

\section{重新認識「物件」：你其實早就在用了}
\label{sec:objects-recap}

回頭看一下你這幾章寫過的程式碼，你可能已經打過類似這樣的東西：

\begin{quote}
\ttfamily
my\_list.append(3) \\
"hello".upper() \\
my\_dict.keys()
\end{quote}

這些寫法有一個共通的形狀：\texttt{東西.動作(...)}。\texttt{my\_list} 這個 list、
\texttt{"hello"} 這個字串、\texttt{my\_dict} 這個 dict，本身就「帶著」一堆你可以直接
呼叫的動作，這些動作直接「長」在這個東西身上，用一個點接著呼叫，不用另外寫一個函式、把
資料傳進去。

像這樣「資料」跟「可以對這份資料做的動作」綁在一起的東西，在 Python 裡就叫做
\textbf{物件}（object）；那些長在物件身上、用點記法呼叫的動作，叫做\textbf{方法}
（method）。所以「物件」不是一個要重新學的陌生東西，你從很早以前就一直在用了，只是還
沒有人跟你點名這件事而已。

那為什麼還需要自己定義新的物件呢？因為 Python 內建的物件種類就那幾種，數字、字串、
list、dict、tuple，但真實世界要處理的東西千奇百怪，「一個學生」「一筆訂單」「一台
車」，這些東西各自帶著好幾種資料，還會被重複建立很多份，內建的物件種類完全不夠用。

想像一下你要辦一場活動，需要每個報名的人填一張報名表：表格本身長什麼樣子（有姓名欄、
有生日欄、有聯絡電話欄）是固定的設計，這張設計只需要做一次；但每個人各自填了自己的
姓名、自己的生日，彼此互不影響，也不會因為小美填了她的姓名，阿凱那張表上的姓名欄就跟著
變。

「類別」就是那張空白表格的設計，你只寫一次，說清楚「這種東西應該有哪些欄位、能做哪些
事」；「物件」則是照著這張設計、每次實際造出來的一份，各自獨立。這一章接下來要做的事，
就是學會怎麼自己設計這種「表格」，以及怎麼照著表格造出一份份獨立的「物件」。
\index{類別（class）}

你在第 \ref{ch:variables} 章學過，變數不是裝著值的盒子，是貼在物件上的名字牌。這一章要教的，是怎麼自己
造一種新的「可以被貼名字牌的東西」，不管是內建的 list、字串，還是你等一下要自己定義的
「學生」，都遵守同一套規則。

\section{動手寫第一個類別：\texttt{class}、\texttt{\_\_init\_\_}、屬性}
\label{sec:class-init}

要造一種新的物件，第一步是用 \texttt{class} 關鍵字定義它的「表格設計」長什麼樣子。
一個類別裡最常見不過的一段，是一個叫 \texttt{\_\_init\_\_} 的方法，它的名字裡那兩對
底線是 Python 的固定寫法，意思大致是「初始化」：每次你照著這個類別造一個新物件，Python
都會自動呼叫這個方法一次，負責幫這個新物件把該有的資料填好。
\index{\_\_init\_\_}

\begin{code}{python}{定義最基本的類別}
class ClassName:
    def __init__(self, param1, param2):
        self.attr1 = param1
        self.attr2 = param2
\end{code}

這裡有兩個第一次看到的東西特別值得說明。第一個是 \texttt{self}：它永遠是
\texttt{\_\_init\_\_}（以及等一下會教的其他方法）的第一個參數，代表「這個正在被造出來
的物件自己」，你不用自己傳這個參數，Python 會自動幫你補上。
\index{self}
第二個是
\texttt{self.attr1 = param1} 這種寫法：等號左邊多了一個 \texttt{self.}，意思是「把
這個值存到\textit{這個物件自己身上}的 \texttt{attr1} 這個位置」，這種存在物件身上的
資料就叫做\textbf{屬性}（attribute）。
\index{屬性（attribute）}

定義好類別之後，用「類別名稱(...)」的形式就能造出一個新物件，這個動作叫做
\textbf{建立實例}（instantiation）；造出來的物件則叫做這個類別的一個\textbf{實例}
（instance）。
\index{建立實例（instantiation）}
\index{實例（instance）}

\begin{example}{python}{(student.py) 記錄一個學生的姓名與分數}
class Student:
    def __init__(self, name, score):
        self.name = name
        self.score = score


s1 = Student("Alice", 88)
s2 = Student("Bob", 76)

print(s1.name, s1.score)
print(s2.name, s2.score)
\end{example}

\iosample{}{Alice 88 \\ Bob 76 \\}

\texttt{Student("Alice", 88)} 這一行，就是照著 \texttt{Student} 這個模板造出一個新物件，
Python 自動把 \texttt{s1} 這個還沒建好的物件當成 \texttt{self}、把 \texttt{"Alice"}
當成 \texttt{name}、把 \texttt{88} 當成 \texttt{score}，傳進 \texttt{\_\_init\_\_} 裡
執行，於是 \texttt{s1} 這個物件身上就有了 \texttt{name} 跟 \texttt{score} 兩個屬性。
之後想取用某個物件的某個屬性，就用「物件.屬性名稱」這種點記法，跟你之前呼叫方法的形狀
很像，只是這裡取的是資料，不是動作。

順帶一提，如果你這時候把 \texttt{print(s1.name)} 換成直接 \texttt{print(s1)}，印整個
物件本身，會看到類似這樣的東西：

\begin{verbatim}
<__main__.Student object at 0x1053a2f90>
\end{verbatim}

這串位址（\texttt{0x} 開頭那串數字）每次執行都不一樣，不用背它，也不用理解它，這只是
Python 目前還不知道你想把這個物件印成什麼樣子，所以印出了一個很陽春的預設格式。這行字
確實不好讀，下一章會教你怎麼換成好讀的格式，這裡先知道有這回事就好。

初學者在這一步最容易犯的錯，是漏寫 \texttt{self.} 這個前綴。看看下面這段哪裡怪怪的：

\begin{example}{python}{(student-broken-self.py) 忘記寫 self.\ 前綴會怎樣}
class Student:
    def __init__(self, name, score):
        name = name          # 漏寫 self.，只是一個普通的區域變數
        self.score = score


s1 = Student("Alice", 88)
print(s1.name)
\end{example}

執行後看到的是：

\begin{verbatim}
Traceback (most recent call last):
  File "student-broken-self.py", line 8, in <module>
    print(s1.name)
          ^^^^^^^
AttributeError: 'Student' object has no attribute 'name'
\end{verbatim}

\keypoint{\texttt{self.attr = param} 是「把名字牌 \texttt{attr} 貼到這個物件身上」；
如果漏寫 \texttt{self.}，\texttt{name = name} 只是在 \texttt{\_\_init\_\_} 這個方法
執行完就消失的一個普通區域變數，根本沒有貼在物件上，物件建立完之後，身上壓根沒有
\texttt{name} 這個屬性，所以之後想用 \texttt{s1.name} 去讀，Python 才會抱怨
\texttt{AttributeError}，說這個物件沒有這個屬性。}

\begin{drills}[label=drills:class-init]
\item 這段程式執行後會印出什麼？
\begin{verbatim}
class Item:
    def __init__(self, name, price):
        self.name = name
        self.price = price

item = Item("notebook", 45)
print(item.name, item.price)
\end{verbatim}
\item 這段程式碼裡 \texttt{price} 那一行漏寫了 \texttt{self.}，執行
\texttt{print(item.price)} 時會發生什麼事？錯誤類型是什麼？
\begin{verbatim}
class Item:
    def __init__(self, name, price):
        self.name = name
        price = price

item = Item("notebook", 45)
print(item.price)
\end{verbatim}
\item 幫 \texttt{Item} 類別的 \texttt{\_\_init\_\_} 補上第三個屬性
\texttt{in\_stock}（布林值，代表是否有庫存），並建立一個 \texttt{Item("pencil", 10)}
的物件，把 \texttt{in\_stock} 設成 \texttt{True}。你會怎麼修改
\texttt{\_\_init\_\_} 的參數列跟裡面的兩行 \texttt{self.} 賦值？
\end{drills}

\section{方法與 self：讓物件自己會做事}
\label{sec:methods-and-self}

\texttt{\_\_init\_\_} 只負責在物件剛造出來的時候填好資料，但一個物件通常還需要能「做
點什麼」，像是印出自己的摘要、更新自己的某個屬性。這些額外的動作，就是定義在類別裡、
\texttt{\_\_init\_\_} 以外的一般方法。

\begin{code}{python}{幫類別加上方法}
class ClassName:
    def __init__(self, param):
        self.attr = param

    def method_name(self):
        # 用 self.attr 做點什麼
        ...
\end{code}

跟 \texttt{\_\_init\_\_} 一樣，每一個方法的第一個參數都要是 \texttt{self}，代表「呼叫
這個方法的那個物件自己」。呼叫的時候你不用自己傳它，寫 \texttt{s1.describe()}，Python
會自動把 \texttt{s1} 傳給方法裡的 \texttt{self}，所以你在方法裡才能用
\texttt{self.attr} 存取到「呼叫者自己」的那份資料。

\begin{example}{python}{(student-methods.py) 幫學生類別加上方法}
class Student:
    def __init__(self, name, score):
        self.name = name
        self.score = score

    def describe(self):
        print(f"Name: {self.name}, Score: {self.score}")

    def update_score(self, new_score):
        self.score = new_score


s1 = Student("Alice", 88)
s2 = Student("Bob", 76)

s1.describe()
s2.describe()

s1.update_score(95)
print("After update:")
s1.describe()
s2.describe()
\end{example}

\iosample{}{Name: Alice, Score: 88 \\ Name: Bob, Score: 76 \\ After update: \\ Name: Alice, Score: 95 \\ Name: Bob, Score: 76 \\}

這裡有兩件事值得停下來看清楚。第一，\texttt{update\_score} 除了 \texttt{self} 之外還
多了一個 \texttt{new\_score} 參數，一個方法可以像一般函式一樣，在 \texttt{self} 後面
再接其他參數，呼叫的時候只要補上那些額外的參數就好（\texttt{self} 永遠是自動的），所以
\texttt{s1.update\_score(95)} 實際上是把 \texttt{s1} 當 \texttt{self}、把 \texttt{95}
當 \texttt{new\_score}。

第二，也是這一節最重要的一件事：\texttt{s1.update\_score(95)} 只改了 \texttt{s1} 自己的
\texttt{score}，\texttt{s2} 完全沒被影響，因為 \texttt{s1} 跟 \texttt{s2} 是兩個各自
獨立的物件，各自的屬性存在各自身上，改一個不會動到另一個，即使它們是照著同一個類別造
出來的也一樣。

方法定義時忘記寫 \texttt{self}，是初學者另一個很容易犯的錯：

\begin{example}{python}{(student-broken-method.py) 方法忘記寫 self 參數會怎樣}
class Student:
    def __init__(self, name, score):
        self.name = name
        self.score = score

    def describe():          # 忘記在參數列裡加上 self
        print(f"Name: {self.name}, Score: {self.score}")


s1 = Student("Alice", 88)
s1.describe()
\end{example}

執行後看到的是：

\begin{verbatim}
Traceback (most recent call last):
  File "student-broken-method.py", line 11, in <module>
    s1.describe()
TypeError: Student.describe() takes 0 positional arguments but 1 was given
\end{verbatim}

\keypoint{呼叫 \texttt{s1.describe()} 的時候，Python 一定會自動把 \texttt{s1} 當成
第一個參數傳進去，不管你有沒有在方法定義裡準備好位置接它。上面這個 \texttt{describe}
定義時完全沒寫參數，所以 Python 抱怨「說好 0 個參數，卻收到 1 個」，這個「多收到的
1 個」正是 Python 自動塞進來的 \texttt{s1}。看到這種訊息，第一件事就是回頭檢查方法的
參數列有沒有把 \texttt{self} 放進去。}

\begin{drills}[label=drills:methods-and-self]
\item 這段程式執行後會印出什麼？
\begin{verbatim}
class Car:
    def __init__(self, brand):
        self.brand = brand

    def honk(self):
        print(f"{self.brand} says: Honk honk!")

car = Car("Toyota")
car.honk()
\end{verbatim}
\item 下面這段程式呼叫 \texttt{car.honk()} 時會出現錯誤，錯誤類型是什麼？哪裡需要修正？
\begin{verbatim}
class Car:
    def __init__(self, brand):
        self.brand = brand

    def honk():
        print(f"{self.brand} says: Honk honk!")

car = Car("Toyota")
car.honk()
\end{verbatim}
\item 這段程式執行後會印出什麼？
\begin{verbatim}
class Car:
    def __init__(self, brand):
        self.brand = brand

    def rename(self, new_brand):
        self.brand = new_brand

car1 = Car("Toyota")
car2 = Car("Honda")
car1.rename("BMW")
print(car1.brand, car2.brand)
\end{verbatim}
\end{drills}

\section{方法之間互相呼叫}
\label{sec:method-calls-method}

一個物件通常不會只有一個方法，而一個方法的邏輯，也常常需要用到另一個方法已經算好的
結果。既然 \texttt{self} 代表「這個物件自己」，那在一個方法裡面，只要用
\texttt{self.其他方法名稱(...)} 這種寫法，就能呼叫同一個物件身上的另一個方法，跟你在
外面直接寫 \texttt{s1.其他方法名稱(...)} 呼叫方法的道理完全一樣，只是呼叫的起點換成了
\texttt{self}。

這件事在 \ref{ch:functions} 章學函式的時候就出現過類似的精神：與其把所有邏輯塞進同一個函式裡，不如拆成
幾個各自負責一件事的小函式、讓它們互相呼叫，程式會更好讀、也更好維護；方法之間互相呼叫，
就是同一個精神用在物件身上。

\begin{code}{python}{一個方法呼叫另一個方法}
class ClassName:
    def helper_method(self):
        # 計算一個中間結果
        ...
        return result

    def main_method(self):
        result = self.helper_method()
        # 用 result 做更多事
        ...
\end{code}

\begin{example}{python}{(student-passed.py) 讓 describe 借用 passed 算好的結果}
class Student:
    def __init__(self, name, score):
        self.name = name
        self.score = score

    def passed(self):
        return self.score >= 60

    def describe(self):
        status = "Pass" if self.passed() else "Fail"
        print(f"Name: {self.name}, Score: {self.score}, Status: {status}")


s1 = Student("Alice", 88)
s2 = Student("Charlie", 52)

s1.describe()
s2.describe()
\end{example}

\iosample{}{Name: Alice, Score: 88, Status: Pass \\ Name: Charlie, Score: 52, Status: Fail \\}

\texttt{describe} 這個方法裡用到的 \texttt{self.passed()}，做的事跟你在外面寫
\texttt{s1.passed()} 完全一樣，只不過它是「從物件內部呼叫自己身上的另一個方法」。因為
\texttt{self} 在 \texttt{describe} 執行的當下，代表的正是呼叫它的那個物件（可能是
\texttt{s1}，也可能是 \texttt{s2}），所以 \texttt{self.passed()} 用的自然也是同一個
物件身上的 \texttt{score}，不用把 \texttt{score} 這個值再重新傳一次進去。

這正是把資料跟操作綁在同一個物件上的好處之一：方法之間可以直接互相借用彼此的結果，不用
每次呼叫都重新把資料搬來搬去。

方法互相呼叫時，一樣很容易漏寫 \texttt{self.} 前綴，而且這種錯誤特別容易發生，因為方法
名稱單獨看起來就像一個普通函式：

\begin{example}{python}{(student-broken-method-call.py) 呼叫另一個方法忘記寫 self.\ 會怎樣}
class Student:
    def __init__(self, name, score):
        self.name = name
        self.score = score

    def passed(self):
        return self.score >= 60

    def describe(self):
        status = "Pass" if passed() else "Fail"   # 忘記寫 self.
        print(f"Name: {self.name}, Score: {self.score}, Status: {status}")


s1 = Student("Alice", 88)
s1.describe()
\end{example}

執行後看到的是：

\begin{verbatim}
Traceback (most recent call last):
  File "student-broken-method-call.py", line 15, in <module>
    s1.describe()
    ~~~~~~~~~~~^^
  File "student-broken-method-call.py", line 10, in describe
    status = "Pass" if passed() else "Fail"   # 忘記寫 self.
                       ^^^^^^
NameError: name 'passed' is not defined. Did you mean: 'self.passed'?
\end{verbatim}

\keypoint{在方法裡面呼叫「同一個物件自己的另一個方法」，一定要寫 \texttt{self.}，不能
只寫方法名稱。\texttt{passed()} 這樣裸著呼叫，Python 只會去找一個叫 \texttt{passed} 的
獨立函式或變數，找不到就是 \texttt{NameError}；只有 \texttt{self.passed()} 才是「呼叫
掛在這個物件身上的 \texttt{passed} 方法」。上面這則錯誤訊息其實還很貼心地提示了
\texttt{Did you mean: 'self.passed'?}，看到這種提示，通常就是漏寫 \texttt{self.} 的
訊號。}

\begin{drills}[label=drills:method-calls-method]
\item 這段程式執行後會印出什麼？
\begin{verbatim}
class Car:
    def __init__(self, brand, speed):
        self.brand = brand
        self.speed = speed

    def is_fast(self):
        return self.speed >= 200

    def describe(self):
        label = "Sports Car" if self.is_fast() else "Normal Speed"
        print(f"{self.brand}: {label}")

car = Car("Ferrari", 320)
car.describe()
\end{verbatim}
\item 下面這段程式呼叫 \texttt{car.describe()} 時會出現錯誤，錯誤類型是什麼？哪一行
需要修正？
\begin{verbatim}
class Car:
    def __init__(self, brand, speed):
        self.brand = brand
        self.speed = speed

    def is_fast(self):
        return self.speed >= 200

    def describe(self):
        label = "Sports Car" if is_fast() else "Normal Speed"
        print(f"{self.brand}: {label}")

car = Car("Ferrari", 320)
car.describe()
\end{verbatim}
\item 幫上面的 \texttt{Car} 類別加一個方法 \texttt{warn(self)}：如果
\texttt{self.is\_fast()} 回傳 \texttt{True}，就印出「小心駕駛」；否則什麼都不做。這個
方法要怎麼寫？
\end{drills}

\section{兩個實例互動：讓方法比較另一個物件}
\label{sec:comparing-instances}

到目前為止，方法用到的都只有自己身上的屬性。但很多時候你需要拿兩個物件互相比較，例如
「這兩個學生誰的分數比較高」，這時候一個方法可以像一般函式一樣，在 \texttt{self} 之外
再接一個參數，只是這個參數傳進來的通常\textit{是另一個同類別的物件}，不只是數字或字串。
在方法裡面，你可以直接用點記法讀這個「另一個物件」身上的屬性，跟讀 \texttt{self} 身上
的屬性寫法一模一樣。

\begin{code}{python}{一個方法接住另一個物件當參數}
class ClassName:
    def compare_method(self, other):
        return self.attr > other.attr
\end{code}

\begin{example}{python}{(student-compare.py) 比較兩個學生的分數}
class Student:
    def __init__(self, name, score):
        self.name = name
        self.score = score

    def higher_score_than(self, other):
        return self.score > other.score


s1 = Student("Alice", 88)
s2 = Student("Bob", 76)

if s1.higher_score_than(s2):
    print(f"{s1.name}'s score is higher than {s2.name}'s")
else:
    print(f"{s1.name}'s score is not higher than {s2.name}'s")

print(s2.higher_score_than(s1))
\end{example}

\iosample{}{Alice's score is higher than Bob's \\ False \\}

這裡的 \texttt{other} 不是 Python 的保留字，只是一個常用的慣例名稱，代表「另外那個
物件」，你可以取任何名字，只是 \texttt{other} 讀起來最順、也最不容易跟 \texttt{self}
搞混。

呼叫 \texttt{s1.higher\_score\_than(s2)} 的時候，Python 自動把 \texttt{s1} 傳給
\texttt{self}，把你寫在括號裡的 \texttt{s2} 傳給 \texttt{other}；方法裡面
\texttt{self.score} 讀的是 \texttt{s1} 的分數，\texttt{other.score} 讀的是
\texttt{s2} 的分數。因為 \texttt{self} 跟 \texttt{other} 都是同一個 \texttt{Student}
類別造出來的物件，形狀完全一樣，所以才能用同一套點記法分別取出兩邊的屬性。最後一行把
呼叫者跟參數對調，\texttt{s2.higher\_score\_than(s1)} 這時候 \texttt{self} 換成
\texttt{s2}、\texttt{other} 換成 \texttt{s1}，結果自然也跟著反過來，變成 \texttt{False}。

如果傳進去的 \texttt{other} 不是同一個類別的物件，方法裡想讀 \texttt{other.score} 就
會出問題：

\begin{example}{python}{(student-compare-wrong-type.py) other 傳錯型別會怎樣}
class Student:
    def __init__(self, name, score):
        self.name = name
        self.score = score

    def higher_score_than(self, other):
        return self.score > other.score


s1 = Student("Alice", 88)
s1.higher_score_than(76)
\end{example}

執行後看到的是：

\begin{verbatim}
Traceback (most recent call last):
  File "student-compare-wrong-type.py", line 11, in <module>
    s1.higher_score_than(76)
    ~~~~~~~~~~~~~~~~~~~~^^^^
  File "student-compare-wrong-type.py", line 7, in higher_score_than
    return self.score > other.score
                         ^^^^^^^^^^^
AttributeError: 'int' object has no attribute 'score'
\end{verbatim}

\keypoint{\texttt{76} 這個整數身上根本沒有 \texttt{score} 這個屬性，所以
\texttt{other.score} 一執行就抱怨 \texttt{AttributeError}。呼叫
\texttt{s1.higher\_score\_than(76)} 的當下，Python 並不會先檢查 \texttt{76} 是不是一個
\texttt{Student}，要等到方法裡真的用到 \texttt{other.score} 那一行才會發現不對，這也是
為什麼呼叫這種「接另一個物件當參數」的方法時，要自己留意傳進去的東西是不是正確的類別。}

順帶一提，你可能會在其他人寫的程式碼裡，看到方法參數旁邊多了一小段像
\texttt{other: "Student"} 這樣的文字，這叫「型別提示」，用來跟看程式碼的人（以及編輯器）
說明「這個參數應該是一個 \texttt{Student} 物件」。寫這一行程式碼的當下，
\texttt{Student} 這個類別本身還沒定義完成，Python 還看不到這個名字，所以型別提示才會
用引號包起來，先當一個字串占位。

這只是給人看的說明文字，Python 執行時並不會真的去檢查傳進來的東西是不是
\texttt{Student}，也不會因此擋掉前面示範的那個 \texttt{AttributeError}。這一章不深入
型別提示的細節，這裡知道有這種寫法、看得懂在做什麼就好。

\begin{code}{python}{幫方法參數加上型別提示（選用寫法）}
class ClassName:
    def compare_method(self, other: "ClassName") -> bool:
        return self.attr > other.attr
\end{code}

\begin{drills}[label=drills:comparing-instances]
\item 這段程式執行後會印出什麼？
\begin{verbatim}
class Book:
    def __init__(self, title, pages):
        self.title = title
        self.pages = pages

    def thicker_than(self, other):
        return self.pages > other.pages

b1 = Book("The Three-Body Problem", 302)
b2 = Book("The Little Prince", 96)
print(b1.thicker_than(b2))
print(b2.thicker_than(b1))
\end{verbatim}
\item 這段程式碼執行 \texttt{b1.thicker\_than(100)} 時會發生什麼事？錯誤類型是什麼？
為什麼？
\begin{verbatim}
class Book:
    def __init__(self, title, pages):
        self.title = title
        self.pages = pages

    def thicker_than(self, other):
        return self.pages > other.pages

b1 = Book("The Three-Body Problem", 302)
b1.thicker_than(100)
\end{verbatim}
\item 如果要幫 \texttt{thicker\_than} 的 \texttt{other} 參數加上型別提示，提示
\texttt{other} 應該是一個 \texttt{Book}，這一行方法定義要怎麼寫？為什麼型別要用引號
包起來？
\end{drills}

\section{兩種屬性：instance attribute 與 class attribute}
\label{sec:instance-vs-class-attr}

前面用到的 \texttt{name}、\texttt{score}，都是寫在 \texttt{\_\_init\_\_} 裡、用
\texttt{self.} 存起來的屬性，每個物件各自有一份，改一個不會動到另一個，這種屬性叫
\textbf{instance attribute}（實例屬性）。
\index{instance attribute（實例屬性）}
但有些資料其實所有物件都該共用同一份，例如
「這個班上所有學生的學校名稱」，沒必要每個 \texttt{Student} 物件各自存一份一模一樣的
資料。這種寫在類別本體裡、\texttt{\_\_init\_\_} 外面的屬性，叫做 \textbf{class
attribute}（類別屬性），所有這個類別造出來的物件預設都共用同一份。
\index{class attribute（類別屬性）}

\begin{code}{python}{定義 class attribute}
class ClassName:
    shared_attr = value   # 寫在 __init__ 外面，屬於類別本體

    def __init__(self, param):
        self.own_attr = param
\end{code}

這聽起來簡單，但實際用起來有個地方很容易搞混：如果你透過某個物件去改這個共用屬性，改到
的到底是「這個物件自己的一份」還是「大家共用的那一份」？直接看三個步驟：

\begin{example}{python}{(student-school.py) instance attribute 與 class attribute 的差別}
class Student:
    school_name = "Yangming Middle School"   # 類別屬性

    def __init__(self, name, score):
        self.name = name
        self.score = score


s1 = Student("Alice", 88)
s2 = Student("Bob", 76)

# 步驟一：兩個實例目前共用同一份類別屬性
print("Step 1:", s1.school_name, s2.school_name)

# 步驟二：只改 s1 自己的 school_name
s1.school_name = "Yangming High School"
print("Step 2:", s1.school_name, s2.school_name)

# 步驟三：直接改動類別本身的 school_name
Student.school_name = "Yangming Experimental School"
print("Step 3:", s1.school_name, s2.school_name)
\end{example}

\iosample{}{Step 1: Yangming Middle School Yangming Middle School \\ Step 2: Yangming High School Yangming Middle School \\ Step 3: Yangming High School Yangming Experimental School \\}

一步一步看：

\begin{enumerate}
\item \texttt{s1} 跟 \texttt{s2} 都還沒被個別改過，兩個都直接讀到類別共用的那一份，
所以一樣。
\item \texttt{s1.school\_name = "Yangming High School"} 這一行，實際上是幫
\texttt{s1} 這個物件\textit{另外造一份自己專屬的} \texttt{school\_name}，不是真的去
「修改」類別共用的那一份。之後 \texttt{s1.school\_name} 會優先讀這份自己的，不再去看
類別共用的那份；\texttt{s2} 完全沒被動到，還是讀類別共用的那份。
\item \texttt{Student.school\_name = "..."} 才是真的改動類別本身共用的那份，這時候
\texttt{s2} 因為從來沒被個別覆寫過，會跟著一起變；但 \texttt{s1} 因為在步驟二已經有了
自己專屬的一份，就不會再跟著類別變動了。
\end{enumerate}

這裡有個轉折值得先講清楚：前面說類別是「表格設計」、物件是「照著設計造出來的東西」，
兩者是不同角色；但步驟三的 \texttt{Student.school\_name = "..."}，卻是直接對
\texttt{Student} 這個類別本身賦值。這是因為在 Python 裡，類別本身其實也是一個物件，
只是它是一個特殊的、專門用來造出其他物件的物件。

這就是為什麼你可以像對待一般物件一樣，直接對 \texttt{Student} 這個類別本身設定像
\texttt{school\_name} 這樣的屬性：「表格設計」這張紙本身，原來在 Python 眼裡也是可以
被貼名字牌、可以被賦值的一個東西，不是只有照著它造出來的那些物件才算數。

\keypoint{這跟你在第 \ref{ch:variables} 章學過的「指定是貼名字牌」是同一套邏輯，只是這裡多了「這個名字牌
是貼在『類別』這個物件上，還是貼在『個別實例』這個物件上」的差別。\texttt{s1.school\_name
= "..."} 是幫 \texttt{s1} 這個物件新貼一張屬於它自己的名字牌，蓋過原本「找不到就往類別
身上找」的預設路徑；\texttt{Student.school\_name = "..."} 才是真的改動貼在類別這個物件
本身上的那張名字牌。}

\begin{drills}[label=drills:instance-vs-class-attr]
\item 這段程式執行後會印出什麼？
\begin{verbatim}
class Book:
    genre = "Novel"

    def __init__(self, title):
        self.title = title

b1 = Book("The Three-Body Problem")
b2 = Book("Harry Potter")
b1.genre = "Sci-Fi Novel"
print(b1.genre, b2.genre)
\end{verbatim}
\item 這段程式執行後會印出什麼？
\begin{verbatim}
class Book:
    genre = "Novel"

    def __init__(self, title):
        self.title = title

b1 = Book("The Three-Body Problem")
b2 = Book("Harry Potter")
Book.genre = "Translated Literature"
print(b1.genre, b2.genre)
\end{verbatim}
\item 這段程式執行後會印出什麼？
\begin{verbatim}
class Book:
    genre = "Novel"

    def __init__(self, title):
        self.title = title

b1 = Book("The Three-Body Problem")
b2 = Book("Harry Potter")
b1.genre = "Sci-Fi Novel"
Book.genre = "Translated Literature"
print(b1.genre, b2.genre)
\end{verbatim}
\end{drills}

\section{資料跟操作為什麼要綁在一起：方法與獨立函式的差別}
\label{sec:methods-vs-functions}

你可能會想：既然方法說穿了也是一段程式碼，那為什麼不乾脆繼續寫成 \ref{ch:functions} 章教過的獨立函式，
把資料當參數傳進去就好，何必特地定義一個類別？

答案跟「把資料跟操作綁在一起」這件事本身有關：資料如果放在物件外面，每次呼叫函式，都要
靠你自己記得把正確配對的那幾個值一起傳進去；資料一多、呼叫的地方一多，光靠人工記憶去
配對，出錯的機會就會越來越大。

\begin{example}{python}{(functions-vs-methods-plain.py) 資料跟函式分開寫}
def describe_student(name, score):
    print(f"Name: {name}, Score: {score}")


name1, score1 = "Alice", 88
name2, score2 = "Bob", 76

describe_student(name1, score1)
describe_student(name2, score2)

# 打錯了：把 Bob 的分數跟 Alice 的姓名配在一起傳入
describe_student(name1, score2)
\end{example}

\iosample{}{Name: Alice, Score: 88 \\ Name: Bob, Score: 76 \\ Name: Alice, Score: 76 \\}

最後一行印出「Name: Alice, Score: 76」，這是一個貨真價實的 bug，但 Python 完全不會抱怨，
因為 \texttt{name1} 跟 \texttt{score2} 各自都是型別正確、內容合法的值，只是配對配錯了。
這種錯誤沒有錯誤訊息可以讀，只能靠你自己盯著輸出、發現「怎麼分數兜不起來」才抓得到，比
前面幾節那些會直接跳出 \texttt{AttributeError}、\texttt{NameError} 的錯誤更難察覺。

改用類別，把 \texttt{name} 跟 \texttt{score} 綁進同一個 \texttt{Student} 物件之後，
這種配對錯誤在結構上就不可能發生：

\begin{example}{python}{(functions-vs-methods-bundled.py) 資料跟方法綁在同一個物件裡}
class Student:
    def __init__(self, name, score):
        self.name = name
        self.score = score

    def describe(self):
        print(f"Name: {self.name}, Score: {self.score}")


s1 = Student("Alice", 88)
s2 = Student("Bob", 76)

s1.describe()
s2.describe()
\end{example}

\iosample{}{Name: Alice, Score: 88 \\ Name: Bob, Score: 76 \\}

呼叫 \texttt{s1.describe()} 的時候，方法裡的 \texttt{self.name}、\texttt{self.score}
永遠讀的是同一個 \texttt{s1} 身上的資料，不需要你在呼叫的當下重新指定一次「這次是哪個
姓名配哪個分數」，因為姓名跟分數從 \texttt{Student} 物件被建立的那一刻起，就已經綁在
同一個地方了。想印出 \texttt{s2} 的摘要，就只能呼叫 \texttt{s2.describe()}，沒有機會
不小心把 \texttt{s1} 的分數混進 \texttt{s2} 的姓名裡，因為根本沒有一個步驟需要你手動把
兩者重新配對。

\stylepoint{獨立函式不是不好，\ref{ch:functions} 章學過的函式在很多情境下依然是最直接的寫法。但當一批
資料本來就該綁在同一份、還會被重複造出很多份、又需要好幾個操作圍著它打轉時，把資料跟
操作一起包進一個類別，能讓「這些值本來就是同一份資料」這件事變成程式結構的一部分，不用
只靠你自己在每次呼叫時記得把值配對正確。}

\section{為什麼要用物件組織程式}
\label{sec:why-objects}

前一節看到的是一筆資料的情境；這一節把場景拉大：如果要同時記錄好幾筆長得很像的資料，
用跟不用物件的差別會更明顯。直接對照一次同一個任務最清楚，假設你要記錄三個
學生的姓名跟分數，不用物件的話，可能會寫成一個 dict 疊 dict：

\begin{example}{python}{(record-students-plain.py) 不用物件記錄多筆學生資料}
students = {
    "Alice": {"name": "Alice", "score": 88},
    "Bob": {"name": "Bob", "score": 76},
    "Charlie": {"name": "Charlie", "score": 92},
}

for key in students:
    info = students[key]
    print(f"Name: {info['name']}, Score: {info['score']}")

# 幫 Alice 加 5 分；記得外層的 key 跟內層的 "score" key 都要對上
students["Alice"]["score"] += 5
print(f"Name: {students['Alice']['name']}, Score: {students['Alice']['score']}")
\end{example}

\iosample{}{Name: Alice, Score: 88 \\ Name: Bob, Score: 76 \\ Name: Charlie, Score: 92 \\ Name: Alice, Score: 93 \\}

能動，但每次存取都要記得「先用姓名當 key 找到那一份資料，再用 \texttt{"score"} 這個字串
去內層找分數」，兩層 key 名稱只要哪裡拼錯一個字，Python 在你寫錯的當下不會抱怨，通常要
等到真的執行到那一行才出錯，很容易漏改、漏查。改用前面幾節寫好的 \texttt{Student}
類別，同樣的任務長這樣：

\begin{example}{python}{(record-students-objects.py) 改用 Student 物件記錄多筆學生資料}
class Student:
    def __init__(self, name, score):
        self.name = name
        self.score = score

    def describe(self):
        print(f"Name: {self.name}, Score: {self.score}")


students = [
    Student("Alice", 88),
    Student("Bob", 76),
    Student("Charlie", 92),
]

for student in students:
    student.describe()

# 直接用屬性存取幫 Alice 加 5 分，不用記兩層 key 名稱
students[0].score += 5
students[0].describe()
\end{example}

\iosample{}{Name: Alice, Score: 88 \\ Name: Bob, Score: 76 \\ Name: Charlie, Score: 92 \\ Name: Alice, Score: 93 \\}

兩段程式印出的結果一模一樣，但 \texttt{Student} 版本每個學生只是 list 裡的一個元素，
存取屬性用 \texttt{.name}、\texttt{.score} 這種固定的點記法，不用自己維護「哪個字串是
哪個欄位的 key」這種容易對不上的細節；要印出摘要，直接呼叫 \texttt{describe()}，不用
每次都重新拼一次那行 f-string。資料筆數一多、要做的操作一多，這個差距只會越拉越大。

\stylepoint{不是所有情況都該定義類別。如果你要處理的東西很簡單，只有一兩個值、不會被
重複造出很多份、也不需要自己的方法，直接用變數或 tuple 可能比特地寫一個類別更省事。
物件是用來處理「資料形狀重複出現、還帶著一些相關動作」的情況，不是每個任務的預設答案。}

前面十四章你已經很熟怎麼用內建的物件；這一章你多會了一整套技能：自己設計一種新的物件
形狀、讓方法之間互相合作、讓兩個物件互相比較，還有為什麼值得把資料跟操作綁在一起。下一章
會告訴你，其實「一個類別可以借用另一個類別已經寫好的東西，再加上自己的新花樣」這件事，
跟這一章學的邏輯是同一套延伸，沒有想像中神祕。

\section*{練習解答}
\addcontentsline{toc}{section}{練習解答}

\paragraph{第~\ref{sec:class-init}~節}
\begin{enumerate}
\item 會印出：notebook 45。
\item 印出 \texttt{item.price} 時會出現 \texttt{AttributeError}，訊息是
\texttt{'Item' object has no attribute 'price'}，因為 \texttt{price = price} 漏了
\texttt{self.}，只是 \texttt{\_\_init\_\_} 裡一個執行完就消失的區域變數，沒有真的存到
物件身上。
\item 把 \texttt{\_\_init\_\_} 改成
\texttt{def \_\_init\_\_(self, name, price, in\_stock):}，裡面補上
\texttt{self.in\_stock = in\_stock}；建立實例時寫成
\texttt{Item("pencil", 10, True)}。
\end{enumerate}

\paragraph{第~\ref{sec:methods-and-self}~節}
\begin{enumerate}
\item 會印出：Toyota says: Honk honk!
\item 錯誤類型是 \texttt{TypeError}，訊息類似
\texttt{Car.honk() takes 0 positional arguments but 1 was given}，因為呼叫
\texttt{car.honk()} 時 Python 會自動傳入 \texttt{car} 當第一個參數，但
\texttt{honk} 的定義完全沒有參數可以接；修正方式是把 \texttt{def honk():} 改成
\texttt{def honk(self):}。
\item 會印出：BMW Honda。\texttt{car1.rename("BMW")} 只改了 \texttt{car1} 自己的
\texttt{brand}，\texttt{car2} 是完全獨立的另一個物件，不會被影響。
\end{enumerate}

\paragraph{第~\ref{sec:method-calls-method}~節}
\begin{enumerate}
\item 會印出：Ferrari: Sports Car。
\item 錯誤類型是 \texttt{NameError}，訊息類似
\texttt{name 'is\_fast' is not defined. Did you mean: 'self.is\_fast'?}，因為
\texttt{describe} 裡呼叫 \texttt{is\_fast()} 時漏寫了 \texttt{self.}，Python 只會去找
一個獨立存在、叫 \texttt{is\_fast} 的名字，找不到就丟這個錯；修正方式是把
\texttt{is\_fast()} 改成 \texttt{self.is\_fast()}。
\item 大致寫法：
\begin{verbatim}
def warn(self):
    if self.is_fast():
        print("小心駕駛")
\end{verbatim}
\end{enumerate}

\paragraph{第~\ref{sec:comparing-instances}~節}
\begin{enumerate}
\item 會印出：True，接著 False。\texttt{b1.thicker\_than(b2)} 裡 \texttt{self} 是
\texttt{b1}（302 頁）、\texttt{other} 是 \texttt{b2}（96 頁），302 大於 96 所以是
\texttt{True}；反過來呼叫 \texttt{b2.thicker\_than(b1)} 時 \texttt{self} 換成
\texttt{b2}，96 不大於 302，所以是 \texttt{False}。
\item 會出現 \texttt{AttributeError}，訊息類似
\texttt{'int' object has no attribute 'pages'}，因為 \texttt{100} 這個整數身上沒有
\texttt{pages} 這個屬性，方法裡執行到 \texttt{other.pages} 時就出錯，Python 不會事先
檢查傳進去的 \texttt{100} 是不是一個 \texttt{Book}。
\item 寫成
\texttt{def thicker\_than(self, other: "Book") -> bool:}；型別要用引號包起來，是因為
寫這一行程式碼的當下，\texttt{Book} 這個類別本身還沒定義完成，Python 還看不到這個名字，
只能先用字串占位。
\end{enumerate}

\paragraph{第~\ref{sec:instance-vs-class-attr}~節}
\begin{enumerate}
\item 會印出：Sci-Fi Novel Novel。\texttt{b1.genre = "Sci-Fi Novel"} 幫 \texttt{b1} 造了一份
自己專屬的 \texttt{genre}，\texttt{b2} 沒被改過，還是讀類別共用的那份。
\item 會印出：Translated Literature Translated Literature。\texttt{b1}、\texttt{b2} 都還沒被個別覆寫過，
\texttt{Book.genre = "Translated Literature"} 直接改動的是類別共用的那份，兩個都跟著變。
\item 會印出：Sci-Fi Novel Translated Literature。\texttt{b1} 在第一行已經有了自己專屬的
\texttt{genre}，之後改類別本身的 \texttt{genre} 不會再影響到它；\texttt{b2} 因為從來
沒被個別覆寫過，會跟著類別一起變。
\end{enumerate}

\chapter{繼承與更多 OOP}
\label{ch:inheritance}

上一章你已經學會怎麼用 class 打造自己的物件，寫 \texttt{\_\_init\_\_}、定義屬性、寫方法。這
一章接著問幾個很自然會冒出來的問題。第一個問題是：如果你已經寫好一個 class，現在想要一個
「幾乎一樣、只是多一點東西」的新 class，難道要把整段程式碼複製貼上再改嗎？這一章前幾節要教你
一個更好的做法，這就是「繼承」，讓一個新 class 直接站在既有 class 的肩膀上，只寫「多出來、不
一樣」的那一部分；順便你也會看到繼承其實可以疊好幾層，還會認識另一種跟繼承並列、常常拿來互相
比較的做法，叫做「組合」。

第二個問題是：你有沒有發現，把自己寫的物件直接 \texttt{print()} 出
來，結果常常長得很奇怪，看不出裡面裝了什麼？接下來兩節要教你兩個特殊方法，讓物件學會「自我
介紹」：一種是講給一般使用者聽的版本，一種是講給寫程式的人（包括未來的你自己）聽、方便除錯的
版本。

最後，這一章要帶你回頭重新看一次第 \ref{ch:variables} 章講過的「名字牌模型」。物件比第 \ref{ch:variables} 章的數字、字串
複雜得多，複製一個物件、比較兩個物件是不是「一樣」，這些問題在物件的世界裡有一些新的細節值得
留意，最後一節會用一個你在第 \ref{ch:functions} 章就見過的老朋友（購物車）重新示範一次。

整章的份量刻意抓得比較
輕，這是點到為止的一章，不是要把你訓練成 OOP 專家，你只要知道這些工具存在、大概什麼時候會用
到就夠了，之後如果工作上真的需要深入某一塊，再回頭查資料也不遲。

\section{繼承：新 class 自動擁有父類別的一切}
\label{sec:inheritance-intro}

想像你已經用第 \ref{ch:objects} 章學的方式，定義好一個 \texttt{Animal} class，裡面存了動物的名字、也有一個
\texttt{speak()} 方法讓動物「發出聲音」。現在你想再定義一個 \texttt{Dog} class，狗當然也是
動物，也有名字、也會發出聲音，但狗還有一些自己專屬的行為（例如去把球撿回來）。如果把
\texttt{Animal} 整段程式碼複製一份、貼進 \texttt{Dog}、再手動加東西，你等於在維護兩份幾乎一
模一樣的程式碼，以後 \texttt{Animal} 改一個地方，你得記得兩邊都要改，這是很容易漏掉的坑。

Python 給了一個更好的做法，叫做「繼承」（inheritance）：讓 \texttt{Dog} 直接站在
\texttt{Animal} 的肩膀上，寫一次「\texttt{Dog} 是一種 \texttt{Animal}」，\texttt{Dog} 就自動
擁有 \texttt{Animal} 已經有的一切，你完全不用重寫。
\index{繼承}

\begin{code}{python}{定義子類別}
class ParentClass:
    def __init__(self, param):
        self.attr = param

    def method(self):
        print(self.attr)


class ChildClass(ParentClass):
    pass
\end{code}

括號裡放的是「父類別」（parent class，也叫 base class）的名字，這一行的意思是「\texttt{Child
Class} 繼承自 \texttt{ParentClass}」。
\index{父類別}
反過來，\texttt{ChildClass} 就叫「子類別」（child
class，也叫 derived class）：繼承自父類別、可以使用父類別的屬性和方法，也可以新增或覆寫。
\index{子類別}
上面這個 \texttt{ChildClass} 裡面除了一個
\texttt{pass}（表示「這裡先什麼都不做」）之外什麼都沒寫，但因為它繼承了 \texttt{ParentClass}，
它自動擁有 \texttt{ParentClass} 的 \texttt{\_\_init\_\_} 跟 \texttt{method}，一行都不用自己
重寫。

\begin{example}{python}{(dog-inherit.py) 狗是一種動物}
class Animal:
    def __init__(self, name):
        self.name = name

    def speak(self):
        print(f"{self.name} made a sound.")


class Dog(Animal):
    pass


dog = Dog("Max")
dog.speak()
\end{example}

\iosample{}{Max made a sound. \\}

\texttt{Dog} class 的本體裡只有一個 \texttt{pass}，沒有寫任何 \texttt{\_\_init\_\_}、也沒有
寫任何 \texttt{speak()}，但 \texttt{Dog("Max")} 一樣能正常建立物件、\texttt{dog.speak()} 一樣
能正常執行。這是因為 \texttt{class Dog(Animal):} 這一行，已經讓 \texttt{Dog} 把 \texttt{Animal}
的 \texttt{\_\_init\_\_} 跟 \texttt{speak()} 全部「白拿」過來了。

\keypoint{繼承的第一印象請記住這句話：「什麼都不用寫，子類別就自動擁有父類別的一切。」之後
你會學到子類別可以覆寫、可以新增自己的東西，但那些都是「額外做的事」，不是繼承本身要求你做的
事。繼承本身，就是白拿。}

\section{覆寫與 \texttt{super()}：讓子類別做點不一樣的事}
\label{sec:inheritance-basics}

光是「白拿」還不夠，你想要的通常是「大部分一樣，但有幾個地方不一樣」的 \texttt{Dog}：狗的
叫聲跟一般動物不一樣，狗還有自己專屬的行為（撿球），而且狗在建立的時候，除了名字之外可能還
想多記一個「品種」。這一節要教你三件事：怎麼覆寫（override）父類別的方法、怎麼在子類別裡
新增父類別沒有的方法、還有子類別自己定義 \texttt{\_\_init\_\_} 時該怎麼處理。

\begin{code}{python}{覆寫方法、新增方法}
class ChildClass(ParentClass):
    def __init__(self, param1, param2):
        super().__init__(param1)
        self.extra = param2

    def method(self):        # 覆寫父類別的同名方法
        print("the child class's own behavior")

    def extra_method(self):  # 父類別沒有的新方法
        print(self.extra)
\end{code}

\texttt{def method(self):} 這一段名字跟父類別的方法一模一樣，Python 的規則很直接：子類別
如果定義了跟父類別同名的方法，呼叫的時候一律以子類別自己這份為準，這就叫「覆寫」。
\index{覆寫}
至於
\texttt{extra\_method}，父類別根本沒有這個方法，這純粹是子類別自己多出來的東西。比較特別的是
\texttt{\_\_init\_\_} 裡的 \texttt{super().\_\_init\_\_(param1)}：\texttt{super()} 代表「父
類別」，這一行的意思是「先照父類別原本的 \texttt{\_\_init\_\_} 做一次該做的事」，之後
才接著做子類別自己想額外做的事（這裡是多存一個 \texttt{extra}）。
\index{super()}

\begin{example}{python}{(dog-override.py) 狗有自己的叫法跟拿手絕活}
class Animal:
    def __init__(self, name):
        self.name = name

    def speak(self):
        print(f"{self.name} made a sound.")


class Dog(Animal):
    def __init__(self, name, breed):
        super().__init__(name)
        self.breed = breed

    def speak(self):
        print(f"{self.name} barked!")

    def fetch(self):
        print(f"{self.name} ran off to fetch the ball.")


dog = Dog("Max", "Shiba Inu")
dog.speak()
dog.fetch()
print(dog.breed)
\end{example}

\iosample{}{Max barked! \\ Max ran off to fetch the ball. \\ Shiba Inu \\}

留意 \texttt{Dog.\_\_init\_\_} 接收 \texttt{name} 跟 \texttt{breed} 兩個參數，但只把
\texttt{name} 這一個丟給 \texttt{super().\_\_init\_\_(name)}，這一行實際上就是在跑
\texttt{Animal.\_\_init\_\_(self, name)}，幫你把 \texttt{self.name = name} 設好，
\texttt{self.breed = breed} 才是 \texttt{Dog} 自己額外多做的事。

\keypoint{\texttt{super().\_\_init\_\_(...)} 不是「順便」，是「你要，還是要跑一次父類別的
設定」的明確表態。如果子類別自己定義了 \texttt{\_\_init\_\_}，卻沒有呼叫
\texttt{super().\_\_init\_\_()}，父類別原本 \texttt{\_\_init\_\_} 裡做的事就\textbf{完全不會
執行}。很多初學者會誤以為「反正是繼承，父類別的初始化應該會自動疊上去」，但 Python 不會
幫你這樣猜，子類別的 \texttt{\_\_init\_\_} 一旦自己定義，就整個蓋過父類別那一份，除非你自己
明講「還要跑一次」。}

如果不小心漏寫了 \texttt{super().\_\_init\_\_()}，會發生什麼事？

\begin{example}{python}{(dog-broken.py) 忘記呼叫 super().\_\_init\_\_()}
class Animal:
    def __init__(self, name):
        self.name = name

    def speak(self):
        print(f"{self.name} made a sound.")


class Dog(Animal):
    def __init__(self, breed):
        self.breed = breed  # 忘記呼叫 super().__init__()，name 永遠不會被設定


dog = Dog("Shiba Inu")
dog.speak()
\end{example}

存檔執行後，程式不會印出叫聲，而是直接跳出以下這樣的錯誤：

\begin{verbatim}
Traceback (most recent call last):
  File "dog-broken.py", line 12, in <module>
    dog.speak()
  File "dog-broken.py", line 5, in speak
    print(f"{self.name} made a sound.")
AttributeError: 'Dog' object has no attribute 'name'
\end{verbatim}

由下往上讀：\texttt{AttributeError: 'Dog' object has no attribute 'name'} 是說「這個
\texttt{Dog} 物件身上根本沒有 \texttt{name} 這個屬性」，因為 \texttt{Dog.\_\_init\_\_} 只
設定了 \texttt{self.breed}，從來沒有跑到 \texttt{Animal.\_\_init\_\_} 裡
\texttt{self.name = name} 那一行，物件建立完之後身上自然沒有 \texttt{name}。修法很直接：在
\texttt{Dog.\_\_init\_\_} 裡補上 \texttt{super().\_\_init\_\_(name)}（別忘了 \texttt{Dog}
的 \texttt{\_\_init\_\_} 參數列也要加回 \texttt{name}）。

順帶一提，如果你手上有一個 list，裡面同時裝著幾隻 \texttt{Dog} 跟幾隻普通的 \texttt{Animal}，
你可以用同一段迴圈，對每個元素都呼叫 \texttt{.speak()}，不用管每一個實際上是哪個 class，只
要它們都有同名的 \texttt{speak()} 方法，Python 就會自動呼叫到各自正確的那一份（狗會汪汪叫、
一般動物會發出「一些聲音」）。這個好處以後你會經常用到，這裡先知道它存在就好。

\begin{drills}[label=drills:inheritance-basics]
\item 承接上面的 \texttt{Animal}／\texttt{Dog}，如果再定義一個
\texttt{class Cat(Animal): pass}（沒有覆寫任何東西），執行
\texttt{Cat("Luna").speak()} 會印出什麼？
\item 如果 \texttt{Dog} 覆寫了 \texttt{speak()}，執行
\texttt{Dog("Rex", "Mixed Breed").speak()} 會印出什麼？
\item 如果有人定義了 \texttt{class Cat(Animal):}，裡面的 \texttt{\_\_init\_\_} 是
\texttt{def \_\_init\_\_(self, name, color): self.color = color}（忘記呼叫
\texttt{super().\_\_init\_\_(name)}），接著執行 \texttt{Cat("Luna", "orange").speak()}，
你覺得會發生什麼事？
\end{drills}

\section{多層繼承：子類別也可以再被繼承}
\label{sec:multi-level-inheritance}

前兩節看到的都是「一層」繼承：\texttt{Dog} 繼承自 \texttt{Animal}。但繼承這件事不會只能疊
一層，\texttt{Dog} 本身也是一個普通的 class，只要語法上合法，任何 class 都可以再被別人繼承，
包括已經繼承過別人的 class。也就是說，你可以再定義一個 \texttt{Puppy}（幼犬）繼承自
\texttt{Dog}，這時候 \texttt{Puppy} 不只擁有 \texttt{Dog} 自己定義的東西，連 \texttt{Dog}
從 \texttt{Animal} 那邊繼承來的東西，也會一路往下傳給 \texttt{Puppy}。這一節不需要新的片語
模板，繼承的語法你已經會了，直接看範例最快。

\begin{example}{python}{(dog-puppy.py) 幼犬是一種狗，狗是一種動物}
class Animal:
    def __init__(self, name):
        self.name = name

    def speak(self):
        print(f"{self.name} made a sound.")


class Dog(Animal):
    def __init__(self, name, breed):
        super().__init__(name)
        self.breed = breed

    def speak(self):
        print(f"{self.name} barked!")

    def fetch(self):
        print(f"{self.name} ran off to fetch the ball.")


class Puppy(Dog):
    def __init__(self, name, breed, age_months):
        super().__init__(name, breed)
        self.age_months = age_months

    def speak(self):
        print(f"{self.name} let out a clumsy little bark.")

    def nap(self):
        print(f"{self.name} got tired and lay down for a nap.")


puppy = Puppy("Buddy", "Shiba Inu", 3)
puppy.speak()
puppy.fetch()
puppy.nap()
print(puppy.breed)
\end{example}

\iosample{}{Buddy let out a clumsy little bark. \\ Buddy ran off to fetch the ball. \\ Buddy got
tired and lay down for a nap. \\ Shiba Inu \\}

這段程式一次示範了三件事：

\begin{itemize}
\item 第一，\texttt{Puppy} 自己又覆寫了一次 \texttt{speak()}，改成
「還不太熟練的汪汪叫聲」，覆寫的規則跟前一節一樣，只是這次覆寫的對象是父類別 \texttt{Dog}
（而 \texttt{Dog} 本身也覆寫過 \texttt{Animal} 的版本）。
\item 第二，\texttt{puppy.fetch()} 正常
執行，印出「跑去把球撿回來了」，但 \texttt{fetch()} 這個方法從頭到尾只在 \texttt{Dog} 裡定義
過一次，\texttt{Puppy} 完全沒有自己寫過，這正是「白拿」的概念延伸到多層繼承的樣子：
\texttt{Puppy} 不只白拿了直接父類別 \texttt{Dog} 的東西，也間接白拿到更上層 \texttt{Animal}
的東西。
\item 第三，\texttt{Puppy.\_\_init\_\_} 裡的 \texttt{super().\_\_init\_\_(name, breed)}
呼叫的是 \texttt{Dog.\_\_init\_\_}，而 \texttt{Dog.\_\_init\_\_} 裡面自己又呼叫了一次
\texttt{super().\_\_init\_\_(name)} 去跑 \texttt{Animal.\_\_init\_\_}，於是 \texttt{self.name}、
\texttt{self.breed}、\texttt{self.age\_months} 三個屬性，分別由三層 class 各自負責設定好自己
那一部分，一層一層往上串起來。
\end{itemize}

\keypoint{\texttt{super()} 每次呼叫，只會跳到「自己的直接父類別」那一層，不會一次跳到最上層的
祖先。多層繼承之所以能一路往上生效，是因為每一層自己的 \texttt{\_\_init\_\_} 都老實呼叫了自己
的 \texttt{super()}，一層接一層疊上去，不是 \texttt{Puppy} 一句話就直接連到 \texttt{Animal}。
如果中間有任何一層忘記呼叫 \texttt{super().\_\_init\_\_()}，這條鏈就會在那一層斷掉，跟上一節
看到的 \texttt{AttributeError} 是同一種坑，只是這次可能發生在鏈的中間，不一定是最底層。}

多層繼承還有一個新手容易踩到的坑：每一層的 \texttt{super().\_\_init\_\_(...)} 要傳的引數，
得照著「直接父類別」的 \texttt{\_\_init\_\_} 簽名走，不是照著最上層祖先的簽名走。

\begin{example}{python}{(puppy-broken.py) 傳給 super() 的引數不齊全}
class Animal:
    def __init__(self, name):
        self.name = name

    def speak(self):
        print(f"{self.name} made a sound.")


class Dog(Animal):
    def __init__(self, name, breed):
        super().__init__(name)
        self.breed = breed

    def speak(self):
        print(f"{self.name} barked!")


class Puppy(Dog):
    def __init__(self, name, breed, age_months):
        super().__init__(name)  # 忘記把 breed 也傳給 super().__init__()
        self.age_months = age_months


puppy = Puppy("Buddy", "Shiba Inu", 3)
\end{example}

存檔執行後，會看到這樣的錯誤：

\begin{verbatim}
Traceback (most recent call last):
  File "puppy-broken.py", line 20, in <module>
    puppy = Puppy("Buddy", "Shiba Inu", 3)
  File "puppy-broken.py", line 16, in __init__
    super().__init__(name)
TypeError: Dog.__init__() missing 1 required positional argument: 'breed'
\end{verbatim}

\texttt{Puppy.\_\_init\_\_} 呼叫 \texttt{super().\_\_init\_\_(...)} 的時候，跳到的是它的直接
父類別 \texttt{Dog} 的 \texttt{\_\_init\_\_}，而 \texttt{Dog.\_\_init\_\_} 需要 \texttt{name}
跟 \texttt{breed} 兩個引數，這裡只傳了 \texttt{name}，Python 就直接在這一行報錯，根本還沒機會
執行到 \texttt{self.age\_months = age\_months} 那一行。修法是把 \texttt{breed} 也一起傳進去：
\texttt{super().\_\_init\_\_(name, breed)}。這個錯誤訊息也是由下往上讀最清楚：先看最下面一行
\texttt{TypeError} 講的是哪個 \texttt{\_\_init\_\_} 缺了哪個引數，再往上對照是哪一層的
\texttt{super()} 呼叫出的錯。

繼承鏈當然可以繼續往下疊（例如再定義一個繼承自 \texttt{Puppy} 的 class），但實務上疊得太深、
太複雜的繼承鏈反而會讓程式碼變得不好讀（想像一下要搞清楚某個方法到底是哪一層定義的，得往上翻
好幾層 class 才找得到）。

這一章的定位是點到為止，你只要知道「繼承可以疊好幾層、每一層各自負責
自己那部分」這個事實就夠了，實際設計要疊幾層、什麼時候該停下來，是之後累積經驗才會抓到手感的
事，這裡不深入。

\begin{drills}[label=drills:multi-level-inheritance]
\item 承接上面的 \texttt{Puppy}，如果 \texttt{Puppy} 沒有覆寫 \texttt{speak()}，執行
\texttt{Puppy("Buddy", "Shiba Inu", 3).speak()} 會印出什麼？
\item 如果再定義一個 \texttt{class Toy(Puppy): pass}（沒有覆寫或新增任何東西），執行
\texttt{Toy("Tiny", "Mixed Breed", 2).fetch()} 會發生什麼事？
\item 如果 \texttt{Puppy.\_\_init\_\_} 寫成
\texttt{def \_\_init\_\_(self, name, breed, age\_months): super().\_\_init\_\_(name)}（漏傳
\texttt{breed}），執行 \texttt{Puppy("Buddy", "Shiba Inu", 3)} 你覺得會發生什麼事？
\end{drills}

\section{繼承 vs 組合：另一種重複利用程式碼的方式}
\label{sec:inheritance-vs-composition}

前面幾節解決的都是同一種問題：「這個新 class 幾乎跟舊的一樣，只是多一點東西」，用繼承來處理。
但繼承有一個前提，它適合的情境是「新 class 是一種舊 class」，例如「狗是一種動物」。如果兩個
class 之間的關係比較像是「這個 class 裡面裝了另一個 class」，而不是「這個 class 是另一個
class 的一種」，繼承就不是對的工具了。這種情況下，Python 常常改用另一種做法，叫「組合」
（composition）：讓一個 class 的屬性直接是另一個 class 的物件。
\index{組合}

\begin{example}{python}{(car-composition.py) 車子裡面裝了一顆引擎}
class Engine:
    def __init__(self, horsepower):
        self.horsepower = horsepower

    def start(self):
        print(f"Engine started, horsepower {self.horsepower}")


class Car:
    def __init__(self, brand, horsepower):
        self.brand = brand
        self.engine = Engine(horsepower)

    def start(self):
        print(f"{self.brand} getting ready to start")
        self.engine.start()


car = Car("Toyota", 150)
car.start()
\end{example}

\iosample{}{Toyota getting ready to start \\ Engine started, horsepower 150 \\}

留意 \texttt{class Car:} 這一行的括號是空的，\texttt{Car} 完全沒有繼承 \texttt{Engine}。
\texttt{Car.\_\_init\_\_} 裡的 \texttt{self.engine = Engine(horsepower)}，是直接造一個
\texttt{Engine} 物件，存成 \texttt{Car} 自己的一個屬性。\texttt{Car.start()} 裡面要用到引擎
的功能時，就透過 \texttt{self.engine.start()} 去呼叫，這就是「組合」：一台車「有一個」引擎，
不是「一種」引擎。

跟前面幾節的 \texttt{Dog} 對照著看會更清楚：\texttt{Dog(Animal)} 表達的是
「狗是一種動物」，適合用繼承；這裡的 \texttt{Car} 跟 \texttt{Engine} 表達的是「車子裡面裝了一
顆引擎」，適合用組合。如果硬要把 \texttt{Car} 寫成繼承 \texttt{Engine}（\texttt{class
Car(Engine):}），語法上不會出錯，但意思會變得很奇怪，等於在說「車子是一種引擎」，這跟現實世界
的認知完全不符。

\keypoint{要用繼承還是組合，可以先問自己一句話：這兩個 class 的關係，講起來比較像「A 是一種
B」，還是「A 裡面有一個 B」？前者適合繼承，後者適合組合。這只是一個入門判準，實務上有些情況
會更微妙，需要更多經驗才能判斷，這裡你只要知道除了繼承之外還有組合這個選項存在，之後看到別人
的程式碼裡某個屬性本身又是另一個自訂物件時，不會覺得陌生就好，不需要現在就能精準判斷每種情境
該選哪一個。}

組合也有一個新手容易搞混的地方：因為 \texttt{Car} 沒有繼承 \texttt{Engine}，\texttt{Engine}
身上的屬性、方法不會像繼承那樣自動「白拿」到 \texttt{Car} 身上，一定要透過 \texttt{self.engine}
這一層才拿得到。

\begin{example}{python}{(car-composition-error.py) 以為組合也會白拿屬性}
class Engine:
    def __init__(self, horsepower):
        self.horsepower = horsepower

    def start(self):
        print(f"Engine started, horsepower {self.horsepower}")


class Car:
    def __init__(self, brand, horsepower):
        self.brand = brand
        self.engine = Engine(horsepower)


car = Car("Toyota", 150)
print(car.horsepower)
\end{example}

存檔執行後，會看到這樣的錯誤：

\begin{verbatim}
Traceback (most recent call last):
  File "car-composition-error.py", line 16, in <module>
    print(car.horsepower)
AttributeError: 'Car' object has no attribute 'horsepower'
\end{verbatim}

\texttt{horsepower} 是存在 \texttt{car.engine.horsepower} 這個位置，不是
\texttt{car.horsepower}，\texttt{Car} 物件本身根本沒有這個屬性。

這正好跟前幾節的繼承對照出
差異：如果 \texttt{Car} 是繼承 \texttt{Engine}（\texttt{class Car(Engine):}），
\texttt{car.horsepower} 會直接可以用，因為繼承會把父類別的屬性、方法整個「白拿」過來；但組合
不會，組合只是把另一個物件存成自己的一個屬性，想用裡面的東西，一定要照著 \texttt{self.engine.
horsepower} 這種「一層一層點進去」的路徑去拿，這也是判斷「這裡該用繼承還是組合」時，除了語意
上「是一種」還是「有一個」之外，另一個看得到的實際差異。

順帶用第 \ref{ch:variables} 章的名字牌模型回頭看一次組合：\texttt{self.engine = Engine(horsepower)} 這一行，
做的事情跟任何一次賦值一樣，是造一個 \texttt{Engine} 物件、把名字牌 \texttt{engine} 貼上去，
只是這張名字牌剛好貼在 \texttt{self} 身上、成為 \texttt{Car} 物件的一個屬性，不是貼在一個
獨立的變數上。

這也解釋了為什麼組合不會白拿屬性：\texttt{car.engine} 這張名字牌指向的是那個
\texttt{Engine} 物件，\texttt{car} 這個名字牌本身指向的是另一個完全不同的 \texttt{Car} 物件，
兩者是分開的兩張名字牌，貼在兩個不同的物件上，想從 \texttt{Car} 走到 \texttt{Engine} 的屬性，
自然要先經過 \texttt{engine} 這張名字牌。

\begin{drills}[label=drills:inheritance-vs-composition]
\item 如果把 \texttt{class Car:} 改成 \texttt{class Car(Engine):}（讓 \texttt{Car} 直接繼承
\texttt{Engine}），並把 \texttt{Car.\_\_init\_\_} 也對應改成呼叫
\texttt{super().\_\_init\_\_(horsepower)}，這時候執行 \texttt{car.horsepower} 會不會出錯？
跟原本組合的寫法比起來，語意上合不合理？
\item 承接上面的 \texttt{car}（原本組合的寫法，\texttt{car.engine} 是一個 \texttt{Engine}
物件），執行 \texttt{car.engine.start()} 會印出什麼？
\end{drills}

\section{讓物件學會自我介紹：\texttt{\_\_str\_\_}}
\label{sec:str-method}

延續前面的 \texttt{Animal}／\texttt{Dog}，如果你直接把一個物件丟進 \texttt{print()}，會
看到什麼？

\begin{example}{python}{(dog-print-ugly.py) 直接印出物件，會發生什麼事}
class Animal:
    def __init__(self, name):
        self.name = name

    def speak(self):
        print(f"{self.name} made a sound.")


class Dog(Animal):
    def __init__(self, name, breed):
        super().__init__(name)
        self.breed = breed

    def speak(self):
        print(f"{self.name} barked!")


dog = Dog("Max", "Shiba Inu")
print(dog)
\end{example}

\iosample{}{<\_\_main\_\_.Dog object at 0x10309ef90> \\}

這行字看起來很醜，也完全看不出這隻狗叫什麼名字，這是 Python 對「不知道該怎麼描述自己」的
物件給的預設寫法：類別名稱、加上這個物件在記憶體裡的位置（那串 \texttt{0x...} 的十六進位
數字，每次執行、每一台機器都可能不一樣，沒有任何實際意義）。你當然不希望使用者看到這種東西，
好消息是 Python 讓你可以自己定義「這個物件被 \texttt{print()} 出來時該長什麼樣子」。

\begin{code}{python}{定義 \_\_str\_\_}
class ClassName:
    def __str__(self):
        return "the content you want shown when this object is printed"
\end{code}

\texttt{\_\_str\_\_} 是一個特殊方法：你自己不會直接呼叫它，當你對這個物件用
\texttt{print()} 或 \texttt{str()} 時，Python 會自動幫你呼叫它，並且把它 \texttt{return}
的字串印出來。留意 \texttt{\_\_str\_\_} 一定要回傳一個字串，不能直接在裡面 \texttt{print}。
\index{特殊方法}
\index{\_\_str\_\_}

\begin{example}{python}{(dog-print-str.py) 幫 Animal 加上 \_\_str\_\_}
class Animal:
    def __init__(self, name):
        self.name = name

    def speak(self):
        print(f"{self.name} made a sound.")

    def __str__(self):
        return f"{self.name} (an animal)"


class Dog(Animal):
    def __init__(self, name, breed):
        super().__init__(name)
        self.breed = breed

    def speak(self):
        print(f"{self.name} barked!")


dog = Dog("Max", "Shiba Inu")
print(dog)
\end{example}

\iosample{}{Max (an animal) \\}

留意這裡 \texttt{\_\_str\_\_} 是定義在 \texttt{Animal} 身上的，\texttt{Dog} 完全沒有自己
定義 \texttt{\_\_str\_\_}，但 \texttt{print(dog)} 一樣印出了好讀的內容。這跟上一節學到的
道理是同一件事：\texttt{Dog} 沒有覆寫 \texttt{\_\_str\_\_}，就直接繼承 \texttt{Animal} 那
一份，白拿。

前面提過 \texttt{\_\_str\_\_} 一定要 \texttt{return} 一個字串，不能直接在裡面
\texttt{print}，如果不小心寫錯會發生什麼事？

\begin{example}{python}{(str-return-mistake.py) \_\_str\_\_ 裡忘記用 return}
class Animal:
    def __init__(self, name):
        self.name = name

    def __str__(self):
        print(f"{self.name} (an animal)")  # 忘記用 return，直接印出來了


animal = Animal("Max")
print(animal)
\end{example}

存檔執行後，畫面上先印出一行「Max (an animal)」，緊接著卻還是跳出錯誤：

\begin{verbatim}
Max (an animal)
Traceback (most recent call last):
  File "str-return-mistake.py", line 10, in <module>
    print(animal)
TypeError: __str__ returned non-string (type NoneType)
\end{verbatim}

這裡最容易讓人困惑的地方是：畫面上明明先印出了一行內容看起來很正常的文字，後面卻還是報錯。
原因是 \texttt{\_\_str\_\_} 裡面那句 \texttt{print(...)} 本身確實執行了、把文字顯示到螢幕上，
但它沒有 \texttt{return} 任何東西，一個函式（或方法）如果沒有寫 \texttt{return}，結果就是
\texttt{None}（第 \ref{ch:functions} 章教過的規則）。\texttt{print(animal)} 外層呼叫 \texttt{\_\_str\_\_} 拿到
的是 \texttt{None}，不是字串，Python 才會丟出 \texttt{TypeError: \_\_str\_\_ returned
non-string (type NoneType)}。修法是把 \texttt{print(...)} 換成 \texttt{return ...}，讓
\texttt{\_\_str\_\_} 真的把字串「送回去」，而不是自己搶著印出來。

\stylepoint{只要你定義的 class 之後有機會被 \texttt{print()} 出來（幾乎所有 class 都有這個
機會，例如除錯的時候），養成習慣幫它寫一個 \texttt{\_\_str\_\_}，讀者（包括除錯時的你自己）
會感謝你，這比每次都要另外寫 \texttt{print(f"...{obj.attr1}...")} 這種一長串屬性堆疊的
程式碼方便多了。}

\section{給寫程式的人看的版本：\texttt{\_\_repr\_\_}}
\label{sec:str-vs-repr}

\texttt{\_\_str\_\_} 是講給「一般使用者」聽的版本，但其實 Python 還有另一個很類似的特殊
方法，叫 \texttt{\_\_repr\_\_}，是講給「寫程式的人」（包括除錯時的你自己）聽的版本。
\index{\_\_repr\_\_}
這一節
沒有新的片語模板可以套，直接看範例最快。

\begin{example}{python}{(dog-print-repr.py) \_\_str\_\_ 與 \_\_repr\_\_ 的分工}
class Animal:
    def __init__(self, name):
        self.name = name

    def speak(self):
        print(f"{self.name} made a sound.")

    def __str__(self):
        return f"{self.name} (an animal)"

    def __repr__(self):
        return f"Animal({self.name!r})"


class Dog(Animal):
    def __init__(self, name, breed):
        super().__init__(name)
        self.breed = breed

    def speak(self):
        print(f"{self.name} barked!")


dog = Dog("Max", "Shiba Inu")
print(dog)
print([dog])
\end{example}

\iosample{}{Max (an animal) \\{} [Animal('Max')] \\}

留意 \texttt{\_\_repr\_\_} 裡這一行 \texttt{f"Animal({self.name!r})"}，比之前看過的 f-string
多了一個 \texttt{!r}，這是 f-string 的「轉換旗標」，意思是對 \texttt{self.name} 呼叫
\texttt{repr()}，而不是預設的 \texttt{str()} 來轉成文字。差別看得出來：字串值用
\texttt{repr()} 轉出來的樣子會帶著引號（例如 \texttt{'Max'}），用 \texttt{str()} 轉出來則
不帶引號（例如 \texttt{Max}）。這正好呼應 \texttt{\_\_repr\_\_} 的用途，給開發者看、要盡量
精確，帶引號才看得出「這是一個字串」，跟給使用者看、講求好讀的 \texttt{\_\_str\_\_} 不一樣。

同一個 \texttt{dog} 物件，兩行印出來的內容完全不一樣：\texttt{print(dog)} 走的是
\texttt{\_\_str\_\_}，還是原本那句好讀的「Max (an animal)」；但 \texttt{print([dog])} 印的
是一個 list，Python 在印 list 裡的每個元素時，用的是 \texttt{\_\_repr\_\_}，不是
\texttt{\_\_str\_\_}，所以你看到的是 \texttt{Animal('Max')} 這種格式。如果你在互動式直譯器
（第 \ref{ch:introduction} 章教過的 \texttt{>>>} 那個模式）直接打 \texttt{dog} 這個變數名稱、不加
\texttt{print()}，看到的也會是同樣這行 \texttt{Animal('Max')}，因為直譯器要把一個變數的值
「秀」給你看的時候，走的正是 \texttt{\_\_repr\_\_}，跟丟進 list 裡印出來是同一條路。

你可能注意到，這裡印出來是 \texttt{Animal('Max')}，而不是 \texttt{Dog('Max')}，雖然
\texttt{dog} 其實是一個 \texttt{Dog} 物件。這是因為 \texttt{\_\_repr\_\_} 是定義在
\texttt{Animal} 身上、\texttt{Dog} 直接繼承了它，而這裡的寫法是手動寫死了 \texttt{"Animal("}
這幾個字，沒有另外去問「我實際上是哪個 class」。

這是刻意選的最簡單寫法，老實說出它的限制：
如果想讓子類別印出來也顯示自己的類別名稱，是有辦法做到的，但那已經超出這一章想教的範圍，這裡
你只要知道有這回事、不用被這個小細節絆住就好。

如果一個物件只定義了 \texttt{\_\_repr\_\_}，完全沒有定義 \texttt{\_\_str\_\_}，
\texttt{print()} 又會發生什麼事？

\begin{example}{python}{(repr-fallback.py) 只定義 \_\_repr\_\_，沒定義 \_\_str\_\_}
class Animal:
    def __init__(self, name):
        self.name = name

    def __repr__(self):
        return f"Animal({self.name!r})"


animal = Animal("Max")
print(animal)
\end{example}

\iosample{}{Animal('Max') \\}

這裡的 \texttt{Animal} 完全沒有定義 \texttt{\_\_str\_\_}，但 \texttt{print(animal)} 一樣正常
印出東西，印出來的還是 \texttt{\_\_repr\_\_} 的格式。這是因為 \texttt{print()} 的規則是「優先
找 \texttt{\_\_str\_\_}，找不到才退而求其次改用 \texttt{\_\_repr\_\_}」，不是「沒有
\texttt{\_\_str\_\_} 就出錯」。這也是為什麼很多 class 就算偷懶只定義 \texttt{\_\_repr\_\_}、
沒定義 \texttt{\_\_str\_\_}，\texttt{print()} 出來的結果通常還是看得懂，只是格式會比較像是
給開發者看的那種樣子。

\keypoint{一句話記住兩者的分工：\texttt{\_\_str\_\_} 是說給一般人聽的，\texttt{\_\_repr\_\_}
是說給寫程式的人（包括除錯時的你自己）聽的。\texttt{print()} 優先找 \texttt{\_\_str\_\_}；如果
一個物件根本沒定義 \texttt{\_\_str\_\_}，\texttt{print()} 才會退而求其次改用
\texttt{\_\_repr\_\_} 的結果，但只要你像本節這樣兩個都定義了，各自在各自的場合派上用場。}

\begin{drills}[label=drills:str-vs-repr]
\item 承接上面的 \texttt{dog}（已經同時定義了 \texttt{\_\_str\_\_} 跟
\texttt{\_\_repr\_\_}），\texttt{print(dog)} 這一行實際呼叫的是 \texttt{\_\_str\_\_} 還是
\texttt{\_\_repr\_\_}？
\item 如果你在互動式直譯器裡，先打 \texttt{dog}（不加 \texttt{print()}）按 Enter，接著又打
\texttt{print([dog, dog])}，這兩次分別呼叫的是 \texttt{\_\_str\_\_} 還是
\texttt{\_\_repr\_\_}？
\end{drills}

\section{回頭看名字牌模型：比較與複製物件}
\label{sec:copy-and-equality}

還記得第 \ref{ch:variables} 章的「名字牌模型」嗎？變數不是裝著值的盒子，是貼在物件上的名字牌，\texttt{=} 不是
複製，是幫物件多貼一張名字牌。

物件也是一樣的道理，只是物件比數字、字串複雜得多：物件裡面常常又
貼著別的名字牌（例如一個 list 屬性），這時候「複製」一個物件、或是拿兩個物件互相比較，到底
是在比什麼、複製到哪一層，就是這一節要搞清楚的事。這裡用一個你在第 \ref{ch:functions} 章就見過的老朋友，也就是
購物車，重新示範一次。

\begin{example}{python}{(cart-equality.py) 兩台內容一樣的購物車，== 卻是 False}
class ShoppingCart:
    def __init__(self, items):
        self.items = items


cart1 = ShoppingCart(["apple", "milk"])
cart2 = ShoppingCart(["apple", "milk"])

print(cart1 == cart2)
print(cart1 == cart1)
print(cart1 is cart2)
\end{example}

\iosample{}{False \\ True \\ False \\}

\texttt{cart1} 跟 \texttt{cart2} 裝的東西明明一模一樣，\texttt{cart1 == cart2} 卻是
\texttt{False}，這是因為自訂物件的 \texttt{==} 預設行為，跟 \texttt{is} 其實是同一套邏輯：
比的是「這兩個名字牌是不是貼在同一個物件上」，不是「這兩個物件裝的內容是不是一樣」。
\texttt{cart1 == cart1} 是 \texttt{True}，因為那本來就是同一個物件；\texttt{cart1 is cart2}
是 \texttt{False}，因為它們是兩個分開建立的物件，只是內容剛好一樣。如果你想要的是「比內容」，
要自己定義一個特殊方法叫 \texttt{\_\_eq\_\_}。

\begin{code}{python}{定義 \_\_eq\_\_}
class ClassName:
    def __eq__(self, other):
        return self.attr == other.attr
\end{code}

跟 \texttt{\_\_str\_\_} 一樣，\texttt{\_\_eq\_\_} 也是你不會自己去呼叫的特殊方法，只要
你寫了 \texttt{\_\_eq\_\_}，往後任何用到 \texttt{==} 比較這個 class 的地方，Python 都會
自動改叫它。
\index{\_\_eq\_\_}

\begin{example}{python}{(cart-eq-fixed.py) 定義 \_\_eq\_\_，改成比內容}
class ShoppingCart:
    def __init__(self, items):
        self.items = items

    def __eq__(self, other):
        return self.items == other.items


cart1 = ShoppingCart(["apple", "milk"])
cart2 = ShoppingCart(["apple", "milk"])

print(cart1 == cart2)
\end{example}

\iosample{}{True \\}

除了 \texttt{\_\_eq\_\_} 之外，Python 還有一個很類似的特殊方法叫 \texttt{\_\_lt\_\_}（less
than，「小於」），只要幫一個 class 定義了 \texttt{\_\_lt\_\_}，這個 class 的物件就能直接拿去
用 \texttt{sorted()} 排序，不用另外告訴 Python「依照什麼順序排」。
\index{\_\_lt\_\_}
定義方式跟 \texttt{\_\_eq\_\_}
很像，只是把比較運算子從 \texttt{==} 換成 \texttt{<}。

\begin{example}{python}{(cart-lt.py) 定義 \_\_lt\_\_，讓購物車可以排序}
class ShoppingCart:
    def __init__(self, items):
        self.items = items

    def __eq__(self, other):
        return self.items == other.items

    def __lt__(self, other):
        return len(self.items) < len(other.items)


cart1 = ShoppingCart(["apple", "milk", "bread"])
cart2 = ShoppingCart(["apple"])
cart3 = ShoppingCart(["milk", "egg"])

for cart in sorted([cart1, cart2, cart3]):
    print(cart.items)
\end{example}

\iosample{}{['apple'] \\{} ['milk', 'egg'] \\{} ['apple', 'milk', 'bread'] \\}

這裡的 \texttt{\_\_lt\_\_} 定義成「裝的東西比較少的購物車，比較『小』」，所以 \texttt{sorted()}
會依照每台購物車裝的東西數量，從少排到多。跟第 \ref{ch:lists} 章學過的 \texttt{sorted()} 是同一個內建函式，
差別在於這次排序的對象不是數字或字串，是你自己定義的物件，Python 需要靠 \texttt{\_\_lt\_\_}
才知道「兩個購物車之間，誰比較小」。如果你想原地排序（不造一份新的），也可以對一份裝著購物車
的 list 呼叫 \texttt{.sort()}，第 \ref{ch:lists} 章學過的 \texttt{sort()}／\texttt{sorted()} 差異，套用在
自訂物件上也是同一套規則。

\keypoint{這裡只點到為止地介紹了 \texttt{\_\_lt\_\_}，讓你知道自訂物件也能被排序這件事存在，
不展開 \texttt{\_\_le\_\_}／\texttt{\_\_gt\_\_}／\texttt{\_\_ge\_\_} 這一整組完整的比較特殊
方法（分別對應 \texttt{<=}／\texttt{>}／\texttt{>=}）。之後如果你真的需要一個 class 支援全部
的比較運算子，可以再回頭查文件，這裡先知道 \texttt{sorted()} 對自訂物件排序時，靠的正是
\texttt{\_\_lt\_\_} 就夠了。}

現在來看複製。你可能會想，複製一台購物車給朋友，應該就是「造一台一模一樣、但各自獨立」的
新購物車吧？實際上沒這麼簡單。

\begin{example}{python}{(cart-shallow-copy.py) 淺層複製陷阱}
import copy


class ShoppingCart:
    def __init__(self, items):
        self.items = items


cart1 = ShoppingCart(["apple", "milk"])

# 淺層複製：只複製最外層的 ShoppingCart 物件，items 還是同一個 list
cart2 = copy.copy(cart1)
cart2.items.append("egg")
print("cart1.items:", cart1.items)
print("cart2.items:", cart2.items)

# 深層複製：連 items 這個 list 也複製了一份，兩台車才真正各自獨立
cart3 = copy.deepcopy(cart1)
cart3.items.append("bread")
print("cart1.items:", cart1.items)
print("cart3.items:", cart3.items)
\end{example}

\iosample{}{cart1.items: ['apple', 'milk', 'egg'] \\ cart2.items: ['apple', 'milk', 'egg'] \\
cart1.items: ['apple', 'milk', 'egg'] \\ cart3.items: ['apple', 'milk', 'egg', 'bread'] \\}

往 \texttt{cart2} 裡加東西，\texttt{cart1} 竟然也跟著多了一樣，這是因為
\texttt{copy.copy}（淺層複製，shallow copy）只複製最外層的 \texttt{ShoppingCart} 物件本身，
物件裡面裝的 \texttt{items} 這個 list，兩台車其實還是貼著同一張名字牌，指向同一個 list
物件。
\index{淺層複製}
往 \texttt{cart3} 加東西就不會影響到 \texttt{cart1}，因為 \texttt{copy.deepcopy}（深層
複製，deep copy）連 \texttt{items} 這個 list 也另外造了一份，這次兩台車才是真正各自獨立、
互不影響。
\index{深層複製}

\keypoint{不確定該用 \texttt{copy.copy} 還是 \texttt{copy.deepcopy} 的時候，先問自己一個
問題：這個物件的屬性裡，有沒有 list、dict 這種可變（mutable）的東西？如果有，複製的時候就要
多想一步，淺層複製只複製最外層，內層那些可變屬性還是同一份，改一邊、另一邊也會跟著變；只有
\texttt{copy.deepcopy} 才會連內層也複製一份。}

這個陷阱不是只有 list 屬性才會遇到。只要屬性裡裝的是任何可變物件（list、dict、甚至另一個
自訂物件），淺層複製都可能踩到同一個坑。下面換一個裝著 dict 的例子看看。

\begin{example}{python}{(packing-list-copy.py) 淺層複製陷阱的 dict 版本}
import copy


class PackingList:
    def __init__(self, categories):
        self.categories = categories


trip1 = PackingList({"clothes": ["T-shirt", "jacket"], "toiletries": ["toothbrush"]})

trip2 = copy.copy(trip1)
trip2.categories["clothes"].append("umbrella")
print("trip1:", trip1.categories["clothes"])
print("trip2:", trip2.categories["clothes"])

trip3 = copy.deepcopy(trip1)
trip3.categories["clothes"].append("sunglasses")
print("trip1:", trip1.categories["clothes"])
print("trip3:", trip3.categories["clothes"])
\end{example}

\iosample{}{trip1: ['T-shirt', 'jacket', 'umbrella'] \\ trip2: ['T-shirt', 'jacket', 'umbrella'] \\
trip1: ['T-shirt', 'jacket', 'umbrella'] \\ trip3: ['T-shirt', 'jacket', 'umbrella', 'sunglasses'] \\}

\texttt{PackingList} 這次裝的是一個 dict（`categories`），dict 裡每個 key 對應的 value 又是
一個 list。跟購物車的例子一模一樣：\texttt{copy.copy} 只複製最外層的 \texttt{PackingList}
物件，\texttt{categories} 這個 dict 本身、以及裡面每一個 list，兩趟行程其實都還是同一份，
所以往 \texttt{trip2} 裡加東西，\texttt{trip1} 也會跟著變。

只有 \texttt{copy.deepcopy} 會
連 dict、以及 dict 裡面每一個 list 都遞迴複製一份，這次 \texttt{trip3} 才真正跟 \texttt{trip1}
各自獨立。淺層複製與深層複製這組概念，套用在 list 屬性、dict 屬性，甚至更複雜的巢狀結構上，
道理都是同一套：只複製最外層的物件，還是連內層所有可變物件也一起複製。

最後，回頭看 \texttt{ShoppingCart} 的 \texttt{\_\_init\_\_} 還有一個地方值得注意，如果你把
它寫成這樣：

\begin{example}{python}{(cart-mutable-default.py) 購物車版的可變預設參數陷阱}
class ShoppingCart:
    def __init__(self, items=[]):  # 危險：這個預設值只在定義的時候造一次
        self.items = items


cart_a = ShoppingCart()
cart_a.items.append("apple")

cart_b = ShoppingCart()
print("cart_b.items:", cart_b.items)


class FixedShoppingCart:
    def __init__(self, items=None):
        self.items = items if items is not None else []


cart_c = FixedShoppingCart()
cart_c.items.append("apple")

cart_d = FixedShoppingCart()
print("cart_d.items:", cart_d.items)
\end{example}

\iosample{}{cart\_b.items: ['apple'] \\ cart\_d.items: [] \\}

\texttt{cart\_a} 沒指定 \texttt{items}，往它的購物車加了一顆蘋果；照直覺，另外新造一台沒指定
內容的 \texttt{cart\_b} 應該是空的才對，但它的 \texttt{items} 裡卻也有一顆蘋果。

\keypoint{這就是第 \ref{ch:functions} 章教過的那個坑，只是這次出現在 class 的 \texttt{\_\_init\_\_} 裡：
\texttt{items=[]} 這個預設值，只在 \texttt{def \_\_init\_\_(...)} 這一行執行的當下造一次那個
list 物件，之後每次呼叫 \texttt{ShoppingCart()}、沒有自己傳 \texttt{items} 的話，用的都是
同一個物件，\texttt{cart\_a} 跟 \texttt{cart\_b} 其實共用同一張「預設 items」名字牌。安全
寫法跟第 \ref{ch:functions} 章一樣：預設值改成 \texttt{None}，在 \texttt{\_\_init\_\_} 裡面判斷「如果沒傳，才
新造一個空 list」，上面的 \texttt{FixedShoppingCart} 就是這樣寫的，\texttt{cart\_c} 跟
\texttt{cart\_d} 才會真正各自獨立。}

\begin{drills}[label=drills:copy-and-equality]
\item
\begin{verbatim}
class Box:
    def __init__(self, contents):
        self.contents = contents

box1 = Box(["book", "pen"])
box2 = copy.copy(box1)
box2.contents.append("ruler")
print(box1.contents)
\end{verbatim}
（假設檔案最上面已經 \texttt{import copy}）這段程式會印出什麼？
\item
\begin{verbatim}
class Point:
    def __init__(self, x, y):
        self.x = x
        self.y = y

p1 = Point(1, 2)
p2 = Point(1, 2)
print(p1 == p2)
\end{verbatim}
\texttt{Point} 沒有定義 \texttt{\_\_eq\_\_}，這段程式會印出什麼？
\item
\begin{verbatim}
class Logger:
    def __init__(self, messages=[]):
        self.messages = messages

    def add(self, msg):
        self.messages.append(msg)

log1 = Logger()
log1.add("started")

log2 = Logger()
print(log2.messages)
\end{verbatim}
這段程式會印出什麼？為什麼？
\item 承接 \texttt{PackingList} 的例子，如果改成
\texttt{trip4 = copy.copy(trip1)}，接著執行
\texttt{trip4.categories["toiletries"] = ["toothbrush", "towel"]}（直接把某個 key 對應的值整個
換掉，不是用 \texttt{.append()}），\texttt{trip1.categories["toiletries"]} 會不會也跟著變？
為什麼？
\end{drills}

這一章走了不少東西，回頭串一次：

\begin{itemize}
\item 繼承（可能疊好幾層）讓子類別「白拿」父類別的屬性和方法，需要
不一樣時再覆寫、再用 \texttt{super()} 銜接父類別原本的設定；組合則是另一條路，不談「是一種」，
談「有一個」，把另一個物件存成自己的屬性，換來的是不會自動白拿東西，一定要點名 \texttt{self.
xxx} 才拿得到。
\item \texttt{\_\_str\_\_}／\texttt{\_\_repr\_\_} 讓物件學會怎麼自我介紹，一個對使用者，
一個對開發者。
\item \texttt{\_\_eq\_\_}／\texttt{\_\_lt\_\_}／淺層與深層複製，講的其實都是同一
件事：物件比數字、字串複雜，第 \ref{ch:variables} 章的名字牌模型延伸到物件世界之後，得多想一層「這裡複製到底
複製到哪一層、比較到底在比什麼」。
\end{itemize}

這些工具你不需要每次寫程式都用到，但知道它們存在、知道遇到
奇怪行為時該往哪個方向查，之後真的碰上了就不會兩眼一抹黑。

\section*{練習解答}
\addcontentsline{toc}{section}{練習解答}

\paragraph{第~\ref{sec:inheritance-basics}~節}
\begin{enumerate}
\item Luna made a sound. \texttt{Cat} 沒有覆寫 \texttt{speak()}，直接繼承
\texttt{Animal} 的版本。
\item Rex barked! \texttt{Dog} 覆寫過 \texttt{speak()}，呼叫時一律以子類別這份為準。
\item 會出現 \texttt{AttributeError: 'Cat' object has no attribute 'name'}。因為
\texttt{Cat.\_\_init\_\_} 沒有呼叫 \texttt{super().\_\_init\_\_(name)}，
\texttt{Animal.\_\_init\_\_} 裡 \texttt{self.name = name} 那一行根本沒有被執行過，物件
身上沒有 \texttt{name} 這個屬性，\texttt{speak()} 裡用到 \texttt{self.name} 時就出錯了。
\end{enumerate}

\paragraph{第~\ref{sec:multi-level-inheritance}~節}
\begin{enumerate}
\item Buddy barked! \texttt{Puppy} 如果沒有覆寫 \texttt{speak()}，就會往上找到直接父類別
\texttt{Dog} 定義的版本，跟第~\ref{sec:inheritance-basics}~節看到的覆寫規則是同一套，只是這次
「往上找」多找了一層。
\item Tiny ran off to fetch the ball. \texttt{Toy} 沒有自己定義 \texttt{fetch()}，也沒有繼承它的
\texttt{Puppy} 定義過 \texttt{fetch()}，所以 Python 會繼續往上找到 \texttt{Dog} 定義的那一份，
正常執行。
\item 會出現 \texttt{TypeError: Dog.\_\_init\_\_() missing 1 required positional argument:
'breed'}。\texttt{Puppy.\_\_init\_\_} 呼叫 \texttt{super().\_\_init\_\_(name)} 時只傳了
\texttt{name}，但它跳到的 \texttt{Dog.\_\_init\_\_} 需要 \texttt{name} 跟 \texttt{breed} 兩個
引數，缺一個就會直接報錯。
\end{enumerate}

\paragraph{第~\ref{sec:inheritance-vs-composition}~節}
\begin{enumerate}
\item 不會出錯，\texttt{car.horsepower} 可以直接用，因為改成繼承之後，\texttt{Engine} 的
屬性會整個「白拿」到 \texttt{Car} 身上。但語意上會很奇怪，等於在說「車子是一種引擎」，不符合
現實世界的認知，這裡繼承在語法上可行，不代表這是恰當的設計選擇。
\item Engine started, horsepower 150. \texttt{car.engine} 是一個 \texttt{Engine} 物件，
\texttt{car.engine.start()} 正常呼叫 \texttt{Engine} 的 \texttt{start()} 方法。
\end{enumerate}

\paragraph{第~\ref{sec:str-vs-repr}~節}
\begin{enumerate}
\item \texttt{\_\_str\_\_}，\texttt{print()} 一個物件時優先呼叫的是 \texttt{\_\_str\_\_}。
\item 兩次都是 \texttt{\_\_repr\_\_}，在直譯器直接打變數名稱、以及物件被放進 list 裡顯示，
用的都是 \texttt{\_\_repr\_\_}，不是 \texttt{\_\_str\_\_}。
\end{enumerate}

\paragraph{第~\ref{sec:copy-and-equality}~節}
\begin{enumerate}
\item \texttt{['book', 'pen', 'ruler']}，\texttt{copy.copy} 是淺層複製，\texttt{box1} 跟
\texttt{box2} 的 \texttt{contents} 其實是同一個 list，往 \texttt{box2.contents} 加東西，
\texttt{box1.contents} 也會跟著變。
\item \texttt{False}，\texttt{Point} 沒有定義 \texttt{\_\_eq\_\_}，\texttt{==} 用的是
預設行為，比的是「是不是同一個物件」，\texttt{p1} 跟 \texttt{p2} 是兩個分開建立的物件，即使
內容一樣也是 \texttt{False}。
\item \texttt{['started']}，這是可變預設參數陷阱的 class 版本：\texttt{messages=[]} 這個
list 只在 \texttt{def \_\_init\_\_} 執行時造一次，\texttt{log1} 跟 \texttt{log2} 沒指定
\texttt{messages} 的話用的是同一個 list，\texttt{log1.add("started")} 加進去的內容，
\texttt{log2.messages} 也會看到。
\item 會，\texttt{trip1.categories["toiletries"]} 也會變成 \texttt{['toothbrush', 'towel']}。
\texttt{copy.copy} 只複製最外層的 \texttt{PackingList} 物件，\texttt{trip1.categories} 跟
\texttt{trip4.categories} 其實是同一個 dict 物件，不管是修改 dict 裡某個 key 對應的 list
內容，還是直接把某個 key 重新賦值成一個全新的 list，改的都是同一個 dict，兩邊看到的結果
都會一起變。
\end{enumerate}

\chapter{正規表達式}
\label{ch:regex}

到目前為止，你已經用 \texttt{in}、\texttt{.replace()}、\texttt{.split()} 這些字串方法處理過
不少「文字裡有沒有出現某個東西、要不要換掉某段文字」的問題。這些方法有一個共同的前提：你得
\textbf{確切知道}要找的字串長什麼樣子。

但現實裡很多時候你要找的其實是一種\textbf{格式}：
你可能不知道一支電話號碼的確切數字是什麼，但看到「三個數字、一個連字號、三個數字、一個連字
號、四個數字」這種形狀，你一眼就認得出那是電話號碼；一個 email 地址也是一樣，不管帳號跟網域
是什麼，「一段文字、一個 \texttt{@}、再一段文字、一個點」這個形狀是固定的。正規表達式
（regular expression，常簡稱 regex）就是讓你把這種形狀寫給 Python 看的工具。
\index{正規表達式}

這一章分幾段推進：

\begin{itemize}
\item 先從你已經會的字串方法出發，看看「找格式」這件事跟「找固定字串」有什麼不
一樣，順便學會 \texttt{re} 模組最基本的用法。
\item 接著把「形狀」拆解成幾個最常用的積木（數字、
文字、空白、重複次數），組出電話號碼、email 這類實際會用到的 pattern。
\item 再來學怎麼一次把一段
文字裡所有符合的地方都找出來、怎麼確認一整段文字完全符合某種格式（也就是「驗證」）、怎麼用
regex 做搜尋取代跟切字串。
\end{itemize}

要先說在前面：正規表達式是一個很深的主題，光是語法規則就能寫成一
整本書，但這一章的目標不是讓你精通它，只要看到常見的 pattern 看得懂在比對什麼、需要的時候
自己也能寫出簡單的 pattern（數字、email 這種格式），這一章的任務就完成了，更進階的寫法留給
你以後真的用到時再自己查文件。

\section{為什麼需要正規表達式、\texttt{re} 模組怎麼用}
\label{sec:regex-basics}

假設你在整理一批客服留言，想知道哪些留言裡有留電話號碼、哪些沒有。用你已經會的工具，你會
怎麼寫？如果你知道確切的電話號碼是什麼，\texttt{in} 可以幫你判斷這段文字有沒有出現；但你
不可能事先知道每個客戶會留下哪個號碼，你真正想問的其實是「這裡有沒有出現\textbf{一串數字}」，
而不是「這裡有沒有出現某個特定的字串」。

這正是字串方法做不到、但正規表達式擅長的事：它讓你
可以描述「找長這種樣子的東西」，不只是「找這個字」。

要在 Python 裡用正規表達式，第一步是 \texttt{import re}。\texttt{re} 是 Python 內建的正規
表達式模組，不用另外安裝。最常用的入口是 \texttt{re.search(pattern, text)}：如果 \texttt{text}
裡有一段符合 \texttt{pattern} 這個「形狀」的文字，它會回傳一個 Match 物件；找不到的話，回傳
\texttt{None}。拿到 Match 物件之後，呼叫它的 \texttt{.group()} 方法，就能取出實際比對到的那段
文字。
\index{re 模組}
\index{re.search()}
\index{Match 物件}

\begin{code}{python}{Matching a pattern and extracting the matched text}
import re

match = re.search(pattern, text)
if match:
    print(match.group())
else:
    print("Not found")
\end{code}

這裡 \texttt{pattern} 通常會寫成 \texttt{r'...'} 這種前面多一個 \texttt{r} 的字串，叫 raw
string。正規表達式常常要用到反斜線（像等一下會教的 \texttt{\textbackslash d}），如果不加
這個 \texttt{r}，反斜線可能被 Python 字串本身的跳脫規則誤解成別的意思。之後每一段 pattern，
這本書都會統一用 \texttt{r'...'} 這種寫法，你也照著做就好，不用另外深究背後的規則。
\index{raw string}

比對「一個數字」的 pattern 寫作 \texttt{\textbackslash d}，先只用這一個最基本的元件，其他
常用的元件留到下一節再教。

\begin{example}{python}{(check-digit.py) Checking whether a message contains a digit}
import re

message = "Hello, my phone number is 0912345678, please call back soon."

match = re.search(r'\d', message)
if match:
    print("This message contains a digit, the first one is:", match.group())
else:
    print("This message has no digits")
\end{example}

\iosample{}{This message contains a digit, the first one is: 0 \\}

如果同一個 pattern 之後還要重複用好幾次，你也可以先用 \texttt{re.compile()} 把它變成一個
regex 物件、貼給一個名字，例如 \texttt{pattern = re.compile(r'\textbackslash d')}，之後
每次呼叫 \texttt{pattern.search(text)} 就好。這裡的 \texttt{pattern} 這個名字牌，跟你前面
幾章看過的物件用法完全一樣，只是這次貼的對象是一個 regex 物件而已。

\index{re.compile()}
還記得第 \ref{ch:files-exceptions} 章教的「先做、錯了再處理」（EAFP）嗎？這裡看起來不太一樣：\texttt{search()} 比對
不到的時候，是老實回傳 \texttt{None}，不是丟例外，所以慣用寫法是先用 \texttt{if match:} 檢查
再用，乍看之下比較像「先檢查」（LBYL）。

這其實仍然算是 EAFP 的精神，只是同一個原則的另一面：
EAFP 真正在乎的是操作本身可能失敗時，要用能反映這件事的寫法。\texttt{search()} 的設計本來
就是老實回答「有」或「沒有」，不是「可能出錯會丟例外的操作」（不像開檔案、字典查詢那樣），所以
檢查 \texttt{if match:} 才是正確的寫法，不需要每個地方都硬套 \texttt{try/except}。

如果忘記檢查、直接對比對結果呼叫 \texttt{.group()}，會發生什麼事？

\begin{example}{python}{(check-digit-broken.py) Calling .group() without checking first}
import re

message = "Thank you for your message, a representative will get back to you soon."
match = re.search(r'\d', message)
print(match.group())
\end{example}

這則留言裡沒有任何數字，\texttt{match} 會是 \texttt{None}，執行到最後一行時會看到：

\begin{verbatim}
Traceback (most recent call last):
  File "check-digit-broken.py", line 4, in <module>
    print(match.group())
          ^^^^^^^^^^^
AttributeError: 'NoneType' object has no attribute 'group'
\end{verbatim}

由下往上讀：最後一行 \texttt{AttributeError: 'NoneType' object has no attribute 'group'}
是說「\texttt{None} 這種東西身上沒有 \texttt{group} 這個方法可以呼叫」，因為 \texttt{match}
根本不是 Match 物件，是 \texttt{search()} 老實告訴你「沒找到」的 \texttt{None}。修法很直接：
呼叫 \texttt{.group()} 之前，一定要先用 \texttt{if match:} 確認真的有比對到，這也是為什麼
本節一開始的片語會把 \texttt{if match:} 當成固定樣板的一部分，不是可有可無的裝飾。

\begin{drills}[label=drills:regex-basics]
\item 用 \texttt{re.search(r'\textbackslash d', "Customer ID A007")} 這行程式，你覺得會不會
比對到？如果會，比對到的是哪個字元？
\item 用 \texttt{re.search(r'\textbackslash d', "A sentence without any digits")}，你覺得 \texttt{match}
最後會是 \texttt{None} 還是一個 Match 物件？如果接著不檢查、直接寫
\texttt{match.group()}，會發生什麼事？
\end{drills}

\section{常用 pattern 積木：\texttt{\textbackslash d}、\texttt{\textbackslash w}、
\texttt{\textbackslash s} 與量詞}
\label{sec:regex-patterns}

上一節只用了「比對一個數字」這個最小的積木，但實際情境裡，你通常要描述的是一整段有規律的
格式，例如「三個數字、一個連字號、三個數字、一個連字號、四個數字」。要做到這件事，你需要
兩種積木：一種是「這個位置該是什麼樣的字元」，另一種是「這個位置要重複幾次」。

最常用的三個字元類別：

\begin{itemize}
\item \texttt{\textbackslash d}：一個數字（0 到 9 之間任何一個）。
\item \texttt{\textbackslash w}：一個「文字字元」，包含英文字母、數字、底線。
\item \texttt{\textbackslash s}：一個空白字元，包含空格、換行、tab。
\end{itemize}
\index{字元類別 d, w, s}

拿標點符號試試就能感受到 \texttt{\textbackslash w} 跟你想像的「文字」不完全一樣：
\texttt{re.search(r'\textbackslash w', "！")} 會是 \texttt{None}，因為驚嘆號不算英文
字母、數字或底線；而 \texttt{re.search(r'\textbackslash s', "a b").group()} 會找到
\texttt{"a"} 跟 \texttt{"b"} 中間那個空格。

光有字元類別還不夠，你還需要說「要重複幾次」，這就是量詞的工作：

\begin{itemize}
\item \texttt{*}：前面那個元件出現零次以上（可以完全不出現）。
\item \texttt{+}：前面那個元件出現一次以上（至少要有一個）。
\item \texttt{?}：前面那個元件可有可無（出現零次或一次）。
\item \texttt{\{n\}}：前面那個元件精確出現 \texttt{n} 次。
\item \texttt{\{n,m\}}：前面那個元件出現 \texttt{n} 到 \texttt{m} 次（含頭尾兩端）。
\end{itemize}
\index{量詞}

把字元類別跟量詞組合起來，就能寫出「三個數字、一個連字號、三個數字、一個連字號、四個數字」
這種格式：\texttt{\textbackslash d\{3\}-\textbackslash d\{3\}-\textbackslash d\{4\}}。

\begin{code}{python}{Combining character classes with quantifiers}
import re

pattern = r'\d{3}-\d{3}-\d{4}'
match = re.search(pattern, text)
if match:
    print(match.group())
\end{code}

\begin{example}{python}{(check-phone-format.py) Checking whether a message has a phone number in the required format}
import re

message = "Hello, my phone number is 091-234-5678, please call back soon."

match = re.search(r'\d{3}-\d{3}-\d{4}', message)
if match:
    print("Found a phone number matching the format:", match.group())
else:
    print("No phone number matching the format was found")
\end{example}

\iosample{}{Found a phone number matching the format: 091-234-5678 \\}

如果客戶留的號碼格式不太一樣，例如少了一個連字號、或分段方式不同，就比對不到。把
\texttt{message} 換成 \texttt{"Hello, my phone number is 0912-345-678, please call back soon."} 再跑一次，
會印出「No phone number matching the format was found」，因為 \texttt{4-3-3} 這種分段方式跟
\texttt{\textbackslash d\{3\}-\textbackslash d\{3\}-\textbackslash d\{4\}} 描述的
\texttt{3-3-4} 形狀對不上。

\keypoint{字元類別跟量詞要記住的重點是「規律的東西才需要 regex」，不是要你把每個符號背下
來。如果你要找的字串本身就是固定的（例如一定是 \texttt{"客服專線"} 這四個字），直接用
\texttt{in} 或 \texttt{.replace()} 更簡單、也更好懂，不用每個地方都想著要動用正規表達式。}

\begin{drills}[label=drills:regex-patterns]
\item 下面三個字串，哪些符合 \texttt{\textbackslash d\{3\}-\textbackslash d\{3\}-%
\textbackslash d\{4\}} 這個 pattern？為什麼：
(a) \texttt{"091-234-5678"}
(b) \texttt{"0912-345-678"}
(c) \texttt{"12-345-6789"}
\item 如果要用同樣的精神，寫一個「兩碼-兩碼-兩碼-兩碼」格式的 pattern（例如
\texttt{"12-34-56-78"}），你會怎麼寫？
\item 句子 \texttt{"Customer service line: 091-234-5678, service hours 9:00-18:00"} 用
\texttt{re.search(r'\textbackslash d\{3\}-\textbackslash d\{3\}-\textbackslash d\{4\}',
sentence)} 去找，比對到的會是哪一段文字？
\item 用 \texttt{\{n,m\}} 這個新學到的量詞，寫一個「密碼長度要在 6 到 12 個文字字元之間」
的 pattern（提示：文字字元就是上面教的 \texttt{\textbackslash w}）。
\end{drills}

\section{\texttt{findall()}：一次找出所有符合的地方}
\label{sec:regex-findall}

前面兩節的 \texttt{re.search()} 只會告訴你「有沒有找到、第一個找到的是哪段」，但很多實際
任務要的是這段文字裡\textbf{所有}符合的地方，例如從一份聯絡資訊清單裡，一次把所有 email
地址都抓出來。這種時候要用 \texttt{re.findall(pattern, text)}：它會把 \texttt{text} 裡所有
符合 \texttt{pattern} 的地方，通通收集成一個 list 回傳給你。
\index{re.findall()}

在動手之前，先組出「email 地址」這個 pattern。一個 email 大致長這樣：帳號部分（英文字母、
數字，可能夾雜幾個符號）、一個 \texttt{@}、網域部分（英文字母、數字、點）、最後接一段純
英文字母。把這幾段用上一節的積木拼起來，會得到
\texttt{[a-zA-Z0-9.\_\%+-]+@[a-zA-Z0-9.-]+\textbackslash.[a-zA-Z]+}。

中括號
\index{字元集合}
\texttt{[...]} 是「這個位置可以是括號裡列出的任何一個字元」，\texttt{+} 表示「這樣的字元
至少出現一次」，\texttt{\textbackslash.} 則是用反斜線跳脫，表示「真正的一個點，不是別的
意思」。這個 pattern 已經比前面幾節複雜不少，你不需要自己從零想出這種寫法，記得它長什麼
樣子、看得懂在比對什麼，需要的時候查文件或直接沿用，就已經達到這一章的目標。

「點要記得跳脫」是初學正規表達式很容易忽略的一個細節，因為 \texttt{.} 這個符號在正規表達式
裡有自己的意思：它代表「任何一個字元都可以」，不是專指「一個真正的點」。如果忘記在點前面加
反斜線，pattern 會比你以為的鬆很多：

\begin{example}{python}{(dot-escape.py) Forgetting to escape the dot makes a pattern too loose}
import re

samples = ["3.14", "3x14", "3-14"]

pattern_loose = r'\d.\d'    # 這裡的 . 沒有跳脫，所以可以比對任何字元
pattern_strict = r'\d\.\d'  # 這裡的 \. 才會真正比對到一個點

for s in samples:
    loose = re.search(pattern_loose, s)
    strict = re.search(pattern_strict, s)
    print(s, "loose:", loose.group() if loose else None,
          "strict:", strict.group() if strict else None)
\end{example}

\iosample{}{3.14 loose: 3.1 strict: 3.1 \\
3x14 loose: 3x1 strict: None \\
3-14 loose: 3-1 strict: None \\}

沒有跳脫的 \texttt{pattern\_loose} 對 \texttt{"3x14"}、\texttt{"3-14"} 這兩個字串也照樣
比對成功，因為 \texttt{.} 把中間那個 \texttt{x} 或 \texttt{-} 都當成「隨便一個字元」接受
了，並非只認得真正的點。加上反斜線的 \texttt{pattern\_strict} 才正確地只認得真正的一個點。
本節開頭 email pattern 裡的 \texttt{\textbackslash.}，用的就是這個跳脫寫法。

\begin{code}{python}{Finding all text that matches a pattern at once}
import re

results = re.findall(pattern, text)
print(results)
\end{code}

\begin{example}{python}{(find-emails.py) Finding all email addresses in a block of contact info}
import re

contacts = """Alice - alice.wang@example.com
Bob - bob99@mail-service.net
Cindy hasn't left an email yet, just wrote a note"""

pattern = r'[a-zA-Z0-9._%+-]+@[a-zA-Z0-9.-]+\.[a-zA-Z]+'
emails = re.findall(pattern, contacts)
print(emails)

# 把 pattern 用括號包起來，findall() 就只會回傳括號裡比對到的那一段
username_pattern = r'([a-zA-Z0-9._%+-]+)@[a-zA-Z0-9.-]+\.[a-zA-Z]+'
usernames = re.findall(username_pattern, contacts)
print(usernames)
\end{example}

\iosample{}{['alice.wang@example.com', 'bob99@mail-service.net'] \\{}
['alice.wang', 'bob99'] \\}

留意最後兩行的差別：第一次的 pattern 沒有括號，\texttt{findall()} 回傳的是完整比對到的
email；第二次把帳號的部分用括號 \texttt{()} 包起來，\texttt{findall()} 就只回傳括號裡比對
到的那一小段。這其實跟第一節 \texttt{match.group()} 是同一件事的延伸：括號告訴 Python
「這一段我特別在意，之後只要把這段還我就好」。這個用法叫「分組」，正式的規則（多個括號怎麼
編號、怎麼分別取出）還有更多細節，這裡先知道加了括號、拿到的東西會變就夠用。
\index{分組}

如果一個 pattern 裡有好幾組括號，光靠「第一個括號、第二個括號」這種順序去記，讀起來會很
吃力。這種時候可以把括號寫成\textbf{命名分組}，語法是 \texttt{(?P<名字>...)}，之後就能用
這個名字取出對應的那一段，而不用去數第幾個括號：
\index{命名分組}

\begin{example}{python}{(named-group-phone.py) Using named groups to pull apart a phone number}
import re

text = "Taylor's phone number is 091-234-5678"
pattern = r'(?P<area>\d{3})-(?P<mid>\d{3})-(?P<last>\d{4})'

match = re.search(pattern, text)
if match:
    print(match.group("area"))
    print(match.group("mid"))
    print(match.group("last"))
\end{example}

\iosample{}{091 \\
234 \\
5678 \\}

命名分組的好處是可讀性：一段複雜的 pattern 裡如果有好幾個括號，用名字取比用數字數第幾個
括號清楚很多。這裡先知道有這個寫法即可，之後真的碰到分組較多、需要清楚標示每一段代表什麼的
情境，再回頭查文件細節。

\section{驗證格式：用 \texttt{fullmatch()} 檢查整段字串}
\label{sec:regex-fullmatch}

前面幾節的 \texttt{search()}、\texttt{findall()} 有一個共同的特性：它們只在乎文字裡「有
沒有一段」符合 pattern，不在乎符合的那一段前後還有沒有其他文字。這對「找出留言裡的電話
號碼」這種任務很合適，但如果你的任務換成「驗證使用者填的這整個欄位是不是一支合法的電話
號碼」，\texttt{search()} 反而會讓你上當：

\begin{example}{python}{(search-trap.py) search() is not suitable for validation}
import re

pattern = r'\d{3}-\d{3}-\d{4}'
text = "091-234-5678abc"

match = re.search(pattern, text)
print("search() found:", match.group())
\end{example}

\iosample{}{search() found: 091-234-5678 \\}

這個輸入其實不是一支合法的電話號碼（後面多了 \texttt{"abc"}），但 \texttt{search()} 只顧
找出「有沒有一段符合」，找到 \texttt{"091-234-5678"} 這一段就滿意收工，完全不管後面還跟著
什麼。要正確做「驗證」，你需要的是確認\textbf{整段字串從頭到尾}都符合 pattern，不能有多餘
的部分，也不能少。這正是 \texttt{re.fullmatch(pattern, text)} 的工作：只有整個 \texttt{text}
完全符合 \texttt{pattern}，才會回傳 Match 物件，只要頭尾任何一邊多了或少了字元，都回傳
\texttt{None}。
\index{re.fullmatch()}

\begin{code}{python}{Confirming a whole string fully matches a pattern}
import re

result = re.fullmatch(pattern, text)
if result:
    print("Format is valid")
else:
    print("Format is invalid")
\end{code}

\begin{example}{python}{(validate-phone.py) Validating whether a phone field is in a legal format}
import re

def is_valid_phone(text):
    pattern = r'\d{3}-\d{3}-\d{4}'
    return re.fullmatch(pattern, text) is not None

print(is_valid_phone("091-234-5678"))
print(is_valid_phone("091-234-5678 "))
print(is_valid_phone("my phone number is 091-234-5678"))
\end{example}

\iosample{}{True \\
False \\
False \\}

三個輸入裡，只有第一個從頭到尾完全符合「三碼-三碼-四碼」的格式。第二個結尾多了一個空格、
第三個開頭多了一段文字，用 \texttt{search()} 都還是找得到，但用 \texttt{fullmatch()} 都會
正確判定「不合法」，因為它們符合的都只是字串中間夾著的一段，整段字串本身並不完全符合格式。

\keypoint{\texttt{search()} 跟 \texttt{fullmatch()} 很容易搞混，因為兩者的呼叫方式幾乎
一樣。判斷用哪一個的標準很簡單：如果任務是「這段文字裡有沒有出現某種格式」，用
\texttt{search()}；如果任務是「使用者填的這整個欄位，是不是一個合法的格式」，用
\texttt{fullmatch()}。用錯的後果是實際的資安風險，例如驗證信箱格式時如果誤用
\texttt{search()}，使用者填 \texttt{"attacker@evil.com; DROP TABLE users"} 這種夾帶額外
內容的字串也會被判定成合法，因為裡面確實藏著一段符合 email pattern 的文字。}

同樣的做法拿來驗證 email 也一樣：

\begin{example}{python}{(validate-email.py) Validating whether an email field is in a legal format}
import re

email_pattern = r'[a-zA-Z0-9._%+-]+@[a-zA-Z0-9.-]+\.[a-zA-Z]+'

def is_valid_email(text):
    return re.fullmatch(email_pattern, text) is not None

for candidate in ["alice.wang@example.com", "alice.wang@example.com please inquire",
                  "alice@", "a@b.tw"]:
    print(candidate, "->", is_valid_email(candidate))
\end{example}

\iosample{}{alice.wang@example.com -> True \\
alice.wang@example.com please inquire -> False \\
alice@ -> False \\
a@b.tw -> True \\}

留意 \texttt{"alice@"} 這個輸入：帳號跟 \texttt{@} 都在，但少了網域部分，不符合 pattern
要求的「\texttt{@} 後面要有網域、一個點、再一段英文字母」，所以正確判定成不合法。這正是
regex 驗證的價值：不用一項一項手動檢查「有沒有 \texttt{@}」「有沒有點」，一個 pattern 就
把整套格式規則講清楚。

\begin{drills}[label=drills:regex-fullmatch]
\item 用本節教的 \texttt{is\_valid\_phone()} 函式，\texttt{is\_valid\_phone("0912345678")}
（沒有連字號）的結果會是 \texttt{True} 還是 \texttt{False}？為什麼？
\item 如果只是想確認一段文字「有沒有出現」合法格式的電話號碼，而不是要求使用者填的整個
欄位都得是電話號碼，這種情境該用 \texttt{search()} 還是 \texttt{fullmatch()}？
\item 承上一節的命名分組 pattern，你會怎麼把它跟 \texttt{fullmatch()} 結合，寫一個
\texttt{is\_valid\_phone\_named(text)} 函式？（提示：命名分組不影響 \texttt{fullmatch()}
判斷整段字串符不符合的邏輯，只是額外多了取出各段的能力。）
\end{drills}

\section{\texttt{sub()}：簡單的搜尋取代}
\label{sec:regex-sub}

找到符合格式的文字之後，另一個常見需求是把它換掉，例如把一份客服留言紀錄裡的電話號碼、
email 地址都遮蔽掉再存檔，避免個資外流。跟字串的 \texttt{.replace()} 一樣，正規表達式也有
對應的取代工具，只是它換的是符合某個 pattern 的所有地方，不是固定字串：
\texttt{re.sub(pattern, replacement, text)} 會把 \texttt{text} 裡所有符合 \texttt{pattern}
的部分，統統換成 \texttt{replacement}，回傳換完的新字串。
\index{re.sub()}

\begin{code}{python}{Replacing all text that matches a pattern with something else}
import re

new_text = re.sub(pattern, replacement, text)
print(new_text)
\end{code}

\begin{example}{python}{(mask-contact-info.py) Masking emails and phone numbers in messages}
import re

log = """Customer message log:
Alice: please reply to alice.wang@example.com for me
Bob: my phone number is 091-234-5678, please call back soon
"""

email_pattern = r'[a-zA-Z0-9._%+-]+@[a-zA-Z0-9.-]+\.[a-zA-Z]+'
phone_pattern = r'\d{3}-\d{3}-\d{4}'

masked = re.sub(email_pattern, '****', log)
masked = re.sub(phone_pattern, '****', masked)
print(masked)
\end{example}

\iosample{}{Customer message log: \\
Alice: please reply to **** for me \\
Bob: my phone number is ****, please call back soon \\}

這支程式呼叫了兩次 \texttt{re.sub()}：第一次先把符合 email 格式的部分換掉，第二次接著在
換過的文字上再處理電話號碼格式，兩次取代互不干擾。像 \texttt{sub()} 用 \texttt{\textbackslash
1}、\texttt{\textbackslash 2} 這種寫法把比對到的分組內容原封不動塞回換掉的文字裡，是更進階
的用法，這裡不展開。

\section{\texttt{split()}：用正規表達式切字串}
\label{sec:regex-split}

第 \ref{ch:strings} 章教過的 \texttt{.split()} 方法，只能按照一種固定的分隔符號切字串，例如
\texttt{"a,b,c".split(",")}。但有些真實資料的分隔方式並不統一，同一份清單裡逗號、分號、
空格可能混著用，甚至同一種符號還連續出現好幾次（像打字時不小心多打了幾個空格）。

這種時候
用固定分隔符號的 \texttt{.split()} 就不夠用了，你需要的是「照某種格式切」，而不是「照某個
固定字元切」，這正是 \texttt{re.split(pattern, text)} 的工作：它會把 \texttt{text} 裡所有
符合 \texttt{pattern} 的地方當成分隔點，切成好幾段，回傳一個 list。
\index{re.split()}

\begin{code}{python}{Splitting a string using a pattern-described separator}
import re

parts = re.split(pattern, text)
print(parts)
\end{code}

\begin{example}{python}{(split-messy-list.py) Splitting a list with inconsistent separators}
import re

text = "apple, banana;  cherry   durian"

# 字元集合 [,;\s] 代表「一個逗號、分號、或任何空白字元」，
# 後面的 + 代表「連續出現一次以上的這些分隔符號，都算成同一個分隔點」
words = re.split(r'[,;\s]+', text)
print(words)
\end{example}

\iosample{}{['apple', 'banana', 'cherry', 'durian'] \\}

這個 pattern 裡的 \texttt{[,;\textbackslash s]} 是字元集合 \texttt{[...]} 的另一種用法：
不只可以放英文字母、數字這種範圍，也可以直接列出幾個具體的字元（逗號、分號）跟一個字元類別
（\texttt{\textbackslash s}）混在一起，意思是「這個位置可以是逗號、分號、或一個空白字元」。
後面的 \texttt{+} 讓連續好幾個分隔符號（例如中間那三個空格）合併當成一個分隔點，不會因為
分隔符號重複出現就切出一堆空字串。

\begin{example}{python}{(split-record.py) Using a regular expression to split a custom-format record}
import re

record = "Taylor:091-234-5678|Jordan:092-111-2222|Morgan:093-333-4444"
entries = re.split(r'[:|]', record)
print(entries)
\end{example}

\iosample{}{['Taylor', '091-234-5678', 'Jordan', '092-111-2222', 'Morgan', '093-333-4444'] \\}

這裡的分隔符號有兩種（冒號跟直線），如果用 \texttt{.split()} 只能處理其中一種，剩下的還得
再對每一段跑第二次 \texttt{.split()}。用 \texttt{re.split(r'[:|]', record)} 一次描述「冒號
或直線都算分隔點」，一步就切完，不用分兩階段處理。

\begin{drills}[label=drills:regex-split]
\item 給定 \texttt{text = "red/green,blue   yellow"}（分隔符號有斜線、逗號、空格三
種），你會怎麼寫一個 \texttt{re.split()} 的 pattern，把它切成
\texttt{['red', 'green', 'blue', 'yellow']}？
\item 用 \texttt{re.split(r'\textbackslash s+', "restore   multiple    spaces")} 這行程式，你覺得
回傳的 list 長什麼樣子？
\item \texttt{.split()} 跟 \texttt{re.split()} 該用哪一個，判斷的標準是什麼？
\end{drills}

\section{除錯正規表達式的小技巧}
\label{sec:regex-debug}

正規表達式有一個共同的痛點：一個 pattern 寫錯，通常不會出現 \texttt{SyntaxError} 這種明確
的錯誤訊息，程式照樣能執行，只是比對結果不是你要的（比對不到、或比對到不對的那一段）。這種
「沒有報錯、但結果不對」的狀況，比看得到錯誤訊息的問題更難抓，這裡整理幾個實用的除錯習慣。

第一個習慣是\textbf{從小處著手，不要一次寫出完整 pattern}。與其一口氣寫出一整段又長又複雜
的 pattern、丟給一段很長的真實資料去測，不如先寫一個最簡單的版本、丟一個你自己確定答案的短
字串去測，確認這一小段沒問題，再一步步加上其他部分。下面用一個日期格式的例子示範這個過程：

\begin{example}{python}{(debug-step-by-step.py) Debugging a date pattern step by step}
import re

text = "The registration deadline is 2026-09-01, act soon"

# 第一次嘗試：\d 只代表一個數字，忘了加上量詞
pattern_v1 = r'\d-\d-\d'
match = re.search(pattern_v1, text)
print("v1:", match)

# 加上量詞，改成「四位數字、兩位數字、兩位數字」
pattern_v2 = r'\d{4}-\d{2}-\d{2}'
match = re.search(pattern_v2, text)
print("v2:", match.group())
\end{example}

\iosample{}{v1: None \\
v2: 2026-09-01 \\}

\texttt{pattern\_v1} 比對不到任何東西，因為 \texttt{\textbackslash d} 預設只代表「一個」
數字，寫 \texttt{2026} 這種四位數完全不符合「一個數字、一個連字號、一個數字」的形狀。這種
時候與其對著長長的 pattern 一直瞪，不如像上面這樣先印出中間結果，確認到底是哪一段假設錯了
（這裡是漏了量詞），再逐步修正。

第二個習慣是\textbf{善用 \texttt{match.group()} 檢查「比對到的到底是哪一段」}。有時候
pattern 確實比對到東西了，但比對到的不是你以為的那一段，例如量詞寫太寬鬆，把不該算進去的
字元也吃了進去。與其用眼睛猜，直接印出 \texttt{match.group()} 看實際結果，比憑感覺判斷可靠
得多。

第三個習慣是\textbf{借助線上工具}。網路上有像 regex101.com 這類網站，讓你把 pattern 跟測試
字串貼上去，畫面會即時標示出比對到的每一段，也會逐條解釋 pattern 裡每個符號在做什麼。這種
視覺化的回饋，通常比在終端機裡印文字更快抓出問題所在，尤其是 pattern 開始變長、變複雜的
時候。

\keypoint{正規表達式沒有比對到東西，問題通常出在 pattern 描述的形狀跟你以為的不一樣，很少
是 Python 或 \texttt{re} 模組本身壞了。除錯的第一步永遠是回頭檢查 pattern 本身：量詞有沒有
漏加、字元類別選對了沒有、raw string 的 \texttt{r} 前綴有沒有忘記加，把力氣花在懷疑
\texttt{re.search()} 這個函式本身有問題，通常是走錯方向。}

\begin{drills}[label=drills:regex-debug]
\item 一個 pattern \texttt{r'\textbackslash d+'} 原本是想比對「一段連續的數字」，拿去比對
\texttt{"Order number AB-123"} 卻只比對到 \texttt{"123"}，沒有把 \texttt{"AB"} 算進去，這是
pattern 寫錯了嗎？還是 pattern 本來就只該比對到 \texttt{"123"}？
\item 如果一個 pattern 在你自己準備的短測試字串上跑起來正常，換到真實資料上卻常常找不到，
你會優先懷疑哪裡出了問題？
\end{drills}

正規表達式其實還有很多這一章沒教到的語法，像用 \texttt{|} 表示「這個或那個」、用
\texttt{\^{}} 和 \texttt{\$} 限定「一定要出現在開頭／結尾」（本章改用 \texttt{fullmatch()}
達到類似的驗證效果）、\texttt{re.IGNORECASE} 這種忽略英文字母大小寫的選項、還有更精細的
非貪婪比對寫法。這些留給你以後真的用到時，再查 Python 官方文件裡 \texttt{re} 模組的說明
即可，不需要現在就全部記住。

\section*{練習解答}
\addcontentsline{toc}{section}{練習解答}

\paragraph{第~\ref{sec:regex-basics}~節}
\begin{enumerate}
\item 會比對到，字串裡的 \texttt{"007"} 含有數字，\texttt{\textbackslash d} 會比對到第一個
出現的數字字元，也就是 \texttt{"0"}。
\item \texttt{match} 會是 \texttt{None}，因為這句話裡沒有任何數字字元。如果不檢查就直接呼叫
\texttt{match.group()}，會出現
\texttt{AttributeError: 'NoneType' object has no attribute 'group'}，因為 \texttt{None}
身上沒有 \texttt{group} 這個方法可以呼叫。
\end{enumerate}

\paragraph{第~\ref{sec:regex-patterns}~節}
\begin{enumerate}
\item 只有 (a) \texttt{"091-234-5678"} 符合，它剛好是「三碼-三碼-四碼」的形狀。(b) 是
「四碼-三碼-三碼」，(c) 是「兩碼-三碼-四碼」，分段的數字個數都對不上 pattern 要求的
\texttt{3-3-4}。
\item \texttt{r'\textbackslash d\{2\}-\textbackslash d\{2\}-\textbackslash d\{2\}-%
\textbackslash d\{2\}'}，把重複次數從 \texttt{\{3\}}／\texttt{\{4\}} 換成
\texttt{\{2\}}，段數也跟著改成四段。
\item 比對到的是 \texttt{"091-234-5678"}，\texttt{search()} 只會找出文字裡第一段符合
pattern 的部分，後面的 \texttt{"9:00-18:00"} 因為格式不符（不是三碼-三碼-四碼），不會被
算進去。
\item \texttt{r'\textbackslash w\{6,12\}'}，用 \texttt{\{6,12\}} 表示「6 到 12 次」，
\texttt{\textbackslash w} 涵蓋英文字母、數字、底線這些常見的密碼字元。
\end{enumerate}

\paragraph{第~\ref{sec:regex-findall}~節}
\begin{enumerate}
\item \texttt{['amy@school.edu.tw', 'david.chen@company.com']}，兩個 email 地址都符合
pattern，依出現順序收集成一個 list。
\end{enumerate}

\paragraph{第~\ref{sec:regex-fullmatch}~節}
\begin{enumerate}
\item \texttt{False}，因為 \texttt{is\_valid\_phone()} 用的 pattern 要求整段字串一定要有
兩個連字號分隔的「三碼-三碼-四碼」形狀，\texttt{"0912345678"} 雖然數字個數對，但完全沒有
連字號，不符合 pattern 描述的形狀。
\item 用 \texttt{search()}，因為這裡只在乎「有沒有出現」，不要求整個欄位都得是電話號碼。
\item 把 pattern 換成命名分組版本、其他邏輯不變即可：
\texttt{re.fullmatch(r'(?P<area>\textbackslash d\{3\})-(?P<mid>\textbackslash d\{3\})-%
(?P<last>\textbackslash d\{4\})', text) is not None}，命名分組不影響 \texttt{fullmatch()}
判斷整段字串符不符合的邏輯，只是讓你多了事後取出各段內容的能力。
\end{enumerate}

\paragraph{第~\ref{sec:regex-sub}~節}
\begin{enumerate}
\item 給定 \texttt{text2 = "Customer service email: service@shop.com, complaints go to complain@shop.com"}，
執行 \texttt{re.sub(pattern, "****", text2)}（\texttt{pattern} 是本章教的 email
pattern）之後，會印出 \texttt{"Customer service email: ****, complaints go to ****"}，兩個 email 地址都被換成
\texttt{"****"}，其他文字維持原樣。
\end{enumerate}

\paragraph{第~\ref{sec:regex-split}~節}
\begin{enumerate}
\item \texttt{r'[/,\textbackslash s]+'}，字元集合裡列出斜線、逗號，再搭配
\texttt{\textbackslash s} 涵蓋空格，後面接 \texttt{+} 讓連續出現的分隔符號合併成一個
分隔點。
\item \texttt{['restore', 'multiple', 'spaces']}，\texttt{\textbackslash s+} 表示「一個以上連續的
空白字元」，所以中間那三個空格會被當成同一個分隔點，不會切出空字串。
\item 如果分隔符號固定只有一種、也不會重複出現，用 \texttt{.split()} 就夠了，寫法更簡單；
如果分隔符號有好幾種、或同一種符號可能連續出現好幾次，才需要 \texttt{re.split()} 描述
「這是一種格式」而不是「這是一個固定字元」。
\end{enumerate}

\paragraph{第~\ref{sec:regex-debug}~節}
\begin{enumerate}
\item pattern 本身沒有寫錯，只是它本來就只描述「一段連續的數字」，\texttt{"AB"} 是英文
字母，不符合 \texttt{\textbackslash d+} 的定義，所以理所當然不會被比對進去。如果原本的
目的是想把 \texttt{"AB-123"} 整個編號都抓出來，那才是 pattern 該加上 \texttt{\textbackslash
w+} 或類似的元件去涵蓋英文字母的部分，不是 \texttt{\textbackslash d+} 本身出錯。
\item 優先懷疑真實資料裡藏著測試字串沒有涵蓋到的變化，例如多了空格、少了連字號、大小寫
不同、或多了測試時沒想到的標點符號。這種時候可以把幾筆比對失敗的真實資料印出來，當成新的
測試字串，重複「從小處著手」的除錯習慣，逐步調整 pattern，而不是直接懷疑工具本身壞了。
\end{enumerate}

\chapter{整合應用專案}
\label{ch:automating}

前面十七章，你已經一塊一塊學會了寫程式需要的零件：怎麼存資料、怎麼讓程式重複做事、
怎麼讓程式自己判斷、怎麼把一段邏輯包成可以重複呼叫的函式、怎麼讀寫檔案、怎麼在事情
出錯時不讓程式整個當掉。這一章不會教你任何新的語法，它要做的事情，是帶你把這些零件
一次全部裝進同一台機器裡，看看它們接在一起會變成什麼樣子。

這一章分三段推進。

\begin{itemize}
\item 第一段，我們從一個你我都有過的煩惱開始：桌面或下載資料夾裡堆滿了亂七八糟
命名、亂七八糟分類的檔案，光用滑鼠一個一個整理，可能要花掉一整個下午。我們會先
養成一個專業的習慣，寫程式去「動」使用者的檔案之前，先讓程式「假裝」動一次，只
印出它打算做什麼，讓你在真的按下執行鍵之前，先看清楚後果；接著我們讓分類規則可以
放進一份設定檔，不用改程式碼本身就能讓使用者自己調整規則。
\item 第二段，我們把「假裝」升級成「真的做」，讓程式真的去分類、搬移、甚至幫你把
整理好的檔案打包備份，過程中你會看到迴圈、條件判斷、函式、字串處理、例外處理怎麼
像齒輪一樣咬合在一起，不再是各自孤立的知識點；這一段也會花不少篇幅處理「真實世界
比想像中亂」這件事，例如兩個檔案剛好同名、或者處理到一半有個檔案不見了，一個像樣
的工具得撐得住這些意外，不能一遇到就整個罷工。
\item 第三段，我們把視野從你自己電腦裡的檔案，稍微往外拉一點，寫程式解決瑣事，也
可以主動去網路上抓一點資料回來，變成你整理流程的一部分，而且一次抓好幾個檔案時，
其中一個失敗也不該拖累其他檔案。這一段只是輕輕點一下，讓你有個印象「原來程式也
可以伸手到網路上」，真正深入「軟體怎麼跟軟體對話」這件事，會留到後面第 \ref{ch:api} 章的
API 小專案再展開。
\end{itemize}

這一章的重點始終是「把你已經會的東西串起來」，不是學新東西。

翻到下一節之前先說一句：接下來看到的每一段程式碼，幾乎都是你在前面章節見過的東西
重新排列組合，沒有一個是你沒學過的全新概念。如果讀到一半覺得眼熟，那是應該的。

\section{把散亂的檔案分類列出來：先只印出，不要真的動手}
\label{sec:list-preview}

假設你的下載資料夾裡混著圖片、文件、音樂檔，命名也毫無規則，你想寫一支程式幫你看看
「如果照副檔名分類，每個檔案該歸到哪一類」。這裡刻意先不去搬動任何檔案，只是把資料
夾裡的檔案一個一個列出來，判斷它的副檔名，把「這個檔案應該歸到哪一類」印出來看看。

這聽起來像是在偷懶，但其實是專業工程師處理任何會影響到使用者資料的操作時，共同的
第一步：先看清楚程式打算做什麼，確認邏輯是對的，才放心讓它真的去動手。

\begin{code}{python}{列出資料夾裡的檔案並判斷分類}
import os

for filename in sorted(os.listdir(folder_path)):
    if filename.endswith(extension):
        category = "category_name"
    else:
        category = "Other"
    print(filename, "->", category)
\end{code}

\texttt{os.listdir(folder\_path)} 會回傳一個 list，裡面是這個資料夾底下所有檔案（跟
子資料夾）的名字；外層再包一層 \texttt{sorted()}，是為了讓每次執行的順序都一樣，
資料夾裡檔案的原始排列順序其實跟作業系統有關，不一定照你建立的先後順序，先排序過，
你（跟這本書上印出來的結果）看到的順序才會是可預期的。

\begin{example}{python}{(classify-preview.py) 預覽檔案分類結果}
import os

folder = "downloads_demo"

for filename in sorted(os.listdir(folder)):
    if filename.endswith(".jpg") or filename.endswith(".png"):
        category = "Image"
    elif filename.endswith(".pdf") or filename.endswith(".docx") \
            or filename.endswith(".txt"):
        category = "Document"
    elif filename.endswith(".mp3") or filename.endswith(".wav"):
        category = "Music"
    else:
        category = "Other"
    print(f"{filename} -> {category}")
\end{example}

假設 \texttt{downloads\_demo} 資料夾裡有 \texttt{archive.zip}、\texttt{notes.txt}、
\texttt{report.pdf}、\texttt{song.mp3}、\texttt{vacation.jpg} 這五個檔案：

\iosample{}{archive.zip -> Other \\ notes.txt -> Document \\ report.pdf -> Document \\ song.mp3 -> Music \\ vacation.jpg -> Image \\}

注意 \texttt{archive.zip} 沒有符合任何一個 \texttt{elif} 條件，掉進最後的
\texttt{else}，被歸類成「Other」，這是故意設計的，不用每個副檔名都窮舉，沒對到的
就先放進一個大分類，之後真的有需要再細分。

\keypoint{先印出來看看、不真的動手，是處理任何「會影響使用者資料」的操作（搬移、
刪除、覆寫檔案）之前的專業習慣，跟你是不是新手完全無關。等你以後寫的程式會刪除
檔案、清空資料庫，這個「先乾跑一次確認邏輯，再打開真正執行的開關」的節奏，跟今天在
這裡練的一模一樣。}

\begin{drills}[label=drills:list-preview]
\item 承接 \texttt{classify-preview.py} 的分類邏輯，如果 \texttt{downloads\_demo}
資料夾裡的檔案換成 \texttt{clip.mp4}、\texttt{cover.pdf}、\texttt{photo.png}、
\texttt{todo.txt} 這四個，這段程式會印出什麼（照 \texttt{sorted()} 之後的順序寫出
完整四行）？
\item 如果資料夾裡多了一個檔名整個是 \texttt{mystery}、完全沒有副檔名（連 \texttt{.}
都沒有）的檔案，這段程式會把它歸類到哪一類？為什麼？
\end{drills}

\section{把分類邏輯包成函式，方便重複使用}
\label{sec:classify-function}

上一節的 \texttt{if}/\texttt{elif} 鏈能動，但如果類別變多，會愈疊愈長、愈來愈難讀；
而且「判斷一個檔名該歸到哪一類」這件事，接下來這一章會反覆用到，把它包成一個函式，
之後不管在哪裡都能直接呼叫，不用把同一段判斷邏輯重複貼好幾次。這裡換一種寫法，用
一個對照表（dict）記錄「副檔名對應分類」，比一長串 \texttt{elif} 更容易增減項目。

\begin{code}{python}{把分類邏輯包成函式}
def function_name(filename):
    extension = filename.split(".")[-1].lower()
    return lookup_table.get(extension, "Other")
\end{code}

\texttt{filename.split(".")} 會把檔名以 \texttt{.} 為界切成好幾段，取最後一段
（\texttt{[-1]}）就是副檔名；\texttt{.lower()} 是為了讓 \texttt{Photo.JPG} 這種
大寫副檔名也能對到對照表裡的 \texttt{jpg}（對照表通常都用小寫存）。\texttt{lookup\_table.get(extension, "Other")} 是你在第 \ref{ch:dict-set} 章看過的寫法，查得到就回傳對應的分類，查不到
就回傳預設值 \texttt{"Other"}，不用自己先寫 \texttt{if extension in lookup\_table} 再判斷。

\begin{example}{python}{(classify-function.py) 用函式判斷檔案分類}
import os

EXTENSION_MAP = {
    "jpg": "Image", "png": "Image",
    "pdf": "Document", "docx": "Document", "txt": "Document",
    "mp3": "Music", "wav": "Music",
}


def classify_file(filename):
    extension = filename.split(".")[-1].lower()
    return EXTENSION_MAP.get(extension, "Other")


folder = "downloads_demo"

for filename in sorted(os.listdir(folder)):
    category = classify_file(filename)
    print(f"{filename} -> {category}")
\end{example}

用同樣的五個檔案跑一次，結果跟上一節的 \texttt{classify-preview.py} 完全一樣，這
是應該的，我們只是換了一種「寫」分類邏輯的方式，判斷的結果沒有改變：

\iosample{}{archive.zip -> Other \\ notes.txt -> Document \\ report.pdf -> Document \\ song.mp3 -> Music \\ vacation.jpg -> Image \\}

\stylepoint{像 \texttt{EXTENSION\_MAP} 這種「查表比一長串 \texttt{if}/\texttt{elif}
好維護」的情境，在真實世界的程式碼裡很常見，之後如果想多支援一種副檔名，只要在
dict 裡多加一行 \texttt{"gif": "Image"}，不用去改函式本體的邏輯，這也是把邏輯包成
函式、資料另外放進 dict 的好處之一：程式的「規則」和「怎麼套用規則」分開放。}

\begin{drills}[label=drills:classify-function]
\item 如果想讓 \texttt{EXTENSION\_MAP} 多支援一種副檔名 \texttt{gif}、歸類成
「Image」，只需要改動哪一行？需不需要改動 \texttt{classify\_file()} 函式本體？
\item \texttt{classify\_file("REPORT.PDF")} 會回傳什麼？哪一行程式碼負責讓大寫
副檔名也能正確分類？
\end{drills}

\section{用設定檔讓分類規則可以由使用者調整}
\label{sec:config-file}

上一節把分類規則放進 \texttt{EXTENSION\_MAP} 這個 dict，已經比一長串 \texttt{elif}
好維護很多，但規則還是寫死在程式碼裡，想加一種新的副檔名，得打開 \texttt{.py} 檔案
改程式碼、存檔、重新執行。如果這支工具以後要給不太熟 Python 的人用（例如同事、家人），
你大概不希望他們為了多分類一種檔案，就得打開程式碼修改。這一節要做的事情，是把「規則」
搬到程式碼外面，放進一份使用者自己就能編輯的設定檔。

第 \ref{ch:csv} 章教過 JSON 是一種用巢狀 dict／list 表示資料的純文字格式，這裡剛好派上用場：
我們讓使用者可以自己寫一份 \texttt{config.json}，裡面列出想額外加上的「副檔名對應
分類」，程式啟動時把這份設定檔讀進來，跟原本內建的 \texttt{EXTENSION\_MAP} 合併。

\begin{code}{text}{(config.json) 一份客製化的分類規則}
{
    "gif": "Image",
    "csv": "Data File"
}
\end{code}

這份設定檔的內容，讀者不需要懂任何 Python 語法就看得懂：它在說「副檔名是
\texttt{gif} 的檔案，歸到『Image』；副檔名是 \texttt{csv} 的檔案，歸到『Data File』」，
比起打開 \texttt{.py} 檔案改程式碼，這種格式對非工程背景的使用者友善得多。

\begin{code}{python}{讀取設定檔並跟內建規則合併}
import json

def load_config():
    try:
        with open(config_path, "r") as f:
            return json.loads(f.read())
    except FileNotFoundError:
        return {}


for extension, category in load_config().items():
    lookup_table[extension] = category
\end{code}

\texttt{load\_config()} 這個函式直接嘗試打開設定檔，不先用 \texttt{os.path.exists()}
檢查存不存在，找不到檔案就讓 \texttt{try}/\texttt{except} 接住
\texttt{FileNotFoundError}，回傳一個空的 dict。這正是第 \ref{ch:files-exceptions} 章教過的 EAFP 精神：
設定檔本來就是「有也可以、沒有也可以」的東西，沒有的話就當作使用者沒有要客製化任何
規則，程式照樣能跑，不用因為少了一份設定檔就整支程式當掉。

讀進設定檔之後，用
\texttt{for extension, category in load\_config().items()} 把裡面的每一組規則，一條一條
寫進原本的 \texttt{lookup\_table}，跟第 \ref{ch:dict-set} 章教過的 dict 賦值語法完全一樣，只是這次新增的
key 是從設定檔讀出來的，不是寫死在程式碼裡的常數。

\begin{example}{python}{(config-classify.py) 讓分類規則可以被設定檔調整}
import json
import os

CONFIG_PATH = "config.json"

EXTENSION_MAP = {
    "jpg": "Image", "png": "Image",
    "pdf": "Document", "docx": "Document", "txt": "Document",
    "mp3": "Music", "wav": "Music",
}


def load_config():
    try:
        with open(CONFIG_PATH, "r") as f:
            return json.loads(f.read())
    except FileNotFoundError:
        return {}


for extension, category in load_config().items():
    EXTENSION_MAP[extension] = category


def classify_file(filename):
    extension = filename.split(".")[-1].lower()
    return EXTENSION_MAP.get(extension, "Other")


for filename in sorted(os.listdir("downloads_demo")):
    print(f"{filename} -> {classify_file(filename)}")
\end{example}

假設 \texttt{downloads\_demo} 資料夾裡有 \texttt{archive.zip}、\texttt{cover.gif}、
\texttt{data.csv}、\texttt{notes.txt} 四個檔案，跟上面那份 \texttt{config.json}
放在同一個資料夾底下執行：

\iosample{}{archive.zip -> Other \\ cover.gif -> Image \\ data.csv -> Data File \\ notes.txt -> Document \\}

\texttt{cover.gif} 跟 \texttt{data.csv} 用的分類，都是 \texttt{config.json} 裡新增
進來的規則；\texttt{archive.zip} 的副檔名沒有出現在設定檔或內建的
\texttt{EXTENSION\_MAP} 裡，還是照舊掉進「Other」。如果把 \texttt{config.json} 整個
刪掉再執行一次，\texttt{load\_config()} 會因為 \texttt{FileNotFoundError} 回傳空
dict，\texttt{cover.gif} 跟 \texttt{data.csv} 就會一起變回「Other」，只有
\texttt{notes.txt} 還維持內建規則的「Document」：

\iosample{}{archive.zip -> Other \\ cover.gif -> Other \\ data.csv -> Other \\ notes.txt -> Document \\}

\keypoint{把「規則」（設定檔）跟「怎麼套用規則」（程式邏輯）分開放，是一個很值得
養成的設計習慣。程式邏輯（讀檔案、分類、搬移）通常不太需要常常改，但規則（哪個副
檔名歸哪一類）常常要因人因時調整；分開放之後，調整規則不需要碰程式碼，也不太可能
因為手滑改壞了原本能動的邏輯。}

\begin{drills}[label=drills:config-file]
\item 如果 \texttt{config.json} 的內容改成 \texttt{\{"txt": "Memo"\}}，重新執行
\texttt{config-classify.py} 之後，\texttt{notes.txt} 會被歸類到哪一類？這說明了
設定檔裡的規則跟內建 \texttt{EXTENSION\_MAP} 衝突時，哪一邊會贏？
\item \texttt{load\_config()} 用 \texttt{try}/\texttt{except FileNotFoundError}
接住「設定檔不存在」的狀況，而不是先用 \texttt{os.path.exists()} 檢查一次。這樣寫
除了跟第 \ref{ch:files-exceptions} 章教過的 EAFP 風格一致之外，還有什麼實際好處（提示：想想檢查完到真正
打開檔案之間，狀態可能發生什麼事）？
\end{drills}

\section{真的動手搬移檔案，並接住可能發生的錯誤}
\label{sec:move-files}

前面幾節都只是印出來看看，這一節要把「假裝」升級成「真的做」，用
\texttt{shutil.move()} 把檔案真的搬到分類好的資料夾裡。但只要程式開始真的去動使用者
的檔案，就一定會遇到「這件事可能失敗」的狀況：目的地資料夾可能還不存在、檔案可能正
被別的程式開著占用。

還記得第 \ref{ch:files-exceptions} 章教過的模型嗎，Python 比較習慣「先做，錯了再處理」
（EAFP），跟「先檢查一輪，確定安全才動手」（LBYL）不一樣。這裡就是那個模型正式派上
用場的地方：不用先寫一堆 \texttt{if} 檢查資料夾在不在，直接呼叫 \texttt{shutil.move()}，
如果目的地資料夾不存在，Python 自己會告訴你，我們再用 \texttt{try}/\texttt{except}
接住這個狀況。
\index{shutil 模組／shutil.move()}

\begin{code}{python}{搬移檔案並接住可能發生的錯誤}
import shutil

try:
    shutil.move(source_path, dest_path)
except FileNotFoundError:
    os.makedirs(dest_folder)
    shutil.move(source_path, dest_path)
\end{code}

\texttt{os.makedirs(路徑)} 是 \texttt{os} 模組另一個函式，用來建立一個原本不存在的
新資料夾，如果路徑中間有好幾層資料夾都還不存在，它會一次全部建好，不用自己一層
一層檢查、一層一層建。
\index{os.makedirs()}

如果你想要的是「複製一份、原本的檔案還留著」，而不是「搬過去、原地不留」，把
\texttt{shutil.move()} 換成 \texttt{shutil.copy()} 就好，其他寫法完全一樣，這一章的
整理情境比較適合用「搬移」，把檔案真正歸位，不留兩份佔空間。
\index{shutil.copy()}

\begin{example}{python}{(organize-files.py) 真的搬移檔案並統計結果}
import os
import shutil

SOURCE_FOLDER = "downloads_demo"
DEST_ROOT = "organized_demo"

EXTENSION_MAP = {
    "jpg": "Image", "png": "Image",
    "pdf": "Document", "docx": "Document", "txt": "Document",
    "mp3": "Music", "wav": "Music",
}


def classify_file(filename):
    extension = filename.split(".")[-1].lower()
    return EXTENSION_MAP.get(extension, "Other")


def record_result(counts, category):
    # counts 是呼叫端傳進來的同一個 dict，不是複製一份，
    # 所以在這裡修改也會讓呼叫端看到。
    counts[category] = counts.get(category, 0) + 1


counts = {}

for filename in sorted(os.listdir(SOURCE_FOLDER)):
    source_path = os.path.join(SOURCE_FOLDER, filename)
    category = classify_file(filename)
    dest_folder = os.path.join(DEST_ROOT, category)
    dest_path = os.path.join(dest_folder, filename)

    try:
        shutil.move(source_path, dest_path)
    except FileNotFoundError:
        # 目的地資料夾還不存在，所以先建立資料夾
        # 再重試一次，而不是事先檢查。
        os.makedirs(dest_folder)
        shutil.move(source_path, dest_path)

    record_result(counts, category)

for category, count in counts.items():
    print(f"{category}: {count} file(s)")
\end{example}

第一次執行時，\texttt{organized\_demo} 底下什麼分類資料夾都還沒有，每一種分類第一次
用到都會先撞進 \texttt{except} 區塊、建好資料夾，再搬一次：

\iosample{}{Other: 1 file(s) \\ Document: 2 file(s) \\ Music: 1 file(s) \\ Image: 1 file(s) \\}

\stylepoint{這裡再次示範 Python 偏好「先做，錯了再處理」：直接呼叫
\texttt{shutil.move()}，讓 Python 自己告訴你資料夾不存在，不是每次都先用
\texttt{os.path.exists()} 檢查一輪才動手。檢查完到你真的搬檔案之間，狀態本來就可能
已經變了（資料夾被刪、被搬走），先做再處理反而更可靠。}

\texttt{record\_result()} 這個小函式也值得停下來看一眼：它接收 \texttt{counts} 這個
dict，直接在裡面改東西，沒有 \texttt{return} 任何新的 dict 回去，呼叫端的
\texttt{counts} 卻已經被改過了。這正是第 \ref{ch:functions} 章講過的「名字牌模型」：傳進函式的是同
一個物件的名字牌，不是複製一份，函式內對它做的修改，呼叫端看得到。

如果忘記包 \texttt{try}/\texttt{except}，直接寫 \texttt{shutil.move(source\_path,
dest\_path)}，目的地資料夾又還不存在，會發生什麼事？我們老實示範一次：

\begin{example}{python}{(broken-organize.py) 忘記接住錯誤會發生的事}
import os
import shutil

SOURCE_FOLDER = "downloads_demo"
DEST_ROOT = "organized_demo"

EXTENSION_MAP = {
    "jpg": "Image", "png": "Image",
    "pdf": "Document", "docx": "Document", "txt": "Document",
    "mp3": "Music", "wav": "Music",
}


def classify_file(filename):
    extension = filename.split(".")[-1].lower()
    return EXTENSION_MAP.get(extension, "Other")


for filename in sorted(os.listdir(SOURCE_FOLDER)):
    source_path = os.path.join(SOURCE_FOLDER, filename)
    category = classify_file(filename)
    dest_path = os.path.join(DEST_ROOT, category, filename)
    shutil.move(source_path, dest_path)
    print(f"Moved {filename}")
\end{example}

存檔執行後（\texttt{organized\_demo} 資料夾還是空的），螢幕上只印出一片紅字，一個
檔案都沒有真的被搬到目的地，程式在第一個檔案就整個停住：

\begin{verbatim}
Traceback (most recent call last):
  ...（中間幾行是 shutil 內部呼叫的細節，可以先跳過）
  File "broken_organize.py", line 21, in <module>
    shutil.move(source_path, dest_path)
FileNotFoundError: [Errno 2] No such file or directory:
'organized_demo/Document/notes.txt'
\end{verbatim}

由下往上讀：最後一行 \texttt{FileNotFoundError} 說「找不到這個檔案或路徑」，真正的
原因是目的地那條路徑裡的 \texttt{organized\_demo/Document} 這個資料夾根本還沒被建出來，
\texttt{notes.txt} 這個來源檔案本身好好地在，Python 沒辦法把檔案搬進一個不存在的
資料夾。這正是
\texttt{organize-files.py} 用 \texttt{try}/\texttt{except} 接住、順手用
\texttt{os.makedirs()} 補建資料夾要解決的問題；少了它，程式就直接在第一個檔案當掉，
後面的檔案連試都沒試就整批放棄了。

\begin{drills}[label=drills:move-files]
\item 承接 \texttt{broken-organize.py}：如果 \texttt{organized\_demo} 資料夾（跟底下
所有分類子資料夾）本來就已經存在，這段沒有 \texttt{try}/\texttt{except} 的程式碼還會
出錯嗎？為什麼？
\item 用 \texttt{organize-files.py} 的邏輯，如果 \texttt{downloads\_demo} 資料夾裡只有
\texttt{a.txt}、\texttt{b.jpg}、\texttt{c.txt} 三個檔案，執行後會印出什麼（照
\texttt{sorted()} 之後遇到的順序寫出完整幾行）？
\item \texttt{record\_result()} 裡面用的是 \texttt{counts.get(category, 0) + 1}，
不是先寫 \texttt{if category in counts: ... else: ...}。這兩種寫法哪一種比較接近
這一章、也是第 \ref{ch:files-exceptions} 章教過的 Python 慣用風格？為什麼？
\end{drills}

\section{避免同名檔案被悄悄覆蓋}
\label{sec:avoid-overwrite}

上一節的 \texttt{organize-files.py} 有一個沒說出口的假設：每個要搬移的檔案，在目的
地都還沒有同名的檔案存在。第一次整理一個乾淨的下載資料夾時，這個假設通常成立；但只
要你不是只跑一次，這個假設就很容易被打破，例如你今天整理過一次，過幾天下載資料夾裡
又多了一個檔名剛好也叫 \texttt{notes.txt} 的新檔案，再跑一次 \texttt{organize-files.py}
會發生什麼事？我們老實測試一次。

\begin{code}{python}{測試同名檔案會不會被覆蓋}
import shutil

shutil.move("downloads_demo/notes.txt", "organized_demo/Document/notes.txt")
\end{code}

假設目的地那份舊檔案 \texttt{organized\_demo/Document/notes.txt} 裡存的是「上次整理
留下的舊檔案」，來源那份新檔案 \texttt{downloads\_demo/notes.txt} 裡存的是「這次新
下載的檔案，剛好同名」。執行上面這一行之後，目的地的 \texttt{notes.txt} 會變成什麼
內容？答案是新檔案的內容，舊檔案的內容直接消失，而且\textbf{沒有任何錯誤訊息}，
程式安安靜靜地把舊檔案覆蓋掉了。

這比前面幾節看到的 \texttt{FileNotFoundError} 危險
得多：\texttt{try}/\texttt{except} 能接住的，是 Python 主動丟出來的例外，但覆蓋
同名檔案這件事，Python 完全不覺得有問題，不會丟出任何例外讓你接，資料就這樣悄悄
不見了。

\keypoint{不是每一種「壞事」都會讓程式當掉、丟出例外讓你接住。覆蓋掉一份還有用的舊
檔案，對 Python 來說是完全合法、成功執行的操作，沒有錯誤訊息不代表沒有問題。處理
使用者資料的程式，除了要接住「會噴錯」的狀況，也要主動想清楚「會不會安靜地把東西
弄丟」這種更難察覺的狀況。}

解法是在真正搬移之前，先確認目的地會不會撞名，如果會，就自己想一個不會撞名的新
檔名，例如在檔名後面加上一個編號，變成 \texttt{notes (1).txt}。這裡雖然是先「檢查」
再決定檔名，但跟前面幾節不建議的 LBYL 檢查不太一樣：這不是在檢查「等一下這個操作
會不會失敗」，是在計算「這個檔案最後應該叫什麼名字」，本身就是這個功能需要的邏輯，
不是繞過 EAFP 的偷懶寫法。

\begin{code}{python}{算出一個不會撞名的檔名}
def unique_filename(dest_folder, filename):
    if not os.path.exists(os.path.join(dest_folder, filename)):
        return filename

    parts = filename.split(".")
    if len(parts) > 1:
        name = ".".join(parts[:-1])
        extension = "." + parts[-1]
    else:
        name, extension = filename, ""

    counter = 1
    while True:
        candidate = f"{name} ({counter}){extension}"
        if not os.path.exists(os.path.join(dest_folder, candidate)):
            return candidate
        counter += 1
\end{code}

先用 \texttt{os.path.exists()} 檢查目的地資料夾裡有沒有同名檔案，沒有的話直接回傳
原本的檔名，什麼都不用改。有的話，用跟 \texttt{classify\_file()} 類似的手法把檔名
拆成「名字」跟「副檔名」兩段（例如 \texttt{"notes.txt"} 拆成 \texttt{"notes"} 跟
\texttt{".txt"}），用一個 \texttt{while True} 迴圈，從編號 1 開始一路往上試
（\texttt{notes (1).txt}、\texttt{notes (2).txt}……），每次都檢查這個新檔名存不存
在，第一個「不存在」的候選檔名，就是可以安全使用的新檔名。

\begin{example}{python}{(organize-files-safe.py) 搬移檔案時避免覆蓋同名檔案}
import os
import shutil

SOURCE_FOLDER = "downloads_demo"
DEST_ROOT = "organized_demo"

EXTENSION_MAP = {
    "jpg": "Image", "png": "Image",
    "pdf": "Document", "docx": "Document", "txt": "Document",
    "mp3": "Music", "wav": "Music",
}


def classify_file(filename):
    extension = filename.split(".")[-1].lower()
    return EXTENSION_MAP.get(extension, "Other")


def unique_filename(dest_folder, filename):
    if not os.path.exists(os.path.join(dest_folder, filename)):
        return filename

    parts = filename.split(".")
    if len(parts) > 1:
        name = ".".join(parts[:-1])
        extension = "." + parts[-1]
    else:
        name, extension = filename, ""

    counter = 1
    while True:
        candidate = f"{name} ({counter}){extension}"
        if not os.path.exists(os.path.join(dest_folder, candidate)):
            return candidate
        counter += 1


for filename in sorted(os.listdir(SOURCE_FOLDER)):
    source_path = os.path.join(SOURCE_FOLDER, filename)
    category = classify_file(filename)
    dest_folder = os.path.join(DEST_ROOT, category)

    final_filename = unique_filename(dest_folder, filename)
    if final_filename != filename:
        print(f"{filename} already exists, saving as {final_filename}")
    dest_path = os.path.join(dest_folder, final_filename)

    try:
        shutil.move(source_path, dest_path)
    except FileNotFoundError:
        os.makedirs(dest_folder)
        shutil.move(source_path, dest_path)

print("Done organizing")
\end{example}

延續前面「目的地已經有一個舊的 \texttt{notes.txt}」的情境，\texttt{downloads\_demo}
裡這次還多了一個新的 \texttt{vacation.jpg}（目的地還沒有同名檔案）：

\iosample{}{notes.txt already exists, saving as notes (1).txt \\ Done organizing \\}

執行完之後，\texttt{organized\_demo/Document} 底下會同時有 \texttt{notes.txt}（舊檔案，
內容沒被動過）跟 \texttt{notes (1).txt}（這次新搬過去的檔案），兩份內容都保留了下
來，沒有任何一份被悄悄蓋掉。

\begin{drills}[label=drills:avoid-overwrite]
\item 如果目的地資料夾裡已經同時有 \texttt{notes.txt} 跟 \texttt{notes (1).txt}，
這次又要搬一個新的 \texttt{notes.txt} 過去，\texttt{unique\_filename()} 最後會回傳
什麼檔名？
\item \texttt{unique\_filename()} 一開始用 \texttt{os.path.exists()} 檢查檔案存不
存在，這是不是違反了本章前面一直強調的 EAFP 原則？為什麼？
\end{drills}

\section{一個檔案的失敗，不該讓整批當機}
\label{sec:robust-errors}

\ref{sec:move-files}~節的 \texttt{organize-files.py} 已經會用 \texttt{try}/
\texttt{except} 接住「目的地資料夾不存在」這一種狀況，但那個 \texttt{except}
區塊背後藏著一個沒有明講的假設：只要看到 \texttt{FileNotFoundError}，就一定是
「目的地資料夾不存在」造成的。

如果同一種例外，其實是別的原因造成的呢？例如，處理
到一半，其中一個來源檔案剛好被別的程式（或別的使用者）搬走或刪除了，這時候
\texttt{shutil.move()} 一樣會丟出 \texttt{FileNotFoundError}，但原因完全不是資料
夾不存在。我們示範一次這個情境會怎麼拖垮整支程式。

\begin{example}{python}{(organize-race-broken.py) 同一種例外，兩種完全不同的原因}
import os
import shutil

SOURCE_FOLDER = "downloads_demo"
DEST_ROOT = "organized_demo"

EXTENSION_MAP = {
    "jpg": "Image", "png": "Image",
    "pdf": "Document", "docx": "Document", "txt": "Document",
    "mp3": "Music", "wav": "Music",
}


def classify_file(filename):
    extension = filename.split(".")[-1].lower()
    return EXTENSION_MAP.get(extension, "Other")


for filename in sorted(os.listdir(SOURCE_FOLDER)):
    source_path = os.path.join(SOURCE_FOLDER, filename)
    category = classify_file(filename)
    dest_folder = os.path.join(DEST_ROOT, category)
    dest_path = os.path.join(dest_folder, filename)

    if filename == "song.mp3":
        # 模擬：另一個程式已經先把這個檔案搬走或刪除了
        os.remove(source_path)

    try:
        shutil.move(source_path, dest_path)
    except FileNotFoundError:
        os.makedirs(dest_folder)
        shutil.move(source_path, dest_path)

    print(f"Moved {filename}")
\end{example}

延續前面用過的五個檔案，這裡故意在程式跑到 \texttt{song.mp3} 之前，就先把它從
\texttt{downloads\_demo} 資料夾裡刪掉，模擬「檔案在處理前就消失了」：

\begin{verbatim}
Moved archive.zip
Moved notes.txt
Moved report.pdf
Traceback (most recent call last):
  ...（中間有好幾層「處理某個錯誤的過程中，又發生了新的錯誤」的訊息，
      都是 shutil 內部呼叫的細節，可以先跳過）
  File "organize_race_broken.py", line 20, in <module>
    shutil.move(source_path, dest_path)
FileNotFoundError: [Errno 2] No such file or directory: 'downloads_demo/song.mp3'
\end{verbatim}

\texttt{archive.zip}、\texttt{notes.txt}、\texttt{report.pdf} 都順利搬完了，輪到
\texttt{song.mp3} 時，第一次 \texttt{shutil.move()} 因為來源檔案不見了而丟出
\texttt{FileNotFoundError}，程式照原本的邏輯認定「這一定是資料夾不存在」，呼叫
\texttt{os.makedirs()} 補建資料夾之後重試，但重試的這一次還是失敗，這次的問題根本
出在來源檔案本身，跟資料夾無關。這次的 \texttt{FileNotFoundError} 沒有被任何
\texttt{except} 接住，整支程式就在這裡中止，剩下還沒處理的 \texttt{vacation.jpg}
連試都沒試就被放棄了，跟本章一開始那個完全沒包 \texttt{try}/\texttt{except} 的
\texttt{broken-organize.py} 落到一樣的下場。

這裡的教訓是：光靠例外的「類型」判斷該怎麼處理是不夠的，同一個
\texttt{FileNotFoundError} 背後可能是完全不同的原因。修正方式有兩個重點：

\begin{itemize}
\item 把第二次嘗試也包進 \texttt{try}/\texttt{except}，這次如果還是失敗，代表問題
出在來源檔案，記錄下來就好，不用再重試。
\item 把整段搬移邏輯放進迴圈本身的錯誤處理裡，讓一個檔案的失敗只影響那一個檔案，
用 \texttt{continue} 跳到下一輪，繼續處理其他檔案，不要讓整支程式因為一個檔案而
停擺。
\end{itemize}

\begin{example}{python}{(organize-robust.py) 一個檔案失敗不影響其他檔案}
import os
import shutil

SOURCE_FOLDER = "downloads_demo"
DEST_ROOT = "organized_demo"

EXTENSION_MAP = {
    "jpg": "Image", "png": "Image",
    "pdf": "Document", "docx": "Document", "txt": "Document",
    "mp3": "Music", "wav": "Music",
}


def classify_file(filename):
    extension = filename.split(".")[-1].lower()
    return EXTENSION_MAP.get(extension, "Other")


def record_result(counts, category):
    counts[category] = counts.get(category, 0) + 1


def record_error(errors, filename, reason):
    errors.append((filename, reason))


counts = {}
errors = []

for filename in sorted(os.listdir(SOURCE_FOLDER)):
    source_path = os.path.join(SOURCE_FOLDER, filename)
    category = classify_file(filename)
    dest_folder = os.path.join(DEST_ROOT, category)
    dest_path = os.path.join(dest_folder, filename)

    if filename == "song.mp3":
        os.remove(source_path)  # 模擬檔案在搬移前就消失了

    try:
        shutil.move(source_path, dest_path)
    except FileNotFoundError:
        try:
            os.makedirs(dest_folder)
            shutil.move(source_path, dest_path)
        except FileNotFoundError:
            record_error(errors, filename, "source file disappeared before it could be moved")
            continue

    record_result(counts, category)

for category, count in counts.items():
    print(f"{category}: {count} file(s)")

if errors:
    print(f"{len(errors)} file(s) failed to move:")
    for filename, reason in errors:
        print(f"  {filename}: {reason}")
\end{example}

\iosample{}{Other: 1 file(s) \\ Document: 2 file(s) \\ Image: 1 file(s) \\ 1 file(s) failed to move: \\ \ \ song.mp3: source file disappeared before it could be moved \\}

這次 \texttt{song.mp3} 失敗了，但 \texttt{archive.zip}、\texttt{notes.txt}、
\texttt{report.pdf}、\texttt{vacation.jpg} 全部都正常整理完成，程式最後還老實印出
一份「哪些檔案失敗了、為什麼失敗」的清單，讓使用者知道要自己去確認那個檔案的下落，
而不是打開一片 Traceback 卻完全不知道整理到哪裡就停了。

\texttt{record\_error()} 這個函式跟 \ref{sec:move-files}~節的 \texttt{record\_result()}
是同一種寫法：接收一個 list，直接在裡面加東西，不用回傳新的 list。這裡的巢狀
\texttt{try}/\texttt{except} 也值得停下來看一眼：外層 \texttt{try} 對應「第一次嘗試
搬移」，內層 \texttt{try} 對應「補建資料夾之後再試一次」，兩層各自處理各自可能發生
的 \texttt{FileNotFoundError}，內層失敗才是真正束手無策的狀況，這時候才記錄錯誤、
放棄這個檔案。

\stylepoint{程式裡的 \texttt{except FileNotFoundError} 一直都寫得很明確，指名要接
的是哪一種例外，不是籠統地寫一個空白的 \texttt{except:} 把所有錯誤都吞掉。籠統接住
所有錯誤聽起來很安全，實際上很危險：它連你自己程式碼裡打錯字造成的 \texttt{NameError}
都會一起吞掉，讓你以為程式「沒問題」，其實是某個地方悄悄壞掉了卻沒人發現。只接你真
的預期會發生、也知道怎麼處理的例外類型，其他的讓它照樣噴出來，你才看得到。}

\begin{drills}[label=drills:robust-errors]
\item \texttt{organize-robust.py} 裡的內層 \texttt{try}/\texttt{except} 如果拿掉，
直接讓內層的 \texttt{os.makedirs()} 跟 \texttt{shutil.move()} 兩行沒有包在任何
\texttt{try} 裡面，重跑一次會發生什麼事？
\item 如果 \texttt{downloads\_demo} 資料夾裡有兩個檔案在處理前就消失了（例如
\texttt{song.mp3} 跟 \texttt{vacation.jpg}），\texttt{organize-robust.py} 執行完
之後，\texttt{errors} 這個 list 裡會有幾筆紀錄？
\end{drills}

\section{把整理結果打包備份}
\label{sec:backup-zip}

檔案整理到一半，很多人會多做一件事，把整理好的成果打包成一份備份，避免之後不小心
弄丟。Python 標準函式庫的 \texttt{zipfile} 模組可以直接在程式裡建立一個 ZIP 壓縮檔，
不用另外開檔案總管手動壓縮。這裡只示範最常用的寫法：建立一個新的 ZIP 檔，把資料夾
裡的檔案一個一個加進去；\texttt{zipfile} 其實還有解壓縮、讀取既有 ZIP 內容等其他功能，
用得到的時候查文件就好，不用在這裡一次學完。
\index{zipfile 模組／zipfile.ZipFile()}

\begin{code}{python}{用 zipfile 打包資料夾}
import zipfile

with zipfile.ZipFile(zip_name, "w") as zip_file:
    for filename in os.listdir(folder_path):
        zip_file.write(os.path.join(folder_path, filename))
\end{code}

跟第 \ref{ch:files-exceptions} 章讀寫檔案時一樣，這裡也用 \texttt{with} 包住整個操作，離開 \texttt{with}
區塊時，Python 會自動幫你把 ZIP 檔正確收尾關好，不用自己記得呼叫關檔的方法。

\begin{example}{python}{(backup-zip.py) 打包整理好的資料夾}
import os
import zipfile

SOURCE_FOLDER = "organized_demo"
ZIP_NAME = "organized_backup.zip"

with zipfile.ZipFile(ZIP_NAME, "w") as backup_zip:
    for category in sorted(os.listdir(SOURCE_FOLDER)):
        category_folder = os.path.join(SOURCE_FOLDER, category)
        for filename in sorted(os.listdir(category_folder)):
            file_path = os.path.join(category_folder, filename)
            backup_zip.write(file_path)
            print(f"Added to backup: {file_path}")

print(f"Backup complete: {ZIP_NAME}")
\end{example}

承接前面整理好的 \texttt{organized\_demo} 資料夾（裡面已經分好「Document」「Image」
「Music」「Other」四個子資料夾）：

\iosample{}{Added to backup: organized\_demo/Document/notes.txt \\ Added to backup: organized\_demo/Document/report.pdf \\ Added to backup: organized\_demo/Image/vacation.jpg \\ Added to backup: organized\_demo/Music/song.mp3 \\ Added to backup: organized\_demo/Other/archive.zip \\ Backup complete: organized\_backup.zip \\}

這裡刻意用兩層 \texttt{os.listdir()}（先列出分類資料夾，再列出每個分類資料夾裡的
檔案），不是因為 \texttt{os} 模組沒有更直接的寫法，它其實還有一個
\texttt{os.walk()} 可以一次走訪整棵資料夾樹，但這一章只會用到兩層深度，先用你已經
會的 \texttt{os.listdir()} 疊兩層就夠了，不用為了這一個場景多學一個新函式。

\begin{drills}[label=drills:backup-zip]
\item 如果 \texttt{organized\_demo} 底下有一個分類資料夾是空的（一個檔案都沒有），
\texttt{backup-zip.py} 執行時會不會出錯？備份檔裡會不會有這個分類的紀錄？
\end{drills}

\section{避免每次備份互相覆蓋}
\label{sec:backup-timestamp}

上一節的 \texttt{backup-zip.py} 每次都把備份存成同一個檔名
\texttt{organized\_backup.zip}，如果你今天整理一次、備份一次，過幾天又整理一次、
又備份一次，第二次的備份會直接把第一次的備份整個覆蓋掉，跟 \ref{sec:avoid-overwrite}~
節看到的問題一樣：Python 不會因為覆蓋一個舊的 ZIP 檔而丟出任何例外，你可能過了很久
才發現，早就想留著的那份舊備份，其實已經不在了。

第 \ref{ch:modules} 章教過的 \texttt{datetime} 模組這裡剛好能派上用場：把今天的日期放進備份檔名
裡，每天最多產生一份同名的備份檔，不同天的備份自然不會互相覆蓋。

\begin{code}{python}{把日期放進備份檔名}
import datetime

today = datetime.date.today()
zip_name = f"organized_backup_{today}.zip"
\end{code}

\texttt{datetime.date.today()} 回傳一個代表「今天」的日期物件，直接放進 f-string
裡，會自動印成 \texttt{YYYY-MM-DD} 的格式，例如 \texttt{2026-09-02}，剛好可以當成
檔名的一部分，不用自己再手動組字串。

\begin{example}{python}{(backup-zip-dated.py) 幫每次備份加上日期}
import datetime
import os
import zipfile

SOURCE_FOLDER = "organized_demo"

today = datetime.date.today()
zip_name = f"organized_backup_{today}.zip"

with zipfile.ZipFile(zip_name, "w") as backup_zip:
    for category in sorted(os.listdir(SOURCE_FOLDER)):
        category_folder = os.path.join(SOURCE_FOLDER, category)
        for filename in sorted(os.listdir(category_folder)):
            backup_zip.write(os.path.join(category_folder, filename))

print(f"Backup complete: {zip_name}")
\end{example}

下面是作者示範當下（2026 年 9 月 2 日）執行的結果，你自己執行時看到的日期，會是你
執行當天的日期：

\iosample{}{Backup complete: organized\_backup\_2026-09-02.zip \\}

\keypoint{這裡刻意檢查了「輸出裡有沒有隨執行當下改變的內容」這件事：日期一定會隨著
你執行的那一天而不同，展示這種範例時，最好像這裡一樣講清楚「這是作者示範當下看到的
結果」，讀者自己執行時看到不一樣的日期是正常的，不是程式壞了。}

\begin{drills}[label=drills:backup-timestamp]
\item 如果同一天執行了兩次 \texttt{backup-zip-dated.py}，會發生什麼事？這樣的檔名
規則有沒有辦法完全避免同一天的兩次備份互相覆蓋？（提示：可以參考
\ref{sec:avoid-overwrite}~節的做法。）
\end{drills}

\section{把視野拉到網路上：用 requests.get() 下載檔案}
\label{sec:requests-download}

到目前為止，程式處理的都是你電腦裡本來就有的檔案。但寫程式解決瑣事，除了處理硬碟
裡已經有的東西，也可以主動去網路上抓一點資料回來，變成整理流程的一部分。這裡只示
範最基本的用法：用 \texttt{requests} 套件發一個請求，把下載回來的內容存成本機檔案。

至於伺服器回傳的狀態碼代表什麼、怎麼跟一個正式的 API 對話，這些留給第 \ref{ch:api} 章的專案
再展開，這裡只是讓你有個印象「原來程式也可以伸手到網路上」。
\index{requests 套件／requests.get()}

\begin{code}{python}{用 requests.get() 下載檔案}
import requests

try:
    response = requests.get(url, timeout=5)
except requests.exceptions.RequestException as e:
    print("Download failed:", e)
else:
    with open(save_filename, "wb") as f:
        f.write(response.content)
\end{code}

跟前面幾節接住檔案操作錯誤的道理一樣，網路請求也是一種「跟外部世界打交道、本來就
可能失敗」的操作，網址可能打錯字、對方伺服器可能沒回應、你的網路可能剛好斷線。
一樣不用先想辦法「檢查網路正不正常」，直接發請求，失敗了讓 \texttt{try}/
\texttt{except} 接住就好。

這裡的 \texttt{except 例外類型 as e} 多做了一件事：
\texttt{as e} 會把捕捉到的這個例外物件存進一個叫 \texttt{e} 的變數裡，之後可以用
\texttt{print(e)} 印出更詳細的錯誤訊息，不只是知道「出錯了」，還能看到 Python
具體回報了什麼。
\index{except ... as e}

注意這裡開檔案用的是 \texttt{"wb"} 模式（\texttt{w} 是
寫入、\texttt{b} 代表二進位），跟你在第 \ref{ch:files-exceptions} 章看過的純文字 \texttt{"w"} 不一樣，
下載回來的內容不一定是純文字（圖片、壓縮檔都不是），用 \texttt{"wb"} 才不會把資料
寫壞。

\begin{example}{python}{(download-file.py) 下載一張圖片存到本機}
import requests

URL = "https://www.python.org/static/img/python-logo.png"
FILENAME = "python-logo.png"

try:
    response = requests.get(URL, timeout=5)
except requests.exceptions.RequestException as e:
    print("Download failed:", e)
else:
    with open(FILENAME, "wb") as f:
        f.write(response.content)
    print(f"Downloaded {FILENAME}, size {len(response.content)} bytes")
\end{example}

\iosample{}{Downloaded python-logo.png, size 15770 bytes \\}

下載回來的 \texttt{python-logo.png} 存進資料夾之後，就是一個跟其他檔案沒有兩樣的
本機檔案，你完全可以把它丟進前面 \ref{sec:classify-function}、\ref{sec:backup-zip}~
節建立的分類、整理流程裡，跟其他檔案一起被分類、搬移、打包備份。這也是這一節故意
做得很淺的原因：它只負責「把網路上的東西變成一個本機檔案」，剩下的整理工作，交給你
已經會的技巧接手就好。

\begin{drills}[label=drills:requests-download]
\item 如果把 \texttt{download-file.py} 的 \texttt{try}/\texttt{except} 拿掉，直接寫
\texttt{response = requests.get(URL, timeout=5)}，而這次 \texttt{URL} 打錯字（指向
一個根本不存在的網址），執行後會發生什麼事？跟前面 \ref{sec:move-files}~節忘記包
\texttt{try}/\texttt{except} 時看到的情況，有什麼共同點？
\end{drills}

\section{一次下載好幾個檔案，個別失敗不影響其他檔案}
\label{sec:download-multiple}

上一節只下載了一個檔案，但真實情境裡，你常常一次要抓好幾個檔案回來，例如好幾張圖片、
好幾份報表。如果把上一節的寫法直接放進一個迴圈，逐一下載一份清單裡的每個網址，會遇到
跟 \ref{sec:robust-errors}~節一樣的問題：只要其中一個網址失敗，整支程式就會在那裡
中止，後面清單裡還沒下載的網址，連試都沒試就被放棄了。

這一節示範怎麼讓下載這件事也
擁有跟搬移檔案一樣的韌性：一個網址失敗，跳過它，繼續處理下一個。

\begin{example}{python}{(download-multiple.py) 依序下載多個檔案，個別失敗就跳過}
import os
import requests

DOWNLOAD_URLS = [
    "https://www.python.org/static/img/python-logo.png",
    "https://no-such-site.invalid/x.png",
    "https://www.python.org/static/favicon.ico",
]
DEST_FOLDER = "downloads_demo"


def download_file(url, dest_folder):
    filename = url.split("/")[-1]
    dest_path = os.path.join(dest_folder, filename)
    response = requests.get(url, timeout=5)
    with open(dest_path, "wb") as f:
        f.write(response.content)
    return dest_path


downloaded = 0
skipped = 0

for url in DOWNLOAD_URLS:
    try:
        dest_path = download_file(url, DEST_FOLDER)
    except requests.exceptions.RequestException as e:
        print(f"Download failed, skipping {url} ({type(e).__name__})")
        skipped += 1
        continue
    print(f"Downloaded {dest_path}")
    downloaded += 1

print(f"Done: {downloaded} succeeded, {skipped} skipped")
\end{example}

\iosample{}{Downloaded downloads\_demo/python-logo.png \\ Download failed, skipping https://no-such-site.invalid/x.png (ConnectionError) \\ Downloaded downloads\_demo/favicon.ico \\ Done: 2 succeeded, 1 skipped \\}

清單裡第二個網址刻意寫錯，\texttt{requests.get()} 連不上去，會丟出一種叫
\texttt{ConnectionError} 的例外。這裡的 \texttt{except} 接的是比較廣的
\texttt{requests.exceptions.RequestException}，第 \ref{ch:api} 章會再提到，
\texttt{ConnectionError} 只是 \texttt{RequestException} 底下好幾種更具體的錯誤
之一，用比較廣的類型接住，可以一次涵蓋連線失敗、逾時等好幾種網路相關的錯誤。這裡
印錯誤訊息時用的是 \texttt{type(e).\_\_name\_\_}（例外物件的型別名稱），不是把
整段冗長的錯誤細節都印出來，讓輸出保持簡短好讀，同時還是能看出「是哪一種錯誤」。
\index{type(e).\_\_name\_\_}

\begin{drills}[label=drills:download-multiple]
\item 如果 \texttt{DOWNLOAD\_URLS} 這份清單裡有三個網址全部都連不上，
\texttt{download-multiple.py} 執行完之後，\texttt{downloaded} 跟 \texttt{skipped}
分別會是多少？
\item \texttt{download\_file()} 函式本身沒有包任何 \texttt{try}/\texttt{except}，
是呼叫它的迴圈負責接住錯誤。這樣設計跟把 \texttt{try}/\texttt{except} 直接寫進
\texttt{download\_file()} 函式內部相比，有什麼好處？
\end{drills}

\section{完整走一次：把所有片段串成一支程式}
\label{sec:full-integration}

這一節不教任何新技巧，只是把前面幾節的片段接起來，讓你看看整套流程真正跑起來長什麼
樣子：先去網路上抓一個檔案回來，放進待整理的資料夾；把資料夾裡所有檔案分類、搬移，
順手統計每一類有幾個檔案；最後把整理好的結果打包成一份備份。

\begin{example}{python}{(organize-and-backup.py) 完整的檔案整理流程}
import os
import shutil
import zipfile
import requests

DOWNLOAD_URL = "https://www.python.org/static/img/python-logo.png"
SOURCE_FOLDER = "downloads_demo"
DEST_ROOT = "organized_demo"
ZIP_NAME = "organized_backup.zip"

EXTENSION_MAP = {
    "jpg": "Image", "png": "Image",
    "pdf": "Document", "docx": "Document", "txt": "Document",
    "mp3": "Music", "wav": "Music",
}


def classify_file(filename):
    extension = filename.split(".")[-1].lower()
    return EXTENSION_MAP.get(extension, "Other")


def record_result(counts, category):
    counts[category] = counts.get(category, 0) + 1


# 第一步：先從網路上下載一個檔案，加進待整理的資料夾
try:
    response = requests.get(DOWNLOAD_URL, timeout=5)
except requests.exceptions.RequestException as e:
    print("Download failed, skipping this file:", e)
else:
    downloaded_path = os.path.join(SOURCE_FOLDER, "python-logo.png")
    with open(downloaded_path, "wb") as f:
        f.write(response.content)
    print(f"Downloaded {downloaded_path}")

# 第二步：把資料夾裡所有檔案分類、搬移，並統計結果
counts = {}

for filename in sorted(os.listdir(SOURCE_FOLDER)):
    source_path = os.path.join(SOURCE_FOLDER, filename)
    category = classify_file(filename)
    dest_folder = os.path.join(DEST_ROOT, category)
    dest_path = os.path.join(dest_folder, filename)

    try:
        shutil.move(source_path, dest_path)
    except FileNotFoundError:
        os.makedirs(dest_folder)
        shutil.move(source_path, dest_path)

    record_result(counts, category)

for category, count in counts.items():
    print(f"{category}: {count} file(s)")

# 第三步：把整理好的結果打包成備份
with zipfile.ZipFile(ZIP_NAME, "w") as backup_zip:
    for category in sorted(os.listdir(DEST_ROOT)):
        category_folder = os.path.join(DEST_ROOT, category)
        for filename in sorted(os.listdir(category_folder)):
            backup_zip.write(os.path.join(category_folder, filename))

print(f"Backup complete: {ZIP_NAME}")
\end{example}

假設一開始 \texttt{downloads\_demo} 資料夾裡有 \texttt{vacation.jpg}、
\texttt{notes.txt}、\texttt{report.pdf}、\texttt{song.mp3} 四個檔案：

\iosample{}{Downloaded downloads\_demo/python-logo.png \\ Document: 2 file(s) \\ Image: 2 file(s) \\ Music: 1 file(s) \\ Backup complete: organized\_backup.zip \\}

從頭看一遍會發現：這支程式裡沒有一行是這一章第一次出現的全新語法，迴圈、條件判斷、
函式、字串方法、\texttt{try}/\texttt{except}、\texttt{with}，全部都是前面某一章教過
的東西，只是這次它們一起解決同一個真實的任務，不再是各自獨立示範。這就是這一章
真正想讓你看到的畫面：你已經有的技巧，接在一起可以做到比任何單一技巧看起來都更完整
的事情。

\section{進階版：把新學到的韌性也一起裝進去}
\label{sec:advanced-integration}

上一節的 \texttt{organize-and-backup.py} 已經是一個完整能動的工具，但這一章後半段
額外學到的三件事，設定檔客製化、避免同名覆蓋、單一檔案失敗不拖累整批，都還沒被放
進這支完整流程裡。這一節做最後一次整合：在保留原本三步驟骨架的前提下，把這一章學到
的所有加強都疊上去，順便也把 \ref{sec:download-multiple}~節「一次下載好幾個檔案」的
能力接進來，變成一個更貼近真實世界會怎麼寫的版本。

\begin{example}{python}{(organize-and-backup-plus.py) 整合所有新功能的完整流程}
import datetime
import json
import os
import shutil
import zipfile

import requests

DOWNLOAD_URLS = [
    "https://www.python.org/static/img/python-logo.png",
    "https://no-such-site.invalid/x.png",
]
SOURCE_FOLDER = "downloads_demo"
DEST_ROOT = "organized_demo"
CONFIG_PATH = "config.json"

EXTENSION_MAP = {
    "jpg": "Image", "png": "Image",
    "pdf": "Document", "docx": "Document", "txt": "Document",
    "mp3": "Music", "wav": "Music",
}


def load_config():
    try:
        with open(CONFIG_PATH, "r") as f:
            return json.loads(f.read())
    except FileNotFoundError:
        return {}


for extension, category in load_config().items():
    EXTENSION_MAP[extension] = category


def classify_file(filename):
    extension = filename.split(".")[-1].lower()
    return EXTENSION_MAP.get(extension, "Other")


def unique_filename(dest_folder, filename):
    if not os.path.exists(os.path.join(dest_folder, filename)):
        return filename
    parts = filename.split(".")
    if len(parts) > 1:
        name = ".".join(parts[:-1])
        extension = "." + parts[-1]
    else:
        name, extension = filename, ""
    counter = 1
    while True:
        candidate = f"{name} ({counter}){extension}"
        if not os.path.exists(os.path.join(dest_folder, candidate)):
            return candidate
        counter += 1


def record_result(counts, category):
    counts[category] = counts.get(category, 0) + 1


def record_error(errors, filename, reason):
    errors.append((filename, reason))


def download_file(url, dest_folder):
    filename = url.split("/")[-1]
    dest_path = os.path.join(dest_folder, filename)
    response = requests.get(url, timeout=5)
    with open(dest_path, "wb") as f:
        f.write(response.content)
    return dest_path


# 第一步：下載要處理的檔案，單一失敗不影響其他檔案
for url in DOWNLOAD_URLS:
    try:
        downloaded_path = download_file(url, SOURCE_FOLDER)
    except requests.exceptions.RequestException as e:
        print(f"Download failed, skipping {url} ({type(e).__name__})")
        continue
    print(f"Downloaded {downloaded_path}")

# 第二步：分類並搬移，避免覆蓋同名檔案，且單一檔案失敗不影響其他檔案
counts = {}
errors = []

for filename in sorted(os.listdir(SOURCE_FOLDER)):
    source_path = os.path.join(SOURCE_FOLDER, filename)
    category = classify_file(filename)
    dest_folder = os.path.join(DEST_ROOT, category)

    final_filename = unique_filename(dest_folder, filename)
    dest_path = os.path.join(dest_folder, final_filename)

    try:
        shutil.move(source_path, dest_path)
    except FileNotFoundError:
        try:
            os.makedirs(dest_folder)
            shutil.move(source_path, dest_path)
        except FileNotFoundError:
            record_error(errors, filename, "source file disappeared before it could be moved")
            continue

    record_result(counts, category)

for category, count in counts.items():
    print(f"{category}: {count} file(s)")

if errors:
    print(f"{len(errors)} file(s) failed to move:")
    for filename, reason in errors:
        print(f"  {filename}: {reason}")

# 第三步：把整理好的結果打包成有日期的備份，避免覆蓋前一次的備份
today = datetime.date.today()
zip_name = f"organized_backup_{today}.zip"

with zipfile.ZipFile(zip_name, "w") as backup_zip:
    for category in sorted(os.listdir(DEST_ROOT)):
        category_folder = os.path.join(DEST_ROOT, category)
        for filename in sorted(os.listdir(category_folder)):
            backup_zip.write(os.path.join(category_folder, filename))

print(f"Backup complete: {zip_name}")
\end{example}

假設 \texttt{downloads\_demo} 資料夾裡一開始有 \texttt{vacation.jpg}、
\texttt{notes.txt}、\texttt{report.pdf}、\texttt{song.mp3} 四個檔案，目的地資料夾
還沒有任何同名檔案，下面是作者示範當下（2026 年 9 月 2 日）執行的結果：

\iosample{}{Downloaded downloads\_demo/python-logo.png \\ Download failed, skipping https://no-such-site.invalid/x.png (ConnectionError) \\ Document: 2 file(s) \\ Image: 2 file(s) \\ Music: 1 file(s) \\ Backup complete: organized\_backup\_2026-09-02.zip \\}

比較這支程式跟上一節的 \texttt{organize-and-backup.py}，骨架完全沒有變，一樣是「先
下載、再分類搬移、最後打包備份」三步驟，變的只是每一步底下多墊了幾層防護：

\begin{itemize}
\item 下載那一步現在能處理一份清單、個別失敗不影響其他網址。
\item 分類搬移那一步現在會避開同名覆蓋、單一檔案失敗也不會拖累整批。
\item 打包備份那一步現在會自動避開覆蓋前一天的備份。
\end{itemize}

這正是這一章從頭到尾想傳達的重點：一個能用的工具，跟一個經得起現實世界各種意外考驗的工具，中間
的差距通常不是全新的知識，是願不願意多花力氣，把「萬一發生了呢」一項一項想過、接住。

\begin{drills}[label=drills:advanced-integration]
\item 如果 \texttt{organized\_demo} 資料夾裡已經有一個 \texttt{vacation.jpg}（跟這次
要下載並整理的檔案剛好同名），\texttt{organize-and-backup-plus.py} 執行後，這個檔案
最後會被存成什麼名字？
\item 這支程式跟 \ref{sec:full-integration}~節的 \texttt{organize-and-backup.py}
相比，多用到了哪三個本章新學的函式？各自對應解決哪一種「真實世界比較亂」的狀況？
\end{drills}

\section*{練習解答}
\addcontentsline{toc}{section}{練習解答}

\paragraph{第~\ref{sec:list-preview}~節}
\begin{enumerate}
\item 照 \texttt{sorted()} 之後的順序（\texttt{clip.mp4}、\texttt{cover.pdf}、
\texttt{photo.png}、\texttt{todo.txt}），會印出：\texttt{clip.mp4 -> Other}、
\texttt{cover.pdf -> Document}、\texttt{photo.png -> Image}、\texttt{todo.txt -> Document}。
\item 會被歸類到「Other」，因為 \texttt{mystery} 這個檔名沒有結尾符合
\texttt{.jpg}／\texttt{.png}／\texttt{.pdf}／\texttt{.docx}／\texttt{.txt}／
\texttt{.mp3}／\texttt{.wav} 任何一個，程式一路掉進最後的 \texttt{else}。
\end{enumerate}

\paragraph{第~\ref{sec:classify-function}~節}
\begin{enumerate}
\item 只需要在 \texttt{EXTENSION\_MAP} 這個 dict 裡多加一行 \texttt{"gif": "Image",}，
不用改動 \texttt{classify\_file()} 函式本體，這正是把規則放進 dict、邏輯放進函式
分開處理的好處。
\item 會回傳 \texttt{"Document"}。負責的是 \texttt{extension = filename.split(".")[-1]
.lower()} 這一行裡的 \texttt{.lower()}，把切出來的副檔名統一轉成小寫，才能對到
\texttt{EXTENSION\_MAP} 裡小寫的 key。
\end{enumerate}

\paragraph{第~\ref{sec:config-file}~節}
\begin{enumerate}
\item 會被歸類到「Memo」。設定檔的規則是在內建的 \texttt{EXTENSION\_MAP} 建立好之
後才用 \texttt{for} 迴圈一條一條寫進去，同一個 key（這裡是 \texttt{"txt"}）後寫入的
值會覆蓋掉先前的值，所以設定檔的規則會贏過內建規則。
\item 好處是就算設定檔在檢查完到真正打開之間被刪除（例如使用者手滑刪掉了），程式也
不會因為一個檢查和實際操作之間的時間差而出錯，直接嘗試打開、失敗了再處理，比先檢查
一次更可靠。
\end{enumerate}

\paragraph{第~\ref{sec:move-files}~節}
\begin{enumerate}
\item 不會出錯，\texttt{FileNotFoundError} 只有在目的地資料夾（例如
\texttt{organized\_demo/Document}）不存在時才會發生；如果所有分類子資料夾都已經先建好，
\texttt{shutil.move()} 每次都能直接成功，程式會正常跑完，只是少了自動補建資料夾這個
保險。
\item 依序處理 \texttt{a.txt}（Document）、\texttt{b.jpg}（Image）、\texttt{c.txt}
（Document），會印出：\texttt{Document: 2 file(s)}、\texttt{Image: 1 file(s)}。
\item \texttt{counts.get(category, 0) + 1} 比較接近 Python 慣用風格，這是 EAFP
精神的另一種展現，把「查得到就用，查不到就用預設值」包成一個方法呼叫，不用自己先寫
\texttt{if category in counts} 再分兩條路徑處理。
\end{enumerate}

\paragraph{第~\ref{sec:avoid-overwrite}~節}
\begin{enumerate}
\item 會回傳 \texttt{"notes (2).txt"}，因為 \texttt{unique\_filename()} 用
\texttt{while True} 從編號 1 開始一路往上試，\texttt{notes (1).txt} 已經被佔用了，
下一個沒被佔用的候選就是編號 2。
\item 不算違反。EAFP 原則要避免的，是用檢查來「猜測」某個本來就可能失敗的操作會不會
成功；這裡的檢查不是為了猜測會不會失敗，是這個功能本身需要的邏輯，用來計算出一個
還沒被使用過的檔名，跟避免失敗無關。
\end{enumerate}

\paragraph{第~\ref{sec:robust-errors}~節}
\begin{enumerate}
\item 內層的 \texttt{os.makedirs()} 或 \texttt{shutil.move()} 如果再次失敗，會丟出
沒有被任何 \texttt{except} 接住的例外，程式會在處理 \texttt{song.mp3} 時整個中止，
跟本節一開始的 \texttt{organize-race-broken.py} 一樣，\texttt{vacation.jpg} 也不會
被處理到。
\item 會有兩筆紀錄，一筆是 \texttt{song.mp3}，一筆是 \texttt{vacation.jpg}，因為
每個檔案的搬移都各自包在自己的 \texttt{try}/\texttt{except} 裡，一個檔案失敗只會
讓那一個檔案被記進 \texttt{errors}、用 \texttt{continue} 跳過，不會影響其他檔案繼續
被正常處理。
\end{enumerate}

\paragraph{第~\ref{sec:backup-zip}~節}
\begin{enumerate}
\item 不會出錯，內層的 \texttt{for filename in sorted(os.listdir(category\_folder))}
在空資料夾裡就是跑零次，直接跳過；備份檔裡不會有這個分類資料夾的任何紀錄，因為
\texttt{zipfile} 只會寫入你明確呼叫 \texttt{.write()} 加進去的檔案，沒有檔案就沒有
東西可以加。
\end{enumerate}

\paragraph{第~\ref{sec:backup-timestamp}~節}
\begin{enumerate}
\item 同一天執行兩次，兩次算出來的 \texttt{today} 是同一個日期，檔名
\texttt{organized\_backup\_2026-09-02.zip} 完全相同，第二次還是會把第一次的備份
整個覆蓋掉。這個檔名規則只能避免「不同天」的備份互相覆蓋，沒辦法避免「同一天」的多次
備份互相覆蓋；真的要完全避免，可以搭配 \ref{sec:avoid-overwrite}~節教的
\texttt{unique\_filename()}，讓同一天的第二次備份自動變成
\texttt{organized\_backup\_2026-09-02 (1).zip}。
\end{enumerate}

\paragraph{第~\ref{sec:requests-download}~節}
\begin{enumerate}
\item 程式會在 \texttt{requests.get()} 那一行丟出一個沒被接住的例外（例如連線失敗
相關的錯誤），螢幕上出現一片 Traceback，程式直接中止，後面統計、備份的程式碼完全不
會被執行到。這跟 \ref{sec:move-files}~節忘記包 \texttt{try}/\texttt{except} 時的
狀況一樣：只要是「跟外部世界打交道、本來就可能失敗」的操作，沒有用 \texttt{try}/
\texttt{except} 接住，一失敗就會讓整支程式當場停擺，不會只是那一步失敗、其他部分
照跑。
\end{enumerate}

\paragraph{第~\ref{sec:download-multiple}~節}
\begin{enumerate}
\item \texttt{downloaded} 會是 \texttt{0}，\texttt{skipped} 會是 \texttt{3}，三個
網址都會各自被 \texttt{except} 接住、印出失敗訊息、\texttt{skipped} 累加，最後迴圈
正常跑完，不會中途中止。
\item 好處是 \texttt{download\_file()} 保持單純，只負責「下載這一個檔案」，不用管
「失敗了要怎麼辦」；決定失敗後要跳過、記錄還是重試，是呼叫端（迴圈）依情境決定的
策略，同一個 \texttt{download\_file()} 函式可以被不同的呼叫端用不同方式處理錯誤，
不用為了不同情境寫好幾個版本的下載函式。
\end{enumerate}

\paragraph{第~\ref{sec:advanced-integration}~節}
\begin{enumerate}
\item 會被存成 \texttt{vacation (1).jpg}，因為 \texttt{unique\_filename()} 發現
\texttt{organized\_demo/Image} 底下已經有同名的 \texttt{vacation.jpg}，改用加了編號
的新檔名，不會覆蓋掉原本已經存在的檔案。
\item 多用到 \texttt{load\_config()}（讓分類規則可以被 \texttt{config.json} 客製化，
對應 \ref{sec:config-file}~節）、\texttt{unique\_filename()}（避免同名檔案被悄悄
覆蓋，對應 \ref{sec:avoid-overwrite}~節）、\texttt{record\_error()}（讓單一檔案的
搬移失敗不拖累整批、還能留下紀錄，對應 \ref{sec:robust-errors}~節）。
\end{enumerate}

\chapter{軟體業詞彙與運作方式}
\label{ch:industry}

前面十八章，你一直在跟「怎麼寫出一支能動的程式」打交道。這一章要暫時放下語法，回頭問一個
不太一樣的問題：你寫的這些程式，在真正的軟體公司裡，是怎麼被一群人一起做出來、一起維護、
最後真的送到使用者手上的？這個問題不管你有沒有職場經驗都值得認識。

如果你還沒進過職場，這一
章帶你認識一個新世界，讓你知道「軟體工程師的日常」大概長什麼樣，之後不管是準備面試、選志向，
還是單純好奇，都有一個具體的畫面可以想。

如果你已經在工作，尤其是常常要跟工程師開會、卻聽不太
懂「這個 API 還在 review」「這張 ticket 排到下個 sprint」這類話的人（PM、營運、行銷都很常見
遇到這種情況），這一章的目標是讓你讀完之後，真的能拿這些詞彙去跟工程師同事對話，聽得懂、也
問得出對的問題。

整章分四段推進，順序是刻意設計過的：

\begin{itemize}
\item 先從你在第 \ref{ch:introduction} 章已經碰過的終端機出發，往下多學幾個指令，這是最小、最具體、你馬上就能
自己動手驗證的部分。
\item 接著把視角從「單一指令」拉高到「一個專案長什麼樣」：真正的軟體專案通常是一個資料夾，
裡面裝著一堆檔案、一份說明文件、一個獨立的執行環境，很少只有單一一個 \texttt{.py} 檔案，
這一段直接呼應你在第 \ref{ch:automating} 章自己動手做出來的整合應用專案。
\item 有了「專案」這個具體的東西之後，第三段才問一個更根本的問題：好幾個工程師同時改同一個
專案時，要怎麼避免互相蓋掉對方的成果，這就是版本控制存在的理由。
\item 最後一段拉到最高、最抽象的視角，把剛剛看到的「一個專案」放進一整間軟體公司的脈絡裡，
把你以後開會會聽到的詞彙一次講清楚，包括程式合併之後自動發生的事、軟體實際運作的樣子，以及
團隊安排工作進度時常用的說法。
\end{itemize}

先講清楚這一章\textbf{不}做什麼：接下來不會出現任何一行 Git 指令，也不會把終端機操作教到
精熟，這一章的目標是「聽得懂、問得出問題」，不是「會操作」。

真正動手把一個小專案從頭做到尾
的部分，留給下一章用一個實際的 API 串接小專案讓你完整體驗一次；這一章新增的幾個小節（例如
\texttt{.gitignore}、CI/CD、以及會議情境模擬）也遵守同一個原則，目標是聽懂、看懂，不是現在
就動手操作。

\section{你現在在哪裡、怎麼移動}
\label{sec:terminal-location}

回想第 \ref{ch:introduction} 章，你已經打開過終端機，也已經在裡面打過 \texttt{python3 hello.py} 這樣的指令，讓
Python 執行一支程式。但那時候你做的事情很單純：終端機已經幫你「站」在正確的資料夾裡了，你
只需要打對執行指令。這一節要補的，是終端機裡更基本、卻更常用的一件事：你怎麼知道自己現在
「人在哪個資料夾」，又怎麼「走」到別的資料夾去。

你可以把終端機想成一個只能用打字操控的資料夾總管，跟你平常用滑鼠點兩下打開資料夾的那個總管
做的是同一件事，只是這裡沒有圖示可以點，你得用指令告訴它該往哪裡走。以下三個指令，你接下來
每次打開終端機都會用到。

\textbf{第一步：確認自己在哪裡。} macOS／Linux 打 \texttt{pwd}（print working directory 的
縮寫），Windows 的命令提示字元打 \texttt{cd}（不加任何東西，直接按 Enter），PowerShell 兩種
都能用。畫面會印出一長串路徑，就是你現在所在資料夾的完整位置。
\index{pwd}

\textbf{第二步：切換到別的資料夾。} 用 \texttt{cd 資料夾名稱}。例如你在第 \ref{ch:introduction} 章新建的那個
\texttt{python-practice} 資料夾，如果它就在你目前的位置底下，打
\texttt{cd python-practice} 就能走進去。想往回走一層，打 \texttt{cd ..}，這兩個點是
「上一層資料夾」的意思，不管在哪個作業系統都通用。
\index{cd}

\textbf{第三步：看看這個資料夾裡有什麼。} macOS／Linux 打 \texttt{ls}（list 的縮寫），
Windows 命令提示字元打 \texttt{dir}，PowerShell 兩種都能用。畫面會列出目前資料夾底下的所有
檔案跟子資料夾。
\index{ls／dir}

把這三個指令連起來用，會像這樣（以下用 macOS／Linux 的指令示範，PowerShell 使用者可以直接
照打）：

\begin{verbatim}
$ pwd
/Users/xiaomei/python-practice
$ ls
hello.py
$ cd ..
$ pwd
/Users/xiaomei
$ cd python-practice
$ pwd
/Users/xiaomei/python-practice
\end{verbatim}

這段互動在做的事情，翻成白話就是：先確認「我現在在 \texttt{python-practice} 這個資料夾」，
接著看看裡面有什麼（喔，有一個 \texttt{hello.py}），然後往上走一層，再確認一次「我現在在
\texttt{xiaomei} 這個資料夾」，走回 \texttt{python-practice}，最後再確認一次「我又回來了」。
之後不管你的專案資料夾在電腦裡的哪個位置，這三個指令都能帶你走到你要去的地方，不需要離開
終端機、切去用滑鼠點資料夾總管再切回來。

\keypoint{終端機裡的「路徑」永遠是相對於你現在站的位置。\texttt{cd python-practice} 能不能
成功，取決於 \texttt{python-practice} 是不是真的在你現在所在的資料夾底下。如果不確定，先
\texttt{pwd}／\texttt{ls} 看清楚自己在哪裡、旁邊有什麼，比直接憑印象打指令可靠得多。}

再補兩個好用的捷徑，省得你每次都要一層一層 \texttt{cd ..} 走回去。第一個是「回家」：
macOS／Linux／PowerShell 打 \texttt{cd \textasciitilde}（波浪符號，通常在鍵盤數字鍵 1 左邊），
不管你目前站在哪裡，都會直接跳回你這個使用者帳號的「家目錄」，也就是你剛打開終端機時預設所在
的位置。

第二個是「回到剛剛那裡」：打 \texttt{cd -}（減號），會跳回你「上一次所在」的資料夾，
不用先想清楚完整路徑該怎麼寫，在兩個資料夾之間來回切換時特別方便。
\index{cd \textasciitilde／cd -}

\begin{drills}[label=drills:terminal-location]
\item 你在 \texttt{python-practice} 資料夾裡，打了 \texttt{cd ..} 走到上一層，這時候想直接
跳回剛剛那個 \texttt{python-practice}，不想再重新打一次完整路徑，該打哪個指令？
\item 你不確定自己現在在電腦裡的哪個很深的資料夾裡，想先跳回一個熟悉的起點重新開始，該打
哪個指令？
\end{drills}

\section{建立、複製與整理檔案}
\label{sec:cli-basics}

上一節你學會了「看」跟「走」，這一節補上「動手」的部分：怎麼在終端機裡建立、複製、搬移
檔案，以及兩個能讓你打指令的手速變快、變準的小技巧，最後補上一個所有終端機初學者都該提早
知道的風險提醒。

\textbf{建立資料夾：} \texttt{mkdir 資料夾名稱}（make directory 的縮寫），macOS／Linux／
Windows 都通用。建好之後用上一節學的 \texttt{cd} 走進去，就能開始在裡面工作。
\index{mkdir}

\textbf{複習：執行 \texttt{.py} 檔案。} 這件事你在第 \ref{ch:introduction} 章已經做過很多次了，確定自己在正確
的資料夾裡（用 \texttt{pwd}／\texttt{ls} 確認），打 \texttt{python3 檔名.py}（macOS）或
\texttt{python 檔名.py}（Windows）執行。這裡不重新教一次，只是提醒你：這一節後面的示範，都
建立在你已經很熟悉這個動作的前提上。

\textbf{複製檔案：} macOS／Linux／PowerShell 用 \texttt{cp 來源檔名 目的地}（copy 的縮寫），
Windows 命令提示字元用 \texttt{copy 來源檔名 目的地}。複製完成後，原本的檔案還在原地，
資料夾裡多了一份一模一樣的副本。舉例來說，\texttt{cp hello.py backup\_hello.py} 會在同一個
資料夾裡多一個叫 \texttt{backup\_hello.py} 的檔案，內容跟 \texttt{hello.py} 一模一樣。
\index{cp／mv／touch／cat}

\textbf{搬移或改名檔案：} macOS／Linux／PowerShell 用 \texttt{mv 來源檔名 目的地}（move 的
縮寫），Windows 命令提示字元用 \texttt{move}。跟 \texttt{cp} 不一樣的地方在於：\texttt{mv}
執行完之後，原本的檔案就不在原地了，只有目的地那一份存在。\texttt{mv} 還有一個容易被忽略的
用法：如果目的地只是「同一個資料夾裡的另一個檔名」，效果就等於「改名」，例如
\texttt{mv draft.py final.py} 會把 \texttt{draft.py} 直接改名成 \texttt{final.py}，資料夾
裡並不會同時出現兩個檔案。

\keypoint{\texttt{cp} 是複製，原本的檔案還在，多一份新的；\texttt{mv} 是搬移或改名，原本的
檔案不見了，只剩目的地那一份。兩個指令長得很像，功能卻完全不同，打錯的後果也不一樣：複製
打錯最多是多出一個垃圾檔案，搬移打錯可能會讓你以為某個檔案不見了，其實它只是被改名或搬到
別的地方，仍然可以用 \texttt{ls}／\texttt{cd} 找回來。}

\begin{code}{text}{複製與搬移檔案}
$ pwd
/Users/xiaomei/python-practice
$ mkdir archive
$ cp hello.py backup_hello.py
$ ls
archive  backup_hello.py  hello.py
$ mv backup_hello.py archive/backup_hello.py
$ ls
archive  hello.py
$ ls archive
backup_hello.py
\end{code}

這段記錄示範了兩件事：\texttt{cp} 之後資料夾裡多了一個 \texttt{backup\_hello.py}，原本的
\texttt{hello.py} 完全沒有受影響；\texttt{mv} 之後，\texttt{backup\_hello.py} 從目前資料夾
的列表裡消失了，改成出現在 \texttt{archive} 這個子資料夾底下，並不是又多存了一份到別的
地方，是同一個檔案換了位置。

\textbf{清空畫面：} 終端機用久了，畫面會被前面打過的指令跟輸出結果塞滿，不好閱讀。macOS／
Linux／PowerShell 打 \texttt{clear}，Windows 命令提示字元打 \texttt{cls}，會把畫面清乾淨，
只留下一行新的提示字元等你輸入，之前執行過的指令並沒有真的消失，只是「看起來」乾淨了。
\index{clear／cls}

\textbf{兩個讓手速變快的小技巧：} 終端機不是只能一個字一個字慢慢打出來。第一個技巧是自動
完成：打資料夾或檔案名稱時，打前面幾個字母，按一下 \texttt{Tab} 鍵，終端機通常會自動幫你把
剩下的部分補完（如果同時有好幾個檔案開頭一樣，可能要按兩下 \texttt{Tab} 才會列出所有候選），
可以少打很多字，也比較不會打錯字。

第二個技巧是叫出歷史指令：按鍵盤上的方向鍵「上」，終端機
會把你剛剛打過的上一行指令整行叫回來，不用重新打一次，這在你需要連續執行好幾次同樣（或很像）
的指令時特別省事。
\index{Tab 自動完成／方向鍵歷史指令}

\textbf{建立一個空檔案：} 有時候你只是想先佔一個檔名、之後再回頭填內容，不需要真的打開編輯器。
macOS／Linux／PowerShell 打 \texttt{touch 檔名}（原本的用途是「更新檔案的最後修改時間」，
但如果檔案不存在，順帶效果就是幫你建一個內容全空的新檔案），Windows 命令提示字元沒有直接對應
的指令，通常改用 \texttt{type nul > 檔名} 達到同樣效果。建好之後用你已經很熟的編輯器打開它，
把程式碼寫進去就好。

\textbf{不開編輯器、直接看檔案內容：} 有時候你只是想快速確認一個檔案裡寫了什麼，不需要為了
看一眼就整個打開編輯器。macOS／Linux 打 \texttt{cat 檔名}（concatenate 的縮寫），Windows
命令提示字元打 \texttt{type 檔名}，PowerShell 兩種都能用，畫面會把整個檔案的內容直接印在
終端機裡。這個指令適合快速確認一個小檔案（例如上一節提過的 \texttt{requirements.txt}）裡
寫了什麼，檔案內容一多，印出來反而會洗掉畫面，這時候還是打開編輯器看比較方便。

\textbf{刪除檔案的提醒：} macOS／Linux 用 \texttt{rm 檔名}，Windows 命令提示字元用
\texttt{del 檔名}，PowerShell 兩種都能用。這個指令本身不難，難的是它帶來的風險：\textbf{終端
機刪除的檔案不會進垃圾桶，也沒有「復原」這個選項，指令一按下去，檔案就真的消失了。} 你平常用
滑鼠把檔案拖進垃圾桶，就算手滑了還能救回來；終端機沒有這層保護，打錯資料夾、打錯檔名，後果是
直接、立即、不可逆的。
\index{rm／del}

下面示範一段完整的終端機互動記錄，把這一節跟上一節學到的指令串起來，延續第 \ref{ch:introduction} 章的
\texttt{python-practice} 資料夾：

\begin{code}{text}{在終端機裡建立、確認、清理}
$ pwd
/Users/xiaomei/python-practice
$ mkdir scratch
$ cd scratch
$ ls
$ python3 hello.py
python3: can't open file 'hello.py': [Errno 2] No such file or directory
$ cd ..
$ python3 hello.py
Hello, world!
$ rm hello.py
$ ls
\end{code}

這段記錄裡藏了一個很典型的小意外：在剛建立、還是空的 \texttt{scratch} 資料夾裡打
\texttt{python3 hello.py}，Python 老實地回報「找不到這個檔案」，因為 \texttt{hello.py}
根本不在這裡，它還留在上一層的 \texttt{python-practice}。這正是上一節〈學習要點〉提醒的事：
指令對不對，取決於你現在站在哪裡；\texttt{cd ..} 回到正確的位置之後，同一個指令就順利執行了。

最後 \texttt{rm hello.py} 把檔案刪掉、\texttt{ls} 確認資料夾已經空了，這一步在自己的練習
資料夾裡做沒什麼風險，但養成「刪之前先用 \texttt{pwd}／\texttt{ls} 確認自己在哪裡」的習慣，
會讓你在正式的專案資料夾裡也不容易失手。

到這裡，你已經知道怎麼用終端機在資料夾之間走動、確認位置、建立資料夾、複製與搬移檔案、
自動完成與叫回歷史指令，還有謹慎地刪除檔案。這幾個指令看起來瑣碎，但它們是接下來要做的每
一件事的地基，下一節要帶你用這些指令，走進一個「正常的軟體專案」長什麼樣。

為了方便之後查閱，這裡把這兩節用到的指令整理成一張對照表：

\begin{center}
\begin{tabular}{@{}lll@{}}
\toprule
作用 & macOS／Linux & Windows \\
\midrule
顯示目前所在路徑 & \texttt{pwd} & \texttt{cd}（不加參數）／PowerShell 可用 \texttt{pwd} \\
切換資料夾 & \texttt{cd 資料夾名稱} & 同左 \\
回到上一層資料夾 & \texttt{cd ..} & 同左 \\
列出資料夾內容 & \texttt{ls} & \texttt{dir}／PowerShell 可用 \texttt{ls} \\
建立新資料夾 & \texttt{mkdir 資料夾名稱} & 同左 \\
建立空檔案 & \texttt{touch 檔名} & \texttt{type nul > 檔名} \\
印出檔案內容 & \texttt{cat 檔名} & \texttt{type 檔名}／PowerShell 可用 \texttt{cat} \\
複製檔案 & \texttt{cp 來源 目的地} & \texttt{copy 來源 目的地} \\
搬移／改名檔案 & \texttt{mv 來源 目的地} & \texttt{move 來源 目的地} \\
清空畫面 & \texttt{clear} & \texttt{cls}／PowerShell 可用 \texttt{clear} \\
刪除檔案 & \texttt{rm 檔名}（不會進垃圾桶！） & \texttt{del 檔名}／PowerShell 可用 \texttt{rm} \\
\bottomrule
\end{tabular}
\end{center}

\begin{drills}[label=drills:cli-basics]
\item 你剛打開終端機，想先確認自己現在「人在哪個資料夾」，該打哪個指令？
\item 你用 \texttt{mkdir} 建了一個新資料夾，想確認裡面目前是不是真的空的，該打哪個指令？
\item 你想把 \texttt{report.py} 複製一份、取名叫 \texttt{report\_backup.py}，而且原本的
\texttt{report.py} 還要留著，該打哪個指令？如果改成用 \texttt{mv} 打同樣的指令，結果會有
什麼不一樣？
\end{drills}

\section{一個「正常專案」長什麼樣}
\label{sec:project-anatomy}

到目前為止，你寫過的程式大多是單獨一個 \texttt{.py} 檔案，寫完、存檔、執行，事情就結束了。
但你在第 \ref{ch:automating} 章做的那個檔案整理工具，已經稍微不一樣了：裡面有分類邏輯、有真正搬移檔案的部分、
還接了一段下載檔案的程式碼，而且你需要 \texttt{requests} 這個第三方套件才能讓下載的部分動
起來。你當時可能沒有特別意識到，但你已經在做業界說的「一個專案」了，只是規模比較小。

業界所謂的「一個軟體專案」，指的通常是一個資料夾，裡面裝著好幾個互相搭配的檔案，各自負責
不同的部分，很少只是單一一個檔案。以你第 \ref{ch:automating} 章那個檔案整理工具收尾後的樣子為例（實際檔名
依你當時的寫法可能略有出入，這裡示範的是典型的結構）：

\begin{verbatim}
file-organizer/
|-- organize-and-backup.py   main script: classify, move, back up, download
|-- requirements.txt         packages to install (requests)
|-- README.md                 documentation
|-- .gitignore                files excluded from version control (see Section 19.6)
|-- downloads_demo/           unsorted source files go here
`-- organized_demo/           sorted files, organized by category, go here
    |-- images/
    |-- documents/
    `-- music/
\end{verbatim}

這棵樹狀圖要傳達的重點是：\texttt{organize-and-backup.py} 不是孤零零一個檔案，它旁邊還有
其他東西一起組成「這個專案」，套件清單、說明文件、版本控制設定、甚至程式執行後產生的輸出
資料夾，全部放在同一個 \texttt{file-organizer/} 底下。如果你曾經打開終端機，用
\texttt{cd}／\texttt{ls} 這些指令在自己的專案資料夾裡走來走去、確認檔案都在，做的其實就是
業界工程師每天在做的事。

其中特別值得說一次的是 \texttt{README.md}：它不是任何程式語言規定要有的東西，Python 執行
你的程式時完全不會去讀它，單純是寫給「人」看的說明書，這個專案在做什麼、怎麼執行、需要先
裝什麼套件。
\index{README.md}

你現在可能覺得自己寫的東西自己當然知道在幹嘛，不需要特別寫，但三個月後你回頭
打開這個資料夾，或是換一台電腦想重新執行一次，你會感謝三個月前的自己留了這幾行字，更不用說，
如果這個專案以後要交給別人接手，一份清楚的 README 常常是對方唯一的入口。

規模再大一點的專案，資料夾裡常常還會多幾樣東西，這裡先點名兩個你以後很可能碰到、但這一章
不會深入的名字，讓你之後看到時不會完全陌生。

一個叫 \texttt{tests/} 的資料夾，裡面放的不是
給使用者用的程式，是工程師自己寫來「檢查程式有沒有寫對」的程式，一跑出錯就代表某個地方壞掉
了。

另一個常見的慣例是把主程式放進一個叫 \texttt{src/}（source 的縮寫）的子資料夾裡，跟
README、requirements.txt 這些「關於這個專案」的檔案分開放，資料夾一打開就能一眼分清楚
「這是程式本體」跟「這是給人看的說明」。這兩個名字現在只需要有印象，不需要現在就自己動手建。
\index{tests/／src/}

把這一節提到的所有東西放進同一個資料夾，一個規模再大一點、更接近業界實際樣子的專案，樹狀圖
大概會長這樣：

\begin{verbatim}
file-organizer/
|-- src/
|   `-- organize_and_backup.py   main program
|-- tests/
|   `-- test_classify.py         test script that checks the classification logic
|-- requirements.txt             packages to install
|-- README.md                     documentation
`-- .gitignore                    files excluded from version control
\end{verbatim}

跟前面那個比較陽春的版本比起來，這裡多了 \texttt{src/} 跟 \texttt{tests/} 兩層資料夾，程式
本體被收進 \texttt{src/} 底下，跟頂層那些「關於這個專案」的檔案分開；\texttt{tests/} 底下的
\texttt{test\_classify.py} 不是給使用者執行的程式，是工程師自己寫來檢查
\texttt{classify\_file()} 這種函式邏輯有沒有寫對用的，這一章不會教你怎麼寫測試程式，只需要
知道看到 \texttt{tests/} 資料夾時，裡面裝的是這一類東西。

兩種樹狀圖沒有哪一種「才是正確
答案」，規模小的專案用第一種也完全夠用，資料夾長什麼樣，取決於團隊的習慣跟專案的規模。

\begin{drills}[label=drills:project-anatomy]
\item 你打開一個同事傳來的專案資料夾，看到裡面有 \texttt{src/}、\texttt{tests/}、
\texttt{README.md} 三個東西，你會先打開哪一個，才能最快知道這個專案在做什麼、要怎麼執行？
\item \texttt{tests/} 資料夾裡的程式，是寫給誰用、檢查什麼東西的？
\end{drills}

\section{虛擬環境與套件管理}
\label{sec:project-structure}

你在第 \ref{ch:data-tools} 章第一次用 \texttt{pip} 裝過第三方套件，在第 \ref{ch:automating} 章可能又因為 \texttt{requests}
再裝了一次。這一節要解決一個你可能還沒碰過、但只要電腦上專案一多就一定會遇到的問題：如果你
同時在做兩個專案，一個需要某個套件的舊版本，另一個需要同一個套件的新版本，會發生什麼事？

想像兩個家庭合租一間房子、共用同一台冰箱：甲家庭把冰箱裡的鮮奶換成豆漿，是因為家裡有人不能
喝奶製品，結果乙家庭隔天想煮的那道菜，需要的正好是鮮奶，東西卻被換掉了。兩邊都沒有做錯
事，問題出在他們共用了同一個空間，而彼此的需求不一樣。

你電腦裡的 Python 套件，預設情況下
也是這樣共用的：如果 A 專案需要某個套件的 1.0 版、B 專案需要同一個套件的 2.0 版，直接
\texttt{pip install} 只會有一個版本裝在電腦上，兩個專案不可能同時滿意。

「虛擬環境」（virtual environment）就是在解決這個問題：它讓每個專案擁有一份自己專屬的套件
清單，彼此互相隔離，A 專案裝的版本不會影響到 B 專案，就像幫每個家庭各自準備一台冰箱。你不
需要現在就自己動手建一個虛擬環境，這一節的目標是讓你知道「這是在幹嘛、為什麼需要」，真正
動手操作留給你以後真的碰到版本衝突時再學。
\index{虛擬環境（virtual environment）}

跟虛擬環境經常一起出現的，是上一節樹狀圖裡看到的 \texttt{requirements.txt}，一份純文字
檔案，記錄「這個專案需要裝哪些套件、哪個版本」。
\index{requirements.txt}
以第 \ref{ch:automating} 章那個檔案整理工具為例，它的內容
可能長這樣：

\begin{code}{text}{(requirements.txt) 一個專案的套件清單}
requests==2.31.0
\end{code}

這份清單存在的意義，是讓別人（或是你自己換一台電腦）不用一個一個猜「這個專案到底需要裝什麼
套件、裝哪個版本才不會出錯」，直接照著這份清單裝，就能重建出一模一樣的環境，這也是為什麼
它幾乎是每個「正常專案」資料夾裡的標準配備，跟 README 一樣，是寫給未來的自己或別人看的。

反過來想，如果你已經用 \texttt{pip install} 陸續裝了好幾個套件，想幫這個專案生出一份
\texttt{requirements.txt}，業界工程師通常會用 \texttt{pip freeze} 這個指令，讓 \texttt{pip}
自己把目前環境裡裝了哪些套件、哪個版本整理成一份清單印出來，直接存成
\texttt{requirements.txt} 就好，很少手動一行一行慢慢打。這裡只需要知道有這個指令、它在幫你
做什麼，不需要現在就自己動手用。

\begin{drills}[label=drills:project-structure]
\item 你收到同事傳來的一個專案資料夾，裡面有一份 \texttt{requirements.txt}，內容寫著
\texttt{pandas==2.1.0}。這份檔案大概在告訴你什麼事？
\item 如果一個資料夾裡只有一支 \texttt{.py} 檔案，沒有 README、沒有
\texttt{requirements.txt}，這樣還算不算業界說的「一個專案」？（提示：想想 \ref{sec:project-anatomy} 節的定義，
「專案」看的是什麼，不是看檔案數量。）
\end{drills}

\section{版本控制在解決什麼問題}
\label{sec:version-control-why}

想像你跟同事一起維護一份共用的文件，不是程式碼，就是一份普通的說明文件或報告。你們兩個人
都可以編輯它，你今天下午改了一段，同事晚上也改了同一段，隔天早上你打開檔案，發現自己那段
改動整個不見了，取而代之的是同事的版本。

或者更糟：你上禮拜寫的內容，這禮拜想找回來看看，卻
發現檔案已經被改到面目全非，沒有人記得原本長什麼樣、是誰在什麼時候改的。這種經驗，只要跟
別人共同編輯過同一份東西，幾乎每個人都碰過。

軟體工程師面對的是同一個問題的加強版：一個專案可能同時有好幾個工程師在改同一批程式碼檔案，
而且程式碼一旦改壞，往往不像「這句話讀起來怪怪的」那麼容易被發現，很可能是整支程式直接
動不了。

他們解決這個問題的工具，叫做「版本控制系統」，目前業界最主流的實作是一套叫做 Git
的工具，但這一節要先建立的是「這一類工具在做什麼」的概念，不是 Git 的操作細節。
\index{版本控制系統（version control system）}

版本控制系統做的核心事情，是幫你的專案「拍照存證」：每一次你覺得「這是一個值得記住的
狀態」，就對整個專案的目前狀態拍一張快照，並且寫一句話說明「這次改了什麼、為什麼改」。這個
動作在版本控制的世界裡有個專門的名字，叫做 \textbf{commit}。
\index{commit}

跟你平常用 Ctrl+S 存檔不一樣的
地方在於：一般存檔每次都會蓋掉上一個版本，你回不去了；但每一次 commit 都會被完整保留下來，
你隨時可以回頭翻出三天前、三個禮拜前的某一次 commit，看看那個時間點的程式碼長什麼樣，甚至
把整個專案還原回那個狀態。

\keypoint{commit 不只是「存檔」的另一個說法，是「存檔加上一句話說明這次改了什麼、為什麼改」
的組合。少了那句說明，你三個月後回頭看一長串「更新」「修正」這種沒有內容的紀錄，會完全不
知道自己當初在改什麼。}

有了這套「拍照存證」的機制，前面提過的兩個困擾就都有解法了：檔案改壞了，能翻回某一次 commit
的狀態找回舊版本；兩人同時改同一份檔案，版本控制系統至少能告訴你「這兩處改動衝突了，需要
你自己決定要保留哪一邊」，不會像單純共用一份檔案那樣，其中一人的改動悄悄地、無聲無息地被
蓋掉。

這個「告訴你哪裡衝突了」具體長什麼樣，用一個小例子感受一下。假設 \texttt{organize.py} 裡
原本有這一行：

\begin{verbatim}
EXTENSION_MAP = {"jpg": "images", "pdf": "documents"}
\end{verbatim}

小美想多支援 \texttt{.png}，把它改成 \texttt{\{"jpg": "images", "png": "images", "pdf": "documents"\}}
並且 commit；同一時間，小華想多支援 \texttt{.docx}，也在自己的分支上把同一行改成
\texttt{\{"jpg": "images", "pdf": "documents", "docx": "documents"\}} 並且 commit。

等小華想把自己的分支
合併回主線（此時小美的改動已經先進去了）時，版本控制系統會發現「這一行在兩邊被改成不一樣的
內容」，沒辦法自己決定該用哪一種，於是停下來、把兩個版本都攤開給小華看，請小華自己決定最後
這一行要長什麼樣（通常答案是把兩邊都要的改動合在一起，變成同時支援 png 跟 docx）。這種
「兩邊都改了同一個地方、系統沒辦法自動判斷該用哪個」的狀況，就叫做\textbf{衝突（conflict）}。
\index{衝突（conflict）}

衝突本身不是壞事，甚至該說是版本控制系統在盡它的本分：與其悄悄選一邊、讓另一邊的改動憑空
消失，不如老實告訴你「這裡兩邊打架了，你自己來裁決」。

\section{不該被拍進快照裡的東西：.gitignore}
\label{sec:gitignore}

上一節的「拍照存證」比喻很好用，但也帶出一個新問題：如果每一次 commit 都要把整個專案目前的
狀態拍下來，是不是資料夾裡每一樣東西都該被拍進去？答案是不一定。有些東西刻意不該被記錄進
版本控制系統裡，業界工程師會準備一份叫 \texttt{.gitignore} 的檔案，明確列出「這些東西不用
拍、每次都跳過」。
\index{.gitignore}

最常見的三類「不該拍」的東西：

\begin{enumerate}
\item \ref{sec:project-structure} 節提過的虛擬環境資料夾。虛擬環境裡裝的套件檔案數量往往非常龐大，而且內容完全
可以靠 \texttt{requirements.txt} 重新裝出來，沒有必要把整個虛擬環境資料夾本身也塞進版本
控制系統裡，那樣做只會讓專案的紀錄變得又大又難管理。
\item Python 執行時自動產生的暫存檔案，最典型的是一個叫 \texttt{\_\_pycache\_\_} 的資料
夾。你目前的程式還不會產生這種東西，但等你的專案變大、開始用到 \texttt{import} 自己寫的
模組時，Python 會自動幫你的模組建立一份預先編譯好的版本，放在這個資料夾裡加快下次執行的
速度。這類檔案是 Python 自己重新產生就好的暫存品，不算你動手寫出來的原始碼，同樣沒有記錄的
必要。
\index{\_\_pycache\_\_}
\item 最重要、也最容易被新手忽略的一類：帳號密碼、API 金鑰這類機密資訊。如果你的程式裡有
一段像 \texttt{API\_KEY = "abcd1234"} 這樣直接寫死金鑰的程式碼，一旦這個專案被記錄進版本
控制系統，這串金鑰就會永久留在每一次 commit 的紀錄裡，就算你之後刪掉那一行重新 commit，
舊的紀錄裡還是找得到，這是業界公認的高風險失誤。之後如果這個專案要分享給別人或公開在網路上，
機密資訊一定要事先排除在外。
\end{enumerate}

一份典型的 \texttt{.gitignore} 內容大概長這樣：

\begin{code}{text}{(.gitignore) 排除不該被記錄的東西}
venv/
__pycache__/
*.pyc
.env
\end{code}

每一行代表一種要排除的檔案或資料夾名稱模式：\texttt{venv/} 排除整個虛擬環境資料夾（資料夾
名稱不一定叫 \texttt{venv}，你的專案實際用什麼名稱就寫什麼名稱）；\texttt{\_\_pycache\_\_/}
排除剛剛提到的暫存資料夾；\texttt{*.pyc} 排除所有以 \texttt{.pyc} 結尾的暫存檔案，星號代表
「不管檔名前面是什麼都算」；\texttt{.env} 則是業界慣用來存放機密設定值（像 API 金鑰、資料庫
密碼）的檔案名稱，同樣排除在外。

\keypoint{\texttt{.gitignore} 只是「請版本控制系統跳過這些東西」的名單，不是刪除，也不是
加密。真正保護機密資訊的做法，是一開始就不要把金鑰直接寫死在程式碼裡；\texttt{.gitignore}
只是最後一道防線，避免有人不小心把不該提交的檔案也一起拍進快照裡。}

\begin{drills}[label=drills:gitignore]
\item 以下四樣東西，哪些通常該放進 \texttt{.gitignore}：\texttt{organize-and-backup.py}、
\texttt{venv} 資料夾、\texttt{requirements.txt}、\texttt{\_\_pycache\_\_} 資料夾？
\item 如果一個專案完全沒有 \texttt{.gitignore}，會發生什麼具體的問題？
\end{drills}

\section{一群人一起改同一個專案}
\label{sec:version-control}

上一節解決的是「怎麼記錄改動、怎麼找回舊版本」，但還沒回答一個更根本的問題：如果團隊裡有五
個工程師，同時都在改同一個專案，難道大家都直接在同一份「正式版本」上動手嗎？如果有人改到
一半、程式碼還不能動，其他人豈不是也會被拖累？

答案是不會。版本控制系統讓每個人可以先在「自己的分支」上工作，互不干擾。

這個分支的概念叫做
\textbf{branch}：你可以把整個專案目前的狀態想成一條主線，這條主線任何時候都應該是「能正常
運作的版本」；每個工程師要動手改東西之前，先從主線上「岔出」一條屬於自己的分支，在那條分支
上盡情實驗、修改，就算改到一半、程式碼暫時是壞的，也完全不會影響主線，也不會影響其他人正在
做的事。用文字描述這個關係大概像這樣：
\index{branch}

\begin{verbatim}
main   --o---------------o---o-->
            \               /
mei      o---o---o---o
\end{verbatim}

小美從主線上的某個點岔出自己的分支，在上面做了好幾次 commit（圖上的每個 \texttt{o}），等
她確定這些改動都沒問題了，才把整條分支合併回主線。這個「我這邊改完了，想合併回主線，請大家
先看一下」的請求，有個專門的名字，叫做 \textbf{pull request}（簡稱 PR）。
\index{pull request（PR）}

但主線的守門員不會照單全收。在正式合併之前，通常會有其他工程師先看過這些改動，確認邏輯沒
問題、沒有明顯的錯誤、風格跟團隊的習慣一致，這個「別人先看過、給意見，才決定要不要合併」的
把關步驟，叫做 \textbf{code review}。等 review 通過，這條分支才會真正併回主線，變成主線的
一部分。
\index{code review}

\stylepoint{code review 不是在挑剔或不信任誰的能力，多一雙眼睛看過，本來就比較容易抓到
寫程式的人自己盲點裡看不到的問題（例如漏掉一種特殊情況、或跟專案其他地方的寫法衝突），這是
團隊維持程式碼品質的例行機制，幾乎每個正式的軟體團隊都會走這一道流程。}

這整套 branch、pull request、code review 的流程存在的理由很直接：團隊規模一旦大到超過
一兩個人，光靠口頭溝通「我改了這裡、你先別碰那裡」根本不可行，需要一套所有人都遵守、不用
互相口頭確認也能安心協作的流程，這正是版本控制系統要解決的問題。

再補一個常聽到、卻常被搞混的說法：把小美的分支併回主線這個動作，業界工程師平常講話會直接
說「merge」（合併），跟前面說的 pull request 是兩件事，但緊密相關。pull request 是「提出
合併請求、附上要 review 的內容」；merge 是「review 通過之後，改動真的被接進主線」那個動作
本身。你會聽到「這支 PR 我 merge 了」，說的就是這條分支的改動已經正式進到主線了。
\index{merge}

\begin{drills}[label=drills:version-control]
\item 小明把自己剛剛改好的程式碼推上去，並且提出「想請同事先看過再合併回主線」的請求，這個
動作叫什麼？
\item 小華還沒動手改任何東西之前，先從主線岔出一條屬於自己的分支開始工作，這麼做的目的是
什麼？
\item 「review 通過」跟「merge」是同一件事嗎？請用這一節的說法解釋兩者的差別。
\end{drills}

\section{軟體長什麼樣：前端、後端、部署}
\label{sec:software-anatomy}

前三段學到的「終端機」「專案」「版本控制」，都是工程師平常怎麼「做」軟體；這一節開始，要把
視角拉高，看看一個做出來的軟體，實際上是由哪些零件組成的，這是你以後開會最常聽到的一組
詞彙。

用一個大家都有的生活經驗切入：去餐廳吃飯。

\begin{itemize}
\item 你走進餐廳，跟外場服務生點餐、看菜單、拿到上桌的菜，這整個「你看得到、直接互動的
部分」，對應到軟體世界裡叫做\textbf{前端（front-end）}，就是使用者實際點得到、看得到的
畫面，例如一個網頁或手機 App 的介面。
\index{前端（front-end）}
\item 你看不到的內場廚房，負責真正把菜煮出來，對應到\textbf{後端（back-end）}，使用者看
不到，負責處理邏輯與資料的部分。
\index{後端（back-end）}
\item 廚房裡的食材儲藏室，對應到\textbf{資料庫（database）}，東西實際存放的地方。
\index{資料庫（database）}
\item 外場跟內場之間，靠一張點餐單傳遞訊息：外場不需要知道廚房怎麼炒菜，廚房也不需要知道
客人坐在哪一桌，兩邊只靠這張單子溝通，這正是\textbf{API（Application Programming
Interface）}在做的事：軟體跟軟體之間講話的接口，一方送出請求（「這桌要一份炒飯」），
另一方回覆結果。
\index{API（Application Programming Interface）}
\end{itemize}

這裡只需要
記住這個白話說明，程式碼層次的實作留到第 \ref{ch:api} 章實際串接一次 API 的時候再展開。

這個類比也能延伸到「一道新菜怎麼正式上菜單」：主廚先在後場試做一次，端給店長跟幾個內部員工
試吃，確認味道跟出餐流程都沒問題，這個「先在沒人看到的地方確認沒問題」的階段，對應到
\textbf{staging}，正式上線前，先給團隊內部測試確認沒問題的環境。
\index{staging}

等一切都確認過了，這道菜
才正式端給客人，對應到\textbf{production}，使用者實際在用的正式環境。把程式從工程師的
電腦、經過 staging 確認，一路送到 production 讓使用者真正用得到的整個動作，叫做
\textbf{部署（deploy）}。
\index{production}
\index{部署（deploy）}

這裡也順便帶一句你以後可能會聽到的說法：如果整套後端邏輯全部塞在同一支程式裡一起部署，
業界稱這種做法是「monolith」（單體式架構）；如果拆成好幾個各自獨立、各自部署的小服務，
彼此再透過 API 互相溝通，則稱為「微服務（microservice）」。這兩種做法各有適合的情境，這
一章不深入比較，你只需要知道以後聽到這兩個詞，指的是後端架構要不要拆成很多小塊這個選擇。
\index{monolith／微服務（microservice）}

這個「餐廳」的比喻只在這一節用，接下來不會反覆出現，它的任務只是幫你把這幾個詞彙在腦中
掛上一個具體的畫面，之後你在職場上聽到這些詞，直接對照白話意思即可，不需要再繞回餐廳。整理
成一張表方便你之後查閱：

\begin{center}
\begin{tabular}{@{}ll@{}}
\toprule
詞彙 & 白話意思 \\
\midrule
前端（front-end） & 使用者看得到、點得到的部分，例如網頁或 App 的畫面 \\
後端（back-end） & 使用者看不到，負責處理邏輯與資料的部分 \\
資料庫（database） & 東西實際存放的地方 \\
API & 軟體跟軟體之間講話的接口，一方送出請求、另一方回覆 \\
部署（deploy） & 把程式從工程師的電腦送到使用者用得到的地方 \\
staging & 正式上線前，先給團隊內部測試確認沒問題的環境 \\
production & 使用者實際在用的正式環境 \\
\bottomrule
\end{tabular}
\end{center}

\begin{drills}[label=drills:software-anatomy]
\item 一個網站的畫面看起來完全正常，按鈕點下去也有反應，但你填的資料每次重新整理都不見了，
比較可能是前端出問題，還是後端或資料庫出問題？
\item 「這個功能已經在 staging 上跑得很順，但 production 上還沒開放」，這句話代表使用者
現在看不看得到這個功能？為什麼？
\end{drills}

光是知道「deploy」這個動作的終點在哪裡還不夠，下一節要接著回答一個很實際的問題：合併之後，
這些步驟是每次都要有人手動做一遍，還是有辦法讓電腦自動幫你做？

\section{合併之後，自動發生的事：CI/CD}
\label{sec:cicd}

上一節提到，程式要經過 staging 確認，才會真正 deploy 到 production。如果每一次合併都要
工程師自己記得「該測試了、該部署了」，一忙起來很容易漏掉某個步驟，甚至忘記測試就直接把有
問題的程式碼送上 production，讓使用者遇到壞掉的功能。業界的解法，是讓電腦自動幫你做這些
重複、容易忘記的事，這一整套自動化機制有個統稱，叫做 \textbf{CI/CD}。
\index{CI/CD}

CI 是 Continuous Integration（持續整合）的縮寫：每一次有人把改動合併回主線，就自動觸發一
連串檢查，最常見的是自動把整個專案的測試程式碼跑一遍，確認這次改動沒有把原本能動的功能弄
壞。如果測試沒過，這套系統會立刻回報「這次合併有問題」，讓工程師在使用者發現之前就先發現。

CD 則有兩種常見的意思，都跟「後續動作自動化」有關：Continuous Delivery（持續交付）指的是
自動把通過測試的版本準備好、隨時可以部署；Continuous Deployment（持續部署）則更進一步，
測試一通過就自動直接 deploy 到 production，中間完全不需要人手動按下部署的按鈕。

你可以把 CI/CD 想成一條裝在合併動作後面的輸送帶：改動一旦被 merge 進主線，就自動被送上這條
輸送帶，依序經過測試、打包、部署到 staging、甚至部署到 production 這幾道關卡，每一道關卡都
是電腦自動執行，不需要工程師手動一步一步操作。這也是為什麼你可能會聽到工程師說「這次
pipeline 跑失敗了」，這裡的 pipeline（管線）指的就是這整條自動化輸送帶，「跑失敗」代表某一
道關卡（通常是測試那一關）沒有通過。
\index{pipeline（工作流程）}

一條 pipeline 實際跑完之後，留下的紀錄大概長這樣，工程師平常就是看這種畫面判斷「這次合併
到底有沒有問題」：

\begin{verbatim}
[pipeline #482] trigger: merge into main
[1/3] install packages.............. passed
[2/3] run tests...................... failed (3 tests did not pass)
[3/3] deploy to staging.............. skipped (previous stage failed)
result: pipeline failed, not deployed
\end{verbatim}

由上往下讀：前兩關寫著「passed」，第三關「deploy to staging」寫著「skipped」，不是它自己也失敗了，
是因為它的前一關（run tests）已經先失敗，系統判斷沒有必要再往下做，直接停在這裡。這正是這條
「輸送帶」的重點：只要有一道關卡沒過，後面的關卡就不會執行，壞掉的程式碼不會有機會一路溜到
production。工程師看到這種紀錄，會先點開「run tests」那一關的細節，找出是哪三項測試沒通過、
為什麼沒通過，修好之後再重新觸發一次 pipeline。

\keypoint{CI/CD 自動化的是「合併之後該做的重複性工作」，不是取代 code review。code review
仍然需要人親自看過改動的邏輯合不合理；CI/CD 負責的是那些可以寫成規則、機械式檢查的部分
（測試有沒有過、程式碼風格符不符合規範），兩者是互補的把關機制，不是二選一。}

\begin{drills}[label=drills:cicd]
\item 同事跟你說：「我剛剛 merge 的那次 pipeline 跑失敗了，我先去看一下。」請用白話翻譯這
句話在說什麼。
\item CI 跟 CD 各自處理的是「合併之後」流程裡的哪一段？
\end{drills}

\section{職場詞彙地圖：sprint、backlog、ticket}
\label{sec:workplace-vocab}

上一節看的是「軟體長什麼樣」，這一節把視角轉到「工程師團隊怎麼安排『這禮拜要做什麼』」，
這是你開會時最常聽到、也最容易一頭霧水的一組詞彙。

想像一個團隊共同維護一份待辦清單，上面記錄著所有「將來想做、但還沒排進行程」的事情：新
功能的點子、使用者回報的 bug、需要優化的地方。這份清單，業界叫它 \textbf{backlog}：已經
記錄下來、但還沒排進任何工作區間的待辦事項清單。
\index{backlog}

清單上每一件具體的事，會被寫成一張獨立的
工作卡片，叫做 \textbf{ticket}，上面寫著要做什麼、可能還有誰負責、目前進度到哪裡，可能是
一個新功能，也可能是要修的一個 bug。
\index{ticket}

團隊通常不會一次把 backlog 上的所有事做完，會每隔一段固定的時間（常見的是兩個星期），開會
決定接下來這段時間要從 backlog 裡挑哪些 ticket 出來做，這段約定要完成一批工作的固定時間
區間，叫做 \textbf{sprint}。sprint 結束時，團隊通常會檢視這段時間完成了什麼、還有什麼沒做
完，沒做完的就退回 backlog，或排進下一個 sprint。
\index{sprint}

同一組詞彙，PM 跟工程師常用的語境會有點不一樣，整理成一張對照表，讓你以後開會聽到時能立刻
對上：

\begin{center}
\begin{tabular}{@{}lp{4cm}p{4cm}@{}}
\toprule
詞彙 & PM／營運常見講法 & 工程師常見講法 \\
\midrule
backlog & 「這個先丟進 backlog，之後再排」 & 「這張票在 backlog 躺很久了」 \\
sprint & 「這個 sprint 要交付什麼」 & 「這個 sprint 排不進去，下一輪再看」 \\
ticket & 「我先開一張 ticket 追蹤這件事」 & 「這張 ticket 指派給我了，我這週處理」 \\
\bottomrule
\end{tabular}
\end{center}

\begin{drills}[label=drills:workplace-vocab]
\item 同事在會議上說：「這張 ticket 我先丟回 backlog，這個 sprint 排不進去了。」這句話裡的
ticket、backlog、sprint 分別在指什麼？
\item 「這張 ticket 還躺在 backlog 裡」跟「這張 ticket 已經排進這個 sprint 了」，這兩句話
在講的是同一張工作卡片的什麼差異？
\end{drills}

\section{情境模擬：聽懂會議裡的那句話}
\label{sec:workplace-scenarios}

這一章介紹了不少詞彙，這一節不再介紹新東西，是把它們放回真實的對話情境裡，看看你能不能
聽懂。

\textbf{情境一：「這支 PR 還在 review。」} 你可能在會議上聽過工程師這樣回你「這個功能做好
了嗎」的問題。翻成白話，這句話代表程式碼已經寫完、已經提出 pull request，正在等別的工程師
做 code review，還沒有人看完、確認沒問題，所以還沒 merge 回主線，當然也還沒開始跑後面
CI/CD 那條自動化輸送帶，更不用說部署上線了。

如果你想追問「大概還要多久」，比較有意義的問法
通常是「已經有人開始看了嗎」或「大概還要等誰」，直接問「什麼時候寫完」常常問不出重點，因為
程式碼可能早就寫完了，卡住的是「還沒有人看過」這一關。

\textbf{情境二：「pipeline 跑失敗了，我在查。」} 如果工程師這樣回你，代表 \ref{sec:cicd} 節說的
CI/CD 自動化流程在測試（或其他自動檢查）那一關卡住了，工程師正在排查是哪裡出了問題，不一
定跟你關心的那個功能有關（也可能是別人改動的另一段程式碼恰好把測試弄壞了）。

這種情況下追問
「是不是我要的這個功能有 bug」通常問不出什麼，比較好的問法是「這次卡關會不會影響到我們原本
排定的時間」。

\textbf{情境三：PM 問「這個功能什麼時候能上線」，工程師回答「程式碼都寫完了，正在 staging
上測試，確認沒問題就會 deploy」。} 翻成白話：程式碼已經寫完、也已經通過 code review 合併
進主線了，目前正放在 staging（團隊內部測試用的環境）上做最後確認，一旦確認沒問題，就會
deploy（部署），把程式正式送到使用者用得到的 production（正式環境）上線。

這句話裡沒有一個
字提到「還沒寫完」，卻常常被誤會成「還在寫」，聽懂 staging 跟 production 的差別，才不會
反覆追問一個其實已經有明確答案的問題。

\textbf{情境四：「這支 API 前端已經串好了，後端那邊資料庫欄位還沒改完。」} 對照 \ref{sec:software-anatomy} 節的
餐廳比喻：API 對應點餐單，是前端跟後端之間傳遞請求的窗口；前端對應外場，是使用者看得到、
互動的部分；後端對應內場廚房，處理邏輯的部分；資料庫對應食材儲藏室，是東西實際存放的地方。

這句話合起來的意思是：畫面（前端）那一側已經準備好可以送出請求了，但廚房（後端）那一側要用
的食材儲藏室（資料庫）的欄位結構還沒調整完，所以這個功能整體還不能真正動起來，卡的是後端跟
資料庫，不是前端。

走到這裡，整章其實可以串成一幅完整的畫面：有人先開了一張 \textbf{ticket}，描述要做的事；
這張 ticket 被排進某個 \textbf{sprint}；負責的工程師從主線岔出一條 \textbf{branch}，動手
改一個「正常專案」資料夾裡的檔案，過程中資料夾裡有些東西（像虛擬環境、暫存檔）刻意被
\texttt{.gitignore} 排除在版本控制之外；

改完之後發出 \textbf{pull request}，讓同事先
\textbf{code review} 過；review 通過、merge 回主線之後，交給 \textbf{CI/CD} 這條自動化
輸送帶跑測試、跑部署，經過 \textbf{staging} 的確認，最後 \textbf{deploy} 到
\textbf{production}，使用者才真正用得到這個新功能或修好的 bug。

你在職場上會聽到的每一句
行話，說的都是這整幅畫面裡的某一小段，這一章的目標，就是讓你下次聽到時，能立刻在腦中對上
是哪一段。

\begin{drills}[label=drills:workplace-scenarios]
\item PM 問工程師：「這個功能什麼時候能上線？」工程師回答：「程式碼都寫完了，正在 staging
上測試，確認沒問題就會 deploy。」請用白話翻譯工程師這句回答在說什麼。
\item 如果有人跟你說「這支 API 前端已經串好了，但後端那邊資料庫欄位還沒改完」，這句話裡
API、前端、後端、資料庫，分別對應到 \ref{sec:software-anatomy} 節餐廳比喻裡的哪個角色？
\item 同事說「這支 PR 還在 review」，跟同事說「pipeline 跑失敗了」，這兩句話卡住的分別是
整個流程裡的哪一關（提示：一個發生在 merge 之前，一個發生在 merge 之後）？
\end{drills}

\section*{練習解答}
\addcontentsline{toc}{section}{練習解答}

\paragraph{第~\ref{sec:terminal-location}~節}
\begin{enumerate}
\item \texttt{cd -}，會跳回你上一次所在的那個資料夾，不用重新打完整路徑。
\item \texttt{cd \textasciitilde}，不管目前站在哪個資料夾，都會直接跳回自己的家目錄。
\end{enumerate}

\paragraph{第~\ref{sec:cli-basics}~節}
\begin{enumerate}
\item \texttt{pwd}（macOS／Linux）或不加參數的 \texttt{cd}（Windows 命令提示字元，
PowerShell 也可以打 \texttt{pwd}）。
\item \texttt{ls}（macOS／Linux）或 \texttt{dir}（Windows）。
\item \texttt{cp report.py report\_backup.py}。改用 \texttt{mv} 打同樣的指令
（\texttt{mv report.py report\_backup.py}）的話，資料夾裡不會同時留下兩個檔案，
\texttt{report.py} 會直接消失，只剩下改名後的 \texttt{report\_backup.py}。
\end{enumerate}

\paragraph{第~\ref{sec:project-anatomy}~節}
\begin{enumerate}
\item \texttt{README.md}，它是寫給人看的說明文件，用來快速知道這個專案在做什麼、怎麼執行；
\texttt{src/} 跟 \texttt{tests/} 裡放的是程式碼本身，不看過 README 直接讀程式碼，通常比較
難快速抓到整體脈絡。
\item 寫給工程師自己用的，用來檢查專案裡的某段邏輯（例如分類函式）有沒有寫對，不是給使用者
執行的程式。
\end{enumerate}

\paragraph{第~\ref{sec:project-structure}~節}
\begin{enumerate}
\item 這個專案需要裝 \texttt{pandas} 這個套件，版本要是 2.1.0。別人（或你自己在另一台電腦
上）想重建這個專案的執行環境時，照著這份清單裝，就能裝出一模一樣的套件版本。
\item 算。\ref{sec:project-anatomy} 節說的「專案」，看的是有沒有一份說明文件、有沒有記錄清楚需要什麼套件這種
完整程度，不是資料夾裡塞了幾支 \texttt{.py} 檔案。只有一支檔案、什麼說明都沒有，仍然可以是
一個專案，只是還沒有具備「正常專案」該有的說明文件跟套件清單，比較不方便別人（或未來的自己）
接手。
\end{enumerate}

\paragraph{第~\ref{sec:gitignore}~節}
\begin{enumerate}
\item \texttt{venv} 資料夾跟 \texttt{\_\_pycache\_\_} 資料夾通常該放進
\texttt{.gitignore}；\texttt{organize-and-backup.py} 是你自己寫的原始碼、
\texttt{requirements.txt} 是給別人重建環境用的清單，這兩個都應該被版本控制記錄下來，不該
排除。
\item 版本控制系統會把 \texttt{venv}、\texttt{\_\_pycache\_\_} 這類龐大又可以重新產生的
東西也整個記錄進每一次 commit，讓專案的紀錄變得又大又難管理；如果程式碼裡剛好寫死了帳號
密碼或 API 金鑰，少了 \texttt{.gitignore} 排除機密設定檔，這些機密資訊也會被永久留在版本
紀錄裡，就算之後刪掉也還找得到。
\end{enumerate}

\paragraph{第~\ref{sec:version-control}~節}
\begin{enumerate}
\item pull request（簡稱 PR）。
\item 讓自己可以放心實驗、修改，就算改到一半程式碼暫時是壞的，也不會影響主線、也不會影響
其他人正在進行的工作。
\item 不是同一件事。review 通過只代表「有人看過這些改動、覺得沒問題」，merge 才是「改動真的
被接進主線」那個動作本身，通常是 review 通過之後才會發生 merge，但兩者是流程裡先後不同的
兩個步驟。
\end{enumerate}

\paragraph{第~\ref{sec:software-anatomy}~節}
\begin{enumerate}
\item 比較可能是後端或資料庫出問題。前端負責畫面顯示跟互動，看起來正常代表前端這一側大致
沒問題；資料填了卻存不住，通常是後端處理資料的邏輯，或資料庫實際存放資料的那一步出了狀況。
\item 看不到。staging 是團隊內部測試用的環境，不是使用者看得到的地方；使用者實際用的是
production，這句話代表這個功能還沒 deploy 到 production，所以還沒有任何使用者看得到它。
\end{enumerate}

\paragraph{第~\ref{sec:cicd}~節}
\begin{enumerate}
\item 這位同事剛剛把自己的改動 merge 進主線，但 CI/CD 自動跑的檢查（通常是測試）沒有通過，
他正在查是哪裡出了問題，這件事目前還沒有解決。
\item CI（持續整合）處理的是「merge 之後、自動確認這次改動有沒有把原本能動的功能弄壞」這一
段，主要是跑測試；CD（持續交付／持續部署）處理的是「測試通過之後，怎麼把改動送到 staging
甚至 production」這一段。
\end{enumerate}

\paragraph{第~\ref{sec:workplace-vocab}~節}
\begin{enumerate}
\item ticket 是那件具體要做的工作（可能是新功能或 bug）；backlog 是「已經記錄下來、但還沒
排進任何 sprint」的待辦清單，這張 ticket 現在被放回這份清單裡；sprint 是團隊約定要完成一批
工作的固定時間區間，這張 ticket 沒能排進目前這一輪。
\item 這是同一張 ticket 在時間安排上的兩種狀態：「躺在 backlog 裡」代表這件事已經被記錄
下來，但還沒被排進任何一個具體的時間區間去做；「排進這個 sprint」代表團隊已經決定接下來這
段固定時間要處理這件事了。
\end{enumerate}

\paragraph{第~\ref{sec:workplace-scenarios}~節}
\begin{enumerate}
\item 程式碼已經寫完了，目前正放在 staging（團隊內部測試用的環境）上確認沒有問題；一旦
確認沒問題，就會 deploy（部署），把程式正式送到使用者用得到的 production（正式環境）上線。
\item API 對應點餐單（前端跟後端之間傳遞請求的窗口）；前端對應外場（使用者看得到、互動的
部分）；後端對應內場廚房（處理邏輯的部分）；資料庫對應食材儲藏室（東西實際存放的地方）。
\item 「這支 PR 還在 review」卡在 merge 之前，等的是有人先看過這些改動、確認沒問題；
「pipeline 跑失敗了」發生在 merge 之後，是 merge 進主線之後 CI/CD 自動跑的檢查沒有通過。
\end{enumerate}

\chapter{API 串接小專案}
\label{ch:api}

到目前為止，這本書裡「資料從哪裡來」這件事，答案幾乎都是同一種：你自己在程式裡打出來
的，或是存在電腦裡的一個檔案，第 \ref{ch:csv} 章讀的 CSV 也是這樣，那是一份早就寫好、存在硬碟
上的靜態檔案，你打開它、讀它，它不會自己變。

軟體要用到的資料，很多時候是這一秒才剛剛
發生、存在別人家伺服器上的東西，例如現在台北的氣溫幾度、現在一美元換多少日圓，不會乖乖
躺在你電腦裡等你讀取。

這一章要讓你動手做一件新事：讓你自己寫的程式，直接開口跟另一支
（跑在世界某個角落的）程式要資料，即時拿到答案。第 \ref{ch:industry} 章你已經認識了「API」這個詞，
知道它大致上是「軟體跟軟體溝通的窗口」；這一章要讓這個詞第一次真正被你自己的手打過
一次。

這一章分成六段推進：

\begin{itemize}
\item 第一段讓你實際打一次公開的 API，看到最原始的回應長什麼樣子，建立
「軟體真的可以直接跟另一支軟體要資料」的具體體感，順便把網址後面那一長串「問號加參數」
的寫法講清楚。
\item 第二段把那包看起來雜亂的回應，拆解成你早就很熟的字典跟串列，變成真正能
拿來用的資訊。
\item 第三段刻意換一個完全不同的服務再打一次，這次查的是即時匯率，用來確認你
剛剛學的技巧不是天氣查詢專屬的招式，換一個 API 一樣管用，也順便提一下有些 API 不像這
兩個，必須先申請一把專屬的金鑰才能查詢，那把金鑰又該怎麼收好。
\item 第四段把焦點轉到「這中間
會出什麼錯」，網路斷線、伺服器忙線、網址打錯，這些狀況跟你在第 \ref{ch:files-exceptions} 章學過的檔案讀取一樣，
都是「本來就可能失敗的操作」，正好是 \texttt{try}/\texttt{except} 的主場。
\item 第五段把抓回來
的資料存成 CSV 檔、畫成圖表，跟第 \ref{ch:csv}、\ref{ch:data-tools} 章學過的工具接起來。
\item 第六段把前面學到的東西全部
串成一支小工具，讓你輸入城市或幣別，就能查到即時的天氣或匯率，親手完整跑一次「一支程式
去跟另一支程式要東西」的流程。
\end{itemize}

跟第 \ref{ch:automating} 章的整合應用專案不一樣，\ref{ch:automating} 章的任務是把前面十七章
學過的程式技巧全部兜在一起，服務的是這本書「自己動手寫程式解決瑣事」這個目標；這一章
刻意窄很多，只聚焦在一件事，具體示範軟體之間怎麼溝通，服務的是「看懂軟體世界怎麼運作」
這個目標。

\section{第一次打 API：看到「軟體跟軟體對話」長什麼樣子}
\label{sec:first-api-call}

要跟另一支程式要資料，第一步得先知道去哪裡要。跟你在瀏覽器打網址很像，軟體之間溝通
用的也是一個網址，只是這個網址背後跑的是一支專門用來回答「查詢」的程式，回來的東西
是一包給程式讀的資料。這一節要用的是 Open-Meteo 這個公開的天氣查詢服務
（\texttt{api.open-meteo.com}），它有一個很適合初學者的優點：不用註冊帳號、不用申請
金鑰，只要帶上經緯度，就能直接查到那個地點的天氣。

要在 Python 裡發出這種請求，得先借助別人寫好的工具，也就是 \texttt{requests} 這個套件。
還記得第 \ref{ch:data-tools} 章提過，要用別人寫好的套件，得先用 \texttt{pip install} 把它裝進你的環境嗎？
\texttt{requests} 就是一個很典型的例子：在終端機打 \texttt{pip install requests}
（有些系統要打 \texttt{pip3 install requests}），裝好之後，就能在程式裡
\texttt{import requests} 直接使用了。

\begin{code}{python}{Making a GET Request}
import requests

response = requests.get(url)   # url 是你要查詢的網址字串
print(response.status_code)    # 伺服器的狀態碼，200 代表成功
print(response.text)           # 伺服器實際回傳的內容，是一段文字
\end{code}

\texttt{requests.get(url)} 會真的連到那個網址，把伺服器回應的東西整包帶回來，貼上一張
名字牌叫 \texttt{response}，跟你在第 \ref{ch:variables} 章學過的一樣，這只是一張貼在「這次請求結果」這個
物件上的名字牌而已。這個物件裡有兩個你現在用得到的欄位：\texttt{status\_code} 是一個
數字，用來告訴你這次請求順不順利（\texttt{200} 是最常見的「一切正常」）；\texttt{text}
則是伺服器實際回傳的內容，本質上就是一段文字。
\index{response 物件}

\begin{example}{python}{(first-weather-request.py) Checking Taipei's Weather}
import requests

url = (
    "https://api.open-meteo.com/v1/forecast"
    "?latitude=25.0330&longitude=121.5654"
    "&current_weather=true"
)
response = requests.get(url)

print(response.status_code)
print(response.text)
\end{example}

\iosample{}{200 \\}

狀態碼印出 \texttt{200}，代表請求成功。但如果你接著把 \texttt{response.text} 印出來，
會看到這樣一長串、擠成一坨的文字（下面這一整段其實是同一行，紙本排版沒辦法印出一行
一百多個字，這裡手動斷開，實際在你的螢幕上會看到它自動換行，但邏輯上仍然是同一行；
其中 \texttt{generationtime\_ms} 是伺服器處理這次請求花的時間，每次呼叫都不一樣，這裡
只保留三位小數方便閱讀）：

\begin{verbatim}
{"latitude":25.06151,"longitude":121.5194,
"generationtime_ms":0.137,"utc_offset_seconds":0,
"timezone":"GMT","timezone_abbreviation":"GMT",
"elevation":12.0,"current_weather_units":{
"time":"iso8601","interval":"seconds",
"temperature":"°C","windspeed":"km/h",
"winddirection":"°","is_day":"","weathercode":
"wmo code"},"current_weather":{
"time":"2026-09-01T15:00","interval":900,
"temperature":25.6,"windspeed":1.0,
"winddirection":239,"is_day":0,"weathercode":80}}
\end{verbatim}

有兩件事值得注意。第一，這包東西看起來很像你已經學過的字典跟串列：大括號、方括號、
冒號、逗號，只是全部擠在一起、沒有任何排版，這不是巧合，下一節會正式解釋為什麼。第二，
仔細看你會發現回應裡的經緯度（\texttt{25.06151}、\texttt{121.5194}）跟你查詢時帶的
（\texttt{25.0330}、\texttt{121.5654}）不完全一樣，原因是 Open-Meteo 的天氣資料來自
固定間距的氣象網格，它會自動幫你抓最接近的那個網格點，不是查詢地點本身沒有天氣資料。

\begin{example}[label=ex:browser-check]{text}{(Paste into a Browser) The Same URL}
https://api.open-meteo.com/v1/forecast?latitude=25.0330
&longitude=121.5654&current_weather=true
\end{example}

這其實是同一個網址，只是紙本排版沒辦法印出一行一百多個字，這裡在
\texttt{latitude} 跟 \texttt{longitude} 之間手動斷開，真的要貼進瀏覽器時，記得把它
接回同一行、中間不要留空格或換行。把這個網址整段複製、貼到瀏覽器網址列，按下 Enter，
你會看到瀏覽器畫面上出現的內容，
跟剛剛 \texttt{print(response.text)} 印出來的其實是同一包東西（多數瀏覽器還會貼心地
幫你加上顏色、縮排，方便閱讀，但骨子裡的資料一模一樣）。

這件事值得停下來體會一下：所謂
「打 API」，說穿了就是拿程式去發一個網址請求，跟你手動在瀏覽器貼網址是同一件事，只是
瀏覽器把結果畫成一個網頁給\emph{人}看，你的程式則是把結果讀進變數裡，給\emph{另一支
程式}接著處理。

\keypoint{API 本質上就是一個網址，差別只在於它回來的資料是設計給程式讀的，不是設計給人
看的排版網頁。你在瀏覽器貼過的每一個普通網址，跟你剛剛用 \texttt{requests.get()} 打的
這個網址，走的其實是同一套機制。}

眼尖的話，你可能已經注意到剛剛那個網址裡，問號後面接了一串用 \& 連起來的
\texttt{latitude=...}、\texttt{longitude=...}、\texttt{current\_weather=...}，這種寫法
有一個正式名字，叫做「查詢字串」（query string）。下一節要把這件事講清楚，順便介紹一個
比手動拼字串更保險的寫法。

\section{網址裡的問號和 \&：query string 參數怎麼運作}
\label{sec:query-strings}

一個網址其實可以拆成兩段。問號前面那一段，是你要查詢哪一支程式，例如
\texttt{https://api.open-meteo.com/v1/forecast} 就是「Open-Meteo 這個服務裡負責回答
天氣預報的那一支程式」。問號後面那一段，叫做查詢字串，是你要把哪些條件交給那支程式：
每一組條件寫成「鍵=值」的形式，好幾組條件之間用 \texttt{\&} 串起來。
\index{查詢字串 (query string)}

這個形狀，其實你
一點都不陌生：查詢字串本質上就是把一個字典壓平、塞進網址裡的寫法，鍵就是字典的 key、
值就是字典的 value，只是字典原本用大括號、逗號分隔，查詢字串改用問號、\texttt{\&}、
等號來分隔而已。

前一節的範例是自己手動把這些「鍵=值」拼成一個長字串，遇到參數不多的時候還好，一旦
參數變多，手動拼字串很容易漏加一個 \texttt{\&}、或是打錯一個等號的位置，程式碼也會
變得很難讀。\texttt{requests} 提供了一個更好的做法：把所有查詢條件包成一個 Python
字典，交給 \texttt{requests.get()} 的 \texttt{params} 這個關鍵字引數，剩下的組字串
工作交給套件自己處理。
\index{params 引數}

\begin{code}{python}{Passing Query Parameters with a params Dictionary}
params = {
    "key1": value1,
    "key2": value2,
}
response = requests.get(url, params=params)
\end{code}

\begin{example}{python}{(query-string-params.py) Querying Weather with a Parameter Dictionary}
import requests

url = "https://api.open-meteo.com/v1/forecast"
params = {
    "latitude": 25.0330,
    "longitude": 121.5654,
    "current_weather": "true",
    "timezone": "Asia/Taipei",
}
response = requests.get(url, params=params, timeout=5)

print(response.url)
data = response.json()
temperature = data["current_weather"]["temperature"]
print(f"Current temperature: {temperature}°C")
\end{example}

執行後 \texttt{response.url} 印出的內容如下（這是 \texttt{requests} 幫你組好、真正送出去
的完整網址，跟你自己寫程式裡的 \texttt{params} 字典比對一下，會發現每一組鍵值都乖乖
出現在網址裡，字典裡的順序也保留下來了）：
\index{response.url}

\begin{verbatim}
https://api.open-meteo.com/v1/forecast?latitude=25.033&longitude=121.5654&current_weather=true&timezone=Asia%2FTaipei
\end{verbatim}

\iosample{}{Current temperature: 25.4°C \\}

值得多看一眼的是 \texttt{timezone} 這個值：你在 Python 字典裡寫的是
\texttt{"Asia/Taipei"}，中間有一個斜線，但網址裡印出來的是 \texttt{Asia\%2FTaipei}。

斜線在網址裡有特殊意義（用來分隔路徑），如果直接把它原封不動塞進查詢字串，某些伺服器
或某些工具可能會誤判網址的結構，所以正式的做法是把這類特殊符號轉換成「百分號加兩位
十六進位數字」的形式，這叫做 URL 編碼，\texttt{\%2F} 就是斜線編碼過後的樣子。空白字元、
中文字，甚至逗號，都有各自對應的編碼結果。
\index{URL 編碼}

這件事你不需要自己動手轉換，這正是把參數包成
\texttt{params} 字典交給 \texttt{requests} 處理的好處：只要值來自使用者輸入、有可能包含
空白或特殊符號，交給 \texttt{params} 遠比自己手動拼字串可靠，該轉換的它會自動幫你轉換好。

\stylepoint{查詢參數固定不變、只有一兩個的時候，手動拼字串（像本章第一節那樣）還算
看得懂；但只要參數有機會來自使用者輸入、或是參數一多，優先用 \texttt{params} 字典交給
\texttt{requests} 處理，比自己動手拼字串更不容易出錯，也更容易在程式裡新增或刪掉某一個
查詢條件。}

\begin{drills}[label=drills:query-strings]
\item 如果 \texttt{params = \{"city": "New York"\}}，你猜 \texttt{requests} 組出來的
網址裡，\texttt{city} 這個值會被轉換成什麼樣子（提示：空白字元在查詢字串裡通常會被轉成
加號或 \texttt{\%20}）？
\item 承上題，如果你自己手動把網址寫成 \texttt{...?city=New York}（中間留一個沒有轉換
過的空白），實際貼到瀏覽器網址列時通常還是能用，因為多數瀏覽器會自動幫你補做轉換。但
在程式裡呢？為什麼書裡建議優先用 \texttt{params} 字典，而不是依賴這種「剛好可以動」的
手動寫法？
\end{drills}

\section{把回應拆開：JSON 其實就是你已經很熟的字典和串列}
\label{sec:parsing-json}

上一節印出來的那一整坨文字，格式其實有名字，叫做 JSON（發音類似「傑森」），這是網路
服務之間交換資料時最常用的格式，而它長的樣子，剛好就是 Python 字典跟串列的樣子：大括號
包起來的是字典，方括號包起來的是串列，裡面可以一層包一層。這代表你不需要學一套全新的
語法來讀懂 API 回來的東西，只要把這包文字轉成真正的 Python 字典跟串列，接下來就是你
已經很熟的操作。

\begin{code}{python}{Converting a Response to Python Data and Extracting Fields}
data = response.json()             # 把 JSON 文字轉成字典或串列
value = data["key"]                # 確定這個欄位一定存在時使用
value = data.get("key", default_value) # 欄位可能不存在時，給一個預設值
\end{code}

\texttt{response.json()} 會把 \texttt{response.text} 那包文字，直接轉成對應的 Python
字典（或串列，要看最外層是 \texttt{\{\}} 還是 \texttt{[]}）。轉完之後，接下來要拿哪個
欄位，就跟你平常操作字典完全一樣，用中括號、或用 \texttt{.get()}。
\index{response.json()}

\begin{example}{python}{(current-weather-fields.py) Extracting Fields from Weather Data}
import requests

url = (
    "https://api.open-meteo.com/v1/forecast"
    "?latitude=25.0330&longitude=121.5654"
    "&current_weather=true"
)
response = requests.get(url)
data = response.json()

temperature = data["current_weather"]["temperature"]
print(f"Current temperature: {temperature}°C")

wind_speed = data.get("current_weather", {}).get("windspeed", "no data")
print(f"Current wind speed: {wind_speed} km/h")

# 刻意查一個這包資料裡沒有的欄位，用來示範 .get() 的預設值
sunshine = data.get("current_weather", {}).get("sunshine_hours", "this field is not provided")
print(f"Sunshine duration: {sunshine}")
\end{example}

\iosample{}{Current temperature: 25.6°C \\ Current wind speed: 1.0 km/h \\ Sunshine duration: this field is not provided \\}

\texttt{data["current\_weather"]["temperature"]} 這種一路用中括號往下挖的寫法，前提是
你很確定每一層都真的存在，如果查的欄位剛好不存在，中括號會直接讓程式當掉。假設你把上面
範例最後一段改成直接用中括號查 \texttt{sunshine\_hours}（而不是用 \texttt{.get()}）：

\begin{verbatim}
data = {
    "current_weather": {"temperature": 25.6, "windspeed": 1.0}
}
print(data["current_weather"]["sunshine_hours"])
\end{verbatim}

執行後會看到：

\begin{verbatim}
Traceback (most recent call last):
  File "bracket-keyerror-demo.py", line 4, in <module>
    print(data["current_weather"]["sunshine_hours"])
          ~~~~~~~~~~~~~~~~~~~~~~~^^^^^^^^^^^^^^^^^^
KeyError: 'sunshine_hours'
\end{verbatim}

這就是為什麼上面的範例查 \texttt{sunshine\_hours} 時改用 \texttt{.get()}，它在欄位不
存在的時候，能安全地回傳你指定的預設值，不會讓程式當掉。但這裡先留一個問題：如果連
「查詢這個動作本身」都失敗了呢？比如網路斷線、伺服器忙線，這種情況根本沒能拿到任何一包
資料，也就沒有哪個欄位存在或不存在的問題，\texttt{.get()} 這時候完全派不上用場，答案
留到第四節用 \texttt{try}/\texttt{except} 正式處理。

API 回來的資料常常沒有「一個城市配一個溫度」這麼單純，經常是一個地點底下還包著一長串
按時間排列的資料，這種「字典包串列、串列裡每個元素又是字典」的巢狀結構，在真實世界的
API 裡非常常見，值得單獨練習一次。

\begin{example}{python}{(hourly-forecast.py) Unpacking Nested Hourly Forecast Data}
import requests

url = (
    "https://api.open-meteo.com/v1/forecast"
    "?latitude=25.0330&longitude=121.5654"
    "&hourly=temperature_2m"
    "&timezone=Asia%2FTaipei"
    "&forecast_days=1"
)
response = requests.get(url)
data = response.json()

hourly = data["hourly"]            # 這是一個字典，裡面裝著好幾個等長的串列
times = hourly["time"]             # 一串時間戳記
temperatures = hourly["temperature_2m"]  # 一串對應的氣溫，跟 times 一一對齊

for time, temperature in zip(times[:5], temperatures[:5]):
    print(f"{time}: {temperature}°C")
\end{example}

\iosample{}{2026-09-01T00:00: 26.7°C \\ 2026-09-01T01:00: 26.3°C \\ 2026-09-01T02:00: 26.6°C \\ 2026-09-01T03:00: 26.3°C \\ 2026-09-01T04:00: 26.2°C \\}

拆解這種巢狀結構的訣竅，是一層一層剝，不用急著一次想通整包資料的形狀：先確認最外層
\texttt{data} 是字典還是串列、印出它的 \texttt{.keys()} 或長度看看裡面有什麼；再往下一
層做一樣的事，直到挖到你要的那個欄位為止。這裡 \texttt{data["hourly"]} 底下裝的是好幾個
「等長的串列」，\texttt{time} 跟 \texttt{temperature\_2m} 的長度永遠一樣，第 $i$ 個時間
對應第 $i$ 個氣溫，所以才能用 \texttt{zip()} 把它們兩個串列一組一組配對起來，交給
\texttt{for} 迴圈一起走一遍。
\index{zip()}

\stylepoint{面對一包沒看過的 JSON，不要急著憑猜測寫程式碼，先把它轉成 Python 物件、
用 \texttt{print()} 或互動模式一層一層看它的形狀，確認每一層是字典還是串列、鍵長什麼
樣子，再動手寫抓取欄位的程式碼。這比對照文件憑空想像可靠得多。}

\section{換一個 API 對照著看：查詢即時匯率}
\label{sec:currency-api}

前三節用的都是 Open-Meteo 這一個服務，這裡刻意換一個完全不同的服務再打一次，看看剛剛
學的那一套流程（發請求、轉 JSON、拆字典）是不是換個場景照樣管用。這一節要用的是
Frankfurter 這個公開的匯率查詢服務（\texttt{api.frankfurter.dev}），一樣不用註冊帳號、
不用申請金鑰，資料來源是歐洲央行公布的參考匯率，每個工作日更新一次。

Frankfurter 查最新匯率的網址固定是 \texttt{https://api.frankfurter.dev/v1/latest}，
常用的查詢參數有三個：\texttt{base} 指定「以哪個幣別為基準」（例如 \texttt{"USD"}
代表美元）、\texttt{symbols} 指定「想看哪個幣別的匯率」（可以用逗號隔開一次查好幾個）、
\texttt{amount} 指定「換算多少金額」，不帶這個參數時預設等於 1。

\begin{code}{python}{The Standard Pattern for Querying an Exchange Rate}
url = "https://api.frankfurter.dev/v1/latest"
params = {"base": base_currency, "symbols": target_currency}
response = requests.get(url, params=params, timeout=5)
data = response.json()
rate = data["rates"][target_currency]
\end{code}

\begin{example}{python}{(usd-to-jpy-rate.py) Checking the USD-to-JPY Exchange Rate}
import requests

url = "https://api.frankfurter.dev/v1/latest"
params = {"base": "USD", "symbols": "JPY"}
response = requests.get(url, params=params, timeout=5)
response.raise_for_status()
data = response.json()

print(data)

rate = data["rates"]["JPY"]
print(f"1 USD is about {rate} JPY")
\end{example}

\iosample{}{\{'amount': 1.0, 'base': 'USD', 'date': '2026-09-01', 'rates': \{'JPY': 160.16\}\} \\ 1 USD is about 160.16 JPY \\}

先看 \texttt{print(data)} 印出來的整包字典：最外層有 \texttt{amount}、\texttt{base}、
\texttt{date}、\texttt{rates} 四個鍵，跟第一節看到的天氣資料形狀不一樣（那邊最外層是
\texttt{latitude}、\texttt{current\_weather} 這些鍵），但骨子裡的道理完全相同：一樣是
先確認最外層字典有哪些鍵，再往你要的那個鍵一路挖下去。這裡 \texttt{data["rates"]} 本身
又是一個字典，鍵是幣別代碼、值是匯率，跟天氣資料裡 \texttt{current\_weather} 底下又包著
一個字典的巢狀寫法，是同一套結構。

\texttt{symbols} 也可以一次帶好幾個幣別，用逗號分隔：

\begin{example}{python}{(multi-currency-rates.py) Checking Rates for Three Currencies at Once}
import requests

url = "https://api.frankfurter.dev/v1/latest"
params = {"base": "USD", "symbols": "JPY,EUR,GBP"}
response = requests.get(url, params=params, timeout=5)
response.raise_for_status()
data = response.json()

for currency, rate in data["rates"].items():
    print(f"1 USD = {rate} {currency}")
\end{example}

\iosample{}{1 USD = 0.86281 EUR \\ 1 USD = 0.73904 GBP \\ 1 USD = 160.16 JPY \\}

這裡 \texttt{data["rates"]} 是一個裝著三組鍵值的字典，用第 \ref{ch:dict-set} 章學過的
\texttt{.items()} 就能同時走訪鍵（幣別代碼）跟值（匯率），不用另外分兩次分別取鍵跟值。

再加上 \texttt{amount} 這個參數，就能直接算好換算後的金額，不用自己另外乘一次：

\begin{example}{python}{(currency-converter.py) Converting an Amount}
import requests


def get_converted_amount(base, target, amount):
    url = "https://api.frankfurter.dev/v1/latest"
    params = {"base": base, "symbols": target, "amount": amount}
    response = requests.get(url, params=params, timeout=5)
    response.raise_for_status()
    data = response.json()
    return data["rates"][target]


amount = 100
converted = get_converted_amount("USD", "JPY", amount)
print(f"{amount} USD is about {round(converted, 1)} JPY")
\end{example}

\iosample{}{100 USD is about 16016 JPY \\}

這支函式把「組網址、發請求、拆欄位」包在 \texttt{get\_converted\_amount()} 裡面，呼叫端
只要給基準幣別、目標幣別、金額三個值，不用管底下實際打了哪一支 API，這是第 \ref{ch:functions} 章學過的
函式最基本的好處：把細節封裝起來，呼叫的人不用重複看一次實作。這裡也用到了第 \ref{ch:data-tools} 章介紹
過的 \texttt{round()}，把換算結果四捨五入到小數點下一位，畢竟匯率乘出來的小數位數常常
很長，直接印出來反而不好讀。

\keypoint{把 Open-Meteo 換成 Frankfurter，程式碼的骨架幾乎沒有變：一樣是
\texttt{requests.get(url, params=params)}，一樣是 \texttt{response.raise\_for\_status()}
加上 \texttt{response.json()}，一樣是照著回應的形狀一層一層挖出你要的欄位。真正會變的
只有網址跟參數名稱，這正是這一章想讓你親手體會的事：「軟體跟軟體對話」有一套通用的骨架，
換一個服務，套用同一套骨架，換掉細節就好。}

\begin{drills}[label=drills:currency-api]
\item 如果把 \texttt{params} 裡的 \texttt{symbols} 改成 \texttt{"SGD"}（新加坡幣），
你覺得 \texttt{data["rates"]} 這個字典裡會有幾組鍵值？ 那如果改成 \texttt{"TWD"}
（新台幣）呢？（提示：這一題請你真的跑跑看，兩次的結果不一樣。）
\item 承〈換算一筆金額〉那支範例：如果呼叫變成
\texttt{get\_converted\_amount("EUR", "GBP", 50)}，這支函式實際發出去的查詢參數
\texttt{base} 跟 \texttt{symbols} 分別會是什麼？
\item 如果完全不帶 \texttt{amount} 這個參數，你猜 \texttt{data["amount"]} 印出來會是
多少？（提示：回頭看〈查一次美金兌日圓的匯率〉那支範例印出的整包字典。）
\end{drills}

\section{會失敗，就是 EAFP 的主場：處理網路請求出錯}
\label{sec:eafp-requests}

打開檔案這件事，你在第 \ref{ch:files-exceptions} 章學過一個重要的心法：直接嘗試打開它，出了問題再用
\texttt{try}/\texttt{except} 接住，不用一路寫一堆「先確認檔案存不存在」的檢查，這叫
EAFP（先做，錯了再處理），是 Python 慣用的風格。

發網路請求，是這個心法最貼切的另一個
應用場景，甚至比檔案更極端：檔案至少還在你自己的電腦上，網路請求要仰賴的東西完全不在
你的掌控範圍，對方的伺服器可能當機、你的網路可能斷線、封包可能在路上就不見了。

就算你
事先花力氣「檢查」網路正不正常，檢查完到你真正發出請求之間，狀態也隨時可能又變了。
Python 面對這種「本來就可能失敗」的操作，一貫的態度是：直接做，錯了再處理。

\begin{code}{python}{Wrapping a Request in try/except}
try:
    response = requests.get(url, timeout=5)
    response.raise_for_status()   # 如果狀態碼不是 2xx，就丟出例外
except requests.exceptions.ConnectionError:
    # 處理連不上網路的情況
    ...
except requests.exceptions.Timeout:
    # 處理伺服器回應太慢的情況
    ...
except requests.exceptions.HTTPError:
    # 處理伺服器回傳錯誤狀態碼的情況
    ...
else:
    # 只有在完全沒出錯時才會執行到這裡，可以安心處理 response
    ...
\end{code}

\texttt{timeout=5} 是告訴 \texttt{requests}：最多等 5 秒，等太久就直接放棄、丟出
\texttt{Timeout} 例外，不要讓程式卡死在那裡乾等。
\index{requests.get(url, timeout=秒數)}
\texttt{raise\_for\_status()} 則是一個
方便的檢查：如果狀態碼代表失敗（例如 404、500），它會主動丟出 \texttt{HTTPError}，讓你
可以跟其他錯誤一樣，用同一套 \texttt{try}/\texttt{except} 接住，不用自己另外寫
\texttt{if response.status\_code != 200} 這種判斷。
\index{response.raise\_for\_status()}

\begin{example}{python}{(safe-weather-request.py) Safely Checking the Weather}
import requests

url = (
    "https://api.open-meteo.com/v1/forecast"
    "?latitude=25.0330&longitude=121.5654"
    "&current_weather=true"
)

try:
    response = requests.get(url, timeout=5)
    response.raise_for_status()
except requests.exceptions.ConnectionError:
    print("Can't reach the network, please check your connection")
except requests.exceptions.Timeout:
    print("The server took too long to respond, timed out, try again later")
except requests.exceptions.HTTPError:
    print(f"The server returned an error, status code: {response.status_code}")
else:
    data = response.json()
    temperature = data["current_weather"]["temperature"]
    print(f"Current temperature: {temperature}°C")
\end{example}

\iosample{}{Current temperature: 25.6°C \\}

\keypoint{檔案可能不存在、網路可能斷線，這兩件事的心法完全一樣：先做，錯了再處理，不要
一路寫一堆「先檢查網路正不正常」「先確認伺服器有沒有回應」的程式碼。你在第 \ref{ch:files-exceptions} 章學檔案
讀取、第 \ref{ch:debugging} 章學讀懂錯誤訊息時已經看過這個心法兩次，網路請求是它最不受你控制、也最能
體現「為什麼 Python 偏好先做再處理」的一個場景，因為根本沒有「先檢查」這種選項，你檢查
完到真正發出請求之間，網路狀況隨時可能又不一樣了。}

事實上，只要網址打錯或伺服器回了非預期的東西，錯誤訊息長什麼樣子、該怎麼讀，也值得
實際看一次，不要只靠想像。

\begin{example}{python}{(check-url-typo.py) What Happens with a Typo in the URL}
import requests

# 情況一：網域打錯字，requests 連不上伺服器
bad_url = "https://api.open-meteoo.com/v1/forecast"
try:
    requests.get(bad_url, timeout=5)
except requests.exceptions.ConnectionError as error:
    print("Connection failed:")
    print(error)

# 情況二：路徑打錯字，伺服器收到請求，但回傳 404
url = (
    "https://api.open-meteo.com/v1/foercast"
    "?latitude=25.03&longitude=121.56"
    "&current_weather=true"
)
response = requests.get(url, timeout=5)
print("Status code:", response.status_code)

try:
    response.raise_for_status()
except requests.exceptions.HTTPError as error:
    print("Error raised by raise_for_status():")
    print(error)
\end{example}

情況一（網域打成 \texttt{open-meteoo.com}，多了一個 \texttt{o}）印出的內容，實際在這台
機器上跑起來是這樣（原本是同一行，這裡一樣手動斷開方便印在紙上）：

\begin{verbatim}
Connection failed:
HTTPSConnectionPool(host='api.open-meteoo.com', port=443):
Max retries exceeded with url: /v1/forecast (Caused by
NameResolutionError("HTTPSConnection(
host='api.open-meteoo.com', port=443): Failed to resolve
'api.open-meteoo.com' ([Errno 8] nodename nor servname
provided, or not known)"))
\end{verbatim}

這段訊息看起來嚇人，但抓重點就好：關鍵字 \texttt{Failed to resolve} 是說「找不到這個
網域」，這正是網址打錯字最典型的癥狀，跟你在瀏覽器打錯網址會看到
「找不到這個網站」是同一回事。情況二（路徑打成 \texttt{foercast}）的結果不一樣：網域是
對的，所以連得上，伺服器也確實回了東西，只是內容是「找不到這個路徑」：

\begin{verbatim}
Status code: 404
Error raised by raise_for_status():
404 Client Error: Not Found for url:
https://api.open-meteo.com/v1/foercast?latitude=25.03
&longitude=121.56&current_weather=true
\end{verbatim}

這裡的關鍵差異，正好對應到前面片語裡的兩個不同 \texttt{except}：連不上網域，是
\texttt{ConnectionError}，通常代表網址本身打錯了、或裝置根本沒有網路；連得上但伺服器
說「這不是我認得的東西」，是 \texttt{HTTPError}，代表網址有一部分（路徑、參數）不對，
或是你查詢的東西伺服器那邊剛好沒有。
\index{requests.exceptions.ConnectionError／Timeout／HTTPError}

舉個例子，如果你把緯度打成一個超出地球範圍的數字
（例如 \texttt{999}），Open-Meteo 也會用同一種方式告訴你哪裡不對，回應狀態碼 400，
內容是 \texttt{"reason":"Latitude must be in range of -90 to 90°. Given: 999.0."}，
一樣可以被 \texttt{raise\_for\_status()} 抓到、當成 \texttt{HTTPError} 處理。

換一個服務結果也一樣：如果你把上一節查匯率用的 \texttt{symbols} 改成一個根本不存在的
幣別代碼（例如 \texttt{XXX}），Frankfurter 回應的狀態碼是 404，內容是
\texttt{\{"message":"not found"\}}，跟前面 Open-Meteo 回的 400 加上一句人話說明的格式
完全不一樣。

\begin{example}{python}{(bad-currency-code.py) Querying a Nonexistent Currency Code}
import requests

url = "https://api.frankfurter.dev/v1/latest"
params = {"base": "USD", "symbols": "XXX"}

try:
    response = requests.get(url, params=params, timeout=5)
    response.raise_for_status()
except requests.exceptions.HTTPError as error:
    print("Error raised by raise_for_status():")
    print(error)
\end{example}

\begin{verbatim}
Error raised by raise_for_status():
404 Client Error: Not Found for url: https://api.frankfurter.dev/v1/latest?base=USD&symbols=XXX
\end{verbatim}

不同服務「查詢失敗」的長相不一定一樣，有的用 400、有的用 404，錯誤訊息的措辭也各自
不同，但只要統一用 \texttt{raise\_for\_status()} 搭配 \texttt{try}/\texttt{except
HTTPError} 接住，就不用為每一個服務各自另外記一套判斷邏輯，這也是這一節想讓你體會的：
EAFP 這套心法不只省了「先檢查」的力氣，還讓你的程式碼不用管背後那支服務用哪個狀態碼、
怎麼措辭，反正失敗了就是丟例外，接住、處理，這件事永遠是同一套寫法。

\begin{drills}[label=drills:eafp-requests]
\item 用前面〈安全地查一次天氣〉那支程式當範本：如果這次網域打錯字，導致連線失敗，
螢幕上會印出哪一行？
\item 如果這次伺服器在 5 秒內沒有回應（觸發逾時），螢幕上會印出哪一行？
\item 如果這次請求成功連上伺服器，但伺服器回了狀態碼 404，螢幕上會印出哪一行？（提示：
別忘了程式裡呼叫了 \texttt{response.raise\_for\_status()}。）
\item 如果這次查匯率、幣別代碼打錯字，導致伺服器回了 404，〈查一個不存在的幣別代碼〉
這支程式螢幕上會印出哪兩行？
\end{drills}

\section{把抓到的資料存下來：寫成 CSV、畫成圖表}
\label{sec:save-and-chart}

前面每一節查完一次資料，程式一結束，那包資料也就跟著消失了，下次要看還得重新打一次
API。實際上，API 抓回來的資料常常值得保存下來，之後不用重新查詢就能直接拿來用，或是
畫成圖表讓人一眼看懂。這一節要做的事，其實就是把第 \ref{ch:csv} 章學過的 \texttt{csv} 模組跟
第 \ref{ch:data-tools} 章學過的 matplotlib，接到這一章新學的 \texttt{requests} 後面：查一次 API，把
結果存成一份 CSV 檔，再從那份 CSV 檔畫出一張圖。

\begin{code}{python}{The Standard Pattern for Saving Fetched Data as CSV}
import csv

with open("filename.csv", "w", newline="") as f:
    writer = csv.writer(f)
    writer.writerow([field1, field2])
    for value1, value2 in zip(list1, list2):
        writer.writerow([value1, value2])
\end{code}

\begin{example}{python}{(fetch-hourly-to-csv.py) Saving a Day's Hourly Temperatures as CSV}
import csv
import requests

url = "https://api.open-meteo.com/v1/forecast"
params = {
    "latitude": 25.0330,
    "longitude": 121.5654,
    "hourly": "temperature_2m",
    "timezone": "Asia/Taipei",
    "forecast_days": 1,
}
response = requests.get(url, params=params, timeout=5)
response.raise_for_status()
data = response.json()

hourly = data["hourly"]
times = hourly["time"]
temperatures = hourly["temperature_2m"]

with open("taipei_hourly.csv", "w", newline="") as f:
    writer = csv.writer(f)
    writer.writerow(["time", "temperature"])
    for time, temperature in zip(times, temperatures):
        writer.writerow([time, temperature])

print(f"Wrote {len(times)} rows to taipei_hourly.csv")
\end{example}

\iosample{}{Wrote 24 rows to taipei\_hourly.csv \\}

留意 \texttt{open()} 那一行多寫了 \texttt{newline=""} 這個參數，第 \ref{ch:csv} 章已經解釋過，
這是使用 \texttt{csv.writer()} 時官方文件明講的固定寫法，忘記加的話，某些系統上可能會
在每一列中間多出一行空白，不是程式邏輯錯了，養成每次都加上它的習慣就好。執行完之後，
打開 \texttt{taipei\_hourly.csv} 會看到（這裡只列出前五列）：

\begin{verbatim}
time,temperature
2026-09-01T00:00,25.6
2026-09-01T01:00,25.6
2026-09-01T02:00,25.4
2026-09-01T03:00,25.2
2026-09-01T04:00,25.1
\end{verbatim}

這份檔案存下來之後，就跟第 \ref{ch:csv} 章讀過的任何一份 CSV 沒有兩樣了，可以用
\texttt{csv.DictReader} 讀回來，接著用第 \ref{ch:data-tools} 章學過的 \texttt{plt.bar()} 畫成圖：

\begin{example}{python}{(chart-from-csv.py) Reading the CSV Back and Plotting a Bar Chart}
import csv
import matplotlib.pyplot as plt

times = []
temperatures = []
with open("taipei_hourly.csv") as f:
    reader = csv.DictReader(f)
    for row in reader:
        times.append(row["time"][-5:])            # 只留下最後 5 個字元，"HH:MM"
        temperatures.append(float(row["temperature"]))

plt.bar(times[:8], temperatures[:8])
plt.xlabel("Hour")
plt.ylabel("Temperature (C)")
plt.title("Taipei Hourly Temperature")
plt.show()
\end{example}

這裡有兩個地方值得多看一眼。第一，\texttt{row["time"][-5:]} 用的是第 \ref{ch:strings} 章學過的切片，
從像 \texttt{"2026-09-01T00:00"} 這樣完整的時間字串裡，只留下最後五個字元
（\texttt{"00:00"}），圖表的橫軸才不會擠得看不清楚。第二，\texttt{float(row["temperature"])}
不能省略：CSV 檔案裡存的所有內容，讀回來都是字串，就算看起來像數字，第 \ref{ch:csv} 章已經強調過
這件事，這裡的溫度也不例外，要拿來畫圖或做任何數學運算之前，得先轉型別。

執行完你會看到
一個彈出的視窗，裡面是 8 根長條，代表凌晨前八個小時的氣溫變化，跟第 \ref{ch:data-tools} 章畫班級平均分數
長條圖的結果長得很像，只是這次橫軸換成了時間、資料來源換成了剛剛存下來的 CSV 檔，不是
你自己手打的清單。

\begin{drills}[label=drills:save-and-chart]
\item 如果 \texttt{open("taipei\_hourly.csv", "w", newline="")} 那一行不小心漏寫
\texttt{newline=""}，在某些系統上會有什麼後果？（提示：回頭看第 \ref{ch:csv} 章學過的坑。）
\item 如果只想存下前 12 小時的資料，〈把一整天的每小時氣溫存成 CSV〉那支程式裡的
\texttt{zip(times, temperatures)} 該怎麼改？
\item 〈讀回 CSV、畫成長條圖〉這支程式裡，如果漏寫 \texttt{float()}，直接把
\texttt{row["temperature"]} 這個字串傳給 \texttt{plt.bar()}，你覺得畫出來的圖會有什麼
問題？
\end{drills}

\section{API 金鑰：需要認證的服務，怎麼安全存放}
\label{sec:api-keys}

這一章從頭到尾用的 Open-Meteo 跟 Frankfurter，刻意挑的都是不需要帳號、不需要金鑰的
服務，方便你專心體會「打 API」這件事本身。

但實際情況裡，更多公開 API（地圖服務、部分
天氣服務、翻譯服務都很常見）要求你先註冊一個帳號，換到一組專屬於你的「API 金鑰」（API
key）才能查詢。金鑰本質上是一段文字，用來讓伺服器知道「這次請求是誰發出來的」，通常是
為了控制每個人能查詢的次數，或是把使用量跟你的帳號對起來計費。

大部分服務的做法，是把這組金鑰當成第二節學過的查詢參數之一，跟其他條件一起放進
\texttt{params} 字典：

\begin{code}{python}{The Standard Pattern for Querying with an API Key}
API_KEY = "the-long-string-you-registered-for"
params = {
    "q": query,
    "appid": API_KEY,
}
response = requests.get(url, params=params, timeout=5)
\end{code}

第 \ref{ch:industry} 章已經完整講過這裡真正要注意的事：金鑰是機密資訊，絕對不能像上面片語那樣，直接
把一長串文字寫死在會被 commit 進版本控制系統的程式碼裡，一旦某次 commit 留下紀錄，之後
就算刪掉那一行重新 commit，舊的紀錄裡還是找得到。

第 \ref{ch:industry} 章介紹過的 \texttt{.env} 檔案跟
\texttt{.gitignore}，就是為了應付這種情況：把金鑰放進一個獨立、不會被記錄進版本控制
系統的檔案裡，程式執行的時候再讀進來用，這樣即使你把程式碼分享出去、甚至公開在網路上，
金鑰也不會跟著外流。

這一章不深入示範讀取 \texttt{.env} 的細節，記住這個原則就好：申請到
金鑰的那一刻起，就把它當成密碼對待，不要圖方便直接寫死在程式碼裡。

\section{小專案：串成一支查詢小工具}
\label{sec:weather-tool}

前面每一節分別學了怎麼發出請求、怎麼帶查詢參數、怎麼拆解回應、怎麼處理失敗、怎麼存下
結果，這一節要把這些東西串成一支真正能用的小工具：讓使用者選擇要查天氣還是查匯率，輸入
對應的城市或幣別，程式就去查詢、把結果講成一句人看得懂的話，查完可以繼續查下一筆，直到
使用者想離開為止。

跟第 \ref{ch:automating} 章的整合應用專案不一樣，\ref{ch:automating} 章的目標是把整本書學過的程式技巧
全部兜在一起、做出一支功能完整的小工具；這裡刻意窄很多，不做快取、不做介面、也不處理
需要申請帳號才能查到的項目，唯一目的就是完整走一次「你的程式主動跟外部服務要資料」的
流程，

順便親手體會一次：同一支程式裡，接兩支完全不同的 API，用的是同一套骨架，讓「軟體
怎麼跟軟體溝通」這件事，變成你自己動手做過、不是只在書上讀過的東西。

\begin{example}{python}{(weather-currency-tool.py) A Weather-or-Exchange-Rate Lookup Tool}
import requests

CITIES = {
    "Taipei": (25.0330, 121.5654),
    "Taichung": (24.1477, 120.6736),
    "Kaohsiung": (22.6273, 120.3014),
    "Hualien": (23.9871, 121.6015),
}

CURRENCIES = ["JPY", "EUR", "GBP"]


def get_weather(latitude, longitude):
    url = "https://api.open-meteo.com/v1/forecast"
    params = {
        "latitude": latitude,
        "longitude": longitude,
        "current_weather": "true",
    }
    response = requests.get(url, params=params, timeout=5)
    response.raise_for_status()
    return response.json()["current_weather"]


def get_rate(target):
    url = "https://api.frankfurter.dev/v1/latest"
    params = {"base": "USD", "symbols": target}
    response = requests.get(url, params=params, timeout=5)
    response.raise_for_status()
    return response.json()["rates"][target]


while True:
    choice = input("Enter w for weather, e for exchange rate, q to quit: ")
    if choice == "q":
        print("Goodbye!")
        break

    if choice == "w":
        city = input(f"Which city? ({'/'.join(CITIES)}): ")
        if city not in CITIES:
            print("Only Taipei, Taichung, Kaohsiung, and Hualien are supported, please try again")
            continue
        latitude, longitude = CITIES[city]
        try:
            weather = get_weather(latitude, longitude)
        except requests.exceptions.ConnectionError:
            print("Can't reach the network, please check your connection")
            continue
        except requests.exceptions.Timeout:
            print("The server took too long to respond, try again later")
            continue
        except requests.exceptions.HTTPError:
            print("The server returned an error, please try again later")
            continue
        temperature = weather["temperature"]
        wind_speed = weather["windspeed"]
        print(f"{city}: temperature {temperature}°C, wind speed {wind_speed} km/h")

    elif choice == "e":
        target = input(f"Convert USD to which currency? ({'/'.join(CURRENCIES)}): ")
        if target not in CURRENCIES:
            print("Only JPY, EUR, and GBP are supported, please try again")
            continue
        try:
            rate = get_rate(target)
        except requests.exceptions.ConnectionError:
            print("Can't reach the network, please check your connection")
            continue
        except requests.exceptions.Timeout:
            print("The server took too long to respond, try again later")
            continue
        except requests.exceptions.HTTPError:
            print("The server returned an error, please try again later")
            continue
        print(f"1 USD = {rate} {target}")

    else:
        print("Unrecognized command, please try again")
\end{example}

\iosample{w \\ Taipei \\ e \\ JPY \\ x \\ q \\}{Enter w for weather, e for exchange rate, q to quit: Which city? (Taipei/Taichung/Kaohsiung/Hualien): Taipei: temperature 25.4°C, wind speed 3.5 km/h \\ Enter w for weather, e for exchange rate, q to quit: Convert USD to which currency? (JPY/EUR/GBP): 1 USD = 160.16 JPY \\ Enter w for weather, e for exchange rate, q to quit: Unrecognized command, please try again \\ Enter w for weather, e for exchange rate, q to quit: Goodbye! \\}

這支程式把前面每一節都用上了：

\begin{itemize}
\item \texttt{CITIES} 跟 \texttt{CURRENCIES} 兩個容器把使用者
能查詢的範圍固定下來，省去每次都要自己輸入座標或幣別代碼的麻煩。
\item \texttt{get\_weather()}
跟 \texttt{get\_rate()} 兩支函式，各自包著第一節學的 \texttt{requests.get()} 加上第二節
學的 \texttt{params} 字典，還有第四節的 \texttt{raise\_for\_status()}。
\item 拿到回應之後，用
第三節學的中括號取出需要的欄位。
\item 每一種查詢外層都用 \texttt{try}/\texttt{except} 接住
所有可能的失敗，讓程式只印一句話、讓使用者可以繼續查下一筆，不會因為一次查詢失敗就整個
當掉。
\item 最外層的 \texttt{while True} 搭配 \texttt{if}/\texttt{elif}/\texttt{else}，讓
使用者不用每查一次就重新執行一次程式，也能在天氣跟匯率之間自由切換。
\end{itemize}

\begin{drills}[label=drills:weather-tool]
\item 開放式練習：幫這支程式加一個功能，例如多印一句更口語化的天氣描述（提示：
\texttt{weather["weathercode"]} 是一個數字代碼，代表當下的天氣狀況，你可以查
Open-Meteo 文件裡代碼跟天氣描述的對照表，寫一個小字典把常見代碼轉成「晴天」「多雲」
「下雨」這類文字）；或是在 \texttt{CITIES} 字典裡多加幾個你自己感興趣的城市。
\item 開放式練習：幫這支程式加一個「記錄」功能，不管使用者查的是天氣還是匯率，每查
一次就用第五節學過的 \texttt{csv.writer()} 加一列記錄，寫進一個 \texttt{log.csv} 檔案
裡（提示：記錄的欄位不完全一樣，天氣要記城市跟氣溫，匯率要記幣別跟匯率，可以用一個共通
的「查詢類型」欄位區分兩種紀錄）。
\end{drills}

\section*{練習解答}
\addcontentsline{toc}{section}{練習解答}

\paragraph{第~\ref{sec:query-strings}~節}
\begin{enumerate}
\item 空白字元通常會被轉成加號 \texttt{+} 或 \texttt{\%20}，實際結果依 \texttt{requests}
內部使用的編碼規則而定，但重點是它不會原封不動地留下一個空白字元。
\item 手動寫的那個網址「剛好能動」是因為瀏覽器或這台伺服器剛好願意寬容處理沒編碼過的
空白，不代表所有伺服器都會這樣。程式裡用 \texttt{params} 字典交給 \texttt{requests}
處理，可以確保不管值裡出現什麼特殊符號，都會照標準規則轉換好，不用賭對方伺服器夠不夠
寬容。
\end{enumerate}

\paragraph{第~\ref{sec:currency-api}~節}
\begin{enumerate}
\item 改成 \texttt{"SGD"} 是一組——\texttt{symbols} 只帶一個幣別代碼，
\texttt{data["rates"]} 就只會有 \texttt{"SGD"} 這一個鍵。 但改成 \texttt{"TWD"} 你會拿到
一個 \texttt{HTTPError}：Frankfurter 用的是歐洲央行公布的參考匯率，總共只有三十種幣別，
裡面\textbf{沒有新台幣}，查不到的幣別它回的是 404，於是
\texttt{response.raise\_for\_status()} 當場把程式擋下來，根本走不到
\texttt{data["rates"]} 那一行。 這件事值得記住：一支 API 能回答什麼，是它自己決定的，
不會因為那個資料「聽起來理所當然存在」就一定查得到。 遇到這種情況，要嘛換一個查得到的
服務，要嘛接受這個服務的範圍——第 \ref{ch:capstone} 章的專案就會真的碰上這個決定。
\item \texttt{base} 會是 \texttt{"EUR"}，\texttt{symbols} 會是 \texttt{"GBP"}，因為
函式呼叫時第一個引數對應 \texttt{base}、第二個引數對應 \texttt{target}（也就是
\texttt{symbols}）。
\item 會是 \texttt{1.0}。不帶 \texttt{amount} 這個參數時，Frankfurter 預設把它當成 1，
所以 \texttt{data["amount"]} 印出來是 \texttt{1.0}。
\end{enumerate}

\paragraph{第~\ref{sec:eafp-requests}~節}
\begin{enumerate}
\item 印出「Can't reach the network, please check your connection」，網域打錯字會讓
\texttt{requests.get()} 丟出 \texttt{ConnectionError}，被對應的 \texttt{except} 接住。
\item 印出「The server took too long to respond, timed out, try again later」，超過
\texttt{timeout=5} 秒沒有回應會丟出 \texttt{Timeout}，被對應的 \texttt{except} 接住。
\item 印出「The server returned an error, status code: 404」，雖然請求本身成功連上了
伺服器，但 \texttt{response.raise\_for\_status()} 看到狀態碼是 404，會主動丟出
\texttt{HTTPError}，被對應的 \texttt{except} 接住。
\item 印出「Error raised by raise\_for\_status():」，接著印出
\texttt{404 Client Error: Not Found for url: https://api.frankfurter.dev/v1/latest?base=USD\&symbols=XXX}
這一行，道理跟前一題一樣：狀態碼 404 讓 \texttt{raise\_for\_status()} 丟出
\texttt{HTTPError}，被對應的 \texttt{except} 接住並印出來。
\end{enumerate}

\paragraph{第~\ref{sec:save-and-chart}~節}
\begin{enumerate}
\item 在某些系統上，每一列中間可能會多出一行空白，這是文字模式處理換行符號的方式跟
\texttt{csv} 模組預期的不一致所導致，不是程式邏輯本身有錯誤，但仍然會讓輸出的 CSV 檔案
格式跑掉，養成一律加上 \texttt{newline=""} 的習慣才不用擔心這件事。
\item 改成 \texttt{zip(times[:12], temperatures[:12])}，用切片先把兩個串列都只留下前
12 個元素，再交給 \texttt{zip()} 配對。
\item \texttt{plt.bar()} 會拿到一串字串而不是數字，長條圖的高度沒辦法正確反映氣溫的
大小關係，甚至可能直接因為型別不對而報錯，這正是第 \ref{ch:csv} 章一直強調「CSV 讀回來的內容全部
是字串」這件事在這裡的具體後果。
\end{enumerate}

\paragraph{第~\ref{sec:weather-tool}~節}
\begin{enumerate}
\item 開放性練習，只要程式仍然能正常執行、確實新增了你描述的功能即可，例如把
\texttt{weathercode} 轉成文字描述時，記得處理「代碼不在你自己寫的對照表裡」的情況，
可以用 \texttt{.get()} 給一個「未知天氣狀況」之類的預設值。
\item 開放性練習，重點是每次查詢完（不管成功或失敗）都要記得幫 CSV 加一列記錄，並且
留意 \texttt{open()} 要用附加模式（第 \ref{ch:files-exceptions} 章學過的 \texttt{'a'} 模式）打開記錄檔，而不是
每次都用 \texttt{'w'} 模式，不然前面查過的紀錄會被整個清空重寫。
\end{enumerate}

\chapter{Capstone：做一個像樣的專案}
\label{ch:capstone}

你在第 \ref{ch:automating} 章已經做過一次整合了。那一章把你前面學過的技巧接成一支能整理
檔案的腳本，能分類、能搬移、能打包備份，甚至能上網抓東西回來。那支程式很好用，但它是
\textbf{一支腳本}：所有東西擠在同一個檔案裡，跑完就結束，下一次執行的時候，它完全不記得
上一次做了什麼。

這一章要做的是另一件事：把同一批技巧，組織成一個\textbf{像真的軟體的專案}。差別不在難度，
在「要同時顧的事情變多了」——程式要拆成好幾個檔案、要有人記得資料放在哪、要能重複執行而且
記得上次的結果、要有辦法確認自己沒有改壞、還要讓別人拿到這個資料夾就知道怎麼跑。

第 \ref{ch:api} 章那支查詢小工具也一樣，它做的是「打一次 API、把答案講給你聽」，查完就忘。
這一章的 API 只是整個專案裡的\textbf{一塊零件}：它負責去外面拿一個數字回來，拿回來之後
怎麼存、怎麼算、怎麼變成一份報表，才是這一章真正的重點。

還有一件事要先講清楚，而且它比上面幾段都重要：\textbf{這一章不會教你任何新語法。一個都
沒有。} 你接下來會用到的每一樣東西，前面的章節都已經教過了。如果你在某一節覺得吃力，
那個吃力不會是「這個語法我沒學過」，而是「東西太多，我抓不住」——這正是這一章要練的東西，
也是所有初學者第一次面對一個真正的專案時，唯一會卡住的地方。

這一章分三段推進：

\begin{itemize}
\item 第一段從一句人話開始。我們會先寫下一段需求，當場把它拆成幾個檔案，把專案資料夾
建起來，再決定「一筆花費」在程式裡該長什麼樣子
（第 \ref{sec:capstone-spec}～\ref{sec:capstone-model} 節）。
\item 第二段把零件一個一個做出來：把紀錄存進檔案再讀回來、去外面問一次匯率、把一堆紀錄
變成一份看得懂的報表。每做完一個零件，整個專案都會比上一節多做一點事，而且立刻跑得起來
（第 \ref{sec:capstone-storage}～\ref{sec:capstone-report} 節）。
\item 第三段是「交出去」該做的事：寫一支檢查腳本確認自己沒有改壞、把六個檔案接成一支
真正能用的工具、最後對照第 \ref{ch:industry} 章那些詞彙，看看你剛剛做的每一件事在業界
叫什麼名字（第 \ref{sec:capstone-tests}～\ref{sec:capstone-handoff} 節）。
\end{itemize}

還有一個承諾要先給你：\textbf{每一節結束，你手上都會有一份能執行的程式}。這一章沒有任何
一節是「先忍耐、最後一節才跑得起來」的。

\section{從一句需求開始：把它拆成幾塊}
\label{sec:capstone-spec}

真實世界裡，一個專案很少是從「我想寫一個 class」開始的，通常是從一句抱怨開始的。這一章
的抱怨是這樣：

\begin{quote}
我常常有一些花費不是台幣——網購日本的東西、訂閱美金計價的服務、出國旅遊刷卡。我想要一支
工具：每次花錢就記一筆下來，之後隨時可以問它「我這陣子總共花了多少台幣、哪一類花最多」，
外幣的部分請它自己去查匯率換算好。
\end{quote}

這段話就是這個專案的\textbf{規格}（spec），也就是「這個東西該做到什麼」的白紙黑字版本。
\index{規格（spec）}
它不長，但它已經把要做的事情全部講完了，接下來的工作只有一件：一句一句讀它，問「這句話
要程式做什麼事」，然後把答案分堆。

\begin{center}
\begin{tabular}{@{}p{5.2cm}p{3.6cm}p{4.2cm}@{}}
\toprule
規格裡的話 & 要程式做什麼 & 變成哪個檔案 \\
\midrule
「每次花錢就記一筆下來」 & 描述一筆花費 & \texttt{expense.py} \\
「記下來」「之後隨時可以問它」 & 存進檔案、讀回來 & \texttt{storage.py} \\
「請它自己去查匯率換算」 & 去外面拿匯率 & \texttt{rates.py} \\
「總共花了多少、哪一類花最多」 & 彙總、印成報表 & \texttt{report.py} \\
「我想要一支工具」 & 跟使用者互動 & \texttt{main.py} \\
（規格裡沒寫，但你會需要） & 確認自己沒改壞 & \texttt{check\_moneylog.py} \\
\bottomrule
\end{tabular}
\end{center}

這張表就是整個專案的地圖，也是這一章接下來六節的順序。它看起來像是憑感覺分的，其實只用了
一個判準：\textbf{每一個檔案負責一件事，而且那件事講得出一句完整的話}。「描述一筆花費」是
一句話，「去外面拿匯率」也是一句話；如果某個檔案要用「而且⋯⋯還要⋯⋯」才講得完，那它八成
該被拆成兩個。

順著同一段規格，還可以再多決定一件事：一筆花費要記哪些欄位？「每次花錢就記一筆」至少要有
日期跟項目；「哪一類花最多」需要一個分類；「總共花了多少台幣」需要金額，而且因為有外幣，
還需要知道這個金額是哪一國的錢。五個欄位：\textbf{日期、項目、分類、金額、幣別}。

這個專案叫 \texttt{moneylog}。開工之前先把資料夾準備好，這幾個步驟你在第 \ref{ch:industry}
章都讀過名字，這次要真的動手做一次：

\begin{enumerate}
\item 建一個叫 \texttt{moneylog} 的資料夾，之後所有檔案都放在裡面。
\item 在這個資料夾裡建一個虛擬環境（第 \ref{sec:project-structure} 節解釋過它在解決什麼
問題：讓每個專案有自己一份套件清單，不會互相打架）。
\item 在虛擬環境裡執行 \texttt{pip install requests}，這是這個專案唯一需要的第三方套件，
第 \ref{ch:api} 章已經裝過一次。
\item 把它記進 \texttt{requirements.txt}。
\item 寫一份 \texttt{.gitignore}。
\item 開一個空的 \texttt{README.md}，先放標題就好，完整內容留到第 \ref{sec:capstone-wiring}
節再補。
\end{enumerate}

\texttt{requirements.txt} 只有一行。版本號請照你自己裝到的那一版寫，不用跟這裡一樣：

\begin{code}{text}{(moneylog/requirements.txt) 這個專案需要哪些套件}
requests==2.32.5
\end{code}

\texttt{.gitignore} 三行：

\begin{code}{text}{(moneylog/.gitignore) 哪些東西不進版本控制}
venv/
__pycache__/
*.csv
\end{code}

前兩行跟第 \ref{sec:gitignore} 節那份範例一模一樣：虛擬環境資料夾可以靠
\texttt{requirements.txt} 重建，\texttt{\_\_pycache\_\_} 是 Python 自己產生的暫存品——
而且這一次你真的會看到它出現，因為這個專案開始用 \texttt{import} 匯入自己寫的模組了。

第三行 \texttt{*.csv} 是這個專案特有的。這個資料夾裡的 CSV 檔案全部都是\textbf{程式跑出來
的資料}，不是你手寫的程式碼：\texttt{expenses.csv} 是使用者自己的花費紀錄，
\texttt{summary.csv} 是每次出報表都會整份覆蓋掉的產出。版本控制要保存的是「你寫了什麼
程式」，不是「這台電腦上的人花了多少錢」，所以兩份都排除掉。星號的意思跟第
\ref{sec:gitignore} 節那行 \texttt{*.pyc} 一樣：不管檔名前面是什麼都算。

最後寫第一支 \texttt{main.py}。它現在什麼都不做，只印兩行字，目的只有一個：\textbf{確認
你的環境是通的}，Python 跑得起來、資料夾位置對。

\begin{example}{python}{(moneylog/main.py 第 0 版) 先確認跑得起來}
TOOL_NAME = "moneylog"
DATA_FILE = "expenses.csv"

print(f"{TOOL_NAME}: a multi-currency expense log")
print(f"Records will be stored in {DATA_FILE}")
\end{example}

\iosample{}{moneylog: a multi-currency expense log \\ Records will be stored in expenses.csv \\}

在 \texttt{moneylog} 資料夾裡執行 \texttt{python main.py}，看到這兩行，這一節就結束了。
五行程式碼看起來很可笑，但它證明了一件事：從現在開始，任何一次執行失敗都是你剛剛寫的東西
造成的，不是環境的問題。

\keypoint{把規格寫下來，不是形式主義。它有兩個非常實際的用途：第一，三個月後你回頭打開
這個資料夾，你會忘記自己當初想做什麼，那段文字是你唯一的線索；第二，這個專案如果要交給
別人接手，一份寫清楚「這是什麼、怎麼跑」的說明，常常是對方唯一的入口——這正是第
\ref{sec:project-anatomy} 節講 README 時說過的那件事。差別只在於，那時候你是讀到它，
這次你要自己寫一份。}

\stylepoint{這個專案刻意把所有 \texttt{.py} 檔案放在同一層，沒有像第
\ref{sec:project-anatomy} 節第二張樹狀圖那樣分成 \texttt{src/} 跟 \texttt{tests/} 兩個子
資料夾。理由很誠實：模組一旦放進子資料夾，\texttt{import} 就要多一套規則才找得到它，而第
\ref{ch:modules} 章明確講過「自己寫的模組要跟呼叫它的程式放在同一個資料夾裡」，那套規則
這本書沒有教。真實的專案確實常常會分層，你以後一定會遇到；但分層不是這一章要練的東西，
這一章要練的是「這件事該放進哪一個檔案」。這兩件事沒有關係——扁平的資料夾一樣可以分工
分得很清楚。}

\begin{drills}[label=drills:capstone-spec]
\item 使用者又提了一個需求：「我還想知道每個月分別花了多少。」照這一節那張表的分法，
這個功能主要應該加進哪一個檔案？為什麼不是 \texttt{storage.py}？
\item 下面四樣東西，哪些該被寫進 \texttt{.gitignore}：\texttt{main.py}、虛擬環境資料夾、
\texttt{requirements.txt}、\texttt{expenses.csv}？
\item 如果把 \texttt{rates.py} 跟 \texttt{report.py} 兩個檔案合併成一個叫
\texttt{money.py} 的檔案，程式還是會動。用這一節那個「講得出一句完整的話」的判準看，
合併之後的檔案要怎麼描述它在做什麼？
\end{drills}

\section{先決定資料長什麼樣：\texttt{expense.py}}
\label{sec:capstone-model}

在寫任何一個功能之前，要先決定一件事：一筆花費在程式裡，到底是什麼東西？

你手上有兩個選項，兩個都學過。一個是 dict：\texttt{\{"date": "2026-09-01", "item":
"bento", ...\}}，五個欄位五個鍵，寫起來最快。另一個是 class（第 \ref{ch:objects} 章）。

怎麼選？第 \ref{sec:methods-vs-functions} 節給過判準，這裡把它換成一個更好問的問題：
\textbf{你最常對這份資料做什麼？}

如果答案是「照順序走一遍、印出來」，那 dict 就夠了，你在第 \ref{ch:csv} 章整章都是這樣
用的。但這裡的答案不是。這裡每一筆花費除了「有五個欄位」以外，還要\textbf{會做一件事}：
把自己換算成台幣。這件事需要用到它自己的 \texttt{amount} 跟 \texttt{currency}，換句話說，
資料跟動作是綁在一起的——這正是第 \ref{sec:why-objects} 節說「該用物件了」的那個訊號。

還有一個決定要先講：\textbf{這裡不需要繼承}。第 \ref{sec:inheritance-vs-composition} 節
教過一個判準：「A 是一種 B」用繼承，「A 裡面有一個 B」用組合。這個專案裡只有一種花費，
沒有「日圓花費是一種花費」這種階層，也沒有「花費裡面裝著另一個物件」。學過的技巧不代表
每次都要用，這是 capstone 該練的判斷之一。

\begin{example}{python}{(moneylog/expense.py) 一筆花費長什麼樣}
CATEGORIES = {
    "ramen": "Food",
    "bento": "Food",
    "coffee": "Food",
    "metro": "Transport",
    "taxi": "Transport",
    "movie": "Entertainment",
    "book": "Learning",
    "course": "Learning",
}

DEFAULT_CATEGORY = "Other"


def guess_category(item):
    return CATEGORIES.get(item.strip().lower(), DEFAULT_CATEGORY)


class Expense:
    def __init__(self, date, item, category, amount, currency):
        self.date = date
        self.item = item
        self.category = category
        self.amount = amount
        self.currency = currency

    def amount_in_twd(self, rate):
        return self.amount * rate

    def __str__(self):
        return f"{self.date} {self.item} [{self.category}] {self.amount} {self.currency}"
\end{example}

這個檔案裡沒有一行是新的。\texttt{CATEGORIES} 是一張對照表，加上
\texttt{.get(key, 預設值)} 查表——你在第 \ref{sec:classify-function} 節用同一招做過
\texttt{EXTENSION\_MAP}，把「規則」放進 dict、把「怎麼套用規則」放進函式，之後要多支援
一個項目只要加一行，不用動函式本體。\texttt{item.strip().lower()} 是第 \ref{ch:strings}
章的兩個字串方法，讓使用者打成 \texttt{"Ramen "} 也查得到。

\texttt{class Expense} 是第 \ref{sec:class-init} 節的 \texttt{\_\_init\_\_} 加五個屬性、
第 \ref{sec:methods-and-self} 節的方法與 \texttt{self}、第 \ref{sec:str-method} 節的
\texttt{\_\_str\_\_}，一樣都沒有多。

值得停下來看一眼的只有 \texttt{amount\_in\_twd(self, rate)} 這一行的\textbf{參數}：它
收下一個已經查好的匯率，然後拿自己的金額乘上去。它\textbf{不會自己去查匯率}。這個看起來
很小的決定，會在第 \ref{sec:capstone-rates} 節跟第 \ref{sec:capstone-tests} 節各兌現
一次，現在先記著它。

\subsection*{接下來每一節，你會看到同一個檔名出現好幾次}

在寫下一支程式之前，先講一下這一章的閱讀規則，不然你等一下會困惑。

這個專案有六個 \texttt{.py} 檔案，其中五個（\texttt{expense.py}、\texttt{storage.py}、
\texttt{rates.py}、\texttt{report.py}、\texttt{check\_moneylog.py}）\textbf{各自只會在
自己那一節出現一次完整版，之後就不再改動}。寫完就定案，你可以放心把它們存起來。

只有 \texttt{main.py} 例外。它是整個工具的門面——「這個專案目前會做什麼」，就寫在這個
檔案裡，所以每加一塊新零件，它就會跟著長大一點。接下來每一節，這本書都會把
\texttt{main.py} \textbf{完整重列一次}，並在標題裡標明是第幾版。它一直很短，重列比要你
自己拼湊清楚得多。

現在是第 1 版。它造兩筆花費、印出來，其中一筆是日圓，用一個先寫死的假匯率換算成台幣：

\begin{example}{python}{(moneylog/main.py 第 1 版) 造兩筆花費看看}
from expense import Expense, guess_category

FAKE_JPY_RATE = 0.21

lunch = Expense("2026-09-01", "bento", guess_category("bento"), 100, "TWD")
ramen = Expense("2026-09-02", "ramen", guess_category("ramen"), 1200, "JPY")

print(lunch)
print(ramen)
print(f"{ramen.item} costs about {round(ramen.amount_in_twd(FAKE_JPY_RATE), 2)} TWD")
\end{example}

\iosample{}{2026-09-01 bento [Food] 100 TWD \\ 2026-09-02 ramen [Food] 1200 JPY \\ ramen costs about 252.0 TWD \\}

第一行 \texttt{from expense import Expense, guess\_category} 就是第
\ref{sec:write-your-own-module} 節那個動作：\texttt{expense.py} 跟 \texttt{main.py} 放在
同一個資料夾裡，所以 \texttt{import} 找得到它。前兩行輸出來自 \texttt{\_\_str\_\_}，
第三行是 \texttt{amount\_in\_twd()} 收下 \texttt{0.21} 之後算出來的。

那個 \texttt{FAKE\_JPY\_RATE = 0.21} 是暫時的，第 \ref{sec:capstone-rates} 節會換成真的
去查。這裡先寫死一個數字，是為了讓這一節能立刻跑起來——先讓程式動，再讓它正確，是這一章
會反覆用到的節奏。

現在專案長這樣：

\begin{verbatim}
moneylog/
|-- main.py            工具的門面（這一節改過）
|-- expense.py         一筆花費長什麼樣        <- 這一節新增
|-- requirements.txt   要裝哪些套件
|-- README.md          給人看的說明（還是空的）
`-- .gitignore         哪些東西不進版本控制
\end{verbatim}

\keypoint{同一批資料，在記憶體裡是一個 \texttt{Expense} 物件，在硬碟上會是 CSV 的一列
文字。這是兩種完全不同的形狀，中間需要有人負責翻譯——那就是下一節 \texttt{storage.py}
的全部工作。只要一個專案開始需要「跨執行記得事情」，就一定會出現這樣一層。}

\begin{drills}[label=drills:capstone-model]
\item 執行 \texttt{print(Expense("2026-09-05", "taxi", "Transport", 250, "TWD"))}
會印出什麼？負責決定這個格式的是 \texttt{expense.py} 裡的哪一個方法？
\item \texttt{guess\_category("hovercraft")} 會回傳什麼？是哪一行程式碼決定的？
\item 如果之後想多記一個「付款方式」欄位（現金／刷卡），\texttt{expense.py} 裡有哪幾個
地方要改？
\end{drills}

\section{把紀錄存下來、再讀回來：\texttt{storage.py}}
\label{sec:capstone-storage}

上一節那兩筆花費有一個致命的問題：程式一結束它們就消失了。規格裡「之後隨時可以問它」這
半句話，要求的正是「這次記的東西，下次還在」。這一節要做的就是那層翻譯：把
\texttt{Expense} 物件寫成 CSV 的一列，再把 CSV 的一列讀回成 \texttt{Expense} 物件。

\begin{example}{python}{(moneylog/storage.py) 存進 CSV、再讀回來}
import csv

from expense import Expense

FIELDS = ["date", "item", "category", "amount", "currency"]


def append_expense(path, expense):
    try:
        with open(path, encoding="utf-8") as f:
            file_is_empty = f.read() == ""
    except FileNotFoundError:
        file_is_empty = True

    with open(path, "a", newline="", encoding="utf-8") as f:
        writer = csv.DictWriter(f, fieldnames=FIELDS)
        if file_is_empty:
            writer.writeheader()
        writer.writerow({
            "date": expense.date,
            "item": expense.item,
            "category": expense.category,
            "amount": expense.amount,
            "currency": expense.currency,
        })


def load_expenses(path):
    with open(path, newline="", encoding="utf-8") as f:
        reader = csv.DictReader(f)
        for row in reader:
            yield Expense(
                row["date"],
                row["item"],
                row["category"],
                float(row["amount"]),
                row["currency"],
            )
\end{example}

三個地方值得停下來看。

\textbf{第一，開頭那段 \texttt{try}/\texttt{except} 在解決一個很具體的問題}：CSV 的第一列
是欄位名稱，只該寫一次。但「這個檔案存在嗎、裡面有東西嗎」這種問題，第
\ref{sec:lbyl-vs-eafp} 節講得很清楚：與其先檢查再動手，不如直接開開看，開不起來再說。
所以這裡直接 \texttt{open()}，接到 \texttt{FileNotFoundError} 就知道「還沒有這個檔案，
待會要寫標題列」。這就是 EAFP。

\textbf{第二，\texttt{append\_expense()} 用的是 \texttt{"a"} 模式}（第
\ref{sec:appending-files} 節）。這是整個專案「會累積」的關鍵：如果用 \texttt{"w"}，每記
一筆就會把前面全部清空。\texttt{csv.DictWriter} 加 \texttt{writeheader()} 加
\texttt{writerow()} 是第 \ref{sec:csv-dict} 節的寫法，餵給它的那個 dict，就是把
\texttt{Expense} 物件的五個屬性攤平——這就是「翻譯」這件事真正的樣子。

\textbf{第三，\texttt{load\_expenses()} 用的是 \texttt{yield}，不是 \texttt{return}}。
它每讀一列就吐出一個 \texttt{Expense}，是一台第 \ref{sec:yield-basics} 節那樣的發號機。
另外注意 \texttt{float(row["amount"])}：第 \ref{sec:csv-module-read} 節強調過「CSV 讀回來
的東西全部是字串」，不轉成數字，後面所有加總都會出事。

\stylepoint{\texttt{load\_expenses()} 大可以寫成「造一個空 list、把每一筆
\texttt{append} 進去、最後 \texttt{return}」，那樣也完全正確。用 \texttt{yield} 的理由是
第 \ref{sec:lazy-memory} 節那一條：記帳這件事只會愈記愈多，一份記了三年的
\texttt{expenses.csv} 沒有理由整份先攤在記憶體裡，才開始處理第一筆。這個選擇在檔案還只有
兩列的時候看不出差別，但它是「這個程式以後會長大」這個判斷的具體結果。}

現在 \texttt{main.py} 升到第 2 版：造兩筆、存進去、再把檔案裡的紀錄全部讀出來。既然讀
出來了，順便走第二趟，數數看總共有幾筆。

\begin{example}{python}{(moneylog/main.py 第 2 版) 存下來，再讀回來}
from expense import Expense, guess_category
from storage import append_expense, load_expenses

DATA_FILE = "expenses.csv"

lunch = Expense("2026-09-01", "bento", guess_category("bento"), 100, "TWD")
ramen = Expense("2026-09-02", "ramen", guess_category("ramen"), 1200, "JPY")

append_expense(DATA_FILE, lunch)
append_expense(DATA_FILE, ramen)

records = load_expenses(DATA_FILE)

print("Everything in the log:")
for expense in records:
    print(expense)

print("Walking through it again:")
for expense in records:
    print(expense)

print("(that's the end of the second walk)")
\end{example}

\iosample{}{Everything in the log: \\ 2026-09-01 bento [Food] 100.0 TWD \\ 2026-09-02 ramen [Food] 1200.0 JPY \\ Walking through it again: \\ (that's the end of the second walk) \\}

第一趟兩筆都在。第二趟\textbf{什麼都沒有}。

沒有 traceback、沒有紅字、沒有任何警告，第二個 \texttt{for} 迴圈確確實實執行了，只是跑了
零輪。這就是第 \ref{sec:one-shot} 節那個坑，在一個真正的專案裡長出來的樣子：
\texttt{load\_expenses()} 是一台發號機，號碼發完就沒了。它不是壞掉，它是用完了。

修法也是第 \ref{sec:one-shot} 節給過的：如果你打算走不只一次，就先用
\texttt{list()} 把它整個收下來。整份程式只要改一行：

\begin{example}{python}{(moneylog/main.py 第 2 版，修好之後) 先收成 list 再走}
from expense import Expense, guess_category
from storage import append_expense, load_expenses

DATA_FILE = "expenses.csv"

lunch = Expense("2026-09-01", "bento", guess_category("bento"), 100, "TWD")
ramen = Expense("2026-09-02", "ramen", guess_category("ramen"), 1200, "JPY")

append_expense(DATA_FILE, lunch)
append_expense(DATA_FILE, ramen)

records = list(load_expenses(DATA_FILE))

print("Everything in the log:")
for expense in records:
    print(expense)

print("Walking through it again:")
for expense in records:
    print(expense)

print("(that's the end of the second walk)")
\end{example}

\iosample{}{Everything in the log: \\ 2026-09-01 bento [Food] 100.0 TWD \\ 2026-09-02 ramen [Food] 1200.0 JPY \\ Walking through it again: \\ 2026-09-01 bento [Food] 100.0 TWD \\ 2026-09-02 ramen [Food] 1200.0 JPY \\ (that's the end of the second walk) \\}

順帶留意一個小細節：第 1 版印出來的金額是 \texttt{100}，這一版變成 \texttt{100.0}。因為
這一次的資料是從 CSV 讀回來的，中間經過了 \texttt{float()}。同一筆花費，繞了硬碟一圈
回來，型別變了——這正是「兩種形狀之間有一層翻譯」最好的證據。

打開 \texttt{expenses.csv}，會看到這個：

\begin{verbatim}
date,item,category,amount,currency
2026-09-01,bento,Food,100,TWD
2026-09-02,ramen,Food,1200,JPY
\end{verbatim}

還有一件事要提醒：這一版每執行一次就會多記兩筆。跑第二次，檔案裡就會有四列。這不是
bug，是 \texttt{"a"} 模式該有的行為；示範完之後把 \texttt{expenses.csv} 刪掉，再往下一節
走，你看到的數字才會跟書上一樣。

\begin{verbatim}
moneylog/
|-- main.py            工具的門面（這一節改過）
|-- expense.py         一筆花費長什麼樣
|-- storage.py         存進 CSV、讀回來        <- 這一節新增
|-- requirements.txt   要裝哪些套件
|-- README.md          給人看的說明（還是空的）
|-- .gitignore         哪些東西不進版本控制
`-- expenses.csv       你的紀錄（程式產生）    <- 這一節開始出現
\end{verbatim}

\keypoint{\texttt{Expense} 是資料在記憶體裡的形狀，CSV 的一列是它在硬碟上的形狀，
\texttt{storage.py} 是這兩種形狀之間唯一的翻譯層。這樣安排有一個具體好處：以後如果決定
改用別的儲存方式，要動的只有這一個檔案，\texttt{expense.py} 跟後面的
\texttt{report.py} 一行都不用改，因為它們從來就不知道資料是怎麼存的。}

\begin{drills}[label=drills:capstone-storage]
\item 如果把 \texttt{append\_expense()} 裡的 \texttt{"a"} 改成 \texttt{"w"}，執行第 2 版
的 \texttt{main.py} 之後，\texttt{expenses.csv} 裡會剩下幾列？
\item 如果把 \texttt{load\_expenses()} 裡的 \texttt{float(row["amount"])} 改成
\texttt{row["amount"]}，第 2 版的 \texttt{main.py} 印出來會有什麼不同？哪一節之後會真的
出事？
\item \texttt{load\_expenses()} 裡的 \texttt{yield} 如果改成 \texttt{return}（其他都不
動），第 2 版修好之後的那支 \texttt{main.py} 會怎樣？
\end{drills}

\section{跟外面的世界要匯率：\texttt{rates.py}}
\label{sec:capstone-rates}

到目前為止，那個 \texttt{0.21} 都是硬寫在程式裡的假數字。這一節要把它換成真的。

匯率是第 \ref{ch:api} 章開頭講的那種資料：這一秒才剛剛決定、存在別人家的伺服器上，不會
乖乖躺在你的電腦裡等你讀。所以要去問。

要問的骨架你已經完整學過了，第 \ref{sec:currency-api} 節的學習要點也講得很白：換一個
服務，套用同一套骨架，換掉細節就好。這一節就是那句話第一次被兌現。

\begin{code}{python}{去問一個匯率的通用骨架}
url = "the service URL"
params = {"base": base_currency}
response = requests.get(url, params=params, timeout=5)
response.raise_for_status()
data = response.json()
rate = data["rates"][target_currency]
\end{code}

第 \ref{ch:api} 章用的 Frankfurter 有一個限制：它的資料來源是歐洲央行，那份清單裡\textbf{
沒有台幣}。這個專案要把所有東西換成台幣，所以要換一個涵蓋台幣的服務。這裡用
\texttt{open.er-api.com}，一樣公開、一樣不用註冊、不用金鑰。看看程式碼就知道，換服務對
你的影響有多小：網址換掉、參數少一個，其餘每一行都跟第 \ref{sec:currency-api} 節一模
一樣。

\begin{example}{python}{(moneylog/rates.py) 去外面拿匯率，順便記起來}
import requests

URL = "https://open.er-api.com/v6/latest"


def fetch_rate(base, target):
    params = {"base": base}
    response = requests.get(URL, params=params, timeout=5)
    response.raise_for_status()
    data = response.json()
    return data["rates"][target]


class RateBook:
    def __init__(self, known_rates=None):
        if known_rates is None:
            self.rates = {}
        else:
            self.rates = known_rates

    def get(self, currency):
        if currency == "TWD":
            return 1.0
        if currency in self.rates:
            return self.rates[currency]
        print(f"(looking up the {currency} rate...)")
        rate = fetch_rate(currency, "TWD")
        self.rates[currency] = rate
        return rate
\end{example}

\texttt{fetch\_rate()} 就是那個骨架，沒有別的。\texttt{RateBook} 才是這一節的重點：它是
一本「匯率記事本」，裡面就一個 dict。\texttt{get()} 做三件事——台幣直接回 \texttt{1.0}、
記事本裡有就直接給、都沒有才真的上網問一次，問到之後\textbf{記進 dict}，下次就不用再問。

\texttt{known\_rates=None} 這個預設值不是隨便寫的。第 \ref{sec:mutable-default} 節講過，
預設值\textbf{不能}直接寫成 \texttt{\{\}}，那個坑會讓所有 \texttt{RateBook} 共用同一本
記事本。所以預設寫 \texttt{None}，進到 \texttt{\_\_init\_\_} 裡再造一個新的空 dict。
至於「為什麼要留這個參數，讓別人可以塞一份現成的匯率進來」——第 \ref{sec:capstone-tests}
節會揭曉，那是這個參數存在的唯一理由。

\subsection*{為什麼 \texttt{amount\_in\_twd()} 不自己去查匯率}

現在可以回答第 \ref{sec:capstone-model} 節埋的那個問題了。

看起來，讓 \texttt{Expense.amount\_in\_twd()} 自己去打一次 API 會方便得多——呼叫端連匯率
都不用管。但這麼寫會同時壞掉三件事：

\begin{enumerate}
\item \textbf{慢，而且失禮。} 一百筆日圓花費就會打一百次 API，問一百次同一個問題。你等得
很久，對方的伺服器也白白被你敲了九十九次。
\item \textbf{沒有必要。} 匯率一天只變一次，同一次執行裡查一次就夠了。「查過的答案先記
起來，下次直接用」是 dict 最經典的用途之一。
\item \textbf{最嚴重的一件：測試會跑不起來。} 「一千日圓乘上 0.21 等於多少」是一個純粹的
算術，答案永遠一樣；但如果這個方法裡藏著一次 API 呼叫，那它就變成「要有網路才能回答、
而且答案每天不一樣」。第 \ref{sec:capstone-tests} 節那支檢查腳本，會因此一個檢查都寫不
出來。
\end{enumerate}

所以這個專案把「去拿資料」跟「拿資料來算」切成兩層：\texttt{rates.py} 負責跟外面的世界
打交道，\texttt{expense.py} 負責算術，它們之間傳的是一個 \texttt{float}。

\keypoint{這是整章唯一一個「不是語法、是判斷」的教學點，也是 capstone 真正的難度所在：
\textbf{一個小決定會同時影響效能、可靠度、跟能不能被測試。} \texttt{amount\_in\_twd(rate)}
跟 \texttt{amount\_in\_twd()} 只差一個參數，語法上沒有誰對誰錯，但前者讓這個專案的三個
問題一次解決。你以後看別人的程式碼時，「這個函式要自己去拿它需要的東西，還是等別人餵給
它」會是一個一直反覆出現的問題。}

還有一個決定值得說一句：\texttt{RateBook.get()} 裡面\textbf{沒有} \texttt{try}/
\texttt{except}。網路掛掉的時候，它就讓例外往外丟。這跟第 \ref{sec:download-multiple} 節
那題練習的答案是同一個道理：接不接得住、接住之後要做什麼，呼叫的人比較清楚。一支只想印
報表的程式可能想印一句話就算了，另一支程式可能想換個備援來源再試一次。真正接住它的地方
在第 \ref{sec:capstone-wiring} 節的 \texttt{main.py}。

\texttt{main.py} 升到第 3 版：多記兩筆（一筆美金、一筆台幣），然後把檔案裡所有紀錄都
換算成台幣印出來。

\begin{example}{python}{(moneylog/main.py 第 3 版) 換成真的匯率}
from expense import Expense, guess_category
from rates import RateBook
from storage import append_expense, load_expenses

DATA_FILE = "expenses.csv"

new_expenses = [
    Expense("2026-09-03", "movie", guess_category("movie"), 12.99, "USD"),
    Expense("2026-09-03", "metro", guess_category("metro"), 30, "TWD"),
]
for expense in new_expenses:
    append_expense(DATA_FILE, expense)

rate_book = RateBook()

for expense in list(load_expenses(DATA_FILE)):
    rate = rate_book.get(expense.currency)
    print(f"{expense} -> {round(expense.amount_in_twd(rate), 2)} TWD")

print("Rates I ended up looking up:", rate_book.rates)
\end{example}

\iosample{}{2026-09-01 bento [Food] 100.0 TWD -> 100.0 TWD \\ (looking up the JPY rate...) \\ 2026-09-02 ramen [Food] 1200.0 JPY -> 237.44 TWD \\ (looking up the USD rate...) \\ 2026-09-03 movie [Entertainment] 12.99 USD -> 412.02 TWD \\ 2026-09-03 metro [Transport] 30.0 TWD -> 30.0 TWD \\ Rates I ended up looking up: \{'JPY': 0.197868, 'USD': 31.71834\} \\}

\textbf{你跑的時候，看到的數字一定跟這裡不一樣。} 上面這些是某一天真的跑出來的結果，
匯率每天都在動，這正是「這份資料存在別人家的伺服器上」的意思。看數字有沒有跟書上一致
沒有意義，該看的是形狀對不對。

有兩件事在這份輸出裡看得特別清楚。第一，台幣那兩筆完全沒有觸發任何一次查詢——
\texttt{get()} 第一行就把它們攔下來了。第二，\texttt{(looking up ...)} 這句話對 JPY 跟
USD 各只印了一次，就算後面還有更多同幣別的紀錄也不會再印，因為第二次開始都從 dict 裡拿。
最後那一行 \texttt{rate\_book.rates} 把記事本整個攤開來，裡面就只有兩筆。

\begin{verbatim}
moneylog/
|-- main.py            工具的門面（這一節改過）
|-- expense.py         一筆花費長什麼樣
|-- storage.py         存進 CSV、讀回來
|-- rates.py           去外面拿匯率            <- 這一節新增
|-- requirements.txt   要裝哪些套件
|-- README.md          給人看的說明（還是空的）
|-- .gitignore         哪些東西不進版本控制
`-- expenses.csv       你的紀錄（程式產生）
\end{verbatim}

\begin{drills}[label=drills:capstone-rates]
\item 如果把 \texttt{RateBook.get()} 裡的 \texttt{self.rates[currency] = rate} 這一行
刪掉（其他都不動），第 3 版的 \texttt{main.py} 執行起來會有什麼不同？
\item \texttt{RateBook.get()} 第一行為什麼要特別處理 \texttt{"TWD"}？如果拿掉這兩行，
會發生什麼事？
\item 執行第 3 版的時候把網路關掉，會看到什麼？是哪一行程式碼丟出例外的？（提示：
\texttt{rates.py} 裡沒有任何 \texttt{try}/\texttt{except}。）
\end{drills}

\section{把資料變成一份報表：\texttt{report.py}}
\label{sec:capstone-report}

規格的最後半句還沒做：「總共花了多少台幣、哪一類花最多」。零件都到齊了——紀錄讀得回來、
匯率查得到、每一筆都會把自己換算成台幣——剩下的只是把它們加起來、排個序、印成一張看得
懂的表。

\begin{example}{python}{(moneylog/report.py) 彙總、排版、輸出}
import csv


def summarize(expenses, rate_book):
    categories = sorted({expense.category for expense in expenses})
    return {
        category: sum(
            expense.amount_in_twd(rate_book.get(expense.currency))
            for expense in expenses
            if expense.category == category
        )
        for category in categories
    }


def format_table(summary):
    ranked = sorted(summary.items(), key=lambda item: item[1], reverse=True)
    lines = [f"{'Category':<16}{'TWD':>12}", "-" * 28]
    for category, total in ranked:
        lines.append(f"{category:<16}{total:>12,.2f}")
    lines.append("-" * 28)
    lines.append(f"{'Total':<16}{sum(summary.values()):>12,.2f}")
    return "\n".join(lines)


def write_summary_csv(path, summary):
    ranked = sorted(summary.items(), key=lambda item: item[1], reverse=True)
    with open(path, "w", newline="", encoding="utf-8") as f:
        writer = csv.writer(f)
        writer.writerow(["category", "total_twd"])
        for category, total in ranked:
            writer.writerow([category, round(total, 2)])
\end{example}

三個函式，主力是第 \ref{ch:advanced-functions}、\ref{ch:comprehensions} 兩章那些「更
Python 的寫法」，再加上第 \ref{ch:strings} 章的 f-string 跟第 \ref{ch:csv} 章的
\texttt{csv} 模組。

\texttt{summarize()} 最密，拆開來看只有兩層。裡面那個
\texttt{\{expense.category for expense in expenses\}} 是一個 set comprehension（第
\ref{sec:comp-dict-set} 節），把「總共有哪些分類」去重複整理出來，再用
\texttt{sorted()} 固定順序。外面那一整包是 dict comprehension，對每一個分類算一個總額，
而總額用的是 \texttt{sum(...)} 吃一個「還沒算好的東西」（第 \ref{sec:reducing-builtins}
節），中間那個 \texttt{if expense.category == category} 就是第 \ref{sec:comp-filter} 節
的篩選。

\texttt{format\_table()} 用的是第 \ref{sec:f-string-format} 節那些格式規格：
\texttt{:<16} 靠左補到 16 格、\texttt{:>12} 靠右補到 12 格、\texttt{,} 加千分位逗號、
\texttt{.2f} 固定兩位小數。四個符號拼起來就是一張對齊的表。排序用
\texttt{sorted(..., key=lambda item: item[1], reverse=True)}，也就是第 \ref{sec:lambda}
節的 \texttt{lambda} 加第 \ref{sec:key-function} 節的 \texttt{key=}：\texttt{.items()}
給出來的每一項是 \texttt{(分類, 金額)}，取 \texttt{item[1]} 就是照金額排，
\texttt{reverse=True} 讓花最多的排最上面。

\texttt{write\_summary\_csv()} 是第 \ref{sec:csv-module-write} 節的 \texttt{csv.writer}，
用 \texttt{"w"} 模式——這裡跟 \texttt{storage.py} 相反，報表本來就該是「最新的那一份」，
每次整份覆蓋才對。

\subsection*{同一件事，用你三章前的寫法會長什麼樣}

\texttt{summarize()} 那一大包看起來很陌生的話，把它換成第 \ref{ch:csv} 章那種寫法，就
會發現它其實在做很簡單的事：

\begin{code}{python}{同一個 summarize()，改用造空 dict、跑迴圈、逐一累加}
def summarize(expenses, rate_book):
    summary = {}
    for expense in expenses:
        rate = rate_book.get(expense.currency)
        twd = expense.amount_in_twd(rate)
        if expense.category in summary:
            summary[expense.category] = summary[expense.category] + twd
        else:
            summary[expense.category] = twd
    return summary
\end{code}

兩種寫法算出來的東西一模一樣，都對，換上去程式照跑。差別在讀的時候：迴圈版要一路讀到
那個 \texttt{if}/\texttt{else}，才拼得出「喔，原來是在按分類加總」；comprehension 版
第一眼就看得出「我要一個 dict，key 是分類、value 是那個分類的總和」。第
\ref{ch:comprehensions} 章開場說的「一行就說完我要什麼」，在這裡第一次有了實際的對照組。

\stylepoint{回頭翻一下你在第 \ref{ch:csv} 章寫過的那些程式——造一個空 list、跑一個迴圈、
逐一 \texttt{append}。那時候那樣寫是對的，因為你手上只有那些工具。現在你手上多了
comprehension、\texttt{sum}、\texttt{sorted(key=)}，同一件事寫出來就是不一樣的形狀了。
這不是誰比較高級，是同一個人多學了一整個區塊的寫法之後，講同一件事的方式變簡潔了。}

\texttt{main.py} 第 4 版，第一次印出完整報表：

\begin{example}{python}{(moneylog/main.py 第 4 版) 印出報表、順便存成 CSV}
from rates import RateBook
from report import format_table, summarize, write_summary_csv
from storage import load_expenses

DATA_FILE = "expenses.csv"
SUMMARY_FILE = "summary.csv"

expenses = list(load_expenses(DATA_FILE))
rate_book = RateBook()
summary = summarize(expenses, rate_book)

print(format_table(summary))

write_summary_csv(SUMMARY_FILE, summary)
print(f"The same numbers are now in {SUMMARY_FILE}:")

with open(SUMMARY_FILE, encoding="utf-8") as f:
    print(f.read())
\end{example}

這一次的輸出有一張對齊的表格，欄位寬度是靠空白撐出來的，所以這裡直接照終端機上看到的
樣子印出來：

\begin{verbatim}
(looking up the USD rate...)
(looking up the JPY rate...)
Category                 TWD
----------------------------
Entertainment         412.02
Food                  337.44
Transport              30.00
----------------------------
Total                 779.46
The same numbers are now in summary.csv:
category,total_twd
Entertainment,412.02
Food,337.44
Transport,30.0
\end{verbatim}

（匯率每天都在變，你跑出來的金額不會跟這裡一樣，看的是形狀對不對。）

\textbf{這個專案的功能到這裡就齊了。} 記得下來、存得住、讀得回、換算得了、彙總得出來、
印得出報表，規格裡的每一句話都有東西對應了。剩下兩節做的是另一種事：確認它沒壞，以及讓
別人也能用。

\begin{verbatim}
moneylog/
|-- main.py            工具的門面（這一節改過）
|-- expense.py         一筆花費長什麼樣
|-- storage.py         存進 CSV、讀回來
|-- rates.py           去外面拿匯率
|-- report.py          彙總、排版、輸出        <- 這一節新增
|-- requirements.txt   要裝哪些套件
|-- README.md          給人看的說明（還是空的）
|-- .gitignore         哪些東西不進版本控制
|-- expenses.csv       你的紀錄（程式產生）
`-- summary.csv        最近一次的報表（程式產生）<- 這一節開始出現
\end{verbatim}

\begin{drills}[label=drills:capstone-report]
\item 把 \texttt{format\_table()} 裡的 \texttt{reverse=True} 拿掉，印出來的表會有什麼
不同？
\item \texttt{summarize()} 收到一個空的 \texttt{expenses}（一筆紀錄都沒有）會回傳什麼？
接著 \texttt{format\_table()} 拿到那個東西，會印出什麼？
\item 要讓報表多一欄「這個分類佔總額的百分比」，要改的是 \texttt{report.py} 裡的哪一個
函式？需不需要動到 \texttt{summarize()}？
\end{drills}

\section{怎麼確認自己沒改壞：一支檢查腳本}
\label{sec:capstone-tests}

專案現在有五個 \texttt{.py} 檔案在互相呼叫。這帶來一個第 \ref{ch:automating} 章那種單檔
腳本不會有的煩惱：\textbf{你改 A 的時候，很難確定 B 有沒有跟著壞掉}。

改一行 \texttt{expense.py}，然後手動跑一次 \texttt{main.py}、新增一筆、出一份報表、用
眼睛看數字對不對——這件事你做第一次還可以，做第十次就會開始偷懶，偷懶第三次就會漏掉一個
錯誤。解法是把「用眼睛看數字對不對」這件事寫成程式，讓它自己看。

工具你已經有了：第 \ref{sec:assert} 節的 \texttt{assert}。它做的事情很單純——條件成立就
安靜通過，條件不成立就馬上停下來丟出 \texttt{AssertionError} 並印出你寫的訊息。

\begin{example}{python}{(moneylog/check\_moneylog.py) 一連串檢查點}
from expense import Expense, guess_category
from rates import RateBook
from report import summarize

assert guess_category("ramen") == "Food", "ramen should be Food"
assert guess_category("hovercraft") == "Other", "unknown items fall back"

ramen = Expense("2026-09-02", "ramen", "Food", 1000, "JPY")
twd = round(ramen.amount_in_twd(0.21), 2)
assert twd == 210.0, "1000 JPY at 0.21 should be 210 TWD"

rate_book = RateBook({"JPY": 0.21, "USD": 32.0})
assert rate_book.get("TWD") == 1.0, "TWD to TWD is always 1"
assert summarize([], rate_book) == {}, "no records, empty report"

expenses = [
    Expense("2026-09-01", "bento", "Food", 100, "TWD"),
    Expense("2026-09-02", "ramen", "Food", 1000, "JPY"),
    Expense("2026-09-02", "metro", "Transport", 30, "TWD"),
]
summary = summarize(expenses, rate_book)
assert round(summary["Food"], 2) == 310.0, "Food is 100 + 210 = 310"
assert round(summary["Transport"], 2) == 30.0, "Transport is 30"

print("All checks passed.")
\end{example}

\iosample{}{All checks passed. \\}

這支檔案跟其他五個一樣，就是一支普通的 \texttt{.py}，用 \texttt{python check\_moneylog.py}
直接執行。它沒有用任何測試工具，只有 \texttt{import}、\texttt{assert} 跟一個
\texttt{print}。

七個檢查大致分成三種。前兩個檢查\textbf{一般情況跟例外情況}：表裡有的項目分類要對，表裡
沒有的項目要落到預設分類。中間三個檢查\textbf{最基本的算術跟邊界情況}——特別是
\texttt{summarize([], rate\_book) == \{\}}，「一筆紀錄都沒有的時候會不會爆炸」是最容易
忘記、也最容易出事的情況。最後兩個檢查\textbf{整條線接起來之後對不對}：三筆花費、跨兩個
分類、其中一筆是外幣，加總出來的數字要正確。

每個比較金額的地方都先包了一層 \texttt{round(x, 2)}（第 \ref{sec:arithmetic} 節）。
處理錢的時候，先四捨五入到小數點後兩位再比較，是一個值得養成的習慣。

\subsection*{這支腳本為什麼不需要網路}

看一眼建立 \texttt{RateBook} 的那一行：\texttt{RateBook(\{"JPY": 0.21, "USD": 32.0\})}。

這就是第 \ref{sec:capstone-rates} 節那個 \texttt{known\_rates} 參數存在的唯一理由。這裡
直接塞一本已經填好的匯率記事本進去，\texttt{RateBook.get()} 查得到，就不會去打 API。
整支檢查腳本因此\textbf{完全不碰網路}，可以在飛機上跑、可以跑一百次、每次結果都一樣。

而這件事能成立，是因為第 \ref{sec:capstone-rates} 節做的那個決定：\texttt{amount\_in\_twd()}
收匯率，不自己查匯率。當初那一句「這個小決定等一下會很重要」，兌現的就是這一刻。

用到的東西只有第 \ref{sec:function-params} 節的預設參數跟第 \ref{ch:dict-set} 章的 dict，
沒有任何特殊工具。

\keypoint{\textbf{測試好不好寫，是設計好不好的一支溫度計。} 如果 \texttt{amount\_in\_twd()}
自己去打 API，這一節每一個 assert 都會需要網路才跑得動，而且答案會隨著匯率每天變，根本
寫不出「應該等於 210」這種檢查。反過來說也成立：\textbf{當你發現一個東西很難測，那通常
不是測試的問題，是那段程式的分工出了問題。} 這是一個很好用的直覺，值得從現在開始養。}

\subsection*{故意改壞一行，看它會不會叫}

寫完檢查腳本，一定要做一件事：\textbf{確認它真的會失敗}。一支從來沒有紅過的檢查腳本，
跟沒有檢查腳本是一樣的。

把 \texttt{expense.py} 裡的 \texttt{amount\_in\_twd()} 故意改成加法（這是一個非常真實的
手誤，\texttt{*} 跟 \texttt{+} 就在隔壁）：

\begin{verbatim}
    def amount_in_twd(self, rate):
        return self.amount + rate
\end{verbatim}

再跑一次 \texttt{python check\_moneylog.py}：

\begin{verbatim}
Traceback (most recent call last):
  File "check_moneylog.py", line 10, in <module>
    assert twd == 210.0, "1000 JPY at 0.21 should be 210 TWD"
           ^^^^^^^^^^^^
AssertionError: 1000 JPY at 0.21 should be 210 TWD
\end{verbatim}

跟第 \ref{sec:assert} 節看過的形狀完全一樣：哪一行、哪個條件、你自己寫的那句訊息。改壞的
是 \texttt{expense.py}，跳出來罵人的是 \texttt{check\_moneylog.py}——這正是它存在的意義。
把那一行改回乘法，再跑一次，應該又會看到 \texttt{All checks passed.}。

最後補兩件事。第一，第 \ref{sec:assert} 節提過 \texttt{assert} 的兩個限制（條件誤包成
tuple 會永遠成立、加上 \texttt{-O} 選項執行會被整個跳過），在這裡一樣適用。這也提醒你
\texttt{assert} 的定位：它是\textbf{寫給你自己看的檢查點}，不是用來擋使用者亂輸入的關卡
——那種事該用 \texttt{try}/\texttt{except}。

第二，業界不會只用 \texttt{assert} 做這件事，而是用 \texttt{pytest} 這類專門的測試工具，
第 \ref{sec:project-anatomy} 節那張樹狀圖裡 \texttt{tests/} 資料夾放的就是這種東西的
正式版。它們能做的事多得多——自動找出所有測試、一次跑完之後給你一份報告、失敗的時候告訴
你更多細節。這一章刻意只用 \texttt{assert}，因為它已經足以回答那個最重要的問題：\textbf{
我剛剛那個改動，有沒有把別的地方弄壞？}

\begin{drills}[label=drills:capstone-tests]
\item 有一個檢查是寫不進這支腳本的：「\texttt{fetch\_rate("JPY", "TWD")} 應該回傳
0.21」。為什麼不能寫？（兩個理由。）
\item 如果把 \texttt{RateBook(\{"JPY": 0.21, "USD": 32.0\})} 改成 \texttt{RateBook()}，
這支腳本執行起來會怎樣？
\item 第 \ref{sec:capstone-report} 節那個「造空 dict、跑迴圈、逐一累加」的
\texttt{summarize()} 版本，如果整個換上去，這支檢查腳本會不會有任何一個 assert 失敗？
這說明了什麼？
\end{drills}

\section{把六個檔案接成一支工具：\texttt{main.py} 與 README}
\label{sec:capstone-wiring}

到目前為止，\texttt{main.py} 每一版都是寫給\textbf{你自己}測試用的：資料寫死在程式碼裡，
想換一筆花費就得改程式。這一節要把它變成真正給\textbf{使用者}用的東西——啟動之後有一個
選單，可以一直記帳、隨時出報表，不用每次都重新執行一次程式。

這個形狀你在第 \ref{sec:weather-tool} 節看過：\texttt{while True} 包一個 \texttt{input()}，
用 \texttt{if}/\texttt{elif}/\texttt{else} 分派指令，\texttt{break} 離開、
\texttt{continue} 跳過這一輪。這裡沿用同一套。

\begin{example}{python}{(moneylog/main.py 最終版) 一支真的能用的工具}
import requests

from expense import Expense, guess_category
from rates import RateBook
from report import format_table, summarize, write_summary_csv
from storage import append_expense, load_expenses

DATA_FILE = "expenses.csv"
SUMMARY_FILE = "summary.csv"

rate_book = RateBook()

print("moneylog: a multi-currency expense log")

while True:
    command = input("What now? (add / report / quit): ")

    if command == "quit":
        print("Bye!")
        break

    if command == "add":
        date = input("Date (YYYY-MM-DD): ")
        item = input("Item: ")
        amount = float(input("Amount: "))
        currency = input("Currency (TWD / JPY / USD): ").upper()
        expense = Expense(date, item, guess_category(item), amount, currency)
        append_expense(DATA_FILE, expense)
        print(f"Saved: {expense}")

    elif command == "report":
        try:
            expenses = list(load_expenses(DATA_FILE))
        except FileNotFoundError:
            print("No records yet, use add to write one down first")
            continue
        try:
            summary = summarize(expenses, rate_book)
        except requests.exceptions.RequestException:
            print("Couldn't reach the rate server, check your connection and try again")
            continue
        print(format_table(summary))
        write_summary_csv(SUMMARY_FILE, summary)
        print(f"The same numbers are now in {SUMMARY_FILE}")

    else:
        print("Unrecognized command, please try again")
\end{example}

跑一次完整的操作：從一個乾淨的資料夾開始，先手殘按了 \texttt{report}（那時候還沒有任何
紀錄），再記兩筆（一筆台幣、一筆日圓），出報表，離開。這次直接照終端機上的樣子印出來，
你打的字跟程式印的字會混在同一行——那就是互動式程式實際看起來的樣子：

\begin{verbatim}
moneylog: a multi-currency expense log
What now? (add / report / quit): report
No records yet, use add to write one down first
What now? (add / report / quit): add
Date (YYYY-MM-DD): 2026-09-01
Item: bento
Amount: 100
Currency (TWD / JPY / USD): twd
Saved: 2026-09-01 bento [Food] 100.0 TWD
What now? (add / report / quit): add
Date (YYYY-MM-DD): 2026-09-02
Item: ramen
Amount: 1200
Currency (TWD / JPY / USD): jpy
Saved: 2026-09-02 ramen [Food] 1200.0 JPY
What now? (add / report / quit): report
(looking up the JPY rate...)
Category                 TWD
----------------------------
Food                  337.44
----------------------------
Total                 337.44
The same numbers are now in summary.csv
What now? (add / report / quit): quit
Bye!
\end{verbatim}

注意 \texttt{twd} 打的是小寫，紀錄裡存進去的卻是大寫 \texttt{TWD}——那是
\texttt{.upper()} 做的事，使用者不用管大小寫。

兩個 \texttt{try}/\texttt{except} 各接住一種「本來就可能發生」的狀況。
\texttt{FileNotFoundError} 是第一次使用的人一定會遇到的：還沒記過任何一筆就按了
\texttt{report}。\texttt{requests.exceptions.RequestException} 是第 \ref{ch:automating}
章跟第 \ref{sec:eafp-requests} 節提過的那個家族——連不上、逾時、伺服器回錯，全都是它底下
的成員，一次接住就好。把網路關掉，用一份已經有外幣紀錄的 \texttt{expenses.csv} 再按一次
\texttt{report}，看到的是這個：

\begin{verbatim}
moneylog: a multi-currency expense log
What now? (add / report / quit): report
(looking up the USD rate...)
Couldn't reach the rate server, check your connection and try again
What now? (add / report / quit): quit
Bye!
\end{verbatim}

程式沒有當掉，只印了一句話，然後回到選單等你下一個指令。這就是第
\ref{sec:capstone-rates} 節說的「讓例外往外丟、由呼叫端決定怎麼辦」——決定在這裡做出來
了，而它的內容是「印一句話，繼續」。

\keypoint{把 \texttt{main.py} 從頭讀到尾：它沒有一行在做計算。它只做三件事——問使用者
要什麼、呼叫對的函式、把結果印出來。這也是為什麼第 \ref{sec:capstone-tests} 節那支檢查
腳本完全不用測它：真正需要被檢查的邏輯，全都住在別的檔案裡。\textbf{把「跟人互動」跟
「做事」分開，是這個專案最重要的一條分界線}，也是你以後看任何一個專案時，最值得先問的
一個問題：這個檔案是門面，還是幹活的？}

\subsection*{一句你以後一定會看到的話}

真實專案的主程式最下面，常常會有一行 \texttt{if \_\_name\_\_ == "\_\_main\_\_":}，底下
縮排包著一整段程式碼。它處理的是「這個檔案是被直接執行，還是被別的檔案 \texttt{import}」
的差別。這本書用的是扁平的寫法——五個模組裡只有 \texttt{def} 跟常數，會動的程式碼全部
集中在 \texttt{main.py} 跟 \texttt{check\_moneylog.py} 的頂層——所以完全用不到它。以後在
別人的專案裡看到這一行，知道它在處理什麼問題就夠了，不用怕。

\subsection*{最後幾份檔案}

\texttt{requirements.txt} 跟 \texttt{.gitignore} 從第 \ref{sec:capstone-spec} 節寫下來
之後就沒有改過，那兩份就是定稿。剩下 README，這是唯一一份還空著的檔案，也是別人拿到這個
資料夾之後\textbf{第一個會打開的東西}：

\begin{code}{text}{(moneylog/README.md) 寫給人看的說明書}
# moneylog

一支記帳的命令列小工具。每一筆花費會被附加到 expenses.csv；
要看報表的時候，外幣的金額會用當天的匯率換算成台幣，
再依分類加總、印成一張表，同時輸出一份 summary.csv。

## 需要什麼

- Python 3
- requirements.txt 裡列的套件（只有 requests 一個）

安裝：

    pip install -r requirements.txt

## 怎麼跑

    python main.py

要執行的檔案只有 main.py 一個，其他 .py 都是它 import 進去的零件。
啟動之後有三個指令：

    add      記一筆花費
    report   印出報表，並寫出 summary.csv
    quit     離開

## 資料放在哪

- expenses.csv   你記過的每一筆花費，第一次 add 的時候自動建立
- summary.csv    最近一次的報表，每次出報表都會整份覆蓋

這兩個檔案都不進版本控制，原因見 .gitignore。

## 怎麼確認程式沒壞

    python check_moneylog.py

順利的話會印出 All checks passed.，而且完全不需要網路。
\end{code}

這份 README 只回答五個問題：這是什麼、要裝什麼、怎麼跑、資料在哪、怎麼確認沒壞。不需要
更多了。特別注意第二段裡「要執行的檔案只有 \texttt{main.py} 一個」這一句——一個資料夾裡
有六支 \texttt{.py}，拿到的人第一個困惑一定是「我該執行哪一支」，一句話就能省掉對方十分鐘。

專案完成的樣子：

\begin{verbatim}
moneylog/
|-- main.py              工具的門面 <- 唯一要執行的檔案
|-- expense.py           一筆花費長什麼樣
|-- storage.py           存進 CSV、讀回來
|-- rates.py             去外面拿匯率
|-- report.py            彙總、排版、輸出
|-- check_moneylog.py    確認自己沒改壞
|-- requirements.txt     要裝哪些套件
|-- README.md            給人看的說明          <- 這一節補完
|-- .gitignore           哪些東西不進版本控制
|-- expenses.csv         你的紀錄（程式產生）
`-- summary.csv          最近一次的報表（程式產生）
\end{verbatim}

其餘五個模組維持第 \ref{sec:capstone-model}～\ref{sec:capstone-tests} 節列出來的樣子，
從頭到尾沒有再改過一行。六支 \texttt{.py} 加起來大約 200 行。

\begin{drills}[label=drills:capstone-wiring]
\item 使用者在 \texttt{Amount:} 那一題輸入 \texttt{"一百"}，程式會怎樣？要接住它的話，
該在 \texttt{main.py} 的哪裡加什麼？
\item \texttt{rate\_book = RateBook()} 這一行寫在 \texttt{while} 迴圈\textbf{外面}。如果
把它搬到 \texttt{elif command == "report":} 裡面，程式還是會動，但會有什麼差別？
\end{drills}

\section{如果這份程式碼要交出去：對照第 \ref{ch:industry} 章走一次}
\label{sec:capstone-handoff}

程式寫完了，這一節不寫程式。

你在第 \ref{ch:industry} 章讀過一整章的詞：專案、虛擬環境、\texttt{requirements.txt}、
README、commit、branch、pull request、code review、CI。那一章說的是「業界大概是這樣運作
的」，讀起來多少有點像在看別人的世界。這一節要做的事很簡單：把那些詞，一個一個對回你
前面七節真的做過的事。

\subsection*{你剛剛做的每一件事，業界叫什麼}

\begin{center}
\begin{tabular}{@{}p{5.2cm}p{5.2cm}p{2.4cm}@{}}
\toprule
你剛剛做的事 & 那一章的說法 & 在哪一節 \\
\midrule
建一個資料夾，裡面好幾個檔案 & 一個「正常專案」 & 第 \ref{sec:project-anatomy} 節 \\
建虛擬環境、寫 \texttt{requirements.txt} & 虛擬環境與套件管理 & 第 \ref{sec:project-structure} 節 \\
寫 \texttt{.gitignore} & 不該被拍進快照裡的東西 & 第 \ref{sec:gitignore} 節 \\
寫 README & 寫給人看的說明書 & 第 \ref{sec:project-anatomy} 節 \\
每做完一節就是一個能跑的狀態 & commit & 第 \ref{sec:version-control-why} 節 \\
交出去之前先請人看過 & branch／pull request／code review & 第 \ref{sec:version-control} 節 \\
那支 \texttt{check\_moneylog.py} & CI 那一關自動跑的東西 & 第 \ref{sec:cicd} 節 \\
\bottomrule
\end{tabular}
\end{center}

第五列值得多說一句。第 \ref{sec:version-control-why} 節說 commit 是「存檔加上一句話說明
這次改了什麼」。這一章的八節結構，本身就是八次 commit 的節奏——每一節結束的時候，程式都
能跑，而且比上一節多做一件事，那正是「一個值得記住的狀態」的定義。如果你把這個專案放進
版本控制，最自然的切法就是照著節走：「加上 \texttt{Expense} 這個 class」「紀錄可以存進
CSV 了」「接上匯率查詢」⋯⋯每一句都講得出這次改了什麼。

（這一章從頭到尾沒有出現任何一行 Git 指令，跟第 \ref{ch:industry} 章一樣。那是另一件要
花時間單獨學的事，這裡要建立的只有節奏感。）

\subsection*{模擬一次 code review}

第 \ref{sec:version-control} 節說過，程式碼合併進主線之前，通常會先請同事看過一遍。假設
你把 \texttt{moneylog} 交出去了，同事大概會問這幾個問題。有趣的是，答案幾乎都已經在前面
幾節裡了：

\begin{itemize}
\item \textbf{「金鑰有沒有寫死在程式碼裡？」} 沒有——這個專案刻意挑了不需要金鑰的匯率
服務。萬一以後換成需要金鑰的服務，第 \ref{sec:api-keys} 節跟第 \ref{sec:gitignore} 節
講過該怎麼處理：不要寫進程式碼，而且那份放金鑰的檔案要進 \texttt{.gitignore}。
\item \textbf{「網路掛掉，整支程式會不會跟著掛？」} 不會，第 \ref{sec:capstone-wiring}
節的 \texttt{main.py} 接住了，只印一句話就回到選單。
\item \textbf{「這個函式在做幾件事？」} 每個都只做一件。最值得講的是第
\ref{sec:capstone-rates} 節那個分層：拿資料歸 \texttt{rates.py}，算資料歸
\texttt{expense.py}。
\item \textbf{「別人拿到這個資料夾，知道怎麼跑嗎？」} README 講了要裝什麼、執行哪一支、
資料存在哪；\texttt{requirements.txt} 讓對方一行指令就把環境裝好。
\item \textbf{「你怎麼知道自己沒改壞？」} \texttt{python check\_moneylog.py}，第
\ref{sec:capstone-tests} 節。而且它不需要網路，任何時候都能跑。
\item \textbf{「\texttt{expenses.csv} 該不該進版本控制？」} 不該。判準很清楚：版本控制
保存的是\textbf{程式碼}，不是\textbf{使用者的資料}。這個檔案的內容因人而異、每天都在變，
把它記進版本控制，除了讓紀錄變得又大又吵之外，還等於把某個人的消費紀錄公開出去。所以第
\ref{sec:capstone-spec} 節那份 \texttt{.gitignore} 裡的 \texttt{*.csv} 是刻意的，而且它
順便把 \texttt{summary.csv} 也一起擋掉了——那份是程式每次重跑就會重新產生的東西，同樣沒
有保存的必要。
\end{itemize}

\subsection*{這個專案還沒做的事}

一份誠實的交件，該講的不只是做了什麼，還有\textbf{刻意沒做什麼}。以下四件事都是有意識
留白的，不是忘了：

\begin{itemize}
\item \textbf{沒有分 \texttt{src/}／\texttt{tests/} 子資料夾}，理由見第
\ref{sec:capstone-spec} 節那個風格要點。
\item \textbf{沒有處理「同一筆花費被記兩次」}。連按兩次 \texttt{add} 輸入一樣的內容，
CSV 裡就會有兩列一模一樣的東西，程式不會攔你。
\item \textbf{報表沒有分月份、也沒有算佔比}。目前只能看「從開始記帳到現在的總和」。
\item \textbf{匯率查的是「現在」，不是「花費當天」}。一筆去年的日圓花費，會用今天的匯率
換算。對記帳來說這其實不太對，但要做對就得改成查歷史匯率，那是另一個規格了。
\end{itemize}

一份講得出自己缺什麼的專案，比一份假裝什麼都做到的專案可信得多。這四件事剛好變成這一章
的最後幾題練習。

\subsection*{你剛剛完成的事}

回頭看一下這一章走過的路：從一句抱怨開始，拆成六個檔案，一個一個做出來，中間做了一個
會同時影響效能、可靠度跟可測試性的判斷，寫了一支不需要網路的檢查腳本，最後接成一支真的
能用的工具，還附了一份說明書。整章沒有教任何新語法——所有零件都是你在前面的章節裡一塊
一塊學會的，這一章只是第一次把它們裝在同一台機器上。

這件事的意義不在於 \texttt{moneylog} 這個工具本身，而在於：下一次有人跟你描述一個需求，
你腦子裡會開始自動出現「這句話要程式做什麼、這件事該放進哪一個檔案」這種問題。那就是這
一章真正想留給你的東西。

下一章不教程式，它要回答的是「接下來呢」：這本書沒教到的東西還有哪些、你現在的程度可以
往哪幾個方向走、要怎麼繼續。

\begin{drills}[label=drills:capstone-handoff]
\item 開放式練習：幫報表加上「每個分類佔總額的百分比」。（提示：總額就是
\texttt{sum(summary.values())}，要改的只有 \texttt{report.py} 裡的
\texttt{format\_table()}，\texttt{summarize()} 不用動。）
\item 開放式練習：讓 \texttt{report} 可以只看某一個月的紀錄，例如只算 2026 年 9 月。
（提示：篩選就是一個 comprehension 的 \texttt{if}；每一筆的 \texttt{date} 是
\texttt{"2026-09-01"} 這種字串，第 \ref{ch:strings} 章有現成的方法可以判斷開頭。）
\item 幫你上面兩題新增的功能，各補一個 \texttt{assert} 進 \texttt{check\_moneylog.py}。
補完之後，故意把新功能改壞一行，確認那個 assert 真的會叫。
\end{drills}

\section*{練習解答}
\addcontentsline{toc}{section}{練習解答}

\paragraph{第~\ref{sec:capstone-spec}~節}
\begin{enumerate}
\item 加進 \texttt{report.py}。「每個月分別花了多少」是一個彙總的問題，跟「怎麼把資料
存進檔案、讀回來」無關——日期本來就已經記在每一筆紀錄裡了，\texttt{storage.py} 一行都
不用改。
\item 虛擬環境資料夾跟 \texttt{expenses.csv} 該進 \texttt{.gitignore}：前者可以靠
\texttt{requirements.txt} 重建，後者是使用者的資料不是程式碼。\texttt{main.py} 跟
\texttt{requirements.txt} 都是你親手寫的、需要被保存的東西，不排除。
\item 大概會變成「去外面拿匯率\textbf{而且}把資料彙總成報表」——需要用「而且」才講得完，
就是該拆成兩個檔案的訊號。實際的壞處是：以後如果想換一個匯率服務，你得在一個同時裝著
報表排版邏輯的檔案裡東翻西找。
\end{enumerate}

\paragraph{第~\ref{sec:capstone-model}~節}
\begin{enumerate}
\item 印出 \texttt{2026-09-05 taxi [Transport] 250 TWD}。決定這個格式的是
\texttt{\_\_str\_\_}。
\item 回傳 \texttt{"Other"}。因為 \texttt{"hovercraft"} 不在 \texttt{CATEGORIES} 這張表
裡，\texttt{CATEGORIES.get(..., DEFAULT\_CATEGORY)} 這一行就會給出預設值。
\item 三個地方：\texttt{\_\_init\_\_} 的參數列多一個、裡面多一行
\texttt{self.payment = payment}，以及（如果希望它被印出來）\texttt{\_\_str\_\_} 的
f-string。另外提醒，\texttt{storage.py} 的 \texttt{FIELDS} 跟兩個函式也要跟著改，否則
新欄位存不進 CSV——這就是「改一個地方，別的地方會不會壞」的實際樣子。
\end{enumerate}

\paragraph{第~\ref{sec:capstone-storage}~節}
\begin{enumerate}
\item 只剩一列，而且是 \texttt{2026-09-02,ramen,Food,1200,JPY} 這一列——連標題列都沒了。
\texttt{"w"} 每次開檔都會把內容整個清空，所以第二次呼叫把第一次寫的兩列（標題列加第一筆）
全部蓋掉；而這一次前面那段 EAFP 讀到的檔案「不是空的」，所以它認為標題列已經寫過了，
不會再寫一次。結果就是一份沒有標題列的 CSV，接著 \texttt{load\_expenses()} 會把僅存的那
一列當成標題列，一筆紀錄都讀不出來。一個字的差別，可以讓資料悄悄消失。
\item 印出來看不出差別，因為 \texttt{\_\_str\_\_} 只是把值放進 f-string，字串照樣印得
出來。真正出事的是第 \ref{sec:capstone-rates} 節：\texttt{amount\_in\_twd()} 會變成
「字串乘上一個小數」，直接丟出 \texttt{TypeError}。
\item 完全不會印出任何一筆紀錄。\texttt{return} 會在讀到第一列的時候就把整個函式結束掉，
只回傳一個 \texttt{Expense} 物件（不是一串），\texttt{list()} 拿到一個物件會再對它做
走訪、丟出 \texttt{TypeError}。
\end{enumerate}

\paragraph{第~\ref{sec:capstone-rates}~節}
\begin{enumerate}
\item 最後那行 \texttt{rate\_book.rates} 會印出一個空 dict——查到的匯率沒有被記下來。
不過在這一節這份只有四筆的紀錄裡，日圓跟美金剛好各只出現一次，所以
\texttt{(looking up ...)} 印的次數看起來沒變。等到同一個幣別有第二筆紀錄（例如第
\ref{sec:capstone-wiring} 節那份），你就會看到同一句話被印兩次，那時候快取的價值才會
現形。這也提醒一件事：效能問題常常在資料還很少的時候完全看不出來。
\item 因為「台幣換台幣」的答案一定是 1，根本不需要問任何人。這兩行省下的不只是一次網路
來回：沒有它們的話，一個從頭到尾只記台幣的使用者，也會被迫要有網路才出得了報表——而他
的報表裡根本沒有任何一筆需要換算。拿掉之後程式仍然算得對，只是變慢，而且在沒有網路的
時候會整個失敗。
\item 會看到一個 \texttt{requests.exceptions.ConnectionError} 的 traceback，程式直接
中斷。丟出例外的是 \texttt{fetch\_rate()} 裡的 \texttt{requests.get()}，而
\texttt{rates.py} 沒有接住它，所以它一路往外傳到 \texttt{main.py}——第
\ref{sec:capstone-wiring} 節才會真正接住。
\end{enumerate}

\paragraph{第~\ref{sec:capstone-report}~節}
\begin{enumerate}
\item 會變成金額由小到大排，花最少的排在最上面。
\item \texttt{summarize()} 會回傳一個空 dict \texttt{\{\}}（那正是檢查腳本裡那個 assert
在確認的事）。\texttt{format\_table()} 拿到空 dict，中間的 \texttt{for} 迴圈跑零輪，
所以只會印出標題列、兩條分隔線、跟一行 \texttt{Total} 加 \texttt{0.00}——不會出錯，因為
\texttt{sum(\{\}.values())} 是 0。
\item 改 \texttt{format\_table()} 就好。總額可以用
\texttt{sum(summary.values())} 當場算出來，每一列的佔比就是
\texttt{total / grand\_total}，\texttt{summarize()} 一行都不用動。
\end{enumerate}

\paragraph{第~\ref{sec:capstone-tests}~節}
\begin{enumerate}
\item 兩個理由：第一，它需要網路，飛機上或斷線時就跑不了；第二，匯率每天都在變，今天
寫 0.21 明天就不成立了，這個 assert 會在程式完全沒問題的情況下失敗——一個會亂叫的檢查
腳本，很快就會被大家忽略。
\item 會變成需要網路才能跑，而且會慢很多（每個沒查過的幣別都真的去打一次 API）。更糟的
是那個 \texttt{assert round(summary["Food"], 2) == 310.0} 幾乎一定會失敗，因為 310 這個
答案是用 0.21 這個假匯率算出來的。
\item 不會有任何一個 assert 失敗，因為兩個版本算出來的東西完全一樣。這說明檢查腳本檢查
的是\textbf{行為}，不是\textbf{寫法}——只要對外的結果沒變，你可以放心把裡面重寫成任何
形狀，這正是有檢查腳本之後最大的好處。
\end{enumerate}

\paragraph{第~\ref{sec:capstone-wiring}~節}
\begin{enumerate}
\item \texttt{float("一百")} 會丟出 \texttt{ValueError}，整支程式直接中斷，前面記的紀錄
雖然已經存進 CSV 了，但使用者得重新執行一次。修法是把
\texttt{amount = float(input("Amount: "))} 包進一個 \texttt{try}/\texttt{except
ValueError}，接住之後印一句話再 \texttt{continue}，回到選單。
\item 差別在快取的壽命。寫在迴圈外面，這本匯率記事本活過整支程式，你連出五次報表也只會
查一次匯率；搬進 \texttt{elif} 裡面，每按一次 \texttt{report} 就會造一本全新的空記事本，
每次都得重查。程式仍然正確，只是白白多打了好幾次 API。
\end{enumerate}

\paragraph{第~\ref{sec:capstone-handoff}~節}
\begin{enumerate}
\item 開放式練習。方向：在 \texttt{format\_table()} 裡先算
\texttt{grand\_total = sum(summary.values())}，再把每一列的 f-string 多加一個欄位，
例如 \texttt{f"\{total / grand\_total:>8.1\%\}"}；記得處理 \texttt{grand\_total} 是 0
的情況，不然會除以零。
\item 開放式練習。方向：在 \texttt{summarize()} 之前先篩一次，
\texttt{expenses = [e for e in expenses if e.date.startswith("2026-09")]}，其餘完全不用
改。這一題的重點是體會到：因為每一層各做一件事，加一個篩選只要動一個地方。
\item 開放式練習。方向：百分比那題可以檢查「三筆假資料裡 Food 佔多少」；月份那題可以
準備兩筆不同月份的假資料，確認篩完之後只剩該有的那幾筆。重點是最後那個動作——把新功能
故意改壞一行、確認 assert 真的會叫，一支從來沒紅過的檢查腳本等於沒有。
\end{enumerate}

\chapter{後續學習方向}
\label{ch:where-next}

第 \ref{ch:introduction} 章結尾，你剛寫完人生第一支跑得動的程式，那時候書裡先打了個預防針：接下來會用
兩個方式反覆幫你建立 Python 的直覺，也會帶你去兩個地方，一個是讓你自己動手解決生活
裡的瑣事，一個是讓你看懂軟體這個世界怎麼運作。二十六章之後，這一章要做的事情，是回頭
兌現那個預告。

這一章分幾段推進：

\begin{itemize}
    \item 先回頭看，把你從一支只會印出一行字的程式，走到現在能處理資料檔案、串接真正的
    網路服務、交出一個有好幾個檔案的專案，這一路上兩個心智模型到底怎麼一步步落地，
    一件一件指給你看。
    \item 接著把兩個目標拆開來看，具體點出你在每個目標底下各自練成了哪些看得見、用得上
    的技能。
    \item 再把視角拉高一層，談真正能帶著走的東西，是幾個學習的習慣，讀錯誤訊息、查文件、
    遇到不會的東西先動手試一次，這些習慣以後遇到任何新語言、新工具都用得上。
    \item 然後把地圖攤開，讓你知道這本書刻意沒有走進去的幾片區域長什麼樣子，更深的
    Python 語言本身、更深的資料工具，還有幾種你以後可以自己找到的學習資源。
    \item 最後，依照你自己的背景，給一個具體一點的下一步建議。
\end{itemize}

這一章不會有新的片語、新的範例要你動手照著打，它的工作是幫你把已經學會的東西看清楚，
然後好好地把書闔上。

\section{回頭看：這趟旅程，你其實走了多遠}
\label{sec:journey-recap}

翻到這本書的最後一章之前，你大概不會沒事回頭翻前面二十六章，但現在正是這麼做的時候，
目的很單純：讓你老實對自己說一句，你真的走了這麼遠。這趟旅程從一行只會印出「哈囉，
世界！」的程式開始（\ref{ch:introduction} 章）；沒多久，你寫的程式第一次能自己做決定，依照某個條件走
哪一條分岔路（\ref{ch:conditionals} 章條件判斷）；再過不久，它學會了重複做一件事，不用你自己一行一行
複製貼上（\ref{ch:loops} 章迴圈）。中段，你的程式第一次打開一份真正的資料檔案，把裡面雜亂的內容
整理成看得懂的樣子，處理的對象從你自己手打的幾個數字，變成別人留下來、格式不一定
乾淨的原始資料（\ref{ch:csv} 章 CSV）。

再往後，你的程式碼本身開始變了樣子：四行迴圈加 append 縮成一行
comprehension（\ref{ch:comprehensions} 章）、一個函式學會「要一個才給一個」而不是先把
全部算完（\ref{ch:iterators-generators} 章）、一個函式敢在自己的主體裡呼叫自己
（\ref{ch:recursion} 章）。這些不是新功能，是同一件事寫得更像 Python 的樣子。

最後，它甚至跨出了你的電腦，主動去問另一個系統要資料，
再把回應接回自己的程式裡（\ref{ch:api} 章 API 串接）；而在 \ref{ch:capstone} 章，你把這一路
學到的東西全部放進同一個資料夾，交出了一個有好幾個檔案、有 \texttt{README}、有檢查腳本的
專案——那已經不是「一支程式」，是一個小小的軟體。這幾個里程碑，你是自己一行一行敲出來、
看著螢幕真的照你的話動起來的，這一段記憶，比任何一條語法規則都更值得留下來。

把這些具體的記憶收好之後，值得回頭看看，這一路你反覆撞見的到底是什麼。這本書一開始
就打過預防針：接下來會用兩個方式反覆幫你建立 Python 的直覺，現在是把這兩條線攤開
來看的時候。第一個是「名字牌模型」：\ref{ch:variables} 章第一次完整講給你聽，指定的意思是把一張名字
牌貼到一個物件上，跟把值複製進一個盒子完全是兩回事；之後你在 \ref{ch:lists} 章看 list 被兩個名字
同時指著、改一個兩個都動的時候撞見它一次，在 \ref{ch:functions} 章看函式怎麼「借」物件給參數的時候又
撞見一次，在 \ref{ch:tuples} 章看 tuple 為什麼不能被原地修改、名字牌模型跟盒子模型在這裡難得一致
的時候再一次，最後在 \ref{ch:advanced-functions} 章，你看見一個函式被回傳出來之後，居然還
記得出生時外層那張名字牌。

第二個是 EAFP：\ref{ch:files-exceptions} 章第一次完整鋪陳，遇到可能失敗的操作，
Python 比較習慣先做、錯了再處理，跟先檢查再動手的直覺剛好相反；你在 \ref{ch:dict-set} 章查字典之前
先埋過伏筆，在 \ref{ch:debugging} 章讀錯誤訊息、真正用 \texttt{try}/\texttt{except} 接住例外的時候，
親手用過它。

這兩條線在 \ref{ch:pythonic} 章已經合流過一次了——那一章把「以為 \texttt{b = a} 會複製」
跟「以為動手前應該先檢查」擺在一起，指出它們其實是同一個直覺，而且都猜錯了同一件事。
所以這裡不再重講一遍。 這裡只補一句那時候還不能說的話：那一章結束時，這兩個模型從
「用來解釋怪事的圖」變成了「你打字時會用到的規則」，而現在，隔了好幾章之後，你大概
已經不太會意識到自己在用它們了——這正是一個心智模型該有的下場。

如果你想幫這段回顧留下一個具體的動作，不妨真的翻回 \ref{ch:variables} 章你寫的第一個關於變數的小範例，
想想現在的你會怎麼重新寫一次它，會不會多包一個函式，讓它更整齊一點？或者翻回 \ref{ch:files-exceptions} 章
第一次寫的 \texttt{try}/\texttt{except}，想想現在的你，是不是已經很自然地先想到用它，
而不會先寫一堆 \texttt{if} 檢查？答案不用寫下來，你自己心裡有數就夠了，這正是這一路
上真正留下來的東西。

\section{兩個目標，你各自帶走了什麼技能}
\label{sec:goals-recap}

這本書一開始講過兩個並重的目標，這裡把它們拆開來看，具體點出你在每個目標底下各自練
成了哪些技能，一件一件指給你看，不停在「你進步了」這種空泛的講法上。

第一個目標，讓你能自己動手解決生活裡的瑣事。回頭數數看，你其實已經湊齊了一整套處理
資料的基本功：

\begin{itemize}
    \item \ref{ch:strings} 章教你用字串方法把一段雜亂的文字清乾淨。
    \item \ref{ch:lists} 章跟 \ref{ch:dict-set} 章教你用 list 跟 dict 把一批資料組織起來，\ref{ch:functions} 章教你把重複用得到的
    邏輯包成函式，不用每次都重寫一遍。
    \item \ref{ch:files-exceptions} 章跟 \ref{ch:csv} 章把這些串起來，讓你真的打得開一份檔案、讀得懂一份 CSV，處理完
    再存回去。
    \item \ref{ch:data-tools} 章讓你認識 numpy、pandas、matplotlib 這幾個常見的資料工具長什麼樣子，
    就算現在還不熟練，至少以後看到別人用這幾個名字，你知道它們在做什麼。
    \item \ref{ch:debugging} 章給你一套讀錯誤訊息、抓蟲的方法，讓「程式不動」這件事從一場災難變成
    一個可以按部就班解決的小問題。
    \item \ref{ch:regex} 章的正規表達式，讓你多一種在一大段文字裡揪出符合特定樣式內容的能力。
    \item \ref{ch:advanced-functions} 章到 \ref{ch:pythonic} 章這五章，讓同樣一件事開始寫得
    像 Python：把函式當成資料傳來傳去、用一行 comprehension 取代四行迴圈、讓資料「要一個
    才給一個」、遇到層層包住的東西時用遞迴拆開，以及一整套「別人一看就懂」的寫法慣例。
    這幾章沒有讓你多會做一件新的事，但讓你寫出來的東西，別人讀得下去。
    \item \ref{ch:automating} 章把這些技巧全部兜在一起，做出一個真正能跑的小專案，讓你體驗過一次
    「從一個模糊的需求，走到一支真的能用的程式」完整是什麼感覺，不再只是零散的片段
    練習。
\end{itemize}

這一路走下來，下次生活裡出現一件瑣碎、重複、要處理一堆資料的事情，你的第一反應會是
「這個我可以寫程式做」，而且真的知道從哪裡下手。

第二個目標，讓你看懂「軟體」這個世界怎麼運作。

\begin{itemize}
    \item \ref{ch:objects} 章跟 \ref{ch:inheritance} 章讓你知道物件、類別、繼承這些詞在講什麼，這是理解大部分現成
    程式庫、框架內部長相的基礎，就算你以後不自己寫大型系統，看得懂這套語彙也有幫助。
    \item \ref{ch:industry} 章一口氣把版本控制、前端、後端、資料庫、staging、production、backlog、
    sprint 這些工程師團隊平常掛在嘴邊的詞攤開講清楚，第一次讓你聽懂了他們平常在講的
    那些詞。
    \item \ref{ch:api} 章讓你親手讓自己的程式跟另一個系統打了一次交道，體會到「軟體跟軟體對話」
    實際上長什麼樣子，這件事光用聽的很難有感覺，只有自己打過一次 API 才會真的懂。
    \item \ref{ch:capstone} 章讓你把 \ref{ch:industry} 章那些只是「聽懂」的詞，親手做了一次：
    多個檔案怎麼分、\texttt{requirements.txt} 寫什麼、\texttt{README} 該交代什麼、
    改壞了要怎麼發現。 那一章交出來的東西，形狀已經跟工程師交出來的東西是同一種了。
\end{itemize}

這兩個目標從頭到尾沒有誰讓位給誰，你這一路走來，兩邊都真的碰過一次，兩個
都是走到今天、真正握在手上的東西。

\section{把方法帶走：比語法更重要的，是幾個習慣}
\label{sec:habits-to-keep}

技巧會忘、語法會查，但有幾件事，你已經在不知不覺中練成了習慣，而且以後不管換到哪
一種語言、哪一套工具，都還是用得上。

\textbf{讀錯誤訊息的習慣。} \ref{ch:introduction} 章你第一次看到一片紅字的時候，這本書就跟你說過：由下往
上讀，先看最後一行的錯誤類型跟訊息，那才是真正的線索；\ref{ch:debugging} 章你又把這套習慣練得更熟。
這件事情不限於 Python，以後不管你換到哪一種語言、用哪一套工具，錯誤訊息的樣子會不
一樣，但「看到紅字先別慌，由下往上找線索」這套習慣，你已經帶在身上了，搬到哪裡都用
得上。

\textbf{查文件的習慣。} \ref{ch:modules} 章這本書刻意沒有要你把每個標準函式庫的函式都背下來，改成
帶你練習「查文件」這件事本身，因為沒有人真的把整個 Python 標準函式庫背在腦子裡，連
寫了很多年程式的人也不例外。以後遇到書裡沒教過的模組、沒教過的套件，官方文件或套件
自己附的說明，永遠是你第一個該去的地方，不用繼續憑猜。

\textbf{遇到不會的東西先動手試一次的習慣。} \ref{ch:introduction} 章教你用互動模式「隨手驗算」的時候，這
本書就已經在示範這個習慣了，不確定某段語法會怎麼動，與其在腦子裡空想，不如直接打開
來試一次，答案馬上就出來了。

這件事其實也是 EAFP 精神的另一種面貌，EAFP 教你的是
「先做、錯了再處理」的程式寫法，但拉高一層來看，它其實也是一種面對陌生事物的態度：
遇到沒學過的東西，先動手碰一次，錯了再回頭修正，往往比坐在那裡把所有可能情況都想過
一遍更快找到答案。這裡是把 EAFP 這個詞往外延伸成一種學習態度，用同一種精神套用在
「學新東西」這件更大的事情上，一樣說得通，它原本在程式碼裡的意思還是那個意思。

\textbf{把大問題拆成小步驟的習慣。} \ref{ch:automating} 章那個整合專案，如果從頭到尾當成一件事來想，
很容易讓人不知道從何下手；但只要把它拆成幾個小任務，一個一個各個擊破，最後接起來就
是完整的成品。

這件事其實從第一次寫函式（\ref{ch:functions} 章）就開始練了，函式本身就是「把一段完整
邏輯包成一個獨立、可以單獨測試的小塊」的做法。以後遇到任何看起來很龐大的專案或任務，
先問自己「這件事可以拆成哪幾個更小的步驟」，往往比想著「怎麼一次寫出整個解法」更容易
真正開始動手。

\textbf{以後如果要用 AI 工具輔助寫程式或除錯。} 這本書沒有教你怎麼用 ChatGPT 之類的
工具生成程式碼，但值得老實說一句：這幾個你已經練出來的習慣，讀錯誤訊息、查文件、
動手先試、把問題拆小，以後一樣用得上，甚至是你判斷「AI 給的答案到底對不對」的基礎。

AI 工具能幫你很快生出一段程式碼，但那段程式碼有沒有真的做了你要的事，有沒有藏著一個
你沒發現的錯誤，還是得靠你自己讀得懂、查得到、試得出來。這本書真正想讓你帶走的，是
有能力判斷 AI 給的東西靠不靠得住，不是把以後所有的事都交給它。

\section{地圖上還有更大的世界}
\label{sec:beyond-this-book}

前面幾節整理的是你已經走過的路；這一節要做的事情不一樣，把地圖攤開，指給你看這本書
刻意沒有走進去的幾片區域長什麼樣子。目的不是要你現在就衝進去，只是讓你知道它們在那裡，
以後想去的時候認得路。

\textbf{更深的 Python 語言本身。} 這本書教的是 Python 能做到的一部分，全貌還大得多。
有幾件事，你在這本書裡碰到的是入門的那一層，往下還有好幾層；也有幾個詞，你完全沒碰過，
但以後跟程式圈的人聊起 Python，大概率會聽到：

\begin{itemize}
    \item \texttt{async}／\texttt{await}，讓程式在等待網路回應這種「發呆時間」裡，先去
    做別的事，不用乾等。第 \ref{ch:api} 章你打 API 的時候，如果一次要打很多個網址，
    就會用得上這個。
    \item \textbf{自己寫一個 context manager}。第 \ref{ch:pythonic} 章你已經知道
    \texttt{with} 在解決什麼問題了，但那一章只教你怎麼用別人做好的。等你學會怎麼替自己的
    東西加上「進去做什麼、出來做什麼」這一層，很多重複的收尾動作就可以收起來。
    \item \textbf{decorator 更完整的樣子}。第 \ref{ch:advanced-functions} 章教的是最
    基本的形狀，真正的專案裡你會看到帶參數的 decorator、疊好幾層的 decorator，以及
    \texttt{functools} 裡幫你把這些寫對的工具。
    \item \textbf{型別檢查工具}。第 \ref{ch:advanced-functions} 章說過，型別標註不會
    影響程式執行，Python 不會幫你檢查——但有另外一類工具會。在多人維護的專案裡，它們
    能在程式跑起來之前就抓出「這裡傳錯東西了」。
    \item \textbf{測試框架}。第 \ref{ch:capstone} 章你用幾個 \texttt{assert} 寫過一支
    檢查腳本，那是測試最素樸的樣子。真正的專案會用專門的測試工具，可以一次跑幾百個檢查、
    告訴你哪一個掛了、還能算出你的程式碼有多少比例被測到。
    \item \textbf{把程式碼分裝成套件}。第 \ref{ch:modules} 章教的自訂模組，前提是檔案跟
    呼叫它的程式放在同一個資料夾。專案再大一點就會需要子資料夾、需要讓別人可以
    \texttt{pip install} 你寫的東西，那是另一套規則。
\end{itemize}

這幾個詞，這裡只給你一句話的印象，不寫成可以動手照打的片語
或範例，這是這一節跟這本書其他每一節最大的不同。

這些詞具體會落在哪裡，可以看一個實際
的例子：像《Fluent Python》這樣的書，會把 iterator 與 generator、context managers、
\texttt{async}／\texttt{await} 這些放進同一個部分，書裡稱為「Control Flow」，用完整
一個部分的份量系統性地教——你在第 \ref{ch:iterators-generators} 章學到的 \texttt{yield}，
在那本書裡只是那一整個部分的開頭。那本書的作者甚至在一開頭就老實提醒讀者：「如果你才剛開始
學 Python，這本書會很難跟上，甚至可能讓你誤以為每一支 Python 程式都得用上特殊方法跟
metaprogramming 技巧。」這句話送給正在看這本書最後一章的你剛好合適：那本書現在還不
適合你，但它在地圖上，以後你想往「更深的 Python 語言本身」走的時候，認得路就好。

\textbf{更嚴謹的開發習慣。} 這本書裡的每支程式，寫完能跑、輸出結果對，任務就算完成。
但一支程式如果要讓很多人一起維護、要長期跑在正式環境裡，光是「能跑」還不夠。
\texttt{testing}（測試），是指另外寫一段程式碼，專門用來自動檢查你原本那段程式碼
在各種情況下是不是都給出正確的結果，這樣以後改動程式碼時，只要重新跑一次測試，就能
立刻知道有沒有不小心弄壞了原本能動的部分，不用每次都靠人工從頭試一遍。

前面提到的
\texttt{type hints}，也常常被歸在這一類「讓程式碼更可靠」的做法裡。這兩件事這本書
完全沒有觸碰，因為它們的價值要在程式碼變大、變得需要跟別人協作之後才真正顯現出來，
現在的你還不需要，但以後往這個方向走的時候，知道有這兩個詞可以查，會省下不少繞路
的時間。

\textbf{更大的應用範圍。} \ref{ch:api} 章你打 API 的時候，是你的程式主動去問別人的伺服器要
資料，扮演「使用者」的角色。

如果反過來，你想讓自己寫的程式變成別人可以在瀏覽器裡
打開、或是別的程式可以打過來要資料的服務，你會用到「網頁框架」這一類工具，它們幫你
把接收請求、整理回應這些底層雜事先打包好，讓你可以專心寫「收到請求之後要做什麼」這
段邏輯。這一塊跟 \ref{ch:industry} 章講的前端、後端、API 這些詞是同一條線上的東西，只是這本書只帶
你從「使用者」的角度碰過一次，還沒帶你走到「自己就是提供服務的那一方」這一步。

\textbf{更深的資料工具。} \ref{ch:data-tools} 章你已經跟 pandas 打過照面，用它讀 CSV、選欄、篩出符合
條件的資料列。但那裡刻意留了兩個詞沒有教：

\begin{itemize}
    \item \texttt{groupby} 做的事情是「把資料依某個欄位分組，再對每組分別算統計量」，
    例如「這份氣溫資料裡，每個城市個別的平均溫度是多少」；\ref{ch:data-tools} 章你如果想做這件事，得
    自己寫迴圈或條件判斷一步步湊出來，\texttt{groupby} 能把整段邏輯壓縮成一行。
    \item \texttt{merge} 做的是「把兩份資料依共同的欄位接在一起」，例如你手上有一份
    訂單資料、一份客戶資料，兩邊都有「客戶編號」這個欄位，\texttt{merge} 能幫你把兩份
    資料合成一份，不用自己一筆一筆對照著找。
\end{itemize}

這兩個詞，這裡一樣只給你一句話的印象，不會像前面章節那樣寫成完整的片語跟
範例。

\section{持續學習的資源}
\label{sec:learning-resources}

如果你想繼續往下走，不管是想多學一點寫給自己用的小工具，還是真的想認真往這個方向
發展，你不用一個人埋頭摸索。市面上跟網路上能找到的資源大致可以分成幾類，各有各的
用途，挑適合自己現在階段的就好，不用每一種都追。

第一類是書。

\begin{itemize}
    \item 有的書走「系統性參考」路線，把一個主題的來龍去脈講得很完整，適合你已經
    對某個方向有點興趣、想深入紮根的時候翻。
    \item 有的書走「動手做專案」路線，帶你一步步做出幾個完整的小成品，適合你想累積
    實作經驗、把書上的東西真正變成自己能重現的能力的時候用。
    \item 還有一種比較像速查手冊，平常不會從頭到尾讀，但寫程式寫到一半忘記某個語法時，
    翻兩下就能找到答案，適合放在手邊當工具書用。
\end{itemize}

第二類是線上課程。

\begin{itemize}
    \item 有的是一支支影片，老師帶著你邊講邊打字，適合喜歡「有人在旁邊帶」的學習方式。
    \item 有的是互動式的練習平台，網頁上直接寫程式碼、馬上看到對不對，適合想邊學邊
    動手、不想只是看的人。
    \item 還有一種比較像結構完整的大學課程搬到線上，會有作業、有進度安排，適合想要
    有人幫你排好學習路徑、按表操課的人。
\end{itemize}

這幾種形式教的內容常常重疊，選哪
一種主要看你自己比較適應哪種學習節奏，不是哪一種天生比較高級。

第三類是練習網站。

\begin{itemize}
    \item 有的是一題一題的程式解題平台，出一個小問題、你寫程式解掉它，適合拿來磨
    基本功、練習「看到問題馬上想到怎麼用程式解」的反射。
    \item 有的偏向真實世界的小專案練習，給你一個比較貼近實際工作場景的任務，讓你
    練習把好幾個技巧兜在一起用。
    \item 還有一種是參與別人已經寫好、公開讓大家一起維護的專案，從讀懂別人的程式碼
    開始，慢慢學會怎麼跟其他人一起改動同一份程式碼，這也是很多工程師真正精進技巧的
    地方。
\end{itemize}

這一節刻意沒有點名任何一本書或任何一個網站，因為書會改版、網站會關站，但上面這幾類
資源的分工不太會變。 如果你現在就想要具體的書名，書末的〈參考資料〉列了這本書寫作時
實際參考的七本 Python 入門書，每一本都註明了適合在什麼時候讀，以及這本書的哪幾章借用
過它。

除了這些，Python 官方文件本身就寫得很完整，大部分你以後會用到的模組、套件，第一手
的說明都在那裡，\ref{ch:modules} 章已經帶你練過怎麼讀這種文件。

如果你住的城市附近剛好有 Python
的使用者社群，他們通常對新手很友善，不用怕自己「還不夠格」加入；找不到實體社群的話，
線上的討論區、留言板一樣有很多人在，遇到問題丟出去問，通常都會有人願意回你。這幾類
資源不用照順序一個一個走完，遇到卡關的地方，哪一種能幫上忙就先用哪一種。

\section{接下來，你是誰，就從哪裡先走}
\label{sec:next-steps-by-reader}

前面幾節給的建議比較通用，這一節具體一點，依照你自己的背景，給一個下一步的方向。你
不用完全對號入座，這幾種只是常見的起點，挑一個離你目前狀況最近的參考就好。

\begin{itemize}
    \item \textbf{國中生或高中生：} 現在最值得做的事情，不是急著衝去學更深的東西，是
    把這本書已經學過的東西多用幾次，變成真正內化的手感。試著幫自己寫一個小工具，整理
    作業進度、算一下成績、幫喜歡的遊戲寫個小腳本，題目越貼近你自己的生活越好。如果
    學校或住的城市附近有程式相關的社團或營隊，去看看也無妨，重點是持續動手，比急著
    往前衝更重要。
    \item \textbf{文組背景的學生，或原本從來沒接觸過理工訓練的人：} 這本書想讓你確認
    的一件事是：你不需要轉去讀資訊工程，也能真的把 Python 用起來。可以把它當成原本
    專業裡多出來的一項工具，用來整理研究資料、自動化一些重複的文書工作，或是單純
    因為好奇多一項技能。下一步不用是「換一個更難的領域」，可以是「把 Python 用回你
    原本在乎的那件事情上」。
    \item \textbf{已經在職場工作，尤其是 PM、營運、行銷這類常常要跟工程師對接、但
    自己不寫程式的角色：} \ref{ch:industry} 章跟 \ref{ch:api} 章大概是這本書對你最直接有用的兩章。下一步可以
    試著在下次開會前，先把想問工程師的問題用這本書教過的詞彙組織一遍，你會發現溝通
    明顯順暢很多。如果手邊剛好有重複性高的報表整理或資料彙整工作，也很值得試著自己
    寫一支小程式接手，不需要變成工程師，只要能看懂、能跟工程師討論、偶爾自己動手
    解決一點小麻煩，這本書設定的目標就已經達成了。
\end{itemize}

不管你屬於哪一種背景，這幾條路徑最後其實會匯合到同一個地方：持續用、持續碰、持續在
真實情境裡解決一件具體的小問題，比一口氣把地圖上剩下的每個區域都學完更重要。

這本書開頭第 \ref{ch:introduction} 章，你剛寫完人生第一支跑得動的程式，那是第一章的尾聲，而且已經有一支
自己跑得動的程式。

現在，你在整本書的尾聲，你手上有的，不只是一些程式碼片段，是兩種
看世界的方式：一種讓你看到生活裡的瑣事，會反射性想「這個我可以寫程式做」；另一種讓
你看懂「軟體」這個世界平常怎麼運作，不再被術語唬住。這兩件事，從第一頁開始就沒有誰
比較重要，現在，它們都真的是你的了。


% 書末參考資料（不編章號，跟索引一樣屬於 back matter）
\chapter*{參考資料}
\label{ch:references}
\addcontentsline{toc}{chapter}{參考資料}
\markboth{參考資料}{參考資料}

這本書不是憑空寫出來的。 決定每一章要教什麼、用什麼順序教、哪裡該慢下來的時候，我們
橫向比較了七本已經存在的 Python 入門書——不是照抄任何一本，而是看它們在同一個主題上
各自怎麼處理，再決定這本書要走哪一條路。

這一章把那七本書列出來，另外附上寫作時實際查證過的官方文件與線上服務。 它跟第
\ref{sec:learning-resources} 節不一樣：那一節談的是「有哪幾類學習資源、各自適合什麼樣的
學習節奏」，刻意不點名任何一本書，因為那些分類不會過期；這一章反過來，給的是具體的書名
與網址，並且老實告訴你每一本適合你在什麼時候拿起來。

\section*{七本參考書}
\addcontentsline{toc}{section}{七本參考書}

下面按照「你讀完這本書之後，大概會照什麼順序用得上」排列，不是按重要性。 每一本後面
標的「本書第 X 章參考」，是這本書寫作時真的拿它當素材的章節。 這七本都改版過好幾次，
挑你買得到的最新版就好。

\subsection*{想繼續動手做小工具}

\paragraph{Automate the Boring Stuff with Python}（Al Sweigart 著，No Starch Press）
——本書第 \ref{ch:introduction}、\ref{ch:variables}、\ref{ch:strings}、\ref{ch:lists}、\ref{ch:dict-set}、\ref{ch:functions}、\ref{ch:files-exceptions}、\ref{ch:csv}、\ref{ch:debugging}、\ref{ch:modules}、\ref{ch:regex}、\ref{ch:automating}、\ref{ch:capstone} 章參考。

跟這本書的「目標一」距離最近的一本：它整本都在回答「這件煩人的瑣事，能不能叫電腦做」。
檔案整理、抓網頁、寄信、操作試算表，一章一個實用主題。 如果你讀完這本書之後最想做的是
「把我每週都要重複做的那件事自動化」，這本是最直接的下一步。 作者把全文免費放在
\url{https://automatetheboringstuff.com}，可以先讀再決定要不要買紙本。

\subsection*{想把語言基礎再打紮實一點}

\paragraph{Python Crash Course}（Eric Matthes 著，No Starch Press）
——本書第 \ref{ch:introduction}、\ref{ch:variables}、\ref{ch:conditionals}、\ref{ch:loops}、\ref{ch:lists}、\ref{ch:dict-set}、\ref{ch:functions}、\ref{ch:files-exceptions}、\ref{ch:debugging}、\ref{ch:pythonic}、\ref{ch:objects}、\ref{ch:capstone} 章參考。

前半是語法基礎，後半是三個從頭做到尾的專案。 它最值得學的地方是「怎麼把一個專案一步步
長出來」——幾乎每個小節結束都會叫你立刻執行看看，讓你一路上都有能跑的東西，這本書第
\ref{ch:capstone} 章的節奏就是參考它的做法。 如果你覺得自己語法都懂了、但一要動手做完整
的東西就卡住，問題通常不在語法，這本書治的正是那個症狀。

\paragraph{Python for Everybody}（Charles Severance 著）
——本書第 \ref{ch:introduction}、\ref{ch:variables}、\ref{ch:strings}、\ref{ch:conditionals}、\ref{ch:loops}、\ref{ch:lists}、\ref{ch:dict-set}、\ref{ch:functions}、\ref{ch:tuples}、\ref{ch:files-exceptions}、\ref{ch:objects}、\ref{ch:regex} 章參考。

步調最慢、預設你知道得最少的一本，特別適合當「同一個主題換一個人再講一次」的備援。
它完全免費，連配套的線上課程與練習都是，網址是 \url{https://www.py4e.com}。 如果這本書
某一章你讀了兩遍還是覺得卡，去看它同一個主題怎麼講，常常就通了。

\subsection*{想弄清楚「為什麼是這樣」}

\paragraph{Think Python}（Allen B. Downey 著，O'Reilly）
——本書第 \ref{ch:variables}、\ref{ch:strings}、\ref{ch:conditionals}、\ref{ch:loops}、\ref{ch:lists}、\ref{ch:dict-set}、\ref{ch:functions}、\ref{ch:tuples}、\ref{ch:files-exceptions}、\ref{ch:modules}、\ref{ch:comprehensions}、\ref{ch:iterators-generators}、\ref{ch:recursion}、\ref{ch:pythonic}、\ref{ch:objects}、\ref{ch:inheritance}、\ref{ch:regex}、\ref{ch:where-next} 章參考，是
七本裡被引用最多的一本。

它關心的不只是「怎麼寫」，還有「怎麼想」。 這本書第 \ref{ch:recursion} 章用縮排追蹤
呼叫堆疊的那個手法，就是從它借來的。 如果你想從「會照著寫」進到「能自己推理出程式為什麼
這樣跑」，這本是最好的橋。 作者同樣提供免費線上版：
\url{https://allendowney.github.io/ThinkPython/}。

\subsection*{想換一種完全不同的講法}

\paragraph{Head First Python}（Paul Barry 著，O'Reilly）
——本書第 \ref{ch:introduction}、\ref{ch:csv}、\ref{ch:comprehensions}、\ref{ch:automating}、\ref{ch:capstone} 章參考。

版面充滿圖、手寫字、對話與練習，跟一般教科書長得完全不一樣。 有人非常吃這一套，有人
完全不吃。 如果你發現自己讀純文字的技術書很容易分心，值得去書店翻個五分鐘，你會很快
知道自己是哪一種人。

\paragraph{Introducing Python}（Bill Lubanovic 著，O'Reilly）
——本書第 \ref{ch:variables}、\ref{ch:strings}、\ref{ch:conditionals}、\ref{ch:loops}、\ref{ch:lists}、\ref{ch:dict-set}、\ref{ch:functions}、\ref{ch:tuples}、\ref{ch:files-exceptions}、\ref{ch:advanced-functions}、\ref{ch:comprehensions}、\ref{ch:iterators-generators}、\ref{ch:recursion}、\ref{ch:pythonic}、\ref{ch:objects}、\ref{ch:where-next} 章參考。

份量厚、覆蓋面廣，從基礎語法一路走到網路、資料庫、並行處理，比較像「一本可以一直放在
手邊翻」的書。 這本書第 \ref{sec:mutable-default} 節那個可變預設參數的坑，七本參考書裡
只有它完整教過——這種「別人都略過、它認真講」的地方，在它書裡還有不少。

\subsection*{等你準備好了再去}

\paragraph{Fluent Python}（Luciano Ramalho 著，O'Reilly）
——本書第 \ref{ch:advanced-functions}、\ref{ch:comprehensions}、\ref{ch:iterators-generators}、\ref{ch:pythonic}、\ref{ch:inheritance}、\ref{ch:where-next} 章參考。

第 \ref{sec:beyond-this-book} 節提過它。 這是一本寫給\textbf{已經在寫 Python 的工程師}看的
書，作者在序言裡就自己提醒過：剛開始學的人讀它會很難跟上，甚至可能誤以為每支 Python
程式都得用上一堆特殊技巧。 這本書第 \ref{ch:advanced-functions} 到 \ref{ch:pythonic} 章
教的 \texttt{lambda}、closure、decorator、generator、comprehension，在它書裡都有深得多
的版本。 不建議現在就讀，但值得知道它在那裡——你哪一天覺得這本書已經餵不飽你了，它就是
下一站。

\section*{官方文件與規範}
\addcontentsline{toc}{section}{官方文件與規範}

\paragraph{Python 官方文件}\url{https://docs.python.org/3/}

第 \ref{ch:modules} 章教過怎麼讀它。 你以後會用到的標準函式庫模組，第一手說明都在這裡，
而且它是唯一一份一定跟你手上這版 Python 對得起來的資料——網路上的教學文章可能是三年前
寫的，官方文件不會。

\paragraph{PEP 8：Python 程式碼風格指南}\url{https://peps.python.org/pep-0008/}

第 \ref{sec:pep8} 節介紹過的那份文件。 它比你想像中好讀，而且會親口告訴你「什麼時候
該不遵守它」。

\paragraph{PEP 20：The Zen of Python}\url{https://peps.python.org/pep-0020/}

十九句話，講 Python 這個語言在意什麼。 在你的 Python 裡輸入 \texttt{import this} 也會
把它整份印出來。 現在讀跟一年後再讀，感受會很不一樣。

\section*{本書用到的公開服務}
\addcontentsline{toc}{section}{本書用到的公開服務}

第 \ref{ch:api} 章與第 \ref{ch:capstone} 章串接的三個 API，都不需要註冊、不需要金鑰，
你可以直接照著書裡的範例打打看。 寫作時每一支都實際呼叫過，但線上服務會改版、會停用，
如果哪天書裡的範例打不通，先去它的官網看看規格有沒有變。

\begin{itemize}
    \item \textbf{Open-Meteo}（天氣）\quad \url{https://open-meteo.com}\\
    第 \ref{ch:api} 章查天氣用的服務。
    \item \textbf{Frankfurter}（匯率）\quad \url{https://frankfurter.dev}\\
    第 \ref{ch:api} 章查匯率用的服務，資料來自歐洲央行的參考匯率。 要注意它只涵蓋
    三十種左右的幣別，\textbf{裡面沒有新台幣}——第 \ref{sec:currency-api} 節的練習會讓你
    親自撞見這件事。
    \item \textbf{ExchangeRate-API（開放版）}\quad \url{https://open.er-api.com/v6/latest}\\
    第 \ref{ch:capstone} 章的記帳本改用的服務，因為它查得到新台幣。 它跟 Frankfurter
    的請求與回應形狀幾乎一樣，這正是第 \ref{ch:api} 章說過的「換一個服務，套用同一套
    骨架，換掉細節就好」。
\end{itemize}


\clearpage
\phantomsection
\addcontentsline{toc}{chapter}{索引}
\printindex

\end{document}
